% !TEX root = ../Elements.tex
%%%%%%
% BOOK 10
%%%%%%
\pdfbookmark[0]{Book 10}{book10}
\pagestyle{empty}
\begin{center}
{\Huge ELEMENTS BOOK 10}\\
\spa\spa\spa
{\huge\it Incommensurable Magnitudes\symbolfootnote[2]{The theory of incommensurable magntidues set out in this
book is generally attributed to Theaetetus of Athens. In the footnotes
throughout this book, $k$, $k'$, {\rm etc.}\ stand for distinct ratios
of positive integers.}}
\end{center}
\vspace{3cm}
\epsfysize=2.5in
\centerline{\epsffile{Book10/front.eps}}
\newpage

%%%%%%%
% Definitions
%%%%%%%
\pdfbookmark[1]{Definitions I}{def10.1}
\pagestyle{fancy}
\cfoot{{\greekfont τ̀ηεπαγε}}
\lhead{\large{\greekfont ΣΤΟΙΘΕΙΩΝ γ̀γν{10}.}}
\rhead{\large ELEMENTS BOOK 10}
\begin{Parallel}{}{} 
\ParallelLText{
\begin{center}
\large{{\greekfont Ὅροι}.}
\end{center}\vspace*{-7pt}

\ggn{1}.~{\greekfont Σύμμετρα μεγέθη λέγεται τὰ τῷ α ὐτῷ μετρῳ
μετρούμενα, ἀσύμμετρα δέ, ὧν μ ηδὲν ἐνδέθεται κοινὸν
μέτρον γενέσθαι.}

\ggn{2}.~{\greekfont Εὐθεῖαι δυνάμει σύμμετροί εἰσιν, ὅταν τὰ ἀπ
α ὐτῶν τετράγωνα τῷ α ὐτῷ θωρίῳ μετρῆται, ἀσύμμετροι
δέ, ὅταν τοῖς ἀπ α ὐτῶν τετραγώνοις μ ηδὲν ἐνδέθηται θωρίον
κοινὸν μέτρον γενέσθαι.}

\ggn{3}.~{\greekfont Τούτων ὑποκειμένων δείκνυται, ὅτι τῇ προτεθείσῃ
εὐθείᾳ ὑπάρθουσιν εὐθεῖαι πλήθει ἄπειροι σύμμετροί
τε καὶ ἀσύμμετροι αἱ μὲν μήκει μόνον, αἱ δὲ καὶ δυνάμει.
καλείσθω οὖν ἡ μὲν προτεθεῖσα εὐθεῖα ῥητή, καὶ
αἱ ταύτῃ σύμμετροι εἴτε μήκει καὶ δυνάμει εἴτε
δυνάμει μόνον ῥηταί, αἱ δὲ ταύτῃ ἀσύμμετροι ἄλογοι καλείσθωσαν.}

\ggn{4}.~{\greekfont Καὶ τὸ μὲν ἀπὸ τῆς προτεθείσης εὐθείας τετράγω-νον ῥητόν,
καὶ τὰ τούτῳ σύμμετρα ῥητά, τὰ δὲ τούτῳ ἀσύμμετρα ἄλογα
καλείσθω, καὶ αἱ δυνάμεναι α ὐτὰ ἄλογοι, εἰ μὲν τετράγωνα
εἴη, α ὐταὶ αἱ πλευραί, εἰ δὲ ἕτερά τινα εὐθύγραμμα, αἱ
ἴσα α ὐτοῖς τετράγωνα ἀναγράφουσαι.}}

\ParallelRText{
\begin{center}
{\large Definitions}
\end{center}

1.~Those magnitudes   measured by the
same measure are said (to be) commensurable, but (those)
of which no (magnitude) admits to be  a common measure (are said to be)
incommensurable.$^\dag$

2.~(Two) straight-lines are commensurable in square$^\ddag$ when the squares on
them are measured by the same area, but (are) incommensurable (in square)
when no area admits to be a common measure of the squares on them.$^\S$

3.~These things being assumed, it is proved that there
exist an infinite multitude of straight-lines commensurable and
incommensurable with an assigned straight-line---those  (incommensurable) in length only, and
those also (commensurable or incommensurable) in square.$^\P$ 
Therefore, let
the assigned straight-line be called rational. And (let) the (straight-lines) commensurable
with it, either in length and square, or in square only, (also be called) rational. But let the (straight-lines) incommensurable with it be called
irrational.$^\ast$ 

4.~And let the square on the assigned straight-line be called  rational. And (let areas) commensurable with it (also be called) rational. But  (let areas) incommensurable with it (be called) irrational, and (let) their square-roots$^\$$ (also be called) irrational---the sides themselves, if the (areas)
are squares, and the (straight-lines) describing squares equal to them, if the (areas) are some other rectilinear (figure).$^\flat$}
\end{Parallel}

{\footnotesize\noindent $^\dag$ In other words, two magnitudes $\alpha$ and
$\beta$ are commensurable if $\alpha:\beta::1:k$, and incommensurable otherwise.\\[0.5ex]
$^\ddag$ Literally, ``in power''.\\[0.5ex]
$^\S$  In other words, two straight-lines of length $\alpha$ and $\beta$ are commensurable in square if $\alpha : \beta::1:k^{1/2}$,  and incommensurable in square otherwise. Likewise, the straight-lines
are commensurable in length if $\alpha:\beta::1:k$, and incommensurable in length otherwise.\\[0.5ex]
$^\P$ To be more exact, straight-lines can either be commensurable in square only, incommensurable in length only, or commenusrable/incommensurable in both
length and square, with an assigned straight-line.\\[0.5ex]
$^\ast$ Let the length of the assigned straight-line be unity. Then rational
straight-lines have lengths expressible as $k$ or $k^{1/2}$, depending on whether the lengths are commensurable in length, or in square only, respectively,
with unity.  All other straight-lines are irrational.\\[0.5ex]
$^\$$ The square-root of an area is the length of the side of an equal area square.\\[0.5ex]
$^\flat$ The area of the  square 
on the assigned straight-line is unity. Rational areas are expressible as $k$.  All other areas are irrational. Thus, squares whose
sides are of rational length have  rational areas, and {\rm vice versa}.}

%%%%%%
% Prop 10.1
%%%%%%
\pdfbookmark[1]{Proposition 10.1}{pdf10.1}
\begin{Parallel}{}{} 
\ParallelLText{
\begin{center}
{\large \ggn{1}.}
\end{center}\vspace*{-7pt}

{\greekfont Δύο μεγεθῶν ἀνίσων ἐκκειμένων, ἐὰν ἀπὸ τοῦ μείζονος ἀφαιρεθῇ μεῖζον ἢ τὸ ἥμισυ καὶ τοῦ καταλειπομένου μεῖζον ἢ
τὸ ἥμισυ, καὶ τοῦτο ἀεὶ γίγνηται, λειφθήσεταί τι μέγεθος, ὃ ἔσται
ἔλασσον τοῦ ἐκκειμένου ἐλάσσονος μεγέθους.}

{\greekfont Ἔστω δύο μεγέθη ἄνισα τὰ ΑΒ, Γ, ὧν μεῖζον τὸ ΑΒ; λέγω, ὅτι,
ἐαν ἀπὸ τοῦ ΑΒ ἀφαιρεθῇ μεῖζον ἢ τὸ ἥμισυ
καὶ τοῦ καταλειπομένου μεῖζον ἢ τὸ ἥμισυ, καὶ τοῦτο ἀεὶ
γίγνηται, λειφθήσεταί τι μέγεθος, ὃ ἔσται ἔλασσον τοῦ Γ μεγέθους.}\\~\\~\\~\\

\epsfysize=1.1in
\centerline{\epsffile{Book10/fig001g.eps}}

{\greekfont Τὸ Γ γὰρ πολλαπλασιαζόμενον ἔσται ποτὲ τοῦ ΑΒ μεῖζον. πεπολλαπλασιάσθω, καὶ ἔστω τὸ ΔΕ τοῦ μὲν Γ πολλαπλάσιον, τοῦ δὲ ΑΒ
μεῖζον, καὶ διῃρήσθω τὸ ΔΕ εἰς τὰ τῷ Γ ἴσα τὰ ΔΖ, ΖΗ, ΗΕ,
καὶ ἀφῃρήσθω ἀπὸ μὲν τοῦ ΑΒ μεῖζον ἢ τὸ ἥμισυ τὸ ΒΘ,
ἀπὸ δὲ τοῦ ΑΘ μεῖζον ἢ τὸ ἥμισυ τὸ ΘΚ, καὶ τοῦτο
ἀεὶ γιγνέσθω, ἕως ἂν αἱ ἐν τῷ ΑΒ διαιρέσεις ἰσοπληθεῖς
γένωνται ταῖς ἐν τῷ ΔΕ διαιρέσεσιν.}

{\greekfont Ἔστωσαν οὖν αἱ ΑΚ, ΚΘ, ΘΒ διαιρέσεις ἰσοπληθεῖς οὖσαι ταῖς
ΔΖ, ΖΗ, ΗΕ; καὶ ἐπεὶ μεῖζόν ἐστι τὸ ΔΕ τοῦ ΑΒ, καὶ ἀφῄρηται
ἀπὸ μὲν τοῦ ΔΕ ἔλασσον τοῦ ἡμίσεως τὸ ΕΗ, ἀπὸ δὲ τοῦ ΑΒ
μεῖζον ἢ τὸ ἥμισυ τὸ ΒΘ, λοιπὸν ἄρα τὸ ΗΔ λοιποῦ τοῦ ΘΑ
μεῖζόν ἐστιν. καὶ ἐπεὶ μεῖζόν ἐστι τὸ ΗΔ τοῦ ΘΑ, καὶ
ἀφῄρηται τοῦ μὲν ΗΔ ἥμισυ τὸ ΗΖ, τοῦ δὲ ΘΑ μεῖζον ἢ
τὸ ἥμισυ τὸ ΘΚ, λοιπὸν ἄρα τὸ ΔΖ λοιποῦ τοῦ ΑΚ μεῖζόν ἐστιν.
ἴσον δὲ τὸ ΔΖ τῷ Γ; καὶ τὸ Γ ἄρα τοῦ ΑΚ μεῖζόν ἐστιν.
ἔλασσον ἄρα τὸ ΑΚ τοῦ Γ.}

{\greekfont Καταλείπεται ἄρα ἀπὸ τοῦ ΑΒ μεγέθους τὸ ΑΚ μέγεθος ἔλασσον ὂν
τοῦ ἐκκειμένου ἐλάσσονος μεγέθους τοῦ Γ; ὅπερ ἔδει δεῖχαι.}\,---\,{\greekfont ὁμοίως δὲ δειθθήσεται, κἂν ἡμίση ᾖ τὰ ἀφαιρούμενα.}}

\ParallelRText{
\begin{center}
{\large Proposition 1}$^\dag$
\end{center}

If, from
the greater of two unequal magnitudes (which are) laid out,  (a part) greater than half is subtracted, and (if from) the remainder (a part) greater than
half (is subtracted), and (if) this happens continually, then some magnitude will (eventually) be left which
will be less than the  lesser laid out magnitude.

Let $AB$ and $C$ be two unequal magnitudes, of which (let) $AB$ (be) the greater. I say that if (a part) greater than half is subtracted  from $AB$,
and (if a part) greater than half (is subtracted) from the remainder, and (if) this happens
continually, then some magnitude will (eventually) be left which will be less
than the magnitude $C$. 

\epsfysize=1.1in
\centerline{\epsffile{Book10/fig001e.eps}}

For $C$, when multiplied (by some number), will sometimes be greater than $AB$ [Def.~5.4]. Let it 
 be (so) multiplied. And let $DE$ be (both) a multiple of $C$, and greater
than $AB$. And let $DE$ be divided into the (divisions) $DF$, $FG$,
$GE$, equal to $C$. And   let $BH$, (which is) greater than half, be subtracted from $AB$. And  (let) $HK$, (which is) greater than half, (be subtracted) from $AH$. And let this happen continually, until the divisions in
$AB$ become equal in number to the divisions in $DE$.

Therefore, let the divisions (in $AB$) be $AK$, $KH$, $HB$, being equal in number to $DF$, $FG$, $GE$.  And since $DE$ is greater than $AB$, and $EG$, (which is) less than half, has been subtracted from $DE$, and $BH$, (which is) greater
than half, from $AB$, the remainder $GD$ is thus greater than the remainder
$HA$. And since $GD$ is greater than $HA$, and the half $GF$ has been
subtracted from $GD$, and $HK$, (which is) greater than half, from $HA$,
the remainder $DF$ is thus greater than the remainder $AK$. And $DF$
(is) equal to $C$. $C$ is thus also greater than $AK$.
Thus, $AK$ (is) less than $C$.

Thus, the magnitude $AK$,  which is less
than the lesser laid out magnitude $C$, is left over from the magnitude $AB$. (Which is) the very thing it was required to show. --- (The theorem) can similarly be proved even if  the
(parts) subtracted are halves.}
\end{Parallel}
{\footnotesize\noindent $^\dag$ This theorem is the basis of the so-called {\rm method of exhaustion}, and is generally attributed to Eudoxus of Cnidus.}

%%%%%%
% Prop 10.2
%%%%%%
\pdfbookmark[1]{Proposition 10.2}{pdf10.2}
\begin{Parallel}{}{} 
\ParallelLText{
\begin{center}
{\large \ggn{2}.}
\end{center}\vspace*{-7pt}

{\greekfont Ἐὰν δύο μεγεθῶν [ἐκκειμένων] ἀνίσων ἀνθυφαιρουμένου ἀεὶ
τοῦ ἐλάσσονος ἀπὸ τοῦ μείζονος τὸ καταλειπόμενον μ ηδέποτε
καταμετρῇ τὸ πρὸ ἑαυτοῦ, ἀσύμμετρα ἔσται τὰ μεγέθη.}\\

\epsfysize=0.9in
\centerline{\epsffile{Book10/fig002g.eps}}

{\greekfont Δύο γὰρ μεγεθῶν ὄντων ἀνίσων τῶν ΑΒ, ΓΔ καὶ ἐλάσσονος
τοῦ ΑΒ ἀνθυφαιρουμένου ἀεὶ τοῦ ἐλάσσονος ἀπὸ τοῦ μείζονος τὸ
περιλειπόμενον μ ηδέποτε καταμετρείτω τὸ πρὸ ἑαυτοῦ; λέγω, ὅτι
ἀσύμμετρά ἐστι τὰ ΑΒ, ΓΔ μεγέθη.}

{\greekfont Εἰ γάρ ἐστι σύμμετρα, μετρήσει τι α ὐτὰ μέγεθος. μετρείτω, εἰ δυνατόν,
καὶ ἔστω τὸ Ε; καὶ τὸ μὲν ΑΒ τὸ ΖΔ καταμετροῦν λειπέτω ἑαυτοῦ
ἔλασσον τὸ ΓΖ, τὸ δὲ ΓΖ τὸ ΒΗ καταμετροῦν λειπέτω ἑαυτοῦ
ἔλασσον τὸ ΑΗ, καὶ τοῦτο ἀεὶ γινέσθω, ἕως οὗ λειφθῇ τι
μέγεθος, ὅ ἐστιν ἔλασσον τοῦ Ε. γεγονέτω, καὶ λελείφθω τὸ ΑΗ
ἔλασσον τοῦ Ε. ἐπεὶ οὖν τὸ Ε τὸ ΑΒ μετρεῖ, ἀλλὰ τὸ ΑΒ τὸ
ΔΖ μετρεῖ, καὶ τὸ Ε ἄρα τὸ ΖΔ μετρήσει. μετρεῖ δὲ καὶ ὅλον
τὸ ΓΔ; καὶ λοιπὸν ἄρα τὸ ΓΖ μετρήσει. ἀλλὰ τὸ ΓΖ τὸ ΒΗ μετρεῖ;
καὶ τὸ Ε ἄρα τὸ ΒΗ μετρεῖ. μετρεῖ δὲ καὶ ὅλον τὸ ΑΒ; καὶ λοιπὸν
ἄρα τὸ ΑΗ μετρήσει, τὸ μεῖζον τὸ ἔλασσον; ὅπερ ἐστὶν ἀδύνατον.
οὐκ ἄρα τὰ ΑΒ, ΓΔ μεγέθη μετρήσει τι μέγεθος; ἀσύμμετρα ἄρα
ἐστὶ τὰ ΑΒ, ΓΔ μεγέθη.}

{\greekfont Ἐὰν ἄρα δύο μεγεθῶν ἀνίσων, καὶ τὰ ἑχῆς.}}

\ParallelRText{
\begin{center}
{\large Proposition 2}
\end{center}

If the remainder of two unequal magnitudes (which are) [laid out] never measures the (magnitude) before it, (when) the lesser (magnitude is) continually subtracted in turn from the greater, then the (original) magnitudes will be incommensurable.

\epsfysize=0.9in
\centerline{\epsffile{Book10/fig002e.eps}}

For, $AB$ and $CD$ being two unequal magnitudes, and $AB$ (being) the lesser, let the remainder
never measure the (magnitude) before it,
(when) the lesser (magnitude is) continually subtracted in turn from the greater. I say that
the magnitudes $AB$ and $CD$ are incommensurable.

For if they are commensurable then some magnitude will measure them (both). If possible, let it
(so) measure (them),  and let it be $E$. And let $AB$ leave $CF$
less than itself (in) measuring $FD$, and let $CF$ leave $AG$ less than
itself (in) measuring $BG$, and let this happen continually, until  some magnitude which is less than  $E$ is left. Let (this) have occurred,$^\dag$ and
let $AG$, (which is) less than $E$, be left. Therefore, since
$E$ measures $AB$, but $AB$ measures $DF$, $E$ will thus also measure $FD$. And it also measures the whole (of) $CD$. Thus, it will
also measure the remainder  $CF$. But, $CF$ measures $BG$. Thus,
$E$ also measures $BG$. And it also measures the whole (of) $AB$. 
Thus, it will also measure the remainder $AG$, the greater (measuring)
the lesser. The very thing is impossible. Thus, some magnitude cannot
measure (both) the magnitudes $AB$ and $CD$. Thus, the
magnitudes $AB$ and $CD$ are incommensurable [Def.~10.1].

Thus, if \ldots of two unequal magnitudes, and so on \ldots.}
\end{Parallel}
{\footnotesize\noindent$^\dag$ The fact that this will eventually occur is guaranteed by
Prop.~10.1.}

%%%%%%
% Prop 10.3
%%%%%%
\pdfbookmark[1]{Proposition 10.3}{pdf10.3}
\begin{Parallel}{}{} 
\ParallelLText{
\begin{center}
{\large \ggn{3}.}
\end{center}\vspace*{-7pt}

{\greekfont Δύο μεγεθῶν συμμέτρων δοθέντων τὸ μέγιστον α ὐτῶν κοινὸν
μέτρον εὑρεῖν.}

\epsfysize=1.1in
\centerline{\epsffile{Book10/fig003g.eps}}

{\greekfont Ἔστω τὰ δοθέντα δύο μεγέθη σύμμετρα τὰ ΑΒ, ΓΔ, ὧν
ἔλασσον τὸ ΑΒ; δεῖ δὴ τῶν ΑΒ, ΓΔ τὸ μέγιστον κοινὸν μέτρον
εὑρεῖν.}

{\greekfont Τὸ ΑΒ γὰρ μέγεθος ἤτοι μετρεῖ τὸ ΓΔ ἢ οὔ. εἰ
μὲν οὖν μετρεῖ, μετρεῖ δὲ καὶ ἑαυτό, τὸ ΑΒ ἄρα τῶν
ΑΒ, ΓΔ κοινὸν μέτρον ἐστίν; καὶ φανερόν, ὅτι καὶ
μέγιστον. μεῖζον γὰρ τοῦ ΑΒ μεγέθους τὸ ΑΒ οὐ
μετρήσει.}

{\greekfont μὴ μετρείτω δὴ τὸ ΑΒ τὸ ΓΔ. καὶ ἀνθυφαιρουμένου ἀεὶ τοῦ ἐλάσσονος ἀπὸ τοῦ  μείζονος, τὸ περιλειπόμενον μετρήσει
ποτὲ τὸ πρὸ ἑαυτοῦ διὰ τὸ μὴ εἶναι ἀσύμμετρα τὰ
ΑΒ, ΓΔ; καὶ τὸ μὲν ΑΒ τὸ ΕΔ καταμετροῦν λειπέτω ἑαυτοῦ
ἔλασσον τὸ ΕΓ, τὸ δὲ ΕΓ τὸ ΖΒ καταμετροῦν λειπέτω ἑαυτοῦ
ἔλασσον τὸ ΑΖ, τὸ δὲ ΑΖ τὸ ΓΕ μετρείτω.}

{\greekfont Ἐπεὶ οὖν τὸ ΑΖ τὸ ΓΕ μετρεῖ, ἀλλὰ τὸ ΓΕ τὸ ΖΒ μετρεῖ, καὶ
τὸ ΑΖ ἄρα τὸ ΖΒ μετρήσει. μετρεῖ δὲ καὶ ἑαυτό; καὶ ὅλον ἄρα
τὸ ΑΒ μετρήσει τὸ ΑΖ. ἀλλὰ τὸ ΑΒ τὸ ΔΕ μετρεῖ; καὶ τὸ ΑΖ ἄρα
τὸ ΕΔ μετρήσει. μετρεῖ δὲ καὶ τὸ ΓΕ; καί ὅλον ἄρα τὸ ΓΔ
μετρεῖ; τὸ ΑΖ ἄρα τῶν ΑΒ, ΓΔ κοινὸν μέτρον ἐστίν. λέγω δή,
ὅτι καὶ μέγιστον. εἰ γὰρ μή, ἔσται τι μέγεθος μεῖζον τοῦ ΑΖ,
ὃ
μετρήσει τὰ ΑΒ, ΓΔ. ἔστω τὸ Η. ἐπεὶ οὖν τὸ Η τὸ ΑΒ
μετρεῖ, ἀλλὰ τὸ ΑΒ τὸ ΕΔ μετρεῖ, καὶ τὸ Η ἄρα τὸ ΕΔ μετρήσει. μετρεῖ δὲ καὶ ὅλον τὸ ΓΔ; καὶ λοιπὸν ἄρα τὸ ΓΕ μετρήσει
τὸ Η. ἀλλὰ τὸ ΓΕ τὸ ΖΒ μετρεῖ; καὶ τὸ Η ἄρα
τὸ ΖΒ μετρήσει. μετρεῖ δὲ καὶ ὅλον τὸ ΑΒ, καὶ λοιπὸν τὸ ΑΖ
μετρήσει, τὸ μεῖζον τὸ ἔλασσον; ὅπερ ἐστὶν ἀδύνατον. οὐκ
ἄρα μεῖζόν τι μέγεθος τοῦ ΑΖ τὰ ΑΒ, ΓΔ μετρήσει; τὸ ΑΖ ἄρα
τῶν ΑΒ, ΓΔ τὸ μέγιστον κοινὸν μέτρον ἐστίν.}

{\greekfont Δύο ἄρα μεγεθῶν συμμέτρων δοθέντων τῶν ΑΒ, ΓΔ τὸ μέγιστον
κοινὸν μέτρον ηὕρηται; ὅπερ
ἔδει δεῖχαι.}\\~\\~\\~\\~\\~\\~\\

\begin{center}
{\large {\greekfont Πόρισμα}.}
\end{center}\vspace*{-7pt}

{\greekfont Ἐκ δὴ τούτου φανερόν, ὅτι, ἐὰν μέγεθος δύο μεγέθη μετρῇ, καὶ
τὸ μέγιστον α ὐτῶν κοινὸν μέτρον μετρήσει.}}

\ParallelRText{
\begin{center}
{\large Proposition 3}
\end{center}

To find the greatest common measure of two
given commensurable magnitudes.

\epsfysize=1.1in
\centerline{\epsffile{Book10/fig003e.eps}}

Let $AB$ and $CD$ be the two given magnitudes, of which (let) $AB$ (be)
the lesser. So, it is required to find the greatest common measure of
$AB$ and $CD$.

For the magnitude $AB$ either measures, or (does) not (measure), $CD$. Therefore, if
it measures ($CD$), and (since) it also measures itself, $AB$ is thus
a common measure of $AB$ and $CD$. And (it is) clear that (it is) also
(the) greatest. For a (magnitude) greater than magnitude $AB$ cannot
measure $AB$.

So let $AB$ not measure $CD$. And continually subtracting in turn the lesser 
(magnitude) from the greater, the remaining (magnitude) will (at) some time measure the (magnitude)
before it, on account of $AB$ and $CD$ not being incommensurable [Prop.~10.2]. And let $AB$ leave $EC$ less than
itself (in) measuring $ED$, and let $EC$ leave $AF$ less than itself (in)
measuring $FB$, and let $AF$ measure $CE$.

Therefore, since $AF$ measures $CE$, but $CE$ measures $FB$, $AF$
will thus also measure $FB$. And it also measures itself. Thus, $AF$
will also measure the whole (of) $AB$. But, $AB$ measures $DE$. Thus, $AF$
will also measure $ED$. And it also measures $CE$. Thus, it also measures the whole of $CD$. Thus, $AF$ is a common measure of $AB$ and $CD$.
So I say that (it is) also (the) greatest (common measure). For, if not,
there will be some magnitude, greater than $AF$, which will measure
(both) $AB$ and $CD$. Let it be $G$. Therefore, since $G$ measures $AB$,
but $AB$ measures $ED$, $G$ will thus also measure $ED$. And
it also measures the whole of $CD$. Thus, $G$ will also measure the remainder $CE$.
But $CE$ measures $FB$. Thus, $G$ will also measure $FB$. And
it also measures the whole (of) $AB$. And (so) it will measure the
remainder $AF$, the greater (measuring) the lesser. The very thing is impossible. Thus, some magnitude greater than $AF$ cannot measure
(both) $AB$ and $CD$. Thus, $AF$ is the greatest common measure of
$AB$ and $CD$.

Thus, the greatest common measure of two given commensurable magnitudes, $AB$ and $CD$, has been found. (Which is) the very thing
it was required to show.\\

\begin{center}
{\large Corollary}
\end{center}\vspace*{-7pt}

So (it is) clear, from this, that if a magnitude measures two magnitudes
then it will also measure their greatest common measure.}
\end{Parallel}

%%%%%%
% Prop 10.4
%%%%%%
\pdfbookmark[1]{Proposition 10.4}{pdf10.4}
\begin{Parallel}{}{} 
\ParallelLText{
\begin{center}
{\large \ggn{4}.}
\end{center}\vspace*{-7pt}

{\greekfont Τριῶν μεγεθῶν συμμέτρων δοθέντων τὸ μέγιστον α ὐτῶν
κοινὸν μέτρον εὑρεῖν.}

\epsfysize=1.2in
\centerline{\epsffile{Book10/fig004g.eps}}

{\greekfont Ἔστω τὰ δοθέντα τρία μεγέθη σύμμετρα τὰ Α, Β, Γ; δεῖ δὴ τῶν
Α, Β, Γ τὸ μέγιστον κοινὸν μέτρον εὑρεῖν.}

{\greekfont Εἰλήφθω γὰρ δύο τῶν Α, Β τὸ μέγιστον κοινὸν μέτρον, καὶ ἔστω
τὸ Δ; τὸ δὴ Δ τὸ Γ ἤτοι μετρεῖ ἢ οὔ [μετρεῖ]. μετρείτω πρότερον.
ἐπεὶ οὖν τὸ Δ τὸ Γ μετρεῖ, μετρεῖ δὲ καὶ τὰ Α, Β, τὸ Δ ἄρα
τὰ Α, Β, Γ μετρεῖ; τὸ Δ ἄρα τῶν Α, Β, Γ κοινὸν μέτρον ἐστίν.
καὶ φανερόν, ὅτι καὶ μέγιστον; μεῖζον γὰρ τοῦ Δ μεγέθους τὰ Α, Β
οὐ μετρεῖ.}

{\greekfont μὴ μετρείτω δὴ τὸ Δ τὸ Γ. λέγω πρῶτον, ὅτι
σύμμετρά ἐστι τὰ Γ, Δ. ἐπεὶ γὰρ σύμμετρά ἐστι τὰ Α, Β, Γ, μετρήσει τι 
α ὐτὰ μέγεθος, ὃ δηλαδὴ καὶ τὰ Α, Β μετρήσει; ὥστε καὶ τὸ τῶν
Α, Β μέγιστον κοινὸν μέτρον τὸ Δ μετρήσει. μετρεῖ δὲ καὶ τὸ Γ;
ὥστε τὸ εἰρημένον μέγεθος μετρήσει τὰ Γ, Δ; σύμμετρα ἄρα 
ἐστὶ τὰ Γ, Δ. εἰλήφθω οὖν α ὐτῶν τὸ μέγιστον κοινὸν μέτρον,
καὶ ἔστω τὸ Ε. ἐπεὶ οὖν τὸ Ε τὸ Δ μετρεῖ, ἀλλὰ τὸ Δ τὰ Α, Β
μετρεῖ, καὶ τὸ Ε ἄρα τὰ Α, Β μετρήσει. μετρεῖ δὲ καὶ τὸ Γ.
τὸ Ε ἄρα τὰ Α, Β, Γ μετρεῖ; τὸ Ε ἄρα τῶν Α, Β, Γ κοινόν ἐστι
μέτρον. λέγω δή, ὅτι καὶ μέγιστον. εἰ γὰρ δυνατόν, ἔστω τι τοῦ
Ε μεῖζον μέγεθος τὸ Ζ, καὶ μετρείτω τὰ Α, Β, Γ. καὶ ἐπεὶ τὸ
Ζ τὰ Α, Β, Γ μετρεῖ, καὶ τὰ Α, Β ἄρα μετρήσει καὶ τὸ τῶν Α, Β
μέγιστον κοινὸν μέτρον μετρήσει. τὸ δὲ τῶν Α, Β μέγιστον κοινὸν
μέτρον ἐστὶ τὸ Δ; τὸ Ζ ἄρα τὸ Δ μετρεῖ. μετρεῖ δὲ καὶ τὸ Γ; τὸ Ζ
ἄρα τὰ Γ, Δ μετρεῖ; καὶ τὸ τῶν Γ, Δ ἄρα μέγιστον κοινὸν
μέτρον μετρήσει τὸ Ζ. ἔστι δὲ τὸ Ε; τὸ Ζ ἄρα τὸ Ε μετρήσει,
τὸ μεῖζον τὸ ἔλασσον; ὅπερ ἐστὶν ἀδύνατον. οὐκ ἄρα μεῖζόν
τι τοῦ Ε μεγέθους [μέγεθος] τὰ Α, Β, Γ μετρεῖ; τὸ Ε ἄρα τῶν Α, Β, Γ τὸ μέγιστον κοινὸν μέτρον ἐστίν, ἐὰν μὴ μετρῇ τὸ Δ τὸ Γ, ἐὰν
δὲ μετρῇ, α ὐτὸ τὸ Δ.}

{\greekfont Τριῶν ἄρα μεγεθῶν συμμέτρων δοθέντων τὸ μέγιστον κοινὸν
μέτρον ηὕρηται [ὅπερ ἔδει δεῖχαι].}\\~\\~\\~\\~\\~\\~\\~\\~\\~\\

\begin{center}
{\large {\greekfont Πόρισμα}.}
\end{center}\vspace*{-7pt}

{\greekfont Ἐκ δὴ τούτου φανερόν, ὅτι, ἐὰν μέγεθος τρία μεγέθη μετρῇ, καὶ
τὸ μέγιστον α ὐτῶν κοινὸν μέτρον μετρήσει.}

{\greekfont Ὁμοίως δὴ καὶ ἐπὶ πλειόνων τὸ μέγιστον κοινὸν
μέτρον ληφθήσεται, καὶ τὸ πόρισμα προθωρήσει. ὅπερ ἔδει δεῖχαι.}}

\ParallelRText{
\begin{center}
{\large Proposition 4}
\end{center}

To find the greatest common measure of three
given commensurable magnitudes.

\epsfysize=1.2in
\centerline{\epsffile{Book10/fig004e.eps}}

Let $A$, $B$, $C$ be the three given commensurable magnitudes. So
it is required to find the greatest common measure of $A$, $B$, $C$.

For let the greatest common measure of the two (magnitudes)
$A$ and $B$ be taken [Prop.~10.3],
and let it be $D$. So $D$ either measures, or [does] not
[measure], $C$. Let it, first of all, measure ($C$). Therefore, since
$D$ measures $C$, and it also measures $A$ and $B$, $D$ thus
measures $A$, $B$, $C$. Thus, $D$ is a common measure of $A$, $B$, $C$.
And (it is) clear that (it is) also (the) greatest (common measure). For
no magnitude larger than $D$ measures (both) $A$ and $B$.

So let $D$ not measure $C$. I say, first, that $C$ and $D$ are
commensurable. For if $A$, $B$, $C$ are commensurable then some
magnitude will measure them which will clearly also measure $A$ and $B$.
Hence, it will also measure $D$, the greatest common measure of $A$ and $B$ [Prop.~10.3~corr.]. And it also measures
$C$. Hence, the aforementioned magnitude will measure (both)
$C$ and $D$. Thus, $C$ and $D$ are commensurable [Def.~10.1]. Therefore, let their
greatest common measure be taken [Prop.~10.3], and let it be $E$. Therefore, since
$E$ measures $D$, but $D$ measures (both) $A$ and $B$, $E$ will
thus also measure $A$ and $B$. And it also measures $C$. Thus,
$E$ measures $A$, $B$, $C$. Thus, $E$ is a common measure of $A$, $B$, $C$.  So I say that (it is) also (the) greatest (common measure). For, if
possible, let $F$ be some magnitude greater than $E$, and let it
measure $A$, $B$, $C$. And since $F$ measures $A$, $B$, $C$, it will
thus also measure $A$ and $B$, and  will (thus) measure the greatest common measure of $A$ and $B$ [Prop.~10.3~corr.].  And
$D$ is the greatest common measure of $A$ and $B$.  Thus, $F$ measures $D$. And it also measures $C$. Thus, $F$ measures (both) $C$ and $D$.
Thus, $F$ will also measure the greatest common measure of $C$ and $D$  [Prop.~10.3~corr.].
And it is $E$. Thus, $F$ will measure $E$, the greater (measuring)
the lesser. The very thing is impossible. Thus, some [magnitude]
greater than the magnitude $E$ cannot measure $A$, $B$, $C$. Thus,
if $D$ does not
measure $C$ then
$E$ is the greatest common measure of $A$, $B$, $C$.  And if it does measure ($C$) then $D$ itself (is the
greatest common measure).

Thus, the greatest common measure  of three given
commensurable magnitudes has been found.  [(Which is) the very thing it was required to
show.]\\

\begin{center}
{\large Corollary}
\end{center}\vspace*{-7pt}

So (it is) clear, from this, that if a magnitude measures three
magnitudes then it will also measure their greatest common measure.

So, similarly, the greatest common measure of more (magnitudes) can also be
taken, and the (above) corollary will go forward. (Which is) the very thing
it was required to show.}
\end{Parallel}

%%%%%%
% Prop 10.5
%%%%%%
\pdfbookmark[1]{Proposition 10.5}{pdf10.5}
\begin{Parallel}{}{} 
\ParallelLText{
\begin{center}
{\large \ggn{5}.}
\end{center}\vspace*{-7pt}

{\greekfont Τὰ σύμμετρα μεγέθη πρὸς ἄλληλα λόγον ἔχει, ὃν ἀριθμὸς πρὸς
ἀριθμόν.}

\epsfysize=0.8in
\centerline{\epsffile{Book10/fig005g.eps}}

{\greekfont Ἔστω σύμμετρα μεγέθη τὰ Α, Β; λέγω, ὅτι τὸ Α πρὸς τὸ Β λόγον ἔχει,
ὃν ἀριθμὸς πρὸς ἀριθμόν.}

{\greekfont Ἐπεὶ γὰρ σύμμετρά ἐστι τὰ Α, Β, μετρήσει τι α ὐτὰ μέγεθος. μετρείτω,
καὶ ἔστω τὸ Γ. καὶ ὁσάκις τὸ Γ τὸ Α μετρεῖ, τοσαῦται μονάδες
ἔστωσαν ἐν τῷ Δ, ὁσάκις δὲ τὸ Γ τὸ Β μετρεῖ, τοσαῦται μονάδες
ἔστωσαν ἐν τῷ Ε.}

{\greekfont Ἐπεὶ οὖν τὸ Γ τὸ Α μετρεῖ κατὰ τὰς ἐν τῷ Δ μονάδας, 
μετρεῖ δὲ καὶ ἡ μονὰς τὸν Δ κατὰ τὰς ἐν α ὐτῷ μονάδας,
ἰσάκις
ἄρα ἡ μονὰς τὸν Δ μετρεῖ ἀριθμὸν καὶ τὸ Γ μέγεθος τὸ Α; ἔστιν
ἄρα ὡς τὸ Γ πρὸς τὸ Α, οὕτως ἡ μονὰς πρὸς τὸν Δ; ἀνάπαλιν
ἄρα, ὡς τὸ Α πρὸς τὸ Γ, οὕτως ὁ Δ πρὸς τὴν μονάδα. πάλιν
ἐπεὶ τὸ Γ τὸ Β μετρεῖ κατὰ τὰς ἐν τῷ Ε μονάδας, μετρεῖ
δὲ καὶ ἡ μονὰς τὸν Ε κατὰ τὰς ἐν α ὐτῷ μονάδας, ἰσάκις
ἄρα ἡ μονὰς τὸν Ε μετρεῖ καὶ τὸ Γ τὸ Β; ἔστιν ἄρα ὡς τὸ Γ
πρὸς τὸ Β, οὕτως ἡ μονὰς πρὸς τὸν Ε. ἐδείθθη δὲ καὶ ὡς τὸ
Α πρὸς τὸ Γ, ὁ Δ πρὸς τὴν μονάδα; δι ἴσου ἄρα ἐστὶν ὡς τὸ
Α πρὸς τὸ Β, οὕτως ὁ Δ ἀριθμὸς πρὸς τὸν Ε.}

{\greekfont Τὰ ἄρα σύμμετρα μεγέθη τὰ Α, Β πρὸς ἄλληλα λόγον ἔχει,
ὃν ἀριθμὸς ὁ Δ πρὸς ἀριθμὸν τὸν Ε; ὅπερ ἔδει δεῖχαι.}}

\ParallelRText{
\begin{center}
{\large Proposition 5}
\end{center}

Commensurable magnitudes have to one another the
ratio which (some) number (has) to (some) number.

\epsfysize=0.8in
\centerline{\epsffile{Book10/fig005e.eps}}

Let $A$ and $B$ be commensurable magnitudes. I say that $A$ has to $B$
the ratio which (some) number (has) to (some) number.

For if $A$ and $B$ are commensurable (magnitudes) then some magnitude
will measure them. Let it (so) measure (them), and let it be $C$. And
as many times as $C$ measures $A$, so many units let there be in
$D$. And as many times as $C$ measures $B$, so many units let there be
in $E$.

Therefore, since $C$ measures $A$ according to the units in $D$, and
a unit also measures $D$ according to the units in it,  a unit thus measures the
number $D$ as many times as the magnitude $C$ (measures) $A$. 
Thus, as $C$ is to $A$, so a unit (is) to $D$ [Def.~7.20].$^\dag$ Thus, inversely, as $A$ (is) to $C$, so
$D$ (is) to a unit [Prop.~5.7~corr.]. Again, since
$C$ measures $B$ according to the units in $E$, and a unit also measures $E$ according to the units in it,   a unit thus measures $E$ the
same number of times that $C$ (measures) $B$. Thus, as $C$ is to $B$, so
a unit (is) to $E$ [Def.~7.20]. And it was also
shown that as $A$ (is) to $C$, so $D$ (is) to a unit. Thus, via equality,
as $A$ is to $B$, so the number $D$ (is) to the (number) $E$ [Prop.~5.22].

Thus, the commensurable magnitudes $A$ and $B$ have to one another
the ratio which the number $D$ (has) to the number $E$. (Which is) the
very thing it was required to show.}
\end{Parallel}
{\footnotesize\noindent$^\dag$ There is a slight logical gap here, since Def.~7.20  applies to four numbers, rather than two number and two magnitudes.}

%%%%%%
% Prop 10.6
%%%%%%
\pdfbookmark[1]{Proposition 10.6}{pdf10.6}
\begin{Parallel}{}{} 
\ParallelLText{
\begin{center}
{\large \ggn{6}.}
\end{center}\vspace*{-7pt}

{\greekfont Ἐὰν δύο μεγέθη πρὸς ἄλληλα λόγον ἔθῃ, ὃν ἀριθμὸς πρὸς
ἀριθμόν, σύμμετρα ἔσται τὰ μεγέθη.}\\

\epsfysize=0.8in
\centerline{\epsffile{Book10/fig006g.eps}}

{\greekfont Δύο γὰρ μεγέθη τὰ Α, Β πρὸς ἄλληλα λόγον ἐθέτω, ὃν ἀριθμὸς
ὁ Δ πρὸς ἀριθμὸν τὸν Ε; λέγω, ὅτι σύμμετρά ἐστι τὰ Α, Β
μεγέθη.}

{\greekfont Ὅσαι γάρ εἰσιν ἐν τῷ Δ μονάδες, εἰς τοσαῦτα ἴσα
διῃρήσθω τὸ Α, καὶ ἑνὶ α ὐτῶν ἴσον ἔστω τὸ Γ; ὅσαι
δέ εἰσιν ἐν τῷ Ε μονάδες, ἐκ τοσούτων μεγεθῶν ἴσων
τῷ Γ συγκείσθω τὸ Ζ.}

{\greekfont Ἐπεὶ οὖν, ὅσαι εἰσὶν ἐν τῷ Δ μονάδες, τοσαῦτά εἰσι
καὶ ἐν τῷ Α μεγέθη ἴσα τῷ Γ, ὃ ἄρα μέρος ἐστὶν ἡ μονὰς
τοῦ Δ, τὸ α ὐτὸ μέρος ἐστὶ καὶ τὸ Γ τοῦ Α; ἔστιν ἄρα ὡς τὸ Γ
πρὸς τὸ Α, οὕτως ἡ μονὰς πρὸς τὸν Δ. μετρεῖ δὲ ἡ
μονὰς τὸν Δ ἀριθμόν; μετρεῖ ἄρα καὶ τὸ Γ τὸ Α. καὶ ἐπεί ἐστιν
ὡς τὸ Γ πρὸς τὸ Α, οὕτως ἡ μονὰς πρὸς τὸν Δ [ἀριθμόν],
ἀνάπαλιν ἄρα ὡς τὸ Α πρὸς τὸ Γ, οὕτως ὁ Δ ἀριθμὸς πρὸς
τὴν μονάδα. πάλιν ἐπεί, ὅσαι εἰσὶν ἐν τῷ Ε μονάδες,
τοσαῦτά εἰσι καὶ ἐν τῷ Ζ ἴσα τῷ Γ, ἔστιν ἄρα ὡς
τὸ Γ πρὸς τὸ Ζ, οὕτως ἡ μονὰς πρὸς τὸν Ε [ἀριθμόν]. ἐδείθθη
δὲ καὶ ὡς τὸ Α πρὸς τὸ Γ, οὕτως ὁ Δ πρὸς τὴν μονάδα; δι
ἴσου ἄρα ἐστὶν ὡς τὸ Α πρὸς τὸ Ζ, οὕτως ὁ Δ πρὸς τὸν
Ε. ἀλλ ὡς ὁ Δ πρὸς τὸν Ε, οὕτως ἐστὶ τὸ Α πρὸς τὸ Β; καὶ
ὡς ἄρα τὸ Α πρὸς τὸ Β, οὕτως καὶ πρὸς τὸ Ζ. τὸ Α ἄρα πρὸς
ἑκάτερον τῶν Β, Ζ τὸν α ὐτὸν ἔχει λόγον; ἴσον ἄρα ἐστὶ
τὸ Β τῷ Ζ. μετρεῖ δὲ τὸ Γ τὸ Ζ; μετρεῖ ἄρα καὶ τὸ Β. ἀλλὰ
μὴν καὶ τὸ Α; τὸ Γ ἄρα τὰ Α, Β μετρεῖ. σύμμετρον ἄρα ἐστὶ
τὸ Α τῷ Β.}

{\greekfont Ἐὰν ἄρα δύο μεγέθη πρὸς ἄλληλα, καὶ τὰ ἑχῆς.}\\

\begin{center}
{\large {\greekfont Πόρισμα}.}
\end{center}\vspace*{-7pt}

{\greekfont Ἐκ δὴ τούτου φανερόν, ὅτι, ἐὰν ὦσι δύο ἀριθμοί, ὡς οἱ
Δ, Ε, καὶ εὐθεῖα, ὡς ἡ Α, δύνατόν ἐστι ποιῆσαι ὡς ὁ Δ ἀριθμὸς
πρὸς τὸν Ε ἀριθμόν, οὕτως τὴν εὐθεῖαν πρὸς εὐθεῖαν. ἐὰν δὲ καὶ
τῶν 
Α, Ζ μέση ἀνάλογον ληφθῇ, ὡς ἡ Β, ἔσται ὡς ἡ Α πρὸς τὴν Ζ,
οὕτως τὸ ἀπὸ τῆς Α πρὸς τὸ ἀπὸ τῆς Β, τουτέστιν ὡς ἡ
πρώτη πρὸς τὴν τρίτην, οὕτως τὸ ἀπὸ τῆς πρώτης πρὸς τὸ ἀπὸ
τῆς δευτέρας τὸ ὅμοιον καὶ ὁμοίως ἀναγραφόμενον. ἀλλ ὡς
ἡ Α πρὸς τὴν Ζ, οὕτως ἐστὶν ὁ Δ ἀριθμος πρὸς τὸν Ε ἀριθμόν;
γέγονεν ἄρα καὶ ὡς ὁ 
Δ ἀριθμὸς πρὸς τὸν Ε ἀριθμόν, οὕτως
τὸ ἀπὸ τῆς Α εὐθείας πρὸς τὸ ἀπὸ τῆς Β εὐθείας; ὅπερ
ἔδει δεῖχαι.}}

\ParallelRText{
\begin{center}
{\large Proposition 6}
\end{center}

If two magnitudes have to one another the ratio
which (some) number (has) to (some) number then the magnitudes
will be commensurable.

\epsfysize=0.8in
\centerline{\epsffile{Book10/fig006e.eps}}

For let the two magnitudes $A$ and $B$ have to one another the ratio which
the number $D$ (has) to the number $E$. I say that the magnitudes $A$ and $B$ are commensurable.

For, as many units as there are in $D$, let $A$ be divided into so many
equal (divisions). And let $C$ be equal to one of them. And as many
units as there are in $E$, let $F$ be the sum  of so many magnitudes equal to
$C$.

Therefore, since as many units as there are in $D$, so many magnitudes equal to $C$ are also in $A$, therefore whichever part a unit is of $D$, $C$ is also
the same part of $A$. Thus, as $C$ is to $A$, so a unit (is) to $D$ [Def.~7.20]. And a unit measures the number $D$.
Thus, $C$ also measures $A$. And since as $C$ is to $A$, so a unit (is) to
the [number] $D$, thus, inversely, as $A$ (is) to $C$, so the number $D$
(is) to a unit [Prop.~5.7~corr.]. Again, since as many 
units as there are in $E$,  so many (magnitudes) equal to $C$ are also in $F$,
thus as $C$ is to $F$, so a unit (is) to the [number] $E$ [Def.~7.20]. And it was also shown that as $A$ (is) to $C$, so $D$ (is) to a unit. Thus, via equality, as $A$ is to $F$, so $D$ (is)
to $E$ [Prop.~5.22]. But, as $D$ (is) to $E$, 
so $A$ is to $B$. And thus as $A$ (is) to $B$, so (it) also is to $F$ [Prop.~5.11]. Thus,
$A$ has the same ratio to each of $B$ and $F$. Thus, $B$ is equal to
$F$ [Prop.~5.9]. And $C$ measures $F$. Thus,
it also measures $B$. But, in fact, (it) also (measures) $A$. Thus, $C$ measures (both) $A$ and $B$. Thus, $A$ is commensurable with $B$ [Def.~10.1].

Thus,  if two magnitudes \ldots to one another, and so on \ldots.\\

\begin{center}
{\large Corollary}
\end{center}\vspace*{-7pt}

So it is clear, from this, that if there are two numbers, like $D$ and $E$, and
a straight-line, like $A$, then it is possible to contrive that as the number $D$ (is)
to the number $E$, so the straight-line (is) to (another) straight-line ({\rm i.e.,} $F$). And if the mean proportion, (say) $B$,  is taken of $A$ and $F$, 
then as $A$ is to $F$, so the (square) on $A$ (will be) to the (square)
on $B$. That is to say, as the first (is) to the third, so the
(figure) on the first (is) to the similar, and similarly described, (figure) on the second [Prop.~6.19~corr.]. But, as $A$ (is) to
$F$, so the number $D$ is to the number $E$. Thus, it has also been contrived
that as the number $D$ (is) to the number $E$, so the (figure) on the
straight-line $A$ (is) to the (similar figure) on the straight-line $B$. (Which is)
the very thing it was required to show.}
\end{Parallel}

%%%%%%
% Prop 10.7
%%%%%%
\pdfbookmark[1]{Proposition 10.7}{pdf10.7}
\begin{Parallel}{}{} 
\ParallelLText{
\begin{center}
{\large\ggn{7}.}
\end{center}\vspace*{-7pt}

{\greekfont Τὰ ἀσύμμετρα μεγέθη πρὸς ἄλληλα λόγον οὐκ ἔχει, ὃν ἀριθμὸς
πρὸς ἀριθμόν.}

{\greekfont Ἔστω ἀσύμμετρα μεγέθη τὰ Α, Β; λέγω, ὅτι τὸ Α πρὸς τὸ Β
λόγον οὐκ ἔχει, ὃν ἀριθμὸς πρὸς ἀριθμόν.}\\

\epsfysize=0.5in
\centerline{\epsffile{Book10/fig007g.eps}}

{\greekfont Εἰ γὰρ ἔχει τὸ Α πρὸς τὸ Β λόγον, ὃν ἀριθμὸς πρὸς ἀριθμόν,
σύμμετρον ἔσται τὸ Α τῷ Β. οὐκ ἔστι δέ; οὐκ ἄρα τὸ Α πρὸς
τὸ Β λόγον ἔχει, ὃν ἀριθμὸς πρὸς ἀριθμόν.}

{\greekfont Τὰ ἄρα ἀσύμμετρα μεγέθη πρὸς ἄλληλα λόγον οὐκ ἔχει, καὶ
τὰ ἑχῆς.}}

\ParallelRText{
\begin{center}
{\large Proposition 7}
\end{center}

Incommensurable magnitudes do not
have to one another the ratio which (some) number (has) to (some)
number.

Let $A$ and $B$ be incommensurable magnitudes. I say that $A$ does not
have to $B$ the ratio which (some) number (has) to (some) number.

\epsfysize=0.5in
\centerline{\epsffile{Book10/fig007e.eps}}

For if $A$ has to $B$ the ratio which (some) number (has) to (some)
number then $A$ will be commensurable with $B$ [Prop.~10.6]. But it is not. Thus, $A$ does not
have to $B$ the ratio which (some) number (has) to (some) number.

Thus, incommensurable numbers do not have to one another, and
so on \ldots.}
\end{Parallel}

%%%%%%
% Prop 10.8
%%%%%%
\pdfbookmark[1]{Proposition 10.8}{pdf10.8}
\begin{Parallel}{}{} 
\ParallelLText{
\begin{center}
{\large \ggn{8}.}
\end{center}\vspace*{-7pt}

{\greekfont Ἐὰν δύο μεγέθη πρὸς ἄλληλα λόγον μὴ ἔθῃ, ὃν ἀριθμὸς πρὸς
ἀριθμόν, ἀσύμμετρα ἔσται τὰ μεγέθη.}\\

\epsfysize=0.5in
\centerline{\epsffile{Book10/fig007g.eps}}

{\greekfont Δύο γὰρ μεγέθη τὰ Α, Β πρὸς ἄλληλα λόγον μὴ ἐθέτω, ὃν 
ἀριθμὸς πρὸς ἀριθμόν; λέγω, ὅτι ἀσύμμετρά ἐστι τὰ Α, Β
μεγέθη.}

{\greekfont Εἰ γὰρ ἔσται σύμμετρα, τὸ Α πρὸς τὸ Β λόγον ἕχει, ὃν ἀριθμὸς
πρὸς ἀριθμόν. οὐκ ἔχει δέ. ἀσύμμετρα ἄρα ἐστὶ τὰ Α, Β μεγέθη.}

{\greekfont Ἐὰν ἄρα δύο μεγέθη πρὸς ἄλληλα, καὶ τὰ ἑχῆς.}}

\ParallelRText{
\begin{center}
{\large Proposition 8}
\end{center}

If two magnitudes do not have to one another the
ratio which (some) number (has) to (some) number then the magnitudes
will be incommensurable.

\epsfysize=0.5in
\centerline{\epsffile{Book10/fig007e.eps}}

For let the two magnitudes $A$ and $B$ not have to one another the ratio
which (some) number (has) to (some) number. I say that the magnitudes
$A$ and $B$ are incommensurable.

For if they are commensurable, $A$ will have to $B$ the ratio which (some)
number (has) to (some) number [Prop.~10.5]. But it does not have (such a ratio).
Thus, the magnitudes $A$ and $B$ are incommensurable.

Thus, if two magnitudes \ldots to one another, and so on \ldots.\\~\\}
\end{Parallel}

%%%%%%
% Prop 10.9
%%%%%%
\pdfbookmark[1]{Proposition 10.9}{pdf10.9}
\begin{Parallel}{}{} 
\ParallelLText{
\begin{center}
{\large \ggn{9}.}
\end{center}\vspace*{-7pt}

{\greekfont Τὰ ἀπὸ τῶν μήκει συμμέτρων εὐθειῶν τετράγωνα πρὸς ἄλληλα
λόγον ἔχει, ὃν τετράγωνος ἀριθμὸς πρὸς τετράγωνον ἀριθμόν;
καὶ τὰ τετράγωνα τὰ πρὸς ἄλληλα λόγον ἔθοντα, ὃν τετράγωνος ἀριθμὸς πρὸς τετράγωνον ἀριθμόν, καὶ τὰς πλευρὰς ἕχει μήκει
συμμέτρους. τὰ δὲ ἀπὸ τῶν μήκει ἀσυμμέτρων εὐθειῶν
τετράγωνα πρὸς ἄλληλα λόγον οὐκ ἔχει,  ὅνπερ
τετράγωνος ἀριθμὸς πρὸς τετράγωνον ἀριθμόν; καὶ τὰ τετράγωνα
τὰ πρὸς ἄλληλα λόγον μὴ ἔθοντα, ὃν τετράγωνος ἀριθμὸς
πρὸς τετράγωνον ἀριθμόν, οὐδὲ τὰς πλευρὰς ἕχει μήκει
συμμέτρους.}\\

\epsfysize=0.5in
\centerline{\epsffile{Book10/fig009g.eps}}

{\greekfont Ἔστωσαν γὰρ αἱ Α, Β μήκει σύμμετροι; λέγω, ὅτι τὸ ἀπὸ τῆς
Α τετράγωνον πρὸς τὸ ἀπὸ τῆς Β τετράγωνον λόγον ἔχει,
ὃν τετράγωνος ἀριθμὸς πρὸς τετράγωνον ἀριθμόν.}

{\greekfont Ἐπεὶ γὰρ σύμμετρός ἐστιν ἡ Α τῇ Β μήκει, ἡ Α ἄρα πρὸς τὴν
Β λόγον ἔχει, ὃν ἀριθμὸς πρὸς ἀριθμόν. ἐθέτω, ὃν ὁ Γ πρὸς
τὸν Δ. ἐπεὶ οὖν ἐστιν ὡς ἡ Α πρὸς τὴν Β, οὕτως ὁ Γ
πρὸς τὸν Δ, ἀλλὰ τοῦ μὲν τῆς Α πρὸς τὴν Β λόγου διπλασίων
ἐστὶν ὁ τοῦ ἀπὸ τῆς Α τετραγώνου πρὸς τὸ ἀπὸ τῆς Β τετράγωνον;
τὰ γὰρ ὅμοια σθήματα ἐν διπλασίονι λόγῳ ἐστὶ τῶν
ὁμολόγων πλευρῶν; τοῦ δὲ τοῦ Γ [ἀριθμοῦ] πρὸς τὸν Δ
[ἀριθμὸν] λόγου διπλασίων ἐστὶν ὁ τοῦ ἀπὸ τοῦ Γ
τετραγώνου πρὸς τὸν ἀπὸ τοῦ Δ τετράγωνον; δύο γὰρ τετραγώνων
ἀριθμῶν εἷς μέσος ἀνάλογόν ἐστιν ἀριθμός, καί ὁ τετράγωνος
πρὸς τὸν τετράγωνον [ἀριθμὸν] διπλασίονα λόγον ἔχει, ἤπερ
ἡ πλευρὰ πρὸς τὴν πλευράν; ἔστιν ἄρα καὶ ὡς τὸ ἀπὸ τῆς
Α τετράγωνον πρὸς τὸ ἀπὸ τῆς Β τετράγωνον, οὕτως ὁ 
ἀπὸ τοῦ Γ τετράγωνος [ἀριθμὸς] πρὸς τὸν ἀπὸ τοῦ Δ [ἀριθμοῦ]
τετράγωνον [ἀριθμόν].}

{\greekfont Ἀλλὰ δὴ ἔστω ὡς τὸ ἀπὸ τῆς Α τετράγωνον πρὸς τὸ ἀπὸ
τῆς Β, οὕτως ὁ ἀπὸ τοῦ Γ τετράγωνος πρὸς τὸν ἀπὸ τοῦ
Δ [τετράγωνον]; λέγω, ὅτι σύμμετρός ἐστιν ἡ Α τῇ Β μήκει.}

{\greekfont Ἐπεὶ γάρ ἐστιν ὡς τὸ ἀπὸ τῆς Α τετράγωνον πρὸς τὸ ἀπὸ
τῆς Β [τετράγωνον], οὕτως ὁ ἀπὸ τοῦ Γ τετράγωνος πρὸς
τὸν ἀπὸ τοῦ Δ [τετράγωνον], ἀλλ ὁ μὲν τοῦ ἀπὸ τῆς Α
τετραγώνου πρὸς τὸ ἀπὸ τῆς Β [τετράγωνον] λόγος
διπλασίων ἐστὶ τοῦ τῆς Α πρὸς τὴν Β λόγου, ὁ δὲ τοῦ ἀπὸ τοῦ Γ [ἀριθμοῦ]
τετραγώνου [ἀριθμοῦ] πρὸς τὸν ἀπὸ τοῦ Δ [ἀριθμοῦ]
τετράγωνον [ἀριθμὸν] λόγος διπλασίων ἐστὶ τοῦ τοῦ Γ
[ἀριθμοῦ] πρὸς τὸν Δ [ἀριθμὸν] λόγου, ἔστιν ἄρα καὶ ὡς
ἡ Α πρὸς τὴν Β, οὕτως ὁ Γ [ἀριθμὸς] πρὸς τὸν Δ [ἀριθμόν]. ἡ Α ἄρα πρὸς τὴν Β λόγον ἔχει, ὃν ἀριθμὸς ὁ Γ
πρὸς ἀριθμὸν τὸν Δ; σύμμετρος ἄρα ἐστὶν ἡ Α τῇ Β μήκει.}

{\greekfont Ἀλλὰ δὴ ἀσύμμετρος ἔστω ἡ Α τῇ Β μήκει; λέγω, ὅτι τὸ
ἀπὸ τῆς Α τετράγωνον πρὸς τὸ ἀπὸ τῆς Β [τετράγωνον] λόγον
οὐκ ἔχει, ὃν τετράγωνος ἀριθμὸς πρὸς τετράγωνον ἀριθμόν.}

{\greekfont Εἰ γὰρ ἔχει τὸ ἀπὸ τῆς Α τετράγωνον πρὸς τὸ ἀπὸ τῆς Β
[τετράγωνον] λόγον, ὃν τετράγωνος ἀριθμὸς πρὸς τετράγωνον
ἀριθμόν, σύμμετρος ἔσται ἡ Α τῇ Β. οὐκ ἔστι δέ; οὐκ  ἄρα τὸ
ἀπὸ τῆς Α τετράγωνον πρὸς τὸ ἀπὸ τῆς Β [τετράγωνον] λόγον
ἔχει, ὃν τετράγωνος ἀριθμὸς πρὸς τετράγωνον ἀριθμόν.}

{\greekfont Πάλιν δὴ τὸ ἀπὸ τῆς Α τετράγωνον πρὸς τὸ ἀπὸ τῆς Β
[τετράγωνον] λόγον μὴ ἐθέτω, ὃν τετράγωνος ἀριθμὸς πρὸς
τετράγωνον ἀριθμόν; λέγω, ὅτι ἀσύμμετρός ἐστιν ἡ Α τῇ Β
μήκει.}

{\greekfont Εἰ γάρ ἐστι σύμμετρος ἡ Α τῇ Β, ἕχει τὸ ἀπὸ τῆς Α πρὸς
τὸ ἀπὸ τῆς Β λόγον, ὃν τετράγωνος ἀριθμὸς πρὸς
τετράγωνον ἀριθμόν. οὐκ ἔχει δέ; οὐκ ἄρα σύμμετρός ἐστιν
ἡ Α τῇ Β μήκει.}

{\greekfont Τὰ ἄρα ἀπὸ τῶν μήκει συμμέτρων, καὶ τὰ ἑχῆς.}\\~\\~\\

\begin{center}
{\large {\greekfont Πόρισμα}.}
\end{center}

{\greekfont Καὶ φανερὸν ἐκ τῶν δεδειγμένων ἔσται, ὅτι αἱ μήκει σύμμετροι
πάντως καὶ δυνάμει, αἱ δὲ δυνάμει οὐ πάντως καὶ μήκει.}}

\ParallelRText{
\begin{center}
{\large Proposition 9}
\end{center}

Squares on straight-lines (which are) commensurable in length have to one another the ratio which (some) square number (has) to (some)
square number. And squares having to one another the ratio which (some)
square number (has) to (some) square number will also have sides (which are)
commensurable in length. But squares on straight-lines (which are) incommensurable
in length do not have to one another the ratio which (some) square
number (has) to (some) square number. And squares not having
to one another the ratio which (some) square number (has) to (some)
square number will not have sides (which are) commensurable in length either.

\epsfysize=0.5in
\centerline{\epsffile{Book10/fig009e.eps}}

For let $A$ and $B$ be (straight-lines which are) commensurable in length.
I say that the square on $A$ has to the square on $B$ the ratio
which (some) square number (has) to (some) square number.

For since $A$ is commensurable in length with $B$, $A$ thus
has to $B$ the ratio which (some) number (has) to (some) number [Prop.~10.5]. Let it have (that) which $C$ (has) to $D$.
Therefore, since as $A$ is to $B$, so $C$ (is) to $D$. But the (ratio)
of the square on $A$ to the square on $B$ is the square of the
ratio of $A$ to $B$. For similar figures are in the squared ratio of (their) corresponding sides [Prop.~6.20~corr.].
And the (ratio) of the square on $C$ to the square on $D$ is the
square of the ratio of the [number] $C$ to the [number] $D$.
For there exits one number in mean proportion to two square numbers,
and (one) square (number) has to the (other) square [number] a squared ratio
with respect to (that) the side (of the former has) to the side (of the latter)
[Prop.~8.11]. And, thus,  as the square on $A$
is to the square on $B$, so the square  [number] on the (number) $C$ (is) to
the square [number] on the [number] $D$.$^\dag$

And so let the square on $A$ be to the (square) on $B$ as the
square (number) on $C$ (is) to the [square] (number) on $D$. I say
that $A$ is commensurable in length  with $B$.

For since as the square on $A$ is to the [square] on $B$, so the
square (number) on $C$ (is) to the [square] (number) on $D$. But,
the ratio of the square on $A$ to the (square) on $B$ is the square
of the (ratio) of $A$ to $B$ [Prop.~6.20 corr.].
And the (ratio) of the square [number] on the [number] $C$ to
the square [number] on the [number] $D$ is the square of the ratio
of the [number] $C$ to the [number] $D$ [Prop.~8.11].  Thus, as $A$ is to $B$, so the
[number] $C$ also (is) to the [number] $D$.
$A$, thus,  has to $B$ the ratio
which the number $C$ has to the number $D$. Thus, $A$ is
commensurable in length with $B$ [Prop.~10.6].$^\ddag$

And so let $A$ be incommensurable in length with $B$. I say that
 the square on $A$ does not have to the [square] on $B$ the ratio
which (some) square number (has) to (some) square number.

For if the square on $A$ has to the [square] on $B$ the ratio which
(some) square number (has) to (some) square number then $A$
will be commensurable (in length) with $B$. But it is not. Thus, the square on $A$
does not have to the [square] on the $B$ the ratio which (some)
square number (has) to (some) square number.

So, again, let the square on $A$ not have to the [square] on $B$
the ratio which (some) square number (has) to (some) square number.
I say that $A$ is incommensurable in length with $B$.

For if $A$ is commensurable (in length) with $B$ then  the 
(square) on $A$ will have to the (square) on $B$ the ratio which (some) square
number (has) to (some) square number. But it does not have (such a ratio). 
Thus, $A$ is not commensurable in length with $B$.

Thus, (squares) on (straight-lines which are) commensurable in length,
and so on \ldots.\\

\begin{center}
{\large Corollary}
\end{center}\vspace*{-7pt}

And it will be clear, from (what) has been demonstrated, that (straight-lines)
commensurable in length (are) always also (commensurable) in square, 
but (straight-lines commensurable) in square (are) not always also (commensurable)
in length.}
\end{Parallel}
{\footnotesize\noindent$^\dag$ There is an unstated assumption here that if $\alpha:\beta::\gamma:\delta$ then $\alpha^2:\beta^2::\gamma^2:\delta^2$.\\[0.5ex]
$^\ddag$ There is an unstated
assumption here that if $\alpha^2:\beta^2::\gamma^2:\delta^2$ then $\alpha:\beta::\gamma:\delta$.}

%%%%%%
% Prop 10.10
%%%%%%
\pdfbookmark[1]{Proposition 10.10}{pdf10.10}
\begin{Parallel}{}{} 
\ParallelLText{
\begin{center}
{\large \ggn{10}.}
\end{center}\vspace*{-7pt}

{\greekfont Τῇ προτεθείσῃ εὐθείᾳ προσευρεῖν δύο εὐθείας ἀσυμμέτ-ρους,
τὴν μὲν μήκει μόνον, τὴν δὲ καὶ δυνάμει.}\\

\epsfysize=1.5in
\centerline{\epsffile{Book10/fig010g.eps}}

{\greekfont Ἔστω ἡ προτεθεῖσα εὐθεῖα ἡ Α; δεῖ δὴ τῇ Α προσευρεῖν
δύο εὐθείας ἀσυμμέτρους, τὴν μὲν μήκει μόνον, τὴν δὲ
καὶ δυνάμει.}

{\greekfont Ἐκκείσθωσαν γὰρ δύο αριθμοὶ οἱ Β, Γ πρὸς ἀλλήλους λόγον
μὴ ἔθοντες, ὃν τετράγωνος ἀριθμὸς πρὸς τετράγωνον
ἀριθμόν, τουτέστι μὴ ὅμοιοι ἐπίπεδοι, καὶ γεγονέτω
ὡς ὁ Β πρὸς τὸν Γ, οὕτως τὸ ἀπὸ τῆς Α τετράγωνον πρὸς
τὸ ἀπὸ τῆς Δ τετράγωνον; ἐμάθομεν γάρ; σύμμετρον ἄρα τὸ
ἀπὸ τῆς Α τῷ ἀπὸ τῆς Δ. καὶ ἐπεὶ ὁ Β πρὸς τὸν Γ λόγον
οὐκ ἔχει, ὃν τετράγωνος ἀριθμὸς πρὸς τετράγωνον ἀριθμόν,
οὐδ ἄρα τὸ ἀπὸ τῆς Α πρὸς τὸ ἀπὸ τῆς Δ λόγον ἔχει,
ὃν τετράγωνος ἀριθμὸς πρὸς τετράγωνον ἀριθμόν; ἀσύμμετρος
 ἄρα ἐστὶν
ἡ Α τῇ Δ μήκει. εἰλήφθω τῶν Α, Δ μέση ἀνάλογον ἡ Ε; ἔστιν
ἄρα ὡς ἡ Α πρὸς τὴν Δ, οὕτως τὸ ἀπὸ τῆς Α τετράγωνον
πρὸς τὸ ἀπὸ τῆς Ε. ἀσύμμετρος δέ ἐστιν ἡ Α τῇ Δ μήκει;
ἀσύμμετρον ἄρα ἐστὶ καὶ τὸ ἀπὸ τῆς Α τετράγωνον
τῷ ἀπὸ τῆς Ε τετραγώνῳ; ἀσύμμετρος ἄρα ἐστὶν ἡ Α τῇ
Ε δυνάμει.}

{\greekfont Τῂ ἄρα προτεθείσῃ εὐθείᾳ τῇ Α προσεύρηνται δύο εὐθεῖαι
ἀσύμμετροι αἱ Δ, Ε, μήκει μὲν μόνον ἡ Δ,
δυνάμει δὲ καὶ μήκει δηλαδὴ ἡ Ε [ὅπερ ἔδει δεῖχαι].}}

\ParallelRText{
\begin{center}
{\large Proposition 10}$^\dag$
\end{center}

To find two straight-lines incommensurable
with a given straight-line, the one (incommensurable) in length only,
the other also (incommensurable) in square.

\epsfysize=1.5in
\centerline{\epsffile{Book10/fig010e.eps}}

Let $A$ be the given straight-line. So it is required to find two straight-lines
incommensurable with $A$, the one (incommensurable) in length only,
the other also (incommensurable) in square.

For let two numbers, $B$ and $C$,  not having to one another
the ratio which (some) square number (has) to (some) square number---that
is to say, not (being) similar plane (numbers)---be taken. And let it be contrived
that as $B$ (is) to $C$, so the square on $A$ (is) to the square on $D$.
For we learned (how to do this) [Prop.~10.6 corr.].
Thus, the (square) on $A$ (is) commensurable with the (square) on $D$
[Prop.~10.6]. And since $B$ does not have to $C$ the
ratio which (some) square number (has) to (some) square number, the
(square) on $A$ thus does not have to the (square) on $D$ the ratio which (some) square number (has) to (some) square number either. Thus, $A$ is
incommensurable in length with $D$ [Prop.~10.9].
Let the (straight-line) $E$ (which is) in mean proportion to $A$ and $D$ 
be taken [Prop.~6.13]. Thus, as $A$ is to $D$, so the square on $A$ (is) to the (square)
on $E$ [Def.~5.9]. And $A$ is incommensurable in length with $D$. Thus, the square on $A$ is also incommensurble
with the square on $E$ [Prop.~10.11]. 
Thus, $A$ is incommensurable in square with $E$.

Thus, two straight-lines, $D$ and $E$, (which are) incommensurable with the given
straight-line $A$, be found, the one, $D$, (incommensurable)
in length only, the other, $E$, (incommensurable) in square, and, clearly, also in length. [(Which is) the very thing it was required to show.]}
\end{Parallel}
{\footnotesize\noindent$^\dag$ This whole proposition is regarded by Heiberg as an interpolation into the original text.}

%%%%%%
% Prop 10.11
%%%%%%
\pdfbookmark[1]{Proposition 10.11}{pdf10.11}
\begin{Parallel}{}{} 
\ParallelLText{
\begin{center}
{\large \ggn{11}.}
\end{center}\vspace*{-7pt}

{\greekfont Ἐὰν τέσσαρα μεγέθη ἀνάλογον ᾖ, τὸ δὲ πρῶτον τῷ δευτέρῳ
σύμμετρον ᾖ, καὶ τὸ τρίτον τῷ τετάρτῳ σύμμετρον ἔσται; κἂν
τὸ πρῶτον τῷ δευτέρῳ ἀσύμμετρον ᾖ, καὶ τὸ τρίτον τῷ
τετάρτῳ ἀσύμμετρον ἔσται.}\\

\epsfysize=0.5in
\centerline{\epsffile{Book10/fig011g.eps}}

{\greekfont Ἔστωσαν τέσσαρα μεγέθη ἀνάλογον τὰ Α, Β, Γ, Δ, ὡς τὸ Α
πρὸς τὸ Β, οὕτως τὸ Γ πρὸς τὸ Δ, τὸ Α δὲ τῷ Β σύμμετρον
ἔστω; λέγω,  ὅτι καὶ τὸ Γ τῷ Δ σύμμετρον ἔσται.}

{\greekfont Ἐπεὶ γὰρ σύμμετρόν ἐστι τὸ Α τῷ Β, τὸ Α ἄρα πρὸς τὸ Β
λόγον ἔχει, ὃν ἀριθμὸς πρὸς ἀριθμόν. καί ἐστιν ὡς τὸ Α
πρὸς τὸ Β, οὕτως τὸ Γ πρὸς τὸ Δ; καὶ τὸ Γ ἄρα πρὸς τὸ Δ
λόγον ἔχει, ὃν ἀριθμὸς πρὸς ἀριθμόν; σύμμετρον ἄρα
ἐστὶ τὸ Γ τῷ Δ.}

{\greekfont Ἀλλὰ δὴ τὸ Α τῷ Β ἀσύμμετρον ἔστω; λέγω, ὅτι καὶ τὸ
Γ τῷ Δ ἀσύμμετρον ἔσται. ἐπεὶ γὰρ ἀσύμμετρόν ἐστι
τὸ Α τῷ Β, τὸ Α ἄρα πρὸς τὸ Β λόγον οὐκ ἔχει, ὃν ἀριθμὸς
πρὸς ἀριθμόν. καί ἐστιν ὡς τὸ Α πρὸς τὸ Β, οὕτως τὸ Γ πρὸς τὸ Δ;
οὐδὲ τὸ Γ ἄρα πρὸς τὸ Δ λόγον ἔχει, ὃν ἀριθμὸς πρὸς
ἀριθμόν; ἀσύμμετρον ἄρα ἐστὶ τὸ Γ τῷ Δ.}

{\greekfont Ἐὰν ἄρα τέσσαρα μεγέθη, καὶ τὰ ἑχῆς.}}

\ParallelRText{
\begin{center}
{\large Proposition 11}
\end{center}

If four magnitudes are proportional, and the
first is commensurable with the second, then the third will also be
commensurable with the fourth. And if the first is incommensurable with
the second, then the third will also be incommensurable with the fourth.

\epsfysize=0.5in
\centerline{\epsffile{Book10/fig011e.eps}}

Let $A$, $B$, $C$, $D$ be four proportional magnitudes, (such that) as
$A$ (is) to $B$, so $C$ (is) to $D$. And let $A$ be commensurable
with $B$. I say that $C$ will also be commensurable with $D$.

For since $A$ is commensurable with $B$, $A$ thus has to $B$ the ratio
which (some) number (has) to (some) number [Prop.~10.5]. And as $A$ is to $B$, so $C$ (is) to $D$. Thus, $C$ also has to $D$ the ratio which (some) number (has) to
(some) number. Thus, $C$ is commensurable with $D$ [Prop.~10.6].

And so let $A$ be incommensurable with $B$. I say that $C$ will also
be incommensurable with $D$. For since $A$ is incommensurable with
$B$, $A$ thus does not have to $B$ the ratio which  (some) number (has) to
(some) number [Prop.~10.7]. And as $A$ is to $B$, so $C$ (is) to $D$. Thus, $C$ does not have to $D$ the ratio which (some) number (has) to
(some) number either. Thus, $C$ is incommensurable with $D$ [Prop.~10.8].

Thus, if four magnitudes, and so on \ldots.}
\end{Parallel}

%%%%%%
% Prop 10.12
%%%%%%
\pdfbookmark[1]{Proposition 10.12}{pdf10.12}
\begin{Parallel}{}{} 
\ParallelLText{
\begin{center}
{\large \ggn{12}.}
\end{center}\vspace*{-7pt}

{\greekfont Τὰ τῷ α ὐτῷ μεγέθει σύμμετρα καὶ ἀλλήλοις ἐστὶ σύμμετρα.}

{\greekfont Ἑκάτερον γὰρ τῶν Α, Β τῷ Γ ἔστω σύμμετρον. λέγω, ὅτι
καὶ τὸ Α τῷ Β ἐστι σύμμετρον.}

{\greekfont Ἐπεὶ γὰρ σύμμετρόν ἐστι τὸ Α τῷ Γ, τὸ Α ἄρα πρὸς τὸ Γ λόγον
ἔχει, ὃν ἀριθμὸς πρὸς ἀριθμόν. ἐθέτω, ὃν ὁ Δ πρὸς τὸν Ε.
πάλιν, ἐπεὶ σύμμετρόν ἐστι τὸ Γ τῷ Β, τὸ Γ ἄρα πρὸς τὸ Β λόγον
ἔχει, ὃν ἀριθμὸς πρὸς ἀριθμόν. ἐθέτω, ὃν ὁ Ζ πρὸς τὸν Η.
καὶ λόγων δοθέντων ὁποσωνοῦν τοῦ τε, ὃν ἔχει ὁ Δ πρὸς τὸν
Ε, καὶ ὁ Ζ πρὸς τὸν Η εἰλήφθωσαν ἀριθμοὶ ἑχῆς ἐν τοῖς δοθεῖσι
λόγοις οἱ Θ, Κ, Λ; ὥστε εἶναι ὡς μὲν τὸν Δ πρὸς τὸν Ε, οὕτως
τὸν Θ πρὸς τὸν Κ, ὡς δὲ τὸν Ζ πρὸς τὸν Η, οὕτως  τὸν Κ πρὸς τὸν Λ.}\\~\\~\\

\epsfysize=1.in
\centerline{\epsffile{Book10/fig012g.eps}}

{\greekfont Ἐπεὶ οὖν ἐστιν ὡς τὸ Α πρὸς τὸ Γ, οὕτως ὁ Δ πρὸς τὸν Ε,
ἀλλ ὡς ὁ Δ πρὸς τὸν Ε, οὕτως ὁ Θ πρὸς τὸν Κ, ἔστιν ἄρα
καὶ ὡς τὸ Α πρὸς τὸ Γ, οὕτως ὁ Θ πρὸς τὸν Κ. πάλιν, ἐπεί ἐστιν
ὡς τὸ Γ πρὸς τὸ Β, οὕτως ὁ Ζ πρὸς τὸν Η, ἀλλ ὡς ὁ Ζ πρὸς
τὸν Η, [οὕτως] ὁ Κ πρὸς τὸν Λ, καὶ ὡς ἄρα τὸ Γ πρὸς
τὸ Β, οὕτως ὁ Κ πρὸς τὸν Λ. ἔστι δὲ καὶ ὡς τὸ Α πρὸς τὸ
Γ, οὕτως ὁ Θ πρὸς τὸν Κ; δι ἴσου ἄρα ἐστὶν ὡς τὸ Α πρὸς
τὸ Β, οὕτως ὁ Θ πρὸς τὸν Λ. τὸ Α ἄρα πρὸς τὸ Β λόγον
ἔχει, ὃν ἀριθμὸς ὁ Θ πρὸς ἀριθμὸν τὸν Λ;
σύμμετρον ἄρα ἐστὶ τὸ Α τῷ Β.}

{\greekfont Τὰ ἄρα τῷ α ὐτῷ μεγέθει σύμμετρα καὶ ἀλλήλοις ἐστὶ
σύμμετρα; ὅπερ ἔδει δεῖχαι.}}

\ParallelRText{
\begin{center}
{\large Proposition 12}
\end{center}

(Magnitudes) commensurable with the same magnitude are also commensurable with one another.

For let $A$ and $B$ each be commensurable with $C$. I say that
$A$ is also commensurable with $B$.

For since $A$ is commensurable with $C$, $A$ thus has to $C$ the
ratio which (some) number (has) to (some) number [Prop.~10.5]. Let it have (the ratio) which $D$ (has)
to $E$. Again, since $C$ is commensurable with $B$, $C$ thus has to $B$
the ratio which (some) number (has) to (some) number [Prop.~10.5].  Let it have (the ratio) which $F$ (has) to $G$. And for any multitude whatsoever of given ratios---(namely,) those
which $D$ has to $E$, and $F$ to $G$---let the numbers $H$, $K$, $L$
(which are) continuously (proportional) in the(se) given ratios be taken [Prop.~8.4]. Hence, as $D$ is to $E$, so $H$ (is) to $K$, and as $F$ (is) to $G$, so $K$ (is) to $L$.

\epsfysize=1.in
\centerline{\epsffile{Book10/fig012e.eps}}

Therefore, since as $A$ is to $C$, so $D$ (is) to $E$, but as $D$ (is) to 
$E$, so $H$ (is) to $K$, thus also as $A$ is to $C$, so $H$ (is) to $K$
[Prop.~5.11]. Again, since as $C$ is to $B$, 
so $F$ (is) to $G$, but as $F$ (is) to $G$, [so] $K$ (is) to $L$, 
thus also as $C$ (is) to $B$, so $K$ (is) to $L$ [Prop.~5.11]. And also as $A$ is to $C$, so $H$ (is) to $K$. Thus, via equality, as $A$ is to $B$, so $H$ (is) to $L$ [Prop.~5.22]. Thus, $A$ has to $B$ the ratio
which the number $H$ (has) to the number $L$. Thus, $A$ is commensurable
with $B$ [Prop.~10.6].

Thus, (magnitudes) commensurable with the same magnitude are also commensurable with one another. (Which is) the very thing it was required
to show.}
\end{Parallel}

%%%%%%
% Prop 10.13
%%%%%%
\pdfbookmark[1]{Proposition 10.13}{pdf10.13}
\begin{Parallel}{}{} 
\ParallelLText{
\begin{center}
{\large\ggn{13}.}
\end{center}\vspace*{-7pt}

{\greekfont Ἐὰν ᾖ δύο μεγέθη σύμμετρα, τὸ δὲ ἕτερον α ὐτῶν μεγέθει
τινὶ ἀσύμμετρον ᾖ, καὶ τὸ λοιπὸν τῷ α ὐτῷ
ἀσύμμετρ-ον ἔσται.}\\

\epsfysize=0.8in
\centerline{\epsffile{Book10/fig013g.eps}}

{\greekfont Ἔστω δύο μεγέθη σύμμετρα τὰ Α, Β, τὸ δὲ ἕτερον α ὐτῶν
τὸ Α ἄλλῳ τινὶ τῷ Γ ἀσύμμετρον ἔστω; λέγω, ὅτι καὶ τὸ
λοιπὸν τὸ Β τῷ Γ ἀσύμμετρόν ἐστιν.}

{\greekfont Εἰ γάρ ἐστι σύμμετρον τὸ Β τῷ Γ, ἀλλὰ καὶ τὸ Α
τῷ Β σύμμετρόν ἐστιν, καὶ τὸ Α ἄρα τῷ Γ σύμμετρόν
ἐστιν. ἀλλὰ καὶ ἀσύμμετρον; ὅπερ ἀδύνατον. οὐκ ἄρα
σύμμετρόν ἐστι τὸ Β τῷ Γ; ἀσύμμετρον ἄρα.}

{\greekfont Ἐὰν ἄρα ᾖ δύο μεγέθη σύμμετρα, καὶ τὰ ἑχῆς.}}

\ParallelRText{
\begin{center}
{\large Proposition 13}
\end{center}

If two magnitudes are commensurable, and
one of them is incommensurable  with some magnitude, then the remaining (magnitude)
will also be incommensurable with it.

\epsfysize=0.8in
\centerline{\epsffile{Book10/fig013e.eps}}

Let $A$ and $B$ be two commensurable magnitudes, and let one
of them, $A$, be incommensurable with some other (magnitude),
$C$. I say that the remaining (magnitude), $B$, is also incommensurable
with $C$.

For if $B$ is commensurable with $C$, but $A$ is also commensurable
with $B$, $A$ is thus also commensurable with $C$ [Prop.~10.12]. But, (it is) also incommensurable (with
$C$). The very thing (is) impossible. Thus, $B$ is not commensurable
with $C$. Thus, (it is) incommensurable.

Thus, if two magnitudes are commensurable, and so on \ldots.}
\end{Parallel}

%%%%%%
% Prop 10.13a
%%%%%%
\begin{Parallel}{}{} 
\ParallelLText{
\begin{center}
{\large {\greekfont Λῆμμα}.}
\end{center}\vspace*{-7pt}

{\greekfont Δύο δοθεισῶν εὐθειῶν ἀνίσων εὑρεῖν, τίνι μεῖζον δύναται
ἡ μείζων τῆς ἐλάσσονος.}\\

\epsfysize=1.2in
\centerline{\epsffile{Book10/fig013ag.eps}}

{\greekfont Ἔστωσαν αἱ δοθεῖσαι δύο ἄνισοι εὐθεῖαι αἱ ΑΒ, Γ, ὧν μείζων
ἔστω ἡ ΑΒ; δεῖ δὴ εὑρεῖν, τίνι μεῖζον δύναται ἡ ΑΒ τῆς Γ.}

{\greekfont Γεγράφθω ἐπὶ τῆς ΑΒ ἡμικύκλιον τὸ ΑΔΒ, καὶ
εἰς α ὐτὸ ἐνηρμόσθω τῇ Γ ἴση ἡ ΑΔ, καὶ ἐπεζεύθθω ἡ ΔΒ.
φανερὸν δή, ὅτι ὀρθή ἐστιν ἡ ὑπὸ ΑΔΒ γωνία, καὶ ὅτι
ἡ ΑΒ τῆς ΑΔ, τουτέστι τῆς Γ, μεῖζον δύναται τῇ ΔΒ.}

{\greekfont Ὁμοίως δὲ καὶ δύο δοθεισῶν εὐθειῶν ἡ δυναμένη α ὐτὰς εὑρίσκεται οὕτως.}

{\greekfont Ἔστωσαν αἱ δοθεῖσαι δύο εὐθεῖαι αἱ ΑΔ, ΔΒ, καὶ δέον ἔστω
εὑρεῖν τὴν δυναμένην α ὐτάς. κείσθωσαν γάρ, ὥστε ὀρθὴν
γωνίαν περιέθειν τὴν ὑπὸ ΑΔ, ΔΒ, καὶ ἐπεζεύθθω ἡ ΑΒ; φανερὸν
πάλιν, ὅτι ἡ τὰς ΑΔ, ΔΒ δυναμένη ἐστὶν ἡ ΑΒ; ὅπερ
ἔδει δεῖχαι.}}

\ParallelRText{
\begin{center}
{\large Lemma}
\end{center}

For two given unequal straight-lines, to find by (the square on) which (straight-line)
the square on the greater (straight-line is) larger   than (the square on) the lesser.$^\dag$

\epsfysize=1.2in
\centerline{\epsffile{Book10/fig013ae.eps}}

Let $AB$ and $C$ be the two given unequal straight-lines, and let $AB$
be the greater of them. So it is required to find by (the square on) which (straight-line)  the square on $AB$ (is) greater  than (the square on) $C$.

Let the semi-circle $ADB$ be described on $AB$.  And let $AD$,
equal to $C$, be inserted into it [Prop.~4.1]. 
And let $DB$ be joined. So (it is) clear that the angle $ADB$
is  a right-angle [Prop.~3.31], and that the square on $AB$
(is) greater than (the square on) $AD$---that is to say, (the square on) $C$---by (the square on) $DB$ [Prop.~1.47].

And, similarly, the square-root of (the sum of the squares on) two given straight-lines
is also found likeso.

Let $AD$ and $DB$ be the two given straight-lines. And let it
be necessary to find the square-root of (the sum of the squares on) them. For let
them be laid down such as to encompass a right-angle---(namely), that (angle encompassed) by
$AD$ and $DB$. And let $AB$ be joined. (It is) again clear that $AB$ is the square-root of   (the sum of the squares on) $AD$ and $DB$ [Prop.~1.47]. (Which is) the very thing it was required to show.}
\end{Parallel}
{\footnotesize\noindent$^\dag$ That is, if $\alpha$ and $\beta$ are the lengths of two given straight-lines, with $\alpha$ being greater
than $\beta$, to find a straight-line of length $\gamma$ such that $\alpha^2=\beta^2+\gamma^2$. Similarly, we can also find $\gamma$
such that $\gamma^2=\alpha^2+\beta^2$.}

%%%%%%
% Prop 10.14
%%%%%%
\pdfbookmark[1]{Proposition 10.14}{pdf10.14}
\begin{Parallel}{}{} 
\ParallelLText{
\begin{center}
{\large \ggn{14}.}
\end{center}\vspace*{-7pt}

{\greekfont Ἐὰν τέσσαρες εὐθεῖαι ἀνάλογον ὦσιν, δύνηται δὲ ἡ πρώτη
τῆς δευτέρας μεῖζον τῷ ἀπὸ συμμέτρου ἑαυτῇ [μήκει],
καὶ ἡ τρίτη τῆς τετάρτης μεῖζον δυνήσεται τῷ ἀπὸ συμμέτρου
ἑαυτῇ [μήκει]. καὶ ἐὰν ἡ πρώτη τῆς δευτέρας μεῖζον δύνηται
τῷ ἀπὸ ἀσυμμέτρου ἑαυτῇ [μήκει], καὶ ἡ τρίτη τῆς τετάρτης
μεῖζον δυνήσεται τῷ ἀπὸ ἀσυμμέτρου ἑαυτῇ [μήκει].}

{\greekfont Ἔστωσαν τέσσαρες εὐθεῖαι ἀνάλογον αἱ Α, Β, Γ, Δ, ὡς ἡ Α
πρὸς τὴν Β, οὕτως ἡ Γ πρὸς τὴν Δ, καὶ ἡ Α μὲν τῆς Β μεῖζον
δυνάσθω τῷ ἀπὸ τῆς Ε, ἡ δὲ Γ τῆς Δ μεῖζον δυνάσθω
τῷ ἀπὸ τῆς Ζ; λέγω, ὅτι, εἴτε σύμμετρός ἐστιν ἡ Α τῇ Ε,
σύμμετρός ἐστι καὶ ἡ Γ τῇ Ζ, εἴτε ἀσύμμετρός ἐστιν ἡ Α
τῇ Ε, ἀσύμμετρός ἐστι καὶ ὁ Γ τῇ Ζ.}\\~\\~\\~\\~\\~\\

\epsfysize=2.in
\centerline{\epsffile{Book10/fig014g.eps}}

{\greekfont Ἐπεὶ γάρ ἐστιν ὡς ἡ Α πρὸς τὴν Β, οὕτως ἡ Γ πρὸς τὴν
Δ, ἔστιν ἄρα καὶ ὡς τὸ ἀπὸ τῆς Α πρὸς τὸ ἀπὸ τῆς Β,
οὕτως τὸ ἀπὸ τῆς Γ πρὸς τὸ ἀπὸ τῆς Δ. ἀλλὰ τῷ μὲν
ἀπὸ τῆς Α ἴσα ἐστὶ τὰ ἀπὸ τῶν Ε, Β, τῷ δὲ ἀπὸ τῆς Γ
ἴσα ἐστὶ τὰ ἀπὸ τῶν Δ, Ζ. ἔστιν ἄρα ὡς τὰ ἀπὸ τῶν Ε, Β
πρὸς τὸ ἀπὸ τῆς Β, οὕτως τὰ ἀπὸ τῶν Δ, Ζ πρὸς τὸ ἀπὸ
τῆς Δ; διελόντι ἄρα ἐστὶν ὡς τὸ ἀπὸ τῆς Ε πρὸς τὸ ἀπὸ
τῆς Β, οὕτως τὸ ἀπὸ τῆς Ζ πρὸς τὸ ἀπὸ τῆς Δ; ἔστιν ἄρα
καὶ ὡς ἡ Ε πρὸς τὴν Β, οὕτως ἡ Ζ πρὸς τὴν Δ; ἀνάπαλιν
ἄρα ἐστὶν ὡς ἡ Β πρὸς τὴν Ε, οὕτως ἡ Δ πρὸς τὴν Ζ. ἔστι
δὲ καὶ ὡς ἡ Α πρὸς τὴν Β, οὕτως ἡ Γ πρὸς τὴν Δ; δι ἴσου
ἄρα ἐστὶν ὡς ἡ Α πρὸς τὴν Ε, οὕτως ἡ Γ πρὸς τὴν Ζ.
εἴτε οὖν σύμμετρός ἐστιν ἡ Α τῇ Ε, συμμετρός ἐστι
καὶ ἡ Γ τῇ Ζ, εἴτε ἀσύμμετρός ἐστιν ἡ Α τῇ Ε, ἀσύμμετρός
ἐστι καὶ ἡ Γ τῇ Ζ.}

{\greekfont Ἐὰν ἄρα, καὶ τὰ ἑχῆς.}}

\ParallelRText{
\begin{center}
{\large Proposition 14}
\end{center}

If four straight-lines are proportional, and the
square on the first is greater  than (the square on) the second  by the (square) on (some straight-line) commensurable [in length] with the first, then the
square on the  third will also
be greater  than (the square on) the fourth by the (square) on (some straight-line) commensurable [in length]
with the third. And if the
square on the first is greater  than (the  square on) the second  by the (square) on (some straight-line) incommensurable [in length] with the first, then the square on the third will also
be greater than (the  square on) the fourth by the (square) on (some straight-line) incommensurable [in length]
with the third. 

Let $A$, $B$, $C$, $D$ be four proportional straight-lines, (such that)
as $A$ (is) to $B$, so $C$ (is) to $D$. And let the square on $A$ be greater 
than (the square on) $B$ by the (square) on $E$, and let the square on $C$ be greater
than (the square on) $D$ by the (square) on $F$. I say that 
$A$ is either commensurable (in length) with $E$, and $C$ is also commensurable with $F$, or $A$ is incommensurable (in length) with $E$, and $C$ is also incommensurable with $F$.

\epsfysize=2.in
\centerline{\epsffile{Book10/fig014e.eps}}

For since as $A$ is to $B$, so $C$ (is) to $D$, thus as the (square) on $A$
is to the (square) on $B$, so the (square) on $C$  (is) to the (square) on $D$
[Prop.~6.22]. But the (sum of the squares) on $E$ and $B$ is equal to the (square) on $A$, and the (sum of the squares) on 
$D$ and $F$ is equal to the (square) on $C$.
Thus, as the (sum of the squares) on $E$ and $B$ is to the (square) on
$B$, so the (sum of the squares) on $D$ and $F$ (is) to the (square)
on $D$.
 Thus, via separation, as the (square) on $E$ is to the (square) on $B$, so the (square) on $F$ (is) to the (square)
on $D$ [Prop.~5.17]. Thus, also, as $E$ is to $B$,
so $F$ (is) to $D$ [Prop.~6.22]. 
 Thus, inversely,  as $B$ is to $E$, so $D$ (is) to $F$ [Prop.~5.7~corr.]. But, as $A$ is to $B$, so 
 $C$ also (is) to $D$. Thus, via equality, as $A$ is to $E$, so $C$ (is) to $F$
 [Prop.~5.22]. Therefore, $A$ is either commensurable (in length)
 with $E$, and $C$ is also commensurable with $F$, or $A$ is
 incommensurable (in length) with $E$, and $C$ is also incommensurable with $F$ [Prop.~10.11].
 
 Thus, if, and so on \ldots.}
\end{Parallel}

%%%%%%
% Prop 10.15
%%%%%%
\pdfbookmark[1]{Proposition 10.15}{pdf10.15}
\begin{Parallel}{}{} 
\ParallelLText{
\begin{center}
{\large \ggn{15}.}
\end{center}\vspace*{-7pt}

{\greekfont Ἐὰν δύο μεγέθη σύμμετρα συντεθῇ, καὶ τὸ ὅλον ἑκατέρῳ α ὐτῶν
σύμμετρον ἔσται; κἂν τὸ ὅλον ἑνὶ α ὐτῶν σύμμετρον ᾖ, καὶ
τὰ ἐχ ἀρθῆς μεγέθη σύμμετρα ἔσται.}

{\greekfont Συγκείσθω γὰρ δύο μεγέθη σύμμετρα τὰ ΑΒ, ΒΓ;  λέγω, ὅτι καὶ
ὅλον τὸ ΑΓ ἑκατέρῳ τῶν ΑΒ, ΒΓ ἐστι σύμμετρον.}\\~\\~\\

\epsfysize=0.6in
\centerline{\epsffile{Book10/fig015g.eps}}

{\greekfont Ἐπεὶ γὰρ σύμμετρά ἐστι τὰ ΑΒ, ΒΓ, μετρήσει τι α ὐτὰ μέγεθος.
μετρείτω, καὶ ἔστω τὸ Δ. ἐπεὶ οὖν τὸ Δ τὰ ΑΒ, ΒΓ μετρεῖ,
καὶ ὅλον τὸ ΑΓ μετρήσει. μετρεῖ δὲ καὶ τὰ ΑΒ, ΒΓ. τὸ Δ ἄρα
τὰ ΑΒ, ΒΓ, ΑΓ μετρεῖ; σύμμετρον ἄρα ἐστὶ τὸ ΑΓ ἑκατέρῳ
τῶν ΑΒ, ΒΓ.}

{\greekfont Ἀλλὰ δὴ τὸ ΑΓ ἔστω σύμμετρον τῷ ΑΒ; λέγω δή, ὅτι καὶ τὰ
ΑΒ, ΒΓ σύμμετρά ἐστιν.}

{\greekfont Ἐπεὶ γὰρ σύμμετρά ἐστι τὰ ΑΓ, ΑΒ, μετρήσει τι α ὐτὰ μέγεθος.
μετρείτω, καὶ ἔστω τὸ Δ. ἐπεὶ οὖν τὸ Δ τὰ ΓΑ, ΑΒ
μετρεῖ, καὶ λοιπὸν ἄρα τὸ ΒΓ μετρήσει. μετρεῖ δὲ καὶ τὸ ΑΒ;
τὸ Δ ἄρα τὰ ΑΒ, ΒΓ μετρήσει; σύμμετρα ἄρα ἐστὶ τὰ ΑΒ, ΒΓ.}

{\greekfont Ἐὰν ἄρα δύο μεγέθη, καὶ τὰ ἑχῆς.}}

\ParallelRText{
\begin{center}
{\large Proposition 15}
\end{center}

If two commensurable magnitudes are
added together then the whole will also be commensurable
with each of them. And if the whole is commensurable with one of
them then the original magnitudes will also be commensurable (with one another).

For let the two commensurable magnitudes $AB$ and $BC$ be laid
down together. I say that the whole $AC$ is also commensurable with
each of $AB$ and $BC$.

\epsfysize=0.6in
\centerline{\epsffile{Book10/fig015e.eps}}

For since $AB$ and $BC$ are commensurable, some  magnitude will
measure them. Let it (so) measure (them), and let it be $D$. Therefore,
since $D$ measures (both) $AB$ and $BC$, it will also measure the
whole $AC$. And it also measures $AB$ and $BC$. Thus, $D$
measures  $AB$, $BC$, and $AC$. Thus, $AC$ is commensurable
with each of $AB$ and $BC$ [Def.~10.1].

And so let $AC$ be commensurable with $AB$. I say that $AB$ and
$BC$ are also commensurable.

For since $AC$ and $AB$ are commensurable, some  magnitude
will measure them. Let it (so) measure (them), and let it be $D$. Therefore,
since $D$ measures (both) $CA$ and $AB$, it will thus also
measure the remainder $BC$. And it also measures $AB$. Thus, $D$ will
measure (both) $AB$ and $BC$. Thus, $AB$ and $BC$ are commensurable [Def.~10.1].

Thus, if two magnitudes, and so on \ldots.}
\end{Parallel}

%%%%%%
% Prop 10.16
%%%%%%
\pdfbookmark[1]{Proposition 10.16}{pdf10.16}
\begin{Parallel}{}{} 
\ParallelLText{
\begin{center}
{\large \ggn{16}.}
\end{center}\vspace*{-7pt}

{\greekfont Ἐὰν δύο μεγέθη ἀσύμμετρα συντεθῇ, καὶ τὸ ὅλον ἑκατέρῳ
α ὐτῶν ἀσύμμετρον ἔσται; κἂν τὸ ὅλον ἑνὶ α ὐτῶν
ἀσύμμετρον ᾖ, καὶ τὰ ἐχ ἀρθῆς μεγέθη ἀσύμμετρα ἔσται.}\\~\\

\epsfysize=0.6in
\centerline{\epsffile{Book10/fig015g.eps}}

{\greekfont Συγκείσθω γὰρ δύο μεγέθη ἀσύμμετρα τὰ ΑΒ, ΒΓ; λέγω, ὅτι
καὶ ὅλον τὸ ΑΓ ἑκατέρῳ τῶν ΑΒ, ΒΓ ἀσύμμετρόν ἐστιν.}

{\greekfont Εἰ γὰρ μή ἐστιν ἀσύμμετρα τὰ ΓΑ, ΑΒ, μετρήσει τι [α ὐτὰ]
μέγεθος. μετρείτω, εἰ δυνατόν, καὶ ἔστω τὸ Δ. ἐπεὶ οὖν
τὸ Δ τὰ ΓΑ, ΑΒ μετρεῖ, καὶ λοιπὸν ἄρα τὸ ΒΓ μετρήσει. μετρεῖ
δὲ καὶ τὸ ΑΒ; τὸ Δ ἄρα τὰ ΑΒ, ΒΓ μετρεῖ. σύμμετρα ἄρα ἐστὶ
τὰ ΑΒ, ΒΓ; ὑπέκειντο δὲ καὶ ἀσύμμετρα; ὅπερ ἐστὶν ἀδύνατον.
οὐκ ἄρα τὰ ΓΑ, ΑΒ μετρήσει τι μέγεθος; ἀσύμμετρα ἄρα
ἐστὶ τὰ ΓΑ, ΑΒ. ὁμοίως δὴ δείχομεν,  ὅτι καὶ τὰ ΑΓ, ΓΒ
ἀσύμμετρά ἐστιν. τὸ ΑΓ ἄρα ἑκατέρῳ τῶν ΑΒ, ΒΓ ἀσύμμετρόν
ἐστιν.}

{\greekfont Ἀλλὰ δὴ τὸ ΑΓ ἑνὶ τῶν ΑΒ, ΒΓ ἀσύμμετρον ἔστω. ἔστω δὴ
πρότερον τῷ ΑΒ; λέγω, ὅτι καὶ τὰ ΑΒ, ΒΓ ἀσύμμετρά ἐστιν.
εἰ γὰρ ἔσται σύμμετρα, μετρήσει τι α ὐτὰ μέγεθος. μετρείτω,
καὶ ἔστω τὸ Δ. ἐπεὶ οὖν τὸ Δ τὰ ΑΒ, ΒΓ μετρεῖ, καὶ
ὅλον ἄρα τὸ ΑΓ μετρήσει. μετρεῖ δὲ καὶ τὸ ΑΒ; τὸ Δ ἄρα
τὰ ΓΑ, ΑΒ μετρεῖ. σύμμετρα ἄρα ἐστὶ τὰ ΓΑ, ΑΒ; ὑπέκειτο
δὲ καὶ ἀσύμμετρα; ὅπερ ἐστὶν ἀδύνατον. οὐκ ἄρα τὰ ΑΒ, ΒΓ
μετρήσει τι μέγεθος; ἀσύμμετρα ἄρα ἐστὶ τὰ ΑΒ, ΒΓ.}

{\greekfont Ἐὰν ἄρα δύο μεγέθη, καὶ τὰ ἑχῆς.}}

\ParallelRText{
\begin{center}
{\large Proposition 16}
\end{center}

If two incommensurable magnitudes are
added together then the whole will also be incommensurable with
each of them. And if the whole is incommensurable with one of them then
the original magnitudes will also be incommensurable (with one another).

\epsfysize=0.6in
\centerline{\epsffile{Book10/fig015e.eps}}

For let the two incommensurable magnitudes $AB$ and $BC$ be
laid down together. I say that that the whole $AC$ is also
incommensurable with each of $AB$ and $BC$.

For if $CA$ and $AB$ are not incommensurable then some magnitude
will measure [them]. If possible, let it (so) measure (them), and let it
be $D$. Therefore, since $D$ measures (both) $CA$ and $AB$, it will
thus also measure the remainder $BC$. And it also measures $AB$.
Thus, $D$ measures (both) $AB$ and $BC$. Thus, $AB$ and $BC$ are
commensurable [Def.~10.1]. But they were also assumed (to be) incommensurable.
The very thing is impossible. Thus, some magnitude cannot
measure (both) $CA$ and $AB$. Thus, $CA$ and $AB$ are incommensurable [Def.~10.1]. So, similarly, we can show that $AC$ and
$CB$ are also incommensurable. Thus, $AC$ is incommensurable with each
of $AB$ and $BC$.

And so let $AC$ be incommensurable with one of $AB$ and $BC$. So
let it, first of all, be incommensurable with $AB$. I say that $AB$ and
$BC$ are also incommensurable. For if they are commensurable then
some magnitude will measure them. Let it (so) measure (them), and let it
be $D$. Therefore, since $D$ measures (both) $AB$ and $BC$, 
it will thus also measure the whole  $AC$. And it also measures $AB$. 
Thus, $D$ measures (both) $CA$ and $AB$. Thus, $CA$ and $AB$
are commensurable  [Def.~10.1]. But they
were also assumed (to be) incommensurable. The very thing is
impossible. Thus, some magnitude cannot measure (both) $AB$ and
$BC$. Thus, $AB$ and $BC$ are incommensurable  [Def.~10.1].

Thus, if two\ldots magnitudes, and so on \ldots.}
\end{Parallel}

%%%%%%
% Prop 10.16a
%%%%%%
\begin{Parallel}{}{} 
\ParallelLText{
\begin{center}
{\large {\greekfont Λῆμμα}.}
\end{center}\vspace*{-7pt}

{\greekfont Ἐὰν παρά τινα εὐθεῖαν παραβληθῇ παραλληλόγραμ-μον ἐλλεῖπον
εἴδει τετραγώνῳ, τὸ παραβληθὲν ἴσον ἐστὶ τῷ ὑπὸ τῶν
ἐκ τῆς παραβολής γενομένων τμ ημάτων τῆς εὐθείας.}\\

\epsfysize=1.3in
\centerline{\epsffile{Book10/fig016ag.eps}}

{\greekfont Παρὰ γὰρ εὐθεῖαν τὴν ΑΒ παραβεβλήσθω παραλληλόγραμμον τὸ
ΑΔ ἐλλεῖπον εἴδει τετραγώνῳ τῷ ΔΒ;
λέγω, ὅτι ἴσον ἐστὶ τὸ ΑΔ τῷ ὑπὸ τῶν ΑΓ, ΓΒ.}

{\greekfont Καί ἐστιν α ὐτόθεν φανερόν; ἐπεὶ γὰρ τετράγωνόν ἐστι τὸ ΔΒ,
ἴση ἐστὶν ἡ ΔΓ τῇ ΓΒ, καί ἐστι τὸ ΑΔ τὸ ὑπὸ τῶν
ΑΓ, ΓΔ, τουτέστι τὸ ὑπὸ τῶν ΑΓ, ΓΒ.}

{\greekfont Ἐὰν ἄρα παρά τινα εὐθεῖαν, καὶ τὰ ἑχῆς.}}

\ParallelRText{
\begin{center}
{\large Lemma}
\end{center}

If a parallelogram,$^\dag$
 falling short by a square figure, is applied to some straight-line then the applied (parallelogram) is equal (in area) to the (rectangle contained) by the
pieces of the straight-line created via the application (of the parallelogram).

\epsfysize=1.3in
\centerline{\epsffile{Book10/fig016ae.eps}}

For let the parallelogram $AD$, falling short by the square figure $DB$, be applied to the straight-line $AB$. I say that $AD$ is equal to
the (rectangle contained) by $AC$ and $CB$.

And it is immediately obvious. For since $DB$ is a square, $DC$ is
equal to $CB$. And $AD$ is the (rectangle contained) by $AC$ and
$CD$---that is to say, by $AC$ and $CB$.

Thus, if \ldots to some straight-line, and so on \ldots.}
\end{Parallel}
{\footnotesize\noindent$^\dag$  Note that this lemma only applies to rectangular parallelograms.}

%%%%%%
% Prop 10.17
%%%%%%
\pdfbookmark[1]{Proposition 10.17}{pdf10.17}
\begin{Parallel}{}{} 
\ParallelLText{
\begin{center}
{\large \ggn{17}.}
\end{center}\vspace*{-7pt}

{\greekfont Ἐὰν ὦσι δύο εὐθεῖαι ἄνισοι, τῷ δὲ τετράτῳ μέρει τοῦ ἀπὸ
τῆς ἐλάσσονος ἴσον παρὰ τὴν μείζονα παραβληθῇ ἐλλεῖπον
εἴδει τετραγώνῳ καὶ εἰς σύμμετρα α ὐτὴν διαιρῇ μήκει, ἡ
μείζων τῆς ἐλάσσονος μεῖζον δυνήσεται τῷ ἀπὸ συμμέτου
ἑαυτῇ [μήκει]. καὶ ἐὰν ἡ μείζων τῆς ἐλάσσονος μεῖζον
δύνηται τῷ ἀπὸ συμμέτρου ἑαυτῇ [μήκει], τῷ δὲ τετράρτῳ τοῦ
ἀπὸ τῆς ἐλάσσονος ἴσον παρὰ τὴν μείζονα παραβληθῇ ἐλλεῖπον
εἴδει τετραγώνῳ, εἰς σύμμετρα α ὐτὴν διαιρεῖ μήκει.}

{\greekfont Ἔστωσαν δύο εὐθεῖαι ἄνισοι αἱ Α, ΒΓ, ὧν μείζων
ἡ ΒΓ, τῷ δὲ τετράρτῳ μέρει τοῦ ἀπὸ ἐλάσσονος
τῆς Α, τουτέστι τῷ ἀπὸ τῆς ἡμισείας τῆς Α, ἴσον παρὰ
τὴν ΒΓ παραβεβλήσθω ἐλλεῖπον εἴδει τετραγώνῳ, καὶ ἔστω
τὸ ὑπὸ τῶν ΒΔ, ΔΓ, σύμμετρος δὲ ἔστω ἡ ΒΔ τῇ ΔΓ
μήκει; λέγω, ὃτι ἡ ΒΓ τῆς Α μεῖζον δύναται τῷ ἀπὸ
συμμέτρου ἑαυτῇ.}\\~\\~\\~\\~\\~\\~\\~\\

\epsfysize=1.1in
\centerline{\epsffile{Book10/fig017g.eps}}

{\greekfont Τετμήσθω γὰρ ἡ ΒΓ δίθα κατὰ τὸ Ε σημεῖον, καὶ κείσθω τῇ ΔΕ
ἴση ἡ ΕΖ. λοιπὴ ἄρα ἡ ΔΓ ἴση ἐστὶ τῇ ΒΖ. καὶ ἐπεὶ εὐθεῖα
ἡ ΒΓ τέτμ ηται εἰς μὲν ἴσα κατὰ τὸ Ε, εἰς δὲ ἄνισα κατὰ τὸ Δ,
τὸ ἄρα ὑπὸ ΒΔ, ΔΓ περειθόμενον ὀρθογώνιον μετὰ τοῦ ἀπὸ
τῆς ΕΔ τετραγώνου ἴσον ἐστὶ τῷ ἀπὸ τῆς ΕΓ τετραγώνῳ;
καὶ τὰ τετραπλάσια; τὸ ἄρα τετράκις ὑπὸ τῶν ΒΔ, ΔΓ μετὰ τοῦ
τετραπλασίου τοῦ ἀπὸ τῆς ΔΕ ἴσον ἐστὶ τῷ τετράκις ἀπὸ
τῆς ΕΓ τετραγώνῳ.
ἀλλὰ τῷ μὲν τετραπλασίῳ τοῦ ὑπὸ τῶν ΒΔ, ΔΓ ἴσον ἐστὶ
τὸ ἀπὸ τῆς Α τετράγωνον, τῷ δὲ τετραπλασίῳ τοῦ ἀπὸ τῆς
ΔΕ ἴσον ἐστὶ τὸ ἀπὸ τῆς ΔΖ τετράγωνον; διπλασίων γάρ
ἐστιν ἡ ΔΖ τῆς ΔΕ. τῷ δὲ τετραπλασίῳ τοῦ ἀπὸ τῆς ΕΓ ἴσον
ἐστὶ τὸ ἀπὸ τῆς ΒΓ τετράγωνον; διπλασίων γάρ ἐστι πάλιν
ἡ ΒΓ τῆς ΓΕ. τὰ ἄρα ἀπὸ τῶν Α, ΔΖ τετράγωνα ἴσα ἐστὶ
τῷ ἀπὸ τῆς ΒΓ τετράγωνῳ; ὥστε τὸ ἀπὸ τῆς ΒΓ τοῦ ἀπὸ
τῆς Α μεῖζόν ἐστι τῷ ἀπὸ τῆς ΔΖ; ἡ ΒΓ ἄρα τῆς Α μεῖζον
δύναται τῇ ΔΖ. δεικτέον, ὅτι καὶ σύμμετρός ἐστιν ἡ ΒΓ τῇ
ΔΖ. ἐπεὶ γὰρ σύμμετρός ἐστιν ἡ ΒΔ τῇ
ΔΓ μήκει, σύμμετρος ἄρα ἐστὶ καὶ ἡ ΒΓ τῇ ΓΔ μήκει. ἀλλὰ
ἡ ΓΔ ταῖς ΓΔ, ΒΖ ἐστι σύμμετρος μήκει; ἴση γάρ
ἐστιν ἡ ΓΔ τῇ ΒΖ. καὶ ἡ ΒΓ ἄρα σύμμετρός ἐστι ταῖς ΒΖ, ΓΔ
μήκει; ὥστε καὶ λοιπῇ τῇ ΖΔ σύμμετρός ἐστιν ἡ ΒΓ μήκει;
ἡ ΒΓ ἄρα τῆς Α μεῖζον δύναται τῷ ἀπὸ
συμμέτρου ἑαυτῇ.}

{\greekfont Ἀλλὰ δὴ ἡ ΒΓ τῆς Α μεῖζον δυνάσθω τῷ ἀπὸ συμμέτρου
ἑαυτῇ, τῷ δὲ τετράτρῳ τοῦ ἀπὸ τῆς Α ἴσον παρὰ τὴν ΒΓ
παραβεβλήσθω ἐλλεῖπον εἴδει τετραγώνῳ, καὶ ἔστω τὸ ὑπὸ
τῶν ΒΔ, ΔΓ. δεικτέον, ὅτι σύμμετρός ἐστιν ἡ ΒΔ τῇ ΔΓ
μήκει.}

{\greekfont Τῶν γὰρ α ὐτῶν κατασκευασθέντων ὁμοίως δείχομεν, ὅτι
ἡ ΒΓ τῆς Α μεῖζον δύναται τῷ ἀπὸ τῆς ΖΔ.
δύναται δὲ ἡ ΒΓ τῆς Α μεῖζον τῷ ἀπὸ συμμέτρου ἑαυτῇ.
σύμμετρος ἄρα ἐστὶν ἡ ΒΓ τῇ ΖΔ μήκει;
ὥστε καὶ λοιπῇ συναμφοτέρῳ τῇ ΒΖ, ΔΓ σύμμετρός
ἐστιν ἡ ΒΓ μήκει. ἀλλὰ συναμφότερος ἡ ΒΖ, ΔΓ σύμμετρός
ἐστι τῇ ΔΓ [μήκει]. ὥστε καὶ ἡ ΒΓ τῇ ΓΔ σύμμετρός
ἐστι μήκει; καὶ διελόντι ἄρα ἡ ΒΔ τῇ ΔΓ ἐστι σύμμετρος
μήκει.}

{\greekfont Ἐὰν ἄρα ὦσι δύο εὐθεῖαι ἄνισοι, καὶ τὰ ἑχῆς.}}

\ParallelRText{
\begin{center}
{\large Proposition 17}$^\dag$
\end{center}

If there are two unequal straight-lines,
and a (rectangle) equal to the fourth part of the (square) on the
lesser, falling short by a square figure, is applied to the greater, and
divides it into (parts which are) commensurable in length, then
 the  square on the greater  will be  larger than (the  square on) the lesser  by the (square)
on (some straight-line) commensurable [in length] with the greater.
And if the square on the greater is larger than (the square on) the lesser  by the (square)
on (some straight-line) commensurable [in length] with the greater, and
a (rectangle) equal to the fourth (part) of the (square) on the lesser,
falling short by a square figure, is applied to the greater, then it divides it into (parts
which are) commensurable in length.

Let $A$ and $BC$ be two unequal straight-lines, of which (let) $BC$ (be) the
greater. And let a (rectangle) equal to the fourth part of the (square) on the
lesser, $A$---that is, (equal) to the (square) on half of $A$---falling short by a
square figure, be applied to $BC$. And let it be the (rectangle contained) by $BD$ and $DC$ [see previous lemma]. And let $BD$ be commensurable in length with
$DC$. I say that that the square on $BC$ is greater than the (square on) $A$ by (the square on some straight-line) commensurable (in length) with ($BC$).

\epsfysize=1.1in
\centerline{\epsffile{Book10/fig017e.eps}}

For let $BC$ be cut in  half at the point $E$ [Prop. 1.10]. And let $EF$ be made equal to
$DE$ [Prop.~1.3]. Thus, the remainder $DC$
is equal to $BF$. And since the straight-line $BC$ has been cut into equal
(pieces) at $E$, and into unequal (pieces) at $D$, the rectangle contained by
$BD$ and $DC$, plus the square on $ED$, is thus equal to the square on $EC$ [Prop.~2.5]. (The same) also (for) the quadruples.  
Thus, four times the (rectangle contained) by $BD$ and $DC$, plus
the quadruple of the (square) on $DE$, is equal to four times the square on 
$EC$. But, the square on $A$ is equal to the quadruple of the (rectangle
contained) by $BD$ and $DC$, and the square on $DF$ is equal to
the quadruple of the (square) on $DE$. For $DF$ is double $DE$. 
And the square on  $BC$ is equal to the quadruple of the (square) on $EC$.
For, again, $BC$ is double $CE$. Thus, the (sum of the) squares on $A$ and
$DF$ is equal to the square on $BC$. Hence, the (square) on $BC$
is greater than the (square) on $A$ by the (square) on $DF$. Thus,
$BC$ is greater in square than $A$ by $DF$. It must also be shown that $BC$
is commensurable (in length) with $DF$. For since $BD$ is commensurable in length
with $DC$, $BC$ is thus also commensurable in length with $CD$ [Prop.~10.15]. 
But, $CD$ is commensurable in length with $CD$ plus $BF$. For $CD$ is equal
to $BF$ [Prop.~10.6].  Thus, $BC$ is also commensurable in length with $BF$ plus $CD$
[Prop.~10.12]. Hence, $BC$ is also commensurable
in length with the remainder $FD$ [Prop.~10.15]. Thus, the square on $BC$ is
greater than (the square on) $A$ by the (square) on (some straight-line)
commensurable (in length) with ($BC$).

And so let the square on $BC$ be greater than the (square on) $A$ by the
(square) on (some straight-line) commensurable (in length) with ($BC$). And let
a (rectangle) equal to the fourth (part) of the (square) on $A$, falling
short by a square figure, be applied to $BC$. And let it be
the (rectangle contained) by $BD$ and $DC$. It must be shown that $BD$
is commensurable in length with $DC$.

For, similarly, by the same construction, we can  show that the
square on $BC$ is greater than the (square on) $A$ by the (square) on $FD$.
And the square on $BC$ is greater than the (square on) $A$ by the
(square) on (some straight-line) commensurable (in length) with ($BC$).  Thus, $BC$ is
commensurable in length with $FD$. Hence, $BC$ is also commensurable in length with the remaining sum of $BF$ and $DC$ [Prop.~10.15]. But, the sum of $BF$ and $DC$
is commensurable [in length] with $DC$ [Prop.~10.6]. Hence, $BC$ is also
commensurable in length with $CD$ [Prop.~10.12]. 
Thus, via separation, $BD$ is also commensurable in length with $DC$ [Prop.~10.15].

Thus,  if there are two unequal straight-lines, and so on \ldots.}
\end{Parallel}
{\footnotesize\noindent$^\dag$ This proposition states that if $\alpha\,x-x^2=\beta^2/4$ (where $\alpha=BC$, $x=DC$,
and $\beta=A$) then $\alpha$ and $\sqrt{\alpha^2-\beta^2}$ are commensurable when $\alpha-x$ are $x$ are commensurable, and {\rm vice versa}.}

%%%%%%
% Prop 10.18
%%%%%%
\pdfbookmark[1]{Proposition 10.18}{pdf10.18}
\begin{Parallel}{}{} 
\ParallelLText{
\begin{center}
{\large\ggn{18}.}
\end{center}\vspace*{-7pt}

{\greekfont Ἐὰν ὦσι δύο εὐθεῖαι ἄνισοι, τῷ δὲ τετάρτῳ μέρει τοῦ ἀπὸ
τῆς ἐλάσσονος ἴσον παρὰ τὴν μείζονα παραβληθῇ ἐλλεῖπον εἴδει
τετραγώνῳ, καὶ εἰς ἀσυμμετρα α ὐτὴν διαιρῇ [μήκει], ἡ μείζων
τῆς ἐλάσσονος μεῖζον δυνήσεται τῷ ἀπὸ ἀσυμμέτρου ἑαυτῇ.
καὶ ἐὰν ἡ μείζων τῆς ἐλάσσονος μεῖζον δύνηται τῷ ἀπὸ
ἀσυμμέτρου ἑαυτῇ, τῷ δὲ τετράρτῳ τοῦ ἀπὸ τῆς ἐλάσσονος
ἴσον παρὰ τὴν μείζονα παραβληθῇ ἐλλεῖπον εἴδει τετραγώνῳ,
εἰς ἀσύμμετρα α ὐτὴν διαιρεῖ [μήκει].}

{\greekfont Ἔστωσαν δύο εὐθεῖαι ἄνισοι αἱ Α, ΒΓ, ὧν μείζων ἡ ΒΓ, τῷ
δὲ τετάρτῳ [μέρει] τοῦ ἀπὸ τῆς ἐλάσσονος τῆς Α ἴσον παρὰ
τὴν ΒΓ παραβεβλήσθω ἐλλεῖπον εἴδει τετραγώνῳ, καὶ ἔστω
τὸ ὑπὸ τῶν ΒΔΓ, ἀσύμμετρος δὲ ἔστω ἡ ΒΔ τῇ ΔΓ μήκει;
λέγω, ὅτι ἡ ΒΓ τῆς Α μεῖζον δύναται τῷ ἀπὸ ἀσυμμέτρου
ἑαυτῇ.}\\~\\~\\~\\~\\~\\~\\

\epsfysize=0.8in
\centerline{\epsffile{Book10/fig018g.eps}}

{\greekfont Τῶν γὰρ α ὐτῶν κατασκευασθέντων τῷ πρότερον ὁμοίως δείχομεν,
ὅτι ἡ ΒΓ τῆς Α μεῖζον δύναται τῷ ἀπὸ τῆς ΖΔ. δεικτέον
[οὖν], ὅτι ἀσύμμετρός ἐστιν ἡ ΒΓ τῇ ΔΖ μήκει. ἐπεὶ γὰρ
ἀσύμμετρός ἐστιν ἡ  ΒΔ τῇ ΔΓ μήκει, ἀσύμμετρος ἄρα ἐστὶ καὶ
ἡ ΒΓ τῇ ΓΔ μήκει. ἀλλὰ ἡ ΔΓ σύμμετρός ἐστι συναμφοτέραις
ταῖς ΒΖ, ΔΓ; καὶ ἡ ΒΓ ἄρα ἀσύμμετρός ἐστι συναμφοτέραις
ταῖς ΒΖ, ΔΓ. ὥστε καὶ λοιπῇ τῇ ΖΔ ἀσύμμετρός  ἔστιν ἡ ΒΓ
μήκει. καὶ ἡ ΒΓ τῆς Α μεῖζον δύναται τῷ ἀπὸ τῆς ΖΔ; ἡ ΒΓ
ἄρα τῆς Α μεῖζον δύναται τῷ ἀπὸ ἀσυμμέτρου ἑαυτῇ.}

{\greekfont Δυνάσθω δὴ πάλιν ἡ ΒΓ τῆς Α μεῖζον τῷ ἀπὸ ἀσυμμέτρ-ου
ἑαυτῇ, τῷ δὲ τετάρτῳ τοῦ ἀπὸ τῆς Α ἴσον παρὰ τὴν
ΒΓ παραβεβλήσθω ἐλλεῖπον εἴδει τετραγώνῳ, καὶ ἔστω τὸ ὑπὸ τῶν
ΒΔ, ΔΓ. δεικτέον, ὅτι ἀσύμμετρός ἐστιν ἡ ΒΔ τῇ ΔΓ μήκει.}

{\greekfont Τῶν γὰρ α ὐτῶν κατασκευασθέντων ὁμοίως δείχομεν, ὅτι
ἡ ΒΓ τῆς Α μεῖζον δύναται τῷ ἀπὸ τῆς ΖΔ. ἀλλὰ ἡ ΒΓ
τῆς Α μεῖζον δύναται τῷ ἀπὸ ἀσυμμέτρου ἑαυτῇ. ἀσύμμετρος
ἄρα ἐστὶν ἡ ΒΓ τῇ ΖΔ μήκει; ὥστε καὶ λοιπῇ συναμφοτέρῳ
τῇ ΒΖ, ΔΓ ἀσύμμετρός ἐστιν ἡ ΒΓ. ἀλλὰ συναμφότερος ἡ ΒΖ,
ΔΓ τῇ ΔΓ σύμμετρός ἐστι μήκει; καὶ ἡ ΒΓ ἄρα τῇ ΔΓ
ἀσύμμετρός ἐστι μήκει; ὥστε καὶ διελόντι ἡ ΒΔ τῇ ΔΓ
ἀσύμμετρός ἐστι μήκει.}

{\greekfont Ἐὰν ἄρα ὦσι δύο εὐθεῖαι, καὶ τὰ ἑχῆς.}}

\ParallelRText{
\begin{center}
{\large Proposition 18}$^\dag$
\end{center}

If there are two unequal straight-lines,
and a (rectangle) equal to the fourth part of the (square) on the
lesser, falling short by a square figure, is applied to the greater, and
divides it into (parts which are) incommensurable [in length], then
 the  square on the greater  will be  larger than the  (square on the) lesser  by the (square)
on (some straight-line) incommensurable  (in length) with the greater.
And if the square on the greater is larger than the (square on the) lesser  by the (square)
on (some straight-line) incommensurable  (in length) with the greater, and
a (rectangle) equal to the fourth (part) of the (square) on the lesser,
falling short by a square figure, is applied to the greater, then it divides it into (parts
which are) incommensurable [in length].

Let $A$ and $BC$ be two unequal straight-lines, of which (let) $BC$ (be) the
greater. And let a (rectangle) equal to the fourth [part] of the (square) on the
lesser, $A$, falling short by a
square figure, be applied to $BC$. And let it be the (rectangle contained) by $BDC$. And let $BD$ be incommensurable in length with
$DC$. I say that that the square on $BC$ is greater than the (square on) $A$ by the (square) on (some straight-line) incommensurable (in length) with ($BC$).

\epsfysize=0.8in
\centerline{\epsffile{Book10/fig018e.eps}}

For, similarly, by the same construction as before, we can show that
the square on $BC$ is greater than the (square on) $A$ by the (square)
on $FD$. [Therefore] it must be shown that $BC$ is incommensurable
in length
with $DF$. For since $BD$ is incommensurable in length with $DC$, $BC$
is thus also incommensurable in length with $CD$ [Prop.~10.16]. But, $DC$ is commensurable (in length) with the sum of 
$BF$ and $DC$ [Prop.~10.6]. And, thus,
$BC$ is incommensurable (in  length) with the sum of $BF$ and $DC$ [Prop.~10.13]. Hence, $BC$ is also incommensurable in length with the remainder $FD$ [Prop.~10.16]. And the square on $BC$ is greater than the (square on) $A$ by the (square) on $FD$. Thus, the square on $BC$
is greater than the (square on) $A$ by the (square) on (some straight-line) incommensurable
(in length) 
with ($BC$).

So, again, let the square on $BC$ be greater than the (square on) $A$ by
the (square) on (some straight-line) incommensurable (in length) with ($BC$). And let a (rectangle) equal to the fourth [part] of the (square) on  $A$, falling short by a
square figure, be applied to $BC$.  And let it be the (rectangle contained) by $BD$ and $DC$. It must be shown that $BD$ is incommensurable in length with $DC$.

For, similarly, by the same construction, we can show that the square on $BC$ is greater than the (square) on $A$ by the (square) on $FD$. But,
the square on $BC$ is greater than the (square) on $A$ by the (square) on (some straight-line)
incommensurable (in length) with ($BC$). Thus, $BC$ is incommensurable in length with
$FD$. Hence, $BC$ is also incommensurable (in length) with the remaining sum of $BF$ and
$DC$ [Prop.~10.16]. But, the sum of
$BF$ and $DC$ is commensurable in length with $DC$ [Prop.~10.6]. Thus, $BC$ is also incommensurable
in length with $DC$ [Prop.~10.13]. Hence, via
separation, $BD$ is
also incommensurable in length with $DC$ [Prop.~10.16].

Thus,  if there are two \ldots straight-lines, and so on \ldots.}
\end{Parallel}
{\footnotesize\noindent$^\dag$ This proposition states that if $\alpha\,x-x^2=\beta^2/4$ (where $\alpha=BC$, $x=DC$,
and $\beta=A$) then $\alpha$ and $\sqrt{\alpha^2-\beta^2}$ are incommensurable when $\alpha-x$ are $x$ are incommensurable, and {\rm vice versa}.}

%%%%%%
% Prop 10.19
%%%%%%
\pdfbookmark[1]{Proposition 10.19}{pdf10.19}
\begin{Parallel}{}{} 
\ParallelLText{
\begin{center}
{\large \ggn{19}.}
\end{center}\vspace*{-7pt}

{\greekfont Τὸ ὑπὸ ῥητῶν μήκει συμμέτρων  εὐθειῶν περιεχόμενον ὀρθογώνιον ῥητόν ἐστιν.}

{\greekfont Ὑπὸ γὰρ ῥητῶν μήκει συμμέτρων εὐθειῶν τῶν
ΑΒ, ΒΓ ὀρθογώνιον περιεθέσθω τὸ ΑΓ; λέγω, ὅτι ῥητόν ἐστι
τὸ ΑΓ.}\\

\epsfysize=2in
\centerline{\epsffile{Book10/fig019g.eps}}

{\greekfont Ἀναγεγράφθω γὰρ ἀπὸ τῆς ΑΒ τετράγωνον τὸ ΑΔ;
ῥητὸν ἄρα ἐστὶ τὸ ΑΔ. καὶ ἐπεὶ σύμμετρός ἐστιν ἡ
ΑΒ τῇ ΒΓ μήκει, ἴση δέ ἐστιν ἡ ΑΒ τῇ ΒΔ, σύμμετρος
ἄρα ἐστὶν ἡ ΒΔ τῇ ΒΓ μήκει. καί ἐστιν ὡς ἡ ΒΔ πρὸς τὴν
 ΒΓ, οὕτως τὸ ΔΑ πρὸς τὸ ΑΓ. σύμμετρον
ἄρα ἐστὶ τὸ ΔΑ τῷ ΑΓ. ῥητὸν δὲ τὸ ΔΑ; ῥητὸν ἄρα ἐστὶ
καὶ τὸ ΑΓ.}

{\greekfont Τὸ ἄρα ὑπὸ ῥητῶν μήκει συμμέτρων, καὶ τὰ ἑχῆς.}}

\ParallelRText{
\begin{center}
{\large Proposition 19}
\end{center}

The rectangle contained by rational straight-lines
(which are) commensurable in length is rational.

For let the rectangle $AC$ be enclosed by the rational straight-lines
$AB$ and $BC$ (which are) commensurable in length. I say that
$AC$ is rational.

\epsfysize=2in
\centerline{\epsffile{Book10/fig019e.eps}}

For let the square $AD$ be described on $AB$.  $AD$ is thus rational [Def.~10.4]. And since $AB$ is commensurable in length with $BC$, and $AB$ is equal to $BD$, $BD$ is thus
commensurable in  length with $BC$.  And as $BD$ is to
$BC$, so $DA$ (is) to $AC$ [Prop.~6.1]. Thus,
$DA$ is commensurable with $AC$ [Prop.~10.11]. 
And $DA$ (is) rational. Thus, $AC$ is also rational [Def.~10.4].
Thus,  the \ldots by rational straight-lines
\ldots commensurable, and so on \ldots.}
\end{Parallel}

%%%%%%
% Prop 10.20
%%%%%%
\pdfbookmark[1]{Proposition 10.20}{pdf10.20}
\begin{Parallel}{}{} 
\ParallelLText{
\begin{center}
{\large \ggn{20}.}
\end{center}\vspace*{-7pt}

{\greekfont Ἐὰν ῥητὸν παρὰ ῥητὴν παραβληθῇ, πλάτος ποιεῖ ῥητὴν καὶ σύμμετρον τῇ, παρ ἣν παράκειται, μήκει.}\\~\\

\epsfysize=2.4in
\centerline{\epsffile{Book10/fig020g.eps}}

{\greekfont Ῥητὸν γὰρ τὸ ΑΓ παρὰ ῥητὴν τὴν ΑΒ παραβεβλήσθω πλάτος ποιοῦν τὴν ΒΓ; λέγω, ὅτι
ῥητή ἐστιν ἡ ΒΓ καὶ σύμμετρος τῇ ΒΑ μήκει.}

{\greekfont Ἀναγεγράφθω γὰρ ἀπὸ τῆς ΑΒ τετράγωνον τὸ ΑΔ;
ῥητὸν ἄρα ἐστὶ τὸ ΑΔ. ῥητὸν δὲ καὶ τὸ ΑΓ; σύμμετρον ἄρα
ἐστὶ τὸ ΔΑ τῷ ΑΓ. καί ἐστιν ὡς τὸ ΔΑ πρὸς τὸ ΑΓ, οὕτως
ἡ ΔΒ πρὸς τὴν ΒΓ. σύμμετρος ἄρα ἐστὶ καὶ ἡ ΔΒ τῇ ΒΓ;
ἴση δὲ ἡ ΔΒ τῇ ΒΑ; σύμμετρος ἄρα καὶ ἡ
ΑΒ τῇ ΒΓ.
ῥητὴ δέ ἐστιν ἡ ΑΒ; ῥητὴ ἄρα ἐστὶ καὶ ἡ ΒΓ καὶ
σύμμετρος τῇ ΑΒ μήκει.}

{\greekfont Ἐὰν ἄρα ῥητὸν παρὰ ῥητὴν παραβληθῇ, καὶ τὰ ἑχῆς.}}

\ParallelRText{
\begin{center}
{\large Proposition 20}
\end{center}

If a rational (area) is applied to a rational (straight-line) then it produces as breadth  a (straight-line which is) rational, and commensurable in length
with the (straight-line) to which it is applied.

\epsfysize=2.4in
\centerline{\epsffile{Book10/fig020e.eps}}

For let the rational (area) $AC$ be applied to the rational (straight-line) $AB$, producing the (straight-line) $BC$ as breadth. I say that
$BC$ is rational, and commensurable in length with $BA$.

For let the square $AD$ be described on $AB$. $AD$ is thus rational
[Def.~10.4]. And $AC$ (is) also rational. $DA$
is thus commensurable  with $AC$. And as $DA$ is to $AC$, so $DB$ (is)
to $BC$ [Prop.~6.1]. Thus, $DB$ is also
commensurable (in length) with $BC$ [Prop.~10.11]. And $DB$ (is) equal to $BA$. Thus, $AB$ (is)
also commensurable (in length) with $BC$. And $AB$ is rational. Thus, $BC$ is
also rational, and commensurable in length with $AB$ [Def.~10.3]. 

Thus, if  a rational (area) is applied to a rational (straight-line), and so
on \ldots.}
\end{Parallel}

%%%%%%
% Prop 10.21
%%%%%%
\pdfbookmark[1]{Proposition 10.21}{pdf10.21}
\begin{Parallel}{}{} 
\ParallelLText{
\begin{center}
{\large \ggn{21}.}
\end{center}\vspace*{-7pt}

{\greekfont Τὸ ὑπὸ ῥητῶν δυνάμει μόνον συμμέτρων εὐθειῶν περιεχόμενον
ὀρθογώνιον ἄλογόν ἐστιν, καὶ ἡ δυναμένη α ὐτὸ ἄλογός ἐστιν,
καλείσθω δὲ μέση.}

\epsfysize=2.3in
\centerline{\epsffile{Book10/fig020g.eps}}

{\greekfont Ὑπὸ γὰρ ῥητῶν δυνάμει μόνον συμμέτρων εὐθειῶν τῶν
ΑΒ, ΒΓ ὀρθογώνιον περιεθέσθω τὸ ΑΓ; λέγω, ὅτι ἄλογόν ἐστι
τὸ ΑΓ, καὶ ἡ δυναμένη α ὐτὸ ἄλογός ἐστιν, καλείσθω δὲ μέση.}

{\greekfont Ἀναγεγράφθω γὰρ ἀπὸ τῆς ΑΒ τετράγωνον τὸ ΑΔ; ῥητὸν ἄρα
ἐστὶ τὸ ΑΔ. καὶ ἐπεὶ ἀσύμμετρός ἐστιν ἡ ΑΒ τῇ ΒΓ μήκει;
δυνάμει γὰρ μόνον ὑπόκεινται σύμμετροι; ἴση δὲ ἡ ΑΒ τῇ
ΒΔ, ἀσύμμετρος ἄρα ἐστὶ καὶ ἡ ΔΒ τῇ ΒΓ μήκει. καί ἐστιν
ὡς ἡ ΔΒ πρὸς τὴν ΒΓ, οὕτως τὸ ΑΔ πρὸς τὸ ΑΓ; ἀσύμμετρον
ἄρα [ἐστὶ] τὸ ΔΑ τῷ ΑΓ. ῥητὸν δὲ τὸ ΔΑ; ἄλογον ἄρα ἐστὶ
τὸ ΑΓ; ὥστε καὶ ἡ δυναμένη τὸ ΑΓ [τουτέστιν ἡ ἴσον α ὐτῷ
τετράγωνον δυναμένη] ἄλογός ἐστιν, καλείσθω δε μέση; ὅπερ ἔδει
δεῖχαι.}}

\ParallelRText{
\begin{center}
{\large Proposition 21}
\end{center}

The rectangle contained by rational straight-lines (which are) commensurable in square only is irrational, and its square-root  is irrational---let it be called medial.$^\dag$

\epsfysize=2.3in
\centerline{\epsffile{Book10/fig020e.eps}}

For let the rectangle $AC$ be contained by the rational straight-lines $AB$ and $BC$  (which are) commensurable in square only. I say that $AC$ is irrational, and
its square-root is irrational---let it be called
medial.

For let the square $AD$ be described on  $AB$. $AD$ is thus
rational [Def.~10.4]. And since $AB$ is incommensurable in length with $BC$. For they were assumed to be commensurable in square only. And $AB$ (is) equal to $BD$. $DB$ is thus
also incommensurable in length with $BC$. And as $DB$ is to $BC$, so
$AD$ (is) to $AC$ [Prop.~6.1]. Thus, $DA$ [is]
incommensurable with $AC$ [Prop.~10.11]. And $DA$ (is) rational. Thus, $AC$ is irrational [Def.~10.4]. Hence, its square-root [that is to say, the square-root of the square equal to it] is also irrational [Def.~10.4]---let it be called medial. (Which is) the
very thing it was required to show.}
\end{Parallel}
{\footnotesize\noindent$^\dag$ 
Thus, a medial straight-line has a length expressible as $k^{1/4}$.}

%%%%%%
% Prop 10.21a
%%%%%%
\begin{Parallel}{}{} 
\ParallelLText{
\begin{center}
{\large {\greekfont Λῆμμα}.}
\end{center}\vspace*{-7pt}

{\greekfont Ἐὰν ὦσι δύο εὐθεῖαι, ἔστιν ὡς ἡ πρώτη πρὸς τὴν
δευτέραν, οὕτως τὸ ἀπὸ τῆς πρώτης πρὸς τὸ ὑπὸ
τῶν δύο εὐθειῶν.}\\

\epsfysize=1.1in
\centerline{\epsffile{Book10/fig021ag.eps}}

{\greekfont Ἔστωσαν δύο εὐθεῖαι αἱ ΖΕ, ΕΗ. λέγω, ὅτι ἐστὶν ὡς ἡ ΖΕ
πρὸς τὴν ΕΗ, οὕτως τὸ ἀπὸ τῆς ΖΕ πρὸς τὸ ὑπὸ τῶν
ΖΕ, ΕΗ.}

{\greekfont Ἀναγεγράφθω γὰρ ἀπὸ τῆς ΖΕ τετράγωνον τὸ ΔΖ, καὶ συμπεπληρώσθω
τὸ ΗΔ. ἐπεὶ οὖν ἐστιν ὡς ἡ ΖΕ πρὸς τὴν ΕΗ, οὕτως
τὸ ΖΔ πρὸς τὸ ΔΗ, καί ἐστι τὸ μὲν ΖΔ τὸ ἀπὸ τῆς ΖΕ, τὸ δὲ
ΔΗ τὸ ὑπὸ τῶν ΔΕ,  ΕΗ, τουτέστι τὸ ὑπὸ τῶν ΖΕ, ΕΗ, ἔστιν
ἄρα ὡς ἡ ΖΕ πρὸς τὴν ΕΗ, οὕτως τὸ ἀπὸ τῆς ΖΕ πρὸς
τὸ ὑπὸ τῶν ΖΕ, ΕΗ. ὁμοίως δὲ καὶ ὡς τὸ ὑπὸ
τῶν ΗΕ, ΕΖ πρὸς τὸ ἀπὸ τῆς ΕΖ, τουτέστιν ὡς τὸ ΗΔ πρὸς τὸ ΖΔ,
οὕτως ἡ ΗΕ πρὸς τὴν ΕΖ; ὅπερ ἔδει δεῖχαι.}}

\ParallelRText{
\begin{center}
{\large Lemma}
\end{center}

If there are two straight-lines then as the first is to the second, so the
(square) on the first (is) to the (rectangle contained) by the two straight-lines.

\epsfysize=1.1in
\centerline{\epsffile{Book10/fig021ae.eps}}

Let $FE$ and $EG$ be two straight-lines. I say that as $FE$ is to $EG$,
so the (square) on $FE$ (is) to the (rectangle contained) by $FE$ and $EG$.

For let the square $DF$ be described on $FE$. And let $GD$ 
be completed. Therefore, since as $FE$ is to $EG$, so
$FD$ (is) to $DG$ [Prop.~6.1], and
$FD$ is the (square) on  $FE$, and $DG$ the (rectangle contained) by $DE$ and $EG$---that is to say,
the (rectangle contained) by $FE$ and $EG$---thus as $FE$ is to $EG$,
so the (square) on $FE$ (is) to the (rectangle contained) by $FE$ and $EG$. 
And also, similarly, as the (rectangle contained) by $GE$ and $EF$ is to the
(square on) $EF$---that is to say, as $GD$ (is) to $FD$---so $GE$
(is) to $EF$. (Which is) the very thing it was required to show.}
\end{Parallel}

%%%%%%
% Prop 10.22
%%%%%%
\pdfbookmark[1]{Proposition 10.22}{pdf10.22}
\begin{Parallel}{}{} 
\ParallelLText{
\begin{center}
{\large \ggn{22}.}
\end{center}\vspace*{-7pt}

{\greekfont Τὸ ἀπὸ μέσης παρὰ ῥητὴν παραβαλλόμενον πλάτος ποιεῖ ῥητὴν καὶ
ἀσύμμετρον τῇ, παρ ἣν παράκειται, μήκει.}

{\greekfont Ἔστω μέση μὲν ἡ Α, ῥητὴ δὲ ἡ ΓΒ, καὶ τῷ ἀπὸ
τῆς Α ἴσον παρὰ τὴν ΒΓ παραβεβλήσθω θωρίον ὀρθογώνιον τὸ ΒΔ
πλάτος ποιοῦν τὴν ΓΔ; λέγω, ὅτι ῥητή ἐστιν ἡ ΓΔ καὶ
ἀσύμμετρος τῇ ΓΒ μήκει.}\\~\\~\\

\epsfysize=2in
\centerline{\epsffile{Book10/fig022g.eps}}

{\greekfont Ἐπεὶ γὰρ μέση ἐστὶν ἡ Α, δύναται θωρίον περιεχόμενον ὑπὸ
ῥητῶν δυνάμει μόνον συμμέτρων. δυνάσθω
τὸ ΗΖ. δύναται δὲ καὶ τὸ ΒΔ; ἴσον ἄρα ἐστὶ τὸ ΒΔ τῷ ΗΖ.
ἔστι δὲ α ὐτῷ καὶ ἰσογώνιον; τῶν δὲ ἴσων τε καὶ ἰσογωνίων
παραλληλογράμμων ἀντιπεπόνθασιν αἱ πλευραὶ αἱ περὶ τὰς
ἴσας  γωνίας; ἀνάλογον ἄρα ἐστὶν ὡς ἡ ΒΓ πρὸς τὴν
ΕΗ, οὕτως ἡ ΕΖ πρὸς τὴν ΓΔ. ἔστιν ἄρα καὶ ὡς τὸ ἀπὸ
τῆς ΒΓ πρὸς τὸ ἀπὸ τῆς ΕΗ, οὕτως τὸ ἀπὸ
 τῆς ΕΖ πρὸς τὸ ἀπὸ τῆς ΓΔ. σύμμετρον δέ ἐστι τὸ ἀπὸ
τῆς ΓΒ τῷ ἀπὸ τῆς ΕΗ; ῥητὴ γάρ ἐστιν ἑκατέρα α ὐτῶν;
σύμμετρον ἄρα ἐστὶ καὶ τὸ ἀπὸ τῆς ΕΖ τῷ ἀπὸ
τῆς ΓΔ. ῥητὸν δέ ἐστι τὸ ἀπὸ τῆς ΕΖ; ῥητὸν ἄρα ἐστὶ
καὶ τὸ ἀπὸ τῆς ΓΔ; ῥητὴ ἄρα ἐστὶν ἡ ΓΔ. καὶ ἐπεὶ
ἀσύμμετρός ἐστιν ἡ ΕΖ τῇ ΕΗ μήκει; δυνάμει
γὰρ μόνον εἰσὶ σύμμετροι; ὡς δὲ ἡ ΕΖ πρὸς τὴν ΕΗ,
οὕτως τὸ ἀπὸ τῆς ΕΖ πρὸς τὸ ὑπὸ τῶν ΖΕ, ΕΗ, ἀσύμμετρον
ἄρα [ἐστὶ] τὸ ἀπὸ τῆς ΕΖ τῷ ὑπὸ τῶν ΖΕ, ΕΗ. ἀλλὰ τῷ
μὲν ἀπὸ τῆς ΕΖ σύμμετρόν ἐστι τὸ ἀπὸ τῆς ΓΔ; ῥηταὶ
γάρ εἰσι δυνάμει; τῷ δὲ ὑπὸ τῶν ΖΕ, ΕΗ σύμμετρόν ἐστι
τὸ ὑπὸ τῶν ΔΓ, ΓΒ; ἴσα γάρ ἐστι τῷ ἀπὸ τῆς Α; ἀσύμμετρον
ἄρα ἐστὶ καὶ τὸ ἀπὸ τῆς ΓΔ τῷ ὑπὸ τῶν ΔΓ, ΓΒ. ὡς δὲ
τὸ ἀπὸ τῆς ΓΔ πρὸς τὸ ὑπὸ τῶν ΔΓ, ΓΒ, οὕτως
ἐστὶν ἡ ΔΓ πρὸς τὴν ΓΒ; ἀσύμμετρος ἄρα ἐστὶν
ἡ ΔΓ τῇ ΓΒ μήκει. ῥητὴ ἄρα ἐστὶν ἡ ΓΔ καὶ ἀσύμμετρος
τῇ ΓΒ μήκει; ὅπερ ἔδει δεῖχαι.}}

\ParallelRText{
\begin{center}
{\large Proposition 22}
\end{center}

The square on a medial (straight-line), being
applied to a rational (straight-line), produces as breadth a (straight-line which
is) rational, and incommensurable in length with the (straight-line) to which it is applied.

Let $A$ be a medial (straight-line), and $CB$ a rational (straight-line),
and let the rectangular area $BD$, equal to the (square) on $A$, 
 be applied to $BC$, producing $CD$ as breadth. I say that
$CD$ is rational, and incommensurable in  length with $CB$.

\epsfysize=2in
\centerline{\epsffile{Book10/fig022e.eps}}

For since $A$ is medial, the square on it is equal to a (rectangular) area
contained by rational (straight-lines which are) commensurable in square
only [Prop.~10.21]. Let the square on ($A$) be equal
to $GF$. And the square on ($A$) is also equal to $BD$. Thus, $BD$
is equal to $GF$. And ($BD$) is also equiangular with ($GF$). And for equal and
equiangular parallelograms, the sides about the equal angles are
reciprocally proportional [Prop.~6.14]. 
Thus, proportionally,  as $BC$ is to $EG$, so $EF$ (is) to $CD$.
And, also, as the (square) on $BC$ is to the (square) on $EG$, so
the (square) on $EF$ (is) to the (square) on $CD$ [Prop.~6.22]. 
And the (square) on $CB$ is commensurable with the (square) on $EG$.
For they are each rational. Thus, the (square) on $EF$ is also commensurable
with the (square) on $CD$ [Prop.~10.11]. And
the (square) on $EF$ is rational. Thus, the (square) on $CD$ is also rational
[Def.~10.4]. Thus, $CD$ is rational.
And since $EF$ is incommensurable in length with
$EG$. For they are commensurable in square only. And as $EF$ (is) to $EG$,
so the (square) on $EF$ (is) to the (rectangle contained) by $FE$ and $EG$ 
[see previous lemma]. The (square) on $EF$ [is] thus incommensurable with the (rectangle contained) by $FE$ and $EG$ [Prop.~10.11]. But, the (square) on $CD$
is commensurable with the (square) on $EF$. For they are rational
in square.
And the (rectangle contained) by $DC$ and $CB$ is commensurable
with the (rectangle contained) by $FE$ and $EG$. For they are (both)
equal to the (square) on $A$. Thus, the (square) on $CD$ is also
incommensurable with the (rectangle contained) by $DC$ and $CB$
[Prop.~10.13]. And as the (square) on $CD$
(is) to the (rectangle contained) by $DC$ and $CB$, so $DC$ is to $CB$ [see
previous lemma]. Thus, $DC$ is incommensurable in length with $CB$
[Prop.~10.11]. Thus, $CD$ is rational, and
incommensurable in length with $CB$. (Which is) the very thing it was required to show.}
\end{Parallel}
{\footnotesize\noindent$^\dag$  Literally, ``rational''.\\}

%%%%%%%
% Prop 10.23
%%%%%%%
\pdfbookmark[1]{Proposition 10.23}{pdf10.23}
\begin{Parallel}{}{} 
\ParallelLText{
\begin{center}
{\large \ggn{23}.}
\end{center}\vspace*{-7pt}

{\greekfont Ἡ  τῇ μέσῃ σύμμετρος μέση ἐστίν.}

{\greekfont Ἔστω μέση ἡ Α, καὶ τῇ Α σύμμετρος ἔστω ἡ Β; λέγω, ὅτι καὶ
ἡ Β μέση ἐστίν.}

{\greekfont Ἐκκείσθω γὰρ ῥητὴ ἡ ΓΔ, καὶ τῷ μὲν ἀπὸ τῆς
Α ἴσον παρὰ τὴν ΓΔ παραβεβλήσθω θωρίον ὀρθογώνιον τὸ ΓΕ πλάτος
ποιοῦν τὴν ΕΔ; ῥητὴ ἄρα ἐστὶν ἡ ΕΔ καὶ ἀσύμμετρος τῂ
ΓΔ μήκει. τῷ δὲ ἀπὸ τῆς Β ἴσον παρὰ τὴν ΓΔ παραβεβλήσθω
θωρίον ὀρθογώνιον τὸ ΓΖ πλάτος ποιοῦν τὴν ΔΖ. ἐπεὶ οὖν
σύμμετρός ἐστιν ἡ Α τῇ Β, σύμμετρόν ἐστι καὶ τὸ ἀπὸ τῆς Α τῷ
ἀπὸ τῆς Β. ἀλλὰ τῷ μὲν ἀπὸ τῆς Α ἴσον ἐστὶ τὸ ΕΓ, τῷ
δὲ ἀπὸ τῆς Β ἴσον ἐστὶ τὸ ΓΖ; σύμμετρον ἄρα ἐστὶ τὸ ΕΓ τῷ ΓΖ. καί ἐστιν ὡς τὸ ΕΓ πρὸς τὸ ΓΖ, οὕτως ἡ ΕΔ πρὸς τὴν
ΔΖ; σύμμετρος ἄρα ἐστὶν ἡ ΕΔ τῇ ΔΖ μήκει. ῥητὴ δέ ἐστιν
ἡ ΕΔ καὶ ἀσύμμετρος τῇ ΔΓ μήκει; ῥητὴ ἄρα ἐστὶ καὶ ἡ ΔΖ
καὶ ἀσύμμετρος τῇ ΔΓ μήκει; αἱ ΓΔ, ΔΖ ἄρα ῥηταί εἰσι
δυνάμει μόνον σύμμετροι. ἡ δὲ τὸ ὑπὸ ῥητῶν δυνάμει
μόνον συμμέτρων δυναμένη μέση ἐστίν. ἡ ἄρα τὸ ὑπὸ τῶν
ΓΔ, ΔΖ δυναμένη μέση ἐστίν; καὶ δύναται τὸ ὑπὸ τῶν ΓΔ, ΔΖ
ἡ Β; μέση ἄρα ἐστὶν ἡ Β.}\\~\\~\\~\\~\\

\epsfysize=2.1in
\centerline{\epsffile{Book10/fig023g.eps}}

\begin{center}~\\
{\large {\greekfont Πόρισμα}.}
\end{center}\vspace*{-7pt}

{\greekfont Ἐκ δὴ τούτου φανερόν, ὅτι τὸ τῷ μέσῳ θωρίῳ σύμμετρ-ον μέσον
ἐστίν.}}

\ParallelRText{
\begin{center}
{\large Proposition 23}
\end{center}

A (straight-line) commensurable with
a medial (straight-line) is medial.

Let $A$ be a medial (straight-line), and let $B$ be commensurable with $A$.
I say that $B$ is also a medial (staight-line).

Let the rational (straight-line) $CD$ be set out, and let the rectangular
area $CE$, equal to the (square) on $A$, be applied to $CD$,
producing $ED$ as width. $ED$ is thus rational, and incommensurable
in length with $CD$ [Prop.~10.22]. And
let the rectangular area $CF$, equal to the (square) on $B$, be
applied to $CD$, producing $DF$ as width.
Therefore, since $A$ is commensurable with $B$, the (square) on $A$
is also commensurable with the (square) on $B$.  But, $EC$ is equal
to the (square) on $A$, and $CF$ is equal to the (square) on $B$. Thus,
$EC$ is commensurable with $CF$. And as $EC$ is to $CF$, so $ED$ (is)
to $DF$ [Prop.~6.1]. Thus, $ED$ is commensurable
in  length with $DF$ [Prop.~10.11]. And $ED$
is rational, and incommensurable in length with $CD$.  $DF$ is thus
 also rational [Def.~10.3], and incommensurable in length with $DC$  [Prop.~10.13]. Thus,
 $CD$ and $DF$ are rational, and commensurable in square only.
 And the square-root of a (rectangle contained) by rational (straight-lines which are)
 commensurable in square only is medial [Prop.~10.21]. Thus, the square-root of the
 (rectangle contained) by $CD$ and $DF$ is medial.
 And the square on $B$ is equal to
 the (rectangle contained) by $CD$ and $DF$. Thus, $B$ is a medial (straight-line).
 
\epsfysize=2.1in
\centerline{\epsffile{Book10/fig023e.eps}}
 
\begin{center}~\\
 {\large Corollary}
 \end{center}\vspace*{-7pt}
 
 And (it is) clear, from this, that an (area) commensurable with a medial 
 area$^\dag$ is medial.}
\end{Parallel}
{\footnotesize\noindent$^\dag$ A medial area is equal to the square on some medial straight-line. Hence, a medial area is expressible as $k^{1/2}$.}
 
%%%%%%
% Prop 10.24
%%%%%%
\pdfbookmark[1]{Proposition 10.24}{pdf10.24}
\begin{Parallel}{}{} 
\ParallelLText{
\begin{center}
{\large \ggn{24}.}
\end{center}\vspace*{-7pt}

{\greekfont Τὸ ὑπὸ μέσων μήκει συμμέτρων εὐθειῶν  περιεχόμενον
ὀρθογώνιον μέσον
ἐστίν.}

{\greekfont Ὑπὸ γὰρ μέσων μήκει συμμέτρων εὐθειῶν τῶν ΑΒ, ΒΓ
περιεθέσθω ὀρθογώνιον τὸ ΑΓ; λέγω, ὅτι τὸ ΑΓ μέσον ἐστίν.}

{\greekfont Ἀναγεγράφθω γὰρ ἀπὸ τῆς ΑΒ τετράγωνον τὸ ΑΔ; μέσον ἄρα
ἐστὶ τὸ ΑΔ. καὶ ἐπεὶ σύμμετρός ἐστιν ἡ ΑΒ τῇ ΒΓ
μήκει, ἴση δὲ ἡ ΑΒ τῇ ΒΔ, σύμμετρος ἄρα ἐστὶ καὶ ἡ ΔΒ τῇ
ΒΓ μήκει; ὥστε καὶ τὸ ΔΑ τῷ ΑΓ σύμμετρόν ἐστιν. μέσον
δὲ τὸ ΔΑ; μέσον ἄρα καὶ τὸ ΑΓ; ὅπερ ἔδει δεῖχαι.}\\~\\~\\

\epsfysize=2.2in
\centerline{\epsffile{Book10/fig024g.eps}}
}

\ParallelRText{
\begin{center}
{\large Proposition 24}
\end{center}

A rectangle contained by medial straight-lines
(which are) commensurable in length is medial.

For let the rectangle $AC$ be contained by the medial straight-lines
$AB$ and $BC$ (which are) commensurable in length. I say that $AC$
is medial.

For let the square $AD$ be described on $AB$. $AD$ is thus medial [see previous footnote]. And since $AB$ is commensurable in length
with $BC$, and $AB$ (is) equal to $BD$, $DB$ is thus also commensurable
in length with $BC$. Hence, $DA$ is also commensurable with $AC$
[Props.~6.1, 10.11].
And $DA$ (is) medial. Thus, $AC$ (is) also medial [Prop.~10.23~corr.]. (Which is) the very thing it
was required to show.

\epsfysize=2.2in
\centerline{\epsffile{Book10/fig024e.eps}}}
\end{Parallel}

%%%%%%
% Prop 10.25
%%%%%%
\pdfbookmark[1]{Proposition 10.25}{pdf10.25}
\begin{Parallel}{}{} 
\ParallelLText{
\begin{center}
{\large \ggn{25}.}
\end{center}\vspace*{-7pt}

{\greekfont Τὸ ὑπὸ μέσων δυνάμει μόνον  συμμέτρων εὐθειῶν περιεχόμενον
ὀρθογώνιον ἤτοι ῥητὸν ἢ μέσον ἐστίν.}

\epsfysize=1.8in
\centerline{\epsffile{Book10/fig025g.eps}}

{\greekfont Ὑπὸ γὰρ μέσων δυνάμει μόνον συμμέτρων εὐθειῶν τῶν ΑΒ, ΒΓ
ὀρθογώνιον περιεθέσθω τὸ ΑΓ; λέγω, ὅτι τὸ ΑΓ ἤτοι ῥητὸν
ἢ μέσον ἐστίν.}

{\greekfont Ἀναγεγράφθω γὰρ ἀπὸ τῶν ΑΒ, ΒΓ τετράγωνα τὰ ΑΔ, ΒΕ; μέσον
ἄρα ἐστὶν ἑκάτερον τῶν ΑΔ, ΒΕ. καὶ ἐκκείσθω ῥητὴ ἡ ΖΗ, καὶ
τῷ μὲν ΑΔ ἴσον παρὰ τὴν ΖΗ παραβεβλήσθω ὀρθογώνιον παραλληλόγραμμον τὸ ΗΘ πλάτος ποιοῦν τὴν ΖΘ, τῷ δὲ ΑΓ ἴσον
παρὰ τὴν ΘΜ παραβεβλήσθω ὀρθογώνιον παραλληλόγραμμον τὸ ΜΚ πλάτος
ποιοῦν τὴν ΘΚ, καὶ ἔτι τῷ ΒΕ ἴσον ὁμοίως παρὰ τὴν
ΚΝ παραβεβλήσθω τὸ ΝΛ πλάτος ποιοῦν τὴν ΚΛ; ἐπ εὐθείας
ἄρα εἰσὶν αἱ ΖΘ, ΘΚ, ΚΛ. ἐπεὶ οὖν μέσον ἐστὶν ἑκάτερον
τῶν ΑΔ, ΒΕ, καί ἐστιν ἴσον τὸ μὲν ΑΔ τῷ ΗΘ, τὸ δὲ ΒΕ τῷ
ΝΛ, μέσον ἄρα καὶ ἑκάτερον τῶν ΗΘ, ΝΛ. καὶ παρὰ ῥητὴν
τὴν ΖΗ παράκειται; ῥητὴ ἄρα ἐστὶν ἑκατέρα τῶν ΖΘ, ΚΛ καὶ
ἀσύμμετρος τῇ ΖΗ μήκει. καὶ ἐπεὶ σύμμετρόν ἐστι τὸ ΑΔ
τῷ ΒΕ, σύμμετρον ἄρα ἐστὶ καὶ τὸ ΗΘ τῷ ΝΛ. καί ἐστιν
ὡς τὸ ΗΘ πρὸς τὸ ΝΛ, οὕτως ἡ ΖΘ πρὸς τὴν ΚΛ;
σύμμετρος ἄρα ἐστὶν ἡ ΖΘ τῇ ΚΛ μήκει. αἱ ΖΘ, ΚΛ
ἄρα ῥηταί εἰσι μήκει σύμμετροι; ῥητὸν ἄρα ἐστὶ τὸ ὑπὸ
τῶν ΖΘ, ΚΛ. καὶ ἐπεὶ ἴση ἐστὶν ἡ μὲν ΔΒ τῇ ΒΑ, ἡ δὲ
ΧΒ τῇ ΒΓ, ἔστιν ἄρα ὡς ἡ ΔΒ πρὸς τὴν ΒΓ, οὕτως ἡ ΑΒ
πρὸς τὴν ΒΧ. ἀλλ ὡς μὲν ἡ ΔΒ πρὸς τὴν ΒΓ, οὕτως
τὸ ΔΑ πρὸς τὸ ΑΓ; ὡς δὲ ἡ ΑΒ πρὸς τὴν ΒΧ, οὕτως 
τὸ ΑΓ πρὸς τὸ ΓΧ; ἔστιν ἄρα ὡς τὸ ΔΑ πρὸς τὸ ΑΓ, οὕτως
τὸ ΑΓ πρὸς τὸ ΓΧ. ἴσον δέ ἐστι τὸ μὲν ΑΔ τῷ ΗΘ, τὸ δὲ
ΑΓ τῷ ΜΚ, τὸ δὲ ΓΧ τῷ ΝΛ; ἔστιν ἄρα ὡς τὸ ΗΘ πρὸς τὸ ΜΚ,
οὕτως τὸ ΜΚ πρὸς τὸ ΝΛ; ἔστιν ἄρα καὶ ὡς
ἡ ΖΘ πρὸς τὴν ΘΚ, οὕτως ἡ ΘΚ πρὸς τὴν ΚΛ; τὸ ἄρα ὑπὸ
τῶν ΖΘ, ΚΛ ἴσον ἐστὶ τῷ ἀπὸ τῆς ΘΚ. ῥητὸν δὲ τὸ ὑπὸ
τῶν ΖΘ, ΚΛ; ῥητὸν ἄρα ἐστὶ καὶ τὸ ἀπὸ τῆς ΘΚ; ῥητὴ ἄρα ἐστὶν
ἡ ΘΚ. καὶ εἰ μὲν σύμμετρός ἐστι τῇ ΖΗ μήκει, ῥητόν ἐστι
τὸ ΘΝ; εἰ δὲ ἀσύμμετρός ἐστι τῇ ΖΗ μήκει, αἱ ΚΘ, ΘΜ ῥηταί
εἰσι δυνάμει μόνον σύμμετροι; μέσον ἄρα τὸ ΘΝ. τὸ ΘΝ
ἄρα ἤτοι ῥητὸν ἢ μέσον ἐστίν. ἴσον δὲ τὸ ΘΝ τῷ ΑΓ; τὸ
ΑΓ ἄρα ἤτοι ῥητὸν ἢ μέσον ἐστίν.}

{\greekfont Τὸ ἄρα ὑπὸ μέσων δυνάμει μόνον συμμέτρων, καὶ τὰ εχῆς.}}

\ParallelRText{
\begin{center}
{\large Proposition 25}
\end{center}

The rectangle contained by medial
straight-lines (which are) commensurable in square only is either
rational or medial.

\epsfysize=1.8in
\centerline{\epsffile{Book10/fig025e.eps}}

For let the rectangle $AC$ be contained by the medial straight-lines
$AB$ and $BC$ (which are) commensurable in square only. I say that
$AC$ is either rational or medial.

For let the squares $AD$ and $BE$ be described on (the straight-lines) $AB$ and
$BC$ (respectively). $AD$ and $BE$ are thus each medial. And let
the rational (straight-line) $FG$ be laid out. And let the rectangular parallelogram $GH$, equal to $AD$, be applied to $FG$, producing
$FH$ as breadth. And let the rectangular parallelogram $MK$, equal to $AC$, be applied to $HM$, producing $HK$ as breadth. And,
finally, let $NL$, equal to $BE$, have similarly been applied to $KN$,
producing $KL$ as breadth. Thus, $FH$, $HK$, and $KL$ are in
a straight-line.  Therefore, since $AD$ and $BE$ are each medial, and
$AD$ is equal to $GH$, and $BE$ to $NL$, $GH$ and $NL$ (are)
thus each also medial. And they are applied to the rational (straight-line)
$FG$. $FH$ and $KL$ are thus each rational, and incommensurable in length
with $FG$ [Prop.~10.22]. And since
$AD$ is commensurable with $BE$, $GH$ is thus also commensurable
with $NL$. And as $GH$ is to $NL$, so $FH$ (is) to $KL$ [Prop.~6.1]. Thus, $FH$ is commensurable
in length with $KL$ [Prop.~10.11]. Thus, 
$FH$ and $KL$ are rational (straight-lines which are) commensurable in length. Thus, the (rectangle contained) by $FH$ and $KL$ is rational
[Prop.~10.19]. And since $DB$ is equal to $BA$,
and $OB$ to $BC$, thus as $DB$ is to $BC$, so $AB$ (is) to $BO$.
But, as $DB$ (is) to $BC$, so $DA$ (is) to $AC$ [Props.~6.1]. 
And as $AB$ (is) to $BO$, so $AC$ (is)
to $CO$ [Prop.~6.1]. Thus, as $DA$ is to $AC$, so $AC$ (is)
to $CO$. And $AD$ is equal to $GH$,
and $AC$ to $MK$, and $CO$ to $NL$. Thus, as $GH$ is to $MK$,
so $MK$ (is) to $NL$. Thus, also, as $FH$ is to $HK$, so $HK$ (is) to
$KL$ [Props.~6.1, 5.11]. Thus, the (rectangle contained) by
$FH$ and $KL$ is equal to the (square) on $HK$ [Prop.~6.17]. And the (rectangle contained) by
$FH$ and $KL$ (is) rational. Thus, the (square) on $HK$ is also
rational. Thus, $HK$ is rational. And if it is commensurable in length with
$FG$ then $HN$ is rational [Prop.~10.19]. 
And if it is incommensurable in length with $FG$ then $KH$ and $HM$
are rational (straight-lines which are) commensurable in square only: thus,
$HN$ is medial [Prop.~10.21]. Thus, $HN$ is
either rational or medial. And $HN$ (is) equal to $AC$. Thus, $AC$
is either rational or medial.

Thus, the \ldots by medial
straight-lines (which are) commensurable in square only, and so on \ldots.}
\end{Parallel}

%%%%%%
% Prop 10.26
%%%%%%
\pdfbookmark[1]{Proposition 10.26}{pdf10.26}
\begin{Parallel}{}{} 
\ParallelLText{
\begin{center}
{\large \ggn{26}.}
\end{center}\vspace*{-7pt}

{\greekfont Μέσον μέσου οὐθ ὑπερέθει ῥητῷ.}\\

\epsfysize=1.8in
\centerline{\epsffile{Book10/fig026g.eps}}

{\greekfont Εἰ γὰρ δυνατόν, μέσον τὸ ΑΒ μέσου τοῦ ΑΓ ὑπερεθέτω
ῥητῷ τῷ ΔΒ, καὶ ἐκκείσθω ῥητὴ ἡ ΕΖ, καὶ τῷ ΑΒ
ἴσον παρὰ τὴν ΕΖ παραβεβλήσθω παραλληλόγραμμον ὀρθογώνιον
τὸ ΖΘ πλάτος ποιοῦν τὴν ΕΘ, τῷ δὲ ΑΓ ἴσον ἀφῃρήσθω τὸ ΖΗ;
λοιπὸν ἄρα τὸ ΒΔ λοιπῷ τῷ ΚΘ ἐστιν ἴσον. ῥητὸν δέ ἐστι
τὸ ΔΒ; ῥητὸν ἄρα ἐστὶ καὶ τὸ ΚΘ. ἐπεὶ οὖν μέσον ἐστὶν
ἑκάτερον τῶν ΑΒ, ΑΓ, καί ἐστι τὸ μὲν ΑΒ τῷ ΖΘ ἴσον,
τὸ δὲ ΑΓ τῷ ΖΗ, μέσον ἄρα καὶ ἑκάτερον τῶν ΖΘ, ΖΗ. καὶ
παρὰ ῥητὴν τὴν ΕΖ παράκειται; ῥητὴ ἄρα ἐστὶν ἑκατέρα
τῶν ΘΕ, ΕΗ καὶ ἀσύμμετρος τῇ ΕΖ μήκει. καὶ ἐπεὶ ῥητόν
ἐστι τὸ ΔΒ καί ἐστιν ἴσον τῷ ΚΘ, ῥητὸν ἄρα ἐστὶ καὶ τὸ
ΚΘ. καὶ παρὰ ῥητὴν τὴν ΕΖ παράκειται; ῥητὴ ἄρα ἐστὶν ἡ ΗΘ
καὶ σύμμετρος τῇ ΕΖ μήκει. ἀλλά καὶ ἡ ΕΗ ῥητή ἐστι καὶ
ἀσύμμετρος τῇ ΕΖ μήκει; ἀσύμμετρος ἄρα ἐστὶν ἡ ΕΗ
τῇ ΗΘ μήκει. καί ἐστιν ὡς ἡ ΕΗ πρὸς τὴν ΗΘ, οὕτως τὸ
ἀπὸ τῆς ΕΗ πρὸς τὸ ὑπὸ τῶν ΕΗ, ΗΘ; ἀσύμμετρον
ἄρα ἐστὶ τὸ ἀπὸ τῆς ΕΗ τῷ ὑπὸ τῶν ΕΗ, ΗΘ. ἀλλὰ
τῷ μὲν ἀπὸ τῆς ΕΗ σύμμετρά ἐστι τὰ ἀπὸ τῶν ΕΗ, ΗΘ 
τετράγωνα; ῥητὰ γὰρ ἀμφότερα;
τῷ δὲ ὑπὸ τῶν ΕΗ, ΗΘ σύμμετρόν ἐστι τὸ δὶς ὑπὸ τῶν
ΕΗ, ΗΘ; διπλάσιον γάρ ἐστιν α ὐτοῦ; ἀσύμμετρα ἄρα ἐστὶ
τὰ ἀπὸ τῶν ΕΗ, ΗΘ τῷ δὶς ὑπὸ τῶν ΕΗ, ΗΘ; καὶ συναμφότερα ἄρα
τά τε ἀπὸ τῶν ΕΗ, ΗΘ καὶ τὸ δὶς ὑπὸ τῶν ΕΗ, ΗΘ, ὅπερ
ἐστὶ τὸ ἀπὸ τῆς ΕΘ, ἀσύμμετρόν ἐστι τοῖς ἀπὸ τῶν ΕΗ, ΗΘ.
ῥητὰ δὲ τὰ ἀπὸ τῶν ΕΗ, ΗΘ; ἄλογον ἄρα τὸ ἀπὸ
τῆς ΕΘ. ἄλογος ἄρα ἐστὶν ἡ ΕΘ. ἀλλὰ καὶ ῥηρή;
ὅπερ ἐστὶν ἀδύνατον.}

{\greekfont Μέσον ἄρα μέσου οὐθ ὑπερέθει ῥητῷ; ὅπερ ἔδει
δεῖχαι.}}

\ParallelRText{
\begin{center}
{\large Proposition 26}
\end{center}

A medial (area) does not exceed a medial
(area) by a rational (area).$^\dag$

\epsfysize=1.8in
\centerline{\epsffile{Book10/fig026e.eps}}

For, if possible, let the medial (area) $AB$ exceed the medial (area) $AC$
by the rational (area) $DB$. And let the rational (straight-line)
$EF$ be laid down. And let the rectangular parallelogram $FH$, equal to $AB$, be
applied to to $EF$, producing $EH$ as breadth. And let $FG$, equal to
$AC$, be cut off (from $FH$). Thus, the remainder $BD$
is equal to the remainder $KH$. And $DB$ is rational. Thus, $KH$
is also rational. Therefore, since $AB$ and $AC$ are each medial, and
$AB$ is equal to $FH$, and $AC$ to $FG$, $FH$ and $FG$ are thus each also medial. And they are applied to the rational (straight-line)
$EF$. Thus, $HE$ and $EG$ are each rational, and incommensurable
in length with $EF$ [Prop.~10.22]. 
And since $DB$ is rational, and is equal to $KH$, $KH$ is thus  also
rational. And ($KH$) is  applied to the rational (straight-line)
$EF$. $GH$ is thus rational, and commensurable in length with
$EF$ [Prop.~10.20]. But, $EG$ is also
rational, and incommensurable in length with $EF$. Thus, $EG$ is
incommensurable in length with $GH$ [Prop.~10.13]. And as $EG$ is to $GH$, so the
(square) on $EG$ (is) to the (rectangle contained) by $EG$ and
$GH$ [Prop.~10.13~lem.]. Thus, the (square)
on $EG$ is incommensurable with the (rectangle contained) by 
$EG$ and $GH$ [Prop.~10.11]. 
But, the (sum of the) squares on $EG$ and $GH$ is commensurable with
the (square) on $EG$. For ($EG$ and $GH$ are) both rational.
And twice the (rectangle contained) by $EG$ and $GH$ is commensurable
with the (rectangle contained) by $EG$ and $GH$ [Prop.~10.6]. For (the former) is double the latter.
Thus, the (sum of the squares) on $EG$ and $GH$ is incommensurable
with twice the (rectangle contained) by $EG$ and $GH$
[Prop.~10.13]. 
And thus the sum of the (squares) on $EG$ and $GH$ plus twice the (rectangle contained) by
$EG$ and $GH$, that is the (square) on $EH$ [Prop.~2.4], is incommensurable with the (sum
of the squares) on $EG$ and $GH$ [Prop.~10.16].
And the (sum of the squares) on $EG$ and $GH$ (is) rational. 
Thus, the (square) on $EH$ is irrational [Def.~10.4].
Thus, $EH$ is irrational  [Def.~10.4]. But,
(it is) also rational. The very thing is impossible.

Thus, a medial (area) does not exceed a medial
(area) by a rational (area). (Which is) the very thing it was required to show.}
\end{Parallel}
{\footnotesize\noindent$^\dag$ In other words, $\sqrt{k}-\sqrt{k'}\neq k''$.}

%%%%%%
% Prop 10.27
%%%%%%
\pdfbookmark[1]{Proposition 10.27}{pdf10.27}
\begin{Parallel}{}{} 
\ParallelLText{
\begin{center}
{\large \ggn{27}.}
\end{center}\vspace*{-7pt}

{\greekfont Μέσας εὑρεῖν δυνάμει μόνον συμμέτρους ῥητὸν περιεθούσας.}\\

\epsfysize=1.6in
\centerline{\epsffile{Book10/fig027g.eps}}

{\greekfont Ἐκκείσθωσαν δύο ῥηταὶ δυνάμει μόνον σύμμετροι αἱ
Α, Β, καὶ εἰλήφθω τῶν Α, Β μέση ἀνάλογον ἡ Γ, καὶ γεγονέτω
ὡς ἡ Α πρὸς τὴν Β, οὕτως ἡ Γ πρὸς τὴν Δ.}

{\greekfont Καὶ ἐπεὶ αἱ Α, Β ῥηταί εἰσι δυνάμει μόνον σύμμετροι, τὸ
ἄρα ὑπὸ τῶν Α, Β, τουτέστι τὸ ἀπὸ τῆς Γ, μέσον ἐστίν.
μέση ἄρα ἡ Γ. καὶ ἐπεί ἐστιν ὡς ἡ Α πρὸς τὴν Β,
[οὕτως] ἡ Γ πρὸς τὴν Δ, αἱ δὲ Α, Β δυνάμει μόνον [εἰσὶ]
σύμμετροι, καὶ αἱ Γ, Δ ἄρα δυνάμει μόνον εἰσὶ σύμμετροι. καί
ἐστι μέση ἡ Γ; μέση ἄρα καὶ ἡ Δ. αἱ Γ, Δ ἄρα μέσαι εἰσὶ
δυνάμει μόνον σύμμετροι. λέγω, ὅτι καὶ ῥητὸν περιέθουσιν.
ἐπεὶ γάρ ἐστιν ὡς ἡ Α πρὸς τὴν Β, οὕτως ἡ Γ πρὸς
τὴν Δ, ἐναλλὰχ ἄρα ἐστὶν ὡς ἡ Α πρὸς τὴν Γ, ἡ Β
πρὸς τὴν Δ. ἀλλ ὡς ἡ Α πρὸς τὴν Γ, ἡ Γ πρὸς τὴν Β; καὶ ὡς
ἄρα ἡ Γ πρὸς τὴν Β, οὕτως ἡ Β πρὸς τὴν Δ; τὸ ἄρα
ὑπὸ τῶν Γ, Δ ἴσον ἐστὶ τῷ ἀπὸ τῆς Β. ῥητὸν δὲ τὸ
ἀπὸ τῆς Β; ῥητὸν ἄρα [ἐστὶ] καὶ τὸ ὑπὸ τῶν
Γ, Δ.}

{\greekfont Εὕρηνται ἄρα μέσαι δυνάμει μόνον σύμμετροι ῥητὸν περιέθουσαι;
ὅπερ ἔδει δεῖχαι.}}

\ParallelRText{
\begin{center}
{\large Proposition 27}
\end{center}

To find (two) medial (straight-lines), containing a rational (area), (which are) commensurable in square only.

\epsfysize=1.6in
\centerline{\epsffile{Book10/fig027e.eps}}

Let the two rational (straight-lines) $A$ and $B$, (which are) commensurable
in square only, be laid down. And let $C$---the mean proportional (straight-line) to
$A$ and $B$---be taken [Prop.~6.13]. 
And let it be contrived that as $A$ (is) to $B$, so $C$ (is) to $D$
[Prop.~6.12].

And since the rational (straight-lines) $A$ and $B$ are commensurable in square only,  the (rectangle contained) by $A$ and $B$---that is to
say, the (square) on $C$ [Prop.~6.17]---is thus medial
[Prop~10.21]. Thus, $C$ is medial
[Prop.~10.21]. And since as $A$ is to $B$, [so]
$C$ (is) to $D$, and $A$ and $B$ [are] commensurable in square only, 
$C$ and $D$ are thus also commensurable in square only [Prop.~10.11]. And $C$ is medial.
Thus, $D$ is also medial [Prop.~10.23]. 
Thus, $C$ and $D$ are medial (straight-lines which are) commensurable
in square only. I say that they also contain a rational (area).
For since as $A$ is to $B$, so $C$ (is) to $D$, thus, alternately,
as $A$ is to $C$, so $B$ (is) to $D$ [Prop.~5.16]. 
But, as $A$ (is) to $C$, (so) $C$ (is) to $B$. 
And thus as $C$ (is) to $B$, so $B$ (is) to $D$ [Prop.~5.11].
 Thus, the (rectangle contained) by $C$ and $D$ is equal to the (square) on $B$ [Prop.~6.17]. And the (square) on $B$ (is) rational. Thus, the (rectangle contained) by $C$ and $D$ [is] also rational.
 
Thus, (two) medial (straight-lines, $C$ and $D$), containing a rational (area), (which are) commensurable in square only, be found.$^\dag$ (Which is) the very thing it was required to show.}
\end{Parallel}
{\footnotesize\noindent$^\dag$ $C$ and $D$ have lengths $k^{1/4}$ and $k^{3/4}$ times that of $A$, respectively, where the length of $B$ is $k^{1/2}$ times that of $A$.}

%%%%%%
% Prop 10.28
%%%%%%
\pdfbookmark[1]{Proposition 10.28}{pdf10.28}
\begin{Parallel}{}{} 
\ParallelLText{
\begin{center}
{\large\ggn{28}.}
\end{center}\vspace*{-7pt}

{\greekfont Μέσας εὑρεῖν δυνάμει μόνον συμμέτρους μέσον πειριεθούσας.}

\epsfysize=0.7in
\centerline{\epsffile{Book10/fig028g.eps}}

{\greekfont Ἐκκείσθωσαν [τρεῖς] ῥηταὶ δυνάμει μόνον σύμμετροι αἱ Α, Β, Γ,
καὶ εἰλήφθω τῶν Α, Β μέση ἀνάλογον ἡ Δ, καὶ γεγονέτω ὡς ἡ
Β πρὸς τὴν Γ, ἡ Δ πρὸς τὴν Ε.}

{\greekfont Ἐπεὶ αἱ Α, Β ῥηταί εἰσι δυνάμει μόνον σύμμετροι, τὸ ἄρα
ὑπὸ τῶν Α, Β, τουτέστι τὸ ἀπὸ τῆς Δ, μέσον ἐστὶν. μέση ἄρα
ἡ Δ. καὶ ἐπεὶ αἱ Β, Γ δυνάμει μόνον εἰσὶ σύμμετροι, καί
ἐστιν ὡς ἡ Β πρὸς τὴν Γ,  ἡ Δ πρὸς τὴν Ε, καὶ αἱ Δ, Ε ἄρα
δυνάμει μόνον εἰσὶ σύμμετροι. μέση δὲ ἡ Δ; μέση ἄρα καὶ
ἡ Ε; αἱ Δ, Ε ἄρα μέσαι εἰσὶ δυνάμει μόνον σύμμετροι. λέγω
δή, ὅτι καὶ μέσον περιέθουσιν. ἐπεὶ γάρ ἐστιν ὡς ἡ Β
πρὸς τὴν Γ, ἡ Δ πρὸς τὴν Ε, ἐναλλὰχ ἄρα ὡς ἡ Β πρὸς
τὴν Δ, ἡ Γ πρὸς τὴν Ε. ὡς δὲ ἡ Β πρὸς τὴν Δ, ἡ Δ πρὸς
τὴν Α; καὶ ὡς ἅρα ἡ Δ πρὸς τὴν Α, ἡ Γ πρὸς τὴν Ε;
τὸ ἄρα ὑπὸ τῶν Α, Γ ἴσον ἐστὶ τῷ ὑπὸ τῶν Δ, Ε.
μέσον δὲ τὸ ὑπὸ τῶν Α, Γ; μέσον ἄρα καὶ τὸ ὑπὸ τῶν
Δ, Ε.}

{\greekfont Ἐὕρηνται ἄρα μέσαι δυνάμει μόνον σύμμετροι μέσον περιέθουσαι;
ὅπερ ἔδει δεῖχαι.}}

\ParallelRText{
\begin{center}
{\large Proposition 28}
\end{center}

To find (two) medial (straight-lines),
containing a medial (area), (which are)
commensurable in square only.

\epsfysize=0.7in
\centerline{\epsffile{Book10/fig028e.eps}}

Let the [three] rational (straight-lines) $A$, $B$, and $C$, (which are)
commensurable in square only, be laid down. And let, $D$,  the mean proportional
(straight-line) to $A$ and $B$, be taken [Prop.~6.13]. And let it be contrived that as
$B$ (is) to $C$, (so) $D$ (is) to $E$ [Prop.~6.12].

Since the rational (straight-lines) $A$ and $B$ are commensurable in square only, the (rectangle contained) by $A$ and $B$---that is to say, the
(square) on $D$ [Prop.~6.17]---is medial
[Prop.~10.21]. Thus, $D$ (is) medial [Prop.~10.21]. And
since $B$ and $C$ are commensurable in square only, and as $B$ is to $C$, (so) $D$ (is) to $E$, $D$ and $E$ are thus commensurable in square
only [Prop.~10.11]. And $D$ (is) medial.
$E$ (is) thus also medial [Prop.~10.23].
Thus, $D$ and $E$ are medial (straight-lines which are) commensurable
in square only. So, I say that they also enclose a medial (area). For since
as $B$ is to $C$, (so) $D$ (is) to $E$, thus, alternately, as $B$ (is) to $D$,
(so) $C$ (is) to $E$ [Prop.~5.16]. And as
$B$ (is) to $D$, (so) $D$ (is) to $A$. And thus as $D$ (is) to $A$, (so)
$C$ (is) to $E$. Thus, the (rectangle contained) by $A$ and $C$ is equal
to the (rectangle contained) by $D$ and $E$ [Prop.~6.16]. And the (rectangle contained) by 
$A$ and $C$ is medial [Prop.~10.21]. Thus, the (rectangle contained) by $D$ and $E$ (is)
also medial.

Thus, (two) medial (straight-lines, $D$ and $E$), containing a medial (area),
(which are) commensurable in square only, be found. (Which is)
the very thing it was required to show.}
\end{Parallel}
{\footnotesize\noindent$^\dag$ $D$ and $E$ have lengths $k^{1/4}$ and $k'^{1/2}/k^{1/4}$ times that of $A$, respectively, where the lengths of $B$ and $C$ are
$k^{1/2}$ and $k'^{1/2}$ times that of $A$, respectively.}

%%%%%%
% Prop 10.28a
%%%%%%
\begin{Parallel}{}{} 
\ParallelLText{
\begin{center}{\large {\greekfont Λῆμμα γ̀γν{1}.}}
\end{center}\vspace*{-7pt}

{\greekfont Εὑρειν δύο τετραγώνους ἀριθμούς, ὥστε καὶ τὸν συγκείμενον
ἐχ α ὐτῶν εἶναι τετράγωνον.}

\epsfysize=0.3in
\centerline{\epsffile{Book10/fig028ag.eps}}

{\greekfont Ἐκκείσθωσαν δύο ἀριθμοὶ οἱ ΑΒ, ΒΓ, ἔστωσαν δὲ ἤτοι ἄρτιοι
ἢ περιττοί. καὶ ἐπεὶ, ἐάν τε ἀπὸ ἀρτίου ἄρτιος ἀφαιρεθῇ, ἐάν
τε ἀπὸ περισσοῦ περισσός, ὁ λοιπὸς ἄρτιός ἐστιν, ὁ λοιπὸς
ἄρα ὁ ΑΓ ἄρτιός ἐστιν. τετμήσθω ὁ ΑΓ δίθα κατὰ τὸ Δ. ἔστωσαν
δὲ καὶ οἱ ΑΒ, ΒΓ ἤτοι ὅμοιοι ἐπίπεδοι ἢ τετράγωνοι,
οἳ καὶ α ὐτοὶ ὅμοιοί εἰσιν ἐπίπεδοι;
 ὁ ἄρα ἐκ τῶν
ΑΒ, ΒΓ μετὰ τοῦ ἀπὸ [τοῦ] ΓΔ τετραγώνου ἴσος ἐστὶ τῷ
ἀπὸ τοῦ ΒΔ τετραγώνῳ. καί ἐστι τετράγωνος ὁ ἐκ
τῶν ΑΒ, ΒΓ, ἐπειδήπερ ἐδείθθη, ὅτι, ἐὰν δύο ὅμοιοι
ἐπίπεδοι πολλαπλασιάσαντες ἀλλήλους ποιῶσι τινα, ὁ
γενόμενος τετράγωνός ἐστιν. εὕρηνται ἄρα δύο τετράγωνοι ἀριθμοὶ
ὅ τε ἐκ τῶν ΑΒ, ΒΓ καὶ ὁ ἀπὸ τοῦ ΓΔ, οἳ συντεθέντες
ποιοῦσι τὸν ἀπὸ τοῦ ΒΔ τετράγωνον.}

{\greekfont Καὶ φανερόν, ὅτι εὕρηνται πάλιν δύο τετράγωνοι
ὅ τε ἀπὸ τοῦ ΒΔ καὶ ὁ ἀπὸ τοῦ ΓΔ, ὥστε τὴν ὑπεροθὴν
α ὐτῶν τὸν ὑπὸ ΑΒ, ΒΓ εἶναι τετράγωνον, ὅταν
οἱ ΑΒ, ΒΓ ὅμοιοι ὦσιν ἐπίπεδοι. ὅταν δὲ μὴ ὦσιν
ὅμοιοι ἐπίπεδοι, εὕρηνται δύο τετράγωνοι ὅ τε ἀπὸ
τοῦ ΒΔ καὶ ὁ ἀπὸ τοῦ ΔΓ, ὧν ἡ ὑπεροθὴ ὁ ὑπὸ
τῶν ΑΒ, ΒΓ οὐκ ἔστι τετράγωνος; ὅπερ ἔδει δεῖχαι.}}

\ParallelRText{
\begin{center}
{\large Lemma I}
\end{center}

To find two square numbers such that
the sum of them is also  square.

\epsfysize=0.3in
\centerline{\epsffile{Book10/fig028ae.eps}}

Let the two numbers $AB$ and $BC$ be laid down. And let them
be either (both)  even or (both) odd. And since, if an even (number) is subtracted
from an even (number), or if an odd (number is subtracted) from an odd
(number), then the remainder is even [Props.~9.24, 9.26], the remainder $AC$ is thus even. Let
$AC$ be cut in half at $D$. And let $AB$ and $BC$ also be
either similar plane (numbers), or square (numbers)---which are themselves
also similar plane (numbers). Thus, the (number created) from (multiplying)  $AB$ and
$BC$, plus the square on $CD$, is equal to the square on $BD$ [Prop.~2.6].  And the (number created) from (multiplying) $AB$ and $BC$ is square---inasmuch as it was shown that
if two similar plane (numbers) make some (number) by multiplying
one another then the (number so) created is square [Prop.~9.1]. Thus, two square numbers be
found---(namely,) the (number created) from (multiplying) $AB$ and $BC$,
and the (square) on $CD$---which, (when) added (together),
make the square on $BD$.

And (it is) clear that two square (numbers) have again been found---(namely,) the (square) on $BD$, and the (square) on $CD$---such that their difference---(namely,) the (rectangle) contained by $AB$ and $BC$---is  square
whenever $AB$ and $BC$ are similar plane (numbers). But, when they are
not similar plane numbers, two square (numbers) be found---(namely,) the (square) on $BD$, and the (square) on $DC$---between which the difference---(namely,) the (rectangle) contained by $AB$ and $BC$---is
not  square. (Which is) the very thing it was required to show.}
\end{Parallel}

%%%%%%%
% Prop 10.28b
%%%%%%%
\begin{Parallel}{}{} 
\ParallelLText{
\begin{center}
{\large {\greekfont Λῆμμα γ̀γν{2}.}}
\end{center}\vspace*{-7pt}

{\greekfont Εὑρεῖν δύο τετραγώνους ἀριθμούς, ὥστε τὸν ἐχ α ὐτῶν συγκείμενον μὴ εἶναι τετράγωνον.}

\epsfysize=0.4in
\centerline{\epsffile{Book10/fig028bg.eps}}

{\greekfont Ἔστω γὰρ ὁ ἐκ τῶν ΑΒ, ΒΓ, ὡς ἔφαμεν, τετράγωνος, καὶ
ἄρτιος ὁ ΓΑ, καὶ τετμήσθω ὁ ΓΑ δίθα τῷ Δ. φανερὸν δή,
ὅτι ὁ ἐκ τῶν ΑΒ, ΒΓ τετράγωνος μετὰ τοῦ ἀπὸ [τοῦ]
ΓΔ τετραγώνου ἴσος ἐστὶ τῷ ἀπὸ [τοῦ]
ΒΔ τετραγώνῳ. ἀφῃρήσθω μονὰς ἡ ΔΕ; ὁ ἄρα ἐκ τῶν ΑΒ, ΒΓ
μετὰ τοῦ ἀπὸ [τοῦ] ΓΕ ἐλάσσων ἐστὶ τοῦ ἀπὸ [τοῦ]
ΒΔ τετραγώνου. λέγω οὖν, ὅτι ὁ ἐκ τῶν ΑΒ, ΒΓ τετράγωνος μετὰ
τοῦ ἀπὸ [τοῦ] ΓΕ οὐκ ἔσται τετράγωνος.}

{\greekfont Εἰ γὰρ ἔσται τετράγωνος, ἤτοι ἴσος ἐστὶ τῷ ἀπὸ [τοῦ]
ΒΕ ἢ ἐλάσσων τοῦ ἀπὸ [τοῦ] ΒΕ, οὐκέτι δὲ καὶ μείζων, ἵνα
μὴ τμ ηθῇ ἡ μονάς. ἔστω, εἰ δυνατόν, πρότερον ὁ ἐκ τῶν
ΑΒ, ΒΓ μετὰ τοῦ ἀπὸ ΓΕ ἴσος τῷ ἀπὸ ΒΕ, καὶ ἔστω τῆς ΔΕ μονάδος διπλασίων ὁ ΗΑ. ἐπεὶ οὖν ὅλος ὁ ΑΓ ὅλου
τοῦ ΓΔ ἐστι διπλασίων, ὧν ὁ ΑΗ τοῦ ΔΕ ἐστι διπλασίων,
καὶ λοιπὸς ἄρα ὁ ΗΓ λοιποῦ τοῦ ΕΓ ἐστι διπλασίων;
δίθα ἄρα τέτμ ηται ὁ ΗΓ τῷ Ε. ὁ ἄρα ἐκ τῶν ΗΒ, ΒΓ
μετὰ τοῦ ἀπὸ ΓΕ ἴσος ἐστὶ τῷ ἀπὸ ΒΕ τετραγώνῳ. 
ἀλλὰ καὶ
ὁ ἐκ τῶν ΑΒ, ΒΓ μετὰ τοῦ ἀπὸ ΓΕ ἴσος ὑπόκειται
τῷ ἀπὸ [τοῦ] ΒΕ τετραγώνῳ; ὁ ἄρα ἐκ τῶν ΗΒ, ΒΓ
μετὰ τοῦ ἀπὸ ΓΕ ἴσος ἐστὶ 
τῷ ἐκ τῶν ΑΒ, ΒΓ μετὰ
τοῦ ἀπὸ ΓΕ. καὶ κοινοῦ ἀφαιρεθέντος τοῦ ἀπὸ ΓΕ συνάγεται
ὁ ΑΒ ἴσος τῷ ΗΒ; ὅπερ ἄτοπον. οὐκ ἄρα ὁ ἐκ τῶν ΑΒ, ΒΓ
μετὰ τοῦ ἀπὸ [τοῦ] ΓΕ ἴσος ἐστὶ τῷ ἀπὸ ΒΕ. λέγω δή,
ὅτι οὐδὲ ἐλάσσων τοῦ ἀπὸ ΒΕ. εἰ γὰρ δυνατόν, ἔστω
τῷ ἀπὸ ΒΖ ἴσος, καὶ τοῦ ΔΖ διπλασίων ὁ ΘΑ. καὶ συναθθήσεται
πάλιν διπλασίων ὁ ΘΓ τοῦ ΓΖ; ὥστε
καὶ τὸν ΓΘ δίθα τετμῆσθαι κατὰ τὸ Ζ, καὶ διὰ τοῦτο
τὸν ἐκ τῶν ΘΒ, ΒΓ μετὰ τοῦ ἀπὸ ΖΓ ἴσον γίνεσθαι
τῷ ἀπὸ ΒΖ. ὑπόκειται δὲ καὶ ὁ ἐκ τῶν ΑΒ, ΒΓ
μετὰ τοῦ ἀπὸ ΓΕ ἴσος τῷ ἀπὸ ΒΖ. ὥστε καὶ ὁ ἐκ
τῶν ΘΒ, ΒΓ μετὰ τοῦ ἀπὸ ΓΖ ἴσος ἔσται τῷ ἐκ
τῶν ΑΒ, ΒΓ μετὰ τοῦ ἀπὸ ΓΕ; ὅπερ ἄτοπον. οὐκ
ἄρα ὁ ἐκ τῶν ΑΒ, ΒΓ μετὰ τοῦ ἀπὸ ΓΕ ἴσος
ἐστὶ [τῷ] ἐλάσσονι τοῦ ἀπὸ ΒΕ. ἐδείθθη δέ, ὅτι οὐδὲ
[α ὐτῷ] τῷ ἀπὸ ΒΕ. οὐκ ἄρα ὁ ἐκ τῶν ΑΒ, ΒΓ μετὰ
τοῦ ἀπὸ ΓΕ τετράγωνός ἐστιν. ὅπερ ἔδει δεῖχαι.}}

\ParallelRText{
\begin{center}
{\large Lemma II}
\end{center}

To find two square numbers such that
the sum of them is not square.

\epsfysize=0.4in
\centerline{\epsffile{Book10/fig028be.eps}}

For let the (number created) from (multiplying) $AB$ and $BC$, as we said, be square.  And (let) $CA$ (be) even. And let $CA$ be
cut in half at $D$. So it is clear that the square (number created) from
(multiplying) $AB$ and $BC$, plus the square on $CD$, is equal to the
square on $BD$ [see previous lemma]. Let the unit $DE$ be subtracted (from $BD$).
Thus, the (number created) from (multiplying) $AB$ and $BC$, plus
the (square) on $CE$, is less than the square on $BD$. I say, therefore,
that the square (number created) from (multiplying) $AB$ and $BC$, plus the
(square) on $CE$, is not square.

For if it is square, it is either equal to the (square) on $BE$, or
less than the (square) on $BE$, but cannot    any more be greater  (than the square on $BE$), lest the unit be divided. First of all, if possible, let the
(number created) from (multiplying) $AB$ and $BC$, plus the (square) on
$CE$, be equal to the (square) on $BE$. And let $GA$ be double the unit
$DE$. Therefore, since the whole of $AC$ is double the whole
of $CD$, of which $AG$ is double $DE$, the remainder $GC$
is thus double the remainder $EC$. Thus, $GC$ has been cut in half at $E$.
Thus, the (number created) from (multiplying) $GB$ and $BC$,
plus the (square) on $CE$, is equal to the square on $BE$ [Prop.~2.6]. But, the (number created) from (multiplying) $AB$ and $BC$, plus the (square) on $CE$, was also assumed
 (to be) equal to the square on $BE$. Thus, the (number created) from
 (multiplying) $GB$ and $BC$, plus the (square) on $CE$, is equal to
 the (number created) from (multiplying)  $AB$ and $BC$, plus the
 (square) on $CE$. And subtracting the (square) on $CE$ from both, $AB$ is  inferred (to be) equal to $GB$. The very thing is absurd. Thus, the (number created) from (multiplying) $AB$ and $BC$, plus the (square) on
$CE$, is not equal to the (square) on $BE$. So I say that (it is) not
less than the (square) on $BE$ either. For, if possible, let it be equal to
the (square) on $BF$. And (let) $HA$ (be) double $DF$. And it
can again be inferred that $HC$ (is) double $CF$. Hence, $CH$ has
also been cut in half at $F$. And, on account of this, the (number created)
from (multiplying) $HB$ and $BC$, plus the (square) on $FC$, becomes
equal to the (square) on $BF$ [Prop.~2.6]. 
And the (number created) from (multiplying) $AB$ and $BC$, plus the
(square) on $CE$, was also
assumed (to be) equal to the (square) on $BF$.
Hence,
the (number created) from (multiplying) $HB$ and $BC$, plus the
(square) on $CF$, will also be equal to the
(number created) from (multiplying) $AB$ and $BC$, plus the (square) on $CE$. The very thing is
absurd. Thus, the (number created) from (multiplying) $AB$ and $BC$,
plus the (square) on $CE$, is not equal to less than the
(square) on $BE$. And it was shown that  (is it) not equal to the (square)
on $BE$ either. Thus, the (number created) from (multiplying)
$AB$ and $BC$, plus the  square on $CE$, is not square. (Which is)
the very thing it was required to show.}
\end{Parallel}

%%%%%%
% Prop 10.29
%%%%%%
\pdfbookmark[1]{Proposition 10.29}{pdf10.29}
\begin{Parallel}{}{} 
\ParallelLText{
\begin{center}
{\large \ggn{29}.}
\end{center}\vspace*{-7pt}

{\greekfont Εὑρεῖν δύο ῥητὰς δυνάμει μόνον συμμέτρους, ὥστε τὴν μείζονα
τῆς ἐλάσσονος μεῖζον δύνασθαι τῷ ἀπὸ συμμέτρου ἑαυτῇ
μήκει.}\\~\\

\epsfysize=1.7in
\centerline{\epsffile{Book10/fig029g.eps}}

{\greekfont Ἐκκείσθω γάρ τις ῥητὴ ἡ ΑΒ καὶ δύο τετράγωνοι ἀριθμοὶ
οἱ ΓΔ, ΔΕ, ὥστε τὴν ὑπεροθὴν α ὐτῶν τὸν ΓΕ μὴ εἶναι
τετράγωνον, καὶ γεγράφθω ἐπὶ τῆς ΑΒ ἡμικύκλιον τὸ ΑΖΒ,
καὶ πεποιήσθω ὡς ὁ ΔΓ πρὸς τὸν ΓΕ, οὕτως τὸ ἀπὸ τῆς ΒΑ
τετράγωνον πρὸς τὸ ἀπὸ τῆς ΑΖ τετράγωνον, καὶ ἐπεζεύθθω ἡ ΖΒ.}

{\greekfont Ἐπεὶ [οὖν] ἐστιν ὡς τὸ ἀπὸ τῆς ΒΑ πρὸς τὸ ἀπὸ τῆς ΑΖ,
οὕτως ὁ ΔΓ πρὸς τὸν ΓΕ, τὸ ἀπὸ τῆς ΒΑ ἄρα πρὸς τὸ
ἀπὸ τῆς ΑΖ λόγον ἔχει, ὅν ἀριθμὸς ὁ ΔΓ πρὸς ἀριθμὸν
τὸν ΓΕ; σύμμετρον ἄρα ἐστὶ τὸ ἀπὸ τῆς ΒΑ τῷ ἀπὸ τῆς
ΑΖ. ῥητὸν δὲ τὸ ἀπὸ τῆς ΑΒ; ῥητὸν ἄρα καὶ τὸ ἀπὸ
τῆς ΑΖ; ῥητὴ ἄρα καὶ ἡ ΑΖ. καὶ ἐπεὶ ὁ ΔΓ πρὸς τὸν ΓΕ
λόγον οὐκ ἔχει, ὃν τετράγωνος ἀριθμὸς πρὸς τετράγωνον
ἀριθμόν,  οὐδὲ τὸ ἀπὸ τῆς ΒΑ ἄρα πρὸς τὸ ἀπὸ τῆς
ΑΖ λόγον ἔχει, ὃν τετράγωνος ἀριθμὸς πρὸς τετράγωνον
ἀριθμόν; ἀσύμμετρος ἄρα ἐστὶν ἡ ΑΒ τῇ ΑΖ μήκει; αἱ
ΒΑ, ΑΖ ἄρα ῥηταί εἰσι δυνάμει μόνον σύμμετροι. καὶ ἐπεί
[ἐστιν] ὡς ὁ ΔΓ πρὸς τὸν ΓΕ, οὕτως τὸ ἀπὸ τῆς ΒΑ πρὸς
τὸ ἀπὸ τῆς ΑΖ, ἀναστρέψαντι ἄρα ὡς ὁ ΓΔ πρὸς τὸν
ΔΕ, οὕτως τὸ ἀπὸ τῆς ΑΒ πρὸς τὸ ἀπὸ τῆς ΒΖ. ὁ δὲ
ΓΔ πρὸς τὸν ΔΕ λόγον ἔχει, ὃν τετράγωνος ἀριθμὸς
πρὸς τετράγωνον ἀριθμόν; καὶ τὸ ἀπὸ τῆς ΑΒ ἄρα πρὸς τὸ
ἀπὸ τῆς ΒΖ λόγον ἔχει, ὃν τετράγωνος ἀριθμὸς πρὸς τετράγωνον
ἀριθμόν;
σύμμετρος ἄρα ἐστὶν ἡ ΑΒ
τῇ ΒΖ μήκει. καί ἐστι τὸ ἀπὸ τῆς ΑΒ ἴσον τοῖς ἀπὸ τῶν ΑΖ, ΖΒ; ἡ ΑΒ ἄρα τῆς ΑΖ μεῖζον δύναται τῇ
ΒΖ συμμέτρῳ ἑαυτῇ.}

{\greekfont Εὕρηνται ἄρα δύο ῥηταὶ δυνάμει μόνον σύμμετροι αἱ ΒΑ, ΑΖ,
ὥστε τὴν μεῖζονα τὴν ΑΒ τῆς ἐλάσσονος τῆς ΑΖ μεῖζον δύνασθαι
τῷ ἀπὸ τῆς ΒΖ συμμέτρου ἑαυτῇ μήκει; ὅπερ ἔδει
δεῖχαι.}}

\ParallelRText{
\begin{center}
{\large Proposition 29}
\end{center}

To  find two rational (straight-lines which are)
commensurable in square only, such that the square on the greater is larger
than the (square on the) lesser by the (square) on (some straight-line which is) commensurable in length with the greater.

\epsfysize=1.7in
\centerline{\epsffile{Book10/fig029e.eps}}

For let some rational (straight-line) $AB$ be laid down, and two
square numbers, $CD$ and $DE$, such that the difference between them, $CE$, is not square [Prop.~10.28 lem.~I]. 
And let the semi-circle $AFB$ be drawn on $AB$. And let
it be contrived that as $DC$ (is) to $CE$, so the square on $BA$
(is) to the square on $AF$ [Prop.~10.6~corr.].
And let $FB$ be joined.

\mbox{[}Therefore,] since as the (square) on $BA$ is to the (square) on $AF$,
so $DC$ (is) to $CE$, the (square) on $BA$ thus has to the (square)
on $AF$ the ratio which the number $DC$ (has) to the number $CE$.
Thus, the (square) on $BA$ is commensurable with the (square) on $AF$
[Prop.~10.6]. And the (square) on $AB$ (is)
rational [Def.~10.4]. Thus, the (square) on $AF$
(is) also rational. Thus, $AF$ (is) also rational. 
And since $DC$ does not have to $CE$ the ratio which
(some) square number (has) to (some) square number,  the (square)
on $BA$ thus does not have to the (square) on $AF$ the ratio which (some)
square number has to (some) square number either. Thus, $AB$ is incommensurable in length with $AF$ [Prop.~10.9].
Thus, the rational (straight-lines) $BA$
and $AF$ are commensurable in square only. And since as $DC$ [is] to $CE$, so the
(square) on $BA$ (is) to the (square) on $AF$, thus, via conversion,
as $CD$ (is) to $DE$, so the (square) on $AB$ (is) to the (square) on
$BF$ [Props.~5.19~corr., 3.31, 1.47]. And $CD$
has to $DE$ the ratio which (some) square number (has) to (some)
square number. Thus, the (square) on $AB$ also has to the (square) on
$BF$ the ratio which (some) square number has to (some) square number.
$AB$ is thus commensurable in length with $BF$ [Prop.~10.9]. And the (square) on $AB$
is equal to the (sum of the squares) on $AF$ and $FB$ [Prop.~1.47]. Thus, the
square on $AB$ is greater than (the square on) $AF$ by (the square on)
$BF$, (which is) commensurable (in length) with ($AB$).

Thus, two rational (straight-lines), $BA$ and $AF$, commensurable in square only, be found such that the square on the greater, $AB$, is larger
than (the square on) the lesser, $AF$, by the (square) on $BF$, (which is)
commensurable in length with ($AB$).$^\dag$ (Which is) the very thing it
was required to show.}
\end{Parallel}
{\footnotesize\noindent$^\dag$ $BA$ and $AF$
have lengths $1$ and $\sqrt{1-k^2}$ times that of $AB$, respectively, where $k=\sqrt{DE/CD}$.}

%%%%%%
% Prop 10.30
%%%%%%
\pdfbookmark[1]{Proposition 10.30}{pdf10.30}
\begin{Parallel}{}{} 
\ParallelLText{
\begin{center}
{\large \ggn{30}.}
\end{center}\vspace*{-7pt}

{\greekfont Εὑρεῖν δύο ῥητὰς δυνάμει μόνον συμμέτρους, ὥστε τὴν
μείζονα τῆς ἐλάσσονος μεῖζον δύνασθαι τῷ ἀπὸ ἀσυμμέτρου
ἑαυτῇ μήκει.}\\~\\

\epsfysize=1.7in
\centerline{\epsffile{Book10/fig030g.eps}}

{\greekfont Ἐκκείσθω ῥητὴ ἡ ΑΒ καὶ δύο τετράγωνοι ἀριθμοὶ
οἱ ΓΕ, ΕΔ, ὥστε τὸν συγκείμενον ἐχ α ὐτῶν τὸν ΓΔ μὴ
εἶναι τετράγωνον, καὶ γεγράφθω ἐπὶ τῆς ΑΒ ἡμικύκλιον
τὸ ΑΖΒ, καὶ πεποιήσθω ὡς ὁ ΔΓ πρὸς τὸν ΓΕ, οὕτως
τὸ ἀπὸ τῆς ΒΑ πρὸς τὸ ἀπὸ τῆς ΑΖ, καὶ ἐπεζεύθθω ἡ ΖΒ.}

{\greekfont Ὁμοίως δὴ δείχομεν τῷ πρὸ τούτου, ὅτι αἱ ΒΑ, ΑΖ ῥηταί
εἰσι δυνάμει μόνον σύμμετροι. καὶ ἐπεί ἐστιν ὡς ὁ
ΔΓ πρὸς τὸν ΓΕ, οὕτως τὸ ἀπὸ τῆς ΒΑ πρὸς τὸ ἀπὸ τῆς
ΑΖ, ἀναστρέψαντι ἄρα ὡς ὁ ΓΔ πρὸς τὸν ΔΕ, οὕτως τὸ ἀπὸ
τῆς ΑΒ πρὸς τὸ ἀπὸ τῆς ΒΖ. ὁ δὲ ΓΔ πρὸς τὸν ΔΕ λόγον
οὐκ ἔχει, ὃν τετράγωνος ἀριθμὸς πρὸς τετράγωνον ἀριθμόν;
οὐδ ἄρα τὸ ἀπὸ τῆς ΑΒ πρὸς τὸ ἀπὸ τῆς ΒΖ λόγον
ἔχει, ὃν τετράγωνος ἀριθμὸς πρὸς τετράγωνον ἀριθμόν;
ἀσύμμετρος ἄρα ἐστὶν ἡ ΑΒ τῇ ΒΖ μήκει. καὶ δύναται
ἡ ΑΒ τῆς ΑΖ μεῖζον τῷ ἀπὸ τῆς ΖΒ ἀσυμμέτρου
ἑαυτῇ.}

{\greekfont Αἱ ΑΒ, ΑΖ ἄρα ῥηταί εἰσι δυνάμει μόνον σύμμετροι, καὶ
ἡ ΑΒ τῆς ΑΖ μεῖζον δύναται τῷ ἀπὸ τῆς ΖΒ
ἀσυμμέτρου ἑαυτῇ μήκει; ὅπερ ἔδει δεῖχαι.}}

\ParallelRText{
\begin{center}
{\large Proposition 30}
\end{center}

To find two rational (straight-lines which are)
commensurable in square only, such that the square on the greater
is larger than the (the square on)  lesser by the (square) on (some
straight-line which is) incommensurable in length with the greater.

\epsfysize=1.7in
\centerline{\epsffile{Book10/fig030e.eps}}

Let the rational (straight-line) $AB$ be laid out, and the two square
numbers, $CE$ and $ED$, such that the sum of them, $CD$,
is not square [Prop.~10.28~lem.~II].
And let the semi-circle $AFB$ be drawn on $AB$. And let it
be contrived that as $DC$ (is) to $CE$, so the (square) on
$BA$ (is) to the (square) on $AF$ [Prop.~10.6~corr].
And let $FB$ be joined.

So, similarly to the (proposition) before this, we can show that $BA$ and
$AF$ are rational (straight-lines which are) commensurable in square only.
And since as $DC$ is to $CE$, so the (square) on $BA$ (is) to the (square)
on $AF$, thus, via conversion, as $CD$ (is) to $DE$, so the (square) on
$AB$ (is) to the (square) on $BF$ [Props.~5.19~corr., 3.31, 1.47]. And
$CD$ does not have to $DE$ the ratio which (some) square number
(has) to (some) square number. Thus, the (square) on $AB$
does not have to the (square) on $BF$ the ratio which (some)
square number has to (some) square number either. Thus, $AB$ is
incommensurable in length with $BF$ [Prop.~10.9].
And the square on $AB$ is greater than the (square on) $AF$ by the
(square) on $FB$  [Prop.~1.47], (which is) incommensurable (in length) with ($AB$).

Thus, $AB$ and $AF$ are rational (straight-lines which are) commensurable
in square only, and the square on $AB$ is greater than (the square on) $AF$
by the (square) on $FB$, (which is) incommensurable (in length) with
 $(AB)$.$^\dag$ (Which is) the very thing it was required to show.}
\end{Parallel}
{\footnotesize\noindent$^\dag$ $AB$ and $AF$ have lengths $1$ and $1/\sqrt{1+k^2}$
times that of $AB$, respectively, where $k=\sqrt{DE/CE}$.}

%%%%%%
% Prop 10.31
%%%%%%
\pdfbookmark[1]{Proposition 10.31}{pdf10.31}
\begin{Parallel}{}{} 
\ParallelLText{
\begin{center}
{\large \ggn{31}.}
\end{center}\vspace*{-7pt}

{\greekfont Εὑρεῖν δύο μέσας δυνάμει μόνον συμμέτρους ῥητὸν περιεθούσας,
ὥστε τὴν μείζονα τῆς ἐλάσσονος μεῖζον δύνασθαι τῷ ἀπὸ
συμμέτρου ἑαυτῇ μήκει.}\\~\\

\epsfysize=1.6in
\centerline{\epsffile{Book10/fig031g.eps}}

{\greekfont Ἐκκείσθωσαν δύο ῥηταὶ δυνάμει μόνον σύμμετροι αἱ Α, Β,
ὥστε τὴν Α μείζονα οὖσαν τῆς ἐλάσσονος τῆς Β μεῖζον
δύνασθαι τῷ ἀπὸ συμμέτρου ἑαυτῇ μήκει. καὶ τῷ ὑπὸ
τῶν Α, Β ἴσον ἔστω τὸ ἀπὸ τῆς Γ. μέσον δὲ τὸ ὑπὸ
τῶν Α, Β; μέσον ἄρα καὶ τὸ ἀπὸ τῆς Γ; μέση ἄρα καὶ ἡ Γ.
τῷ δὲ ἀπὸ τῆς Β ἴσον ἔστω τὸ ὑπὸ τῶν Γ, Δ; ῥητὸν
δὲ τὸ ἀπὸ τῆς Β; ῥητὸν ἄρα καὶ τὸ ὑπὸ τῶν Γ, Δ. καὶ
ἐπεί ἐστιν ὡς ἡ Α πρὸς τὴν Β, οὕτως
τὸ ὑπὸ τῶν Α, Β πρὸς τὸ ἀπὸ τῆς Β, ἀλλὰ τῷ μὲν
ὑπὸ τῶν Α, Β ἴσον ἐστὶ τὸ ἀπὸ τῆς Γ, τῷ δὲ ἀπὸ
τῆς Β ἴσον τὸ ὑπὸ τῶν Γ, Δ, ὡς ἄρα ἡ Α πρὸς τὴν Β,
οὕτως τὸ ἀπὸ τῆς Γ πρὸς τὸ ὑπὸ τῶν Γ, Δ. ὡς δὲ
τὸ ἀπὸ τῆς Γ πρὸς τὸ ὑπὸ τῶν Γ, Δ, οὕτως ἡ Γ πρὸς τὴν
Δ; καὶ ὡς ἄρα ἡ Α πρὸς τὴν Β, οὕτως ἡ Γ πρὸς τὴν Δ.
 σύμμετρος δὲ ἡ Α τῇ Β δυνάμει μόνον; σύμμετρος
ἄρα καὶ ἡ Γ τῇ Δ δυνάμει μόνον. καί ἐστι μέση ἡ Γ;
μέση ἄρα καὶ ἡ Δ. καὶ ἐπεί ἐστιν ὡς ἡ Α πρὸς τὴν
Β, ἡ Γ πρὸς τὴν Δ, ἡ δὲ Α τῆς Β μεῖζον δύναται τῷ
ἀπὸ συμμέτρου ἑαυτῇ, καὶ ἡ Γ ἄρα τῆς Δ μεῖζον
δύναται τῷ ἀπὸ συμμέτρου ἑαυτῇ.}

{\greekfont Εὕρηνται ἄρα δύο μέσαι δυνάμει μόνον σύμμετροι
αἱ Γ, Δ ῥητὸν περιέθουσαι, καὶ ἡ Γ τῆς Δ μεῖζον δυνάται
τῷ ἀπὸ συμμέτρου ἑαυτῇ μήκει.}

{\greekfont Ὁμοίως δὴ δειθθήσεται καὶ τῷ ἀπὸ ἀσυμμέτρου,
ὅταν ἡ Α τῆς Β μεῖζον δύνηται τῷ ἀπὸ ἀσυμμέτρου
ἑαυτῇ.}}

\ParallelRText{
\begin{center}
{\large Proposition 31}
\end{center}

To find two medial (straight-lines),
commensurable in square only, (and) containing a rational (area), such that
the square on the greater is larger than the (square on the) lesser by the
(square) on (some straight-line) commensurable in length with the greater.

\epsfysize=1.6in
\centerline{\epsffile{Book10/fig031e.eps}}

Let two rational (straight-lines), $A$ and $B$, commensurable in square only, 
be laid out, such that the square on the greater $A$ is larger than the (square on
the) lesser $B$ by the (square) on (some straight-line) commensurable
in length with ($A$) [Prop.~10.29]. And let the
(square) on $C$ be equal to the (rectangle contained) by $A$ and $B$.
And the (rectangle contained by) $A$ and $B$ (is) medial [Prop.~10.21]. Thus, the (square) on $C$ (is)
also medial. Thus, $C$ (is) also medial [Prop.~10.21]. And let the (rectangle contained) by $C$ and $D$
be equal to the (square) on $B$. And the (square)
on $B$ (is) rational. Thus, the (rectangle contained) by $C$ and $D$ (is)
also rational. And since as $A$ is to $B$, so the (rectangle contained) by
$A$ and $B$ (is) to the (square) on $B$ [Prop.~10.21~lem.], but the (square) on $C$
is equal to the (rectangle contained) by $A$ and $B$,  
and the (rectangle contained) by $C$ and $D$ to the (square) on
$B$, thus as $A$ (is) to $B$,
so the (square) on $C$ (is) to the (rectangle contained) by $C$ and $D$.
And as the (square) on $C$ (is) to the (rectangle contained) by $C$ and $D$,
so $C$ (is) to $D$ [Prop.~10.21~lem.].
And thus as $A$ (is) to $B$, so $C$ (is) to $D$. And $A$ is commensurable
in square only with $B$. Thus, $C$ (is) also commensurable in square
only with $D$ [Prop.~10.11]. And $C$ is medial.
Thus, $D$ (is) also medial [Prop.~10.23]. 
And since as $A$ is to $B$, (so) $C$ (is) to $D$, and the square on $A$
is greater than (the square on) $B$ by the (square) on (some straight-line)
commensurable (in length) with ($A$), the square on $C$
is thus also greater than (the square on) $D$ by the (square) on (some straight-line)
commensurable (in length) with ($C$) [Prop.~10.14]. 

Thus, two medial (straight-lines), $C$ and $D$, 
commensurable in square only, (and) containing a rational (area), be found. And the square on $C$ is greater than (the square on) $D$ by the
(square) on (some straight-line) commensurable in length with ($C$).$^\dag$

So, similarly, (the proposition) can also be demonstrated for  (some straight-line) incommensurable
(in length with $C$), provided that the square on $A$ is greater than (the square on $B$) by the (square) on (some straight-line) incommensurable (in length) with ($A$) [Prop.~10.30].$^\ddag$}
\end{Parallel}
{\footnotesize\noindent$^\dag$ $C$ and $D$ have lengths $(1-k^2)^{1/4}$ and $(1-k^2)^{3/4}$ times that of $A$, respectively, where $k$ is defined in the
footnote to Prop.~10.29.\\[0.5ex]
$^\ddag$ $C$ and $D$ would have lengths $1/(1+k^2)^{1/4}$ and $1/(1+k^2)^{3/4}$ times that of $A$, respectively, where $k$ is defined in the footnote to  Prop.~10.30.}

%%%%%%
% Prop 10.32
%%%%%%
\pdfbookmark[1]{Proposition 10.32}{pdf10.32}
\begin{Parallel}{}{} 
\ParallelLText{
\begin{center}
{\large \ggn{32}.}
\end{center}\vspace*{-7pt}

{\greekfont Εὑρεῖν δύο μέσας δυνάμει μόνον συμμέτρους μέσον περιεθούσας,
ὥστε τὴν μείζονα τῆς ἐλάσσονος μεῖζον δύνασθαι τῷ ἀπὸ
συμμέτρου ἑαυτῇ.}\\~\\

\epsfysize=0.75in
\centerline{\epsffile{Book10/fig032g.eps}}

{\greekfont Ἐκκείσθωσαν τρεῖς ῥηταὶ δυνάμει μόνον σύμμετροι αἱ Α, Β, Γ,
ὥστε τὴν Α τῆς Γ μεῖζον δύνασθαι τῷ ἀπὸ συμμέτρου
ἑαυτῇ, καὶ τῷ μὲν ὑπὸ τῶν Α, Β ἴσον ἔστω τὸ ἀπὸ
τὴς Δ. μέσον ἄρα τὸ ἀπὸ τῆς Δ; καὶ ἡ Δ ἄρα μέση ἐστίν.
τῷ δὲ ὑπὸ τῶν Β, Γ ἴσον ἔστω τὸ ὑπὸ τῶν Δ, Ε. καὶ
ἐπεί ἐστιν ὡς τὸ ὑπὸ τῶν Α, Β πρὸς τὸ ὑπὸ τῶν Β, Γ, οὕτως
ἡ Α πρὸς τὴν Γ, ἀλλὰ τῷ μὲν ὑπὸ τῶν Α, Β ἴσον ἐστὶ τὸ
ἀπὸ τῆς Δ, τῷ δὲ ὑπὸ τῶν Β, Γ ἴσον τὸ ὑπὸ τῶν Δ, Ε, ἔστιν
ἄρα ὡς ἡ Α πρὸς τὴν Γ, οὕτως τὸ ἀπὸ τῆς Δ πρὸς
τὸ ὑπὸ τῶν Δ, Ε. ὡς δὲ τὸ ἀπὸ τῆς Δ πρὸς τὸ ὑπὸ
τῶν Δ, Ε, οὕτως ἡ Δ πρὸς τὴν Ε; καὶ ὡς ἄρα ἡ Α πρὸς τὴν
Γ, οὕτως ἡ Δ πρὸς τὴν Ε. σύμμετρος δὲ ἡ Α τῇ Γ δυνάμει [μόνον].
σύμμετρος ἄρα καὶ ἡ Δ τῇ Ε δυνάμει μόνον. μέση δὲ ἡ Δ;
μέση ἄρα καὶ ἡ Ε. καὶ ἐπεί ἐστιν ὡς ἡ Α πρὸς τὴν Γ, ἡ Δ πρὸς
τὴν Ε, ἡ δὲ Α τῆς Γ μεῖζον δύναται τῷ ἀπὸ συμμέτρου
ἑαυτῇ, καὶ ἡ Δ ἄρα τῆς Ε μεῖζον δυνήσεται τῷ ἀπὸ
συμμέτρου ἑαυτῇ. λέγω δή, ὅτι καὶ μέσον ἐστὶ τὸ ὑπὸ
τῶν Δ, Ε. ἐπεὶ γὰρ ἴσον ἐστὶ τὸ ὑπὸ τῶν
Β, Γ τῷ ὑπὸ τῶν Δ, Ε, μέσον δὲ τὸ ὑπὸ τῶν Β, Γ [αἱ γὰρ
Β, Γ ῥηταί εἰσι δυνάμει μόνον σύμμετροι], μέσον ἄρα καὶ
τὸ ὑπὸ τῶν Δ, Ε.}

{\greekfont Εὕρηνται ἄρα δύο μέσαι δυνάμει μόνον σύμμετροι αἱ Δ, Ε
μέσον περιέθουσαι, ὥστε τὴν μείζονα τῆς ἐλάσσονος
μεῖζον δύνασθαι τῷ ἀπὸ συμμέτρου ἑαυτῇ.}

{\greekfont Ὁμοίως δὴ πάλιν διεθθήσεται καὶ τῷ ἀπὸ ἀσυμμέτρ-ου, ὅταν
ἡ Α τῆς Γ μεῖζον δύνηται τῷ ἀπὸ ἀσυμμέτρου ἑαυτῃ.}}

\ParallelRText{
\begin{center}
{\large Proposition 32}
\end{center}

To find two medial (straight-lines), commensurable in square only, (and) containing a medial (area), such that
the square on the greater is larger than the (square on the) lesser by the
(square) on (some straight-line) commensurable (in length) with the greater.

\epsfysize=0.75in
\centerline{\epsffile{Book10/fig032e.eps}}

Let three rational (straight-lines), $A$, $B$ and $C$, commensurable
in square only, be laid out such that the square on $A$ is greater
than (the square on $C$) by the (square) on (some straight-line)
commensurable (in length) with ($A$) [Prop.~10.29].
And let the (square) on $D$ be equal to the (rectangle contained) by $A$ and
$B$. Thus, the (square) on $D$ (is) medial. Thus, $D$ is also
medial [Prop.~10.21].  And let the 
(rectangle contained) by $D$ and $E$ be equal to the (rectangle contained)
by $B$ and $C$. And since as the (rectangle contained) by $A$ and $B$
is to the (rectangle contained) by $B$ and $C$, so $A$ (is) to $C$
[Prop.~10.21~lem.], but the (square) on $D$
is equal to the (rectangle contained) by $A$ and $B$, and the
(rectangle contained) by $D$ and $E$ to the (rectangle contained) by
$B$ and $C$, thus as $A$ is to $C$, so the (square) on $D$ (is) to
the (rectangle contained) by $D$ and $E$. And as the (square) on $D$
(is) to the (rectangle contained) by $D$ and $E$, so $D$ (is) to $E$ [Prop.~10.21~lem.]. And thus as
$A$ (is) to $C$, so $D$ (is) to $E$. And $A$ (is) commensurable in square
[only] with $C$. Thus, $D$ (is) also commensurable in square only
with $E$ [Prop.~10.11]. And $D$ (is) medial.
Thus, $E$ (is) also medial [Prop.~10.23]. And
since as $A$ is to $C$, (so) $D$ (is) to $E$, and the square on $A$
is greater than (the square on) $C$ by the (square)
on (some straight-line) commensurable (in length) with ($A$),  the square on $D$
will thus also be greater than (the square on) $E$ by the (square) on (some straight-line)
commensurable (in length) with ($D$) [Prop.~10.14].
So, I also say that the (rectangle contained) by $D$ and $E$ is medial.
For since the (rectangle contained) by $B$ and $C$ is equal to the
(rectangle contained) by $D$ and $E$, and the (rectangle contained)
by $B$ and $C$ (is) medial [for $B$ and $C$ are rational
(straight-lines which are) commensurable in square only] [Prop.~10.21], the (rectangle contained) by $D$ and
$E$ (is) thus also medial.

Thus, two medial (straight-lines), $D$ and $E$, commensurable in square only, (and) containing a medial (area), be found such that
the square on the greater is larger than the (square on the) lesser by the
(square) on (some straight-line) commensurable (in length) with the greater.$^\dag$.

So, similarly, (the proposition) can again also be demonstrated for
(some straight-line) incommensurable (in length with the greater), provided that
the square on $A$ is greater than (the square on) $C$ by the (square)
on (some straight-line) incommensurable (in length) with ($A$) [Prop.~10.30].$^\ddag$}
\end{Parallel}
{\footnotesize\noindent $^\dag$ $D$ and $E$ have lengths $k'^{1/4}$ and $k'^{1/4}\sqrt{1-k^2}$ times that of $A$, respectively, where the length of $B$ is $k'^{1/2}$ times that of $A$, and $k$ is defined in the footnote to
Prop.~10.29.\\[0.5ex]
$^\ddag$ $D$ and $E$ would have lengths $k'^{1/4}$ and $k'^{1/4}/\sqrt{1+k^2}$ times that of $A$, respectively, where the length of $B$ is $k'^{1/2}$ times that of $A$, and $k$ is defined in the footnote to
Prop.~10.30.}

%%%%%%
% Prop 10.32a
%%%%%%
\begin{Parallel}{}{} 
\ParallelLText{
\begin{center}
{\large {\greekfont Λῆμμα}.}
\end{center}\vspace*{-7pt}

{\greekfont Ἔστω τρίγωνον ὀρθογώνιον τὸ ΑΒΓ ὀρθὴν ἔθον τὴν Α, καὶ
ἤθθω κάθετος ἡ ΑΔ; λέγω, ὅτι τὸ μὲν ὑπὸ τῶν
ΓΒΔ ἴσον ἐστὶ τῷ ἀπὸ τῆς ΒΑ, τὸ δὲ ὑπὸ τῶν ΒΓΑ
ἴσον τῷ ἀπὸ τῆς ΓΑ, καὶ τὸ ὑπὸ τῶν ΒΔ, ΔΓ ἴσον τῷ
ἀπὸ τῆς ΑΔ, καὶ ἔτι τὸ ὑπὸ τῶν ΒΓ, ΑΔ ἴσον [ἐστὶ]
τῷ ὑπὸ τῶν ΒΑ, ΑΓ.}

{\greekfont Καὶ πρῶτον, ὅτι τὸ ὑπὸ τῶν ΓΒΔ ἴσον [ἐστὶ] τῷ
ἀπὸ τῆς ΒΑ.}\\~\\~\\

\epsfysize=2in
\centerline{\epsffile{Book10/fig032ag.eps}}

{\greekfont Ἐπεὶ γὰρ ἐν ὀρθογωνίῳ τριγώνῳ ἀπὸ τῆς ὀρθῆς γωνίας
ἐπὶ τὴν βάσιν κάθετος ἦκται ἡ ΑΔ, τὰ ΑΒΔ, ΑΔΓ ἄρα τρίγωνα
ὅμοιά ἐστι τῷ τε ὅλῳ τῷ ΑΒΓ καὶ ἀλλήλοις. καὶ ἐπεὶ
ὅμοιόν ἐστι τὸ ΑΒΓ τρίγωνον τῷ ΑΒΔ τριγώνῳ, ἔστιν
ἄρα ὡς ἡ ΓΒ πρὸς τὴν ΒΑ, οὕτως ἡ ΒΑ πρὸς τὴν ΒΔ;
τὸ ἄρα ὑπὸ τῶν ΓΒΔ ἴσον ἐστὶ τῷ ἀπὸ τῆς ΑΒ.}

{\greekfont Διὰ τὰ α ὐτὰ δὴ καὶ τὸ ὑπὸ τῶν ΒΓΔ ἴσον ἐστὶ τῷ
ἀπὸ τῆς ΑΓ.}

{\greekfont Καὶ ἐπεί, ἐὰν ἐν ὀρθογωνίῳ τριγώνῳ ἀπὸ τῆς ὀρθῆς
γωνίας ἐπὶ τὴν βάσιν κάθετος ἀθθῇ, ἡ ἀθθεῖσα τῶν τῆς
βάσεως τμ ημάτων μέση ἀνάλογόν ἐστιν, ἔστιν ἄρα ὡς ἡ
ΒΑ πρὸς τὴν ΔΑ, οὕτως ἡ ΑΔ πρὸς τὴν ΔΓ; τὸ ἄρα ὑπὸ
τῶν ΒΔ, ΔΓ ἴσον ἐστὶ τῷ ἀπὸ τῆς ΔΑ.}

{\greekfont Λέγω, ὅτι καὶ τὸ ὑπὸ τῶν ΒΓ, ΑΔ ἴσον ἐστὶ τῷ ὑπὸ
τῶν ΒΑ, ΑΓ. ἐπεὶ γὰρ, ὡς ἔφαμεν, ὅμοιόν ἐστι
τὸ ΑΒΓ τῷ ΑΒΔ, ἔστιν ἄρα ὡς ἡ ΒΓ πρὸς τὴν ΓΑ,
οὕτως ἡ ΒΑ πρὸς τὴν ΑΔ. τὸ ἄρα ὑπὸ τῶν ΒΓ, ΑΔ ἴσον
ἐστὶ τῷ ὑπὸ τῶν ΒΑ, ΑΓ; ὅπερ ἔδει δεῖχαι.}}

\ParallelRText{
\begin{center}
{\large Lemma}
\end{center}

Let $ABC$ be a right-angled triangle having
the (angle) $A$ a right-angle. And let the perpendicular $AD$ be
drawn. I say that the (rectangle contained) by $CBD$ is equal to the
(square) on $BA$, and the (rectangle contained) by $BCD$ (is)
equal to the (square) on $CA$, and the (rectangle contained) by
$BD$ and $DC$ (is) equal to the (square) on $AD$, and, further, the
(rectangle contained) by $BC$ and $AD$ [is] equal to the (rectangle
contained) by $BA$ and $AC$.

And, first of all, (let us prove) that the (rectangle contained) by
$CBD$ [is] equal to the (square) on $BA$.

\epsfysize=2in
\centerline{\epsffile{Book10/fig032ae.eps}}

For since $AD$ has been drawn from the right-angle  in a right-angled triangle,
perpendicular to the base, $ABD$ and $ADC$ are thus
triangles (which are) similar to the whole, $ABC$, and to one another
[Prop.~6.8]. And since triangle $ABC$ is
similar to triangle $ABD$, thus as $CB$ is to $BA$, so $BA$ (is)
to $BD$ [Prop.~6.4]. Thus, the
(rectangle contained) by $CBD$ is equal to the (square) on $AB$
[Prop.~6.17].

So, for the same (reasons), the (rectangle contained) by $BCD$
is also equal to the (square) on $AC$.

And since if a (straight-line) is drawn from the right-angle in a right-angled
triangle,  perpendicular to the base, the (straight-line so) drawn
is the mean proportional to the pieces of the base [Prop.~6.8~corr.], thus as $BD$ is to $DA$, so $AD$ (is) to $DC$. Thus, the (rectangle contained) by
$BD$ and $DC$ is equal to the (square) on $DA$ [Prop.~6.17]. 

I also say that the (rectangle contained) by $BC$ and
$AD$ is equal to the (rectangle contained) by $BA$ and $AC$.
For since,  as we said, $ABC$ is similar to $ABD$, thus as $BC$
is to $CA$, so $BA$ (is) to $AD$ [Prop.~6.4].
Thus, the (rectangle contained) by $BC$ and $AD$
is equal to the (rectangle contained) by $BA$ and $AC$ [Prop.~6.16].
(Which is) the very thing it was required to show.}
\end{Parallel}

%%%%%%
% Prop 10.33
%%%%%%
\pdfbookmark[1]{Proposition 10.33}{pdf10.33}
\begin{Parallel}{}{} 
\ParallelLText{
\begin{center}
{\large \ggn{33}.}
\end{center}\vspace*{-7pt}

{\greekfont Εὑρεῖν δύο εὐθείας δυνάμει ἀσυμμέτρους ποιούσας τὸ μὲν συγκείμενον ἐκ τῶν ἀπ α ὐτῶν τετραγώνων ῥητόν, τὸ
δ ὑπ α ὐτῶν μέσον.}

{\greekfont Ἐκκείσθωσαν δύο ῥηταὶ δυνάμει μόνον σύμμετροι αἱ
ΑΒ, ΒΓ, ὥστε τὴν μείζονα τὴν ΑΒ τῆς ἐλάσσονος
τῆς ΒΓ μείζον δύνασθαι τῷ ἀπὸ ἀσυμμέτρου ἑαυτῇ,
καὶ τετμήσθω ἡ ΒΓ δίθα κατὰ τὸ Δ, καὶ τῷ ἀφ
ὁποτέρας τῶν ΒΔ, ΔΓ ἴσον παρὰ τὴν ΑΒ παραβεβλήσθω
παραλληλόγραμμον ἐλλεῖπον εἴδει τετραγώνῳ, καὶ ἔστω τὸ ὑπὸ
τῶν ΑΕΒ, καὶ γεγράφθω ἐπὶ τῆς ΑΒ ημικύκλιον τὸ ΑΖΒ, καὶ
ἤθθω τῇ ΑΒ πρὸς ὀρθὰς ἡ ΕΖ, καὶ ἐπεζεύθθωσαν αἱ ΑΖ, ΖΒ.}\\~\\~\\

\epsfysize=1.in
\centerline{\epsffile{Book10/fig033g.eps}}

{\greekfont Καὶ ἐπεὶ [δύο] εὐθεῖαι ἄνισοί εἰσιν αἱ ΑΒ, ΒΓ, καὶ ἡ ΑΒ
τῆς ΒΓ μεῖζον δύναται τῷ ἀπὸ ἀσυμμέτρου ἑαυτῇ, τῷ
δὲ τετάρτῳ τοῦ ἀπὸ τῆς ΒΓ, τουτέστι τῷ ἀπὸ τῆς ἡμισείας
α ὐτῆς, ἴσον παρὰ τὴν ΑΒ παραβέβληται παραλληλόγραμμον
ἐλλεῖπον εἴδει τετραγώνῳ καὶ ποιεῖ τὸ ὑπὸ τῶν ΑΕΒ,
ἀσύμμετρος ἄρα ἐστὶν ἡ ΑΕ τῇ ΕΒ. καί ἐστιν ὡς ἡ ΑΕ πρὸς
ΕΒ, οὕτως τὸ ὑπὸ τῶν ΒΑ, ΑΕ πρὸς τὸ ὑπὸ τῶν ΑΒ, ΒΕ, ἴσον
δὲ τὸ μὲν ὑπὸ τῶν ΒΑ, ΑΕ τῷ ἀπὸ τῆς ΑΖ, τὸ δὲ ὑπὸ
τῶν ΑΒ, ΒΕ τῷ ἁπὸ τῆς
 ΒΖ; ἀσύμμετρον
ἄρα ἐστὶ τὸ ἀπὸ τῆς ΑΖ τῷ ἀπὸ τῆς ΖΒ; αἱ ΑΖ, ΖΒ
ἄρα δυνάμει εἰσὶν ἀσύμμετροι. καὶ ἐπεὶ ἡ ΑΒ ῥητή ἐστιν,
ῥητὸν ἄρα ἐστὶ καὶ τὸ ἀπὸ τῆς ΑΒ; ὥστε καὶ τὸ συγκείμενον
ἐκ τῶν ἀπὸ τῶν ΑΖ, ΖΒ ῥητόν ἐστιν. καὶ ἐπεὶ πάλιν τὸ
ὑπὸ τῶν ΑΕ, ΕΒ ἴσον ἐστὶ τῷ ἀπὸ τῆς ΕΖ, ὑπόκειται
δὲ τὸ ὑπὸ τῶν ΑΕ, ΕΒ καὶ τῷ ἀπὸ τῆς ΒΔ ἴσον, ἴση
ἄρα ἐστὶν ἡ ΖΕ τῇ ΒΔ; διπλῆ ἄρα ἡ ΒΓ  τὴς ΖΕ; ὥστε
καὶ τὸ ὑπὸ τῶν ΑΒ, ΒΓ σύμμετρόν ἐστι τῷ ὑπὸ τῶν
ΑΒ, ΕΖ. μέσον δὲ τὸ ὑπὸ τῶν ΑΒ, ΒΓ; μέσον ἄρα καὶ τὸ ὑπὸ
τῶν ΑΒ, ΕΖ. ἴσον δὲ τὸ ὑπὸ τῶν ΑΒ, ΕΖ τῷ ὑπὸ τῶν
ΑΖ, ΖΒ; μέσον ἄρα καὶ τὸ ὑπὸ τῶν ΑΖ, ΖΒ. ἐδείθθη
δὲ καὶ ῥητὸν τὸ
 συγκείμενον ἐκ τῶν ἀπ
α ὐτῶν τετραγώνων.}

{\greekfont Εὕρηνται ἄρα δύο εὐθεῖαι δυνάμει ἀσύμμετροι αἱ ΑΖ, ΖΒ ποιοῦσαι τὸ μὲν συγκείμενον ἐκ τῶν ἀπ α ὐτῶν τετραγώνων ῥητόν,
τὸ δὲ ὑπ α ὐτῶν μέσον; ὅπερ ἔδει δεῖχαι.}}

\ParallelRText{
\begin{center}
{\large Proposition 33}
\end{center}

To find two straight-lines (which are) incommensurable
in square, making the sum of the squares on them rational, and the (rectangle contained) by them medial.

Let the two rational (straight-lines) $AB$ and $BC$, (which are) commensurable in square only, be laid out such that the square on the greater, $AB$, is larger
than (the square on) the lesser, $BC$, by the (square) on (some straight-line which is) incommensurable (in length)
with ($AB$) [Prop.~10.30]. And let $BC$ be
cut in half at $D$. And let a parallelogram equal to the (square) on
either of $BD$ or $DC$, (and) falling short by a square figure, 
be applied to $AB$ [Prop.~6.28], and
let it be the (rectangle contained) by $AEB$. And let the semi-circle
$AFB$ be drawn on $AB$. And let $EF$ be
drawn at right-angles to $AB$. And let $AF$ and $FB$ be joined.

\epsfysize=1.in
\centerline{\epsffile{Book10/fig033e.eps}}

And since $AB$ and $BC$ are [two] unequal straight-lines, and the square
on $AB$ is greater than (the square on) $BC$ by the (square) on (some
straight-line which is) incommensurable (in length) with ($AB$). And a parallelogram,
equal to one quarter of the (square) on $BC$---that is to say, (equal) to the
(square) on half of it---(and) falling short by a square figure, has
been applied to $AB$, and  makes the (rectangle contained) by $AEB$.
$AE$ is thus incommensurable (in length) with $EB$ [Prop.~10.18]. And as $AE$ is to $EB$, so the
(rectangle contained) by $BA$ and $AE$ (is) to the (rectangle contained) by
$AB$ and $BE$.
 And the (rectangle contained) by $BA$ and $AE$ (is) equal to the (square) on $AF$, and the (rectangle contained) by $AB$ and $BE$ to the (square) on
$BF$ [Prop.~10.32~lem.]. The (square) on
$AF$ is thus incommensurable with the (square) on $FB$ [Prop.~10.11]. Thus, $AF$ and $FB$ are
incommensurable in square. And since $AB$ is rational,
the (square) on $AB$ is also rational. Hence, the sum of the
(squares) on $AF$ and $FB$ is also rational [Prop.~1.47]. And, again, since the (rectangle
contained) by $AE$ and $EB$ is equal to the (square) on $EF$, and
the (rectangle contained) by $AE$ and $EB$ 
was assumed (to be) equal to the (square) on $BD$, $FE$ is thus equal to $BD$. Thus, $BC$
is double $FE$. And hence the (rectangle contained) by $AB$ and $BC$ is commensurable with the (rectangle contained) by $AB$ and $EF$ [Prop.~10.6]. And the (rectangle contained) by
$AB$ and $BC$ (is) medial [Prop.~10.21]. 
Thus, the (rectangle contained) by $AB$ and $EF$ (is) also medial
[Prop.~10.23~corr.].  And the (rectangle contained)
by $AB$ and $EF$ (is) equal to the (rectangle contained) by $AF$ and $FB$
[Prop.~10.32~lem.]. Thus, the (rectangle
contained) by $AF$ and $FB$ (is) also medial. And the
sum of the squares on them was also shown (to be) rational.

Thus, the two straight-lines, $AF$ and $FB$, (which are)
incommensurable in  square, be found, making the
sum of the squares on them rational, and the (rectangle contained)
by them medial. (Which is) the very thing it was required to show.}
\end{Parallel}
{\footnotesize\noindent$^\dag$ $AF$ and $FB$ have lengths
$\sqrt{[1+k/(1+k^2)^{1/2}]/2}$ and $\sqrt{[1-k/(1+k^2)^{1/2}]/2}$
times that of $AB$, respectively, where $k$ is defined in the footnote
to Prop.~10.30.}

%%%%%%
% Prop 10.34
%%%%%%
\pdfbookmark[1]{Proposition 10.34}{pdf10.34}
\begin{Parallel}{}{} 
\ParallelLText{
\begin{center}
{\large \ggn{34}.}
\end{center}\vspace*{-7pt}

{\greekfont Εὑρεῖν δύο εὐθείας δυνάμει ἀσυμμέτρους ποιούσας τὸ μὲν
συγκείμενον ἐκ τῶν ἀπ α ὐτῶν τετραγώνων μέσον, τὸ δ
ὑπ α ὐτῶν ῥητόν.}

\epsfysize=1.in
\centerline{\epsffile{Book10/fig034g.eps}}

{\greekfont Ἐκκείσθωσαν δύο μέσαι δυνάμει μόνον σύμμετροι αἱ ΑΒ, ΒΓ ῥητὸν
περιέθουσαι τὸ ὑπ α ὐτῶν, ὥστε τὴν ΑΒ τῆς ΒΓ μεῖζον δύνασθαι
τῷ ἀπὸ ἀσυμμέτρου ἑαυτῇ, καὶ γεγράφθω ἐπὶ τῆς ΑΒ τὸ
ΑΔΒ ἡμικύκλιον, καὶ τετμήσθω ἡ ΒΓ δίθα κατὰ τὸ Ε, καὶ
παραβεβλήσθω παρὰ τὴν ΑΒ τῷ ἀπὸ τῆς ΒΕ ἴσον
παραλληλόγραμμον ἐλλεῖπον εἴδει τετραγώνῳ τὸ ὑπὸ τῶν ΑΖΒ;
ἀσύμμετρος ἄρα [ἐστὶν] ἡ ΑΖ τῇ ΖΒ μήκει. καὶ ἤθθω ἀπὸ
τοῦ Ζ τῇ ΑΒ πρὸς ὀρθὰς ἡ ΖΔ, καὶ ἐπεζεύθθωσαν αἱ ΑΔ, ΔΒ.}

{\greekfont Ἐπεὶ ἀσύμμετρός ἐστιν ἡ ΑΖ τῇ ΖΒ, ἀσύμμετρον
ἄρα ἐστὶ καὶ τὸ ὑπὸ τῶν ΒΑ, ΑΖ τῷ ὑπὸ τῶν
ΑΒ, ΒΖ. ἴσον δὲ τὸ μὲν ὑπὸ τῶν ΒΑ, ΑΖ τῷ ἀπὸ τῆς
ΑΔ, τὸ δὲ ὑπὸ τῶν ΑΒ, ΒΖ τῷ ἀπὸ τῆς ΔΒ; ἀσύμμετρον
ἄρα ἐστὶ καὶ τὸ ἀπὸ τῆς ΑΔ τῷ ἀπὸ τῆς ΔΒ.
καὶ ἐπεὶ μέσον ἐστὶ τὸ ἀπὸ τῆς ΑΒ, μέσον ἄρα καὶ τὸ
συγκείμενον ἐκ τῶν ἀπὸ τῶν ΑΔ, ΔΒ. καὶ ἐπεὶ διπλῆ
ἐστιν ἡ ΒΓ τῆς ΔΖ, διπλάσιον ἄρα καὶ τὸ
ὑπὸ τῶν ΑΒ, ΒΓ τοῦ ὑπὸ τῶν ΑΒ, ΖΔ. ῥητὸν δὲ τὸ ὑπὸ
τῶν ΑΒ, ΒΓ; ῥητὸν ἄρα καὶ τὸ ὑπὸ τῶν ΑΒ, ΖΔ. τὸ δὲ
ὑπὸ τῶν ΑΒ, ΖΔ ἴσον τῷ ὑπὸ
τῶν ΑΔ, ΔΒ; ὥστε καὶ τὸ ὑπὸ τῶν ΑΔ, ΔΒ ῥητόν
ἐστιν.}

{\greekfont Εὕρηνται ἄρα δύο εὐθεῖαι δυνάμει ἀσύμμετροι αἱ
ΑΔ, ΔΒ ποιοῦσαι τὸ [μὲν] συγκείμενον ἐκ τῶν ἀπ
α ὐτῶν τετραγώνων μέσον, τὸ δ ὑπ α ὐτῶν ῥητόν;
ὅπερ ἔδει δεῖχαι.}}

\ParallelRText{
\begin{center}
{\large Proposition 34}
\end{center}

To find two straight-lines (which are) incommensurable in square, making the sum of the squares on them medial,
and the (rectangle contained) by them rational.

\epsfysize=1.in
\centerline{\epsffile{Book10/fig034e.eps}}

Let the two medial (straight-lines) $AB$ and $BC$,
(which are) commensurable in square only, be laid out having the
(rectangle contained) by them rational, (and) such that the square on
$AB$ is greater than (the square on) $BC$ by the (square) on (some straight-line)
incommensurable (in length) with ($AB$) [Prop.~10.31].
And let the semi-circle $ADB$ be drawn on $AB$.
And let $BC$ be cut in half at $E$. And let a (rectangular) parallelogram
equal to the (square) on $BE$, (and) falling short by a square figure,
be applied to $AB$, (and let it be) the (rectangle
contained by) $AFB$ [Prop.~6.28].
Thus, $AF$ [is] incommensurable in length with $FB$
[Prop.~10.18]. And let 
$FD$ be drawn from $F$ at right-angles to $AB$.
And let $AD$ and $DB$ be joined.

Since $AF$ is incommensurable (in length) with $FB$, the (rectangle contained) by
$BA$ and $AF$ is thus also incommensurable with the (rectangle contained)
by $AB$ and $BF$ [Prop.~10.11].  And the
(rectangle contained) by $BA$ and $AF$ (is) equal to the (square) on $AD$, and the (rectangle contained) by $AB$ and $BF$ to the (square) on $DB$
[Prop.~10.32~lem.]. Thus, the (square) on
$AD$ is also incommensurable with the (square) on $DB$. And
since the (square) on $AB$ is medial, the sum of the (squares) on
$AD$ and $DB$ (is) thus also medial [Props.~3.31, 1.47]. And since $BC$ is double $DF$ [see previous proposition], the (rectangle
contained) by $AB$ and $BC$ (is) thus also double the (rectangle contained)
by $AB$ and $FD$. And the (rectangle contained) by $AB$ and $BC$ (is)
rational. Thus, the (rectangle contained) by $AB$ and $FD$
(is) also rational [Prop.~10.6, Def.~10.4]. And the (rectangle contained) by
$AB$ and $FD$ (is) equal to the (rectangle contained) by $AD$ and $DB$
[Prop.~10.32~lem.]. And hence the (rectangle
contained) by $AD$ and $DB$ is rational.

Thus, two straight-lines, $AD$ and $DB$, (which are) incommensurable
in square, be found, making the sum of the squares on them medial,
and the (rectangle contained) by them rational.$^\dag$
(Which is) the very thing
it was required to show.}
\end{Parallel}
{\footnotesize\noindent$^\dag$ $AD$ and $DB$ have lengths $\sqrt{[(1+k^2)^{1/2}+k]/[2\,(1+k^2)]}$
and  $\sqrt{[(1+k^2)^{1/2}-k]/[2\,(1+k^2)]}$ times that of $AB$, respectively, where $k$ is defined in the footnote
to Prop.~10.29.}

%%%%%%
% Prop 10.35
%%%%%%
\pdfbookmark[1]{Proposition 10.35}{pdf10.35}
\begin{Parallel}{}{} 
\ParallelLText{
\begin{center}
{\large \ggn{35}.}
\end{center}\vspace*{-7pt}

{\greekfont Εὑρεῖν δύο εὐθείας δυνάμει ἀσυμμέτρους ποιούσας τό τε συγκείμενον ἐκ τῶν ἀπ α ὐτῶν τετραγώνων μέσον καὶ
τὸ ὑπ α ὐτῶν μέσον καὶ ἔτι ἀσύμμετρον τῷ συγκειμένῳ
ἐκ τῶν ἀπ α ὐτῶν τετραγώνῳ.}

\epsfysize=1.in
\centerline{\epsffile{Book10/fig034g.eps}}

{\greekfont Ἐκκείσθωσαν δύο μέσαι δυνάμει μόνον σύμμετροι αἱ ΑΒ, ΒΓ
μέσον περιέθουσαι, ὥστε τὴν ΑΒ τῆς ΒΓ μεῖζον δύνασθαι τῷ
ἀπὸ ἀσυμμέτρου ἑαυτῇ, καὶ γεγράφθω ἐπὶ τῆς ΑΒ
ἡμικύκλιον τὸ ΑΔΒ, καὶ τὰ λοιπὰ γεγονέτω τοῖς ἐπάνω
ὁμοίως.}

{\greekfont Καὶ ἐπεὶ ἀσύμμετρός ἐστιν ἡ ΑΖ τῇ ΖΒ μήκει, ἀσύμμετρ-ός
ἐστι καὶ ἡ ΑΔ τῇ ΔΒ δυνάμει. καὶ ἐπεὶ μέσον ἐστὶ
τὸ ἀπὸ τῆς ΑΒ, μέσον ἄρα καὶ τὸ συγκείμενον ἐκ τῶν
ἀπὸ τῶν ΑΔ, ΔΒ. καὶ ἐπεὶ τὸ ὑπὸ τῶν ΑΖ, ΖΒ ἴσον
ἐστὶ τῷ ἀφ ἑκατέρας τῶν ΒΕ, ΔΖ, ἴση ἄρα
ἐστὶν ἡ ΒΕ τῇ ΔΖ; διπλῆ ἄρα ἡ ΒΓ τῆς ΖΔ; ὥστε
καὶ τὸ ὑπὸ τῶν ΑΒ, ΒΓ διπλάσιόν ἐστι τοῦ ὑπὸ τῶν
ΑΒ, ΖΔ. μέσον δὲ τὸ ὑπὸ τῶν ΑΒ, ΒΓ; μέσον ἄρα καὶ τὸ
ὑπὸ τῶν ΑΒ, ΖΔ. καί ἐστιν ἴσον τῷ ὑπὸ τῶν ΑΔ, ΔΒ;
μέσον ἄρα καὶ τὸ ὑπὸ τῶν ΑΔ, ΔΒ. καὶ ἐπεὶ ἀσύμμετρός
ἐστιν ἡ ΑΒ τῇ ΒΓ μήκει, σύμμετρος δὲ ἡ ΓΒ τῇ ΒΕ,
ἀσύμμετρος ἄρα καὶ ἡ ΑΒ τῇ ΒΕ μήκει; ὥστε καὶ τὸ ἀπὸ
τῆς ΑΒ τῷ ὑπὸ τῶν ΑΒ, ΒΕ ἀσύμμετρόν ἐστιν. ἀλλὰ
τῷ μὲν ἀπὸ τῆς ΑΒ ἴσα ἐστὶ τὰ ἀπὸ τῶν ΑΔ, ΔΒ, τῷ
δὲ ὑπὸ τῶν ΑΒ, ΒΕ ἴσον ἐστὶ τὸ ὑπὸ τῶν ΑΒ,
ΖΔ, τουτέστι τὸ ὑπὸ τῶν ΑΔ, ΔΒ; ἀσύμμετρον ἄρα
ἐστὶ τὸ συγκείμενον ἐκ τῶν ἀπὸ τῶν ΑΔ, ΔΒ
τῷ ὑπὸ τῶν ΑΔ, ΔΒ.}

{\greekfont Εὕρηνται ἄρα δύο εὐθεῖαι αἱ ΑΔ, ΔΒ δυνάμει ἀσύμμετροι
ποιοῦσαι τό τε συγκείμενον ἐκ τῶν ἀπ α ὐτῶν μέσον
καὶ τὸ ὑπ α ὐτῶν μέσον καὶ ἔτι ἀσύμμετρον τῷ
συγκειμένῳ ἐκ τῶν ἀπ α ὐτῶν τετραγώνων;
ὅπερ ἔδει δεῖχαι.}}

\ParallelRText{
\begin{center}
{\large Proposition 35}
\end{center}

To find two straight-lines (which are) incommensurable in square, making the sum of the squares on them medial,
and the (rectangle contained) by them medial, and, moreover,
incommensurable  with the sum of the squares on them.

\epsfysize=1.in
\centerline{\epsffile{Book10/fig034e.eps}}

Let the two medial (straight-lines) $AB$ and $BC$, (which are) commensurable in square only, be laid out containing a medial (area), such
that the square on $AB$ is greater than (the square on) $BC$ by the
(square) on (some straight-line) incommensurable (in length) with ($AB$) [Prop.~10.32].  And let the semi-circle
$ADB$ be drawn on $AB$. And let the remainder (of the figure)
be generated similarly to the above (proposition).

And since $AF$ is incommensurable in length with $FB$ [Prop.~10.18], $AD$
is also incommensurable in square with $DB$ [Prop.~10.11]. And since the (square) on $AB$ is
medial, the sum of the (squares) on $AD$ and $DB$ (is) thus also
medial [Props.~3.31, 1.47]. And since the (rectangle
contained) by $AF$ and $FB$ is equal to the (square) on each of $BE$ and
$DF$, $BE$ is thus equal to $DF$. Thus, $BC$ (is) double $FD$. And
hence the (rectangle contained) by $AB$ and $BC$ is double the
(rectangle) contained by $AB$ and $FD$. And the (rectangle contained)
by $AB$ and $BC$ (is) medial. Thus, the (rectangle contained) by $AB$
and $FD$ (is) also medial. And it is equal to the (rectangle contained) 
by $AD$ and $DB$ [Prop.~10.32~lem.]. Thus, the
(rectangle contained) by $AD$ and $DB$ (is) also medial.
And since $AB$ is incommensurable in length with $BC$, and $CB$
(is) commensurable (in length) with $BE$, $AB$ (is) thus also incommensurable
in length with $BE$ [Prop.~10.13]. And hence the (square) on $AB$ is also incommensurable with the (rectangle contained) by $AB$ and $BE$ [Prop.~10.11]. But the (sum of the squares) on $AD$ and
$DB$ is equal to the (square) on $AB$ [Prop.~1.47].
And the (rectangle contained) by $AB$ and $FD$---that is to say, the (rectangle contained) by $AD$ and $DB$---is equal to the
(rectangle contained) by $AB$ and $BE$. Thus, the sum of the (squares) on
$AD$ and $DB$ is incommensurable with the (rectangle contained) by
$AD$ and $DB$.

Thus, two straight-lines, $AD$ and $DB$, (which are) incommensurable
in square, be found, making the sum of the (squares) on them
medial, and the (rectangle contained) by them medial, and, moreover, 
incommensurable  with the sum of the squares on them.$^\dag$
 (Which is) the
very thing it was required to show.}
\end{Parallel}
{\footnotesize\noindent$^\dag$ $AD$ and $DB$ have lengths
$k'^{1/4}\sqrt{[1+k/(1+k^2)^{1/2}]/2}$ and  $k'^{1/4}\sqrt{[1-k/(1+k^2)^{1/2}]/2}$ times that of $AB$, respectively, where $k$ and $k'$
are defined in the footnote to Prop.~10.32.}

%%%%%%
% Prop 10.36
%%%%%%
\pdfbookmark[1]{Proposition 10.36}{pdf10.36}
\begin{Parallel}{}{} 
\ParallelLText{
\begin{center}
{\large \ggn{36}.}
\end{center}\vspace*{-7pt}

{\greekfont Ἐὰν δύο ῥηταὶ δυνάμει μόνον σύμμετροι συντεθῶσιν, ἡ ὅλη
ἄλογός ἐστιν, καλείσθω δὲ ἐκ δύο ὀνομάτων.}\\

\epsfysize=0.3in
\centerline{\epsffile{Book10/fig036g.eps}}

{\greekfont Συγκείσθωσαν γὰρ δύο ῥηταὶ δυνάμει μόνον σύμμετ-ροι αἱ
ΑΒ, ΒΓ; λέγω, ὅτι ὅλη ἡ ΑΓ ἄλογός ἐστιν.}

{\greekfont Ἐπεὶ γὰρ ἀσύμμετρός ἐστιν ἡ ΑΒ τῇ ΒΓ μήκει; δυνάμει
γὰρ μόνον εἰσὶ σύμμετροι; ὡς δὲ ἡ ΑΒ πρὸς τὴν ΒΓ, οὕτως
τὸ ὑπὸ τῶν ΑΒΓ πρὸς τὸ ἀπὸ τῆς ΒΓ, ἀσύμμετρον ἄρα
ἐστὶ τὸ ὑπὸ τῶν ΑΒ, ΒΓ τῷ ἀπὸ τῆς ΒΓ. ἀλλὰ τῷ μὲν
ὑπὸ τῶν ΑΒ, ΒΓ σύμμετρόν ἐστι τὸ δὶς ὑπὸ τῶν ΑΒ, ΒΓ,
τῷ δὲ ἀπὸ τῆς ΒΓ σύμμετρά ἐστι τὰ ἀπὸ τῶν ΑΒ, ΒΓ;
αἱ γὰρ ΑΒ, ΒΓ ῥηταί εἰσι δυνάμει μόνον σύμμετροι; ἀσύμμετρον
ἄρα ἐστὶ τὸ δὶς ὑπὸ τῶν ΑΒ, ΒΓ τοῖς
ἀπὸ τῶν ΑΒ, ΒΓ. καὶ συνθέντι τὸ δὶς ὑπὸ τῶν ΑΒ, ΒΓ
μετὰ τῶν ἀπὸ τῶν ΑΒ, ΒΓ, τουτέστι τὸ ἀπὸ τῆς ΑΓ,
ἀσύμμετρόν ἐστι τῷ συγκειμένῳ ἐκ τῶν ἀπὸ τῶν ΑΒ, ΒΓ;
ῥητὸν δὲ τὸ συγκείμενον ἐκ τῶν ἀπὸ τῶν ΑΒ, ΒΓ;
ἄλογον ἄρα [ἐστὶ] τὸ ἀπὸ τῆς ΑΓ; ὥστε καὶ ἡ ΑΓ ἄλογός
ἐστιν, καλείσθω δὲ ἐκ δύο ὀνομάτων; ὅπερ ἔδει δεῖχαι.}}

\ParallelRText{
\begin{center}
{\large Proposition 36}
\end{center}

If two rational (straight-lines which are) commensurable
in square only are added together then the whole (straight-line) is irrational---let it
be called a binomial (straight-line).$^\dag$

\epsfysize=0.3in
\centerline{\epsffile{Book10/fig036e.eps}}

For let the two rational (straight-lines), $AB$ and $BC$, (which are)
commensurable in square only, be laid down together. I say that the
whole (straight-line), $AC$, is irrational.
For since $AB$ is incommensurable in length with $BC$---for they
are commensurable in square only---and as $AB$ (is) to $BC$, so
the (rectangle contained) by $ABC$ (is) to the (square) on $BC$, the (rectangle contained)
by $AB$ and $BC$ is thus incommensurable with the
(square) on $BC$ [Prop.~10.11]. But,
twice the (rectangle contained) by $AB$ and $BC$ is commensurable
with the (rectangle contained) by $AB$ and $BC$ [Prop.~10.6]. And (the sum of) the (squares) on
$AB$ and $BC$
 is commensurable with the (square) on $BC$---for the rational (straight-lines) $AB$ and $BC$ are commensurable in square only [Prop.~10.15].
Thus, twice the (rectangle contained) by $AB$ and $BC$ is incommensurable
with (the sum of) the (squares) on $AB$ and $BC$ [Prop.~10.13]. And, via composition, twice
the (rectangle contained) by $AB$ and $BC$, plus (the sum of) the (squares)
on $AB$ and $BC$---that is to say, the (square) on $AC$ [Prop.~2.4]---is incommensurable with the
sum of the (squares) on $AB$ and $BC$ [Prop.~10.16]. And the sum of the (squares) on $AB$ and $BC$ (is) rational. Thus, the (square) on $AC$ [is]
irrational [Def.~10.4]. Hence, $AC$
is also irrational [Def.~10.4]---let it be called
a binomial (straight-line).$^\ddag$  (Which is) the very thing it was required to show.}
\end{Parallel}
{\footnotesize\noindent$^\dag$ Literally, ``from two names".\\[0.5ex]
$^\ddag$ Thus, a binomial straight-line has a length
expressible as  $1 +k^{1/2}$ [or, more generally, $\rho\,(1+k^{1/2})$, where $\rho$ is rational---the same proviso applies to the definitions in the following propositions]. The binomial and the corresponding apotome, whose length is expressible as
$1-k^{1/2}$ (see Prop.~10.73), are
the positive roots of the quartic $x^4-2\,(1+k)\,x^2+ (1-k)^2 = 0$.}

%%%%%%
% Prop 10.37
%%%%%%
\pdfbookmark[1]{Proposition 10.37}{pdf10.37}
\begin{Parallel}{}{} 
\ParallelLText{
\begin{center}
{\large \ggn{37}.}
\end{center}\vspace*{-7pt}

{\greekfont Ἐὰν δύο μέσαι δυνάμει μόνον σύμμετροι συντεθῶσι ῥητὸν περιέθουσαι, ἡ ὅλη ἄλογός ἐστιν, καλείσθω δὲ ἐκ δύο μέσων
πρώτη.}\\

\epsfysize=0.3in
\centerline{\epsffile{Book10/fig036g.eps}}

{\greekfont Συγκείσθωσαν γὰρ δύο μέσαι δυνάμει μόνον σύμμετ-ροι αἱ ΑΒ, ΒΓ
ῥητὸν περιέθουσαι; λέγω, ὅτι ὅλη ἡ ΑΓ ἄλογός ἐστιν.}

{\greekfont Ἐπεὶ γὰρ ἀσύμμετρός ἐστιν ἡ ΑΒ τῇ ΒΓ μήκει, καὶ τὰ ἀπὸ
τῶν ΑΒ, ΒΓ ἄρα ἀσύμμετρά ἐστι τῷ δὶς ὑπὸ τῶν ΑΒ, ΒΓ;
καὶ συνθέντι τὰ ἀπὸ τῶν ΑΒ, ΒΓ μετὰ τοῦ δὶς ὑπὸ τῶν ΑΒ, ΒΓ,
ὅπερ ἐστὶ τὸ ἀπὸ τῆς ΑΓ, ἀσύμμετρόν ἐστι τῷ ὑπὸ τῶν
ΑΒ, ΒΓ.
ῥητὸν δὲ τὸ ὑπὸ τῶν ΑΒ, ΒΓ; ὑπόκεινται γὰρ αἱ ΑΒ, ΒΓ
ῥητὸν περιέθουσαι; ἄλογον ἄρα τὸ ἀπὸ τῆς ΑΓ;
ἄλογος ἄρα ἡ ΑΓ, καλείσθω δὲ ἐκ δύο μέσων πρώτη; ὅπερ
ἔδει δεῖχαι.}}

\ParallelRText{
\begin{center}
{\large Proposition 37}
\end{center}

If two medial (straight-lines), commensurable in square only, which contain a rational (area), are added together then
the whole (straight-line) is irrational---let it be called a first bimedial (straight-line).$^\dag$

\epsfysize=0.3in
\centerline{\epsffile{Book10/fig036e.eps}}

For let the  two medial (straight-lines), $AB$ and $BC$, 
commensurable in square only, (and) containing a rational (area), be laid down together. I say that the
whole (straight-line), $AC$, is irrational.

For since $AB$ is incommensurable in length with $BC$, (the sum of)
the (squares) on $AB$ and $BC$ is thus also incommensurable with twice
the (rectangle contained) by $AB$ and $BC$ [see previous proposition].
And, via composition, (the sum of) the (squares)  on $AB$ and $BC$,
plus twice the (rectangle contained) by $AB$ and $BC$---that is, 
the (square) on $AC$ [Prop.~2.4]---is incommensurable with the (rectangle contained) by $AB$ and $BC$
[Prop.~10.16]. And the (rectangle contained) by
$AB$ and $BC$ (is) rational---for $AB$ and $BC$ were assumed to
enclose a rational (area). Thus, the (square) on $AC$ (is) irrational. Thus, $AC$ (is) irrational [Def.~10.4]---let it be called
a first bimedial (straight-line).$^\ddag$
 (Which is) the very thing it was required to show.}
\end{Parallel}
{\footnotesize\noindent$^\dag$ Literally, ``first from two medials''.\\[0.5ex]
$^\ddag$ Thus, a first bimedial straight-line
has a length expressible as $k^{1/4}+k^{3/4}$. The first bimedial and the corresponding first apotome of a medial, whose length is expressible as
$k^{1/4}-k^{3/4}$ (see Prop.~10.74), are
the positive roots of the quartic $x^4-2\,\sqrt{k}\,(1+k)\,x^2+ k\,(1-k)^2 = 0$.}

%%%%%%
% Prop 10.38
%%%%%%
\pdfbookmark[1]{Proposition 10.38}{pdf10.38}
\begin{Parallel}{}{} 
\ParallelLText{
\begin{center}
{\large \ggn{38}.}
\end{center}\vspace*{-7pt}

{\greekfont Ἐὰν δύο μέσαι δυνάμει μόνον σύμμετροι συντεθῶσι μέσον περιέθουσαι, ἡ ὅλη ἄλογός ἐστιν, καλείσθω δὲ ἐκ δύο μέσων
δυετέρα.}\\

\epsfysize=1.7in
\centerline{\epsffile{Book10/fig038g.eps}}

{\greekfont Συγκείσθωσαν γὰρ δύο μέσαι δυνάμει μόνον σύμμετ-ροι αἱ ΑΒ, ΒΓ
μέσον περιέθουσαι; λέγω, ὅτι ἄλογός ἐστιν ἡ ΑΓ.}

{\greekfont Ἐκκείσθω γὰρ ῥητὴ ἡ ΔΕ, καὶ τῷ ἀπὸ τῆς ΑΓ ἴσον
παρὰ τὴν ΔΕ παραβεβλήσθω τὸ ΔΖ πλάτος ποιοῦν τὴν ΔΗ. καὶ
ἐπεὶ τὸ ἀπὸ τῆς ΑΓ ἴσον ἐστὶ τοῖς τε ἀπὸ τῶν ΑΒ, ΒΓ
καὶ τῷ δὶς ὑπὸ τῶν ΑΒ, ΒΓ, παραβεβλήσθω δὴ τοῖς ἀπὸ τῶν
ΑΒ, ΒΓ παρὰ τὴν ΔΕ ἴσον τὸ ΕΘ; λοιπὸν ἄρα τὸ ΘΖ ἴσον
ἐστὶ τῷ δὶς ὑπὸ τῶν ΑΒ, ΒΓ. καὶ ἐπεὶ μέση ἐστὶν ἑκατέρα
τῶν ΑΒ, ΒΓ, μέσα ἄρα ἐστὶ καὶ τὰ ἀπὸ τῶν ΑΒ, ΒΓ. μέσον
δὲ ὑπόκειται καὶ τὸ δὶς ὑπὸ τῶν ΑΒ, ΒΓ. καί ἐστι τοῖς
μὲν ἀπὸ τῶν ΑΒ, ΒΓ ἴσον τὸ ΕΘ, τῷ δὲ δὶς ὑπὸ τῶν
ΑΒ, ΒΓ ἴσον τὸ ΖΘ; μέσον ἄρα ἑκάτερον τῶν ΕΘ, ΘΖ. καὶ παρὰ
ῥητὴν τὴν ΔΕ παράκειται; ῥητὴ ἄρα ἐστὶν ἑκατέρα τῶν
ΔΘ, ΘΗ καὶ ἀσύμμετρος τῇ ΔΕ μήκει. ἐπεὶ οὖν ἀσύμμετρός
ἐστιν ἡ ΑΒ τῇ ΒΓ μήκει, καί ἐστιν ὡς ἡ ΑΒ πρὸς τὴν ΒΓ, οὕτως
τὸ ἀπὸ τῆς ΑΒ πρὸς τὸ ὑπὸ τῶν ΑΒ, ΒΓ, ἀσύμμετρον ἄρα ἐστὶ
τὸ ἀπὸ τῆς ΑΒ τῷ ὑπὸ τῶν ΑΒ, ΒΓ. ἀλλὰ τῷ μὲν ἀπὸ τῆς
ΑΒ σύμμετρόν ἐστι τὸ συγκείμενον ἐκ τῶν ἀπὸ τῶν ΑΒ, ΒΓ
τετραγώνων, τῷ δὲ ὑπὸ τῶν ΑΒ, ΒΓ σύμμετρόν ἐστι τὸ δὶς
ὑπὸ τῶν ΑΒ, ΒΓ. ἀσύμμετρον ἄρα ἐστὶ τὸ συγκείμενον
ἐκ τῶν ἀπὸ τῶν ΑΒ, ΒΓ τῷ δὶς ὑπὸ τῶν ΑΒ, ΒΓ. ἀλλὰ
τοῖς μὲν ἀπὸ τῶν ΑΒ, ΒΓ ἴσον ἐστὶ τὸ ΕΘ, τῷ δὲ δὶς ὑπὸ
τῶν ΑΒ, ΒΓ ἴσον
ἐστὶ τὸ ΘΖ. ἀσύμμετρον ἄρα ἐστὶ τὸ ΕΘ τῷ ΘΖ; ὥστε καὶ ἡ ΔΘ
τῇ ΘΗ ἐστιν ἀσύμμετρος μήκει. αἱ ΔΘ, ΘΗ ἄρα ῥηταί εἰσι
δυνάμει μόνον σύμμετροι. ὥστε ἡ ΔΗ ἄλογός ἐστιν. ῥητὴ δὲ ἡ ΔΕ; τὸ δὲ ὑπὸ ἀλόγου καὶ ῥητῆς περιεχόμενον ὀρθογώνιον
ἄλογόν ἐστιν; ἄλογον ἄρα ἐστὶ τὸ ΔΖ θωρίον, καὶ ἡ δυναμένη
[α ὐτὸ] ἄλογός ἐστιν. δύναται δὲ τὸ ΔΖ ἡ ΑΓ; ἄλογος ἄρα
ἐστὶν ἡ ΑΓ, καλείσθω δὲ ἐκ δύο μέσων δευτέρα. ὅπερ
ἔδει δεῖχαι.}}

\ParallelRText{
\begin{center}
{\large Proposition 38}
\end{center}
If two medial (straight-lines), commensurable in
square only, which contain a medial (area), are added together then the
whole (straight-line) is irrational---let it be called a second bimedial
(straight-line).

\epsfysize=1.7in
\centerline{\epsffile{Book10/fig038e.eps}}

For let the  two medial (straight-lines), $AB$ and $BC$, 
commensurable in square only, (and) containing a medial (area), be laid down together [Prop.~10.28]. I say that $AC$ is irrational.

For let the rational (straight-line) $DE$ be laid down, and let (the rectangle) $DF$, equal
to the (square) on $AC$, be applied to $DE$, making $DG$ as breadth
[Prop.~1.44]. And since the
(square) on $AC$ is equal to (the sum of) the (squares) on $AB$ and $BC$,
plus twice the (rectangle contained) by $AB$ and $BC$ [Prop.~2.4], so let (the rectangle) $EH$, equal to
(the sum of) the squares on $AB$ and $BC$, be applied
to $DE$. The remainder $HF$ is thus equal to twice the (rectangle contained)
by $AB$ and $BC$. And since $AB$ and $BC$ are each medial, (the
sum of) the squares on $AB$ and $BC$ is thus also medial.$^\ddag$ And twice the (rectangle contained) by $AB$ and $BC$ was also assumed (to be) medial.
And $EH$ is equal to (the sum of) the squares on $AB$ and $BC$, and
$FH$ (is) equal to twice the (rectangle contained) by $AB$ and $BC$. 
Thus, $EH$ and $HF$ (are) each medial. And they were applied to the
rational (straight-line) $DE$. Thus, $DH$ and $HG$ are
each rational, and incommensurable in length with $DE$ [Prop.~10.22]. Therefore, since $AB$ is incommensurable in length with $BC$, and as $AB$ is to $BC$, so the
(square) on $AB$ (is) to the (rectangle contained) by $AB$ and $BC$ [Prop.~10.21~lem.], the (square) on $AB$
is thus incommensurable with the (rectangle contained) by $AB$ and $BC$
[Prop.~10.11]. But, the sum of the squares on $AB$ and $BC$ is commensurable with the (square) on $AB$ [Prop.~10.15], and twice the (rectangle contained)
by $AB$ and $BC$ is commensurable with the (rectangle contained) by $AB$ and $BC$ [Prop.~10.6]. Thus, the sum of
the (squares) on $AB$ and $BC$ is incommensurable with twice the (rectangle
contained) by $AB$ and $BC$ [Prop.~10.13].
But, $EH$ is equal to (the sum of) the squares on $AB$ and $BC$,
and $HF$ is equal to twice the (rectangle) contained by $AB$ and $BC$. Thus, $EH$ is incommensurable with $HF$. Hence, $DH$ is also
incommensurable in length with $HG$ [Props.~6.1, 10.11]. Thus, $DH$ and 
$HG$ are rational (straight-lines which are) commensurable in square only.
Hence, $DG$ is irrational [Prop.~10.36].
And  $DE$ (is) rational. And the rectangle contained by irrational and
rational (straight-lines) is irrational [Prop.~10.20].
The area $DF$ is thus irrational, and (so) the square-root [of it] is irrational
[Def.~10.4]. And  $AC$ is the square-root of $DF$. $AC$ is thus irrational---let it be called a second bimedial (straight-line).$^\S$ (Which is) the very thing it was required to show.}
\end{Parallel}
{\footnotesize\noindent$^\dag$ Literally, 
``second from two medials''.\\[0.5ex]
$^\ddag$ Since, by hypothesis, the squares on $AB$ and $BC$ are commensurable---see Props.~10.15, 10.23.\\[0.5ex]
$^\S$ Thus, a second bimedial straight-line
has a length expressible as $k^{1/4}+k'^{1/2}/k^{1/4}$. The second bimedial and the corresponding second apotome of a medial, whose length is expressible as
$k^{1/4}-k'^{1/2}/k^{1/4}$ (see Prop.~10.75), are
the positive roots of the quartic $x^4-2\,[(k+k')/\sqrt{k}\,]\,x^2+ [(k-k')^2/k] = 0$.}

%%%%%%
% Prop 10.39
%%%%%%
\pdfbookmark[1]{Proposition 10.39}{pdf10.39}
\begin{Parallel}{}{} 
\ParallelLText{
\begin{center}
{\large \ggn{39}.}
\end{center}\vspace*{-7pt}

{\greekfont Ἐὰν δύο εὐθεῖαι δυνάμει ἀσύμμετροι συντεθῶσι ποιοῦς-αι τὸ
μὲν συγκείμενον ἐκ τῶν ἀπ α ὐτῶν τετραγώνων ῥητόν, τὸ
δ ὑπ α ὐτῶν μέσον, ἡ ὅλη εὐθεῖα ἄλογός ἐστιν, καλείσθω
δὲ μείζων.}\\

\epsfysize=0.3in
\centerline{\epsffile{Book10/fig036g.eps}}

{\greekfont Συγκείσθωσαν γὰρ δύο εὐθεῖαι δυνάμει ἀσύμμετροι αἱ ΑΒ, ΒΓ
ποιοῦσαι τὰ προκείμενα; λέγω, ὅτι ἄλογός ἐστιν ἡ ΑΓ.}

{\greekfont Ἐπεὶ γὰρ τὸ ὑπὸ τῶν ΑΒ, ΒΓ μέσον ἐστίν, καὶ τὸ δὶς [ἄρα]
ὑπὸ τῶν ΑΒ, ΒΓ μέσον ἐστίν. τὸ δὲ συγκείμενον ἐκ τῶν
ἀπὸ τῶν ΑΒ, ΒΓ ῥητόν; ἀσύμμετρον ἄρα ἐστὶ τὸ δὶς ὑπὸ
τῶν ΑΒ, ΒΓ τῷ συγκειμένῳ ἐκ τῶν ἀπὸ τῶν ΑΒ, ΒΓ;
 ὥστε καὶ τὰ ἀπὸ τῶν ΑΒ, ΒΓ μετὰ τοῦ δὶς ὑπὸ
τῶν ΑΒ, ΒΓ, ὅπερ ἐστὶ τὸ ἀπὸ τῆς ΑΓ, ἀσύμμετρόν ἐστι
τῷ συγκειμένῳ ἐκ τῶν ἀπὸ τῶν ΑΒ, ΒΓ [ῥητὸν δὲ τὸ συγμείμενον ἐκ τῶν ἀπὸ τῶν ΑΒ, ΒΓ]; ἄλογον ἄρα ἐστὶ
τὸ ἀπὸ τῆς ΑΓ. ὥστε καὶ ἡ ΑΓ ἄλογός ἐστιν, καλείσθω
δὲ μείζων. ὅπερ ἔδει δεῖχαι.}}

\ParallelRText{
\begin{center}
{\large Proposition 39}
\end{center}

If two straight-lines (which are) incommensurable
in square, making the sum of the squares on them
rational, and the (rectangle contained) by them medial,  are added together then the whole
straight-line is irrational---let it be called a major (straight-line).

\epsfysize=0.3in
\centerline{\epsffile{Book10/fig036e.eps}}

For let the two straight-lines, $AB$ and $BC$,  incommensurable
in square, and fulfilling  the  prescribed (conditions), be laid down together [Prop.~10.33]. I say that $AC$ is irrational.

For since the (rectangle contained) by $AB$ and $BC$ is medial, twice
the (rectangle contained) by $AB$ and $BC$ is [thus] also medial [Props.~10.6, 10.23~corr.].
And the sum of the (squares) on $AB$ and $BC$ (is) rational. Thus, twice
the (rectangle contained) by $AB$ and $BC$ is incommensurable with the
sum of the (squares) on $AB$ and $BC$ [Def.~10.4].
Hence, (the sum of) the squares on $AB$ and $BC$, plus twice the
(rectangle contained) by $AB$ and $BC$---that is, the (square) on $AC$ [Prop.~2.4]---is also
incommensurable with the sum of the (squares) on $AB$ and $BC$
[Prop.~10.16] [and the
sum of the (squares) on $AB$ and $BC$ (is) rational]. Thus, the (square) on $AC$ is irrational. Hence, $AC$ is also irrational [Def.~10.4]---let it be called a major (straight-line).$^\dag$ (Which is) the
very thing it was required to show.}
\end{Parallel}
{\footnotesize\noindent$^\dag$ Thus, a major straight-line
has a length expressible as $\sqrt{[1+k/(1+k^2)^{1/2}]/2} + \sqrt{[1-k/(1+k^2)^{1/2}]/2}$. The major and the corresponding minor, whose length is expressible as $\sqrt{[1+k/(1+k^2)^{1/2}]/2} - \sqrt{[1-k/(1+k^2)^{1/2}]/2}$ (see Prop.~10.76), are
the positive roots of the quartic $x^4-2\,x^2+ k^2/(1+k^2) = 0$.}

%%%%%%
% Prop 10.40
%%%%%%
\pdfbookmark[1]{Proposition 10.40}{pdf10.40}
\begin{Parallel}{}{} 
\ParallelLText{
\begin{center}
{\large\ggn{40}.}
\end{center}\vspace*{-7pt}

{\greekfont Ἐὰν δύο εὐθεῖαι δυνάμει ἀσύμμετροι συντεθῶσι ποιοῦς-αι τὸ μὲν
συγκείμενον ἐκ τῶν ἀπ α ὐτῶν τετραγώνων μέσον, τὸ δ ὑπ
α ὐτῶν ῥητόν, ἡ ὅλη εὐθεῖα ἄλογός ἐστιν, καλείσθω δὲ
ῥητὸν καὶ μέσον δυναμένη.}\\~\\

\epsfysize=0.3in
\centerline{\epsffile{Book10/fig036g.eps}}

{\greekfont Συγκείσθωσαν γὰρ δύο εὐθεῖαι δυνάμει ἀσύμμετροι αἱ ΑΒ, ΒΓ
ποιοῦσαι τὰ προκείμενα; λέγω, ὅτι ἄλογός ἐστιν ἡ ΑΓ.}

{\greekfont Ἐπεὶ γὰρ τὸ συγκείμενον ἐκ τῶν ἀπὸ τῶν ΑΒ, ΒΓ μέσον
ἐστίν, τὸ δὲ δὶς ὑπὸ τῶν ΑΒ, ΒΓ ῥητόν, ἀσύμμετρον
ἄρα ἐστὶ τὸ συγκείμενον ἐκ τῶν ἀπὸ τῶν ΑΒ, ΒΓ τῷ δὶς
ὑπὸ τῶν ΑΒ, ΒΓ; ὥστε καὶ τὸ ἁπὸ τῆς ΑΓ ἀσύμμετρόν ἐστι
τῷ δὶς ὑπὸ τῶν ΑΒ, ΒΓ. ῥητὸν δὲ τὸ δὶς ὑπὸ τῶν ΑΒ, ΒΓ;
ἄλογον ἄρα τὸ ἀπὸ τῆς ΑΓ. ἄλογος ἄρα ἡ ΑΓ, καλείσθω δὲ ῥητὸν
καὶ μέσον δυναμένη. ὅπερ ἔδει δεῖχαι.}}

\ParallelRText{
\begin{center}
{\large Proposition 40}
\end{center}

If two straight-lines (which are) incommensurable
in square, making the sum of the squares on them medial, 
and the (rectangle contained) by them rational, are added together then the whole straight-line
is irrational---let it be called the square-root of a rational plus a medial (area).

\epsfysize=0.3in
\centerline{\epsffile{Book10/fig036e.eps}}

For let the two straight-lines, $AB$ and $BC$, incommensurable
in square, (and) fulfilling the prescribed (conditions), be laid down together [Prop.~10.34]. I say that $AC$ is irrational.

For since the sum of the (squares) on $AB$ and $BC$ is medial,
and twice the (rectangle contained) by $AB$ and $BC$ (is) rational, 
the sum of the (squares) on $AB$ and $BC$ is thus incommensurable
with twice the (rectangle contained) by $AB$ and $BC$. Hence, the (square)
on $AC$ is also incommensurable with twice the (rectangle contained)
by $AB$ and $BC$ [Prop.~10.16]. And twice the
(rectangle contained) by $AB$ and $BC$ (is) rational. The (square) on
$AC$ (is) thus irrational. Thus, $AC$ (is) irrational [Def.~10.4]---let it be called
the square-root of a rational plus a medial (area).$^\dag$
(Which is) the very thing it
was required to show.}
\end{Parallel}
{\footnotesize\noindent$^\dag$ Thus, the square-root
of a rational plus a medial (area)
has a length expressible as $\sqrt{[(1+k^2)^{1/2}+k]/[2\,(1+k^2)]} +\sqrt{[(1+k^2)^{1/2}-k]/[2\,(1+k^2)]}$. This  and the corresponding irrational with a minus sign, whose length is expressible as $\sqrt{[(1+k^2)^{1/2}+k]/[2\,(1+k^2)]} -\sqrt{[(1+k^2)^{1/2}-k]/[2\,(1+k^2)]}$
 (see Prop.~10.77), are the positive roots of the quartic $x^4-(2/\sqrt{1+k^2})\,x^2+ k^2/(1+k^2)^2 = 0$.} 
 
%%%%%%
% Prop 10.41
%%%%%%
\pdfbookmark[1]{Proposition 10.41}{pdf10.41}
\begin{Parallel}{}{} 
\ParallelLText{
\begin{center}
{\large\ggn{41}.}
\end{center}\vspace*{-7pt}

{\greekfont Ἐὰν δύο εὐθεῖαι δυνάμει ἀσύμμετροι συντεθῶσι ποιοῦς-αι τό
τε συγκείμενον ἐκ τῶν ἀπ α ὐτῶν τετραγώνων μέσον καὶ τὸ
ὑπ α ὐτῶν μέσον καὶ ἔτι ἀσύμμετρον τῷ συγκειμένῳ
ἐκ τῶν ἀπ α ὐτῶν τετραγώνων, ἡ ὅλη εὐθεῖα ἄλογός
ἐστιν, καλείσθω δὲ δύο μέσα δυναμένη.}\\

\epsfysize=2.2in
\centerline{\epsffile{Book10/fig041g.eps}}

{\greekfont Συγκείσθωσαν γὰρ δύο εὐθεῖαι δυνάμει ἀσύμμετροι
αἱ ΑΒ, ΒΓ ποιοῦσαι τὰ προκείμενα; λέγω, ὅτι ἡ ΑΓ ἄλογός
ἐστιν.}

{\greekfont Ἐκκείσθω ῥητὴ ἡ ΔΕ, καὶ παραβεβλήσθω παρὰ τὴν ΔΕ τοῖς μὲν ἀπὸ
τῶν ΑΒ, ΒΓ ἴσον τὸ ΔΖ, τῷ δὲ δὶς ὑπὸ τῶν ΑΒ, ΒΓ ἴσον
τὸ ΗΘ; ὅλον ἄρα τὸ ΔΘ ἴσον ἐστὶ τῷ ἀπὸ τῆς ΑΓ τετραγώνῳ.
καὶ ἐπεὶ μέσον ἐστὶ τὸ συγκείμενον ἐκ τῶν ἀπὸ τῶν
ΑΒ, ΒΓ, καί ἐστιν ἴσον τῷ ΔΖ, μέσον ἄρα ἐστὶ καὶ τὸ ΔΖ.
καὶ παρὰ ῥητὴν τὴν ΔΕ παράκειται; ῥητὴ ἄρα ἐστὶν ἡ
ΔΗ καὶ ἀσύμμετρος τῇ ΔΕ μήκει. διὰ τὰ α ὐτὰ δὴ
καὶ ἡ ΗΚ ῥητή ἐστι καὶ ἀσύμμετρος τῇ ΗΖ, τουτέστι τῇ ΔΕ, μήκει.
καὶ ἐπεὶ ἀσύμμετρά ἐστι τὰ ἀπὸ τῶν ΑΒ, ΒΓ τῷ δὶς
ὑπὸ τῶν ΑΒ, ΒΓ, ἀσύμμετρόν ἐστι τὸ ΔΖ τῷ ΗΘ; ὥστε καὶ ἡ
ΔΗ τῇ ΗΚ ἀσύμμετρός ἐστιν. καὶ εἰσι ῥηταί; αἱ ΔΗ, ΗΚ
ἄρα ῥηταί εἰσι δυνάμει μόνον σύμμετροι;
ἄλογος ἄρα ἐστὶν ἡ ΔΚ ἡ καλουμένη ἐκ δύο ὀνομάτων.
ῥητὴ δὲ ἡ ΔΕ; ἄλογον ἄρα ἐστὶ τὸ ΔΘ καὶ ἡ δυναμένη
α ὐτὸ ἄλογός ἐστιν. δύναται δὲ τὸ ΘΔ ἡ ΑΓ; ἄλογος ἄρα
ἐστὶν ἡ ΑΓ, καλείσθω δὲ δύο μέσα δυναμένη. ὅπερ ἔδει δεῖχαι.}}

\ParallelRText{
\begin{center}
{\large Proposition 41}
\end{center}

If two straight-lines (which are) incommensurable
in square, making the sum of the squares
on them medial, and the (rectangle contained) by them medial, and,
moreover, incommensurable with the sum of the squares on them, are added together then
the whole straight-line is irrational---let it be called the square-root
of (the sum of) two medial (areas).

\epsfysize=2.2in
\centerline{\epsffile{Book10/fig041e.eps}}

For let the two straight-lines, $AB$ and $BC$,  incommensurable
in square,  (and) fulfilling the prescribed (conditions), be laid down together
[Prop.~10.35]. I say that $AC$ is irrational.

Let the rational (straight-line) $DE$ be laid out, and let (the rectangle) $DF$,
equal to (the sum of) the (squares) on $AB$ and $BC$, and (the rectangle) $GH$, equal to twice the (rectangle contained) by
$AB$ and $BC$, be applied
to $DE$. Thus, the whole of $DH$ is equal to the square on $AC$
[Prop.~2.4]. And since the sum of the (squares)
on $AB$ and $BC$ is medial, and is equal to $DF$, $DF$ is thus also
medial. And it is applied to the rational (straight-line) $DE$. 
Thus, $DG$ is rational, and incommensurable in length with $DE$ [Prop.~10.22]. So, for the same (reasons), $GK$
is also rational, and incommensurable in length with $GF$---that is
to say, $DE$. And since (the sum of) the (squares) on $AB$ and $BC$
is incommensurable with twice the (rectangle contained) by $AB$ and
$BC$, $DF$ is incommensurable with $GH$.
Hence, $DG$ is also incommensurable (in length)
with $GK$ [Props.~6.1, 10.11]. And they are rational. Thus, $DG$ and
$GK$ are rational (straight-lines which are) commensurable in square only. Thus, $DK$ is irrational, and that (straight-line which is) called binomial 
[Prop.~10.36]. And $DE$ (is) rational.
Thus, $DH$ is irrational, and its square-root is irrational [Def.~10.4]. And $AC$ (is) the square-root of $HD$. Thus,
$AC$ is irrational---let it be called the square-root of (the sum of) two medial
(areas).$^\dag$ (Which is) the very thing it was required to show.}
\end{Parallel}
{\footnotesize\noindent$^\dag$ Thus, the square-root
of (the sum of) two medial (areas)
has a length expressible as $k'^{1/4}\left(\sqrt{[1+k/(1+k^2)^{1/2}]/2}
+\sqrt{[1-k/(1+k^2)^{1/2}]/2}\right)$. This  and the corresponding irrational with a minus sign, whose length is expressible as $k'^{1/4}\left(\sqrt{[1+k/(1+k^2)^{1/2}]/2}
-\sqrt{[1-k/(1+k^2)^{1/2}]/2}\right)$
 (see Prop.~10.78), are
the positive roots of the quartic $x^4-2\,k'^{1/2}\,x^2+ k'\, k^2/(1+k^2)= 0$.} 

%%%%%%
% Prop 10.41a
%%%%%%
\begin{Parallel}{}{} 
\ParallelLText{
\begin{center}
{\large {\greekfont Λῆμμα}.}
\end{center}\vspace*{-7pt}

{\greekfont Ὅτι δὲ αἱ εἰρημέναι ἄλογοι μοναθῶς διαιροῦνται εἰς
τὰς εὐθείας, ἐχ ὧν σύγκεινται ποιουσῶν τὰ προκείμενα
εἴδη, δείχομεν ἤδη προεκθέμενοι λημμάτιον τοιοῦτον;}\\

\epsfysize=0.3in
\centerline{\epsffile{Book10/fig041ag.eps}}

{\greekfont Ἐκκείσθω εὐθεῖα ἡ ΑΒ καὶ τετμήσθω ἡ ὅλη εἰς ἄνισα
καθ ἑκάτερον τῶν Γ, Δ, ὑποκείσθω δὲ μείζων ἡ ΑΓ τῆς ΔΒ;
λέγω, ὅτι τὰ ἀπὸ τῶν ΑΓ, ΓΒ μείζονά ἐστι τῶν ἀπὸ
τῶν ΑΔ, ΔΒ.}

{\greekfont Τετμήσθω γὰρ ἡ ΑΒ δίθα κατὰ τὸ Ε. καὶ ἐπεὶ μείζων ἐστὶν
ἡ ΑΓ τῆς ΔΒ, κοινὴ ἀφῃρήσθω ἡ ΔΓ; λοιπὴ ἄρα ἡ ΑΔ λοιπῆς
τῆς ΓΒ μείζων ἐστίν. ἴση δὲ ἡ ΑΕ τῇ ΕΒ; ἐλάττων
ἄρα ἡ ΔΕ τῆς ΕΓ; τὰ Γ, Δ ἄρα σημεῖα οὐκ ἴσον
ἀπέθουσι τῆς διθοτομίας. καὶ ἐπεὶ τὸ ὑπὸ τῶν ΑΓ, ΓΒ
μετὰ τοῦ ἀπὸ τῆς ΕΓ ἴσον ἐστὶ τῷ ἀπὸ τῆς ΕΒ, ἀλλὰ
μὴν καὶ τὸ ὑπὸ τῶν ΑΔ, ΔΒ μετὰ τοῦ ἀπὸ ΔΕ ἴσον
ἐστὶ τῷ ἀπὸ τῆς ΕΒ, τὸ ἄρα ὑπὸ τῶν
ΑΓ, ΓΒ μετὰ τοῦ ἀπὸ τῆς ΕΓ ἴσον ἐστὶ τῷ ὑπὸ τῶν ΑΔ, ΔΒ μετὰ τοῦ ἀπὸ τῆς ΔΕ; ὧν τὸ ἀπὸ τῆς ΔΕ
ἔλασσόν ἐστι τοῦ ἀπὸ τῆς ΕΓ; καὶ λοιπὸν ἄρα τὸ ὑπὸ
τῶν ΑΓ, ΓΒ ἔλασσόν ἐστι τοῦ ὑπὸ τῶν ΑΔ, ΔΒ. ὥστε
καὶ τὸ δὶς ὑπὸ τῶν ΑΓ, ΓΒ ἔλασσόν ἐστι τοῦ δὶς ὑπὸ τῶν
ΑΔ, ΔΒ. καὶ λοιπὸν ἄρα τὸ συγκείμενον ἐκ τῶν ἀπὸ τῶν
ΑΓ, ΓΒ μεῖζόν ἐστι τοῦ συγκειμένου ἐκ τῶν ἀπὸ τῶν
ΑΔ, ΔΒ. ὅπερ ἔδει δεῖχαι.}}

\ParallelRText{
\begin{center}
{\large Lemma}
\end{center}

We will now demonstrate that the aforementioned irrational (straight-lines)
 are uniquely divided into the straight-lines  of which they are the sum,
 and which produce the prescribed types, 
 (after) setting forth the following
 lemma.
 
 \epsfysize=0.3in
\centerline{\epsffile{Book10/fig041ae.eps}}
 
 Let the straight-line $AB$ be laid out, and let the whole (straight-line) be cut
 into unequal parts at each of the (points) $C$ and $D$. And let $AC$ be assumed (to be) greater than $DB$. I say that (the sum of) the (squares) on
 $AC$ and $CB$ is greater than (the sum of) the (squares) on $AD$ and $DB$.
 
 For let $AB$ be cut in half at $E$. And since $AC$ is greater than
 $DB$, let $DC$ be subtracted from both. Thus, the remainder
 $AD$ is greater than the remainder $CB$. And $AE$ (is) equal to $EB$.
 Thus, $DE$ (is) less than $EC$. Thus, points $C$ and $D$ are not
 equally far from the point of bisection. And since the (rectangle
 contained) by $AC$ and $CB$, plus the (square) on $EC$,
 is equal to the (square) on $EB$ [Prop.~2.5], 
 but, moreover, the (rectangle contained) by $AD$ and $DB$, plus the
 (square) on $DE$, is also equal to the (square) on $EB$ [Prop.~2.5], the (rectangle contained) by $AC$ and
 $CB$, plus the (square) on $EC$, is thus equal to the (rectangle contained)
 by $AD$ and $DB$, plus the (square) on $DE$. And, of these, the
 (square) on $DE$ is less than the (square) on $EC$. And, thus, the remaining
 (rectangle contained) by $AC$ and $CB$ is  less than
 the (rectangle contained) by $AD$ and $DB$. And, hence, twice
 the (rectangle contained) by $AC$ and $CB$ is less than twice the
 (rectangle contained) by $AD$ and $DB$. And thus the remaining
 sum of the (squares) on $AC$ and $CB$ is greater than the sum of the
 (squares) on $AD$ and $DB$.$^\dag$  (Which is) the very thing it was required to show.}
\end{Parallel}
{\footnotesize\noindent$^\dag$ Since, $AC^{\,2}+CB^{\,2}+2\,AC\,CB=
 AD^{\,2}+DB^{\,2}
+2\,AD\,DB = AB^{\,2}$.} 

%%%%%%
% Prop 10.42
%%%%%%
\pdfbookmark[1]{Proposition 10.42}{pdf10.42}
\begin{Parallel}{}{} 
\ParallelLText{
\begin{center}
{\large \ggn{42}.}
\end{center}\vspace*{-7pt}

{\greekfont Ἡ ἐκ δύο ὀνομάτων κατὰ ἓν μόνον σημεῖον διαιρεῖται εἰς
τὰ ὀνόματα.}

\epsfysize=0.3in
\centerline{\epsffile{Book10/fig042g.eps}}

{\greekfont Ἔστω ἐκ δύο ὀνομάτων ἡ ΑΒ διῃρημένη εἰς τὰ ὀνόματα
κατὰ τὸ Γ; αἱ ΑΓ, ΓΒ ἄρα ῥηταί εἰσι δυνάμει μόνον σύμμετροι.
λέγω, ὅτι ἡ ΑΒ κατ ἄλλο σημεῖον οὐ διαιρεῖται εἰς δύο
ῥητὰς δυνάμει μόνον συμμέτρους.}

{\greekfont Εἰ γὰρ δυνατόν, διῃρήσθω καὶ κατὰ τὸ Δ, ὥστε καὶ τὰς ΑΔ,
ΔΒ ῥητὰς εἶναι δυνάμει μόνον συμμέτρους. φανερὸν δή,
ὅτι ἡ ΑΓ τῇ ΔΒ οὐκ ἔστιν ἡ α ὐτή. εἰ γὰρ δυνατόν,
ἔστω. ἔσται δὴ καὶ ἡ ΑΔ τῇ ΓΒ ἡ α ὐτή; καὶ ἔσται
ὡς ἡ ΑΓ πρὸς τὴν ΓΒ, οὕτως ἡ ΒΔ πρὸς τὴν ΔΑ,
καὶ ἔσται ἡ ΑΒ κατὰ τὸ α ὐτὸ τῇ κατὰ τὸ Γ διαιρέσει διαιρεθεῖσα
καὶ κατὰ τὸ Δ; ὅπερ οὐθ ὑπόκειται. οὐκ ἄρα ἡ ΑΓ τῇ
ΔΒ ἐστιν ἡ α ὐτή. διὰ δὴ τοῦτο καὶ τὰ Γ, Δ σημεῖα
οὐκ ἴσον ἀπέθουσι τῆς διθοτομίας. ᾧ ἄρα διαφέρει τὰ
ἀπὸ τῶν ΑΓ, ΓΒ τῶν ἀπὸ τῶν ΑΔ, ΔΒ, τούτῳ διαφέρει καὶ
τὸ δὶς ὑπὸ τῶν ΑΔ, ΔΒ τοῦ δὶς ὑπὸ τῶν ΑΓ, ΓΒ
διὰ τὸ καὶ τὰ ἀπὸ τῶν ΑΓ, ΓΒ μετὰ τοῦ δὶς ὑπὸ τῶν
ΑΓ, ΓΒ καὶ τὰ ἀπὸ τῶν ΑΔ, ΔΒ μετὰ τοῦ δὶς ὑπὸ τῶν
ΑΔ, ΔΒ ἴσα εἶναι τῷ ἀπὸ τῆς ΑΒ. ἀλλὰ τὰ ἀπὸ
τῶν ΑΓ, ΓΒ τῶν ἀπὸ τῶν ΑΔ, ΔΒ διαφέρει ῥητῷ;
ῥητὰ γὰρ ἀμφότερα; καὶ τὸ δὶς ἄρα ὑπὸ τῶν ΑΔ, ΔΒ
τοῦ δὶς ὑπὸ τῶν ΑΓ, ΓΒ διαφέρει ῥητῷ μέσα
ὄντα; ὅπερ ἄτοπον; μέσον γὰρ μέσου οὐθ ὑπερέθει
ῥητῷ.}

{\greekfont Οὐθ ἄρα ἡ ἐκ δύο ὀνομάτων κατ ἄλλο καὶ ἄλλο
σημεῖον διαιρεῖται; καθ ἓν ἄρα μόνον; ὅπερ
ἔδει δεῖχαι.}}

\ParallelRText{
\begin{center}
{\large Proposition 42}
\end{center}
A binomial (straight-line) can be divided into
its (component) terms at one
point only.$^\dag$

\epsfysize=0.3in
\centerline{\epsffile{Book10/fig042e.eps}}

Let $AB$ be a binomial (straight-line) which has been divided into its
(component) terms at $C$. $AC$ and $CB$ are thus
rational (straight-lines which are) commensurable in square only [Prop.~10.36]. I say that $AB$ cannot be divided
at another point
into two rational (straight-lines which are) commensurable in square only.

For, if possible, let it also be divided at $D$, such that $AD$ and
$DB$ are also rational (straight-lines which are) commensurable in
square only. So, (it is) clear that $AC$ is not the same as $DB$.
For, if possible, let it be (the same). So, $AD$ will also be the same
as $CB$. And as $AC$ will be to $CB$, so $BD$ (will be) to $DA$.
And $AB$ will (thus) also be divided at $D$ in the same (manner) as the division at $C$. The very opposite was assumed. Thus, $AC$ is not the
same as $DB$. So, on account of this,  points $C$ and $D$ are
not equally far from the point of bisection. Thus, by whatever (amount the sum of) the (squares) on $AC$ and $CB$ differs from (the sum of) the
(squares) on $AD$ and $DB$, twice the (rectangle contained)
by $AD$ and $DB$ also differs from twice the (rectangle contained)
by $AC$ and $CB$ by this (same amount)---on account of both (the sum of) the (squares) on $AC$ and
$CB$, plus twice the (rectangle contained) by $AC$ and $CB$, and
(the sum of) the (squares) on $AD$ and $DB$, plus twice the
(rectangle contained) by $AD$ and $DB$,  being equal to the (square) on $AB$ [Prop.~2.4]. But, (the sum of) the
(squares) on  $AC$ and $CB$ differs from (the sum of) the (squares)
on $AD$ and $DB$ by a rational (area). For (they are) both rational (areas).
Thus, twice the (rectangle contained) by $AD$ and $DB$ also differs from
twice the (rectangle contained) by $AC$ and $CB$ by a
rational (area, despite both) being medial (areas) [Prop.~10.21]. The very thing is absurd. For a
medial (area) cannot exceed a medial (area) by a rational (area) [Prop.~10.26].

Thus, a binomial (straight-line) cannot be divided (into its component terms) at different points.
Thus, (it can be so divided) at one point only. (Which is) the very thing it
was required to show.}
\end{Parallel}
{\footnotesize\noindent$^\dag$ In other words, $k + k'^{1/2} = k'' + k'''^{1/2}$
has only one solution: {i.e.}, $k''=k$ and $k'''=k'$. Likewise,
$k^{1/2} + k'^{1/2} = k''^{1/2}+ k'''^{1/2}$ has only one solution:
{i.e.}, $k''=k$ and $k'''=k'$ (or, equivalently, $k''=k'$ and $k'''=k$).}

%%%%%%
% Prop 10.43
%%%%%%
\pdfbookmark[1]{Proposition 10.43}{pdf10.43}
\begin{Parallel}{}{} 
\ParallelLText{
\begin{center}
{\large \ggn{43}.}
\end{center}\vspace*{-7pt}

{\greekfont Ἡ ἐκ δύο μέσων πρώτη καθ ἓν μόνον σημεῖον διαιρεῖται.}\\

\epsfysize=0.3in
\centerline{\epsffile{Book10/fig042g.eps}}

{\greekfont Ἔστω ἐκ δύο μέσων πρώτη ἡ ΑΒ διῃρημένη κατὰ τὸ Γ, ὥστε
τὰς ΑΓ, ΓΒ μέσας εἶναι δυνάμει μόνον συμμέτρους ῥητὸν
περιεθούσας; λέγω, ὅτι ἡ ΑΒ κατ ἄλλο σημεῖον οὐ
διαιρεῖται.}

{\greekfont Εἰ γὰρ δυνατόν διῃρήσθω καὶ κατὰ τὸ Δ, ὥστε καὶ τὰς ΑΔ, ΔΒ μέσας
εἶναι δυνάμει μόνον συμμέτρους ῥητὸν περιεθούσας. ἐπεὶ
οὖν, ᾧ διαφέρει τὸ δὶς ὑπὸ τῶν ΑΔ, ΔΒ τοῦ δὶς
ὑπὸ τῶν ΑΓ, ΓΒ, τούτῳ διαφέρει τὰ ἀπὸ τῶν ΑΓ, ΓΒ τῶν ἀπὸ
τῶν ΑΔ, ΔΒ, ῥητῷ δὲ διαφέρει τὸ δὶς ὑπὸ τῶν ΑΔ, ΔΒ
τοῦ δὶς ὑπὸ τῶν ΑΓ, ΓΒ; ῥητὰ γὰρ ἀμφότερα; ῥητῷ
ἄρα διαφέρει καὶ τὰ ἀπὸ τῶν ΑΓ, ΓΒ τῶν ἀπὸ τῶν ΑΔ, ΔΒ μέσα
ὄντα; ὅπερ ἄτοπον.}

{\greekfont Οὐκ ἄρα ἡ ἐκ δύο μέσων πρώτη κατ ἄλλο καὶ ἄλλο
σημεῖον διαιρεῖται εἰς τὰ ὀνόματα; καθ ἓν ἄρα
μόνον; ὅπερ ἔδει δεῖχαι.}}

\ParallelRText{
\begin{center}
{\large Proposition 43}
\end{center}

A first bimedial (straight-line) can be divided (into its component terms) at
one point only.$^\dag$

\epsfysize=0.3in
\centerline{\epsffile{Book10/fig042e.eps}}

Let $AB$ be a first bimedial (straight-line) which has been divided at $C$, such
that $AC$ and $CB$ are medial (straight-lines), commensurable
in square only, (and) containing a rational (area) [Prop.~10.37].  I say that $AB$ cannot be (so) divided
at another point.

For, if possible, let it also be divided at $D$, such that $AD$ and
$DB$ are also medial (straight-lines), commensurable in square only, (and)
containing a rational (area). Since, therefore, by whatever (amount) twice the
(rectangle contained) by $AD$ and $DB$ differs from twice the (rectangle
contained) by $AC$ and $CB$, (the sum of) the (squares)
on $AC$ and $CB$ differs from (the sum of) the (squares) on $AD$ and
$DB$ by this (same amount) [Prop.~10.41~lem.]. And twice the (rectangle
contained) by  $AD$ and $DB$ differs from twice the (rectangle contained)
by $AC$ and $CB$ by a rational (area). For (they are) both
rational (areas). (The sum of) the (squares) on $AC$ and $CB$
thus differs from (the sum of) the (squares) on $AD$ and $DB$ by a
rational (area, despite both) being medial (areas). The very thing is
absurd [Prop.~10.26].

Thus,  a first bimedial (straight-line) cannot be divided into its (component)
terms at different points. Thus, (it can be so divided) at one
point only.
(Which is) the very thing it was required to show.}
\end{Parallel}
{\footnotesize\noindent$^\dag$ In other words, $k^{1/4} + k^{3/4} = k'^{1/4} + k'^{3/4}$
has only one solution: {i.e.}, $k'=k$.} 

%%%%%%
% Prop 10.44
%%%%%%
\pdfbookmark[1]{Proposition 10.44}{pdf10.44}
\begin{Parallel}{}{} 
\ParallelLText{
\begin{center}
{\large \ggn{44}.}
\end{center}\vspace*{-7pt}

{\greekfont Ἡ ἐκ δύο μέσων δευτέρα καθ ἓν μόνον σημεῖον διαιρεῖται.}

{\greekfont Ἔστω ἐκ δύο μέσων δευτέρα ἡ ΑΒ διῃρημένη κατὰ τὸ Γ, ὥστε
τὰς ΑΓ, ΓΒ μέσας εἶναι δυνάμει μόνον συμμέτρους μέσον
περιεθούσας; φανερὸν δή, ὅτι τὸ Γ οὐκ ἔστι κατὰ τῆς διθοτομίας,
ὅτι οὐκ εἰσὶ μήκει σύμμετροι. λέγω, ὅτι ἡ ΑΒ κατ ἄλλο
σημεῖον οὐ διαιρεῖται.}\\~\\~\\

\epsfysize=1.5in
\centerline{\epsffile{Book10/fig044g.eps}}

{\greekfont Εἰ γὰρ δυνατόν, διῃρήσθω καὶ κατὰ τὸ Δ, ὥστε τὴν ΑΓ τῇ
ΔΒ μὴ εἶναι τὴν α ὐτήν, ἀλλὰ μείζονα καθ ὑπόθεσιν τὴν ΑΓ;
δῆλον δή, ὅτι καὶ τὰ ἀπὸ τῶν ΑΔ, ΔΒ, ὡς ἐπάνω ἐδείχαμεν,
ἐλάσσονα τῶν ἀπὸ τῶν ΑΓ, ΓΒ; καὶ τὰς ΑΔ, ΔΒ μέσας εἶναι
δυνάμει μόνον συμμέτρους μέσον περιεθούσας. καὶ ἐκκείσθω
ῥητὴ ἡ ΕΖ, καὶ τῷ μὲν ἀπὸ τῆς ΑΒ ἴσον παρὰ τὴν ΕΖ
παραλληλόγραμμον ὀρθογώνιον παραβεβλήσθω τὸ ΕΚ, τοῖς
δὲ ἀπὸ τῶν ΑΓ, ΓΒ ἴσον ἀφῃρήσθω τὸ ΕΗ; λοιπὸν ἄρα τὸ ΘΚ
ἴσον ἐστὶ τῷ δὶς ὑπὸ τῶν ΑΓ, ΓΒ. πάλιν δὴ τοῖς ἀπὸ
τῶν ΑΔ, ΔΒ, ἅπερ ἐλάσσονα ἐδείθθη τῶν ἀπὸ τῶν ΑΓ, ΓΒ, ἴσον
ἀφῃρήσθω τὸ ΕΛ; καὶ λοιπὸν ἄρα τὸ ΜΚ ἴσον τῷ δὶς ὑπὸ τῶν
ΑΔ, ΔΒ. καὶ ἐπεὶ μέσα ἐστὶ τὰ ἀπὸ τῶν ΑΓ, ΓΒ, μέσον ἄρα [καὶ]
τὸ ΕΗ. καὶ παρὰ ῥητὴν τὴν ΕΖ παράκειται; ῥητὴ ἄρα ἐστὶν
ἡ ΕΘ καὶ ἀσύμμετρος τῇ ΕΖ μήκει. διὰ τὰ α ὐτὰ δὴ  καὶ ἡ ΘΝ
ῥητή ἐστι καὶ ἀσύμμετρος τῇ ΕΖ μήκει. καὶ ἐπεὶ αἱ
ΑΓ, ΓΒ μέσαι εἰσὶ δυνάμει μόνον σύμμετροι, ἀσύμμετρος ἄρα
ἐστὶν ἡ ΑΓ τῇ ΓΒ μήκει. ὡς δὲ ἡ ΑΓ πρὸς τὴν ΓΒ, οὕτως
τὸ ἀπὸ τῆς ΑΓ πρὸς τὸ ὑπὸ τῶν ΑΓ, ΓΒ; ἀσύμμετρον
ἄρα ἐστὶ τὸ ἀπὸ τῆς ΑΓ τῷ ὑπὸ τῶν ΑΓ, ΓΒ. ἀλλὰ
τῷ μὲν ἀπὸ τῆς ΑΓ σύμμετρά ἐστι τὰ ἀπὸ τῶν ΑΓ, ΓΒ; δυνάμει
γάρ εἰσι σύμμετροι αἱ ΑΓ, ΓΒ. τῷ δὲ ὑπὸ τῶν ΑΓ, ΓΒ σύμμετρόν
ἐστι τὸ δὶς ὑπὸ τῶν ΑΓ, ΓΒ. καὶ τὰ ἀπὸ τῶν ΑΓ, ΓΒ ἄρα ἀσύμμετρά ἐστι τῷ δὶς ὑπὸ τῶν ΑΓ, ΓΒ. ἀλλὰ τοῖς μὲν ἀπὸ
τῶν ΑΓ, ΓΒ ἴσον ἐστὶ τὸ ΕΗ, τῷ δὲ δὶς ὑπὸ τῶν ΑΓ, ΓΒ
ἴσον τὸ ΘΚ; ἀσύμμετρον ἄρα ἐστὶ τὸ ΕΗ τῷ ΘΚ; ὥστε καὶ ἡ
ΕΘ τῇ ΘΝ ἀσύμμετρός ἐστι μήκει. καί εἰσι ῥηταί; αἱ ΕΘ, ΘΝ
ἄρα ῥηταί εἰσι δυνάμει μόνον σύμμετροι.
ἐὰν δὲ δύο ῥηταὶ δυνάμει μόνον σύμμετροι
 συντεθῶσιν, ἡ ὅλη 
ἄλογός ἐστιν ἡ καλουμένη ἐκ δύο ὀνομάτων;
ἡ ΕΝ ἄρα ἐκ δύο ὀνομάτων ἐστὶ διῃρημένη κατὰ τὸ
Θ. κατὰ τὰ α ὐτὰ δὴ δειθθήσονται καὶ αἱ ΕΜ, ΜΝ
ῥηταὶ δυνάμει μόνον σύμμετροι; καὶ ἔσται ἡ ΕΝ ἐκ δύο
ὀνομάτων κατ ἄλλο καὶ ἄλλο διῃρημένη τό τε Θ καὶ τὸ Μ, καὶ
οὐκ ἔστιν ἡ ΕΘ τῇ ΜΝ ἡ α ὐτή, ὅτι
τὰ ἀπὸ τῶν ΑΓ, ΓΒ μείζονά ἐστι τῶν ἀπὸ τῶν ΑΔ, ΔΒ.
ἀλλὰ τὰ ἀπὸ τῶν ΑΔ, ΔΒ μείζονά ἐστι τοῦ δὶς ὑπὸ
 ΑΔ, ΔΒ; πολλῷ ἄρα καὶ τὰ ἀπὸ τῶν ΑΓ, ΓΒ,
 τουτέστι τὸ ΕΗ, μεῖζόν ἐστι τοῦ δὶς ὑπὸ τῶν ΑΔ, ΔΒ, τουτέστι
 τοῦ 
  ΜΚ; ὥστε καὶ ἡ ΕΘ τῆς ΜΝ
μείζων ἐστίν. ἡ ἄρα ΕΘ τῇ ΜΝ οὐκ ἔστιν ἡ α ὐτή;
ὅπερ ἔδει δεῖχαι.}}

\ParallelRText{
\begin{center}
{\large Proposition 44}
\end{center}

A second bimedial (straight-line)  can be divided
(into its component terms)
at one point only.$^\dag$

Let $AB$ be a second bimedial (straight-line) which has  been divided at $C$, so that $AC$ and $BC$ are medial (straight-lines), commensurable in square
only, (and) containing a medial (area) [Prop.~10.38].
So, (it is) clear that $C$ is not 
(located) at the point of bisection, since ($AC$ and $BC$)
are not commensurable in length. I say that $AB$ cannot be (so) divided at another point.

\epsfysize=1.5in
\centerline{\epsffile{Book10/fig044e.eps}}

For, if possible, let it also have  been (so) divided at $D$, so that $AC$ is not the
same as $DB$, but $AC$ (is), by hypothesis, greater. So, (it is)  clear that
(the sum of) the (squares) on $AD$ and $DB$ is also less than (the sum of)
the (squares) on $AC$ and $CB$, as we showed above [Prop.~10.41~lem.].  And $AD$ and $DB$ are medial (straight-lines), commensurable in square only, (and)
containing a medial (area). And let the rational (straight-line) $EF$
be laid down. And let the rectangular parallelogram $EK$, equal to the (square) on $AB$, be applied to $EF$. And let $EG$, equal to
(the sum of) the (squares) on $AC$ and $CB$, be cut  off (from $EK$).
Thus, the remainder, $HK$, is equal to twice the (rectangle contained)
by $AC$ and $CB$ [Prop.~2.4]. So, again,
let $EL$, equal to  (the sum of) the (squares) on  $AD$ and $DB$---which was shown (to be) less than (the sum of) the (squares) on $AC$ and
$CB$---be cut off (from $EK$). And, thus, the remainder, $MK$,
(is) equal to twice the (rectangle contained) by $AD$ and $DB$. 
And since (the sum of) the (squares) on $AC$ and $CB$ is medial, $EG$ (is) thus [also]
medial. And it is applied to the rational (straight-line) $EF$. Thus, $EH$
is rational, and incommensurable in length with $EF$ [Prop.~10.22]. So, for the same (reasons), $HN$
is also rational, and incommensurable in length with $EF$. And since 
$AC$ and $CB$  are medial (straight-lines which are) commensurable
in square only, $AC$ is thus incommensurable in length with $CB$.
And as $AC$ (is) to $CB$, so the (square) on $AC$ (is) to the
(rectangle contained) by $AC$ and $CB$ [Prop.~10.21~lem.]. Thus, the (square) on $AC$ is
incommensurable with the (rectangle contained) by $AC$ and $CB$
[Prop.~10.11]. But, (the sum of) the (squares)
on $AC$ and $CB$ is commensurable with the (square) on $AC$.
For, $AC$ and $CB$ are commensurable in square [Prop.~10.15]. And
twice the (rectangle contained) by $AC$ and $CB$ is commensurable with the (rectangle contained) by $AC$ and $CB$ [Prop.~10.6]. And thus
(the sum of) the squares on $AC$ and $CB$ is incommensurable with twice
the (rectangle contained) by $AC$ and $CB$ [Prop.~10.13]. But, $EG$ is equal to (the sum
of) the (squares) on $AC$ and $CB$, and $HK$ equal to twice the (rectangle
contained) by $AC$ and $CB$. Thus, $EG$ is incommensurable
with $HK$. Hence, $EH$ is also incommensurable in length with  $HN$
[Props.~6.1, 10.11].
And (they are)  rational (straight-lines). Thus, $EH$ and $HN$ are rational
(straight-lines which are) commensurable in square only. And if two
rational (straight-lines which are) commensurable in square only are
added together then the whole (straight-line) is that irrational called
binomial [Prop.~10.36]. Thus, $EN$
is a binomial (straight-line) which has been divided (into its component terms) at $H$. So, according to
the same (reasoning), $EM$ and $MN$ can be shown (to be) rational
(straight-lines which are) commensurable in square only. And $EN$
will (thus) be a binomial (straight-line) which has been divided (into its component terms) at the different (points)
$H$ and $M$ (which is absurd [Prop.~10.42]).
And $EH$ is not the same as $MN$, since (the sum of) the
(squares) on $AC$ and $CB$ is greater than (the sum of) the (squares) on
$AD$ and $DB$. But, (the sum of) the (squares) on $AD$ and $DB$
is greater than twice the (rectangle contained) by $AD$ and $DB$ [Prop.~10.59~lem.]. Thus,  (the sum of) the (squares) on $AC$ and $CB$---that is to say, $EG$---is also much greater than twice the (rectangle contained) by 
$AD$ and $DB$---that is to say, $MK$. Hence, $EH$ is also greater
than $MN$ [Prop.~6.1]. Thus, $EH$ is not the same as $MN$. (Which is) the very thing
it was required to show.}
\end{Parallel}
{\footnotesize\noindent$^\dag$ In other words, $k^{1/4}+k'^{1/2}/k^{1/4}
= k''^{1/4}+k'''^{1/2}/k''^{1/4}$
has only one solution: {i.e.}, $k''=k$ and $k'''=k'$.}

%%%%%%
% Prop 10.45
%%%%%%
\pdfbookmark[1]{Proposition 10.45}{pdf10.45}
\begin{Parallel}{}{} 
\ParallelLText{\
\begin{center}
{\large\ggn{45}.}
\end{center}\vspace*{-7pt}

{\greekfont Ἡ μείζων κατὰ τὸ α ὐτὸ μόνον σημεῖον διαιρεῖται.}\\

\epsfysize=0.3in
\centerline{\epsffile{Book10/fig042g.eps}}

{\greekfont Ἔστω μείζων ἡ ΑΒ διῃρημένη κατὰ τὸ Γ, ὥστε τὰς ΑΓ, ΓΒ
δυνάμει ἀσυμμέτρους εἶναι ποιούσας τὸ μὲν συγκείμενον
ἐκ τῶν ἀπὸ τῶν ΑΓ, ΓΒ τετραγώνων ῥητόν, τὸ δ ὑπὸ
τῶν ΑΓ, ΓΒ μέσον; λέγω, ὅτι ἡ ΑΒ κατ ἄλλο σημεῖον οὐ
διαιρεῖται.}

{\greekfont Εἰ γὰρ δυνατόν, διῃρήσθω καὶ κατὰ τὸ Δ, ὥστε καὶ τὰς ΑΔ, ΔΒ δυνάμει
ἀσυμμέτρους εἶναι ποιούσας τὸ μὲν συγκείμενον ἐκ τῶν
ἀπὸ τῶν ΑΔ, ΔΒ ῥητόν, τὸ δ ὑπ α ὐτῶν μέσον. καὶ ἐπεί,
ᾧ διαφέρει τὰ ἀπὸ τῶν ΑΓ, ΓΒ τῶν ἀπὸ τῶν ΑΔ, ΔΒ,
τούτῳ διαφέρει καὶ τὸ δὶς ὑπὸ τῶν ΑΔ, ΔΒ τοῦ δὶς ὑπὸ
τῶν ΑΓ, ΓΒ, ἀλλὰ τὰ ἀπὸ τῶν ΑΓ, ΓΒ τῶν ἀπὸ τῶν
ΑΔ, ΔΒ ὑπερέθει ῥητῷ; ῥητὰ γὰρ ἀμφότερα; καὶ τὸ δὶς ὑπὸ
τῶν ΑΔ, ΔΒ ἄρα τοῦ δὶς ὑπὸ τῶν
ΑΓ, ΓΒ ὑπερέθει ῥητῷ μέσα ὄντα; ὅπερ ἐστὶν
ἀδύνατον. οὐκ ἄρα ἡ μείζων κατ ἄλλο καὶ ἄλλο
σημεῖον διαιρεῖται; κατὰ τὸ α ὐτὸ ἄρα μόνον διαιρεῖται;
ὅπερ ἔδει δεῖχαι.}}

\ParallelRText{
\begin{center}
{\large Proposition 45}
\end{center}

A major (straight-line) can only  be divided (into
its component terms) at the same
point.$^\dag$

\epsfysize=0.3in
\centerline{\epsffile{Book10/fig042e.eps}}

Let $AB$ be a major (straight-line) which has been divided at $C$, so
that $AC$ and $CB$ are incommensurable in square, making the sum
of the squares on $AC$ and $CB$ rational, and the (rectangle contained)
by $AC$ and $CD$ medial [Prop.~10.39].
I say that $AB$ cannot be (so) divided at another point.

For, if possible, let it also be divided at $D$, such that $AD$ and
$DB$ are also incommensurable in square, making the
sum of the (squares) on $AD$ and $DB$ rational, and the (rectangle contained) by them medial. And since, by whatever (amount the sum of)
the (squares) on $AC$ and $CB$ differs from (the sum of) the (squares)
on $AD$ and $DB$, 
twice the (rectangle contained) by $AD$ and $DB$ also differs from
twice the (rectangle contained) by $AC$ and $CB$ by this (same amount). But, (the sum of) the
(squares) on $AC$ and $CB$ exceeds (the sum of) the (squares) on
$AD$ and $DB$ by a rational (area). For (they are) both rational (areas). Thus, twice the
(rectangle contained) by $AD$ and $DB$ also exceeds twice the
(rectangle contained) by $AC$ and $CB$ by a rational (area), (despite both) being medial (areas).
The very thing is impossible [Prop.~10.26].
Thus, a major (straight-line) cannot be divided (into its component terms) at different points. Thus,
it can only be (so) divided at the same (point). (Which is) the very thing it was required to show.}
\end{Parallel}
{\footnotesize\noindent$^\dag$ In other words, $\sqrt{[1+k/(1+k^2)^{1/2}]/2} + \sqrt{[1-k/(1+k^2)^{1/2}]/2} =\sqrt{[1+k'/(1+k'^2)^{1/2}]/2} + \sqrt{[1-k'/(1+k'^2)^{1/2}]/2}$ has only one solution: {i.e.}, $k'=k$.}

%%%%%%
% Prop 10.46
%%%%%%
\pdfbookmark[1]{Proposition 10.46}{pdf10.46}
\begin{Parallel}{}{} 
\ParallelLText{
\begin{center}
{\large \ggn{46}.}
\end{center}\vspace*{-7pt}

{\greekfont Ἡ ῥητὸν καὶ μέσον δυναμένη καθ ἓν μόνον σημεῖον διαιρεῖται.}

\epsfysize=0.3in
\centerline{\epsffile{Book10/fig042g.eps}}

{\greekfont Ἔστω ῥητὸν καὶ μέσον δυναμένη ἡ ΑΒ διῃρημένη κατὰ τὸ Γ, ὥστε
τὰς ΑΓ, ΓΒ δυνάμει ἀσυμμέτρους εἶναι ποιούσας τὸ μὲν συγκείμενον ἐκ τῶν ἀπὸ τῶν ΑΓ, ΓΒ μέσον, τὸ δὲ δὶς ὑπὸ
τῶν ΑΓ, ΓΒ ῥητόν; λέγω, ὅτι ἡ ΑΒ κατ ἄλλο σημεῖον
οὐ διαιρεῖται.}

{\greekfont Εἰ γὰρ δυνατόν, διῃρήσθω καὶ κατὰ τὸ Δ, ὥστε καὶ τὰς ΑΔ, ΔΒ
δυνάμει ἀσυμμέτρους εἶναι ποιούσας τὸ μὲν συγκείμενον
ἐκ τῶν ἀπὸ τῶν ΑΔ, ΔΒ μέσον, τὸ δὲ δὶς ὑπὸ
τῶν ΑΔ, ΔΒ ῥητόν. ἐπεὶ οὖν, ᾧ διαφέρει τὸ δὶς ὑπὸ
τῶν ΑΓ, ΓΒ τοῦ δὶς ὑπὸ τῶν ΑΔ, ΔΒ, τούτῳ διαφέρει καὶ
τὰ ἀπὸ τῶν ΑΔ, ΔΒ τῶν ἀπὸ τῶν ΑΓ, ΓΒ, τὸ δὲ δὶς ὑπὸ
τῶν ΑΓ, ΓΒ 
 τοῦ δὶς ὑπὸ τῶν ΑΔ, ΔΒ ὑπερέθει
ῥητῷ, καὶ τὰ ἀπὸ τῶν ΑΔ, ΔΒ ἄρα τῶν ἀπὸ τῶν ΑΓ, ΓΒ
ὑπερέθει ῥητῷ μέσα ὄντα; ὅπερ ἐστὶν ἀδύνατον. οὐκ
ἄρα ἡ ῥητὸν καὶ μέσον δυναμένη κατ ἄλλο
καὶ ἄλλο σημεῖον διαιρεῖται. κατὰ ἓν ἄρα σημεῖον
διαιρεῖται; ὅπερ ἔδει δεῖχαι.}}

\ParallelRText{
\begin{center}
{\large Proposition 46}
\end{center}
The square-root of a rational plus a  medial
(area) can be divided (into its component terms) at  one point only.$^\dag$

\epsfysize=0.3in
\centerline{\epsffile{Book10/fig042e.eps}}

Let $AB$ be the square-root of a rational plus a medial (area) which has
been divided at $C$, so that $AC$ and $CB$ are incommensurable
in square, making the sum of the (squares) on $AC$ and $CB$ medial,
and twice the (rectangle contained) by $AC$ and $CB$ rational
[Prop.~10.40]. I say that $AB$ cannot be (so)
divided at another point.

For, if possible, let it also be divided at $D$, so that
$AD$ and $DB$ are also incommensurable in square, making the sum of 
the (squares) on $AD$ and $DB$ medial, and twice the (rectangle
contained) by $AD$ and $DB$ rational. Therefore, since by whatever
(amount) twice the (rectangle contained) by $AC$ and $CB$
differs from twice the (rectangle contained) by $AD$ and $DB$, (the sum of) the (squares) on  $AD$ and $DB$ also differs from (the sum of) the
(squares) on
$AC$ and $CB$ by this
(same amount). And twice the (rectangle contained) by $AC$ and $CB$
exceeds twice the (rectangle contained) by $AD$ and $DB$ by a
rational (area). (The sum of) the (squares)
on $AD$ and $DB$ thus also exceeds (the sum of) the (squares) on $AC$ and $CB$ by a rational (area),
(despite both) being medial (areas). The very thing is impossible [Prop.~10.26]. Thus, the square-root of a rational
plus a medial (area) cannot be divided (into its component terms) at different points. Thus, it
can  be (so) divided at one point (only). (Which is) the very thing it was required to
show.}
\end{Parallel}
{\footnotesize\noindent$^\dag$ In other words, $\sqrt{[(1+k^2)^{1/2}+k]/[2\,(1+k^2)]} +\sqrt{[(1+k^2)^{1/2}-k]/[2\,(1+k^2)]}=\sqrt{[(1+k'^2)^{1/2}+k']/[2\,(1+k'^2)]}$\\$ +\sqrt{[(1+k'^2)^{1/2}-k']/[2\,(1+k'^2)]}$
has only one solution: {i.e.}, $k'=k$.}

%%%%%%
% Prop 10.47
%%%%%%
\pdfbookmark[1]{Proposition 10.47}{pdf10.47}
\begin{Parallel}{}{} 
\ParallelLText{
\begin{center}
{\large \ggn{47}.}
\end{center}\vspace*{-7pt}

{\greekfont Ἡ δύο μέσα δυναμένη καθ ἓν μόνον σημεῖον διαιρεῖται.}\\

\epsfysize=2in
\centerline{\epsffile{Book10/fig047g.eps}}

{\greekfont Ἔστω [δύο μέσα δυναμένη] ἡ ΑΒ διῃρημένη κατὰ τὸ Γ, ὥστε τὰς
ΑΓ, ΓΒ δυνάμει ἀσυμμέτρους εἶναι ποιούσας τό τε συγκείμενον
ἐκ τῶν ἀπὸ τῶν ΑΓ, ΓΒ μέσον καὶ τὸ ὑπὸ τῶν ΑΓ, ΓΒ μέσον
καὶ ἔτι ἀσύμμετρον τῷ συγκειμένῳ ἐκ τῶν ἀπ α ὐτῶν.
λέγω, ὅτι ἡ ΑΒ κατ ἄλλο σημεῖον οὐ διαιρεῖται ποιοῦσα
τὰ προκείμενα.}

{\greekfont Εἰ γὰρ δυνατόν, διῃρήσθω κατὰ τὸ Δ, ὥστε πάλιν δηλονότι
τὴν ΑΓ τῇ ΔΒ μὴ εἶναι τὴν α ὐτήν, ἀλλὰ μείζονα καθ ὑπόθεσιν
τὴν ΑΓ, καὶ ἐκκείσθω ῥητὴ ἡ ΕΖ, καὶ παραβεβλήσθω παρὰ τὴν ΕΖ
τοῖς μὲν ἀπὸ τῶν ΑΓ, ΓΒ ἴσον τὸ ΕΗ, τῷ δὲ δὶς ὑπὸ τῶν
ΑΓ, ΓΒ ἴσον τὸ ΘΚ; ὅλον ἄρα τὸ ΕΚ ἴσον ἐστὶ τῷ ἀπὸ
τῆς ΑΒ  τετραγώνῳ. πάλιν δὴ παραβεβλήσθω παρὰ τὴν ΕΖ τοῖς ἀπὸ
τῶν ΑΔ, ΔΒ ἴσον τὸ ΕΛ; λοιπὸν ἄρα τὸ δὶς ὑπὸ τῶν ΑΔ, ΔΒ
λοιπῷ τῷ ΜΚ ἴσον ἐστίν. καὶ ἐπεὶ μέσον ὑπόκειται τὸ συγκείμενον ἐκ τῶν ἀπὸ τῶν ΑΓ, ΓΒ, μέσον ἄρα ἐστὶ καὶ
τὸ ΕΗ. καὶ παρὰ ῥητὴν τὴν ΕΖ παράκειται; ῥητὴ ἄρα ἐστὶν
ἡ ΘΕ καὶ ἀσύμμετρος τῇ ΕΖ μήκει. διὰ τὰ α ὐτὰ δὴ καὶ ἡ ΘΝ
ῥητή ἐστι καὶ ἀσύμμετρος τῇ ΕΖ μήκει. καὶ ἐπεὶ ἀσύμμετρόν
ἐστι τὸ συγκείμενον ἐκ τῶν ἀπὸ τῶν ΑΓ, ΓΒ τῷ δὶς ὑπὸ τῶν ΑΓ, ΓΒ, καὶ τὸ ΕΗ ἄρα τῷ ΗΝ ἀσύμμετρόν ἐστιν; ὥστε
καὶ ἡ ΕΘ τῇ ΘΝ ἀσύμμετρός ἐστιν. καί εἰσι ῥηταί; αἱ ΕΘ,
ΘΝ ἄρα ῥηταί εἰσι δυνάμει μόνον σύμμετροι; ἡ ΕΝ ἄρα ἐκ
δύο ὀνομάτων ἐστὶ διῃρημένη κατὰ τὸ Θ. ὁμοίως
δὴ δείχομεν, ὅτι καὶ κατὰ τὸ Μ διῄρηται. καὶ οὐκ ἔστιν ἡ ΕΘ
τῇ ΜΝ ἡ α ὐτή; ἡ ἄρα ἐκ δύο ὀνομάτων κατ ἄλλο καὶ ἄλλο
σημεῖον διῄρηται; ὅπερ ἐστίν ἄτοπον. οὐκ ἄρα ἡ δύο
μέσα δυναμένη κατ ἀλλο καὶ ἄλλο σημεῖον διαιρεῖται;
καθ ἓν ἄρα μόνον [σημεῖον] διαιρεῖται.}}

\ParallelRText{
\begin{center}
{\large Proposition 47}
\end{center}

The square-root of (the sum of) two medial (areas) can be
divided (into its component terms) at one point only.$^\dag$

\epsfysize=2in
\centerline{\epsffile{Book10/fig047e.eps}}

Let $AB$ be [the square-root of (the sum of) two medial (areas)] which has been divided
at $C$, such that $AC$ and $CB$ are incommensurable in square, making the
sum of the (squares) on $AC$ and $CB$ medial, and the (rectangle contained) by $AC$ and $CB$ medial, and, moreover, incommensurable
with the sum of the (squares) on ($AC$ and $CB$) [Prop.~10.41]. I say that $AB$ cannot be divided
at another point fulfilling the prescribed (conditions).

For, if possible, let it be divided at $D$, such that $AC$ is again
manifestly not the same as $DB$, but $AC$ (is), by hypothesis, greater.
And let the rational (straight-line) $EF$ be laid down. And let $EG$,
equal to (the sum of) the (squares) on $AC$ and $CB$, and $HK$, equal to twice the (rectangle contained) by $AC$ and $CB$,
be applied to $EF$. Thus, the whole of $EK$ is equal to the
square on $AB$ [Prop.~2.4]. So, again, let $EL$, equal to (the sum of) the (squares)
on $AD$ and $DB$, be applied to $EF$. Thus, the remainder---twice
the (rectangle contained) by $AD$ and $DB$---is equal to the remainder, $MK$. And since the sum of the (squares) on $AC$ and $CB$
was assumed (to be) medial, $EG$ is also medial. And it is applied
to the rational (straight-line) $EF$. $HE$ is thus rational, and incommensurable in length with $EF$ [Prop.~10.22].  So, for the same (reasons), 
$HN$ is also rational, and incommensurable in length with $EF$.
And since the sum of the (squares) on $AC$ and $CB$ is incommensurable
with twice the (rectangle contained) by $AC$ and $CB$, $EG$ is thus also
incommensurable with $GN$. Hence, $EH$ is also incommensurable with
$HN$ [Props.~6.1, 10.11].
And they are (both) rational (straight-lines). Thus, $EH$ and $HN$ are rational
(straight-lines which are) commensurable in  square only. Thus, $EN$
is a binomial (straight-line) which has been divided (into its component
terms) at $H$ [Prop.~10.36]. So, similarly, we can show that
it has also been (so) divided at $M$. And $EH$ is not the same  as $MN$. Thus,
a binomial (straight-line) has been divided (into its component
terms) at different points. The
very thing is absurd [Prop.~10.42]. Thus,
the square-root of (the sum of) two medial (areas) cannot be divided (into its component terms) at different
points. Thus, it can be (so) divided at one [point] only.}
\end{Parallel}
{\footnotesize\noindent$^\dag$ In other words, $k'^{1/4}\sqrt{[1+k/(1+k^2)^{1/2}]/2}
+k'^{1/4}\sqrt{[1-k/(1+k^2)^{1/2}]/2}=k'''^{1/4}\sqrt{[1+k''/(1+k''^2)^{1/2}]/2}$\\$
+k'''^{1/4}\sqrt{[1-k''/(1+k''^2)^{1/2}]/2}$ has only one solution: {i.e.}, $k''=k$ and $k'''=k'$.}

%%%%%%%%
% Definitions 2
%%%%%%%%
\pdfbookmark[1]{Definitions II}{def10.2}
\begin{Parallel}{}{} 
\ParallelLText{
\begin{center}
\large{{\greekfont Ὅροι δεύτεροι}.}
\end{center}\vspace*{-7pt}

\ggn{5}.~{\greekfont Ὑποκειμένης ῥητῆς καὶ τῆς ἐκ δύο ὀνομάτων
διῃρημένης εἰς τὰ ὀνόματα, ἧς τὸ μεῖζον ὄνομα
τοῦ ἐλάσσονος μεῖζον δύναται τῷ ἀπὸ συμμέτρου ἑαυτῇ
μήκει, ἐὰν μὲν τὸ μεῖζον ὄνομα σύμμετρον ᾖ μήκει
τῇ ἐκκειμένῃ ῥητῇ, καλείσθω [ἡ ὅλη] ἐκ δύο ὀνομάτων
πρώτη.}

\ggn{6}.~{\greekfont Ἐὰν δὲ τὸ ἐλάσσον ὄνομα σύμμετρον ᾖ μήκει τῇ ἐκκειμένῃ ῥητῇ, καλείσθω ἐκ δύο ὀνομάτων δευτέρα.}

\ggn{7}.~{\greekfont Ἐὰν δὲ μ ηδέτερον τῶν ὀνομάτων σύμμετρον ᾖ
μήκει τῇ ἐκκειμένῃ ῥητῇ, καλείσθω ἐκ δύο ὀνομάτων
τρίτη.}

\ggn{8}.~{\greekfont Πάλιν δὴ ἐὰν τὸ μεῖζον ὄνομα [τοῦ ἐλάσσονος]
μεῖζον δύνηται τῷ ἀπὸ ἀσυμμέτρου ἑαυτῇ μήκει,
ἐὰν μὲν τὸ μεῖζον ὄνομα σύμμετρον ᾖ μήκει τῇ
ἐκκειμένῃ ῥητῇ, καλείσθω ἐκ δύο ὀνομάτων τετάρτη.}

\ggn{9}.~{\greekfont Ἐὰν δὲ τὸ ἔλασσον, πέμπτη.}

\ggn{10}.~{\greekfont Ἐὰν δὲ μ ηδέτερον, ἕκτη.}}

\ParallelRText{
\begin{center}
{\large Definitions II}
\end{center}

5.~Given a rational (straight-line), and  a binomial (straight-line) which has been divided into its (component) terms, of which the square on the greater term
is larger than (the square on) the lesser  by the (square) on (some straight-line) commensurable in length with (the greater) then,
if the greater term is commensurable in length with the  rational (straight-line previously) laid out,
 let [the whole] (straight-line) be called a first binomial (straight-line).
 
6.~And if the lesser term is commensurable in length with the  rational (straight-line previously) laid
out then let (the whole straight-line) be called a
second binomial (straight-line).

7.~And if neither of the terms is commensurable in length with the 
rational (straight-line previously) laid out then let (the whole straight-line) be called a third
binomial (straight-line).

8.~So, again, if the square on the greater term is larger than (the
square on) [the lesser] by the (square) on (some straight-line)
incommensurable in length with (the greater) then, if the
greater term is commensurable in length with the  rational
(straight-line previously) laid out, let (the whole straight-line) be called a fourth
binomial (straight-line).

9.~And if the lesser (term is commensurable), a fifth (binomial straight-line).

10.~And if neither (term is commensurable), a sixth (binomial straight-line).}
\end{Parallel}

%%%%
%10.48
%%%%
\pdfbookmark[1]{Proposition 10.48}{pdf10.48}
\begin{Parallel}{}{}
\ParallelLText{
\begin{center}
{\large\ggn{48}.}
\end{center}\vspace*{-7pt}

{\greekfont Εὑρεῖν τὴν ἐκ δύο ὀνομάτων πρώτην.}

{\greekfont Ἐκκείσθωσαν δύο ἀριθμοὶ οἱ ΑΓ, ΓΒ, ὥστε τὸν συγκείμενον
ἐχ α ὐτῶν τὸν ΑΒ πρὸς μὲν τὸν ΒΓ λόγον ἔχειν,
ὃν τετράγωνος ἀριθμὸς πρὸς τετράγωνον ἀριθμόν, πρὸς δὲ τὸν
ΓΑ λόγον μὴ ἔχειν, ὃν τετράγωνος ἀριθμὸς πρὸς τετράγωνον
ἀριθμόν, καὶ ἐκκείσθω τις ῥητὴ ἡ Δ, καὶ τῇ Δ σύμμετρος
ἔστω μήκει ἡ ΕΖ. ῥητὴ ἄρα ἐστὶ καὶ ἡ ΕΖ. καὶ γεγονέτω
ὡς ὁ ΒΑ ἀριθμὸς πρὸς τὸν ΑΓ, οὕτως τὸ ἀπὸ τῆς ΕΖ
πρὸς τὸ ἀπὸ τῆς ΖΗ. ὁ δὲ ΑΒ πρὸς τὸν ΑΓ λόγον ἔχει,
ὃν ἀριθμὸς πρὸς ἀριθμόν; καὶ τὸ ἀπὸ τῆς ΕΖ ἄρα πρὸς τὸ
ἀπὸ τῆς ΖΗ λόγον ἔχει, ὃν ἀριθμὸς πρὸς ἀριθμόν; ὥστε
σύμμετρόν ἐστι τὸ ἀπὸ τῆς ΕΖ τῷ ἀπὸ τῆς ΖΗ. καὶ ἐστι ῥητὴ
ἡ ΕΖ; ῥητὴ ἄρα καὶ ἡ ΖΗ.
καὶ ἐπεὶ
ὁ ΒΑ πρὸς τὸν ΑΓ λόγον οὐκ ἔχει, ὃν τετράγωνος ἀριθμὸς
πρὸς τετράγωνον ἀριθμόν, οὐδὲ τὸ ἀπὸ τῆς ΕΖ ἄρα πρὸς τὸ
ἀπὸ τῆς ΖΗ λόγον ἔχει, ὃν τετράγωνος ἀριθμὸς πρὸς τετράγωνον
ἀριθμόν; ἀσύμμετρος ἄρα ἐστὶν ἡ ΕΖ τῇ ΖΗ μήκει. αἱ ΕΖ,
ΖΗ ἄρα ῥηταί εἰσι δυνάμει μόνον σύμμετροι; ἐκ δύο ἄρα
ὀνομάτων ἐστὶν ἡ ΕΗ.
λέγω, ὅτι καὶ πρώτη.}\\~\\~\\~\\~\\

\epsfysize=1.in
\centerline{\epsffile{Book10/fig048g.eps}}

{\greekfont Ἐπεὶ γάρ ἐστιν ὡς ὁ ΒΑ ἀριθμὸς πρὸς τὸν ΑΓ, οὕτως
τὸ ἀπὸ τῆς ΕΖ πρὸς τὸ ἀπὸ τῆς ΖΗ, μείζων δὲ ὁ ΒΑ τοῦ
ΑΓ, μεῖζον ἄρα καὶ τὸ ἀπὸ τῆς ΕΖ τοῦ ἀπὸ τῆς ΖΗ. ἔστω
οὖν τῷ ἀπὸ τῆς ΕΖ ἴσα τὰ ἀπὸ τῶν ΖΗ, Θ. καὶ
ἐπεί ἐστιν ὡς ὁ ΒΑ πρὸς τὸν ΑΓ, οὕτως τὸ ἀπὸ τῆς
ΕΖ πρὸς τὸ ἀπὸ τῆς ΖΗ, ἀναστρέψαντι ἄρα ἐστὶν ὡς ὁ ΑΒ
πρὸς τὸν ΒΓ, οὕτως τὸ ἀπὸ τῆς ΕΖ πρὸς τὸ ἀπὸ τῆς Θ. ὁ δὲ
ΑΒ πρὸς τὸν ΒΓ λόγον ἔχει, ὃν τετράγωνος ἀριθμὸς πρὸς τετράγωνον ἀριθμόν. καὶ τὸ ἀπὸ τῆς ΕΖ ἄρα πρὸς τὸ ἀπὸ τῆς Θ
λόγον ἔχει, ὃν τετράγωνος ἀριθμὸς πρὸς τετράγωνον ἀριθμόν.
σύμμετρος ἄρα ἐστὶν ἡ ΕΖ τῇ Θ μήκει; ἡ ΕΖ
ἄρα τῆς ΖΗ μεῖζον δύναται τῷ ἀπὸ συμμέτρου ἑαυτῇ. καί
εἰσι ῥηταὶ αἱ ΕΖ, ΖΗ, καὶ σύμμετρος ἡ ΕΖ τῇ Δ μήκει.}

{\greekfont Ἠ ΕΗ ἄρα ἐκ δύο ὀνομάτων ἐστὶ πρώτη; ὅπερ
ἔδει δεῖχαι.}}

\ParallelRText{
\begin{center}
{\large Proposition 48}
\end{center}

To find a first binomial (straight-line).

Let  two numbers $AC$ and $CB$ be laid down such that their sum
$AB$ has to $BC$ the ratio  which (some) square number (has) to (some)
square number, and does not have to $CA$ the ratio which (some)
square number (has) to (some) square number [Prop.~10.28~lem.~I]. And let some rational (straight-line) $D$ be laid down. And let $EF$ be commensurable in length with $D$. $EF$ is thus also rational [Def.~10.3]. 
And let it be contrived that as the number $BA$ (is) to $AC$, so the
(square) on $EF$ (is) to the (square) on $FG$ [Prop.~10.6~corr.]. And $AB$ has to $AC$ the ratio
which (some) number (has) to (some) number. Thus, the (square) on $EF$
also has to the (square) on $FG$ the ratio which (some) number
(has) to (some) number. Hence, the (square) on $EF$ is commensurable
with the (square) on $FG$ [Prop.~10.6].  And
$EF$ is rational. Thus, $FG$ (is) also rational. And since $BA$ does not
have to $AC$ the ratio which (some) square number (has) to (some)
square number, thus the (square) on $EF$ does not have to the (square)
on $FG$ the ratio which (some) square number (has) to (some) square number either. Thus, $EF$ is incommensurable in length with $FG$ [Prop~10.9]. $EF$ and $FG$ are thus rational
(straight-lines which are) commensurable in square only. Thus, $EG$
is a binomial (straight-line) [Prop.~10.36].
I say that (it is) also a first (binomial straight-line).

\epsfysize=1.in
\centerline{\epsffile{Book10/fig048e.eps}}

For since as the number $BA$ is to $AC$, so the (square) on $EF$ (is)
to the (square) on $FG$, and $BA$ (is) greater than $AC$, the
(square) on $EF$ (is) thus also greater than the (square) on $FG$ [Prop.~5.14]. Therefore, let (the sum of) the (squares)
on $FG$ and $H$ be equal to the (square) on $EF$. And since as $BA$
is to $AC$, so the (square) on $EF$ (is) to the (square) on $FG$, thus,
via conversion, as $AB$ is to $BC$, so the (square) on $EF$ (is) to
the (square) on $H$ [Prop.~5.19~corr.]. 
And $AB$ has to $BC$ the ratio which (some) square number (has) to
(some) square number. Thus, the (square) on $EF$ also has to
the (square) on $H$ the ratio which (some) square number (has) to
(some) square number. Thus,
$EF$ is commensurable in length with $H$  [Prop.~10.9]. Thus, the square on $EF$
is greater than (the square on) $FG$ by the (square) on (some straight-line)
commensurable (in length) with ($EF$). And $EF$ and $FG$ are rational (straight-lines). And $EF$ (is) commensurable in length with $D$.

Thus, $EG$ is a first binomial (straight-line) [Def.~10.5].$^\dag$
(Which is) the very thing it was required to show.}
\end{Parallel}
{\footnotesize\noindent $^\dag$If the rational straight-line has unit length then the length of a first binomial straight-line
is  $k+k\sqrt{1-k'^{\,2}}$. This, and the first apotome,
whose length is $k-k\,\sqrt{1-k'^{\,2}}$ [Prop.~10.85],
are the roots of $x^2- 2\,k\,x+k^2\,k'^{\,2}=0$.}  

%%%%
%10.49
%%%%
\pdfbookmark[1]{Proposition 10.49}{pdf10.49}
\begin{Parallel}{}{}
\ParallelLText{
\begin{center}
{\large\ggn{49}.}
\end{center}\vspace*{-7pt}

{\greekfont Εὑρεῖν τὴν ἐκ δύο ὀνομάτων δευτέραν.}

{\greekfont Ἐκκείσθωσαν δύο ἀριθμοὶ οἱ ΑΓ, ΓΒ, ὥστε τὸν συγκείμενον ἐχ
α ὐτῶν τὸν ΑΒ πρὸς μὲν τὸν ΒΓ λόγον ἔχειν, ὃν τετράγωνος
ἀριθμὸς πρὸς τετράγωνον ἀριθμόν, πρὸς δὲ τὸν ΑΓ λόγον
μὴ ἔχειν, ὃν τετράγωνος ἀριθμὸς πρὸς τετράγωνον ἀριθμόν,
καὶ ἐκκείσθω ῥητὴ ἡ Δ, καὶ τῇ Δ σύμμετρος ἔστω ἡ ΕΖ μήκει;
ῥητὴ
ἄρα ἐστὶν ἡ ΕΖ. γεγονέτω δὴ καὶ ὡς
ὁ ΓΑ ἀριθμὸς πρὸς τὸν ΑΒ, οὕτως τὸ ἀπὸ τῆς ΕΖ πρὸς
τὸ ἀπὸ τῆς ΖΗ; σύμμετρον ἄρα ἐστὶ τὸ ἀπὸ τῆς ΕΖ τῷ
ἀπὸ τῆς ΖΗ. ῥητὴ ἄρα ἐστὶ καὶ ἡ ΖΗ. καὶ ἐπεὶ ὁ ΓΑ ἀριθμὸς
πρὸς τὸν ΑΒ λὸγον οὐκ ἔχει, ὃν τετράγωνος ἀριθμὸς πρὸς
τετράγωνον ἀριθμόν, οὐδὲ τὸ ἀπὸ τῆς ΕΖ πρὸς τὸ ἀπὸ
τῆς ΖΗ λόγον ἔχει, ὃν τετράγωνος ἀριθμὸς
πρὸς τετράγωνον ἀριθμόν. ἀσύμμετρος ἄρα ἐστὶν ἡ ΕΖ τῇ ΖΗ
μήκει; αἱ ΕΖ, ΖΗ ἄρα ῥηταί εἰσι δυνάμει μόνον σύμμετροι;
ἐκ δύο ἄρα ὀνομάτων ἐστὶν ἡ ΕΗ. δεικτέον δή, ὅτι καὶ
δευτέρα.}\\~\\~\\~\\~\\

\epsfysize=1.in
\centerline{\epsffile{Book10/fig049g.eps}}

{\greekfont Ἐπεὶ γὰρ ἀνάπαλίν ἐστιν ὡς ὁ ΒΑ ἀριθμὸς πρὸς τὸν ΑΓ,
οὕτως τὸ ἀπὸ τῆς ΗΖ πρὸς τὸ ἀπὸ τῆς ΖΕ, μείζων δὲ ὁ
ΒΑ τοῦ ΑΓ, μεῖζον ἄρα [καὶ] τὸ ἀπὸ τῆς ΗΖ τοῦ ἀπὸ τῆς
ΖΕ. ἔστω τῷ ἀπὸ τῆς ΗΖ ἴσα τὰ ἀπὸ τῶν ΕΖ, Θ; ἀναστρέψαντι
ἄρα ἐστὶν ὡς ὁ ΑΒ πρὸς τὸν ΒΓ, οὕτως τὸ ἀπὸ
τῆς ΖΗ πρὸς τὸ ἀπὸ τῆς Θ. ἀλλ ὁ ΑΒ πρὸς τὸν ΒΓ
λόγον ἔχει, ὃν τετράγωνος ἀριθμὸς πρὸς τετράγωνον ἀριθμόν;
καὶ τὸ ἀπὸ τῆς ΖΗ ἄρα πρὸς τὸ ἀπὸ τῆς Θ λόγον
ἔχει, ὃν τετράγωνος ἀριθμὸς πρὸς τετράγωνον ἀριθμόν.
σύμμετρος ἄρα ἐστὶν ἡ ΖΗ τῇ Θ μήκει; ὥστε ἡ ΖΗ
τῆς ΖΕ μεῖζον δύναται τῷ ἀπὸ συμμέτρου ἑαυτῇ. καί
εἰσι ῥηταὶ αἱ ΖΗ, ΖΕ δυνάμει μόνον σύμμετροι, καὶ τὸ ΕΖ
ἔλασσον ὄνομα τῇ ἐκκειμένῃ ῥητῇ σύμμετρόν
ἐστι τῇ Δ μήκει.}

{\greekfont Ἡ ΕΗ ἄρα ἐκ δύο ὀνομάτων ἐστὶ δευτέρα; ὅπερ ἔδει δεῖχαι.
}}

\ParallelRText{
\begin{center}
{\large Proposition 49}
\end{center}

To find a second binomial (straight-line).

Let the two numbers $AC$ and $CB$ be laid down such that their sum $AB$
has to $BC$ the ratio which (some) square number (has) to (some)
square number, and does not have to $AC$ the ratio which (some)
square number (has) to (some) square number [Prop.~10.28~lem.~I]. And let the rational (straight-line) $D$ be laid down. And let $EF$ be commensurable in length with $D$.
$EF$ is thus a rational (straight-line).
So, let it also be contrived that as the number $CA$ (is) to $AB$,
so the (square) on $EF$ (is) to the (square) on $FG$ [Prop.~10.6~corr.]. Thus, the (square) on $EF$
is commensurable with the (square) on $FG$ [Prop.~10.6]. Thus, $FG$ is also a rational
(straight-line). And since the number $CA$ does not have to $AB$ the ratio
which (some) square number (has) to (some) square number, the
(square) on $EF$ does not have to the (square) on $FG$ the ratio
which (some) square number (has) to (some) square number either.
Thus, $EF$  is incommensurable in length with $FG$ [Prop.~10.9]. $EF$ and $FG$
are thus rational (straight-lines which are) commensurable in square only.
Thus, $EG$ is a binomial (straight-line) [Prop.~10.36]. So, we must show
that (it is) also a second (binomial straight-line).

\epsfysize=1.in
\centerline{\epsffile{Book10/fig049e.eps}}

For since, inversely, as the number $BA$ is to $AC$, so the (square)
on $GF$ (is) to the (square) on $FE$ [Prop.~5.7 corr.],
and $BA$ (is) greater than $AC$, the (square) on $GF$ (is)
thus [also] greater than the (square) on $FE$ [Prop.~5.14]. Let (the sum of) the (squares) on $EF$ and $H$ be equal to the (square) on $GF$. Thus, via conversion, as 
$AB$ is to $BC$, so the (square) on $FG$ (is) to the (square) on $H$
[Prop.~5.19~corr.]. But, $AB$ has to $BC$ the
ratio which (some) square number (has) to (some) square number. Thus,
the (square) on $FG$ also has to the (square) on $H$ the ratio
which (some) square number (has) to (some) square number. Thus,
$FG$ is commensurable in length with $H$ [Prop.~10.9]. Hence, the square on $FG$ is
greater than (the square on) $FE$ by the (square) on (some straight-line)
commensurable in length with ($FG$). And $FG$ and $FE$ are rational (straight-lines which
are) commensurable in square only. And  the lesser term $EF$ is commensurable in length with the rational (straight-line) $D$ (previously)
laid down.

Thus, $EG$ is a second binomial (straight-line) [Def. 10.6].$^\dag$ (Which is) the very thing it
was required to show.}
\end{Parallel}
{\footnotesize\noindent $^\dag$ If the rational straight-line has unit length then the length of a second binomial straight-line
is  $k/\sqrt{1-k'^{\,2}}+k$. This, and the second apotome,
whose length is $k/\sqrt{1-k'^{\,2}}-k$ [Prop.~10.86],
are the roots of $x^2- (2\,k/\sqrt{1-k'^{\,2}})\,x+k^2\,[k'^{\,2}/(1-k'^{\,2})]=0$.}  

%%%%
%10.50
%%%%
\pdfbookmark[1]{Proposition 10.50}{pdf10.50}
\begin{Parallel}{}{}
\ParallelLText{
\begin{center}
{\large \ggn{50}.}
\end{center}\vspace*{-7pt}

{\greekfont Εὑρεῖν τὴν ἐκ δύο ὀνομάτων τρίτην.}

{\greekfont Ἐκκείσθωσαν δύο ἀριθμοὶ οἱ ΑΓ, ΓΒ, ὥστε τὸν συγκείμενον
ἐχ α ὐτῶν τὸν ΑΒ πρὸς μὲν τὸν ΒΓ λόγον ἔχειν, ὃν τετράγωνος
ἀριθμὸς πρὸς τετράγωνον ἀριθμόν, πρὸς δὲ τὸν ΑΓ λόγον μὴ
ἔχειν, ὃν τετράγωνος ἀριθμὸς πρὸς τετράγωνον ἀριθμόν. ἐκκείσθω
δέ τις καὶ ἄλλος μὴ τετράγωνος ἀριθμὸς ὁ Δ, καὶ πρὸς ἑκάτερον
τῶν ΒΑ, ΑΓ λόγον μὴ ἐθέτω, ὃν τετράγωνος ἀριθμὸς πρὸς τετράγωνον ἀριθμόν; καὶ ἐκκείσθω τις ῥητὴ εὐθεῖα ἡ Ε, καὶ
γεγονέτω ὡς ὁ Δ πρὸς τὸν ΑΒ, οὕτως τὸ ἀπὸ τῆς Ε πρὸς τὸ
ἀπὸ τῆς ΖΗ; σύμμετρον ἄρα ἐστὶ τὸ ἀπὸ τῆς Ε τῷ ἀπὸ
τῆς ΖΗ. καί ἐστι ῥητὴ ἡ Ε; ῥητὴ ἄρα ἐστὶ καὶ ἡ ΖΗ.
καὶ ἐπεὶ ὁ Δ πρὸς τὸν ΑΒ λόγον οὐκ ἔχει, ὃν τετράγωνος
ἀριθμὸς πρὸς τετράγωνον ἀριθμόν, οὐδὲ τὸ ἀπὸ τῆς Ε πρὸς
τὸ ἀπὸ τῆς ΖΗ λόγον ἔχει, ὃν τετράγωνος ἀριθμὸς πρὸς
τετράγωνον ἀριθμόν;
 ἀσύμμετρος ἄρα ἐστὶν ἡ Ε τῇ
ΖΗ μήκει. γεγονέτω δὴ πάλιν ὡς ἡ ΒΑ ἀριθμὸς πρὸς τὸν ΑΓ,
οὕτως τὸ ἀπὸ τῆς ΖΗ πρὸς τὸ ἀπὸ τῆς ΗΘ; σύμμετρον
ἄρα ἐστὶ τὸ ἀπὸ τῆς ΖΗ τῷ ἀπὸ τῆς ΗΘ. ῥητὴ δὲ ἡ ΖΗ;
ῥητὴ ἄρα καὶ ἡ ΗΘ. καὶ ἐπεὶ ὁ ΒΑ πρὸς τὸν ΑΓ λόγον οὐκ
ἔχει, ὃν τετράγωνος ἀριθμὸς πρὸς τετράγωνον ἀριθμόν, οὐδὲ
τὸ ἀπὸ τῆς ΖΗ πρὸς τὸ ἀπὸ τῆς ΘΗ λόγον ἔχει, ὃν τετράγωνος
ἀριθμὸς πρὸς τετράγωνον ἀριθμόν;
ἀσύμμετρος ἄρα ἐστὶν ἡ ΖΗ τῇ ΗΘ μήκει. αἱ ΖΗ, ΗΘ
ἄρα ῥηταί εἰσι δυνάμει μόνον σύμμετροι;
ἡ ΖΘ ἄρα ἐκ δύο ὀνομάτων ἐστίν. λέγω δή, ὅτι καὶ τρίτη.}\\~\\~\\~\\~\\~\\~\\

\epsfysize=1.1in
\centerline{\epsffile{Book10/fig050g.eps}}

{\greekfont Ἐπεὶ γάρ ἐστιν ὡς ὁ Δ πρὸς τὸν ΑΒ, οὕτως τὸ ἀπὸ τῆς
Ε πρὸς τὸ ἀπὸ τῆς ΖΗ, ὡς δὲ ὁ ΒΑ πρὸς τὸν ΑΓ, οὕτως
τὸ ἀπὸ τῆς ΖΗ πρὸς τὸ ἀπὸ τῆς ΗΘ, δι ἴσου ἄρα
ἐστὶν ὡς ὁ Δ πρὸς τὸν ΑΓ, οὕτως τὸ ἀπὸ τῆς Ε πρὸς
τὸ ἀπὸ τῆς ΗΘ. ὁ δὲ Δ πρὸς τὸν ΑΓ λόγον οὐκ ἔχει, ὃν
τετράγωνος ἀριθμὸς πρὸς τετράγωνον ἀριθμόν; οὐδὲ τὸ ἀπὸ
τῆς Ε ἄρα πρὸς τὸ ἀπὸ τῆς ΗΘ λόγον ἔχει, ὃν τετράγωνος
ἀριθμὸς πρὸς τετράγωνον ἀριθμόν; ἀσύμμετρος ἄρα ἐστὶν ἡ Ε
τῇ ΗΘ μήκει. καὶ ἐπεί ἐστιν ὡς ὁ ΒΑ πρὸς τὸν ΑΓ, οὕτως
τὸ ἀπὸ τῆς ΖΗ πρὸς τὸ ἀπὸ τῆς ΗΘ, μεῖζον
ἄρα τὸ ἀπὸ τῆς ΖΗ τοῦ ἀπὸ τῆς ΗΘ. ἔστω
οὖν τῷ ἀπὸ τῆς ΖΗ ἴσα τὰ ἀπὸ τῶν ΗΘ, Κ; ἀναστρέψαντι
ἄρα [ἐστὶν] ὡς ὁ ΑΒ πρὸς τὸν ΒΓ, οὕτως τὸ ἀπὸ τῆς ΖΗ
πρὸς τὸ ἀπὸ τῆς Κ. ὁ δὲ ΑΒ πρὸς τὸν ΒΓ λόγον ἔχει,
ὃν τετράγωνος ἀριθμὸς πρὸς τετράγωνον ἀριθμόν; καὶ
τὸ ἀπὸ τῆς ΖΗ ἄρα πρὸς τὸ ἀπὸ τῆς Κ λόγον ἔχει, ὃν
τετράγωνος ἀριθμὸς πρὸς τετράγωνον ἀριθμόν; σύμμετρος
ἄρα [ἐστὶν] ἡ ΖΗ τῇ Κ μήκει. ἡ ΖΗ ἄρα τῆς ΗΘ
μεῖζον δύναται τῷ ἀπὸ συμμέτρου ἑαυτῇ. καί εἰσιν αἱ
ΖΗ, ΗΘ ῥηταὶ δυνάμει μόνον σύμμετροι, καὶ οὐδετέρα α ὐτῶν
σύμμετρός ἐστι τῇ Ε μήκει.}

{\greekfont Ἡ ΖΘ ἄρα ἐκ δύο ὀνομάτων ἐστὶ τρίτη; ὅπερ ἔδει δεῖχαι.}}

\ParallelRText{
\begin{center}
{\large Proposition 50}
\end{center}

To find a third binomial (straight-line).

Let the two numbers $AC$ and $CB$ be laid down such that their sum
$AB$ has to $BC$ the ratio which (some) square number (has) to
(some) square number, and does not have to $AC$ the ratio which (some)
square number (has) to (some) square number. And let some other non-square number
$D$ also be laid down, and let it not have to each of $BA$ and $AC$ the
ratio which (some) square number (has) to (some) square number.
And let some rational straight-line $E$ be laid down, and let it be
contrived that as $D$ (is) to $AB$, so the (square) on $E$ (is) to
the (square) on $FG$ [Prop.~10.6~corr.]. Thus,
the (square) on $E$ is commensurable with the (square) on $FG$ [Prop.~10.6]. And $E$ is a rational (straight-line).
Thus, $FG$ is also a rational (straight-line). And since $D$ does not have to
$AB$ the ratio which (some) square number has to (some) square number,
the (square) on $E$ does not have to the (square) on $FG$ the ratio
which (some) square number (has) to (some) square number either.
$E$ is thus incommensurable in length with $FG$ [Prop.~10.9]. So, again, let it be contrived that
as the number $BA$ (is) to $AC$, so the (square) on  $FG$ (is) to the
(square) on $GH$ [Prop.~10.6~corr.]. Thus,
the (square) on $FG$ is commensurable with the (square) on $GH$
[Prop.~10.6]. And $FG$ (is) a rational (straight-line).
Thus, $GH$ (is) also a rational (straight-line). And since $BA$ does not have to $AC$ the ratio which (some) square number (has) to (some) square number, the (square) on $FG$  does not have to the (square) on $HG$
the ratio which (some) square number (has) to (some) square number either.
Thus, $FG$ is incommensurable in length with $GH$ [Prop.~10.9].  $FG$ and $GH$ are thus rational
(straight-lines which are) commensurable in square only. Thus,
$FH$ is a binomial (straight-line) [Prop.~10.36].
So, I say that (it is) also a third (binomial straight-line).

\epsfysize=1.1in
\centerline{\epsffile{Book10/fig050e.eps}}

For since as $D$ is to $AB$, so the (square) on $E$ (is) to the (square)
on $FG$, and as $BA$ (is) to $AC$, so the (square) on $FG$ (is) to
the (square) on $GH$, thus, via equality, as $D$ (is) to $AC$, so the
(square) on $E$ (is) to the (square) on $GH$ [Prop.~5.22]. And $D$ does not have to $AC$ the
ratio which (some) square number (has) to (some) square number.
Thus, the (square) on $E$ does not have to the (square) on $GH$ the ratio
which (some) square number (has) to (some) square number either.
Thus, $E$ is incommensurable in length with $GH$ [Prop.~10.9].  And since as $BA$ is to $AC$,
so the (square) on $FG$ (is) to the (square) on $GH$, the (square) on 
$FG$ (is) thus greater than the (square) on $GH$ [Prop.~5.14]. Therefore, let (the sum of) the (squares)
on $GH$ and $K$
be equal to the (square) on $FG$. Thus, via
conversion, as $AB$ [is] to $BC$, so the (square) on $FG$ (is) to the
(square) on $K$ [Prop.~5.19~corr.]. And $AB$
has to $BC$ the ratio which (some) square number (has) to (some) square
number. Thus, the (square) on $FG$ also has to the (square) on $K$
the ratio which (some) square number (has) to (some) square number.
Thus,  $FG$ [is] commensurable in length with $K$ [Prop.~10.9]. Thus, the square on
$FG$ is greater than (the square on) $GH$ by the (square) on
(some straight-line) commensurable (in length) with ($FG$). And $FG$ and $GH$
are rational (straight-lines which are) commensurable in square only,
and neither of them is commensurable in length with $E$.

Thus, $FH$ is a third binomial (straight-line) [Def. 10.7].$^\dag$ (Which is) the very thing it
was required to show.}
\end{Parallel}
{\footnotesize\noindent$^\dag$ If the rational straight-line has unit length then the length of a third binomial straight-line
is  $k^{1/2}\,(1+\sqrt{1-k'^{\,2}})$. This, and the third apotome,
whose length is $k^{1/2}\,(1-\sqrt{1-k'^{\,2}})$ [Prop.~10.87],
are the roots of $x^2- 2\,k^{1/2}\,x+k\,k'^{\,2}=0$.}

%%%%
%10.51
%%%%
\pdfbookmark[1]{Proposition 10.51}{pdf10.51}
\begin{Parallel}{}{}
\ParallelLText{
\begin{center}
{\large \ggn{51}.}
\end{center}\vspace*{-7pt}

{\greekfont Εὑρεῖν τὴν ἐκ δύο ὀνομάτων τετάρτην.}

\epsfysize=1.2in
\centerline{\epsffile{Book10/fig051g.eps}}

{\greekfont Ἐκκείσθωσαν δύο ἀριθμοὶ οἱ ΑΓ, ΓΒ, ὥστε τὸν ΑΒ πρὸς τὸν
ΒΓ λόγον μὴ ἔχειν μήτε μὴν πρὸς τὸν ΑΓ, ὃν τετράγωνος
ἀριθμὸς πρὸς τετράγωνον ἀριθμόν. καὶ ἐκκείσθω ῥητὴ ἡ Δ,
καὶ τῇ Δ σύμμετρος ἔστω μήκει ἡ ΕΖ; ῥητὴ ἄρα ἐστὶ
καὶ ἡ ΕΖ. καὶ γεγονέτω ὡς ὁ ΒΑ ἀριθμὸς πρὸς τὸν ΑΓ,
οὕτως τὸ ἀπὸ τῆς ΕΖ πρὸς τὸ ἀπὸ τῆς ΖΗ; σύμμετρον
ἄρα ἐστὶ τὸ ἀπὸ τῆς ΕΖ τῷ ἀπὸ τῆς ΖΗ; ῥητὴ ἄρα ἐστὶ
καὶ ἡ ΖΗ. καὶ ἐπεὶ ὁ ΒΑ πρὸς τὸν ΑΓ λόγον οὐκ ἔχει,
ὃν τετράγωνος ἀριθμὸς πρὸς τετράγωνον ἀριθμόν, οὐδὲ
τὸ ἀπὸ τῆς ΕΖ  πρὸς τὸ ἀπὸ τῆς ΖΗ λόγον ἔχει, ὃν τετράγωνος
ἀριθμὸς πρὸς τετράγωνον ἀριθμόν; ἀσύμμετρος ἄρα ἐστὶν
ἡ ΕΖ τῇ ΖΗ μήκει. αἱ ΕΖ, ΖΗ ἄρα ῥηταί εἰσι δυνάμει μόνον
σύμμετροι; ὥστε ἡ ΕΗ ἐκ δύο ὀνομάτων ἐστίν. λέγω δή,
ὅτι καὶ τετάρτη.}

{\greekfont Ἐπεὶ γάρ ἐστιν ὡς ὁ ΒΑ πρὸς τὸν ΑΓ, οὕτως
τὸ ἀπὸ τῆς ΕΖ πρὸς τὸ ἀπὸ τῆς ΖΗ [μείζων
δὲ ὁ ΒΑ τοῦ ΑΓ], μεῖζον ἄρα τὸ ἀπὸ τῆς ΕΖ τοῦ ἀπὸ
τῆς ΖΗ. ἔστω οὖν τῷ ἀπὸ τῆς ΕΖ ἴσα τὰ ἀπὸ τῶν ΖΗ,
Θ; ἀναστρέψαντι ἄρα ὡς ὁ ΑΒ ἀριθμὸς πρὸς τὸν ΒΓ, οὕτως
τὸ ἀπὸ τῆς ΕΖ πρὸς τὸ ἀπὸ τῆς Θ. ὁ δὲ ΑΒ πρὸς τὸν ΒΓ
λόγον οὐκ ἔχει, ὃν τετράγωνος ἀριθμὸς πρὸς τετράγωνον
ἀριθμόν; οὐδ ἄρα τὸ ἀπὸ τῆς ΕΖ πρὸς τὸ ἀπὸ τῆς Θ
λόγον ἔχει, ὃν τετράγωνος ἀριθμὸς πρὸς τετράγωνον ἀριθμόν.
ἀσύμμετρος ἄρα ἐστὶν ἡ ΕΖ τῇ Θ μήκει; ἡ ΕΖ ἄρα τῆς
ΗΖ μεῖζον δύναται τῷ ἀπὸ ἀσυμμέτρου ἑαυτῇ. καί
εἰσιν αἱ ΕΖ, ΖΗ ῥηταὶ δυνάμει μόνον σύμμετροι, καὶ
ἡ ΕΖ τῇ Δ σύμμετρός ἐστι μήκει.}

{\greekfont Ἡ ΕΗ ἄρα ἐκ δύο ὀνομάτων ἐστὶ τετάρτη; ὅπερ
ἔδει δεῖχαι.}}

\ParallelRText{
\begin{center}
{\large Proposition 51}
\end{center}

To find a fourth binomial (straight-line).

\epsfysize=1.2in
\centerline{\epsffile{Book10/fig051e.eps}}

Let the two numbers $AC$ and $CB$ be laid down such that $AB$ does not have to $BC$,  or to $AC$ either, the ratio which (some) square number
(has) to (some) square number [Prop.~10.28~lem.~I].
And let the rational (straight-line) $D$ be laid down. And let $EF$ be
commensurable in length with $D$. Thus, $EF$ is also a rational (straight-line). And let it be contrived that as the number $BA$ (is) to
$AC$, so the (square) on $EF$ (is) to the (square) on $FG$ [Prop.~10.6~corr.]. Thus, the (square) on $EF$
is commensurable with the (square) on $FG$ [Prop.~10.6]. Thus, $FG$ is also a rational (straight-line). And since $BA$ does not have to $AC$ the ratio which (some)
square number (has) to (some) square number,  the (square) on $EF$
does not have to the (square) on $FG$ the ratio which (some) square number
(has) to (some) square number either. Thus, $EF$ is incommensurable in length with $FG$ [Prop.~10.9]. Thus, $EF$ and
$FG$ are rational (straight-lines which are) commensurable in square only.
Hence, $EG$ is a binomial (straight-line) [Prop.~10.36]. So, I say that (it is) also
a fourth (binomial straight-line).

For since as $BA$ is to $AC$, so the (square) on $EF$ (is) to the (square) on
$FG$ [and $BA$ (is) greater than $AC$], the (square) on $EF$ (is) thus
greater than the (square) on $FG$ [Prop.~5.14]. Therefore, let (the sum of) the squares on
$FG$ and $H$ be equal to the (square) on $EF$. Thus, via conversion,
as the number $AB$ (is) to $BC$, so the (square) on $EF$ (is) to the
(square) on $H$ [Prop.~5.19~corr.]. And $AB$
does not have to $BC$ the ratio which (some) square number (has) to (some)
square number. Thus, the (square) on $EF$ does not have to the
(square) on $H$ the ratio which (some) square number (has) to (some) square
number either. Thus, $EF$ is incommensurable in length with $H$ [Prop.~10.9]. Thus, the square on $EF$ is greater than (the square on) $GF$ by the (square) on (some straight-line)
incommensurable (in length) with ($EF$). And $EF$ and $FG$ are rational (straight-lines which are) commensurable in square only.  And $EF$ is commensurable
in length with $D$.

Thus, $EG$ is a fourth binomial (straight-line) [Def. 10.8].$^\dag$ (Which is) the very thing
it was required to show.}
\end{Parallel}
{\footnotesize\noindent $^\dag$ If the rational straight-line has unit length then the length of a fourth binomial straight-line
is  $k\,(1+1/\sqrt{1+k'})$. This, and the fourth apotome,
whose length is $k\,(1-1/\sqrt{1+k'})$ [Prop.~10.88],
are the roots of $x^2- 2\,k\,x+k^2\,k'/(1+k')=0$.}

%%%%
%10.52
%%%%
\pdfbookmark[1]{Proposition 10.52}{pdf10.52}
\begin{Parallel}{}{}
\ParallelLText{
\begin{center}
{\large \ggn{52}.}
\end{center}\vspace*{-7pt}

{\greekfont Εὑρεῖν τὴν ἐκ δύο ὀνομάτων πέμπτην.}

\epsfysize=1.3in
\centerline{\epsffile{Book10/fig052g.eps}}

{\greekfont Ἐκκείσθωσαν δύο ἀριθμοὶ οἱ ΑΓ, ΓΒ, ὥστε τὸν ΑΒ πρὸς ἑκάτερον
α ὐτῶν λόγον μὴ ἔχειν, ὃν τετράγωνος ἀριθμὸς πρὸς τετράγωνον
ἀριθμόν, καὶ ἐκκείσθω ῥητή τις εὐθεῖα ἡ Δ, καὶ τῇ Δ σύμμετρος
ἔστω [μήκει] ἡ ΕΖ; ῥητὴ ἄρα ἡ ΕΖ. καὶ γεγονέτω ὡς ὁ ΓΑ
πρὸς τὸν ΑΒ, οὕτως τὸ ἀπὸ τῆς ΕΖ πρὸς τὸ ἀπὸ τῆς ΖΗ.
ὁ δὲ ΓΑ πρὸς τὸν ΑΒ λόγον οὐκ ἔχει, ὃν τετράγωνος ἀριθμὸς
πρὸς τετράγωνον ἀριθμόν; οὑδὲ τὸ ἀπὸ τῆς ΕΖ ἄρα πρὸς
τὸ ἀπὸ τῆς ΖΗ λόγον ἔχει, ὃν τετράγωνος ἀριθμὸς πρὸς
τετράγωνον ἀριθμόν. αἱ ΕΖ, ΖΗ ἄρα ῥηταί εἰσι δυνάμει
μόνον σύμμετροι; ἐκ δύο ἄρα ὀνομάτων ἐστὶν ἡ ΕΗ. λέγω δή, ὅτι καὶ πέμπτη.}

{\greekfont Ἐπεὶ γάρ ἐστιν ὡς ὁ ΓΑ πρὸς τὸν ΑΒ, οὕτως τὸ ἀπὸ τῆς ΕΖ
πρὸς τὸ ἀπὸ τῆς ΖΗ, ἀνάπαλιν ὡς ὁ ΒΑ πρὸς τὸν ΑΓ, οὕτως
τὸ ἀπὸ τῆς ΖΗ πρὸς τὸ ἀπὸ τῆς ΖΕ; μεῖζον ἄρα τὸ ἀπὸ
τῆς ΗΖ τοῦ ἀπὸ τῆς ΖΕ. ἔστω οὖν τῷ ἀπὸ τῆς
ΗΖ ἴσα τὰ ἀπὸ τῶν ΕΖ, Θ; ἀναστρέψαντι ἄρα ἐστὶν ὡς ὁ ΑΒ
ἀριθμὸς πρὸς τὸν ΒΓ, οὕτως τὸ ἀπὸ τῆς ΗΖ πρὸς τὸ
ἀπὸ τῆς Θ. ὁ δὲ ΑΒ πρὸς τὸν ΒΓ λόγον οὐκ ἔχει, ὃν
τετράγωνος ἀριθμὸς πρὸς τετράγωνον ἀριθμόν; οὐδ ἄρα τὸ
ἀπὸ τῆς ΖΗ πρὸς τὸ ἀπὸ τῆς Θ λόγον
ἔχει, ὃν τετράγωνος ἀριθμὸς πρὸς τετράγωνον ἀριθμόν.
ἀσύμμετρος ἄρα ἐστὶν ἡ ΖΗ τῇ Θ
μήκει; ὥστε ἡ ΖΗ τῆς ΖΕ μεῖζον δύναται τῷ ἀπὸ ἀσυμμέτρου
ἑαυτῇ. καί εἰσιν αἱ ΗΖ, ΖΕ ῥηταὶ δυνάμει μόνον σύμμετροι,
καὶ τὸ ΕΖ ἔλαττον ὄνομα σύμμετρόν ἐστι τῇ ἐκκειμένῃ
ῥητῇ τῇ Δ μήκει.}

{\greekfont Ἡ ΕΗ ἄρα ἐκ δύο ὀνομάτων ἐστὶ πέμπτη; ὅπερ
ἔδει δεῖχαι.}}

\ParallelRText{
\begin{center}
{\large Proposition 52}
\end{center}

To find a fifth binomial straight-line.

\epsfysize=1.3in
\centerline{\epsffile{Book10/fig052e.eps}}

Let the two numbers $AC$ and $CB$ be laid down such that $AB$ does
not have to either of them the ratio which (some) square number (has)
to (some) square number [Prop.~10.38~lem.]. And
let some rational straight-line $D$ be laid down. And let $EF$ be
commensurable [in length] with $D$. Thus, $EF$ (is) a rational (straight-line).
And let it be contrived that as $CA$ (is) to $AB$, so the (square) on
$EF$ (is) to the (square) on $FG$ [Prop.~10.6~corr.].
 And $CA$ does not have to $AB$ the ratio which (some)
 square number (has) to (some) square number. Thus, the (square) on
 $EF$ does not have to the (square) on $FG$ the ratio which (some)
 square number (has) to (some) square number either. Thus, $EF$ and
 $FG$ are rational (straight-lines which are) commensurable in square only
 [Prop.~10.9].  Thus, $EG$ is a binomial (straight-line) [Prop.~10.36]. So, I say that (it is) also
 a fifth (binomial straight-line).
 
 For since as $CA$ is to $AB$, so the (square) on $EF$ (is) to the (square)
 on $FG$, inversely, as $BA$ (is) to $AC$, so the (square) on $FG$ (is)
 to the (square) on $FE$ [Prop.~5.7~corr.]. Thus,
 the (square) on $GF$ (is) greater than the (square) on $FE$ [Prop.~5.14]. Therefore, let (the sum of) the (squares)
 on $EF$ and $H$ be equal to the (square) on $GF$. Thus, via conversion, 
 as the number $AB$ is to $BC$, so the (square) on $GF$ (is) to
 the (square) on $H$ [Prop.~5.19~corr.]. 
 And $AB$ does not have to $BC$ the ratio which (some) square number
 (has) to (some) square number. Thus, the (square) on $FG$ does not have
 to the (square) on $H$ the ratio which (some) square number (has) to
 (some) square number either. Thus, $FG$ is incommensurable in length
 with $H$ [Prop.~10.9]. Hence, the square on $FG$
 is greater than (the square on) $FE$ by the (square) on (some straight-line)
 incommensurable (in length) with ($FG$). And $GF$ and $FE$ are rational (straight-lines
 which are) commensurable in square only. And the lesser term $EF$ is commensurable in length with the  rational (straight-line previously)  laid down, $D$.
 
 Thus, $EG$ is a fifth binomial (straight-line).$^\dag$  (Which is) the very thing it
 was required to show.}
\end{Parallel}
{\footnotesize\noindent $^\dag$ If the rational straight-line has unit length then the length of a fifth binomial straight-line
is  $k\,(\sqrt{1+k'}+1)$. This, and the fifth apotome,
whose length is $k\,(\sqrt{1+k'}-1)$ [Prop.~10.89],
are the roots of $x^2- 2\,k\,\sqrt{1+k'}\,x+k^2\,k'=0$.}
 
%%%%
%10.53
%%%%
\pdfbookmark[1]{Proposition 10.53}{pdf10.53}
\begin{Parallel}{}{}
\ParallelLText{
\begin{center}
{\large \ggn{53}.}
\end{center}\vspace*{-7pt}

{\greekfont Εὑρεῖν τὴν ἐκ δύο ὀνομάτων ἕκτην.}

\epsfysize=1.in
\centerline{\epsffile{Book10/fig053g.eps}}

{\greekfont Ἐκκείσθωσαν δύο ἀριθμοὶ οἱ ΑΓ, ΓΒ, ὥστε τὸν ΑΒ πρὸς ἑκάτερον
α ὐτῶν λόγον μὴ ἔχειν, ὃν τετράγωνος ἀριθμὸς πρὸς τετράγωνον
ἀριθμόν; ἔστω δὲ καὶ ἕτερος ἀριθμὸς ὁ Δ μὴ τετράγωνος ὢν
μ ηδὲ πρὸς ἑκάτερον τῶν ΒΑ, ΑΓ λόγον ἔθων, ὃν τετράγωνος
ἀριθμὸς πρὸς τετράγωνον ἀριθμόν; καὶ ἐκκείσθω τις
ῥητὴ εὐθεῖα ἡ Ε, καὶ γεγονέτω ὡς ὁ Δ πρὸς τὸν ΑΒ, οὕτως
τὸ ἀπὸ τῆς Ε πρὸς τὸ ἀπὸ τῆς ΖΗ; σύμμετρον ἄρα τὸ ἀπὸ
τῆς Ε τῷ ἀπὸ τῆς ΖΗ. καί ἐστι ῥητὴ ἡ Ε; ῥητὴ ἄρα καὶ ἡ ΖΗ. καὶ ἐπεὶ οὐκ ἔχει ὁ Δ πρὸς τὸν ΑΒ λόγον, ὃν τετράγωνος
ἀριθμὸς πρὸς τετράγωνον ἀριθμόν, οὐδὲ τὸ ἀπὸ τῆς Ε ἄρα πρὸς τὸ ἀπὸ τῆς ΖΗ λόγον ἔχει, ὃν τετράγωνος ἀριθμὸς πρὸς τετράγωνον ἀριθμόν; ἀσύμμετρος ἄρα ἡ Ε τῇ ΖΗ μήκει. γεγονέτω
δὴ πάλιν ὡς ὁ ΒΑ πρὸς τὸν ΑΓ, οὕτως τὸ ἀπὸ τῆς ΖΗ
πρὸς τὸ ἀπὸ τῆς ΗΘ. σύμμετρον ἄρα τὸ ἀπὸ τῆς ΖΗ τῷ
ἀπὸ τῆς ΘΗ. ῥητὸν ἄρα τὸ ἀπὸ τῆς ΘΗ; ῥητὴ ἄρα ἡ ΘΗ. καὶ
ἐπεὶ ὁ ΒΑ πρὸς τὸν ΑΓ λόγον οὐκ ἔχει, ὃν τετράγωνος
ἀριθμὸς πρὸς τετράγωνον ἀριθμόν, οὐδὲ τὸ ἀπὸ τῆς ΖΗ πρὸς
τὸ ἀπὸ τῆς ΗΘ λόγον ἔχει, ὃν τετράγωνος ἀριθμὸς
πρὸς τετράγωνον ἀριθμόν; ἀσύμμετρος ἄρα ἐστὶν ἡ ΖΗ
τῇ ΗΘ μήκει. αἱ ΖΗ, ΗΘ ἄρα ῥηταί εἰσι
δυνάμει μόνον σύμμετροι; ἐκ δύο ἄρα ὀνομάτων ἐστὶν ἡ ΖΘ. δεικτέον δή, ὅτι καὶ ἕκτη.}

{\greekfont Ἐπεὶ γάρ ἐστιν ὡς ὁ Δ πρὸς τὸν ΑΒ, οὕτως τὸ ἀπὸ τῆς Ε πρὸς
τὸ ἀπὸ τῆς ΖΗ, ἔστι δὲ καὶ ὡς ὁ ΒΑ πρὸς τὸν ΑΓ, οὕτως
τὸ ἀπὸ τῆς ΖΗ πρὸς τὸ ἀπὸ τῆς ΗΘ, δι ἴσου ἄρα ἐστὶν
ὡς ὁ Δ πρὸς τὸν ΑΓ, οὕτως τὸ ἀπὸ τῆς Ε πρὸς τὸ ἀπὸ
τῆς ΗΘ. ὁ δὲ Δ πρὸς τὸν ΑΓ λόγον οὐκ ἔχει, ὃν τετράγωνος
ἀριθμὸς πρὸς τετράγωνον ἀριθμόν; οὐδὲ τὸ ἀπὸ τῆς
Ε ἄρα πρὸς τὸ ἀπὸ τῆς ΗΘ λόγον ἔχει, ὃν τετράγωνος
ἀριθμὸς πρὸς τετράγωνον ἀριθμόν;
ἀσύμμετρος ἄρα ἐστὶν
ἡ Ε τῇ ΗΘ μήκει. ἐδείθθη δὲ καὶ τῇ ΖΗ ἀσύμμετρος;
ἑκατέρα ἄρα τῶν ΖΗ, ΗΘ ἀσύμμετρός ἐστι τῇ Ε μήκει.
καὶ ἐπεί ἐστιν ὡς ὁ ΒΑ πρὸς τὸν ΑΓ, οὕτως τὸ ἀπὸ
τῆς ΖΗ πρὸς τὸ ἀπὸ τῆς ΗΘ, μεῖζον ἄρα τὸ ἀπὸ τῆς ΖΗ
τοῦ ἀπὸ τῆς ΗΘ. ἔστω οὖν τῷ ἀπὸ [τῆς] ΖΗ ἴσα
τὰ ἀπὸ τῶν ΗΘ, Κ; ἀναστρέψαντι ἄρα ὡς ὁ ΑΒ πρὸς ΒΓ, οὕτως
τὸ ἀπὸ ΖΗ πρὸς τὸ ἀπὸ τῆς Κ. ὁ δὲ ΑΒ πρὸς τὸν ΒΓ λόγον
οὐκ ἔχει, ὃν τετράγωνος ἀριθμὸς πρὸς τετράγωνον ἀριθμόν; ὥστε
οὐδὲ τὸ ἀπὸ ΖΗ πρὸς τὸ ἀπὸ τῆς Κ λόγον ἔχει, ὃν
τετράγωνος ἀριθμὸς πρὸς τετράγωνον ἀριθμόν. ἀσύμμετρος ἄρα
ἐστὶν ἡ ΖΗ τῇ Κ μήκει; ἡ ΖΗ ἄρα τῆς ΗΘ μεῖζον δύναται
τῷ ἀπὸ ἀσυμμέτρου ἑαυτῇ. καί εἰσιν αἱ ΖΗ, ΗΘ ῥηταὶ δυνάμει
μόνον σύμμετροι, καὶ οὐδετέρα α ὐτῶν σύμμετρός ἐστι μήκει
τῇ ἐκκειμένῃ ῥητῃ τῇ Ε.}

{\greekfont Ἡ ΖΘ ἄρα ἐκ δύο ὀνομάτων ἐστὶν ἕκτη; ὅπερ ἔδει δεῖχαι.}}

\ParallelRText{
\begin{center}
{\large Proposition 53}
\end{center}

To find a sixth binomial (straight-line).

\epsfysize=1.in
\centerline{\epsffile{Book10/fig053e.eps}}

Let the two numbers $AC$ and $CB$ be laid down such that $AB$ does
not have to each of them the ratio which (some) square number (has) to (some) square number.  And let $D$  also be another number, which is not
square, and does not have to each of $BA$ and $AC$ the ratio which (some)
square number (has) to (some) square number either [Prop.~10.28~lem.~I]. And let some rational straight-line $E$ be laid down. And let it be contrived that as $D$ (is) to
$AB$, so the (square) on $E$ (is) to the (square) on $FG$ [Prop.~10.6~corr.].  Thus, the (square) on $E$
(is) commensurable with the (square) on $FG$ [Prop.~10.6]. And $E$ is rational. Thus, $FG$
(is) also rational.  And since $D$ does not have to $AB$ the ratio
which (some) square number (has) to (some) square number, the (square) on $E$ thus does not have to the (square) on $FG$ the ratio which (some) square number (has) to (some) square number either. Thus, $E$ (is) incommensurable in length with $FG$ [Prop.~10.9].
So, again, let it have be contrived that as $BA$ (is) to $AC$, so the (square)
on $FG$ (is) to the (square) on $GH$ [Prop.~10.6~corr.]. The (square) on $FG$ (is)
thus commensurable with the (square) on $HG$ [Prop.~10.6].
 The (square) on $HG$ (is) thus rational.  Thus, $HG$ (is) rational. And since $BA$ does not have to
$AC$ the ratio which (some) square number (has) to (some) square number,
the (square) on $FG$ does not have to the (square) on $GH$ the
ratio which (some) square number (has) to (some) square number either.
Thus, $FG$ is incommensurable in length with $GH$ [Prop.~10.9]. Thus, $FG$ and $GH$ are rational (straight-lines which are) commensurable in square only. Thus, $FH$
is a binomial (straight-line) [Prop.~10.36]. So,
we must show that (it is) also a sixth (binomial straight-line).

For since as $D$ is to $AB$, so the (square) on $E$ (is) to the (square) on
$FG$, and also as $BA$ is to $AC$, so the (square) on $FG$ (is) to
the (square) on $GH$, thus, via equality, as $D$ is to $AC$, so
the (square) on $E$ (is) to the (square) on $GH$ [Prop.~5.22]. And $D$ does not have to $AC$ the
ratio which (some) square number (has) to (some) square number.  Thus, 
the (square) on $E$ does not have to the (square) on $GH$ the ratio
which (some) square number (has) to (some) square number either. $E$
is thus incommensurable in length with $GH$ [Prop.~10.9]. And ($E$) was also shown (to be)
incommensurable (in length) with $FG$.  Thus,  $FG$ and $GH$ are
each incommensurable in length with $E$. And since as $BA$ is to $AC$, 
so the (square) on $FG$ (is) to the (square) on $GH$, the (square) on $FG$ (is) thus greater than the (square) on $GH$ [Prop.~5.14]. Therefore, let (the sum of) the
(squares) on $GH$ and $K$ be equal to the (square) on $FG$. Thus,
via conversion, as $AB$ (is) to $BC$,
 so the (square) on $FG$ (is) to the
(square) on $K$ [Prop.~5.19~corr.].  And $AB$ does not
have to $BC$ the ratio which (some) square number (has) to (some) square number. Hence, the (square) on $FG$ does not have to the (square) on
$K$ the ratio which (some) square number (has) to (some) square number either. Thus, $FG$ is incommensurable in length with $K$ [Prop.~10.9]. The square on $FG$ is thus
greater than (the square on) $GH$ by the (square) on (some straight-line which is) incommensurable (in length) with ($FG$). And  $FG$ and $GH$ are rational
(straight-lines which are) commensurable in square only, and neither of them
is commensurable in length with the  rational (straight-line) $E$ (previously)
laid down.

Thus, $FH$ is a  sixth binomial (straight-line) [Def. 10.10].$^\dag$ (Which is) the very thing it was required to show.}
\end{Parallel}
{\footnotesize\noindent$^\dag$ If the rational straight-line has unit length then the length of a sixth binomial straight-line
is  $\sqrt{k}+\sqrt{k'}$. This, and the sixth apotome,
whose length is $\sqrt{k}-\sqrt{k'}$ [Prop.~10.90],
are the roots of $x^2- 2\,\sqrt{k}\,x+(k-k')=0$.}

%%%%
%10.53lem
%%%%
\begin{Parallel}{}{}
\ParallelLText{

\begin{center}
{\large {\greekfont Λῆμμα}.}
\end{center}\vspace*{-7pt}

{\greekfont Ἔστω δύο τετράγωνα τὰ ΑΒ, ΒΓ καὶ κείσθωσαν ὥστε ἐπ εὐθείας  εἶναι τὴν ΔΒ τῇ ΒΕ; ἐπ εὐθείας ἄρα
ἐστὶ καὶ ἡ ΖΒ τῇ ΒΗ. καὶ συμπεπληρώσθω τὸ ΑΓ παραλληλόγραμμον;
λέγω, ὅτι τετράγωνόν ἐστι τὸ ΑΓ, καὶ ὅτι τῶν ΑΒ, ΒΓ μέσον ἀνάλογόν ἐστι τὸ ΔΗ, καὶ ἔτι τῶν ΑΓ, ΓΒ μέσον ἀνάλογόν ἐστι
τὸ ΔΓ.}

{\greekfont Ἐπεὶ γὰρ ἴση ἐστὶν ἡ μὲν ΔΒ τῇ ΒΖ, ἡ δὲ ΒΕ τῇ ΒΗ, ὅλη
ἄρα ἡ ΔΕ ὅλῃ τῇ ΖΗ ἐστιν ἴση. ἀλλ ἡ μὲν ΔΕ ἑκατέρᾳ
τῶν ΑΘ, ΚΓ ἐστιν ἴση, ἡ δὲ ΖΗ ἑκατέρᾳ τῶν ΑΚ, ΘΓ ἐστιν ἴση;
καὶ ἑκατέρα ἄρα τῶν ΑΘ, ΚΓ ἑκατέρᾳ τῶν ΑΚ, ΘΓ ἐστιν ἴση.
ἰσόπλευρον ἄρα ἐστὶ τὸ ΑΓ παραλληλόγραμμον; ἔστι δὲ καὶ
ὀρθογώνιον; τετράγωνον ἄρα ἐστὶ τὸ ΑΓ.}

{\greekfont Καὶ ἐπεὶ ἐστιν ὡς ἡ ΖΒ πρὸς τὴν ΒΗ, οὕτως ἡ
ΔΒ πρὸς τὴν ΒΕ, ἀλλ ὡς μὲν ἡ ΖΒ πρὸς τὴν ΒΗ, οὕτως τὸ ΑΒ
πρὸς τὸ ΔΗ, ὡς δὲ ἡ ΔΒ πρὸς τὴν ΒΕ, οὕτως τὸ ΔΗ πρὸς
τὸ ΒΓ, καὶ ὡς ἄρα τὸ ΑΒ πρὸς τὸ ΔΗ, οὕτως τὸ ΔΗ πρὸς τὸ
ΒΓ. τῶν ΑΒ, ΒΓ ἄρα μέσον ἀνάλογόν ἐστι τὸ ΔΗ.}

{\greekfont Λέγω δή, ὅτι καὶ τῶν ΑΓ, ΓΒ μέσον ἀνάλογόν [ἐστι] τὸ ΔΓ.}

\epsfysize=1.7in
\centerline{\epsffile{Book10/fig053ag.eps}}

{\greekfont Ἐπεὶ γάρ ἐστιν ὡς ἡ ΑΔ πρὸς τὴν ΔΚ, οὕτως ἡ ΚΗ πρὸς
τὴν ΗΓ; ἴση γάρ [ἐστιν] ἑκατέρα ἑκατέρᾳ; καὶ συνθέντι ὡς ἡ ΑΚ
πρὸς ΚΔ, οὕτως ἡ ΚΓ
πρὸς ΓΗ, ἀλλ ὡς μὲν ἡ ΑΚ πρὸς ΚΔ, οὕτως
τὸ ΑΓ πρὸς τὸ ΓΔ, ὡς δὲ ἡ ΚΓ πρὸς ΓΗ, οὕτως τὸ ΔΓ πρὸς
ΓΒ, καὶ ὡς ἄρα τὸ ΑΓ πρὸς ΔΓ, οὕτως τὸ ΔΓ πρὸς τὸ ΒΓ. τῶν ΑΓ, ΓΒ ἄρα μέσον ἀνάλογόν ἐστι τὸ ΔΓ; ἃ προέκειτο
δεῖχαι.}}

\ParallelRText{
\begin{center}
{\large Lemma}
\end{center}

Let $AB$ and $BC$ be two squares, and let them
be laid down such that $DB$ is straight-on to $BE$. $FB$ is, thus, also
straight-on to $BG$. And let the parallelogram $AC$ be completed.
I say that $AC$ is a square, and that $DG$ is  the mean proportional to $AB$ and
$BC$, and, moreover,  $DC$ is the mean proportional to $AC$ and $CB$.

For since $DB$ is equal to $BF$, and $BE$ to $BG$, the whole of $DE$
is thus equal to the whole of $FG$. But $DE$ is equal to each of $AH$
and $KC$, and $FG$ is equal to each of $AK$ and $HC$ [Prop.~1.34]. Thus, $AH$ and $KC$ are also 
equal to  $AK$ and $HC$, respectively. Thus, the parallelogram $AC$
is equilateral. And (it is) also right-angled. Thus, $AC$ is a square.

And since as $FB$ is to $BG$, so $DB$ (is) to $BE$, but as $FB$
(is) to $BG$, so $AB$ (is) to $DG$, and as $DB$ (is) to $BE$, so
$DG$ (is) to $BC$ [Prop.~6.1], thus also as $AB$ (is) to $DG$, so $DG$  (is) to $BC$ [Prop.~5.11]. Thus, $DG$ is the mean proportional
to $AB$ and $BC$.

So I also say that $DC$ [is] the mean proportional to $AC$ and $CB$.

\epsfysize=1.7in
\centerline{\epsffile{Book10/fig053ae.eps}}

For since as $AD$ is to $DK$, so $KG$ (is) to $GC$. For [they are]  respectively equal. And, via composition, as $AK$ (is) to $KD$, so
$KC$ (is) to $CG$ [Prop.~5.18]. But as
$AK$ (is) to $KD$, so $AC$ (is) to $CD$, and as $KC$ (is) to $CG$,
so $DC$ (is) to $CB$ [Prop.~6.1]. Thus, also, as $AC$ (is) to $DC$, so
$DC$ (is) to $BC$ [Prop.~5.11]. Thus, $DC$ is the mean proportional to $AC$ and
$CB$. Which (is the very thing) it was prescribed to show.}
\end{Parallel}

%%%%
%10.54
%%%%
\pdfbookmark[1]{Proposition 10.54}{pdf10.54}
\begin{Parallel}{}{}
\ParallelLText{
\begin{center}
{\large \ggn{54}.}
\end{center}\vspace*{-7pt}

{\greekfont Ἐὰν θωρίον περιέθηται ὑπὸ ῥητῆς καὶ τῆς ἐκ δύο ὀνομάτων
πρώτης, ἡ τὸ θωρίον δυναμένη ἄλογός ἐστιν ἡ καλουμένη ἐκ δύο
ὀνομάτων.}

\epsfysize=1.3in
\centerline{\epsffile{Book10/fig054g.eps}}

{\greekfont Θωρίον γὰρ τὸ ΑΓ περιεθέσθω ὑπὸ ῥητῆς τῆς ΑΒ καὶ τῆς ἐκ δύο
ὀνομάτων πρώτης τῆς ΑΔ; λέγω, ὅτι ἡ τὸ ΑΓ θωρίον δυναμένη
ἄλογός ἐστιν ἡ καλουμένη ἐκ δύο ὀνομάτων.}

{\greekfont Ἐπεὶ γὰρ ἐκ δύο ὀνομάτων ἐστὶ πρώτη ἡ ΑΔ, διῃρήσθω
εἰς τὰ ὀνόματα κατὰ τὸ Ε, καὶ ἔστω τὸ μεῖζον ὄνομα τὸ
ΑΕ. φανερὸν δή, ὅτι αἱ ΑΕ, ΕΔ ῥηταί εἰσι δυνάμει μόνον σύμμετροι,
καὶ ἡ ΑΕ τῆς ΕΔ μεῖζον δύναται τῷ ἀπὸ συμμέτρου ἑαυτῃ, καὶ
ἡ ΑΕ σύμμετρός ἐστι τῇ ἐκκειμένῃ ῥητῇ τῇ ΑΒ μήκει.
τετμήσθω δὴ ἡ ΕΔ δίθα κατὰ τὸ Ζ σημεῖον. καὶ ἐπεὶ ἡ ΑΕ
τῆς ΕΔ μεῖζον δύναται τῷ ἀπὸ συμμέτρου ἑαυτῇ, ἐὰν ἄρα
τῷ τετάρτῳ μέρει τοῦ ἀπὸ τῆς ἐλάσσονος, τουτέστι τῷ ἀπὸ
τῆς ΕΖ, ἴσον παρὰ τὴν μείζονα τὴν ΑΕ παραβληθῇ ἐλλεῖπον
εἴδει τετραγώνῳ, εἰς σύμμετρα α ὐτὴν διαιρεῖ. παραβεβλήσθω οῦ̓ν
παρὰ τὴν ΑΕ τῷ ἀπὸ τῆς ΕΖ ἴσον τὸ ὑπὸ ΑΗ, ΗΕ; σύμμετρος
ἄρα ἐστὶν ἡ ΑΗ τῇ ΕΗ μήκει. καὶ ἤθθωσαν ἀπὸ τῶν Η, Ε, Ζ
ὁποτέρᾳ τῶν ΑΒ, ΓΔ παράλληλοι αἱ ΗΘ, ΕΚ, ΖΛ; καὶ τῷ μὲν
ΑΘ παραλληλογράμμῳ ἴσον τετράγωνον συνεστάτω τὸ ΣΝ, τῷ δὲ
ΗΚ ἴσον τὸ ΝΠ, καὶ κείσθω ὥστε ἐπ εὐθείας εἶναι τὴν
ΜΝ τῇ ΝΧ; ἐπ εὐθείας ἄρα ἐστὶ καὶ ἡ ΡΝ τῇ ΝΟ. καὶ
συμπεπληρώσθω τὸ ΣΠ παραλληλόγραμμον; τετράγωνον ἄρα ἐστὶ
τὸ ΣΠ. καὶ ἐπεὶ τὸ ὑπὸ τῶν ΑΗ, ΗΕ ἴσον ἐστὶ τῷ ἀπὸ
τῆς ΕΖ, ἔστιν ἄρα ὡς ἡ ΑΗ πρὸς ΕΖ, οὕτως ἡ ΖΕ πρὸς ΕΗ; καὶ
ὡς ἄρα τὸ ΑΘ πρὸς ΕΛ, τὸ ΕΛ πρὸς ΚΗ; τῶν ΑΘ, ΗΚ ἄρα μέσον
ἀνάλογόν ἐστι τὸ ΕΛ. ἀλλὰ τὸ μὲν ΑΘ ἴσον ἐστὶ τῷ ΣΝ, τὸ δὲ
ΗΚ ἴσον τῷ ΝΠ; τῶν ΣΝ, ΝΠ ἄρα μέσον ἀνάλογόν ἐστι τὸ ΕΛ.
ἔστι δὲ τῶν α ὐτῶν τῶν ΣΝ, ΝΠ μέσον ἀνάλογον καὶ τὸ ΜΡ; ἴσον
ἄρα ἐστὶ τὸ ΕΛ τῷ ΜΡ; ὥστε καὶ τῷ ΟΧ ἴσον ἐστίν. ἔστι
δὲ καὶ τὰ ΑΘ, ΗΚ τοῖς ΣΝ, ΝΠ ἴσα; ὅλον ἄρα τὸ ΑΓ ἴσον ἐστὶν
ὅλῳ τῷ ΣΠ, τουτέστι τῷ ἀπὸ τῆς ΜΧ τετραγώνῳ; τὸ ΑΓ
ἄρα δύναται ἡ ΜΧ. λέγω, ὅτι ἡ ΜΧ ἐκ δύο ὀνομάτων ἐστίν.}

{\greekfont Ἐπεὶ γὰρ σύμμετρός ἐστιν ἡ ΑΗ τῇ ΗΕ, σύμμετρός ἐστι
καὶ ἡ ΑΕ ἑκατέρᾳ τῶν ΑΗ, ΗΕ. ὑπόκειται δὲ καὶ ἡ ΑΕ  τῇ ΑΒ σύμμετρος; καὶ αἱ
ΑΗ, ΗΕ ἄρα τῇ ΑΒ σύμμετροί εἰσιν. καί ἐστι ῥητὴ ἡ ΑΒ; ῥητὴ
ἄρα ἐστὶ καὶ ἑκατέρα τῶν ΑΗ, ΗΕ; ῥητὸν ἄρα ἐστὶν ἑκάτερον
τῶν ΑΘ, ΗΚ, καί ἐστι σύμμετρον τὸ ΑΘ τῷ ΗΚ. ἀλλὰ τὸ μὲν ΑΘ
τῷ ΣΝ ἴσον ἐστίν, τὸ δὲ ΗΚ τῷ ΝΠ; καὶ
τὰ ΣΝ, ΝΠ ἄρα,
τουτέστι τὰ ἀπὸ τῶν ΜΝ, ΝΧ, ῥητά ἐστι καὶ σύμμετρα. καὶ
ἐπεὶ ἀσύμμετρός ἐστιν ἡ ΑΕ τῇ ΕΔ μήκει, ἀλλ
ἡ μὲν ΑΕ τῇ ΑΗ ἐστι σύμμετρος, ἡ δὲ ΔΕ τῇ ΕΖ σύμμετρος,
ἀσύμμετρος ἄρα καὶ ἡ ΑΗ τῇ ΕΖ;
ὥστε καὶ τὸ ΑΘ τῷ ΕΛ ἀσύμμετρόν ἐστιν. ἀλλὰ τὸ μὲν
ΑΘ τῷ ΣΝ ἐστιν ἴσον, τὸ δὲ ΕΛ τῷ ΜΡ; καὶ τὸ ΣΝ ἄρα
τῷ ΜΡ ἀσύμμετρόν ἐστιν. ἀλλ ὡς τὸ ΣΝ πρὸς ΜΡ, ἡ ΟΝ
πρὸς τὴν ΝΡ; ἀσύμμετρος ἄρα ἐστὶν ἡ ΟΝ τῇ ΝΡ.
ἴση δὲ ἡ μὲν ΟΝ τῇ ΜΝ, ἡ δὲ ΝΡ
τῇ ΝΧ; ἀσύμμετρος ἄρα ἐστὶν ἡ ΜΝ τῇ ΝΧ. καί
ἐστι τὸ ἀπὸ τῆς ΜΝ σύμμετρον τῷ ἀπὸ τῆς ΝΧ, καὶ
ῥητὸν ἑκάτερον; αἱ ΜΝ, ΝΧ ἄρα ῥηταί εἰσι
δυνάμει μόνον σύμμετροι.}

{\greekfont Ἡ ΜΧ ἄρα ἐκ δύο ὀνομάτων ἐστὶ καὶ δύναται
τὸ ΑΓ; ὅπερ ἔδει δεῖχαι.}}

\ParallelRText{
\begin{center}
{\large Proposition 54}
\end{center}

If an area is contained by a rational (straight-line)
and a first binomial (straight-line) then  the  square-root  of the area is
the irrational (straight-line which is) called binomial.$^\dag$

\epsfysize=1.3in
\centerline{\epsffile{Book10/fig054e.eps}}

For let the area $AC$ be contained by the rational (straight-line) $AB$
and by the first binomial (straight-line) $AD$. I say that square-root of 
area $AC$ is the irrational (straight-line which is) called binomial.

For since $AD$ is a first binomial (straight-line), let it be divided
into its (component) terms at $E$, and let $AE$ be the greater term.
So, (it is) clear that $AE$ and $ED$ are rational (straight-lines which are)
commensurable in square only, and that the square on $AE$ is greater
than (the square on) $ED$ by the (square) on (some straight-line)
commensurable (in length) with ($AE$), and that $AE$ is commensurable 
(in length) with
the rational (straight-line) $AB$ (first) laid out [Def.~10.5]. So, let $ED$ be cut in half
at  point $F$. And since the square on $AE$ is greater than (the square on)
$ED$ by the (square) on (some straight-line) commensurable (in length) with ($AE$),
thus if a (rectangle) equal to the fourth part of the (square) on the lesser (term)---that is to say, the (square) on $EF$---falling short by a square figure,
is applied to the greater (term) $AE$, then it divides it into (terms which are)
commensurable (in length) [Prop~10.17]. 
Therefore, let the (rectangle contained) by $AG$ and $GE$, equal to
the (square) on $EF$, be applied to $AE$. $AG$ 
is thus commensurable in length with $EG$. And let $GH$, $EK$, and $FL$
be drawn from (points) $G$, $E$, and $F$ (respectively), parallel to
either of $AB$ or $CD$. And let the square $SN$, equal to the
parallelogram $AH$, be constructed, and (the square) $NQ$,
equal to (the parallelogram) $GK$ [Prop.~2.14].
And let $MN$ be laid down so as to be straight-on to $NO$. $RN$
is thus also straight-on to $NP$. And let the parallelogram $SQ$ 
be completed. $SQ$ is thus a square [Prop.~10.53~lem.].  And since the (rectangle
contained) by $AG$ and $GE$ is equal to the (square) on $EF$, thus as
$AG$ is to $EF$, so $FE$ (is) to $EG$ [Prop.~6.17]. And thus as $AH$ (is) to
$EL$, (so) $EL$ (is) to $KG$ [Prop.~6.1]. Thus,
$EL$ is the mean proportional to $AH$ and $GK$. But, $AH$ is equal to
$SN$, and $GK$ (is) equal to $NQ$. $EL$ is thus the mean proportional
to $SN$ and $NQ$. And $MR$ is also the mean proportional to the
same---(namely), $SN$ and $NQ$ [Prop.~10.53~lem.]. $EL$ is thus equal to $MR$.
Hence, it is also equal to $PO$ [Prop.~1.43]. And
$AH$ plus $GK$ is equal to $SN$ plus $NQ$. Thus, the whole of $AC$
is equal to the whole of $SQ$---that is to say, to the square on $MO$.
Thus, $MO$ (is) the square-root of (area) $AC$.  I say that $MO$ is a
binomial (straight-line).

For since $AG$ is commensurable (in length) with $GE$, $AE$ is also commensurable (in length)
with each of $AG$ and $GE$ [Prop.~10.15]. 
And $AE$ was also assumed (to be) commensurable (in length) with $AB$.  Thus,
$AG$ and $GE$ are also commensurable (in length) with $AB$
[Prop.~10.12]. And $AB$ is rational. $AG$
and $GE$ are thus each also  rational. Thus, $AH$ and $GK$ are each
rational (areas), and $AH$ is commensurable with $GK$ [Prop.~10.19]. But, $AH$ is equal to $SN$,
and $GK$ to $NQ$. $SN$ and $NQ$---that is to say, the (squares)
on $MN$ and $NO$ (respectively)---are thus also rational and commensurable. And since $AE$ is incommensurable in length with $ED$,
but $AE$ is commensurable (in length) with $AG$, and $DE$ (is) commensurable (in length)
with $EF$, $AG$ (is) thus also incommensurable (in length) with $EF$
[Prop.~10.13]. Hence, $AH$ is also incommensurable with $EL$ [Props.~6.1, 10.11]. But, $AH$ is equal to $SN$, and $EL$
to $MR$. Thus, $SN$ is also incommensurable with $MR$. But, as
$SN$ (is) to $MR$, (so) $PN$ (is) to $NR$ [Prop.~6.1].  $PN$ is thus incommensurable (in length) with $NR$ [Prop.~10.11]. And $PN$ (is) equal to $MN$, and $NR$ to
$NO$. Thus, $MN$ is incommensurable (in length) with $NO$. And the (square)
on $MN$ is commensurable with the (square) on $NO$, and each (is)
rational.  $MN$ and $NO$ are thus rational (straight-lines which are)
commensurable in square only.

Thus, $MO$ is (both) a binomial (straight-line) [Prop. 10.36], and the square-root of $AC$.
(Which is) the very thing it was required to show.}
\end{Parallel}
{\footnotesize\noindent $^\dag$ If the rational straight-line has unit length then this proposition states that the square-root of 
a first binomial straight-line is a binomial straight-line: {i.e.}, 
a first binomial straight-line has a length $k+k\,\sqrt{1-k'^{\,2}}$ whose
square-root can be written $\rho\,(1+\sqrt{k''})$, where $\rho=\sqrt{k\,(1+k')/2}$ and $k''=(1-k')/(1+k')$. This is the length of a binomial straight-line (see Prop.~10.36), since $\rho$ is rational.}

%%%%
%10.55
%%%%
\pdfbookmark[1]{Proposition 10.55}{pdf10.55}
\begin{Parallel}{}{}
\ParallelLText{
\begin{center}
{\large \ggn{55}.}
\end{center}\vspace*{-7pt}

{\greekfont Ἐὰν θωρίον περιέθηται ὑπὸ ῥητῆς καὶ τῆς ἐκ δύο ὁνομάτων
δευτέρας, ἡ τὸ θωρίον δυναμένη ἄλογός ἐστιν ἡ καλουμένη
ἐκ δύο μέσων πρώτη.}

\epsfysize=1.3in
\centerline{\epsffile{Book10/fig054g.eps}}

{\greekfont Περιεθέσθω γὰρ θωρίον τὸ ΑΒΓΔ ὑπὸ ῥητῆς τῆς ΑΒ  καὶ τῆς
ἐκ δύο ὀνομάτων δυετέρας τῆς ΑΔ; λέγω, ὅτι ἡ τὸ ΑΓ θωρίον
δυναμένη ἐκ δύο μέσων πρώτη ἐστίν.}

{\greekfont Ἐπεὶ γὰρ ἐκ δύο ὀνομάτων δευτέρα ἐστὶν ἡ ΑΔ, διῃρήσθω εἰς
τὰ ὀνόματα κατὰ τὸ Ε, ὥστε τὸ μεῖζον ὄνομα εἶναι τὸ ΑΕ;
αἱ ΑΕ, ΕΔ ἄρα ῥηταί εἰσι δυνάμει μόνον σύμμετροι, καὶ ἡ ΑΕ
τῆς ΕΔ μεῖζον δύναται τῷ ἀπὸ συμμέτρου ἑαυτῇ, καὶ τὸ
ἔλαττον ὄνομα ἡ ΕΔ σύμμετρόν ἐστι τῇ ΑΒ μήκει. τετμήσθω
ἡ ΕΔ δίθα κατὰ τὸ Ζ, καὶ τῷ ἀπὸ τῆς ΕΖ ἴσον παρὰ τὴν ΑΕ
παραβεβλήσθω ἐλλεῖπον εἴδει τετραγώνῳ τὸ ὑπὸ τῶν
ΑΗΕ; σύμμετρος ἄρα ἡ ΑΗ τῇ ΗΕ μήκει. καὶ  διὰ τῶν Η, Ε, Ζ
παράλληλοι ἤθθωσαν ταῖς ΑΒ, ΓΔ αἱ ΗΘ, ΕΚ, ΖΛ, καὶ τῷ μὲν
ΑΘ παραλληλογράμμῳ ἴσον τετράγωνον συνεστάτω τὸ ΣΝ, τῷ δὲ ΗΚ
ἴσον τετράγωνον 
τὸ ΝΠ, καὶ κείσθω ὥστε ἐπ εὐθείας εἶναι τὴν ΜΝ τῇ ΝΧ; ἐπ
εὐθείας ἄρα [ἐστὶ] καὶ ἡ ΡΝ τῆ ΝΟ. καὶ συμπεπληρώσθω
τὸ ΣΠ τετράγωνον;  φανερὸν δὴ ἐκ τοῦ προδεδειγμένου, ὅτι τὸ ΜΡ
μέσον ἀνάλογόν ἐστι τῶν ΣΝ, ΝΠ, καὶ ἴσον τῷ ΕΛ, καὶ
ὅτι τὸ ΑΓ θωρίον δύναται ἡ ΜΧ. δεικτέον δή, ὅτι ἡ ΜΧ ἐκ
δύο μέσων ἐστὶ πρώτη.}

{\greekfont Ἐπεὶ ἀσύμμετρός ἐστιν ἡ ΑΕ τῇ ΕΔ
μήκει, σύμμετρος δὲ ἡ ΕΔ τῇ ΑΒ, ἀσύμμετρος ἄρα ἡ ΑΕ τῇ
ΑΒ. καὶ ἐπεὶ σύμμετρός ἐστιν ἡ ΑΗ τῇ ΕΗ, σύμμετρός ἐστι
καὶ ἡ ΑΕ ἑκατέρᾳ τῶν ΑΗ, ΗΕ. ἀλλὰ ἡ ΑΕ ἀσύμμετρος τῇ
ΑΒ μήκει; καὶ αἱ ΑΗ, ΗΕ ἄρα ἀσύμμετροί εἰσι τῇ ΑΒ. αἱ ΒΑ,
ΑΗ, ΗΕ ἄρα ῥηταί εἰσι δυνάμει μόνον σύμμετροι; ὥστε μέσον
ἐστὶν ἑκάτερον τῶν ΑΘ, ΗΚ. ὥστε καὶ ἑκάτερον τῶν ΣΝ, ΝΠ
μέσον ἐστίν. καὶ αἱ ΜΝ, ΝΧ ἄρα μέσαι εἰσίν. καὶ ἐπεὶ
σύμμετρος ἡ ΑΗ τῇ ΗΕ μήκει, σύμμετρόν ἐστι καὶ τὸ ΑΘ τῷ
ΗΚ, τουτέστι τὸ ΣΝ τῷ ΝΠ, τουτέστι τὸ ἀπὸ τῆς ΜΝ τῷ ἀπὸ
τῆς ΝΧ [ὥστε δυνάμει εἰσὶ σύμμετροι αἱ ΜΝ, ΝΧ]. καὶ
ἐπεὶ ἀσύμμετρός ἐστιν ἡ ΑΕ τῇ ΕΔ μήκει, ἀλλ ἡ μὲν ΑΕ
σύμμετρός ἐστι τῇ ΑΗ, ἡ δὲ ΕΔ τῇ ΕΖ σύμμετρος, ἀσύμμετρος ἄρα ἡ ΑΗ τῇ ΕΖ; ὥστε καὶ τὸ ΑΘ τῷ ΕΛ ἀσύμμετρόν ἐστιν,
τουτέστι τὸ ΣΝ τῷ ΜΡ, τουτέστιν ὁ ΟΝ τῇ ΝΡ, τουτέστιν
ἡ ΜΝ τῇ ΝΧ ἀσύμμετρός ἐστι μήκει.  ἐδείθθησαν δὲ αἱ ΜΝ,
ΝΧ καὶ μέσαι οὖσαι καὶ δυνάμει σύμμετροι; αἱ ΜΝ, ΝΧ ἄρα
μέσαι εἰσὶ δυνάμει μόνον σύμμετροι. λέγω δή, ὅτι καὶ ῥητὸν
περιέθουσιν. ἐπεὶ γὰρ ἡ ΔΕ ὑπόκειται ἑκατέρᾳ τῶν ΑΒ, ΕΖ σύμμετρος, σύμμετρος ἄρα καὶ ἡ ΕΖ τῇ ΕΚ. καὶ ῥητὴ ἑκατέρα
α ὐτῶν; ῥητὸν ἄρα τὸ ΕΛ, τουτέστι τὸ ΜΡ; τὸ δὲ ΜΡ ἐστι
τὸ ὑπὸ τῶν ΜΝΧ. ἐὰν δὲ δύο μέσαι δυνάμει μόνον σύμμετροι
συντεθῶσι ῥητὸν περιέθουσαι, ἡ ὅλη ἄλογός ἐστιν, καλεῖται δὲ ἐκ δύο μέσων πρώτη.}

{\greekfont Ἡ ἄρα ΜΧ ἐκ δύο μέσων ἐστὶ πρώτη; ὅπερ ἔδει δεῖχαι.}}

\ParallelRText{
\begin{center}
{\large Proposition 55}
\end{center}

If an area is contained by a rational (straight-line)
and a second binomial (straight-line) then the square-root of the area
is the irrational (straight-line which is) called first bimedial.$^\dag$

\epsfysize=1.3in
\centerline{\epsffile{Book10/fig054e.eps}}

For let the area $ABCD$ be contained by the rational (straight-line) $AB$
and by the second binomial (straight-line)  $AD$. I say that the square-root
of  area $AC$ is a first bimedial (straight-line).\

For since $AD$ is a second binomial (straight-line), let it be divided
into its (component) terms at $E$, such that $AE$ is the greater term. Thus,
$AE$ and $ED$ are rational (straight-lines which are) commensurable
in square only, and the square on $AE$ is greater than (the square on)
$ED$ by the (square) on (some straight-line) commensurable (in length) with ($AE$),
and the lesser term $ED$ is commensurable in length with $AB$ [Def.~10.6].  Let $ED$ be cut in half at $F$.
And let the (rectangle contained) by $AGE$, equal to the (square) on $EF$,
 be applied to $AE$, falling short by a square figure. $AG$
(is) thus commensurable in length with $GE$ [Prop.~10.17]. And let $GH$, $EK$, and
$FL$ be drawn through (points) $G$, $E$, and $F$ (respectively), parallel
to $AB$ and $CD$. And let the square $SN$, equal to the parallelogram
$AH$, be constructed, and the square $NQ$, equal to
$GK$.  And let $MN$ be laid down so as to be straight-on to $NO$. Thus,
$RN$ [is] also straight-on to $NP$. And let the square $SQ$ be
completed. So, (it is) clear from what has been previously
demonstrated [Prop.~10.53~lem.] that
$MR$ is the mean proportional to $SN$ and $NQ$, and (is) equal to
$EL$, and that $MO$ is the square-root of the area $AC$. So, we
must show that $MO$ is a first bimedial (straight-line).

Since $AE$
is incommensurable in length with $ED$, and $ED$ (is) commensurable 
 (in length) with $AB$, $AE$ (is) thus incommensurable (in length) with $AB$ [Prop.~10.13]. And since $AG$
is commensurable (in length) with $EG$, $AE$ is also commensurable
(in length) with each of $AG$ and $GE$ [Prop.~10.15]. 
But, $AE$ is incommensurable in length with $AB$.  Thus, $AG$ and $GE$
are also (both) 
incommensurable (in length) with $AB$ [Prop.~10.13]. 
Thus, $BA$, $AG$,  and ($BA$, and) $GE$ are (pairs of) rational (straight-lines which are)
commensurable in square only. And, hence, each of $AH$ and $GK$ is a medial
(area) [Prop.~10.21]. Hence, each of $SN$ and $NQ$
is also a medial (area). Thus, $MN$ and $NO$ are medial (straight-lines). 
And since $AG$ (is) commensurable in length with $GE$, $AH$
is also commensurable with $GK$---that is to say, $SN$ with $NQ$---that is to say, the (square)
on $MN$  with the (square) on $NO$ [hence, $MN$ and $NO$
are commensurable in square] [Props.~6.1, 10.11]. And since $AE$ is incommensurable
in length with $ED$, but $AE$ is commensurable (in length) with $AG$,
and $ED$ commensurable (in length) with $EF$, $AG$ (is) thus
incommensurable (in length) with $EF$ [Prop.~10.13]. Hence, $AH$ is also incommensurable with $EL$---that is to say, $SN$ with $MR$---that is to
say, $PN$ with $NR$---that is to say, $MN$ is incommensurable
in length with $NO$  [Props.~6.1, 10.11]. But $MN$ and $NO$ have also
been shown to be medial (straight-lines) which are commensurable in square. 
Thus, $MN$ and $NO$ are medial (straight-lines which are) commensurable
in square only.
So,
I say that they also contain a rational (area). For since $DE$ was assumed
(to be) commensurable (in length) with each of $AB$ and $EF$, $EF$
(is) thus also commensurable with $EK$ [Prop.~10.12]. And they (are) each rational.
Thus, $EL$---that is to say, $MR$---(is) rational [Prop.~10.19]. And $MR$ is the (rectangle contained) by $MNO$. And if two medial (straight-lines), commensurable
in square only, which contain a rational (area), are added together, then the
whole is (that) irrational (straight-line which is)
called first bimedial [Prop.~10.37].

Thus, $MO$ is a first bimedial (straight-line). (Which is) the very thing it
was required to show.}
\end{Parallel}
{\footnotesize\noindent$^\dag$ If the rational straight-line has unit length then this proposition states that the square-root of 
a second binomial straight-line is a first bimedial straight-line: {i.e.}, 
a second binomial straight-line has a length $k/\sqrt{1-k'^{\,2}}+k$ whose
square-root can be written $\rho\,(k''^{1/4}+k''^{3/4})$, where $\rho=\sqrt{(k/2)\,(1+k')/(1-k')}$ and $k''=(1-k')/(1+k')$. This is the length of a first bimedial straight-line (see Prop.~10.37), since $\rho$ is rational.}

%%%%
%10.56
%%%%
\pdfbookmark[1]{Proposition 10.56}{pdf10.56}
\begin{Parallel}{}{}
\ParallelLText{
\begin{center}
{\large \ggn{56}.}
\end{center}\vspace*{-7pt}

{\greekfont Ἐὰν θωρίον περιέθηται ὑπὸ ῥητῆς καὶ τῆς ἐκ
δύο ὀνομάτων τρίτης, ἡ τὸ θωρίον δυναμένη ἄλογός ἐστιν
ἡ καλουμένη ἐκ δύο μέσων δευτέρα.}

{\greekfont Θωρίον γὰρ τὸ ΑΒΓΔ περιεθέσθω ὑπὸ ῥητῆς τῆς ΑΒ καὶ τῆς
ἐκ δύο ὀνομάτων τρίτης τῆς ΑΔ διῃρημένης εἰς τὰ ὀνόματα
κατὰ τὸ Ε, ὧν μεῖζόν ἐστι τὸ ΑΕ; λέγω,  ὅτι ἡ τὸ ΑΓ
θωρίον δυναμένη ἄλογός ἐστιν ἡ καλουμένη ἐκ δύο μέσων
δευτέρα.}

\epsfysize=1.3in
\centerline{\epsffile{Book10/fig054g.eps}}

{\greekfont Κατεσκευάσθω γὰρ τὰ α ὐτὰ τοῖς πρότερον. καὶ ἐπεὶ ἐκ δύο
ὀνομάτων ἐστὶ τρίτη ἡ ΑΔ, αἱ ΑΕ, ΕΔ ἄρα ῥηταί εἰσι
δυνάμει μόνον σύμμετροι, καὶ ἡ ΑΕ τῆς ΕΔ μεῖζον δύναται
τῷ ἀπὸ συμμέτρου ἑαυτῇ, καὶ οὐδετέρα τῶν ΑΕ, ΕΔ
σύμμετρός [ἐστι] τῇ ΑΒ μήκει. ὁμοίως δὴ τοῖς προδεδειγμένοις
δείχομεν, ὅτι ἡ ΜΧ ἐστιν ἡ τὸ ΑΓ θωρίον δυναμένη, καὶ αἱ ΜΝ,
ΝΧ μέσαι εἰσὶ δυνάμει μόνον σύμμετροι; ὥστε ἡ ΜΧ ἐκ δύο
μέσων ἐστίν. δεικτέον δή, ὅτι καὶ δευτέρα.}

{\greekfont μ̀βοχ{[}Καὶ] ἐπεὶ ἀσύμμετρός ἐστιν ἡ ΔΕ τῇ ΑΒ μ ήκει, τουτέστι
τῇ ΕΚ, σύμμετρος δὲ ἡ ΔΕ τῇ ΕΖ, ἀσύμμετρος ἄρα ἐστὶν
ἡ ΕΖ τῇ ΕΚ μ ήκει. καί εἰσι ῥηταί; αἱ ΖΕ, ΕΚ ἄρα ῥηταί
εἰσι δυνάμει μόνον σύμμετροι. μέσον ἄρα [ἐστὶ] τὸ ΕΛ,
τουτέστι τὸ ΜΡ; καὶ περιέθεται ὑπὸ τῶν ΜΝΧ; μέσον ἄρα ἐστὶ
τὸ ὑπὸ τῶν ΜΝΧ.}

{\greekfont Ἡ ΜΧ ἄρα ἐκ δύο μέσων ἐστὶ δευτέρα; ὅπερ ἔδει δεῖχαι.}}

\ParallelRText{
\begin{center}
{\large Proposition 56}
\end{center}

If an area is contained by a rational (straight-line)
and a third binomial (straight-line) then the square-root of the area is
the irrational (straight-line which is) called second bimedial.$^\dag$

For let the area $ABCD$ be contained by the rational (straight-line)
$AB$ and by the third binomial (straight-line) $AD$, which has been divided into
its (component) terms at $E$, of which $AE$ is the greater.  I say that
the square-root of area $AC$ is the irrational (straight-line which is)
called second bimedial.

\epsfysize=1.3in
\centerline{\epsffile{Book10/fig054e.eps}}

For let the same construction be made as  previously. And since
$AD$ is a third binomial (straight-line), $AE$ and $ED$
are thus rational (straight-lines which are) commensurable in square only,
and the square on $AE$ is greater than (the square on) $ED$
by the (square) on (some straight-line) commensurable (in length) with
($AE$), and neither of $AE$ and $ED$ [is] commensurable in length
with $AB$ [Def.~10.7]. So, similarly
to that which has been previously demonstrated, we can show that
$MO$ is the square-root of  area $AC$, and $MN$ and $NO$
are medial (straight-lines which are) commensurable in square only.
Hence, $MO$ is bimedial. So, we must show that (it is)
also second (bimedial).

\mbox{[}And] since $DE$ is incommensurable in length with $AB$---that is to say,
with $EK$---and $DE$ (is) commensurable (in length) with $EF$,
$EF$ is thus incommensurable in length with $EK$ [Prop.~10.13]. And they are (both)
rational (straight-lines). Thus, $FE$ and $EK$ are rational (straight-lines which are) commensurable in square only. $EL$---that is to say, $MR$---[is] thus medial [Prop.~10.21]. And it
is contained by $MNO$.  Thus, the (rectangle contained) by $MNO$
is medial.

Thus, $MO$ is a second bimedial (straight-line) [Prop. 10.38]. (Which is) the very thing it was required to show.}
\end{Parallel}
{\footnotesize\noindent $^\dag$ If the rational straight-line has unit length then this proposition states that the square-root of 
a third binomial straight-line is a second bimedial straight-line: {i.e.}, 
a third binomial straight-line has a length $k^{1/2}\,(1+\sqrt{1-k'^{\,2}})$ whose
square-root can be written $\rho\,(k^{1/4}+k''^{1/2}/k^{1/4})$, where $\rho=\sqrt{(1+k')/2}$ and $k''=k\,(1-k')/(1+k')$. This is the length of a second bimedial straight-line (see Prop.~10.38), since $\rho$ is rational.}

%%%%
%10.57
%%%%
\pdfbookmark[1]{Proposition 10.57}{pdf10.57}
\begin{Parallel}{}{}
\ParallelLText{
\begin{center}
{\large \ggn{57}.}
\end{center}\vspace*{-7pt}

{\greekfont Ἐὰν θωρίον περιέθηται ὑπὸ ῥητῆς καὶ τῆς ἐκ δύο ὀνομάτων
τετάρτης, ἡ τὸ θωρίον δυναμένη ἄλογός ἐστιν ἡ καλουμένη 
μείζων.}

\epsfysize=1.3in
\centerline{\epsffile{Book10/fig054g.eps}}

{\greekfont Θωρίον γὰρ τὸ ΑΓ περιεθέσθω ὑπὸ ῥητῆς τῆς ΑΒ καὶ τῆς
ἐκ δύο ὀνομάτων τετάρτης τῆς ΑΔ διῃρημένης εἰς τὰ
ὀνόματα κατὰ τὸ Ε, ὧν μεῖζον ἔστω τὸ ΑΕ; λέγω, ὅτι
ἡ τὸ ΑΓ θωρίον δυναμένη ἄλογός ἐστιν ἡ καλουμένη
μείζων.}

{\greekfont Ἐπεὶ γὰρ ἡ ΑΔ ἐκ δύο ὀνομάτων ἐστὶ τετάρτη, αἱ ΑΕ, ΕΔ
ἄρα ῥηταί εἰσι δυνάμει μόνον σύμμετροι, καὶ ἡ ΑΕ τῆς
ΕΔ μεῖζον δύναται τῷ ἀπὸ ἀσυμμέτρου ἑαυτῇ, καὶ ἡ ΑΕ
τῇ ΑΒ σύμμετρός [ἐστι] μήκει. τετμήσθω ἡ ΔΕ δίθα κατὰ
τὸ Ζ, καὶ τῷ ἀπὸ τῆς ΕΖ ἴσον παρὰ τὴν ΑΕ παραβεβλήσθω
παραλληλόγραμμον τὸ ὑπὸ ΑΗ, ΗΕ; ἀσύμμετρος ἄρα ἐστὶν
ἡ ΑΗ τῇ ΗΕ μήκει. ἤθθωσαν παράλληλοι τῇ ΑΒ αἱ ΗΘ, ΕΚ, 
ΖΛ, καὶ τὰ λοιπὰ τὰ α ὐτὰ τοῖς πρὸ τούτου γεγονέτω;
φανερὸν δή, ὅτι ἡ τὸ ΑΓ θωρίον δυναμένη ἐστὶν ἡ
ΜΧ. δεικτέον δή, ὅτι ἡ ΜΧ ἄλογός ἐστιν ἡ καλουμένη
μείζων.}

{\greekfont Ἐπεὶ ἀσύμμετρός ἐστιν ἡ ΑΗ τῇ ΕΗ μήκει, ἀσύμμετρόν
ἐστι καὶ τὸ ΑΘ τῷ ΗΚ, τουτέστι τὸ ΣΝ τῷ ΝΠ; αἱ ΜΝ, ΝΧ ἄρα δυνάμει εἰσὶν ἀσύμμετροι. καὶ ἐπεὶ σύμμετρός ἐστιν ἡ ΑΕ
τῇ ΑΒ μήκει, ῥητόν ἐστι τὸ ΑΚ; καί ἐστιν ἴσον τοῖς
ἀπὸ τῶν ΜΝ, ΝΧ; ῥητὸν ἄρα [ἐστὶ] καὶ τὸ συγκείμενον
ἐκ τῶν ἀπὸ τῶν ΜΝ, ΝΧ. καὶ ἐπεὶ ἀσύμμετρός
[ἐστιν] ἡ ΔΕ τῇ ΑΒ μήκει, τουτέστι  τῇ ΕΚ, ἀλλὰ
ἡ ΔΕ σύμμετρός ἐστι τῇ ΕΖ, ἀσύμμετρος ἄρα ἡ ΕΖ
τῇ ΕΚ μήκει. αἱ ΕΚ, ΕΖ ἄρα ῥηταί εἰσι δυνάμει μόνον σύμμετροι;
μέσον ἄρα τὸ ΛΕ, τουτέστι τὸ ΜΡ. καὶ περιέθεται ὑπὸ τῶν
ΜΝ, ΝΧ; μέσον ἄρα ἐστὶ τὸ ὑπὸ τῶν ΜΝ, ΝΧ. καὶ ῥητὸν
τὸ [συγκείμενον] ἐκ τῶν ἀπὸ τῶν ΜΝ, ΝΧ, καί εἰσιν ἀσύμμετροι
αἱ ΜΝ, ΝΧ δυνάμει. ἐὰν δὲ δύο εὐθεῖαι δυνάμει ἀσύμμετροι
συντεθῶσι ποιοῦσαι τὸ μὲν συγκείμενον ἐκ τῶν ἀπ α ὐτῶν
τετραγώνων ῥητόν, τὸ δ ὑπ α ὐτῶν μέσον, ἡ ὅλη ἄλογός
ἐστιν, καλεῖται δὲ μείζων.}

{\greekfont Ἡ ΜΧ ἄρα ἄλογός ἐστιν ἡ καλουμένη μείζων, καὶ δύναται
τὸ ΑΓ θωρίον; ὅπερ ἔδει δεῖχαι.}}

\ParallelRText{
\begin{center}
{\large Proposition 57}
\end{center}

If an area is contained by a rational (straight-line)
and a fourth binomial (straight-line) then the square-root of the area
is the irrational (straight-line which is) called major.$^\dag$

\epsfysize=1.3in
\centerline{\epsffile{Book10/fig054e.eps}}

For let the area $AC$ be contained by the rational (straight-line) $AB$
and the fourth binomial (straight-line) $AD$, which has been divided into
its (component) terms at $E$, of which let $AE$ be the greater. I say that
the square-root of $AC$ is the irrational (straight-line which is) called major.

For since $AD$ is a fourth binomial (straight-line), $AE$ and $ED$
are thus rational (straight-lines which are) commensurable in square only,
and the square on $AE$ is greater than (the square on) $ED$ by the
(square) on (some straight-line) incommensurable (in length) with ($AE$),
and $AE$ [is] commensurable in length with $AB$ [Def.~10.8]. Let $DE$ be cut in half at $F$,
and let the parallelogram (contained by) $AG$ and $GE$, equal to the
(square) on $EF$, (and falling short by a square figure) be applied to
$AE$. $AG$ is thus incommensurable in length with $GE$ [Prop.~10.18]. Let $GH$, $EK$, and $FL$
be drawn parallel to $AB$, and let the rest (of the construction) be made the same as the (proposition) before this. So, it is clear that
$MO$ is the square-root of  area $AC$. So, we must show that $MO$
is the irrational (straight-line which is) called major.

Since $AG$ is incommensurable in length with $EG$, $AH$
is also incommensurable with $GK$---that is to say, $SN$
 with $NQ$ [Props.~6.1, 10.11].  Thus, $MN$ and $NO$
 are incommensurable in square. And since $AE$ is
 commensurable in length with $AB$, $AK$ is rational [Prop.~10.19].  And it is equal to the (sum of the squares) on $MN$ and $NO$.  Thus, the sum of the (squares) on 
 $MN$ and $NO$ [is] also rational. And since $DE$ [is]  incommensurable
 in length with $AB$ [Prop.~10.13]---that is to say, with $EK$---but $DE$ is commensurable (in length) with $EF$, $EF$ (is) thus incommensurable
 in length with $EK$ [Prop.~10.13].
 Thus, $EK$ and $EF$ are rational (straight-lines which are) commensurable
 in square only. $LE$---that is to say, $MR$---(is) thus medial
 [Prop.~10.21]. And it is contained by $MN$ and $NO$. The (rectangle contained) by $MN$ and $NO$ is thus medial. And
 the [sum] of the (squares) on $MN$ and $NO$ (is) rational, and
 $MN$ and $NO$ are incommensurable in square. And if two
 straight-lines (which are) incommensurable in square, making the
 sum of the squares on them rational, and the (rectangle contained) by them
 medial, are added together, then the whole is the irrational (straight-line which is) called major [Prop.~10.39].
 
 Thus, $MO$ is the irrational (straight-line which is) called major.
 And (it is) the square-root of area $AC$. (Which is) the very thing it
 was required to show.}
\end{Parallel}
 {\footnotesize\noindent$^\dag$ If the rational straight-line has unit length then this proposition states that the square-root of 
a fourth binomial straight-line is a major straight-line: {i.e.}, 
a fourth binomial straight-line has a length $k\,(1+1/\sqrt{1+k'})$ whose
square-root can be written $\rho\,\sqrt{[1+k''/(1+k''^{\,2})^{1/2}]/2}+\rho\,\sqrt{[1-k''/(1+k''^{\,2})^{1/2}]/2}$, where $\rho=\sqrt{k}$ and $k''^{\,2}=k'$. This is the length of a major straight-line (see Prop.~10.39), since $\rho$ is rational.}

%%%%
%10.58
%%%%
\pdfbookmark[1]{Proposition 10.58}{pdf10.58}
\begin{Parallel}{}{}
\ParallelLText{
\begin{center}
{\large \ggn{58}.}
\end{center}\vspace*{-7pt}

{\greekfont Ἐὰν θωρίον περιέθηται ὑπὸ ῥητῆς καὶ τῆς ἐκ δύο ὀνομάτων
πέμπτης, ἡ τὸ θωρίον δυναμένη ἄλογός ἐστιν ἡ καλουμένη
ῥητὸν καὶ μέσον δυναμένη.}

{\greekfont Θωρίον γὰρ τὸ ΑΓ περιεθέσθω ὑπὸ ῥητῆς τῆς ΑΒ καὶ τῆς ἐκ
δύο ὀνομάτων πέμπτης τῆς ΑΔ διῃρημένης εἰς τὰ ὀνόματα
κατὰ τὸ Ε, ὥστε τὸ μεῖζον ὄνομα εἶναι τὸ ΑΕ; λέγω [δή],
ὅτι ἡ τὸ ΑΓ θωρίον δυναμένη ἄλογός ἐστιν ἡ καλουμένη
ῥητὸν καὶ μέσον δυναμένη.}

{\greekfont Κατεσκευάσθω γὰρ τὰ α ὐτὰ τοῖς πρότερον δεδειγμένοις; φανερὸν
δή, ὅτι ἡ τὸ ΑΓ θωρίον δυναμένη ἐστὶν ἡ ΜΧ. δεικτέον
δή, ὅτι ἡ ΜΧ ἐστιν ἡ ῥητὸν καὶ μέσον δυναμένη.}\\~\\~\\

\epsfysize=1.35in
\centerline{\epsffile{Book10/fig054g.eps}}

{\greekfont Ἐπεὶ
γὰρ ἀσύμμετρός ἐστιν ἡ ΑΗ τῇ ΗΕ, ἀσύμμετ-ρον ἄρα ἐστὶ
καὶ τὸ ΑΘ τῷ ΘΕ, τουτέστι τὸ ἀπὸ τῆς ΜΝ τῷ ἀπὸ
τῆς ΝΧ; αἱ ΜΝ, ΝΧ ἄρα δυνάμει εἰσὶν ἀσύμμετροι. καὶ
ἐπεὶ ἡ ΑΔ ἐκ δύο ὀνομάτων ἐστὶ πέμπτη, καί [ἐστιν]
ἔλασσον α ὐτῆς τμῆμα τὸ ΕΔ, σύμμετρος ἄρα ἡ ΕΔ τῇ ΑΒ
μήκει. ἀλλὰ ἡ ΑΕ τῇ ΕΔ ἐστιν ἀσύμμετρος; καὶ ἡ ΑΒ ἄρα
τῇ ΑΕ ἐστιν ἀσύμμετρος μήκει [αἱ ΒΑ, ΑΕ ῥηταί εἰσι δυνάμει
μόνον σύμμετροι]; μέσον ἄρα ἐστὶ τὸ ΑΚ, τουτέστι τὸ συγκείμενον
ἐκ τῶν ἀπὸ τῶν ΜΝ, ΝΧ. καὶ ἐπεὶ σύμμετρός ἐστιν ἡ ΔΕ
τῇ ΑΒ μήκει, τουτέστι τῇ ΕΚ, ἀλλὰ ἡ ΔΕ τῇ ΕΖ σύμμετρός
ἐστιν, καὶ ἡ ΕΖ ἄρα τῇ ΕΚ σύμμετρός ἐστιν. καὶ ῥητὴ ἡ
ΕΚ; ῥητὸν ἄρα καὶ τὸ ΕΛ, τουτέστι τὸ ΜΡ, τουτέστι τὸ ὑπὸ
ΜΝΧ; αἱ ΜΝ, ΝΧ ἄρα δυνάμει ἀσύμμετροί εἰσι ποιοῦσαι
τὸ μὲν συγκείμενον ἐκ τῶν ἀπ α ὐτῶν τετραγώνων
μέσον, τὸ δ ὑπ α ὐτῶν ῥητόν.}

{\greekfont Ἡ ΜΧ ἄρα ῥητὸν καὶ μέσον δυναμένη ἐστὶ καὶ δύναται
τὸ ΑΓ θωρίον; ὅπερ ἔδει δεῖχαι.}}

\ParallelRText{
\begin{center}
{\large Proposition 58}
\end{center}

If an area is contained by a rational (straight-line)
and a fifth binomial (straight-line) then the square-root of the area is
the irrational (straight-line which is) called the square-root of a rational
plus a medial (area).$^\dag$

For let the area $AC$ be contained by the rational (straight-line) $AB$
and the fifth binomial (straight-line) $AD$, which has been divided into its
(component) terms at $E$, such that $AE$ is the greater term. [So] I say that
the square-root of area $AC$ is the irrational (straight-line which is)
called the square-root of a rational plus a medial (area).

For let the same construction  be made as that shown previously.
So, (it is) clear that $MO$ is the square-root of area $AC$.
So, we must show that $MO$ is the square-root of a rational plus a medial
(area).

\epsfysize=1.35in
\centerline{\epsffile{Book10/fig054e.eps}}

 For since $AG$ is incommensurable (in length) with $GE$ [Prop.~10.18], $AH$ is thus also
 incommensurable with $HE$---that is to say, the (square) on $MN$
 with the (square) on $NO$ [Props.~6.1, 10.11]. Thus, $MN$ and $NO$
 are incommensurable in square. And since $AD$ is a fifth
 binomial (straight-line), and $ED$ [is] its lesser segment, $ED$
 (is) thus commensurable in length with $AB$ [Def.~10.9]. But, $AE$ is incommensurable (in length) with $ED$. Thus, $AB$ is also incommensurable in length with $AE$
 [$BA$ and $AE$ are rational (straight-lines which are) commensurable
 in square only] [Prop.~10.13]. Thus, $AK$---that
 is to say, the sum of the (squares) on $MN$ and $NO$---is medial [Prop.~10.21]. And since $DE$ is
 commensurable in length with $AB$---that is to say, with $EK$---but,
 $DE$ is commensurable (in length) with $EF$, $EF$ is thus
also  commensurable (in length) with $EK$ [Prop.~10.12]. And $EK$ (is) rational. Thus, $EL$---that is to say, $MR$---that is to say, the (rectangle contained) by
$MNO$---(is) also rational [Prop.~10.19]. 
$MN$ and $NO$ are thus (straight-lines which are) incommensurable in square, making
the sum of the squares on them medial, and the (rectangle contained)
by them rational.

Thus, $MO$ is the square-root of a rational plus a medial (area) [Prop.~10.40]. And (it is) the square-root
of area $AC$. (Which is) the very thing it was required to show.}
\end{Parallel}
{\footnotesize\noindent $^\dag$
If the rational straight-line has unit length then this proposition states that the square-root of 
a fifth binomial straight-line is the square root of a rational plus a medial area: {i.e.}, 
a fifth binomial straight-line has a length $k\,(\sqrt{1+k'}+1)$ whose
square-root can be written\\ $\rho\,\sqrt{[(1+k''^{\,2})^{1/2}+k'']/[2\,(1+k''^{\,2})]}+\rho\,\sqrt{[(1+k''^{\,2})^{1/2}-k'']/[2\,(1+k''^{\,2})]}$, where $\rho=\sqrt{k\,(1+k''^{\,2})}$ and $k''^{\,2}=k'$. This is the length of the square root of a rational plus a medial area (see Prop.~10.40), since $\rho$ is rational.}

%%%%
%10.59
%%%%
\pdfbookmark[1]{Proposition 10.59}{pdf10.59}
\begin{Parallel}{}{}
\ParallelLText{
\begin{center}
{\large \ggn{59}.}
\end{center}\vspace*{-7pt}

{\greekfont Ἐὰν θωρίον περιέθηται ὑπὸ ῥητῆς καὶ τῆς ἐκ δύο ὀνομάτων
ἕκτης, ἡ τὸ θωρίον δυναμένη ἄλογός ἐστιν ἡ καλουμένη
δύο μέσα δυναμένη.}\\

\epsfysize=1.35in
\centerline{\epsffile{Book10/fig054g.eps}}

{\greekfont Θωρίον γὰρ τὸ ΑΒΓΔ περιεθέσθω ὑπὸ ῥητῆς τῆς ΑΒ καὶ τῆς
ἐκ δύο ὀνομάτων ἕκτης τῆς ΑΔ διῃρημένης εἰς τὰ ὀνόματα
κατὰ τὸ Ε, ὥστε τὸ μεῖζον ὄνομα εἶναι τὸ ΑΕ; λέγω, ὅτι
ἡ τὸ ΑΓ δυναμένη ἡ δύο μέσα δυναμένη ἐστίν.}

{\greekfont Κατεσκευάσθω [γὰρ] τὰ α ὐτὰ τοῖς προδεδειγμένοις. φανερὸν δή,
ὅτι [ἡ] τὸ ΑΓ δυναμένη ἐστὶν ἡ ΜΧ, καὶ ὅτι ἀσύμμετρός ἐστιν
ἡ ΜΝ τῇ ΝΧ δυνάμει. καὶ ἐπεὶ ἀσύμμετρός ἐστιν ἡ ΕΑ τῇ
ΑΒ μήκει, αἱ ΕΑ, ΑΒ ἄρα ῥηταί εἰσι δυνάμει μόνον σύμμετροι; μέσον ἄρα ἐστὶ τὸ ΑΚ, τουτέστι τὸ συγκείμενον ἐκ τῶν ἀπὸ
τῶν ΜΝ, ΝΧ. πάλιν, ἐπεὶ ἀσύμμετρός ἐστιν ἡ ΕΔ τῇ ΑΒ μήκει,
ἀσύμμετρος ἄρα ἐστὶ καὶ ἡ ΖΕ τῇ ΕΚ; αἱ ΖΕ, ΕΚ ἄρα
ῥηταί εἰσι δυνάμει μόνον σύμμετροι; μέσον ἄρα ἐστὶ τὸ ΕΛ,
τουτέστι τὸ ΜΡ, τουτέστι τὸ ὑπὸ τῶν ΜΝΧ. καὶ ἐπεὶ ἀσύμμετρος
ἡ ΑΕ τῇ ΕΖ, καὶ τὸ ΑΚ τῷ ΕΛ ἀσύμμετρόν ἐστιν. ἀλλὰ τὸ μὲν
ΑΚ ἐστι τὸ συγκείμενον ἐκ τῶν ἀπὸ τῶν ΜΝ, ΝΧ,
τὸ δὲ ΕΛ ἐστι τὸ ὑπὸ τῶν ΜΝΧ; ἀσύμμετρον ἄρα ἐστὶ τὸ
συγκείμενον ἐκ τῶν ἀπὸ τῶν ΜΝΧ τῷ ὑπὸ
τῶν ΜΝΧ. καί ἐστι μέσον ἑκάτερον α ὐτῶν, καὶ αἱ ΜΝ, ΝΧ
δυνάμει εἰσὶν ἀσύμμετροι.}

{\greekfont Ἡ ΜΧ ἄρα δύο μέσα δυναμένη ἐστὶ καὶ δύναται τὸ 
ΑΓ; ὅπερ ἔδει δεῖχαι.}}

\ParallelRText{
\begin{center}
{\large Proposition 59}
\end{center}

If an area is contained by a rational (straight-line)
and a sixth binomial (straight-line) then the square-root of the
area is the irrational (straight-line which is)  called the square-root
of (the sum of) two medial (areas).$^\dag$

\epsfysize=1.35in
\centerline{\epsffile{Book10/fig054e.eps}}

For let the area $ABCD$ be contained by the rational (straight-line)
$AB$ and the sixth binomial (straight-line) $AD$, which has been
divided into its (component) terms at $E$, such that $AE$ is the greater
term. So, I say that the square-root of $AC$ is the square-root
of (the sum of) two medial (areas).

\mbox{[}For] let the same construction  be made as that shown previously. So, (it is)
clear that $MO$ is the square-root of $AC$, and that $MN$ is incommensurable
in square with $NO$. And since $EA$ is incommensurable
in length with $AB$ [Def.~10.10], $EA$ and $AB$ are thus rational (straight-lines
which are) commensurable in square only. Thus, $AK$---that is to
say, the sum of the (squares) on $MN$ and $NO$---is medial [Prop.~10.21]. Again, since  $ED$ is incommensurable in length with $AB$  [Def.~10.10], $FE$ is thus also incommensurable 
(in length) with $EK$ [Prop.~10.13]. 
Thus, $FE$ and $EK$ are rational (straight-lines which are) commensurable
in square only. Thus, $EL$---that is to say,  $MR$---that is to say,
the (rectangle contained) by $MNO$---is medial [Prop.~10.21]. And since $AE$ is incommensurable
(in length) with $EF$, $AK$ is also incommensurable with $EL$
[Props.~6.1, 10.11]. 
But, $AK$ is the sum of the (squares) on $MN$ and $NO$, and $EL$
is the (rectangle contained) by $MNO$. Thus,
the sum of the (squares) on $MNO$ is incommensurable with the (rectangle
contained) by $MNO$. And each of them is medial.  And $MN$ and
$NO$ are incommensurable in square.

Thus, $MO$ is the square-root of (the sum of) two medial (areas) [Prop.~10.41].
And (it is) the square-root of $AC$. (Which is) the very thing it was required
to show.}
\end{Parallel}
{\footnotesize\noindent$^\dag$ If the rational straight-line has unit length then this proposition states that the square-root of 
a sixth binomial straight-line is the square root of the sum of  two  medial areas: {i.e.}, 
a sixth binomial straight-line has a length $\sqrt{k}+\sqrt{k'}$ whose
square-root can be written\\ $k^{1/4}\left(\sqrt{[1+k''/(1+k''^{\,2})^{1/2}]/2}+\sqrt{[1-k''/(1+k''^{\,2})^{1/2}]/2}\right)$, where $k''^{\,2}=(k-k')/k'$. This is the length of the square-root of the sum of two medial areas (see Prop.~10.41).}

%%%%
%10.59a
%%%%
\begin{Parallel}{}{}
\ParallelLText{
\begin{center}
{\large {\greekfont Λῆμμα}.}
\end{center}\vspace*{-7pt}

{\greekfont Ἐὰν εὐθεῖα γραμμὴ τμ ηθῇ εἰς ἄνισα, τὰ ἀπὸ τῶν ἀνίσων
τετράγωνα μείζονά ἐστι τοῦ δὶς ὑπὸ τῶν ἀνίσων περιεχομένου
ὀρθογωνίου.}

\epsfysize=0.3in
\centerline{\epsffile{Book10/fig059ag.eps}}

{\greekfont Ἔστω εὐθεῖα ἡ ΑΒ καὶ τετμήσθω εἰς ἄνισα κατὰ τὸ Γ, καὶ ἔστω
μείζων ἡ ΑΓ; λέγω, ὅτι τὰ ἀπὸ τῶν ΑΓ, ΓΒ μείζονά ἐστι
τοῦ δὶς ὑπὸ τῶν ΑΓ, ΓΒ.}

{\greekfont Τετμήσθω γὰρ ἡ ΑΒ δίθα κατὰ τὸ Δ. ἐπεὶ οὖν εὐθεῖα
γραμμὴ τέτμ ηται εἰς μὲν ἴσα κατὰ τὸ Δ, εἰς
δὲ ἄνισα κατὰ τὸ Γ, τὸ ἄρα ὑπὸ τῶν ΑΓ, ΓΒ μετὰ
τοῦ ἀπὸ ΓΔ ἴσον ἐστὶ τῷ ἀπὸ ΑΔ; ὥστε τὸ ὑπὸ
τῶν ΑΓ, ΓΒ ἔλαττόν ἐστι τοῦ ἀπὸ ΑΔ; τὸ ἄρα δὶς
ὑπὸ τῶν ΑΓ, ΓΒ ἔλαττον
 ἢ διπλάσιόν ἐστι τοῦ ἀπὸ
ΑΔ. ἀλλὰ τὰ ἀπὸ τῶν ΑΓ, ΓΒ διπλάσιά [ἐστι] τῶν
ἀπὸ τῶν ΑΔ, ΔΓ; τὰ ἄρα ἀπὸ τῶν ΑΓ, ΓΒ μείζονά
ἐστι τοῦ δὶς ὑπὸ τῶν ΑΓ, ΓΒ; ὅπερ ἔδει δεῖχαι.}}

\ParallelRText{
\begin{center}
{\large Lemma}
\end{center}

If a straight-line is cut unequally then 
(the sum of) the squares on the unequal (parts) is greater than twice
the rectangle contained by the unequal (parts).

\epsfysize=0.3in
\centerline{\epsffile{Book10/fig059ae.eps}}

Let $AB$ be a straight-line, and let it be cut unequally at $C$, and
let $AC$ be greater (than $CB$). I say that (the sum of) the (squares) on
$AC$ and $CB$ is greater than twice the (rectangle contained) by $AC$ and
$CB$.

For let $AB$ be cut in half at $D$. Therefore, since a straight-line
has been cut into equal (parts) at $D$, and into unequal (parts) at $C$,
 the (rectangle contained) by $AC$ and $CB$, plus the (square)
on $CD$, is thus equal to the (square) on $AD$ [Prop.~2.5]. Hence, the (rectangle contained) by $AC$ and $CB$ is less than the (square) on $AD$. Thus, twice the
(rectangle contained) by $AC$ and $CB$ is less than double the (square) on
$AD$. But, (the sum of) the (squares) on $AC$ and $CB$ [is] double (the sum of) the (squares) on $AD$ and $DC$ [Prop.~2.9].
Thus, (the sum of) the (squares) on $AC$ and $CB$ is greater than twice
the (rectangle contained) by $AC$ and $CB$. (Which is) the very thing it
was required to show.}
\end{Parallel}

%%%%
%10.60
%%%%
\pdfbookmark[1]{Proposition 10.60}{pdf10.60}
\begin{Parallel}{}{}
\ParallelLText{
\begin{center}
{\large \ggn{60}.}
\end{center}\vspace*{-7pt}

{\greekfont Τὸ ἀπὸ τῆς ἐκ δύο ὀνομάτων παρὰ ῥητὴν  παραβαλλόμενον πλάτος
ποιεῖ τὴν ἐκ δύο ὀνομάτων πρώτην.}\\

\epsfysize=1.5in
\centerline{\epsffile{Book10/fig060g.eps}}

{\greekfont Ἔστω ἐκ δύο ὀνομάτων ἡ ΑΒ διῃρημένη εἰς τὰ ὀνόματα κατὰ
τὸ Γ, ὥστε τὸ μεῖζον ὄνομα εἶναι τὸ ΑΓ, καὶ ἐκκείσθω
ῥητὴ ἡ ΔΕ, καὶ τῷ ἀπὸ τῆς ΑΒ ἴσον παρὰ τὴν ΔΕ παραβεβλήσθω
τὸ ΔΕΖΗ πλάτος ποιοῦν τὴν ΔΗ; λέγω, ὅτι ἡ ΔΗ ἐκ δύο ὀνομάτων
ἐστὶ πρώτη.}

{\greekfont Παραβεβλήσθω γὰρ παρὰ τὴν ΔΕ τῷ μὲν ἀπὸ τῆς ΑΓ ἴσον τὸ ΔΘ,
τῷ δὲ ἀπὸ τῆς ΒΓ ἴσον τὸ ΚΛ; λοιπὸν ἄρα τὸ δὶς ὑπὸ τῶν
ΑΓ, ΓΒ ἴσον ἐστὶ τῷ ΜΖ. τετμήσθω ἡ ΜΗ δίθα κατὰ τὸ Ν, καὶ
παράλληλος ἤθθω ἡ ΝΧ [ἑκατέρᾳ τῶν ΜΛ, ΗΖ]. ἑκάτερον ἄρα
τῶν ΜΧ, ΝΖ ἴσον ἐστὶ τῷ ἅπαχ ὑπὸ τῶν ΑΓΒ. καὶ ἐπεὶ
ἐκ δύο ὀνομάτων ἐστὶν ἡ ΑΒ διῃρημένη εἰς τὰ ὀνόματα
κατὰ τὸ Γ, αἱ ΑΓ, ΓΒ ἄρα ῥηταί εἰσι δυνάμει μόνον σύμμετροι;
τὰ ἄρα ἀπὸ τῶν ΑΓ, ΓΒ ῥητά ἐστι καὶ σύμμετρα ἀλλήλοις;
ὥστε καὶ τὸ συγκείμενον ἐκ τῶν ἀπὸ τῶν ΑΓ, ΓΒ.
καί ἐστιν ἴσον τῷ ΔΛ; ῥητὸν ἄρα
ἐστὶ τὸ ΔΛ. καὶ παρὰ ῥητὴν τὴν ΔΕ παράκειται; ῥητὴ ἄρα ἐστὶν
ἡ ΔΜ καὶ σύμμετρος τῇ ΔΕ μήκει. πάλιν, ἐπεὶ αἱ ΑΓ, ΓΒ ῥηταί
εἰσι δυνάμει μόνον σύμμετροι, μέσον ἄρα ἐστὶ τὸ δὶς ὑπὸ
τῶν ΑΓ, ΓΒ, τουτέστι τὸ ΜΖ. καὶ παρὰ ῥητὴν τὴν ΜΛ παράκειται;
ῥητὴ ἄρα καὶ ἡ ΜΗ καὶ ἀσύμμετρος τῇ ΜΛ, τουτέστι
τῇ ΔΕ, μήκει. ἔστι δὲ καὶ ἡ ΜΔ ῥητὴ καὶ τῇ ΔΕ μήκει
σύμμετρος; ἀσύμμετρος ἄρα ἐστὶν ἡ ΔΜ τῇ ΜΗ μήκει. καί εἰσι
ῥηταί; αἱ ΔΜ, ΜΗ ἄρα ῥηταί εἰσι δυνάμει μόνον σύμμετροι;
ἐκ δύο ἄρα ὀνομάτων ἐστὶν ἡ ΔΗ. δεικτέον δή, ὅτι
καὶ πρώτη.}

{\greekfont Ἐπεὶ τῶν ἀπὸ τῶν ΑΓ, ΓΒ μέσον ἀνάλογόν ἐστι τὸ ὑπὸ τῶν
ΑΓΒ, καὶ τῶν ΔΘ, ΚΛ ἄρα μέσον ἀνάλογόν ἐστι τὸ ΜΧ. ἔστιν
ἄρα ὡς τὸ ΔΘ πρὸς τὸ ΜΧ, οὕτως τὸ ΜΧ πρὸς τὸ
ΚΛ, τουτέστιν ὡς ἡ ΔΚ πρὸς τὴν ΜΝ, ἡ ΜΝ πρὸς τὴν ΜΚ;
τὸ ἄρα ὑπὸ τῶν ΔΚ, ΚΜ ἴσον ἐστὶ τῷ ἀπὸ τῆς ΜΝ.
καὶ ἐπεὶ σύμμετρόν ἐστι τὸ ἀπὸ τῆς ΑΓ τῷ ἀπὸ τῆς ΓΒ,
σύμμετρόν ἐστι καὶ τὸ ΔΘ τῷ ΚΛ; ὥστε καὶ ἡ ΔΚ τῇ
ΚΜ σύμμετρός ἐστιν. καὶ ἐπεὶ μείζονά ἐστι τὰ ἀπὸ τῶν ΑΓ,
ΓΒ τοῦ δὶς ὑπὸ τῶν ΑΓ, ΓΒ, μεῖζον ἄρα καὶ τὸ ΔΛ τοῦ
ΜΖ; ὥστε καὶ ἡ ΔΜ τῆς ΜΗ μείζων ἐστίν. καί ἐστιν ἴσον
τὸ ὑπὸ τῶν ΔΚ, ΚΜ τῷ ἀπὸ τῆς ΜΝ, τουτέστι τῷ τετάρτῳ
 τοῦ ἀπὸ τῆς ΜΗ  -.7πτ η, καὶ σύμμετρος ἡ ΔΚ τῇ ΚΜ. ἐὰν
δὲ ὦσι δύο εὐθεῖαι ἄνισοι,  τῷ δὲ τετάρτῳ μέρει τοῦ ἀπὸ τῆς
ἐλάσσονος ἴσον παρὰ τὴν μείζονα παραβληθῇ ἐλλεῖπον
εἴδει τετραγώνῳ  καὶ εἰς σύμμετρα α ὐτὴν διαιρῇ, ἡ μείζων
τῆς ἐλάσσονος μεῖζον δύναται τῷ ἀπὸ συμμέτρου ἑαυτῇ; ἡ
ΔΜ ἄρα τῆς ΜΗ μεῖζον δύναται τῷ ἀπὸ συμμέτρου
ἑαυτῇ. καί εἰσι ῥηταὶ αἱ ΔΜ, ΜΗ , καὶ ἡ ΔΜ μεῖζον
ὄνομα οὖσα σύμμετρός ἐστι τῇ ἐκκειμένῃ ῥητῇ τῇ
ΔΕ μήκει.}

{\greekfont Ἡ ΔΗ ἄρα ἐκ δύο ὀνομάτων ἐστὶ πρώτη; ὅπερ ἔδει
δεῖχαι.}}

\ParallelRText{
\begin{center}
{\large Proposition 60}
\end{center}

The square on a binomial (straight-line) applied to a rational (straight-line) produces as breadth a first binomial (straight-line).$^\dag$

\epsfysize=1.5in
\centerline{\epsffile{Book10/fig060e.eps}}

Let $AB$ be a binomial (straight-line), having been divided into its
(component) terms at $C$, such that $AC$ is the greater term. And let the
rational (straight-line) $DE$ be laid down. And let the (rectangle)
$DEFG$, equal to the (square) on $AB$, be applied to $DE$,
producing $DG$ as breadth. I say that $DG$ is a first binomial (straight-line).

For let $DH$, equal to the (square) on $AC$,  and $KL$, equal to the
(square) on $BC$, be applied to $DE$. Thus, the remaining twice
the (rectangle contained) by $AC$ and $CB$ is equal to $MF$ [Prop.~2.4].  Let $MG$ be cut in half
at $N$, and let $NO$ be drawn parallel [to each of $ML$ and
$GF$]. $MO$ and $NF$ are thus each equal to once the (rectangle contained)
by $ACB$. And since $AB$ is a binomial (straight-line), having been divided into its (component) terms at $C$,  $AC$ and $CB$ are thus rational (straight-lines which are) commensurable in square only [Prop.~10.36]. Thus, the (squares) on $AC$ and $CB$ are rational, and commensurable with one another.  And hence the
sum of the (squares) on $AC$ and $CB$ (is rational) [Prop.~10.15], and is equal to $DL$. Thus, $DL$
is rational. And it is applied to the rational (straight-line) $DE$. $DM$
is thus rational, and commensurable in length with $DE$ [Prop.~10.20]. Again, since $AC$ and $CB$
are rational (straight-lines which are) commensurable in square only, twice the (rectangle contained) by $AC$ and $CB$---that is to say, $MF$---is thus medial [Prop.~10.21]. And it is applied to the
rational (straight-line) $ML$. $MG$ is thus also rational, and incommensurable in length with $ML$---that is to say, with $DE$
[Prop.~10.22].  And $MD$ is also rational, and
commensurable in length with $DE$. Thus, $DM$ is  incommensurable in length with $MG$ [Prop.~10.13]. And
they are rational. $DM$ and $MG$ are thus rational (straight-lines which
are) commensurable in square only. Thus, $DG$ is a binomial (straight-line)
[Prop.~10.36]. So, we must show that (it is)
also a first (binomial straight-line).

Since the (rectangle contained) by $ACB$ is the mean proportional to
 the squares on $AC$ and $CB$ [Prop.~10.53~lem.], $MO$ is thus also the
mean proportional to   $DH$ and $KL$. Thus,
as $DH$ is to $MO$, so $MO$ (is) to $KL$---that is to say, as $DK$
(is) to $MN$, (so) $MN$ (is) to $MK$ [Prop.~6.1].
Thus, the (rectangle contained) by $DK$ and $KM$ is equal
to the (square) on $MN$ [Prop.~6.17]. And since
the (square) on $AC$ is commensurable with the (square) on $CB$,
$DH$ is also commensurable with $KL$. Hence, $DK$
is also commensurable with $KM$ [Props.~6.1, 10.11].
And since (the sum of) the squares on $AC$ and $CB$ is greater
than twice the (rectangle contained) by $AC$ and $CB$ [Prop.~10.59~lem.], $DL$ (is) thus also
greater than $MF$. Hence, $DM$ is also greater than $MG$
[Props.~6.1, 5.14]. And
the (rectangle contained) by $DK$ and $KM$ is equal to
the (square) on $MN$---that is to say, to one quarter the (square) on $MG$.
And $DK$ (is) commensurable (in length) with $KM$. And if there are
two unequal straight-lines, and a (rectangle) equal to the fourth part
of the (square) on the lesser, falling short by a square figure, is applied to
the greater, and divides it into (parts which are) commensurable (in length),
then the square on the greater is larger than (the square on) the lesser by
the (square) on (some straight-line) commensurable (in length) with the greater [Prop.~10.17]. Thus,  the square on $DM$
is greater than (the square on) $MG$ by the (square) on (some straight-line)
commensurable (in length) with ($DM$). And $DM$ and $MG$ are rational. And
$DM$, which is the greater term, is commensurable in length with the (previously) laid down rational (straight-line) $DE$.

Thus, $DG$ is a first binomial (straight-line) [Def.~10.5]. (Which is) the very thing it was required to show.}
\end{Parallel}
{\footnotesize\noindent$^\dag$ In other words, the square of a binomial  is a
first binomial. See Prop.~10.54.}

%%%%
%10.61
%%%%
\pdfbookmark[1]{Proposition 10.61}{pdf10.61}
\begin{Parallel}{}{}
\ParallelLText{
\begin{center}
{\large \ggn{61}.}
\end{center}\vspace*{-7pt}

{\greekfont Τὸ ἀπὸ τῆς ἐκ δύο μέσων πρώτης παρὰ ῥητὴν παραβαλλόμενον πλάτος ποιεῖ τὴν ἐκ δύο ὀνομάτων δευτέραν.}

{\greekfont Ἔστω ἐκ δύο μέσων πρώτη ἡ ΑΒ διῃρημένη εἰς τὰς μέσας κατὰ
τὸ Γ, ὧν μείζων ἡ ΑΓ, καὶ ἐκκείσθω ῥητὴ ἡ ΔΕ, καὶ παραβεβλήσθω παρὰ τὴν ΔΕ τῷ ἀπὸ τῆς ΑΒ ἴσον παραλληλόγραμμον
τὸ ΔΖ πλάτος ποιοῦν τὴν ΔΗ; λέγω, ὅτι ἡ ΔΗ ἐκ δύο ὀνομάτων ἐστὶ δευτέρα.}

{\greekfont Κατεσκευάσθω γὰρ τὰ α ὐτὰ τοῖς πρὸ τούτου. καὶ ἐπεὶ ἡ ΑΒ
ἐκ δύο μέσων ἐστὶ πρώτη διῃρημένη κατὰ τὸ Γ, αἱ ΑΓ, ΓΒ ἄρα
μέσαι εἰσὶ δυνάμει μόνον σύμμετροι ῥητὸν περιέθουσαι; ὥστε
καὶ τὰ ἀπὸ τῶν ΑΓ, ΓΒ μέσα ἐστίν. μέσον ἄρα ἐστὶ τὸ ΔΛ.
καὶ παρὰ ῥητὴν τὴν ΔΕ παραβέβληται; ῥητὴ ἄρα ἐστίν ἡ ΜΔ
καὶ ἀσύμμετρος τῇ ΔΕ μήκει. πάλιν, ἐπεὶ ῥητόν ἐστι τὸ δὶς
ὑπὸ τῶν ΑΓ, ΓΒ, ῥητόν ἐστι καὶ τὸ ΜΖ. καὶ παρὰ ῥητὴν τὴν
ΜΛ παράκειται; ῥητὴ ἄρα [ἐστὶ] καὶ ἡ ΜΗ καὶ μήκει σύμμετρος
τῇ ΜΛ, τουτέστι τῇ ΔΕ; ἀσύμμετρος ἄρα ἐστὶν ἡ ΔΜ τῇ
ΜΗ  μήκει. καί εἰσι ῥηταί; αἱ ΔΜ, ΜΗ ἄρα ῥηταί εἰσι δυνάμει
μόνον σύμμετροι; ἐκ δύο ἄρα ὀνομάτων ἐστὶν ἡ ΔΗ. δεικτέον
δή, ὅτι καὶ δευτέρα.}\\~\\~\\~\\~\\~\\~\\~\\

\epsfysize=1.55in
\centerline{\epsffile{Book10/fig060g.eps}}

{\greekfont Ἐπεὶ γὰρ τὰ ἀπὸ τῶν ΑΓ, ΓΒ μείζονά ἐστι τοῦ δὶς ὑπὸ τῶν
ΑΓ, ΓΒ, μεῖζον ἄρα καὶ τὸ ΔΛ τοῦ ΜΖ; ὥστε καὶ ἡ ΔΜ τῆς
ΜΗ . καὶ ἐπεὶ σύμμετρόν ἐστι τὸ ἀπὸ τῆς ΑΓ τῷ ἀπὸ τῆς
ΓΒ, σύμμετρόν ἐστι καὶ τὸ ΔΘ τῷ ΚΛ; ὥστε καὶ ἡ ΔΚ τῇ ΚΜ
σύμμετρός ἐστιν. καί ἐστι τὸ ὑπὸ τῶν ΔΚΜ ἴσον τῷ ἀπὸ
τῆς ΜΝ; ἡ ΔΜ ἄρα τῆς ΜΗ μεῖζον δύναται τῷ ἀπὸ συμμέτρου
ἑαυτῇ. καί ἐστιν ἡ ΜΗ σύμμετρος τῇ ΔΕ μήκει.}

{\greekfont Ἡ ΔΗ ἄρα ἐκ δύο ὀνομάτων ἐστὶ δευτέρα.}}

\ParallelRText{
\begin{center}
{\large Proposition 61}
\end{center}

The square on a first bimedial (straight-line)
applied to a rational (straight-line) produces as breadth a second binomial
(straight-line).$^\dag$

Let $AB$ be a first bimedial (straight-line) having been divided into its (component)
medial (straight-lines) at $C$, of which $AC$ (is) the greater. And let
the rational (straight-line) $DE$ be laid down. And let the parallelogram
$DF$, equal to the (square) on $AB$, be applied to $DE$, producing
$DG$ as breadth. I say that $DG$ is a second binomial (straight-line).

For let the same construction be made as in the (proposition) before
this. And since $AB$ is a first bimedial (straight-line), having been divided at $C$, 
$AC$ and $CB$ are thus medial (straight-lines) commensurable
in square only, and containing a rational (area) [Prop.~10.37]. Hence, the (squares) on $AC$ and
$CB$ are also medial [Prop.~10.21]. Thus,
$DL$ is medial [Props.~10.15, 10.23~corr.].  And it has been applied to the
rational (straight-line) $DE$. $MD$ is thus rational, and incommensurable
in length with $DE$ [Prop.~10.22]. Again,
since twice the (rectangle contained) by $AC$ and $CB$ is rational,
$MF$ is also rational. And it is applied to the rational (straight-line)
$ML$. Thus, $MG$ [is] also rational, and commensurable in length
with $ML$---that is to say, with $DE$ [Prop.~10.20]. $DM$ is thus incommensurable in length with $MG$ [Prop.~10.13]. And they are rational. $DM$ and $MG$ are thus rational, and commensurable in square only. $DG$ is thus a binomial (straight-line) [Prop.~10.36]. So, we must show that (it
is) also a second (binomial straight-line).

\epsfysize=1.55in
\centerline{\epsffile{Book10/fig060e.eps}}

For since (the sum of) the squares on $AC$ and $CB$ is greater than twice
the (rectangle contained) by $AC$ and $CB$ [Prop.~10.59], $DL$ (is) thus also greater
than $MF$. Hence, $DM$ (is) also (greater) than $MG$ [Prop.~6.1]. And since
the (square) on $AC$ is commensurable with the (square) on $CB$, $DH$
is also commensurable with $KL$. Hence, $DK$ is also commensurable
(in length) with $KM$ [Props.~6.1, 10.11].  And the (rectangle contained) by $DKM$
is equal to the (square) on $MN$. Thus, the square on $DM$ is greater
than (the square on) $MG$ by the (square) on (some straight-line)
commensurable (in length) with ($DM$) [Prop.~10.17]. And $MG$ is commensurable in length
with $DE$.

Thus, $DG$ is a second binomial (straight-line) [Def. 10.6].}
\end{Parallel}
{\footnotesize\noindent $^\dag$In other words, the square of a first bimedial is a
second binomial. See Prop.~10.55.}

%%%%
%10.62
%%%%
\pdfbookmark[1]{Proposition 10.62}{pdf10.62}
\begin{Parallel}{}{}
\ParallelLText{
\begin{center}
{\large \ggn{62}.}
\end{center}\vspace*{-7pt}

{\greekfont Τὸ ἀπὸ τῆς ἐκ δύο μέσων δευτέρας παρὰ ῥητὴν παραβαλλόμενον
πλάτος ποιεῖ τὴν ἐκ δύο ὀνομάτων τρίτην.}

{\greekfont Ἔστω ἐκ δύο μέσων δευτέρα ἡ ΑΒ διῃρημένη εἰς τὰς μέσας
κατὰ τὸ Γ, ὥστε τὸ μεῖζον τμῆμα εἶναι τὸ ΑΓ, ῥητὴ δέ τις
ἔστω ἡ ΔΕ, καὶ παρὰ τὴν ΔΕ τῷ ἀπὸ τῆς ΑΒ ἴσον παραλληλόγραμμον παραβεβλήσθω τὸ ΔΖ πλάτος ποιοῦν τὴν ΔΗ;
λέγω, ὅτι ἡ ΔΗ ἐκ δύο ὀνομάτων ἐστὶ τρίτη.}

{\greekfont Κατεσκευάσθω τὰ α ὐτὰ τοῖς προδεδειγμένοις. καὶ ἐπεὶ ἐκ δύο
μέσων δευτέρα ἐστὶν ἡ ΑΒ διῃρημένη κατὰ τὸ Γ, αἱ ΑΓ, ΓΒ
ἄρα μέσαι εἰσὶ δυνάμει μόνον σύμμετροι μέσον
περιέθουσαι; ὥστε καὶ τὸ συγκείμενον ἐκ τῶν ἀπὸ τῶν ΑΓ, ΓΒ
μέσον ἐστίν. καί ἐστιν ἴσον τῷ ΔΛ; μέσον ἄρα καὶ τὸ ΔΛ.
καὶ παράκειται παρὰ ῥητὴν τὴν ΔΕ; ῥητὴ ἄρα ἐστὶ καὶ ἡ
ΜΔ καὶ ἀσύμμετρος τῇ ΔΕ μήκει. διὰ τὰ α ὐτὰ δὴ καὶ ἡ
ΜΗ  ῥητή ἐστι καὶ ἀσύμμετρος τῇ ΜΛ, τουτέστι τῇ ΔΕ, μήκει; ῥητὴ
ἄρα ἐστὶν ἑκατέρα τῶν ΔΜ, ΜΗ καὶ ἀσύμμετρος τῇ ΔΕ μήκει.
καὶ ἐπεὶ ἀσύμμετρός ἐστιν ἡ ΑΓ τῇ ΓΒ μήκει, ὡς
δὲ ἡ ΑΓ πρὸς τὴν ΓΒ, οὕτως τὸ ἀπὸ τῆς ΑΓ πρὸς τὸ ὑπὸ τῶν
ΑΓΒ, ἀσύμμετρον ἄρα καὶ τὸ ἀπὸ τῆς ΑΓ τῷ ὑπὸ
τῶν ΑΓΒ. ὥστε καὶ τὸ συγκείμενον ἐκ τῶν ἀπὸ τῶν ΑΓ, ΓΒ
τῷ δὶς ὑπὸ τῶν ΑΓΒ ἀσύμμετρόν ἐστιν, τουτέστι τὸ ΔΛ τῷ
ΜΖ;  ὥστε καὶ ἡ ΔΜ τῷ ΜΗ ἀσύμμετρός ἐστιν. καί εἰσι
ῥηταί; ἐκ δύο ἄρα ὀνομάτων ἐστὶν ἡ ΔΗ. δεικτέον
[δή], ὅτι καὶ τρίτη.}\\~\\~\\~\\~\\~\\~\\~\\~\\

\epsfysize=1.5in
\centerline{\epsffile{Book10/fig060g.eps}}

{\greekfont Ὁμοίως δὴ τοῖς προτέροις ἐπιλογιούμεθα, ὅτι μείζων ἐστὶν ἡ ΔΜ
τῆς ΜΗ , καὶ σύμμετρος ἡ ΔΚ τῇ ΚΜ. καί ἐστι τὸ ὑπὸ τῶν
ΔΚΜ ἴσον τῷ ἀπὸ τῆς ΜΝ; ἡ ΔΜ ἄρα τῆς ΜΗ μεῖζον
δύναται τῷ ἀπὸ συμμέτρου ἑαυτῇ. καὶ οὐδετέρα τῶν ΔΜ, ΜΗ 
σύμμετρός ἐστι τῇ ΔΕ μήκει.}

{\greekfont Ἡ ΔΗ ἄρα ἐκ δύο ὀνομάτων ἐστὶ τρίτη; ὅπερ ἔδει
δεῖχαι.}}

\ParallelRText{
\begin{center}
{\large Proposition 62}
\end{center}

The square on a second bimedial (straight-line)
applied to  a rational (straight-line) produces as breadth a third binomial
(straight-line).$^\dag$

Let $AB$ be a second bimedial (straight-line) having been divided into its (component)
medial (straight-lines) at $C$, such that $AC$ is the greater segment. And
let $DE$ be some rational (straight-line). And let the parallelogram
$DF$, equal to the (square) on $AB$, be applied to $DE$,
producing $DG$ as breadth. I say that $DG$ is a third binomial (straight-line).

Let the same construction  be made as that shown previously. And since $AB$ is a second bimedial (straight-line), having been divided at $C$, $AC$ and
$CB$ are thus medial (straight-lines) commensurable in square only,
and containing a medial (area) [Prop.~10.38]. Hence, the sum  of the (squares) on
$AC$ and $CB$ is also medial [Props.~10.15, 10.23~corr.]. And it is equal to $DL$. Thus, $DL$ (is)
also medial. And it is applied to the rational (straight-line) $DE$. $MD$
is thus also rational, and incommensurable in length with $DE$
[Prop.~10.22]. So, for the same (reasons),
$MG$ is also rational, and incommensurable in length with $ML$---that is to say, with $DE$. Thus, $DM$ and $MG$ are each rational, and incommensurable in length with $DE$. And since	$AC$ is incommensurable
in length with $CB$, and as $AC$ (is) to $CB$, so the (square)
on $AC$ (is) to the (rectangle contained) by $ACB$ [Prop.~10.21~lem.], the (square) on $AC$
(is) also incommensurable with the (rectangle contained) by $ACB$
[Prop.~10.11]. And hence the sum of the
(squares) on $AC$ and $CB$ is incommensurable with twice the (rectangle
contained) by $ACB$---that is to say, $DL$ with $MF$ [Props.~10.12, 10.13].  Hence,
$DM$ is also incommensurable (in length) with $MG$ [Props.~6.1, 10.11].
And they are rational.  $DG$ is thus a binomial (straight-line)
[Prop.~10.36]. [So]
we must show that (it is) also a third (binomial straight-line).

\epsfysize=1.5in 
\centerline{\epsffile{Book10/fig060e.eps}}

So, similarly to the previous (propositions), we can conclude that
$DM$ is greater than $MG$, and $DK$ (is) commensurable (in length) 
with $KM$.  And the (rectangle contained) by $DKM$ is equal to the
(square) on $MN$. Thus, the square on $DM$ is greater than (the square on)
$MG$ by the (square) on (some straight-line) commensurable (in length) with ($DM$)
[Prop.~10.17]. And neither of $DM$ and $MG$
is commensurable in length with $DE$.

Thus, $DG$ is a third binomial (straight-line) [Def. 10.7]. (Which is) the very thing it
was required to show.}
\end{Parallel}
{\footnotesize\noindent$^\dag$ In other words, the square of a second bimedial  is a
third binomial. See Prop.~10.56.}

%%%%
%10.63
%%%%
\pdfbookmark[1]{Proposition 10.63}{pdf10.63}
\begin{Parallel}{}{}
\ParallelLText{
\begin{center}
{\large\ggn{63}.}
\end{center}\vspace*{-7pt}

{\greekfont Τὸ ἀπὸ τῆς μείζονος παρὰ ῥητὴν παραβαλλόμεν-ον πλάτος ποιεῖ
τὴν ἐκ δύο ὀνομάτων τετάρτην.}

{\greekfont Ἔστω μείζων ἡ ΑΒ διῃρημένη κατὰ τὸ Γ, ὥστε μείζονα εἶναι
τὴν ΑΓ τῆς ΓΒ, ῥητὴ δὲ ἡ ΔΕ, καὶ τῷ ἀπὸ τῆς ΑΒ ἴσον
παρὰ τὴν ΔΕ παραβεβλήσθω τὸ ΔΖ παραλληλόγραμμον πλάτος ποιοῦν
τὴν ΔΗ; λέγω, ὅτι ἡ ΔΗ ἐκ δύο ὀνομάτων ἐστὶ τετάρτη.}\\

\epsfysize=1.6in
\centerline{\epsffile{Book10/fig060g.eps}}

{\greekfont Κατεσκευάσθω τὰ α ὐτὰ τοῖς προδεδειγμένοις. καὶ ἐπεὶ μείζων
ἐστὶν ἡ ΑΒ διῃρημένη κατὰ τὸ Γ, αἱ ΑΓ, ΓΒ δυνάμει εἰσὶν ἀσύμμετροι ποιοῦσαι τὸ μὲν συγκείμενον ἐκ τῶν ἀπ α ὐτῶν
τετραγώνων ῥητόν, τὸ δὲ ὑπ α ὐτῶν μέσον. ἐπεὶ οὖν ῥητόν
ἐστι τὸ συγκείμενον ἐκ τῶν ἀπὸ τῶν ΑΓ, ΓΒ, ῥητὸν ἄρα ἐστὶ
τὸ ΔΛ; ῥητὴ ἄρα καὶ ἡ ΔΜ καὶ σύμμετρος τῇ ΔΕ μήκει.
πάλιν, ἐπεὶ μέσον ἐστὶ τὸ δὶς ὑπὸ τῶν ΑΓ, ΓΒ, τουτέστι τὸ ΜΖ,
καὶ παρὰ ῥητήν ἐστι τὴν ΜΛ, ῥητὴ ἄρα ἐστὶ καὶ ἡ ΜΗ καὶ
ἀσύμμετρος  τῇ ΔΕ μήκει; ἀσύμμετρος ἄρα ἐστὶ καὶ ἡ ΔΜ
τῇ ΜΗ μήκει. αἱ ΔΜ, ΜΗ ἄρα ῥηταί εἰσι δυνάμει μόνον σύμμετροι; ἐκ δύο ἄρα ὀνομάτων ἐστὶν ἡ ΔΗ. δεικτέον
[δή], ὅτι καὶ τετάρτη.}

{\greekfont Ὁμοίως δὴ δείχομεν τοῖς πρότερον, ὅτι μείζων ἐστὶν ἡ ΔΜ
τῆς ΜΗ , καὶ ὅτι τὸ ὑπὸ ΔΚΜ ἴσον ἐστὶ τῷ ἀπὸ τῆς
ΜΝ. ἐπεὶ οὖν ἀσύμμετρόν ἐστι τὸ ἀπὸ τῆς ΑΓ τῷ ἀπὸ
τῆς ΓΒ, ἀσύμμετρον ἄρα ἐστὶ καὶ τὸ ΔΘ τῷ ΚΛ; ὥστε
ἀσύμμετρος καὶ ἡ ΔΚ τῇ ΚΜ ἐστιν. ἐὰν δὲ ὦσι δύο εὐθεῖαι
ἄνισοι, τῷ δὲ τετάρτῳ μέρει τοῦ ἀπὸ τῆς ἐλάσσονος ἴσον
παραλληλόγραμμον παρὰ τὴν μείζονα παραβληθῇ ἐλλεῖπον εἴδει
τετραγώνῳ καὶ εἰς ἀσύμμετρα α ὐτὴν διαιρῇ, ἡ μείζων τῆς
ἐλάσσονος μεῖζον δυνήσεται τῷ ἀπὸ ἀσύμμέτρου ἑαυτῇ
μήκει; ἡ ΔΜ ἄρα τῆς ΜΗ μεῖζον δύναται τῷ ἀπὸ ἀσυμμέτρου
ἑαυτῇ. καί εἰσιν αἱ ΔΜ, ΜΗ ῥηταὶ δυνάμει μόνον σύμμετροι,
καὶ ἡ ΔΜ σύμμετρός ἐστι τῇ ἐκκειμένῃ ῥητῇ τῇ ΔΕ.}

{\greekfont Ἡ ΔΗ ἄρα ἐκ δύο ὀνομάτων ἐστὶ τετάρτη; ὅπερ ἔδει δεῖχαι.}}

\ParallelRText{
\begin{center}
{\large Proposition 63}
\end{center}

The square on a major (straight-line) applied
to a rational (straight-line)  produces as breadth a fourth binomial (straight-line).$^\dag$

Let $AB$ be a major (straight-line) having been divided at $C$, such that
$AC$ is greater than $CB$, and (let) $DE$ (be)  a rational (straight-line). And
let the parallelogram $DF$, equal to the (square) on $AB$, be
applied to $DE$, producing $DG$ as breadth. I say that
$DG$ is a fourth binomial (straight-line).

\epsfysize=1.6in 
\centerline{\epsffile{Book10/fig060e.eps}}

Let the same construction be made as that shown previously.
And since $AB$ is a major (straight-line), having been divided at $C$, $AC$ and
$CB$ are incommensurable in square, making the sum of the squares
on them rational, and the (rectangle contained) by them medial [Prop.~10.39]. Therefore, since the sum of the
(squares) on $AC$ and $CB$ is rational, $DL$ is thus rational.
Thus, $DM$ (is) also rational, and commensurable in length with $DE$ [Prop.~10.20]. Again, since twice the (rectangle
contained) by $AC$ and $CB$---that is to say, $MF$---is medial,
and is (applied to) the rational (straight-line) $ML$,  $MG$
is thus also rational, and incommensurable in length with $DE$ [Prop.~10.22]. $DM$ is thus also incommensurable
in length with $MG$ [Prop.~10.13]. $DM$
and $MG$ are thus rational (straight-lines which are) commensurable
in square only. Thus, $DG$ is a binomial (straight-line) [Prop.~10.36]. [So] we must show that (it is)
also a fourth (binomial straight-line).

So, similarly to the previous (propositions), we can show that $DM$ is
greater than $MG$, and that the (rectangle contained) by $DKM$
is equal to the (square) on $MN$. Therefore, since the (square) on $AC$
is incommensurable with the (square) on $CB$, $DH$ is also incommensurable
with $KL$. Hence, $DK$ is also incommensurable with $KM$
[Props.~6.1, 10.11].
And if there are two unequal straight-lines, and a parallelogram
equal to the fourth part of the (square) on the lesser, falling short
by a square figure, is applied to the greater, and divides it
into (parts which are) incommensurable (in length), then the square on
the greater will be larger than (the square on) the lesser by the
(square) on (some straight-line) incommensurable in length
with the greater [Prop.~10.18]. Thus,
the square on $DM$ is greater than (the square on) $MG$ by the
(square) on (some straight-line) incommensurable (in length) with ($DM$). And
$DM$ and $MG$ are rational (straight-lines which are) commensurable
in square only. And $DM$ is commensurable (in length)
with the (previously) laid down rational (straight-line) $DE$.

Thus, $DG$ is a fourth binomial (straight-line) [Def. 10.8]. (Which is) the very thing it was required
to show.}
\end{Parallel}
{\footnotesize\noindent$^\dag$ In other words, the square of a major is a
fourth binomial. See Prop.~10.57.\\}

%%%%
%10.64
%%%%
\pdfbookmark[1]{Proposition 10.64}{pdf10.64}
\begin{Parallel}{}{}
\ParallelLText{
\begin{center}
{\large \ggn{64}.}
\end{center}\vspace*{-7pt}

{\greekfont Τὸ ἀπὸ τῆς ῥητὸν καὶ μέσον δυναμένης παρὰ ῥητὴν
παραβαλλόμενον πλάτος ποιεῖ τὴν ἐκ δύο ὀνομάτων
πέμπτην.}

{\greekfont Ἔστω ῥητὸν καὶ μέσον δυναμένη ἡ ΑΒ διῃρημένη εἰς
τὰς εὐθείας κατὰ τὸ Γ, ὥστε μείζονα εἶναι τὴν ΑΓ,
καὶ ἐκκείσθω ῥητὴ ἡ ΔΕ, καὶ τῷ ἀπὸ τῆς ΑΒ ἴσον
παρὰ τὴν ΔΕ παραβεβλήσθω τὸ ΔΖ πλάτος ποιοῦν τὴν ΔΗ;
λέγω, ὅτι ἡ ΔΗ ἐκ δύο ὀνομάτων ἐστὶ πέμπτη.}\\~\\

\epsfysize=1.5in
\centerline{\epsffile{Book10/fig060g.eps}}

{\greekfont Κατεσκευάσθω τὰ α ὐτα τοῖς πρὸ τούτου. ἐπεὶ οὖν
ῥητὸν καὶ μέσον δυναμένη ἐστὶν ἡ ΑΒ διῃρημένη
κατὰ τὸ Γ, αἱ ΑΓ, ΓΒ ἄρα δυνάμει εἰσὶν ἀσύμμετροι
ποιοῦσαι τὸ μὲν συγκείμενον ἐκ τῶν ἀπ α ὐτῶν
τετραγώνων μέσον, τὸ δ ὑπ α ὐτῶν ῥητόν. ἐπεὶ οὖν
μέσον ἐστὶ τὸ συγκείμενον ἐκ τῶν ἀπὸ τῶν ΑΓ, ΓΒ,
μέσον ἄρα ἐστὶ τὸ ΔΛ; ὥστε ῥητή ἐστιν ἡ ΔΜ καὶ
μήκει ἀσύμμετρος τῇ ΔΕ. πάλιν, ἐπεὶ ῥητόν ἐστι τὸ δὶς
ὑπὸ τῶν ΑΓΒ, τουτέστι τὸ ΜΖ, ῥητὴ ἄρα ἡ ΜΗ 
καὶ σύμμετρος τῇ ΔΕ. ἀσύμμετρος ἄρα ἡ ΔΜ τῇ
ΜΗ ; αἱ ΔΜ, ΜΗ ἄρα ῥηταί εἰσι δυνάμει μόνον σύμμετροι;
ἐκ δύο ἄρα ὀνομάτων ἐστὶν ἡ ΔΗ. λέγω δή, ὅτι καὶ
πέμπτη.}

{\greekfont Ὁμοίως γὰρ διεθθήσεται, ὅτι τὸ ὑπὸ τῶν ΔΚΜ ἴσον
ἐστὶ τῷ ἀπὸ τῆς ΜΝ, καὶ ἀσύμμετρος ἡ ΔΚ τῇ ΚΜ
μήκει; ἡ ΔΜ ἄρα τῆς ΜΗ μεῖζον δύναται τῷ ἀπὸ
ἀσυμμέτρου ἑαυτῇ. καί εἰσιν αἱ ΔΜ, ΜΗ [ῥηταὶ] δυνάμει
μόνον σύμμετροι, καὶ ἡ ἐλάσσων ἡ ΜΗ σύμμετρος τῇ
ΔΕ μήκει.}

{\greekfont Ἡ ΔΗ ἄρα ἐκ δύο ὀνομάτων ἐστὶ πέμπτη; ὅπερ
ἔδει δεῖχαι.}}

\ParallelRText{
\begin{center}
{\large Proposition 64}
\end{center}

The square on the square-root of a rational
plus a medial (area) applied to a rational (straight-line) produces as breadth
a fifth binomial (straight-line).$^\dag$

Let $AB$ be the square-root of a rational plus a medial (area) having been divided
into its (component) straight-lines at $C$, such that $AC$ is greater. 
And let the rational (straight-line) $DE$ be laid down. And let the
(parallelogram) $DF$, equal to the (square) on $AB$, be applied
to $DE$, producing $DG$ as breadth. I say that $DG$ is a fifth
binomial straight-line.

\epsfysize=1.5in 
\centerline{\epsffile{Book10/fig060e.eps}}

Let the same construction  be made as in the (propositions)
before this. Therefore, since $AB$ is the square-root of a
rational plus a medial (area), having been  divided at $C$, $AC$ and $CB$ are thus
incommensurable in square, making the sum of the squares on them
medial, and the (rectangle contained) by them rational [Prop.~10.40]. Therefore, since the
sum of the (squares) on $AC$ and $CB$ is medial, $DL$ is thus
medial. Hence, $DM$ is rational and incommensurable
in length with $DE$ [Prop.~10.22]. Again,
since twice the (rectangle contained) by $ACB$---that is to say,
$MF$---is rational, $MG$ (is) thus rational and commensurable
(in length) with $DE$ [Prop.~10.20]. 
$DM$ (is) thus incommensurable (in length) with $MG$ [Prop.~10.13]. Thus, $DM$ and $MG$
are rational (straight-lines which are) commensurable in square only. Thus,
$DG$ is a binomial (straight-line) [Prop.~10.36].
So, I say that (it is) also a fifth (binomial straight-line).

For, similarly (to the previous propositions), it can be shown that the
(rectangle contained) by $DKM$ is equal to the (square) on $MN$,
and $DK$ (is) incommensurable in length with $KM$. Thus, the
square on $DM$ is greater than (the square on) $MG$ by the (square)
on (some straight-line) incommensurable (in length) with ($DM$)
[Prop.~10.18]. And  $DM$ and $MG$
are [rational] (straight-lines which are) commensurable in square only,
and the lesser $MG$ is commensurable in length with $DE$.

Thus, $DG$ is a fifth binomial (straight-line) [Def.~10.9]. (Which is) the very thing it was required to show.}
\end{Parallel}
{\footnotesize\noindent$^\dag$ In other words, the square of the square-root of a rational plus medial  is a
fifth binomial. See Prop.~10.58.}

%%%%
%10.65
%%%%
\pdfbookmark[1]{Proposition 10.65}{pdf10.65}
\begin{Parallel}{}{}
\ParallelLText{
\begin{center}
{\large \ggn{65}.}
\end{center}\vspace*{-7pt}

{\greekfont Τὸ ἀπὸ τῆς δύο μέσα δυναμένης παρὰ ῥητὴν παραβαλλόμενον
πλάτος ποιεῖ τὴν ἐκ δύο ὀνομάτων ἕκτην.}

{\greekfont Ἔστω δύο μέσα δυναμένη ἡ ΑΒ διῃρημένη κατὰ τὸ Γ, ῥητὴ
δὲ ἔστω ἡ ΔΕ, καὶ παρὰ τὴν ΔΕ τῷ ἀπὸ τῆς ΑΒ ἴσον παραβεβλήσθω τὸ ΔΖ πλάτος ποιοῦν τὴν ΔΗ; λέγω, ὅτι ἡ ΔΗ
ἐκ δύο ὀνομάτων ἐστὶν ἕκτη.}\\~\\

\epsfysize=1.6in
\centerline{\epsffile{Book10/fig060g.eps}}

{\greekfont Κατεσκευάσθω γὰρ τὰ α ὐτὰ τοῖς πρότερον. καὶ ἐπεὶ ἡ ΑΒ
δύο μέσα δυναμένη ἐστὶ διῃρημένη κατὰ τὸ Γ, αἱ ΑΓ, ΓΒ
ἄρα δυνάμει εἰσὶν ἀσύμμετροι ποιοῦσαι τό τε συγκείμενον
ἐκ τῶν ἀπ α ὐτῶν τετραγώνων μέσον καὶ τὸ ὑπ α ὐτῶν
μέσον καὶ ἔτι ἀσύμμετρον τὸ ἐκ τῶν ἀπ α ὐτῶν
τετραγώνων συγκείμενον τῷ ὑπ α ὐτῶν; ὥστε κατὰ τὰ
προδεδειγμένα μέσον ἐστὶν ἑκάτερον τῶν ΔΛ, ΜΖ. καὶ παρὰ
ῥητὴν τὴν ΔΕ παράκειται; ῥητὴ ἄρα ἐστὶν ἑκατέρα τῶν
ΔΜ, ΜΗ καὶ ἀσύμμετρος τῇ ΔΕ μήκει. καὶ ἐπεὶ ἀσύμμετρόν
ἐστι τὸ συγκείμενον ἐκ τῶν ἀπὸ τῶν ΑΓ, ΓΒ τῷ δὶς
ὑπὸ τῶν ΑΓ, ΓΒ, ἀσύμμετρον ἄρα ἐστὶ τὸ ΔΛ τῷ ΜΖ. ἀσύμμετρος ἄρα καὶ ἡ ΔΜ τῇ ΜΗ ; αἱ ΔΜ, ΜΗ  ἄρα ῥηταί
εἰσι δυνάμει μόνον σύμμετροι; ἐκ δύο ἄρα ὀνομάτων ἐστὶν
ἡ ΔΗ. λέγω δή, ὅτι καὶ ἕκτη.}

{\greekfont Ὁμοίως δὴ πάλιν δεῖχομεν, ὅτι τὸ ὑπὸ τῶν ΔΚΜ ἴσον
ἐστὶ τῷ ἀπὸ τῆς ΜΝ, καὶ ὅτι ἡ ΔΚ τῇ ΚΜ μήκει ἐστὶν
ἀσύμμετρος; καὶ διὰ τὰ α ὐτὰ δὴ ἡ ΔΜ τῆς ΜΗ μεῖζον
δύναται τῷ ἀπὸ ἀσυμμέτρου ἑαυτῇ μήκει. καὶ
οὐδετέρα τῶν ΔΜ, ΜΗ σύμμετρός ἐστι τῇ ἐκκειμένῃ ῥητῇ
τῇ ΔΕ μήκει.}

{\greekfont Ἡ ΔΗ ἄρα ἐκ δύο ὀνομάτων ἐστὶν ἕκτη; ὅπερ
ἔδει δεῖχαι.}}

\ParallelRText{
\begin{center}
{\large Proposition 65}
\end{center}

The square on the square-root of (the sum of) two medial
(areas) applied to a rational (straight-line) produces as breadth  a sixth binomial (straight-line).$^\dag$

Let $AB$ be the square-root of (the sum of) two medial (areas), having been divided
at $C$. And let $DE$ be a  rational (straight-line). And let the (parallelogram) $DF$, equal to the (square) on $AB$, be applied to $DE$, producing $DG$
as breadth. I say that $DG$ is a sixth binomial (straight-line).

\epsfysize=1.6in 
\centerline{\epsffile{Book10/fig060e.eps}}

For  let the same construction  be made  as in the previous (propositions). And since $AB$ is the square-root of (the sum of) two medial
(areas), having been divided at $C$, $AC$ and $CB$ are thus incommensurable
in square, making the sum of the squares on them medial, and the (rectangle
contained) by them medial, and, moreover, the sum of the squares on them
incommensurable with the (rectangle contained) by them [Prop.~10.41]. Hence, according to
what has been previously  demonstrated, $DL$ and $MF$ are each medial.
And they are applied to the rational (straight-line) $DE$. Thus, $DM$ and
$MG$ are each rational, and incommensurable in length with $DE$ [Prop.~10.22]. And since the sum of the (squares)
on $AC$ and $CB$ is incommensurable with twice the (rectangle contained)
by $AC$ and $CB$, $DL$ is thus incommensurable with $MF$. Thus,
$DM$ (is) also incommensurable (in length) with $MG$
[Props.~6.1, 10.11]. 
$DM$ and $MG$ are thus rational (straight-lines which are) commensurable
in square only. Thus, $DG$ is a binomial (straight-line) [Prop.~10.36]. So, I say that (it is) also a sixth
(binomial straight-line).

So, similarly (to the previous propositions), we can again show that the (rectangle contained) by $DKM$ is equal to the (square) on $MN$, and that
$DK$ is incommensurable in length with $KM$. And so, for the same (reasons), the square on $DM$ is greater than (the square on) $MG$ by the
(square) on (some straight-line) incommensurable in length with ($DM$) [Prop.~10.18]. And neither of $DM$ and $MG$
is commensurable in length with the (previously) laid down rational (straight-line) $DE$.

Thus, $DG$ is a sixth binomial (straight-line) [Def. 10.10]. (Which is) the very thing it
was required to show.}
\end{Parallel}
{\footnotesize\noindent $^\dag$ In other words, the square of the square-root of  two  medials  is a
sixth binomial. See Prop.~10.59.\\}

%%%%
%10.66
%%%%
\pdfbookmark[1]{Proposition 10.66}{pdf10.66}
\begin{Parallel}{}{}
\ParallelLText{
\begin{center}
{\large \ggn{66}.}
\end{center}\vspace*{-7pt}

{\greekfont Ἡ τῇ ἐκ δύο ὀνομάτων μήκει σύμμετρος καὶ α ὐτὴ ἐκ δύο
ὀνομάτων ἐστὶ καὶ τῇ τάχει ἡ α ὐτή.}

{\greekfont Ἔστω ἐκ δύο ὀνομάτων ἡ ΑΒ, καὶ τῇ ΑΒ μήκει σύμμετρος
ἔστω ἡ ΓΔ; λέγω, ὅτι ἡ ΓΔ ἐκ δύο ὀνομάτων ἐστὶ καὶ τῇ
τάχει ἡ α ὐτὴ τῇ ΑΒ.}

\epsfysize=0.7in
\centerline{\epsffile{Book10/fig066g.eps}}

{\greekfont Ἐπεὶ γὰρ ἐκ δύο ὀνομάτων ἐστὶν ἡ ΑΒ, διῃρήσθω εἰς τὰ
ὀνόματα κατὰ τὸ Ε, καὶ ἔστω μεῖζον ὄνομα τὸ ΑΕ; αἱ ΑΕ, ΕΒ ἄρα
ῥηταί εἰσι δυνάμει μόνον σύμμετροι. γεγονέτω ὡς ἡ ΑΒ πρὸς
τὴν ΓΔ, οὕτως ἡ ΑΕ πρὸς τὴν ΓΖ; καὶ λοιπὴ ἄρα ἡ ΕΒ πρὸς
λοιπὴν τὴν ΖΔ ἐστιν, ὡς ἡ ΑΒ πρὸς τὴν ΓΔ. σύμμετρος δὲ ἡ ΑΒ τῇ ΓΔ μήκει; σύμμετρος ἄρα ἐστὶ καὶ ἡ μὲν ΑΕ τῇ ΓΖ, ἡ
δὲ ΕΒ τῇ ΖΔ. καί εἰσι ῥηταὶ αἱ ΑΕ, ΕΒ; ῥηταὶ ἄρα εἰσὶ καὶ αἱ
ΓΖ, ΖΔ. καὶ  ἐστιν ὡς ἡ ΑΕ πρὸς ΓΖ, ἡ ΕΒ πρὸς ΖΔ. ἐναλλὰχ
ἄρα ἐστὶν ὡς ἡ ΑΕ πρὸς ΕΒ, ἡ ΓΖ πρὸς ΖΔ.
αἱ δὲ
ΑΕ, ΕΒ δυνάμει μόνον [εἰσὶ] σύμμετροι; καὶ αἱ ΓΖ, ΖΔ ἄρα
δυνάμει μόνον  εἰσὶ σύμμετροι.
καί εἰσι ῥηταί;
ἐκ δύο ἄρα ὀνομάτων ἐστὶν ἡ ΓΔ. λέγω δή, ὅτι τῇ τάχει
ἐστὶν ἡ α ὐτὴ τῇ ΑΒ.}

{\greekfont Ἡ γὰρ ΑΕ τῆς ΕΒ μεῖζον δύναται ἤτοι τῷ ἀπὸ συμμέτρου
ἑαυτῇ ἢ τῷ ἀπὸ ἀσυμμέτρου. εἰ μὲν οὖν ἡ ΑΕ τῆς
ΕΒ μεῖζον δύναται τῷ ἀπὸ συμμέτρου ἑαυτῇ, καὶ ἡ ΓΖ τῆς
ΖΔ μεῖζον δυνήσεται τῷ ἀπὸ συμμέτρου ἑαυτῇ. καὶ εἰ μὲν
σύμμετρός ἐστιν ἡ ΑΕ τῇ ἐκκειμένῃ ῥητῇ, καὶ ἡ ΓΖ σύμμετρος
α ὐτῇ ἔσται, καὶ διὰ τοῦτο ἑκατέρα τῶν ΑΒ, ΓΔ ἐκ δύο ὀνομάτων
ἐστὶ πρώτη, τουτέστι τῇ τάχει ἡ α ὐτή. εἰ δὲ ἡ ΕΒ σύμμετρός
ἐστι τῇ ἐκκειμένῃ ῥητῇ, καὶ ἡ ΖΔ σύμμετρός ἐστιν α ὐτῇ, καὶ
διὰ τοῦτο πάλιν τῇ τάχει ἡ α ὐτὴ ἔσται τῇ ΑΒ; ἑκατέρα γὰρ
α ὐτῶν ἔσται ἐκ δύο ὀνομάτων δευτέρα. εἰ δὲ οὐδετέρα
τῶν ΑΕ, ΕΒ σύμμετρός ἐστι τῇ ἐκκειμένῃ ῥητῇ, οὐδετέρα
τῶν ΓΖ, ΖΔ σύμμετρος α ὐτῇ ἔσται, καί ἐστιν ἑκατέρα τρίτη. εἰ δὲ
ἡ ΑΕ τῆς ΕΒ μεῖζον δύναται τῷ ἀπὸ ἀσυμμέτρου ἑαυτῇ,
καὶ ἡ ΓΖ τὴς ΖΔ μεῖζον δύναται τῷ ἀπὸ ἀσυμμέτρου
ἑαυτῇ. καὶ εἰ μὲν ἡ ΑΕ σύμμετρός ἐστι τῇ ἐκκειμένῃ ῥητῃ,
καὶ ἡ ΓΖ σύμμετρός ἐστιν α ὐτῇ, καὶ ἐστιν ἑκατέρα τετάρτη.
εἰ δὲ ἡ ΕΒ, καὶ ἡ ΖΔ, καὶ ἔσται ἑκατέρα πέμπτη. εἰ δὲ οὐδετέρα
τῶν ΑΕ, ΕΒ, καὶ τῶν ΓΖ, ΖΔ οὐδετέρα σύμμετρός ἐστι τῇ ἐκκειμένῃ ῥητῇ, καὶ ἔσται ἑκατέρα ἕκτη.}

{\greekfont Ὥστε ἡ τῇ ἐκ δύο ὀνομάτων μήκει σύμμετρος ἐκ δύο
ὀνομάτων ἐστὶ καὶ τῇ τάχει ἡ α ὐτή; ὅπερ ἔδει δεῖχαι.
}}

\ParallelRText{
\begin{center}
{\large Proposition 66}
\end{center}

A (straight-line) commensurable in length with
a binomial (straight-line) is itself also binomial, and the same in order.

Let $AB$ be a binomial (straight-line), and let $CD$ be commensurable
in length with $AB$. I say that $CD$ is a binomial (straight-line),
and (is)  the same in order as $AB$.

\epsfysize=0.7in 
\centerline{\epsffile{Book10/fig066e.eps}}

For since $AB$ is a binomial (straight-line), let it be divided into
its (component) terms at $E$, and let $AE$ be the greater term. $AE$ and
$EB$ are thus rational (straight-lines which are) commensurable in square only [Prop.~10.36]. Let it be contrived
that as $AB$ (is) to $CD$, so $AE$ (is) to $CF$ [Prop.~6.12]. Thus, the remainder $EB$ is also
to the remainder $FD$, as $AB$ (is) to $CD$ [Props.~6.16, 5.19~corr.]. 
And $AB$ (is) commensurable in length with $CD$. Thus, $AE$ is also
commensurable (in length) with $CF$, and $EB$ with $FD$ [Prop.~10.11]. And $AE$ and $EB$ are rational.
Thus, $CF$ and $FD$ are also rational. And  as $AE$
is to $CF$, (so) $EB$ (is) to $FD$ [Prop.~5.11]. 
Thus, alternately, as $AE$ is to $EB$, (so) $CF$ (is) to $FD$ [Prop.~5.16]. And $AE$ and $EB$ [are]
commensurable in square only. Thus, $CF$ and $FD$ are also
commensurable in square only [Prop.~10.11]. 
And they are rational. $CD$ is thus a binomial (straight-line) [Prop.~10.36]. So, I say that it is the same in order as $AB$.

For the square on $AE$ is  greater than (the square on) $EB$ by the
(square) on (some straight-line) either commensurable or incommensurable
(in length) with ($AE$).  Therefore, if the square on $AE$ is greater than (the square on)
$EB$ by the (square) on (some straight-line) commensurable (in length)
with ($AE$) then the square on $CF$ will also be greater than (the square on)
$FD$ by the (square) on (some straight-line) commensurable (in length)
with ($CF$) [Prop.~10.14]. And if $AE$
is commensurable (in length) with (some previously) laid down rational (straight-line) then $CF$ will also be commensurable (in length) with it [Prop.~10.12]. And, on account of this,
$AB$ and $CD$ are each first binomial (straight-lines) [Def.~10.5]---that is to say, the same  in order.
And if $EB$ is commensurable (in length) with the (previously) laid down
rational (straight-line) then $FD$ is also  commensurable (in length) with it [Prop.~10.12], and, again, on account of this,
($CD$) will be  the same in order as $AB$. For each of them will be second
binomial (straight-lines) [Def.~10.6]. And if neither
of $AE$ and $EB$ is commensurable (in length) with the (previously) laid down
rational (straight-line) then neither of $CF$ and $FD$ will be commensurable (in length)
with it [Prop.~10.13], and each (of $AB$ and $CD$) is a third (binomial straight-line) [Def.~10.7]. And if the square on $AE$
is greater than (the square on) $EB$ by the (square) on (some straight-line)
incommensurable (in length) with ($AE$) then the square on $CF$
is also greater than (the square on) $FD$ by the (square) on (some straight-line) incommensurable (in length) with ($CF$) [Prop.~10.14]. And if $AE$ is commensurable
(in length) with the (previously) laid down rational (straight-line) then $CF$ is also
commensurable (in  length) with it [Prop.~10.12], 
and each (of $AB$ and $CD$) is a fourth (binomial straight-line)
[Def.~10.8]. And
if $EB$ (is commensurable in length with the previously laid down rational
straight-line) then $FD$ (is) also (commensurable in length with it), and
each (of $AB$ and $CD$) will be a fifth (binomial straight-line) [Def.~10.9]. And if neither of $AE$ and  $EB$
(is commensurable in length with the previously laid down rational
straight-line) then also neither of $CF$ and $FD$ is   commensurable (in length) with the laid down rational (straight-line), and
each (of $AB$ and $CD$) will be a sixth (binomial straight-line) [Def.~10.10].

Hence, a (straight-line) commensurable in length with a binomial (straight-line) is a binomial (straight-line), and the same in order. (Which is) the
very thing it was required to show.}
\end{Parallel}

%%%%
%10.67
%%%%
\pdfbookmark[1]{Proposition 10.67}{pdf10.67}
\begin{Parallel}{}{}
\ParallelLText{
\begin{center}
{\large\ggn{67}.}
\end{center}\vspace*{-7pt}

{\greekfont Ἡ τῇ ἐκ δύο μέσων μήκει σύμμετρος καὶ α ὐτὴ ἐκ δύο μέσων
ἐστὶ καὶ τῇ τάχει ἡ α ὐτή.}

\epsfysize=0.7in
\centerline{\epsffile{Book10/fig066g.eps}}

{\greekfont Ἔστω ἐκ δύο μέσων ἡ ΑΒ, καὶ τῇ ΑΒ σύμμετρος ἔστω μήκει
ἡ ΓΔ; λέγω, ὅτι ἡ ΓΔ ἐκ δύο μέσων ἐστὶ καὶ τῇ τάχει
ἡ α ὐτὴ τῇ ΑΒ.}

{\greekfont Ἐπεὶ γὰρ ἐκ δύο μέσων ἐστὶν ἡ ΑΒ, διῃρήσθω εἰς τὰς
μέσας κατὰ τὸ Ε; αἱ ΑΕ, ΕΒ ἄρα μέσαι εἰσὶ δυνάμει μόνον
σύμμετροι. καὶ γεγονέτω ὡς ἡ ΑΒ πρὸς ΓΔ, ἡ ΑΕ πρὸς ΓΖ;
καὶ λοιπὴ ἄρα ἡ ΕΒ πρὸς λοιπὴν τὴν ΖΔ ἐστιν, ὡς ἡ ΑΒ πρὸς ΓΔ. σύμμετρος δὲ ἡ ΑΒ τῇ ΓΔ μήκει; σύμμετρος ἄρα καὶ ἑκατέρα τῶν
ΑΕ, ΕΒ ἑκατέρᾳ τῶν ΓΖ, ΖΔ. μέσαι δὲ αἱ ΑΕ, ΕΒ; μέσαι ἄρα καὶ αἱ
ΓΖ, ΖΔ. καὶ ἐπεί ἐστιν ὡς ἡ ΑΕ πρὸς ΕΒ, ἡ ΓΖ πρὸς ΖΔ, αἱ 
δὲ ΑΕ, ΕΒ δυνάμει μόνον σύμμετροί εἰσιν, καὶ αἱ ΓΖ, ΖΔ [ἄρα]
δυνάμει μόνον σύμμετροί εἰσιν, ἐδείθθησαν δὲ καὶ μέσαι; ἡ
ΓΔ ἄρα ἐκ δύο μέσων ἐστίν. λέγω δή, ὅτι καὶ τῇ τάχει ἡ
α ὐτή ἐστι τῇ ΑΒ.}

{\greekfont Ἐπεὶ γάρ ἐστιν ὡς ἡ ΑΕ πρὸς ΕΒ, ἡ ΓΖ πρὸς ΖΔ, καὶ ὡς
ἄρα τὸ ἀπὸ τῆς ΑΕ πρὸς τὸ ὑπὸ τῶν ΑΕΒ, οὕτως
τὸ ἀπὸ τῆς ΓΖ πρὸς τὸ ὑπὸ τῶν ΓΖΔ; ἐναλλὰχ ὡς τὸ
ἀπὸ τῆς ΑΕ πρὸς τὸ ἀπὸ τῆς ΓΖ, οὕτως τὸ ὑπὸ τῶν ΑΕΒ
πρὸς τὸ ὑπὸ τῶν ΓΖΔ. σύμμετρον δὲ τὸ ἀπὸ τῆς
ΑΕ τῷ ἀπὸ τῆς ΓΖ; σύμμετρον
ἄρα καὶ τὸ ὑπὸ τῶν ΑΕΒ τῷ ὑπὸ τῶν ΓΖΔ. εἴτε οὖν ῥητόν
ἐστι τὸ ὑπὸ τῶν ΑΕΒ, καὶ τὸ ὑπὸ τῶν ΓΖΔ ῥητόν ἐστιν
[καὶ διὰ τοῦτό ἐστιν ἐκ δύο μέσων πρώτη]. εἴτε μέσον, μέσον,
καί ἐστιν ἑκατέρα δευτέρα.}

{\greekfont Καὶ διὰ τοῦτο ἔσται ἡ ΓΔ τῇ ΑΒ τῇ τάχει ἡ
α ὐτή; ὅπερ ἔδει δεῖχαι.}}

\ParallelRText{
\begin{center}
{\large Proposition 67}
\end{center}

A (straight-line) commensurable
in length with a bimedial (straight-line) is  itself also bimedial, and
 the same in order.

\epsfysize=0.7in 
\centerline{\epsffile{Book10/fig066e.eps}}
 
Let $AB$ be a bimedial (straight-line), and let $CD$ be commensurable
in length with $AB$. I say that $CD$ is bimedial, and the same in order as
$AB$.

For since $AB$ is a bimedial (straight-line), let it be divided into its
(component) medial (straight-lines) at $E$. Thus, $AE$ and $EB$
are medial (straight-lines which are) commensurable in square only [Props.~10.37, 10.38]. And
let it be contrived that as $AB$ (is) to $CD$, (so) $AE$ (is)
to $CF$ [Prop.~6.12]. And thus as the remainder $EB$ is to the remainder $FD$, so $AB$ (is) to $CD$ [Props.~5.19~corr., 6.16].
And $AB$ (is) commensurable in length with $CD$. Thus, $AE$
and $EB$ are also commensurable (in length) with $CF$ and $FD$, respectively [Prop.~10.11]. And 
$AE$ and $EB$ (are) medial. Thus, $CF$ and $FD$ (are) also
medial [Prop.~10.23]. And since as $AE$
is to $EB$, (so) $CF$ (is) to $FD$, and $AE$ and $EB$ are commensurable
in square only, $CF$ and $FD$ are [thus] also commensurable in square only
[Prop.~10.11]. And they were also shown (to be)
medial. Thus, $CD$ is a bimedial (straight-line). So, I say that it is also the same in order as $AB$.

For since as $AE$ is to $EB$, (so) $CF$ (is) to $FD$, thus also as the
(square) on $AE$ (is) to the (rectangle contained) by $AEB$, so the (square)
on $CF$ (is) to the (rectangle contained) by $CFD$ [Prop.~10.21~lem.]. Alternately, as the (square)
on $AE$ (is) to the (square) on $CF$, so the (rectangle
contained) by $AEB$ (is) to the (rectangle contained) by $CFD$ [Prop.~5.16]. And the (square) on $AE$
(is) commensurable with the (square) on $CF$. Thus, the (rectangle
contained) by $AEB$ (is) also commensurable with the (rectangle contained)
by $CFD$ [Prop.~10.11]. Therefore,  either the
(rectangle contained) by $AEB$ is rational, and  the (rectangle contained) by
$CFD$ is rational [and, on account of this, ($AE$ and $CD$) are first
bimedial (straight-lines)], or (the rectangle contained by $AEB$ is) medial,
and (the rectangle contained by $CFD$ is) medial, and ($AB$ and
$CD$) are each second (bimedial straight-lines) [Props.~10.23, 10.37, 10.38].

And, on account of this, $CD$ will be  the same in order as $AB$.
(Which is) the very thing it was required to show.}
\end{Parallel}

%%%%
%10.68
%%%%
\pdfbookmark[1]{Proposition 10.68}{pdf10.68}
\begin{Parallel}{}{}
\ParallelLText{
\begin{center}
{\large\ggn{68}.}
\end{center}\vspace*{-7pt}

{\greekfont Ἡ τῇ μείζονι σύμμετρος καὶ α ὐτὴ μείζων ἐστίν.}\\

\epsfysize=0.7in
\centerline{\epsffile{Book10/fig066g.eps}}

{\greekfont Ἔστω μείζων ἡ ΑΒ, καὶ τῇ ΑΒ σύμμετρος ἔστω ἡ ΓΔ; λέγω,
ὅτι ἡ ΓΔ μείζων ἐστίν.}

{\greekfont Διῃρήσθω ἡ ΑΒ κατὰ τὸ Ε; αἱ ΑΕ, ΕΒ ἄρα δυνάμει εἰσὶν ἀσύμμετροι ποιοῦσαι τὸ μὲν συγκείμενον ἐκ τῶν ἀπ α ὐτῶν
τετραγώνων ῥητόν, τὸ δ ὑπ α ὐτῶν μέσον; καὶ γεγονέτω
τὰ α ὐτὰ τοῖς πρότερον. καὶ ἐπεί ἐστιν ὡς ἡ ΑΒ πρὸς τὴν
ΓΔ, οὕτως ἥ τε ΑΕ πρὸς τὴν ΓΖ καὶ ἡ ΕΒ πρὸς τὴν ΖΔ, καὶ ὡς ἄρα ἡ ΑΕ πρὸς τὴν ΓΖ, οὕτως ἡ ΕΒ πρὸς τὴν ΖΔ. σύμμετρος
δὲ ἡ ΑΒ τῇ ΓΔ; σύμμετρος ἄρα καὶ ἑκατέρα τῶν ΑΕ, ΕΒ ἑκατέρᾳ
τῶν ΓΖ, ΖΔ. καὶ ἐπεί ἐστιν ὡς ἡ ΑΕ πρὸς τὴν ΓΖ, οὕτως
ἡ ΕΒ πρὸς τὴν ΖΔ,  καὶ
ἐναλλὰχ ὡς ἡ ΑΕ πρὸς ΕΒ, οὕτως ἡ ΓΖ πρὸς ΖΔ, καὶ
συνθέντι ἄρα ἑστὶν ὡς ἡ ΑΒ πρὸς τὴν ΒΕ, οὕτως  ἡ ΓΔ πρὸς
τὴν ΔΖ; καὶ ὡς ἄρα τὸ ἀπὸ τῆς ΑΒ πρὸς τὸ ἀπὸ τῆς ΒΕ, οὕτως
τὸ ἀπὸ τῆς ΓΔ πρὸς τὸ ἀπὸ τῆς ΔΖ. ὁμοίως δὴ δείχομεν,
ὅτι καὶ ὡς τὸ ἀπὸ τῆς ΑΒ πρὸς τὸ ἀπὸ τῆς ΑΕ, οὕτως
τὸ ἀπὸ τῆς ΓΔ πρὸς τὸ ἀπὸ τῆς ΓΖ. καὶ ὡς ἄρα τὸ ἀπὸ
τῆς ΑΒ πρὸς τὰ ἀπὸ τῶν ΑΕ, ΕΒ, οὕτως τὸ ἀπὸ τῆς ΓΔ
πρὸς τὰ ἀπὸ τῶν ΓΖ, ΖΔ; καὶ ἐναλλὰχ ἄρα ἐστὶν ὡς τὸ
ἀπὸ τῆς ΑΒ πρὸς τὸ ἀπὸ τῆς ΓΔ, οὕτως τὰ ἀπὸ τῶν
ΑΕ, ΕΒ πρὸς τὰ ἀπὸ τῶν ΓΖ, ΖΔ. σύμμετρον δὲ τὸ ἀπὸ
τῆς ΑΒ τῷ ἀπὸ τῆς ΓΔ; σύμμετρα ἄρα καὶ τὰ ἀπὸ τῶν ΑΕ, ΕΒ
τοῖς ἀπὸ τῶν ΓΖ, ΖΔ. καί ἐστι τὰ ἀπὸ τῶν ΑΕ, ΕΒ ἅμα
ῥητόν, καὶ τὰ ἀπὸ τῶν ΓΖ, ΖΔ ἅμα ῥητόν ἐστιν. ὁμοίως
δὲ καὶ τὸ δὶς ὑπὸ τῶν ΑΕ, ΕΒ σύμμετρόν ἐστι τῷ δὶς
ὑπὸ τῶν
ΓΖ, ΖΔ. καί ἐστι μέσον τὸ δὶς ὑπὸ τῶν ΑΕ, ΕΒ;
μέσον ἄρα καὶ τὸ δὶς ὑπὸ τῶν ΓΖ, ΖΔ. αἱ ΓΖ, ΖΔ
ἄρα δυνάμει ἀσύμμετροί
εἰσι ποιοῦσαι τὸ μὲν συγκείμενον ἐκ τῶν ἀπ α ὐτῶν
τετραγώνων ἅμα ῥητόν, τὸ δὲ δὶς ὑπ α ὐτῶν μέσον;
ὅλη ἄρα ἡ ΓΔ ἄλογός ἐστιν ἡ καλουμένη μείζων.}

{\greekfont Ἡ ἄρα τῇ μείζονι σύμμετρος μείζων ἐστίν; ὅπερ
ἔδει δεῖχαι.}}

\ParallelRText{
\begin{center}
{\large Proposition 68}
\end{center}

A (straight-line) commensurable (in length)
with a major (straight-line) is  itself also major.

\epsfysize=0.7in 
\centerline{\epsffile{Book10/fig066e.eps}}

Let $AB$ be a major (straight-line), and let $CD$ be commensurable (in
length) with $AB$. I say that $CD$ is a major (straight-line).

Let $AB$ be divided (into its component terms) at $E$.
$AE$ and $EB$ are thus incommensurable in square, making the sum
of the squares on them rational, and the (rectangle contained) by them
medial [Prop.~10.39]. And let (the) same (things) be contrived as in the previous (propositions). And since
as $AB$ is to $CD$, so $AE$ (is) to $CF$ and $EB$ to $FD$, 
thus also as $AE$ (is) to $CF$, so $EB$ (is) to $FD$ [Prop.~5.11]. And $AB$ (is) commensurable
(in length) with $CD$. Thus, $AE$ and $EB$ (are) also commensurable
(in length) with $CF$ and $FD$, respectively [Prop.~10.11].  And since as $AE$ is to $CF$, so
$EB$ (is) to $FD$, also, alternately, as $AE$ (is) to $EB$, so
$CF$ (is) to $FD$ [Prop.~5.16], and thus, via composition, as $AB$ is to $BE$, so $CD$ (is) to $DF$ [Prop.~5.18]. And thus as the (square) on $AB$
(is) to the (square) on $BE$, so the (square) on $CD$ (is) to the (square)
on $DF$ [Prop.~6.20]. So, similarly, we can also show that as the (square) on $AB$ (is) to the (square) on $AE$, so the
(square) on $CD$ (is) to the (square) on $CF$. And thus as the (square)
on $AB$ (is) to (the sum of) the (squares) on $AE$ and $EB$, so
the (square) on $CD$ (is) to (the sum of) the (squares) on $CF$ and $FD$.
And thus, alternately, as the (square) on $AB$ is to the (square) on $CD$,
so (the sum of) the (squares) on $AE$ and $EB$ (is) to (the sum of)
the (squares) on $CF$ and $FD$ [Prop.~5.16]. And the (square) on $AB$
(is) commensurable with the (square) on $CD$. Thus, (the sum of) the
(squares) on $AE$ and $EB$ (is) also commensurable with (the sum of)
the (squares) on $CF$ and $FD$ [Prop.~10.11]. 
And the (squares) on $AE$ and $EB$ (added) together are rational.
The (squares) on $CF$ and $FD$ (added) together (are) thus also
rational. So, similarly, twice the (rectangle contained) by $AE$ and
$EB$ is also commensurable with twice the (rectangle contained)
by $CF$ and $FD$. And twice the (rectangle contained) by $AE$ and $EB$
is medial. Therefore, twice the (rectangle contained) by $CF$ and
$FD$ (is) also medial [Prop.~10.23~corr.]. 
$CF$ and $FD$ are thus (straight-lines which are) incommensurable in square [Prop~10.13], simultaneously making the sum of the
squares on them rational, and twice the (rectangle contained) by them medial.
The whole, $CD$, is thus that irrational (straight-line) called major
[Prop.~10.39].

Thus, a (straight-line) commensurable (in length) with a major
(straight-line) is major. (Which is) the very thing it was required to show.}
\end{Parallel}

%%%%
%10.69
%%%%
\pdfbookmark[1]{Proposition 10.69}{pdf10.69}
\begin{Parallel}{}{}
\ParallelLText{
\begin{center}
{\large \ggn{69}.}
\end{center}\vspace*{-7pt}

{\greekfont Ἡ τῇ ῥητὸν καὶ μέσον δυναμένῃ σύμμετρος [καὶ α ὐτὴ]
ῥητὸν καὶ μέσον δυναμένη ἐστίν.}\\

\epsfysize=0.7in
\centerline{\epsffile{Book10/fig066g.eps}}

{\greekfont Ἔστω ῥητὸν καὶ μέσον δυναμένη ἡ ΑΒ, καὶ τῇ
ΑΒ σύμμετρος ἔστω ἡ ΓΔ; δεικτέον, ὅτι καὶ ἡ ΓΔ
ῥητὸν καὶ μέσον δυναμένη ἐστίν.}

{\greekfont Διῃρήσθω ἡ ΑΒ εἰς τὰς εὐθείας κατὰ τὸ Ε; αἱ ΑΕ, ΕΒ ἄρα
δυνάμει εἰσὶν ἀσύμμετροι ποιοῦσαι
τὸ μὲν συγκείμενον ἐκ τῶν ἀπ α ὐτῶν τετραγώνων
μέσον, τὸ δ ὑπ α ὐτῶν ῥητόν; καὶ τὰ α ὐτὰ κατεσκευάσθω
τοῖς πρότερον. ὁμοίως δὴ δείχομεν, ὅτι καὶ αἱ ΓΖ,
ΖΔ δυνάμει εἰσὶν ἀσύμμετροι, καὶ σύμμετρον τὸ μὲν
συγκείμενον ἐκ τῶν ἀπὸ τῶν ΑΕ, ΕΒ τῷ συγκειμένῳ
ἐκ τῶν ἀπὸ τῶν ΓΖ, ΖΔ, τὸ δὲ ὑπὸ ΑΕ, ΕΒ τῷ
ὑπὸ ΓΖ, ΖΔ; ὥστε καὶ τὸ [μὲν] συγκείμενον ἐκ τῶν
ἀπὸ τῶν ΓΖ, ΖΔ τετραγώνων ἐστὶ μέσον, τὸ δ ὑπὸ
τῶν ΓΖ, ΖΔ ῥητόν.}

{\greekfont Ῥητὸν ἄρα καὶ μέσον δυναμένη ἐστὶν ἡ ΓΔ; ὅπερ
ἔδει δεῖχαι.}}

\ParallelRText{
\begin{center}
{\large Proposition 69}
\end{center}

A (straight-line) commensurable (in length)
with the square-root of a rational plus a medial (area) is [itself also]
the square-root of a rational plus a medial (area).

\epsfysize=0.7in 
\centerline{\epsffile{Book10/fig066e.eps}}

Let $AB$ be the square-root of a rational plus a medial (area), 
and let $CD$ be commensurable (in length) with $AB$. We must  show
that $CD$ is also the square-root of a rational plus a medial (area).

Let $AB$ be divided into its (component) straight-lines at $E$.
$AE$ and $EB$ are thus incommensurable in square, making 
the sum of the squares on them medial, and the (rectangle contained)
by them rational [Prop.~10.40]. And let the
same construction be made as in the previous (propositions).
So, similarly, we can show that $CF$ and $FD$ are also incommensurable
in square, and that the sum of the (squares) on $AE$ and $EB$
(is) commensurable with the sum of the (squares) on $CF$ and $FD$,
and the (rectangle contained) by $AE$ and $EB$ with the
(rectangle contained) by $CF$ and $FD$. And hence the sum of the
squares on $CF$ and $FD$ is medial, and the (rectangle contained) by
$CF$ and $FD$ (is) rational.

Thus, $CD$ is the square-root of a rational plus a medial (area) [Prop.~10.40]. (Which is) the very thing it was required to show.}
\end{Parallel}

%%%%
%10.70
%%%%
\pdfbookmark[1]{Proposition 10.70}{pdf10.70}
\begin{Parallel}{}{}
\ParallelLText{
\begin{center}
{\large \ggn{70}.}
\end{center}\vspace*{-7pt}

{\greekfont Ἡ τῇ δύο μέσα δυναμένῃ σύμμετρος δύο μέσα δυναμένη ἐστίν.}


{\greekfont Ἔστω δύο μέσα δυναμένη ἡ ΑΒ, καὶ τῇ ΑΒ σύμμετρος
ἡ ΓΔ; δεικτέον, ὅτι καὶ ἡ ΓΔ δύο μέσα δυναμένη ἐστίν.}\\~\\~\\

\epsfysize=0.7in
\centerline{\epsffile{Book10/fig066g.eps}}
{\greekfont Ἐπεὶ γὰρ δύο μέσα δυναμένη ἐστὶν ἡ ΑΒ, διῃρήσθω
εἰς τὰς εὐθείας κατὰ τὸ Ε; αἱ ΑΕ, ΕΒ ἄρα δυνάμει εἰσὶν
ἀσύμμετροι ποιοῦσαι τό τε συγκείμενον ἐκ τῶν ἀπ α ὐτῶν
[τετραγώνων] μέσον καὶ τὸ ὑπ α ὐτῶν μέσον καὶ ἔτι
ἀσύμμετρον τὸ συγκείμενον ἐκ τῶν ἀπὸ τῶν ΑΕ, ΕΒ
τετραγώνων τῷ ὑπὸ τῶν ΑΕ, ΕΒ; καὶ κατεσκευάσθω τὰ α ὐτὰ
τοῖς πρότερον. ὁμοίως δὴ δείχομεν, ὅτι καὶ αἱ ΓΖ, ΖΔ δυνάμει
εἰσὶν ἀσύμμετροι καὶ σύμμετρον τὸ μὲν συγκείμενον ἐκ
τῶν ἀπὸ τῶν ΑΕ, ΕΒ τῷ συγκειμένῳ ἐκ τῶν ἀπὸ τῶν
ΓΖ, ΖΔ, τὸ δὲ ὑπὸ τῶν ΑΕ, ΕΒ τῷ ὑπὸ τῶν ΓΖ, ΖΔ;
ὥστε καὶ τὸ συγκείμενον ἐκ τῶν ἀπὸ τῶν ΓΖ, ΖΔ
 τετραγώνων
μέσον ἐστὶ καὶ τὸ ὑπὸ τῶν ΓΖ, ΖΔ μέσον καὶ ἔτι ἀσύμμετρον
τὸ συγκείμενον ἐκ τῶν ἀπὸ τῶν ΓΖ, ΖΔ τετραγώνων τῷ ὑπὸ
τῶν ΓΖ, ΖΔ.}

{\greekfont Ἡ ἄρα ΓΔ δύο μέσα δυναμένη ἐστίν; ὅπερ ἔδει δεῖχαι.}}

\ParallelRText{
\begin{center}
{\large Proposition 70}
\end{center}

A (straight-line) commensurable (in length)
with the square-root of (the sum of) two medial (areas) is (itself also)
the square-root of (the sum of) two medial (areas).

Let $AB$ be the square-root of (the sum of) two medial (areas),
and (let) $CD$ (be) commensurable (in length) with $AB$. We must show that
$CD$ is also the square-root of (the sum of) two medial (areas).

\epsfysize=0.7in 
\centerline{\epsffile{Book10/fig066e.eps}}

For since $AB$ is the square-root of (the sum of) two medial (areas),
let it be divided into its (component) straight-lines at $E$.
Thus, $AE$ and $EB$ are incommensurable in square, making the sum of
the [squares] on them medial, and the (rectangle contained) by them
medial, and, moreover, the sum of the (squares) on $AE$ and $EB$
incommensurable with the (rectangle) contained by $AE$ and $EB$ [Prop.~10.41]. And let the same construction
be made as in the previous (propositions). So, similarly, we
can show that $CF$ and $FD$ are also incommensurable in square,
and (that) the sum of the (squares) on $AE$ and $EB$ (is)
commensurable with the sum of the (squares) on $CF$ and $FD$, and the
(rectangle contained) by $AE$ and $EB$ with the (rectangle contained)
by $CF$ and $FD$. Hence, the sum of the squares on $CF$ and
$FD$ is also medial, and the (rectangle contained) by $CF$ and $FD$
(is) medial, and, moreover, the sum of the squares on $CF$ and $FD$
(is) incommensurable with the (rectangle contained) by $CF$ and $FD$.

Thus, $CD$ is the square-root of (the sum of) two medial (areas) [Prop.~10.41].
(Which is) the very thing it was required to show.}
\end{Parallel}

%%%%
%10.71
%%%%
\pdfbookmark[1]{Proposition 10.71}{pdf10.71}
\begin{Parallel}{}{}
\ParallelLText{
\begin{center}
{\large \ggn{71}.}
\end{center}\vspace*{-7pt}

{\greekfont Ῥητοῦ καὶ μέσου συντιθεμένου τέσσαρες ἄλογοι γίγνον-ται
ἤτοι ἐκ δύο ὀνομάτων ἢ ἐκ δύο μέσων πρώτη ἢ
μείζων ἢ ῥητὸν καὶ μέσον δυναμένη.}

{\greekfont Ἔστω ῥητὸν μὲν τὸ ΑΒ, μέσον δὲ τὸ ΓΔ; λέγω, ὅτι ἡ τὸ
ΑΔ θωρίον δυναμένη ἤτοι ἐκ δύο ὀνομάτων ἐστὶν
ἢ ἐκ δύο μέσων πρώτη ἢ μείζων ἢ ῥητὸν καὶ
μέσον δυναμένη.}

{\greekfont Τὸ γὰρ ΑΒ τοῦ ΓΔ ἤτοι μεῖζόν ἐστιν ἢ ἔλασσον.
ἔστω πρότερον μεῖζον; καὶ ἐκκείσθω ῥητὴ ἡ ΕΖ, καὶ
παραβεβλήσθω παρὰ τὴν ΕΖ τῷ ΑΒ ἴσον τὸ ΕΗ
πλάτος ποιοῦν τὴν ΕΘ; τῷ δὲ ΔΓ ἴσον παρὰ τὴν ΕΖ παραβεβλήσθω
τὸ ΘΙ πλάτος ποιοῦν τὴν ΘΚ. καὶ ἐπεὶ ῥητόν ἐστι τὸ ΑΒ καί
ἐστιν ἴσον τῷ ΕΗ, ῥητὸν ἄρα καὶ τὸ ΕΗ. καὶ παρὰ [ῥητὴν]
τὴν ΕΖ παραβέβληται πλάτος ποιοῦν τὴν ΕΘ; ἡ ΕΘ ἄρα ῥητή ἐστι
καὶ σύμμετρος τῇ ΕΖ μήκει. πάλιν, ἐπεὶ μέσον ἐστὶ τὸ ΓΔ καί
ἐστιν ἴσον τῷ ΘΙ, μέσον ἄρα ἐστὶ καὶ τὸ ΘΙ. καὶ παρὰ ῥητὴν
τὴν ΕΖ παράκειται πλάτος ποιοῦν τὴν ΘΚ; ῥητὴ ἄρα ἐστὶν ἡ ΘΚ
καὶ ἀσύμμετρος τῇ ΕΖ μήκει. καὶ ἐπεὶ μέσον ἐστὶ τὸ ΓΔ,
ῥητὸν δὲ τὸ ΑΒ, ἀσύμμετρον ἄρα ἐστὶ τὸ ΑΒ τῷ ΓΔ; ὥστε
καὶ τὸ ΕΗ ἀσύμμετρόν ἐστι τῷ ΘΙ. ὡς δὲ τὸ ΕΗ πρὸς τὸ ΘΙ,
οὕτως ἐστὶν ἡ ΕΘ πρὸς τὴν ΘΚ; ἀσύμμετρος ἄρα ἐστὶ καὶ
ἡ ΕΘ τῇ ΘΚ μήκει. καί εἰσιν ἀμφότεραι ῥηταί; αἱ ΕΘ, ΘΚ
ἄρα ῥηταί εἰσι δυνάμει μόνον σύμμετροι; ἐκ δύο ἄρα
ὀνομάτων ἐστὶν ἡ ΕΚ διῃρημένη κατὰ τὸ Θ. καὶ ἐπεὶ
μεῖζόν ἐστι τὸ ΑΒ τοῦ ΓΔ, ἴσον δὲ τὸ μὲν ΑΒ τῷ ΕΗ, τὸ δὲ
ΓΔ τῷ ΘΙ, μεῖζον ἄρα καὶ τὸ ΕΗ τοῦ ΘΙ; καὶ ἡ ΕΘ
ἄρα μείζων ἐστὶ τῆς ΘΚ. ἤτοι οὖν ἡ ΕΘ τῆς ΘΚ μεῖζον
δύναται τῷ ἀπὸ συμμέτρου ἑαυτῇ μήκει ἢ τῷ
ἀπὸ ἀσυμμέτρου. δυνάσθω πρότερον τῷ ἀπὸ συμμέτρου
ἑαυτῇ; καί ἐστιν ἡ μείζων ἡ ΘΕ σύμμετρος τῇ ἐκκειμένῃ
ῥητῃ τῇ ΕΖ; ἡ ἄρα ΕΚ ἐκ δύο ὀνομάτων ἐστὶ πρώτη. ῥητὴ
δὲ ἡ ΕΖ; ἐὰν δὲ θωρίον περιέθηται ὑπὸ ῥητῆς καὶ τῆς ἐκ
δύο ὀνομάτων πρώτης, ἡ τὸ θωρίον δυναμένη ἐκ δύο ὀνομάτων
ἐστίν. ἡ ἄρα τὸ ΕΙ δυναμένη ἐκ δύο ὀνομάτων ἐστίν; ὥστε
καὶ ἡ τὸ ΑΔ δυναμένη ἐκ δύο ὀνομάτων ἐστίν. ἀλλὰ δὴ δυνάσθω
ἡ ΕΘ τῆς ΘΚ μεῖζον τῷ ἀπὸ ἀσυμμέτρου ἑαυτῇ; καί
ἐστιν ἡ μείζων ἡ ΕΘ σύμμετρος τῇ ἐκκειμένῃ ῥητῇ
τῇ ΕΖ μήκει; ἡ ἄρα ΕΚ ἐκ δύο ὀνομάτων ἐστὶ τετάρτη.
ῥητὴ δὲ ἡ ΕΖ; ἐὰν δὲ θωρίον περιέθηται ὑπὸ ῥητῆς
καὶ τῆς ἐκ δύο ὀνομάτων τετάρτης, ἡ τὸ θωρίον δυναμένη ἄλογός
ἐστιν ἡ καλουμένη μείζων. ἡ ἄρα τὸ ΕΙ θωρίον δυναμένη μείζων
ἐστίν; ὥστε καὶ ἡ τὸ ΑΔ δυναμένη μείζων ἐστίν.}

{\greekfont Ἀλλὰ δὴ ἔστω ἔλασσον τὸ ΑΒ τοῦ ΓΔ; καὶ τὸ ΕΗ ἄρα ἔλασσόν
ἐστι τοῦ ΘΙ; ὥστε καὶ ἡ ΕΘ ἐλάσσων ἐστὶ τῆς ΘΚ. ἤτοι δὲ ἡ ΘΚ
τῆς ΕΘ
 μεῖζον δύναται τῷ ἀπὸ συμμέτρου
ἑαυτῇ ἢ τῷ ἀπὸ ἀσυμμέτρου. δυνάσθω πρότερον τῷ ἀπὸ
συμμέτρου ἑαυτῇ μήκει; καί ἐστιν ἡ ἐλάσσων ἡ ΕΘ σύμμετρος
τῇ ἐκκειμένῃ ῥητῇ τῇ ΕΖ μήκει; ἡ ἄρα ΕΚ ἐκ δύο
ὀνομάτων ἐστὶ δευτέρα. ῥητὴ δὲ ἡ ΕΖ; ἐὰν δὲ θωρίον
περιέθηται ὑπὸ ῥητῆς καὶ τῆς ἐκ δύο ὀνομάτων δευτέρας,
ἡ τὸ θωρίον δυναμένη ἐκ δύο μέσων ἐστὶ πρώτη. ἡ ἄρα
τὸ ΕΙ θωρίον δυναμένη ἐκ δύο μέσων ἐστὶ πρώτη; ὥστε
καὶ ἡ τὸ ΑΔ δυναμένη ἐκ δύο μέσων ἐστὶ πρώτη. ἀλλὰ δὴ
ἡ ΘΚ τῆς ΘΕ μεῖζον δυνάσθω τῷ ἀπὸ ἀσυμμέτρου ἑαυτῇ.
καί ἐστιν ἡ ἐλάσσων ἡ ΕΘ σύμμετρος τῇ ἐκκειμένῃ
ῥητῇ τῇ ΕΖ; ἡ ἄρα ΕΚ ἐκ δύο ὀνομάτων ἐστὶ πέμπτη.
ῥητὴ δὲ ἡ ΕΖ; ἐὰν δὲ θωρίον περιέθηται ὑπὸ ῥητῆς καὶ τῆς
ἐκ δύο ὀνομάτων πέμπτης, ἡ τὸ θωρίον δυναμένη ῥητὸν καὶ μέσον δυναμένη ἐστίν. ἡ ἄρα τὸ ΕΙ θωρίον δυναμένη ῥητὸν
καὶ μέσον δυναμένη ἐστίν; ὥστε καὶ ἡ τὸ ΑΔ θωρίον δυναμένη
ῥητὸν καὶ μέσον δυναμένη ἐστίν.}\\~\\~\\~\\~\\~\\~\\~\\~\\~\\~\\~\\~\\~\\~\\~\\~\\~\\~\\~\\~\\~\\~\\~\\~\\~\\

\epsfysize=2.in
\centerline{\epsffile{Book10/fig071g.eps}}

{\greekfont Ῥητοῦ ἄρα καὶ μέσου συντιθεμένου τέσσαρες ἄλογοι γίγνονται
ἤτοι ἐκ δύο ὀνομάτων ἢ ἐκ δύο μέσων πρώτη ἢ μείζων
ἢ ῥητὸν καὶ μέσον δυναμένη; ὅπερ ἔδει δεῖχαι.}}

\ParallelRText{
\begin{center}
{\large Proposition 71}
\end{center}

When a rational and a medial (area) are added
together, four irrational (straight-lines) arise (as the square-roots of the
total area)---either a binomial, or a first bimedial, or a major, or the
square-root of a rational plus a medial (area).

Let $AB$ be a rational (area), and $CD$ a medial (area). I say that the
square-root of area $AD$ is either binomial, or first bimedial, or
major, or the square-root of a rational plus a medial (area).

For $AB$ is either greater or less than $CD$. Let it, first of all,
be greater. And let the rational (straight-line) $EF$ be laid down. And
let (the rectangle) $EG$, equal to $AB$, be applied to $EF$,
producing $EH$ as breadth. And let (the recatangle) $HI$, equal to $DC$,
be applied to $EF$, producing $HK$ as breadth. And since
$AB$ is rational, and is equal to $EG$, $EG$ is thus also rational. And it
has been applied to the [rational] (straight-line) $EF$, producing $EH$
as breadth. $EH$ is thus rational, and commensurable in length
with $EF$ [Prop.~10.20]. Again, since
$CD$ is medial, and is equal to $HI$, $HI$ is thus also medial. 
And it is applied to the rational (straight-line) $EF$, producing
$HK$ as breadth. $HK$ is thus rational, and incommensurable
in length with $EF$ [Prop.~10.22]. And since
$CD$ is medial, and $AB$ rational, $AB$ is thus incommensurable
with $CD$. Hence, $EG$ is also incommensurable with $HI$. And
as $EG$ (is) to $HI$, so $EH$ is to $HK$ [Prop.~6.1].
Thus, $EH$ is also incommensurable in length with $HK$
[Prop.~10.11]. And they are both rational. Thus,
$EH$ and $HK$ are rational (straight-lines which are) commensurable in
 square only. $EK$ is thus a binomial (straight-line), having been divided
(into its component terms) at $H$ [Prop.~10.36].
And since $AB$ is greater than $CD$, and $AB$ (is) equal to $EG$, and
$CD$ to $HI$, $EG$ (is) thus also greater than $HI$. Thus, $EH$
is also greater than $HK$ [Prop.~5.14]. Therefore,
 the square on $EH$ is greater than (the square on) $HK$
either by the (square) on (some straight-line) commensurable in length
with ($EH$), or by the (square) on (some straight-line)  incommensurable
(in length with $EH$). 
Let it, first of all, be greater by the (square) on (some straight-line)
commensurable (in length with $EH$). And the greater (of the two components of $EK$) $HE$ is commensurable
(in length) with the (previously) laid down (straight-line) $EF$. $EK$ is thus a first binomial
(straight-line) [Def.~10.5]. And $EF$ (is) rational. And if an area is contained by a rational (straight-line) and a first
binomial (straight-line) then the square-root  of the  area is a binomial
(straight-line)  [Prop.~10.54]. Thus,
the square-root of $EI$ is a binomial (straight-line). Hence the square-root of $AD$ is also a binomial (straight-line).  And, so, let the
 square on $EH$ be greater than (the square on) $HK$ by the (square) on 
 (some straight-line) incommensurable (in length) with ($EH$). And the
 greater (of the two components of $EK$) $EH$ is commensurable in length with the (previously)
 laid down rational (straight-line) $EF$. Thus, $EK$ is a fourth
 binomial (straight-line) [Def.~10.8]. And
 $EF$ (is) rational. And if an area
 is contained by a rational (straight-line) and a fourth binomial (straight-line) 
 then the square-root of the area is the irrational (straight-line) called
 major [Prop.~10.57]. Thus, the square-root
 of area $EI$ is a major (straight-line). Hence, the square-root of
 $AD$ is also major.

 And so, let $AB$ be less than $CD$. Thus, $EG$ is also less
 than $HI$. Hence, $EH$ is also less than $HK$ [Props.~6.1, 5.14].
 And the square on $HK$ is greater than (the square on) $EH$ either
 by the (square) on (some straight-line) commensurable (in length)
 with ($HK$), or by the (square) on (some straight-line) incommensurable
 (in length) with ($HK$). Let it, first of all, be greater by the square
 on (some straight-line) commensurable in length with ($HK$).
 And the lesser (of the two components of $EK$) $EH$ is  commensurable in length with the
 (previously) laid down rational (straight-line) $EF$. Thus, $EK$
 is a second binomial (straight-line) [Def.~10.6]. 
 And $EF$ (is) rational. And if an area is contained by a rational (straight-line)
 and a second binomial (straight-line)  then the square-root of the area is a
 first bimedial (straight-line) [Prop.~10.55]. 
 Thus, the square-root of area $EI$ is a first bimedial (straight-line).
 Hence, the square-root of $AD$ is also a first bimedial (straight-line).
 And so, let the square on $HK$ be greater than (the square on) $HE$
 by the (square) on (some straight-line) incommensurable (in length) with
 ($HK$). 
 And the lesser (of the two components of $EK$) $EH$ is commensurable (in length) with the (previously) laid down rational (straight-line) $EF$.  Thus, $EK$ is a fifth binomial
 (straight-line) [Def.~10.9]. And $EF$ (is) rational. And if an area is
 contained by a rational (straight-line) and a fifth binomial (straight-line) 
 then the square-root of the area is the square-root of a rational plus
 a medial (area) [Prop.~10.58]. Thus, the square-root of area $EI$ is the square-root of a rational plus a medial (area). Hence,
 the square-root of area $AD$ is also the square-root of a rational plus
 a medial (area).
 
\epsfysize=2.in 
\centerline{\epsffile{Book10/fig071e.eps}}

 Thus, when a rational and a medial area are added together, four irrational
(straight-lines) arise (as the square-roots of the
total area)---either a binomial, or a first bimedial, or a major, or the
square-root of a rational plus a medial (area). (Which is) the very thing
it was required to show.}
\end{Parallel}

%%%%
%10.72
%%%%
\pdfbookmark[1]{Proposition 10.72}{pdf10.72}
\begin{Parallel}{}{}
\ParallelLText{
\begin{center}
{\large \ggn{62}.}
\end{center}\vspace*{-7pt}

{\greekfont Δύο μέσων ἀσυμμέτρων ἀλλήλοις συντιθεμένων αἱ λοιπαὶ 
δύο ἄλογοι γίγνονται ἤτοι ἐκ δύο μέσων δευτέρα ἢ [ἡ]
δύο μέσα δυναμένη.}

{\greekfont Συγκείσθω γὰρ δύο μέσα ἀσύμμετρα ἀλλήλοις τὰ ΑΒ, ΓΔ; λέγω,
ὅτι ἡ τὸ ΑΔ θωρίον δυναμένη ἤτοι ἐκ δύο μέσων
ἐστὶ δευτέρα ἢ δύο μέσα δυναμένη.}

{\greekfont Τὸ γὰρ ΑΒ τοῦ ΓΔ ἤτοι μεῖζόν ἐστιν  ἢ ἔλασσον. ἔστω,
εἰ τύθον, πρότερον μεῖζον τὸ ΑΒ τοῦ ΓΔ; καὶ ἐκκείσθω ῥητὴ
ἡ ΕΖ, καὶ τῷ μὲν ΑΒ ἴσον παρὰ τὴν ΕΖ παραβεβλήσθω τὸ ΕΗ
πλάτος ποιοῦν τὴν ΕΘ, τῷ δὲ ΓΔ ἴσον τὸ ΘΙ πλάτος ποιοῦν τὴν ΘΚ.
καὶ ἐπεὶ μέσον ἐστὶν ἑκάτερον τῶν ΑΒ, ΓΔ, μέσον
ἄρα καὶ ἐκάτερον τῶν ΕΗ, ΘΙ. καὶ παρὰ ῥητὴν τὴν ΖΕ παράκειται
πλάτος ποιοῦν τὰς ΕΘ, ΘΚ; ἑκατέρα ἄρα τῶν ΕΘ, ΘΚ ῥητή
ἐστι καὶ ἀσύμμετρος τῇ ΕΖ μήκει. καὶ ἐπεὶ ἀσύμμετρόν
ἐστι τὸ ΑΒ τῷ ΓΔ, καί ἐστιν ἴσον τὸ μὲν ΑΒ τῷ ΕΗ, 
τὸ δὲ ΓΔ τῷ ΘΙ, ἀσύμμετρον ἄρα ἐστὶ καὶ τὸ ΕΗ τῷ ΘΙ.
ὡς δὲ τὸ ΕΗ πρὸς τὸ ΘΙ, οὕτως ἐστὶν ἡ ΕΘ
πρὸς ΘΚ; ἀσύμμετρος ἄρα ἐστὶν ἡ ΕΘ τῇ ΘΚ μήκει. αἱ 
ΕΘ, ΘΚ ἄρα ῥηταί εἰσι δυνάμει μόνον σύμμετροι; ἐκ
δύο ἄρα ὀνομάτων ἐστὶν ἡ ΕΚ. ἤτοι δὲ ἡ ΕΘ τῆς ΘΚ
μεῖζον δύναται τῷ ἀπὸ συμμέτρου ἑαυτῇ ἢ τῷ ἀπὸ
ἀσυμμέτρου. δυνάσθω πρότερον τῷ ἀπὸ συμμέτρου ἑαυτῇ
μήκει; καὶ οὐδετέρα τῶν ΕΘ, ΘΚ σύμμετρός ἐστι τῇ ἐκκειμένῃ
ῥητῇ τῇ ΕΖ μήκει; ἡ ΕΚ ἄρα ἐκ δύο ὀνομάτων ἐστὶ τρίτη.
ῥητὴ δὲ ἡ ΕΖ; ἐὰν δὲ θωρίον περιέθηται ὑπὸ ῥητῆς 
καὶ τῆς ἐκ δύο ὀνομάτων τρίτης, ἡ τὸ θωρίον δυναμένη
ἐκ δύο μέσων ἐστὶ δευτέρα; ἡ ἄρα τὸ ΕΙ, τουτέστι τὸ ΑΔ,
δυναμένη ἐκ δύο μέσων ἐστὶ δευτέρα. ἀλλα δὴ ἡ ΕΘ τῆς ΘΚ
μεῖζον δυνάσθω τῷ ἀπὸ ἀσυμμέτρου ἑαυτῇ μήκει; καὶ
ἀσύμμετρός ἐστιν ἑκατέρα τῶν ΕΘ, ΘΚ τῇ ΕΖ μήκει; ἡ ἄρα
ΕΚ ἐκ δύο ὀνομάτων ἐστὶν ἕκτη. ἐὰν δὲ θωρίον περιέθηται
ὑπὸ ῥητῆς καὶ τῆς ἐκ δύο ὀνομάτων ἕκτης, ἡ τὸ
θωρίον δυναμένη ἡ δύο μέσα δυναμένη ἐστίν; ὥστε καὶ ἡ
τὸ ΑΔ θωρίον δυναμένη ἡ δύο μέσα δυναμένη ἐστίν.}\\~\\~\\~\\~\\~\\~\\~\\~\\~\\~\\~\\~\\~\\

\epsfysize=2in
\centerline{\epsffile{Book10/fig071g.eps}}

{\greekfont μ̀βοχ{[}Ὁμοίως δὴ δείχομεν, ὅτι κἂν ἔλαττον ᾖ τὸ ΑΒ
τοῦ ΓΔ, ἡ τὸ ΑΔ θωρίον δυναμένη ἢ ἐκ δύο μέσων
δευτέρα ἐστὶν ἤτοι δύο μέσα δυναμένη].}

{\greekfont Δύο ἄρα μέσων ἀσυμμέτρων ἀλλήλοις συντιθεμένων αἱ λοιπαὶ
δύο ἄλογοι γίγνονται ἤτοι ἐκ δύο μέσων δευτέρα ἢ δύο μέσα
δυναμένη.}\\

{\greekfont Ἠ ἐκ δύο ὀνομάτων καὶ αἱ μετ α ὐτὴν ἄλογοι οὔτε
τῇ μέσῃ οὔτε ἀλλήλαις εἰσὶν αἱ α ὐταί. τὸ μὲν γὰρ ἀπὸ
μέσης παρὰ ῥητὴν παραβαλλόμενον πλάτος ποιεῖ ῥητὴν καὶ
ἀσύμμετρον τῇ παρ ἣν παράκειται μήκει. τὸ δὲ ἀπὸ τῆς
ἐκ δύο ὀνομάτων παρὰ ῥητὴν παραβαλλόμενον πλάτος ποιεῖ
τὴν ἐκ δύο ὀνομάτων πρώτην. τὸ δὲ ἀπὸ τῆς ἐκ δύο μέσων
πρώτης παρὰ ῥητὴν παραβαλλόμενον πλάτος ποιεῖ τὴν ἐκ δύο
ὀνομάτων δευτέραν. τὸ δὲ ἀπὸ τῆς ἐκ δύο μέσων δευτέρας παρὰ
ῥητὴν παραβαλλόμενον πλάτος ποιεῖ τὴν ἐκ δύο ὀνομάτων
τρίτην. τὸ δὲ ἀπὸ τῆς μείζονος παρὰ ῥητὴν παραβαλλόμενον
πλάτος ποιεῖ τὴν ἐκ δύο ὀνομάτων
τετάρτην. τὸ δὲ ἀπὸ τῆς ῥητὸν καὶ μέσον δυναμένης
παρὰ ῥητὴν παραβαλλόμενον πλάτος ποιεῖ τὴν ἐκ δύο
ὀνομάτων πέμπτην. τὸ δὲ ἀπὸ τῆς δύο μέσα δυναμένης
παρὰ ῥητὴν παραβαλλόμενον πλάτος ποιεῖ τὴν ἐκ δύο ὀνομάτων
ἕκτην. τὰ δ εἰρημένα πλάτη διαφέρει τοῦ τε πρώτου καὶ ἀλλήλων,
τοῦ μὲν πρώτου, ὅτι ῥητή ἐστιν, ἀλλήλων δέ, ὅτι τῇ τάχει
οὐκ εἰσὶν αἱ α ὐταί; ὥστε καὶ α ὐταὶ αἱ ἄλογοι διαφέρουσιν
ἀλλήλων.}}

\ParallelRText{
\begin{center}
{\large Proposition 72}
\end{center}

When two medial (areas which are) incommensurable with one another are added together, the remaining
two irrational (straight-lines) arise (as the square-roots of the total area)---either a second bimedial, or the square-root of (the sum of)
two medial (areas).

For let the two medial (areas) $AB$ and $CD$, (which are) incommensurable
with one another, be added together. I say that the square-root
of area $AD$ is either a second bimedial, or the square-root of (the sum of)
two medial (areas).

For $AB$ is either greater than or less than $CD$. By chance, let $AB$,
first of all,  be greater than $CD$. And let the rational (straight-line)
$EF$ be laid down. And let $EG$, equal to $AB$, be applied
to $EF$, producing $EH$ as breadth, and $HI$, equal to $CD$, producing
$HK$ as breadth. And since $AB$ and $CD$ are each medial,
$EG$ and $HI$ (are) thus also each medial. And they are applied to
the rational straight-line $FE$, producing $EH$ and $HK$ (respectively) as breadth. Thus, $EH$ and $HK$ are  each rational (straight-lines which are) incommensurable in length with $EF$ [Prop.~10.22]. And since $AB$ is incommensurable
with $CD$, and $AB$ is equal to $EG$, and $CD$ to $HI$, $EG$
is thus also incommensurable with $HI$. And as $EG$ (is) to $HI$,
so $EH$ is to $HK$ [Prop.~6.1].  $EH$ is thus
incommensurable in length with $HK$ [Prop.~10.11]. Thus, $EH$ and $HK$ are rational
(straight-lines which are) commensurable in square only. $EK$
is thus a binomial (straight-line) [Prop.~10.36]. 
And the square on $EH$ is  greater than (the square on) $HK$
either by the (square) on (some straight-line) commensurable (in length) with ($EH$), or
by the (square) on (some straight-line) incommensurable (in length
with $EH$). Let it, first of all, be greater by the square on (some
straight-line) commensurable in length with ($EH$). And neither of
$EH$ or $HK$ is commensurable in length with the (previously)
laid down rational (straight-line) $EF$.  Thus, $EK$ is a third
binomial (straight-line) [Def.~10.7]. And $EF$
(is) rational. And if an area is contained by a rational (straight-line) and
a third binomial (straight-line)  then the square-root of the area
is a second bimedial  (straight-line) [Prop.~10.56]. Thus,
the square-root of $EI$---that is to say, of $AD$---is a second bimedial.
And so, let the square on $EH$ be greater than (the square) on $HK$
by the (square)  on (some straight-line) incommensurable in length
with ($EH$).  And $EH$ and $HK$
are each incommensurable
in length with $EF$. Thus, $EK$ is a sixth binomial (straight-line)
[Def.~10.10]. And if an area
is contained by
a rational (straight-line) and a sixth binomial (straight-line)  then the square-root of the area is the square-root of (the sum of) two medial (areas) 
[Prop.~10.59]. Hence, the
square-root of area $AD$ is also the square-root of (the sum of)
two medial (areas).

\epsfysize=2in 
\centerline{\epsffile{Book10/fig071e.eps}}

\mbox{[}So, similarly, we can show that, even if $AB$ is less than $CD$, the
square-root of area $AD$ is either a second bimedial or the square-root
of (the sum of) two medial (areas).]

Thus, when two medial (areas which are) incommensurable with one another are added together, the remaining
two irrational (straight-lines) arise (as the square-roots of the total area)---either a second bimedial, or the square-root of (the sum of)
two medial (areas).\\

A binomial (straight-line), and the (other) irrational (straight-lines) 
after it, are neither the same as a medial (straight-line) nor (the same) as one another.
For the (square) on a medial (straight-line), applied to a rational (straight-line),
produces as breadth a rational (straight-line which is) also incommensurable
in length with (the straight-line) to which it is applied [Prop.~10.22]. And the (square) on a binomial
(straight-line), applied to a rational (straight-line), produces
as breadth a first binomial [Prop.~10.60]. 
And the (square) on a first bimedial (straight-line), applied
to a rational (straight-line), produces as breadth a  second binomial
[Prop.~10.61]. And the (square) on a second bimedial (straight-line), applied
to a rational (straight-line), produces as breadth a  third binomial
[Prop.~10.62]. And the (square) on a major (straight-line), applied
to a rational (straight-line), produces as breadth a  fourth binomial
[Prop.~10.63].  And the (square) on the square-root of a rational plus a medial (area), applied
to a rational (straight-line), produces as breadth a  fifth binomial
[Prop.~10.64]. And the (square) on the square-root of (the sum of) two medial (areas), applied
to a rational (straight-line), produces as breadth a  sixth binomial
[Prop.~10.65]. And the aforementioned breadths
differ from the first (breadth), and from one another---from the first, because it is rational---and from one another, because they are not the same in order. Hence, the (previously mentioned) irrational (straight-lines) themselves
also differ from one another.}
\end{Parallel}

%%%%
%10.73
%%%%
\pdfbookmark[1]{Proposition 10.73}{pdf10.73}
\begin{Parallel}{}{}
\ParallelLText{
\begin{center}
{\large \ggn{73}.}
\end{center}\vspace*{-7pt}

{\greekfont Ἐὰν ἀπὸ ῥητῆς ῥητὴ ἀφαιρεθῇ δυνάμει μόνον σύμμετρος οὖσα
τῇ ὅλῃ, ἡ λοιπὴ ἄλογός ἐστιν; καλείσθω δὲ ἀποτομή.}\\~\\

\epsfysize=0.3in
\centerline{\epsffile{Book10/fig073g.eps}}

{\greekfont Ἀπὸ γὰρ ῥητῆς τῆς ΑΒ ῥητὴ ἀφῃρήσθω ἡ ΒΓ δυνάμει μόνον
σύμμετρος οὖσα τῇ ὅλῃ; λέγω, ὅτι ἡ λοιπὴ ἡ ΑΓ ἄλογός
ἐστιν ἡ καλουμένη ἀποτομή.}

{\greekfont Ἐπεὶ γὰρ ἀσύμμετρός ἐστιν ἡ ΑΒ τῇ ΒΓ μήκει, καί ἐστιν
ὡς ἡ ΑΒ πρὸς τὴν ΒΓ, οὕτως τὸ ἀπὸ τῆς ΑΒ πρὸς τὸ ὑπὸ
τῶν ΑΒ, ΒΓ, ἀσύμμετρον ἄρα ἐστὶ τὸ ἀπὸ τῆς ΑΒ τῷ
ὑπὸ τῶν ΑΒ, ΒΓ. ἀλλὰ τῷ μὲν ἀπὸ τῆς ΑΒ σύμμετρά ἐστι
τὰ ἀπὸ τῶν ΑΒ, ΒΓ τετράγωνα, τῷ δὲ ὑπὸ τῶν ΑΒ, ΒΓ σύμμετρόν
ἐστι τὸ δὶς ὑπὸ τῶν ΑΒ, ΒΓ. καὶ ἐπειδήπερ τὰ ἀπὸ τῶν ΑΒ, ΒΓ
ἴσα ἐστὶ τῷ δὶς ὑπὸ τῶν ΑΒ, ΒΓ μετὰ τοῦ ἀπὸ ΓΑ, καὶ
λοιπῷ ἄρα τῷ ἀπὸ τῆς ΑΓ ἀσύμμετρά ἐστι τὰ ἀπὸ τῶν ΑΒ, ΒΓ.
ῥητὰ δὲ τὰ ἀπὸ τῶν ΑΒ, ΒΓ; ἄλογος ἄρα ἐστὶν ἡ ΑΓ; καλείσθω
δὲ ἀποτομή. ὅπερ ἔδει δεῖχαι.}}

\ParallelRText{
\begin{center}
{\large Proposition 73}
\end{center}

If  a  rational (straight-line), which is commensurable in square only with the
whole, is subtracted from a(nother) rational (straight-line)  then the remainder is an irrational (straight-line). Let it be called an apotome.

\epsfysize=0.3in 
\centerline{\epsffile{Book10/fig073e.eps}}

For let the rational (straight-line) $BC$, which commensurable in square only with the whole,  be subtracted from the
rational (straight-line) $AB$. I say that the remainder $AC$ is that irrational
(straight-line) called an apotome.

For since $AB$ is incommensurable in length with $BC$, and as $AB$ is to
$BC$, so the (square) on $AB$ (is) to the (rectangle contained) by $AB$ and
$BC$ [Prop.~10.21~lem.], the
(square) on $AB$ is thus incommensurable with the (rectangle contained)
by $AB$ and $BC$ [Prop.~10.11]. 
But, the (sum of the) squares on $AB$ and $BC$ is commensurable with the
(square) on $AB$ [Prop.~10.15], and twice
the (rectangle contained) by $AB$ and $BC$ is commensurable with the
(rectangle contained) by $AB$ and $BC$ [Prop.~10.6]. And, inasmuch as the (sum of the squares) on $AB$ and $BC$ is equal to twice the (rectangle contained) by
$AB$ and $BC$ plus the (square) on $CA$ [Prop.~2.7], the (sum of the squares) on $AB$ and $BC$
is thus also incommensurable with the remaining (square) on $AC$ [Props.~10.13, 10.16]. 
And the (sum of the squares) on $AB$ and $BC$ is rational. $AC$ is
thus an irrational (straight-line) [Def.~10.4]. And let it be called an
apotome.$^\dag$ (Which is) the very thing it was required to show.}
\end{Parallel}
{\footnotesize\noindent $^\dag$ See footnote to Prop.~10.36.}

%%%%
%10.74
%%%%
\pdfbookmark[1]{Proposition 10.74}{pdf10.74}
\begin{Parallel}{}{}
\ParallelLText{
\begin{center}
{\large \ggn{74}.}
\end{center}\vspace*{-7pt}

{\greekfont Ἐὰν ἀπὸ μέσης μέση ἀφαιρεθῇ δυνάμει μόνον σύμμετρος
οὖσα τῇ ὅλῃ, μετὰ δὲ τῆς ὅλης ῥητὸν περιέθουσα, ἡ λοιπὴ
ἄλογός ἐστιν; καλείσθω δὲ μέσης ἀποτομὴ πρώτη.}\\~\\~\\

\epsfysize=0.3in
\centerline{\epsffile{Book10/fig073g.eps}}

{\greekfont Ἀπὸ γὰρ μέσης τῆς ΑΒ μέση ἀφῃρήσθω ἡ ΒΓ δυνάμει
μόνον σύμμετρος οὖσα τῇ ΑΒ, μετὰ δὲ τῆς ΑΒ ῥητὸν
ποιοῦσα τὸ ὑπὸ τῶν ΑΒ, ΒΓ; λέγω, ὅτι
ἡ λοιπὴ ἡ ΑΓ ἄλογός ἐστιν; καλείσθω δὲ μέσης ἀποτομὴ
πρώτη.}

{\greekfont Ἐπεὶ γὰρ αἱ  ΑΒ, ΒΓ μέσαι εἰσίν, μέσα ἐστὶ καὶ τὰ ἀπὸ τῶν
ΑΒ, ΒΓ. ῥητὸν δὲ τὸ δὶς ὑπὸ τῶν ΑΒ, ΒΓ; ἀσύμμετρα ἄρα τὰ
ἀπὸ τῶν ΑΒ, ΒΓ τῷ δὶς ὑπὸ
τῶν ΑΒ, ΒΓ; καὶ λοιπῷ ἄρα τῷ ἀπὸ τῆς ΑΓ ἀσύμμετρόν ἐστι
τὸ δὶς ὑπὸ τῶν ΑΒ, ΒΓ, ἐπεὶ κἂν τὸ ὅλον ἑνὶ α ὐτῶν
ἀσύμμετρον ᾖ, καὶ τὰ ἐχ ἀρθῆς μεγέθη ἀσύμμετρα ἔσται.
ῥητὸν δὲ τὸ δὶς ὑπὸ τῶν ΑΒ, ΒΓ;
ἄλογον ἄρα τὸ ἀπὸ τῆς ΑΓ; ἄλογος ἄρα ἐστὶν ἡ ΑΓ;
καλείσθω δὲ μέσης ἀποτομὴ πρώτη.}}

\ParallelRText{
\begin{center}
{\large Proposition 74}
\end{center}

If  a medial (straight-line), which is commensurable in square only
with the whole, and which contains a rational (area) with the whole,  is subtracted from a(nother) medial
(straight-line)   then
the remainder is an irrational (straight-line). Let it be called a
first  apotome of a medial (straight-line).

\epsfysize=0.3in 
\centerline{\epsffile{Book10/fig073e.eps}}

For let the medial (straight-line) $BC$, which is commensurable in square only with $AB$,
and which makes with $AB$ the rational (rectangle contained) by $AB$ and $BC$,
be subtracted from the medial (straight-line) $AB$ [Prop.~10.27]. I say that
the remainder $AC$ is an irrational (straight-line). Let it be
called the first apotome of a medial (straight-line).

For since $AB$ and $BC$ are medial (straight-lines), the (sum of the squares)
on $AB$ and $BC$ is also medial. And twice the (rectangle contained)
by $AB$ and $BC$ (is) rational. The (sum of the squares) on $AB$
and $BC$ (is) thus incommensurable with twice the (rectangle contained)
by $AB$ and $BC$. Thus, twice the (rectangle contained) by
$AB$ and $BC$ is also incommensurable with the remaining (square) on $AC$ [Prop.~2.7], since  if the whole is
incommensurable with one of the (constituent magnitudes)   then the original magnitudes will also be incommensurable (with one another) [Prop.~10.16]. And twice the (rectangle
contained) by $AB$ and $BC$ (is) rational. Thus, the (square) on $AC$
is irrational. Thus, $AC$ is an irrational (straight-line) [Def.~10.4]. Let it be called a first apotome
of a medial (straight-line).$^\dag$}
\end{Parallel}
{\footnotesize\noindent $^\dag$ See footnote to Prop.~10.37.} 

%%%%
%10.75
%%%%
\pdfbookmark[1]{Proposition 10.75}{pdf10.75}
\begin{Parallel}{}{}
\ParallelLText{
\begin{center}
{\large \ggn{75}.}
\end{center}\vspace*{-7pt}

{\greekfont Ἐὰν ἀπὸ μέσης μέση ἀφαιρεθῇ δυνάμει μόνον σύμμετρος
οὖσα τῇ ὅλη, μετὰ δὲ τῆς ὅλης μέσον περιέθουσα, ἡ
λοιπὴ ἄλογός ἐστιν; καλείσθω δὲ μέσης ἀποτομὴ δευτέρα.}

{\greekfont Ἀπὸ γὰρ μέσης τῆς ΑΒ μέση ἀφῃρήσθω ἡ ΓΒ δυνάμει μόνον
σύμμετρος οὖσα τῇ ὅλῃ τῇ ΑΒ, μετὰ δὲ τῆς ὅλης τῆς ΑΒ
μέσον περιέθουσα τὸ ὑπὸ τῶν ΑΒ, ΒΓ; λέγω, ὅτι
ἡ λοιπὴ ἡ ΑΓ ἄλογός ἐστιν; καλείσθω δὲ μέσης ἀποτομὴ
δευτέρα.}\\~\\~\\~\\~\\

\epsfysize=1.7in
\centerline{\epsffile{Book10/fig075g.eps}}

{\greekfont Ἐκκείσθω γὰρ ῥητὴ ἡ ΔΙ, καὶ τοῖς μὲν ἀπὸ τῶν
ΑΒ, ΒΓ ἴσον παρὰ τὴν ΔΙ παραβεβλήσθω τὸ ΔΕ
πλάτος ποιοῦν τὴν ΔΗ, τῷ δὲ δὶς ὑπὸ τῶν ΑΒ, ΒΓ
ἴσον παρὰ τὴν ΔΙ παραβεβλήσθω τὸ ΔΘ πλάτος ποιοῦν τὴν
ΔΖ; λοιπὸν ἄρα τὸ ΖΕ ἴσον ἐστὶ τῷ ἀπὸ τῆς ΑΓ.
καὶ ἐπεὶ μέσα καὶ σύμμετρά ἐστι τὰ ἀπὸ τῶν ΑΒ, ΒΓ,
μέσον ἄρα καὶ τὸ ΔΕ. καὶ παρὰ ῥητὴν τὴν ΔΙ παράκειται
πλάτος ποιοῦν τὴν ΔΗ; ῥητὴ ἄρα ἐστὶν ἡ ΔΗ καὶ ἀσύμμετρος
τῇ ΔΙ μήκει. πάλιν, ἐπεὶ μέσον ἐστὶ τὸ ὑπὸ τῶν ΑΒ, ΒΓ, καὶ
τὸ δὶς ἄρα ὑπὸ τῶν ΑΒ, ΒΓ μέσον ἐστίν. καί ἐστιν ἴσον
τῷ ΔΘ; καὶ τὸ ΔΘ ἄρα μέσον ἐστίν. καὶ παρὰ ῥητὴν τὴν ΔΙ
παραβέβληται πλάτος ποιοῦν τὴν ΔΖ; ῥητὴ ἄρα ἐστὶν ἡ ΔΖ
καὶ ἀσύμμετρος τῇ ΔΙ μήκει. καὶ ἐπεὶ αἱ ΑΒ, ΒΓ δυνάμει
μόνον σύμμετροί εἰσιν, ἀσύμμετρος ἄρα ἐστὶν ἡ ΑΒ
τῇ ΒΓ μήκει; ἀσύμμετρον ἄρα καὶ τὸ ἀπὸ τῆς ΑΒ
τετράγωνον τῷ ὑπὸ τῶν ΑΒ, ΒΓ. ἀλλὰ τῷ μὲν ἀπὸ
τῆς ΑΒ σύμμετρά ἐστι τὰ ἀπὸ τῶν ΑΒ, ΒΓ, τῷ δὲ ὑπὸ
τῶν ΑΒ, ΒΓ σύμμετρόν ἐστι τὸ δὶς ὑπὸ τῶν ΑΒ, ΒΓ;
ἀσύμμετρον ἄρα ἐστὶ τὸ δὶς ὑπὸ τῶν ΑΒ, ΒΓ τοῖς
ἀπὸ τῶν ΑΒ, ΒΓ. ἴσον δὲ τοῖς μὲν ἀπὸ τῶν ΑΒ, ΒΓ
τὸ ΔΕ, τῷ δὲ δὶς ὑπὸ τῶν ΑΒ, ΒΓ τὸ ΔΘ; ἀσύμμετρον
ἄρα [ἐστὶ] τὸ ΔΕ τῷ ΔΘ. ὡς δὲ τὸ ΔΕ πρὸς τὸ ΔΘ,
οὕτως ἡ ΗΔ πρὸς τὴν ΔΖ; ἀσύμμετρος ἄρα ἐστὶν ἡ ΗΔ
τῇ ΔΖ. καί εἰσιν ἀμφότεραι ῥηταί; αἱ ἄρα ΗΔ, ΔΖ
ῥηταί εἰσι δυνάμει μόνον σύμμετροι; ἡ ΖΗ ἄρα ἀποτομή
ἐστιν. ῥητὴ δὲ ἡ ΔΙ; τὸ δὲ ὑπὸ ῥητῆς καὶ ἀλόγου
περιεχόμενον ἄλογόν ἐστιν, καὶ ἡ δυναμένη α ὐτὸ
ἄλογός ἐστιν. καὶ δύναται τὸ ΖΕ ἡ ΑΓ; ἡ ΑΓ ἄρα ἄλογός
ἐστιν; καλείσθω δὲ μέσης ἀποτομὴ δευτέρα. ὅπερ
ἔδει δεῖχαι.}}

\ParallelRText{
\begin{center}
{\large Proposition 75}
\end{center}

If  a medial (straight-line), which is commensurable in square only
with the whole, and which contains a medial (area) with the whole, is subtracted from a(nother) medial
(straight-line)   then
the remainder is an irrational (straight-line). Let it be called a second apotome of a medial (straight-line).

For let the medial (straight-line) $CB$, which is commensurable in square
only with the whole, $AB$, and which contains with the whole, $AB$, the
medial (rectangle contained) by $AB$ and $BC$, be subtracted from the medial (straight-line) $AB$ [Prop.~10.28].  I say that the remainder $AC$ is an
irrational (straight-line). Let it be called a second apotome of a medial (straight-line).

\epsfysize=1.7in 
\centerline{\epsffile{Book10/fig075e.eps}}

For let the rational (straight-line) $DI$ be laid down. And let $DE$, equal
to the (sum of the squares) on $AB$ and $BC$, be applied to
$DI$, producing $DG$ as breadth. And let $DH$, equal to twice the
(rectangle contained) by $AB$ and $BC$, be applied to $DI$,
producing $DF$ as breadth. The remainder $FE$ is thus equal to the
(square) on $AC$ [Prop.~2.7]. And since the
(squares) on $AB$ and $BC$ are medial and commensurable (with one another), $DE$ (is) thus also medial [Props.~10.15, 10.23~corr.]. And it is applied to the
rational (straight-line) $DI$, producing $DG$ as breadth. Thus, $DG$
is rational, and incommensurable in length with $DI$ [Prop.~10.22].  Again, since the (rectangle contained)
by $AB$ and $BC$ is medial, twice the (rectangle contained) by $AB$
and $BC$ is thus also medial [Prop.~10.23~corr.]. 
And it is equal to $DH$. Thus, $DH$ is also medial. And it has been
applied to the rational (straight-line) $DI$, producing $DF$ as breadth.
$DF$ is thus rational, and incommensurable in length with $DI$
[Prop.~10.22].  And since $AB$ and $BC$ are commensurable in square only, $AB$ is thus incommensurable in length
with $BC$. Thus, the square on $AB$ (is) also incommensurable
with the (rectangle contained) by $AB$ and $BC$ [Props.~10.21~lem.,
10.11]. But, the (sum of the squares) on $AB$ and $BC$  is commensurable with the (square) on $AB$  [Prop.~10.15], and twice the (rectangle contained)
by $AB$ and $BC$ is commensurable with the (rectangle contained)
by $AB$ and $BC$ [Prop.~10.6]. Thus, twice
the (rectangle contained) by $AB$ and $BC$ is incommensurable
with the (sum of the squares) on $AB$ and $BC$ [Prop.~10.13]. And $DE$ is equal
to the (sum of the squares) on $AB$ and $BC$, and $DH$ to
twice the (rectangle contained) by $AB$ and $BC$. Thus, $DE$
[is] incommensurable with $DH$. And as $DE$ (is) to $DH$, so
$GD$ (is) to $DF$ [Prop.~6.1]. Thus,
$GD$ is incommensurable with $DF$ [Prop.~10.11]. And they are both rational (straight-lines). Thus, $GD$ and $DF$ are rational (straight-lines which are)
commensurable in square only.
Thus, $FG$ is an apotome [Prop.~10.73]. 
And $DI$ (is) rational. And the (area) contained by a rational and
an irrational (straight-line) is irrational [Prop.~10.20], and its square-root 
is irrational. And $AC$ is the square-root of $FE$. Thus, $AC$ is an
irrational (straight-line) [Def.~10.4]. And
let it be called the second apotome of a medial (straight-line).$^\dag$
(Which is) the very thing it
was required to show.}
\end{Parallel}
{\footnotesize\noindent$^\dag$ See footnote to Prop.~10.38.} 

%%%%
%10.76
%%%%
\pdfbookmark[1]{Proposition 10.76}{pdf10.76}
\begin{Parallel}{}{}
\ParallelLText{
\begin{center}
{\large \ggn{76}.}
\end{center}\vspace*{-7pt}

{\greekfont Ἐὰν ἀπὸ εὐθείας εὐθεῖα ἀφαιρεθῂ δυνάμει ἀσύμμετ-ρος
οὖσα τῇ ὅλῃ, μετὰ δὲ τῆς ὅλης ποιοῦσα τὰ μὲν ἀπ
α ὐτῶν ἅμα ῥητόν, τὸ δ ὑπ α ὐτῶν μέσον, ἡ λοιπὴ
ἄλογός ἐστιν; καλείσθω δὲ ἐλάσσων.}\\~\\

\epsfysize=0.27in
\centerline{\epsffile{Book10/fig073g.eps}}

{\greekfont Ἀπὸ γὰρ εὐθείας τῆς ΑΒ εὐθεῖα ἀφῃρήσθω ἡ ΒΓ δυνάμει
ἀσύμμετρος οὖσα τῇ ὅλῃ ποιοῦσα τὰ προκείμενα. λέγω,
ὅτι ἡ λοιπὴ ἡ ΑΓ ἄλογός ἐστιν ἡ καλουμένη ἐλάσσων.}

{\greekfont Ἐπεὶ γὰρ τὸ μὲν συγκείμενον ἐκ τῶν ἀπὸ τῶν ΑΒ, ΒΓ τετραγώνων ῥητόν ἐστιν, τὸ δὲ δὶς ὑπὸ τῶν ΑΒ, ΒΓ μέσον,
ἀσύμμετρα ἄρα ἐστὶ τὰ ἀπὸ τῶν ΑΒ, ΒΓ τῷ δὶς ὑπὸ
τῶν ΑΒ, ΒΓ; καὶ ἀναστρέψαντι λοιπῷ τῷ ἀπὸ τῆς ΑΓ
ἀσύμμετρά ἐστι τὰ ἀπὸ τῶν ΑΒ, ΒΓ.  ῥητὰ δὲ τὰ
ἀπὸ τῶν ΑΒ, ΒΓ;
ἄλογον ἄρα
τὸ ἀπὸ τῆς ΑΓ;
 ἄλογος
ἄρα ἡ ΑΓ; καλείσθω δὲ ἐλάσσων. ὅπερ ἔδει δεῖχαι.}}

\ParallelRText{
\begin{center}
{\large Proposition 76}
\end{center}

If  a straight-line, which is incommensurable in square with the whole,
and  with the whole makes the (squares) on them (added) together rational, and the
(rectangle contained) by them medial, is subtracted from a(nother) straight-line
 then the remainder is an
irrational (straight-line). Let it be called a minor (straight-line).

\epsfysize=0.27in 
\centerline{\epsffile{Book10/fig073e.eps}}

For let the straight-line $BC$, which is incommensurable in square with the whole, and fulfils the (other) prescribed (conditions), be subtracted from
the straight-line $AB$ [Prop.~10.33]. I say that the remainder $AC$ is that irrational (straight-line)
called minor.

For since the sum of the squares on $AB$ and $BC$ is rational, and twice
the (rectangle contained) by $AB$ and $BC$ (is) medial, the (sum of the squares) on $AB$ and $BC$ is thus incommensurable with twice the
(rectangle contained) by $AB$ and $BC$. And, via conversion, 
the (sum of the squares) on $AB$ and $BC$ is incommensurable with the
remaining (square) on $AC$ [Props.~2.7, 10.16]. And the (sum of the squares) on $AB$
and $BC$ (is) rational. The (square) on $AC$ (is) thus irrational. Thus, 
$AC$ (is) an irrational (straight-line) [Def.~10.4]. Let it be
called a minor (straight-line).$^\dag$ (Which is) the very thing it
was required to show.}
\end{Parallel}
{\footnotesize\noindent$^\dag$ See footnote to Prop.~10.39.}

%%%%
%10.77
%%%%
\pdfbookmark[1]{Proposition 10.77}{pdf10.77}
\begin{Parallel}{}{}
\ParallelLText{
\begin{center}
{\large \ggn{77}.}
\end{center}\vspace*{-7pt}

{\greekfont Ἐὰν ἀπὸ εὐθείας εὐθεῖα ἀφαιρεθῇ δυνάμει ἀσύμμετ-ρος
οὖσα τῇ ὅλῃ, μετὰ δὲ τῆς ὅλης ποιοῦσα τὸ μὲν συγκείμενον
ἐκ τῶν ἀπ α ὐτῶν τετραγώνων μέσον, τὸ δὲ δὶς ὑπ
α ὐτῶν ῥητόν, ἡ λοιπὴ ἄλογός ἐστιν; καλείσθω δὲ ἡ μετὰ ῥητοῦ
μέσον τὸ ὅλον ποιοῦσα.}\\

\epsfysize=0.27in
\centerline{\epsffile{Book10/fig073g.eps}}

{\greekfont Ἀπὸ γὰρ εὐθείας τῆς ΑΒ εὐθεῖα ἀφῃρήσθω ἡ ΒΓ δυνάμει
ἀσύμμετος οὖσα τῇ ΑΒ ποιοῦσα τὰ προκείμενα; λέγω,
ὅτι ἡ λοιπὴ ἡ ΑΓ ἄλογός ἐστιν ἡ προειρημένη.}

{\greekfont Ἐπεὶ γὰρ τὸ μὲν συγκείμενον ἐκ τῶν ἀπὸ τῶν ΑΒ, ΒΓ
τετραγώνων μέσον ἐστίν, τὸ δὲ δὶς ὑπὸ τῶν ΑΒ, ΒΓ ῥητόν,
ἀσύμμετρα ἄρα ἐστὶ τὰ ἀπὸ τῶν ΑΒ, ΒΓ τῷ
δὶς ὑπὸ τῶν ΑΒ, ΒΓ; καὶ λοιπὸν ἄρα τὸ ἀπὸ τῆς ΑΓ ἀσύμμετρόν
ἐστι τῷ δὶς ὑπὸ τῶν ΑΒ, ΒΓ. καί ἐστι τὸ δὶς ὑπὸ τῶν
ΑΒ, ΒΓ ῥητόν; τὸ ἄρα ἀπὸ τῆς ΑΓ ἄλογόν ἐστιν; ἄλογος
ἄρα ἐστὶν ἡ ΑΓ; καλείσθω δὲ ἡ μετὰ ῥητοῦ μέσον τὸ ὅλον
ποιοῦσα. ὅπερ ἔδει δεῖχαι.}}

\ParallelRText{
\begin{center}
{\large Proposition 77}
\end{center}

If  a straight-line, which is incommensurable in  square with the whole, and with the
whole makes the sum of the squares on them medial, and twice
the (rectangle contained) by them rational, is subtracted from a(nother) straight-line 
 then the remainder
is an irrational (straight-line). Let it be called that which makes
with a rational (area) a medial whole.

\epsfysize=0.27in 
\centerline{\epsffile{Book10/fig073e.eps}}

For let the straight-line $BC$, which is incommensurable in
square with $AB$, and fulfils the (other) prescribed (conditions), be
subtracted from the straight-line $AB$ [Prop.~10.34]. I say that the remainder
$AC$ is the aforementioned irrational (straight-line).

For since the sum of the squares on $AB$ and $BC$ is medial, and
twice the (rectangle contained) by $AB$ and $BC$ rational,
the (sum of the squares) on $AB$ and $BC$ is thus incommensurable
with twice the (rectangle contained) by $AB$ and $BC$. Thus, the
remaining (square) on $AC$ is also incommensurable with
twice the (rectangle contained) by $AB$ and $BC$ [Props.~2.7, 10.16]. And twice
the (rectangle contained) by $AB$ and $BC$ is rational. 
Thus, the (square) on $AC$ is irrational. Thus, $AC$ is an
irrational (straight-line) [Def.~10.4]. And let it be
called that which makes
with a rational (area) a medial whole.$^\dag$ (Which is) the very thing it
was required to show.}
\end{Parallel}
{\footnotesize\noindent$^\dag$ See footnote to Prop.~10.40.}

%%%%
%10.78
%%%%
\pdfbookmark[1]{Proposition 10.78}{pdf10.78}
\begin{Parallel}{}{}
\ParallelLText{
\begin{center}
{\large \ggn{78}.}
\end{center}\vspace*{-7pt}

{\greekfont Ἐὰν ἀπὸ εὐθείας εὐθεῖα ἀφαιρεθῇ δυνάμει ἀσύμμ-ετρος
οὖσα τῇ ὅλῃ, μετὰ δὲ τῆς ὅλης ποιοῦσα τό τε συγκείμενον
ἐκ τῶν ἀπ α ὐτῶν τετραγώνων μέσον τό τε δὶς ὑπ α ὐτῶν
μέσον καὶ ἔτι τὰ ἀπ α ὐτῶν τετράγωνα ἀσύμμετρα τῷ
δὶς ὑπ α ὐτῶν, ἡ λοιπὴ ἄλογός ἐστιν; καλείσθω δὲ ἡ μετὰ
μέσου μέσον τὸ ὅλον ποιοῦσα.}\\~\\

\epsfysize=1.7in
\centerline{\epsffile{Book10/fig078g.eps}}

{\greekfont Ἀπὸ γὰρ εὐθείας τῆς ΑΒ εὐθεῖα ἀφῃρήσθω ἡ ΒΓ δυνάμει
ἀσύμμετρος οὖσα τῇ ΑΒ ποιοῦσα τὰ προκείμενα; λέγω,
ὅτι ἡ λοιπὴ ἡ ΑΓ ἄλογός ἐστιν ἡ καλουμένη ἡ μετὰ
μέσου μέσον τὸ ὅλον ποιοῦσα.}

{\greekfont Ἐκκείσθω γὰρ ῥητὴ ἡ ΔΙ, καὶ τοῖς μὲν ἀπὸ τῶν ΑΒ, ΒΓ ἴσον
παρὰ τὴν ΔΙ παραβεβλήσθω τὸ ΔΕ πλάτος ποιοῦν τὴν ΔΗ, τῷ δὲ δὶς
ὑπὸ τῶν ΑΒ, ΒΓ ἴσον ἀφῃρήσθω τὸ ΔΘ [πλάτος ποιοῦν τὴν
ΔΖ]. λοιπὸν ἄρα τὸ ΖΕ ἴσον ἐστὶ τῷ ἀπὸ τῆς ΑΓ;
ὥστε ἡ ΑΓ δύναται τὸ ΖΕ. καὶ ἐπεὶ τὸ συγκείμενον ἐκ
τῶν ἀπὸ τῶν ΑΒ, ΒΓ τετραγώνων μέσον ἐστὶ καί ἐστιν
ἴσον τῷ ΔΕ, μέσον ἄρα [ἐστὶ] τὸ ΔΕ. καὶ παρὰ ῥητὴν
τὴν ΔΙ παράκειται πλάτος ποιοῦν τὴν ΔΗ; ῥητὴ
ἄρα ἐστὶν ἡ ΔΗ καὶ ἀσύμμετρος τῇ ΔΙ μήκει. πάλιν,
ἐπεὶ τὸ δὶς ὑπὸ τῶν ΑΒ, ΒΓ μέσον ἐστὶ καί ἐστιν ἴσον
τῷ ΔΘ, τὸ ἄρα ΔΘ μέσον ἐστίν. καὶ παρὰ ῥητὴν τὴν
ΔΙ παράκειται πλάτος ποιοῦν  τὴν ΔΖ;
ῥητὴ ἄρα ἐστὶ καὶ ἡ ΔΖ καὶ ἀσύμμετρος τῇ ΔΙ μήκει.
 καὶ ἐπεὶ ἀσύμμετρά ἐστι τὰ ἀπὸ τῶν
ΑΒ, ΒΓ τῷ δὶς ὑπὸ τῶν ΑΒ, ΒΓ, ἀσύμμετρον ἄρα καὶ
τὸ ΔΕ τῷ ΔΘ. ὡς δὲ τὸ ΔΕ πρὸς τὸ ΔΘ, οὕτως ἐστὶ
καὶ ἡ ΔΗ πρὸς τὴν ΔΖ; ἀσύμμετρος ἄρα ἡ ΔΗ τῇ ΔΖ.
καί εἰσιν ἀμφότεραι ῥηταί; αἱ ΗΔ, ΔΖ ἄρα ῥηταί
εἰσι δυνάμει μόνον σύμμετροι. ἀποτομὴ ἄρα ἐστίν ἡ
ΖΗ; ῥητὴ δὲ ἡ ΖΘ. τὸ δὲ ὑπὸ ῥητῆς καὶ ἀποτομῆς
περιεχόμενον [ὀρθογώνιον] ἄλογόν ἐστιν, καὶ ἡ δυναμένη
α ὐτὸ ἄλογός ἐστιν; καὶ δύναται τὸ ΖΕ ἡ ΑΓ; ἡ ΑΓ ἄρα ἄλογός ἐστιν;
 καλείσθω δὲ ἡ μετὰ μέσου μέσον τὸ ὅλον
ποιοῦσα. ὅπερ ἔδει δεῖχαι.}}

\ParallelRText{
\begin{center}
{\large Proposition 78}
\end{center}

If a straight-line, which is incommensurable in square with the whole, and with the whole
makes the sum of the squares on them medial, and twice the (rectangle
contained) by them medial, and, moreover, the (sum of the) squares on them
incommensurable with twice the (rectangle contained) by them, is
subtracted from a(nother) straight-line  then the
remainder is an irrational (straight-line). Let it be called that which makes
with a medial (area) a medial whole.

\epsfysize=1.7in 
\centerline{\epsffile{Book10/fig078e.eps}}

For let the straight-line $BC$, which is incommensurable in square
$AB$, and fulfils the (other) prescribed (conditions), be subtracted from the
(straight-line) $AB$ [Prop.~10.35].  I say that the remainder $AC$ is the irrational (straight-line) called  that which makes with a medial (area) a medial whole.

For let the rational (straight-line) $DI$ be laid down. And let $DE$, equal to
the (sum of the squares) on $AB$ and $BC$, be applied to $DI$,
producing $DG$ as breadth. And let $DH$, equal to twice the (rectangle
contained) by $AB$ and $BC$, be subtracted (from $DE$) [producing
$DF$ as breadth]. Thus, the remainder $FE$ is equal to the (square) on 
$AC$ [Prop.~2.7]. Hence, $AC$ is the square-root of $FE$. And since the sum of the squares on $AB$ and $BC$ is medial, and
is equal to $DE$, $DE$ [is] thus medial. And it is applied to the
rational (straight-line) $DI$, producing $DG$ as breadth. Thus, $DG$
is rational, and incommensurable in length with $DI$ [Prop~10.22]. Again, since twice the (rectangle
contained) by $AB$ and $BC$ is medial, and is equal to $DH$, $DH$
is thus medial. And it is applied to the rational (straight-line)
$DI$, producing $DF$ as breadth.  Thus, $DF$ is also rational, and incommensurable in length with $DI$ [Prop.~10.22]. And since  the (sum of the squares)
on $AB$ and $BC$ is incommensurable with twice the (rectangle contained)
by $AB$ and $BC$, $DE$ (is) also incommensurable with $DH$. 
And as $DE$ (is) to $DH$, so $DG$  also is to $DF$ [Prop.~6.1]. Thus, $DG$ (is) incommensurable (in length) with
$DF$ [Prop.~10.11]. And they are both rational. 
Thus, $GD$ and $DF$ are rational (straight-lines which are) commensurable
in square only. Thus, $FG$ is an apotome [Prop.~10.73]. And $FH$ (is) rational. And the
[rectangle] contained by a rational (straight-line) and an apotome is
irrational [Prop.~10.20], and its square-root
is irrational. And $AC$ is the square-root of $FE$.  Thus, $AC$
is irrational. Let it be called that which makes with a medial (area) a medial
whole.$^\dag$ 
(Which is) the very thing it was required to show.}
\end{Parallel}
{\footnotesize\noindent$^\dag$ See footnote to Prop.~10.41.}

%%%%
%10.79
%%%%
\pdfbookmark[1]{Proposition 10.79}{pdf10.79}
\begin{Parallel}{}{}
\ParallelLText{
\begin{center}
{\large \ggn{79}.}
\end{center}\vspace*{-7pt}

{\greekfont Τῇ ἀποτομ ῇ μία [μόνον] προσαρμόζει εὐθεῖα ῥητὴ δυνάμει
μόνον σύμμετρος οὖσα τῇ ὅλῃ.}\\

\epsfysize=0.27in
\centerline{\epsffile{Book10/fig079g.eps}}

{\greekfont Ἔστω ἀποτομὴ ἡ ΑΒ, προσαρμόζουσα δὲ α ὐτῇ ἡ ΒΓ; αἱ ΑΓ, ΓΒ
ἄρα ῥηταί εἰσι δυνάμει μόνον σύμμετροι; λέγω, ὅτι τῇ ΑΒ ἑτέρα οὐ
προσαρμόζει ῥητὴ δυνάμει μόνον σύμμετρος οὖσα τῇ ὅλῇ.}

{\greekfont Εἰ γὰρ δυνατόν, προσαρμοζέτω ἡ ΒΔ; καὶ αἱ ΑΔ, ΔΒ ἄρα ῥηταί
εἰσι δυνάμει μόνον σύμμετροι. καὶ ἐπεί, ᾧ ὑπερέθει τὰ
ἀπὸ τῶν ΑΔ, ΔΒ τοῦ δὶς ὑπὸ τῶν ΑΔ, ΔΒ, τούτῳ ὑπερέθει
καὶ τὰ ἀπὸ τῶν ΑΓ, ΓΒ τοῦ δὶς ὑπὸ τῶν ΑΓ, ΓΒ; τῷ γὰρ
α ὐτῷ τῷ ἀπὸ τῆς ΑΒ ἀμφότερα ὑπερέθει; ἐναλλὰχ ἄρα, ᾧ
ὑπερέθει τὰ ἀπὸ τῶν ΑΔ, ΔΒ τῶν ἀπὸ τῶν ΑΓ, ΓΒ, τούτῳ ὑπερέθει [καὶ] τὸ δὶς ὑπὸ τῶν ΑΔ, ΔΒ
τοῦ δὶς ὑπὸ τῶν ΑΓ, ΓΒ. τὰ δὲ ἀπὸ τῶν ΑΔ, ΔΒ τῶν ἀπὸ τῶν
ΑΓ, ΓΒ ὑπερέθει ῥητῷ; ῥητὰ γὰρ ἀμφότερα. καὶ τὸ δὶς ἄρα ὑπὸ
τῶν ΑΔ, ΔΒ τοῦ δὶς ὑπὸ τῶν ΑΓ, ΓΒ ὑπερέθει ῥητῷ; ὅπερ ἐστὶν ἀδύνατον; μέσα γὰρ ἀμφότερα, μέσον δὲ μέσου οὐθ ὑπερέθει
ῥητῷ. τῇ ἄρα ΑΒ ἑτέρα οὐ προσαρμόζει ῥητὴ δυνάμει
μόνον σύμμετρος οὖσα τῇ ὅλῃ.}

{\greekfont Μία ἄρα μόνη τῇ ἀποτομ ῇ προσαρμόζει ῥητὴ δυνάμει
μόνον σύμμετρος οὖσα τῇ ὅλῃ; ὅπερ ἔδει
δεῖχαι.}}

\ParallelRText{
\begin{center}
{\large Proposition 79}
\end{center}

\mbox{[}Only] one rational straight-line, which is
commensurable in square only with the whole, can be attached to an apotome.$^\dag$

\epsfysize=0.26in 
\centerline{\epsffile{Book10/fig079e.eps}}

Let $AB$ be an apotome, with $BC$  (so) attached to it. $AC$ and $CB$ are thus
rational (straight-lines which are) commensurable in square only [Prop.~10.73]. I say that another rational (straight-line), which is commensurable in  square only with the whole, cannot
be attached to $AB$. 

For, if possible, let $BD$ be (so) attached (to $AB$). Thus, $AD$ and $DB$ are also
rational (straight-lines which are) commensurable in square only [Prop.~10.73]. And since by whatever (area)
the (sum of the squares) on $AD$ and $DB$ exceeds twice the
(rectangle contained) by $AD$ and $DB$, the (sum of the squares)
on $AC$ and $CB$ also exceeds twice the (rectangle contained) by $AC$ and
$CB$ by this (same area). For both exceed by the same (area)---(namely), the (square)
on $AB$ [Prop.~2.7]. Thus, alternately, 
by whatever (area) the (sum of the squares) on $AD$ and 
$DB$ exceeds the (sum of the squares) on $AC$ and $CB$, 
 twice the (rectangle contained) by $AD$ and $DB$
[also] exceeds twice the (rectangle contained) by $AC$ and $CB$ by this (same area).
And the (sum of the squares) on $AD$ and $DB$ exceeds the (sum
of the squares) on $AC$ and $CB$ by a rational (area). For both
(are) rational (areas). Thus, twice the (rectangle contained) by $AD$ and $DB$ also
exceeds twice the (rectangle contained) by $AC$ and $CB$ by a
rational (area). The very thing is impossible. For both are
medial (areas) [Prop.~10.21],
and a medial (area) cannot exceed a(nother) medial (area) by a rational
(area) [Prop.~10.26].  
Thus, another rational (straight-line), which is commensurable in  square only with the whole, cannot
be attached to $AB$.

Thus, only one rational (straight-line), which is
commensurable in square only with the whole, can be attached to an apotome.
(Which is) the very thing it was required to show.}
\end{Parallel}
{\footnotesize\noindent$^\dag$ This proposition is equivalent to 
Prop.~10.42, with minus signs instead of
plus signs.}

%%%%
%10.80
%%%%
\pdfbookmark[1]{Proposition 10.80}{pdf10.80}
\begin{Parallel}{}{}
\ParallelLText{
\begin{center}
{\large \ggn{80}.}
\end{center}\vspace*{-7pt}

{\greekfont Τῇ μέσης ἀποτομ ῇ πρώτῃ μία μόνον προσαρμόζει εὐθεῖα
μέση δυνάμει μόνον σύμμετρος οὖσα τῇ ὅλῃ, μετὰ δὲ τῆς ὅλης
ῥητὸν περιέθουσα.}\\

\epsfysize=0.26in
\centerline{\epsffile{Book10/fig079g.eps}}

{\greekfont Ἔστω γὰρ μέσης ἀποτομὴ πρώτη ἡ ΑΒ, καὶ τῇ ΑΒ προσαρμοζέτω ἡ ΒΓ; αἱ ΑΓ,
ΓΒ ἄρα μέσαι εἰσὶ δυνάμει μόνον σύμμετροι ῥητὸν περιέθουσαι
τὸ ὑπὸ τῶν ΑΓ, ΓΒ; λέγω, ὅτι τῇ ΑΒ ἑτέρα οὐ προσαρμόζει
μέση δυνάμει μόνον σύμμετρος οὖσα τῇ ὅλῃ, μετὰ δὲ
τῆς ὅλης ῥητὸν περιέθουσα.}

{\greekfont Εἰ γὰρ δυνατόν, προσαρμοζέτω καὶ ἡ ΔΒ; αἱ ἄρα 
ΑΔ, ΔΒ μέσαι εἰσὶ δυνάμει μόνον σύμμετροι ῥητὸν περιέθουσαι
τὸ ὑπὸ τῶν ΑΔ, ΔΒ. καὶ ἐπεί, ᾧ ὑπερέθει τὰ ἀπὸ τῶν
ΑΔ, ΔΒ τοῦ δὶς ὑπὸ τῶν ΑΔ, ΔΒ, τούτῳ ὑπερέθει καὶ τὰ ἀπὸ
τῶν ΑΓ, ΓΒ τοῦ δὶς ὑπὸ τῶν ΑΓ, ΓΒ; τῷ γὰρ α ὐτῷ [πάλιν]
ὑπερέθουσι τῷ ἀπὸ τῆς ΑΒ; ἐναλλὰχ ἄρα, ᾧ ὑπερέθει τὰ
ἀπὸ τῶν ΑΔ, ΔΒ τῶν ἀπὸ τῶν ΑΓ, ΓΒ, τούτῳ ὑπερέθει
καὶ τὸ δὶς ὑπὸ τῶν
ΑΔ, ΔΒ τοῦ δὶς ὑπὸ τῶν
ΑΓ, ΓΒ. τὸ δὲ δὶς ὑπὸ τῶν ΑΔ, ΔΒ τοῦ δὶς ὑπὸ τῶν ΑΓ, ΓΒ
ὑπερέθει ῥητῷ; ῥητὰ γὰρ ἀμφότερα. καὶ τὰ ἀπὸ τῶν ΑΔ, ΔΒ
ἄρα τῶν ἀπὸ τῶν ΑΓ, ΓΒ [τετραγώνων]
ὑπερέθει ῥητῷ; ὅπερ ἐστὶν ἀδύνατον; μέσα γάρ ἐστιν ἀμφότερα,
μέσον δὲ μέσου οὐθ ὑπερέθει ῥητῷ.}

{\greekfont Τῇ ἄρα μέσης ἀποτομ ῇ πρώτῃ μία μόνον προσαρμόζει
εὐθεῖα μέση δυνάμει μόνον σύμμετρος οὖσα τῇ ὅλῃ,
μετὰ δὲ τῆς ὅλης ῥητὸν περιέθουσα; ὅπερ ἔδει δεῖχαι.}}

\ParallelRText{
\begin{center}
{\large Proposition 80}
\end{center}

Only one medial straight-line,
which is commensurable in square only with the whole, and
contains a rational (area) with the whole, can be attached
to a first apotome of a medial (straight-line).$^\dag$

\epsfysize=0.26in 
\centerline{\epsffile{Book10/fig079e.eps}}

For let $AB$ be a first apotome of a medial (straight-line), and let $BC$ be (so) attached
to $AB$. Thus, $AC$ and $CB$ are medial (straight-lines which are)
commensurable in square only, containing a rational (area)---(namely, that
contained) by $AC$ and $CB$ [Prop.~10.74]. 
I say that a(nother) medial (straight-line), which is commensurable in square only with the whole, and
contains a rational (area) with the whole, cannot be attached to $AB$.

For, if possible, let $DB$ also be (so) attached to $AB$. Thus, $AD$ and
$DB$ are medial (straight-lines which are) commensurable in square only, containing a rational (area)---(namely, that)
contained by $AD$ and $DB$ [Prop.~10.74]. 
And since by whatever (area)
the (sum of the squares) on $AD$ and $DB$ exceeds twice the
(rectangle contained) by $AD$ and $DB$,  the (sum of the squares)
on $AC$ and $CB$ also exceeds twice the (rectangle contained) by $AC$ and
$CB$ by this (same area). For [again] both exceed by the same (area)---(namely), the (square)
on $AB$ [Prop.~2.7]. Thus, alternately, 
by whatever (area) the (sum of the squares) on $AD$ and 
$DB$ exceeds the (sum of the squares) on $AC$ and $CB$, 
 twice the (rectangle contained) by $AD$ and $DB$
also exceeds twice the (rectangle contained) by $AC$ and $CB$ by this (same area).
And twice the (rectangle contained) by $AD$ and $DB$ exceeds twice the (rectangle contained) by $AC$ and $CB$ by a rational (area). For both
(are) rational (areas). Thus,  the (sum of the squares) on $AD$ and $DB$ also
exceeds the (sum of the) [squares] on $AC$ and $CB$ by a
rational (area). The very thing is impossible. For both are
medial (areas) [Props.~10.15, 10.23~corr.],
and a medial (area) cannot exceed a(nother) medial (area) by a rational
(area) [Prop.~10.26].

Thus, only one medial (straight-line),
which is commensurable in square only with the whole, and
contains a rational (area) with the whole, can be attached
to a first apotome of a medial (straight-line). (Which is) the very thing it was required to
show.}
\end{Parallel}
{\footnotesize\noindent$^\dag$ This proposition is equivalent to 
Prop.~10.43, with minus signs instead of
plus signs.}

%%%%
%10.81
%%%%
\pdfbookmark[1]{Proposition 10.81}{pdf10.81}
\begin{Parallel}{}{}
\ParallelLText{
\begin{center}
{\large \ggn{81}.}
\end{center}\vspace*{-7pt}

{\greekfont Τῇ μέσης ἀποτομ ῇ δευτέρᾳ μία μόνον προσαρμόζει εὐθεῖα
μέση δυνάμει μόνον σύμμετρος τῇ ὅλῃ, μετὰ δὲ τῆς ὅλης
μέσον περιέθουσα.}\\

\epsfysize=1.6in
\centerline{\epsffile{Book10/fig081g.eps}}

{\greekfont Ἔστω μέσης ἀποτομὴ δευτέρα ἡ ΑΒ καὶ τῇ ΑΒ προσαρμόζουσα
ἡ ΒΓ; αἱ ἄρα ΑΓ, ΓΒ μέσαι εἰσὶ δυνάμει μόνον σύμμετροι
μέσον περιέθουσαι τὸ ὑπὸ τῶν ΑΓ, ΓΒ; λέγω, ὅτι τῇ ΑΒ
ἑτέρα οὐ προσαρμόσει εὐθεῖα μέση δυνάμει μόνον σύμμετρος
οὖσα τῇ ὅλῃ, μετὰ δὲ τῆς ὅλης μέσον περιέθουσα.}

{\greekfont Εἰ γὰρ δυνατόν, προσαρμοζέτω ἡ ΒΔ; καὶ αἱ ΑΔ,
ΔΒ ἄρα μέσαι εἰσὶ δυνάμει μόνον σύμμετροι μέσον περιέθουσαι
τὸ ὑπὸ τῶν ΑΔ, ΔΒ. καὶ ἐκκείσθω ῥητὴ ἡ ΕΖ, καὶ τοῖς μὲν
ἀπὸ τῶν ΑΓ, ΓΒ ἴσον παρὰ τὴν ΕΖ παραβεβλήσθω τὸ ΕΗ πλάτος ποιοῦν
τὴν ΕΜ; τῷ δὲ δὶς ὑπὸ τῶν ΑΓ, ΓΒ ἴσον ἀφῃρήσθω τὸ ΘΗ
πλάτος ποιοῦν τὴν ΘΜ; λοιπὸν ἄρα τὸ ΕΛ ἴσον ἐστὶ τῷ
ἀπὸ τῆς ΑΒ; ὥστε ἡ ΑΒ δύναται τὸ ΕΛ. πάλιν δὴ τοῖς
ἀπὸ τῶν ΑΔ, ΔΒ ἴσον παρὰ τὴν ΕΖ παραβεβλήσθω τὸ ΕΙ πλάτος
ποιοῦν τὴν ΕΝ; ἔστι δὲ καὶ τὸ ΕΛ ἴσον τῷ ἀπὸ τῆς ΑΒ
τετραγώνῳ; λοιπὸν ἄρα τὸ ΘΙ ἴσον ἐστὶ τῷ δὶς ὑπὸ τῶν ΑΔ, ΔΒ.
καὶ ἐπεὶ μέσαι εἰσὶν αἱ ΑΓ, ΓΒ, μέσα ἄρα ἐστὶ καὶ τὰ
ἀπὸ τῶν ΑΓ, ΓΒ. καί ἐστιν ἴσα τῷ ΕΗ; μέσον ἄρα καὶ τὸ
ΕΗ. καὶ παρὰ ῥητὴν τὴν ΕΖ παράκειται πλάτος ποιοῦν τὴν ΕΜ;
ῥητὴ ἄρα ἐστὶν ἡ ΕΜ καὶ ἀσύμμετρος τῇ ΕΖ μήκει. 
πάλιν,
ἐπεὶ μέσον ἐστὶ τὸ ὑπὸ τῶν ΑΓ, ΓΒ, καὶ τὸ δὶς ὑπὸ τῶν
ΑΓ, ΓΒ μέσον ἐστίν. καί ἐστιν ἴσον τῷ ΘΗ; καὶ τὸ ΘΗ ἄρα
μέσον ἐστίν. καὶ παρὰ ῥητὴν τὴν ΕΖ παράκειται πλάτος ποιοῦν
τὴν ΘΜ; ῥητὴ ἄρα ἐστὶ καὶ ἡ ΘΜ καὶ ἀσύμμετρος τῇ ΕΖ
μήκει. 
καὶ ἐπεὶ αἱ ΑΓ, ΓΒ δυνάμει μόνον σύμμετροί εἰσιν,
ἀσύμμετρος ἄρα ἐστὶν ἡ ΑΓ τῇ ΓΒ μήκει. ὡς δὲ ἡ ΑΓ πρὸς
τὴν ΓΒ, οὕτως ἐστὶ τὸ ἀπὸ τῆς ΑΓ πρὸς τὸ ὑπὸ τῶν ΑΓ, ΓΒ;
ἀσύμμετρον ἄρα ἐστὶ τὸ ἀπὸ τῆς ΑΓ  τῷ ὑπὸ
τῶν ΑΓ, ΓΒ. ἀλλὰ τῷ μὲν ἀπὸ τῆς ΑΓ
σύμμετρά ἐστι τὰ ἀπὸ
τῶν ΑΓ, ΓΒ, τῷ δὲ ὑπὸ τῶν ΑΓ, ΓΒ σύμμετρόν ἐστι τὸ δὶς
ὑπὸ τῶν ΑΓ, ΓΒ; ἀσύμμετρα ἄρα ἐστὶ τὰ ἀπὸ τῶν ΑΓ, ΓΒ
τῷ δὶς ὑπὸ τῶν ΑΓ, ΓΒ. καί ἐστι τοῖς μὲν ἀπὸ τῶν
ΑΓ, ΓΒ ἴσον τὸ ΕΗ, τῷ δὲ δὶς ὑπὸ τῶν ΑΓ, ΓΒ ἴσον τὸ ΗΘ;
 ἀσύμμετρον ἄρα ἐστὶ τὸ ΕΗ τῷ ΘΗ.
ὡς δὲ τὸ ΕΗ
πρὸς τὸ ΘΗ, οὕτως ἐστὶν ἡ ΕΜ πρὸς τὴν ΘΜ; 
ἀσύμμετρος
ἄρα ἐστὶν ἡ ΕΜ τῇ ΜΘ μήκει. καί εἰσιν ἀμφότεραι ῥηταί;
αἱ ΕΜ, ΜΘ ἄρα ῥηταί εἰσι δυνάμει μόνον σύμμετροι;
ἀποτομὴ ἄρα ἐστὶν ἡ ΕΘ, προσαρμόζουσα δὲ α ὐτῇ 
ἡ ΘΜ. ὁμοίως δὴ δείχομεν, ὅτι καὶ ἡ ΘΝ α ὐτῇ 
προσαρμόζει;
τῇ ἄρα ἀποτομ ῇ ἄλλη καὶ ἄλλη προσαρμόζει εὐθεῖα δυνάμει
μόνον σύμμετρος οὖσα τῇ ὅλῃ; ὅπερ ἐστὶν ἀδύνατον.}

{\greekfont Τῇ ἄρα μέσης ἀποτομ ῇ δευτέρᾳ μία μόνον προσαρμόζει εὐθεῖα
μέση δυνάμει μόνον σύμμετρος οὖσα τῇ ὅλῃ, μετὰ δὲ τῆς
ὅλης μέσον περιέθουσα; ὅπερ ἔδει δεῖχαι.}}

\ParallelRText{
\begin{center}
{\large Proposition 81}
\end{center}

Only one medial straight-line,
which is commensurable in square only with the whole, and
contains a medial (area) with the whole, can be attached
to a second apotome of a medial (straight-line).$^\dag$

\epsfysize=1.6in 
\centerline{\epsffile{Book10/fig081e.eps}}

Let $AB$ be a second apotome of a medial (straight-line), with $BC$ (so) attached
to $AB$. Thus, $AC$ and $CB$ are medial (straight-lines which
are) commensurable in square only, containing a medial (area)---(namely,
that contained) by $AC$ and $CB$ [Prop.~10.75].
I say that a(nother) medial straight-line, which is commensurable in square
only with the whole, and contains a medial (area) with the whole, cannot
be attached to $AB$.

For, if possible, let $BD$ be (so) attached. Thus, $AD$ and $DB$ are also medial (straight-lines which
are) commensurable in square only, containing a medial (area)---(namely,
that contained) by $AD$ and $DB$ [Prop.~10.75].
And let the rational (straight-line) $EF$ be laid down. And let $EG$,
equal to the (sum of the squares) on $AC$ and $CB$, be applied
to $EF$, producing $EM$ as breadth. And let $HG$, equal to twice
the (rectangle contained) by $AC$ and $CB$, be subtracted
(from $EG$), producing $HM$ as breadth. The remainder $EL$ is
thus equal to the (square) on $AB$ [Prop.~2.7]. 
Hence, $AB$ is the square-root of $EL$. So, again, let $EI$, equal to
the (sum of the squares) on $AD$ and $DB$ be applied to $EF$,
producing $EN$ as breadth. And $EL$ is also equal to the square
on $AB$. Thus, the remainder 
$HI$ is equal to twice the (rectangle contained) by $AD$ and $DB$ [Prop.~2.7]. And since $AC$ and $CB$ are
(both) medial (straight-lines), the (sum of the squares) on $AC$ and $CB$ is also medial.
And it is equal to $EG$. Thus, $EG$ is also medial [Props.~10.15, 10.23~corr.]. And it is
applied to the rational (straight-line) $EF$, producing $EM$ as breadth.
Thus, $EM$ is rational, and incommensurable in length with $EF$
[Prop.~10.22]. Again, since the (rectangle
contained) by $AC$ and $CB$ is medial, twice the (rectangle
contained) by $AC$ and $CB$ is also medial [Prop.~10.23~corr.]. And it is equal to $HG$.
Thus, $HG$ is also medial. And it is applied to the rational (straight-line)
$EF$, producing $HM$ as breadth. Thus, $HM$ is also rational, and
incommensurable in length with $EF$ [Prop.~10.22]. And since $AC$ and $CB$ are commensurable in square only, $AC$ is thus incommensurable in length with $CB$.  And as $AC$ (is)
to $CB$, so the (square) on $AC$
is to the (rectangle contained) by $AC$ and $CB$ [Prop.~10.21~corr.]. Thus, the (square) on $AC$
is incommensurable with the (rectangle contained) by $AC$ and $CB$
[Prop.~10.11].
But, the (sum of the squares) on $AC$ and $CB$ is commensurable
with the (square) on $AC$, and twice the (rectangle contained) by $AC$ and
$CB$ is commensurable with the (rectangle contained) by $AC$ and $CB$ [Prop.~10.6].
Thus, the (sum of the squares) on $AC$ and $CB$ is incommensurable
with twice the (rectangle contained) by $AC$ and $CB$ [Prop.~10.13]. And $EG$ is equal to
the (sum of the squares) on $AC$ and $CB$. And $GH$ is equal
to twice the (rectangle contained) by $AC$ and $CB$. Thus, 
$EG$ is incommensurable with $HG$. And as $EG$ (is) to $HG$,
so $EM$ is to $HM$ [Prop.~6.1]. 
Thus, $EM$ is incommensurable in length with $MH $ [Prop.~10.11]. And they are both rational (straight-lines).
Thus, $EM$ and $MH $ are rational (straight-lines which are)
commensurable in square only. Thus, $EH$ is an apotome
[Prop.~10.73], and $HM$ (is) attached to
it. So, similarly, we can show that $HN$ (is) also (commensurable
in square only with $EN$ and is) attached to ($EH$). 
Thus, different straight-lines, which are commensurable
in square only with the whole, are attached to an apotome.
The very thing is impossible [Prop.~10.79].

Thus, only one medial straight-line,
which is commensurable in square only with the whole, and
contains a medial (area) with the whole, can be attached
to a second apotome of a medial (straight-line). (Which is) the very thing it was required to show.}
\end{Parallel}
{\footnotesize\noindent$^\dag$ This proposition is equivalent to 
Prop.~10.44, with minus signs instead of
plus signs.}

%%%%
%10.82
%%%%
\pdfbookmark[1]{Proposition 10.82}{pdf10.82}
\begin{Parallel}{}{}
\ParallelLText{
\begin{center}
{\large \ggn{82}.}
\end{center}\vspace*{-7pt}

{\greekfont Τῇ ἐλάσσονι μία μόνον προσαρμόζει εὐθεῖα δυνάμει
ἀσύμμετρος οὖσα τῇ ὅλῃ ποιοῦσα μετὰ τῆς ὅλης
τὸ μὲν ἐκ τῶν ἀπ α ὐτῶν τετραγώνων ῥητόν,
τὸ δὲ δὶς ὑπ α ὐτῶν μέσον.}\\

\epsfysize=0.25in
\centerline{\epsffile{Book10/fig079g.eps}}

{\greekfont Ἔστω ἡ ἐλάσσων ἡ ΑΒ, καὶ τῇ ΑΒ προσαρμόζουσα ἔστω
ἡ ΒΓ; αἱ ἄρα ΑΓ, ΓΒ δυνάμει εἰσὶν ἀσύμμετροι
ποιοῦσαι τὸ μὲν συγκείμενον ἐκ τῶν ἀπ α ὐτῶν τετραγώνων
ῥητόν, τὸ δὲ δὶς ὑπ α ὐτῶν μέσον; λέγω, ὅτι
τῇ ΑΒ ἑτέρα εὐθεῖα οὐ προσαρμόσει τὰ α ὐτὰ ποιοῦσα.}

{\greekfont Εἰ γὰρ δυνατόν, προσαρμοζέτω ἡ ΒΔ; καὶ αἱ ΑΔ,
ΔΒ ἄρα δυνάμει εἰσὶν ἀσύμμετροι ποιοῦσαι τὰ προειρημένα.
καὶ ἐπεί, ᾧ ὑπερέθει τὰ ἀπὸ τῶν ΑΔ, ΔΒ
τῶν ἀπὸ τῶν ΑΓ, ΓΒ, τούτῳ ὑπερέθει καὶ τὸ δὶς
ὑπὸ τῶν ΑΔ, ΔΒ τοῦ δὶς ὑπὸ τῶν ΑΓ, ΓΒ, τὰ δὲ
ἀπὸ τῶν ΑΔ, ΔΒ τετράγωνα τῶν ἀπὸ τῶν ΑΓ, ΓΒ
 τετραγώνων ὑπερέθει
ῥητῷ; ῥητὰ γάρ ἐστιν ἀμφότερα; καὶ τὸ δὶς ὑπὸ τῶν ΑΔ, ΔΒ
ἄρα τοῦ δὶς ὑπὸ τῶν ΑΓ, ΓΒ ὑπερέθει ῥητῷ; ὅπερ
ἐστὶν ἀδύνατον; μέσα γάρ ἐστιν ἀμφότερα.}

{\greekfont Τῇ ἄρα ἐλάσσονι μία μόνον προσαρμόζει εὐθεῖα δυνάμει
ἀσύμμετρος οὖσα τῇ ὅλῃ καὶ ποιοῦσα τὰ μὲν ἀπ 
α ὐτῶν τετράγωνα ἅμα ῥητόν, τὸ δὲ δὶς ὑπ α ὐτῶν
μέσον; ὅπερ ἔδει δεῖχαι.}}

\ParallelRText{
\begin{center}
{\large Proposition 82}
\end{center}

Only one straight-line, which is incommensurable
in square with the whole, and (together) with the whole makes the (sum of the) squares on them
rational, and twice the (rectangle contained) by them medial, can be
attached to a minor (straight-line).

\epsfysize=0.25in 
\centerline{\epsffile{Book10/fig079e.eps}}

Let $AB$ be a minor (straight-line), and let $BC$ be (so) attached to $AB$.
Thus, $AC$ and $CB$ are (straight-lines which are) incommensurable  in square, making the sum of
the squares on them rational, and twice the (rectangle contained) by them
medial [Prop.~10.76]. I say that another
another straight-line fulfilling the same (conditions) cannot be
attached to $AB$.

For, if possible, let $BD$ be (so) attached (to $AB$). Thus, $AD$
and $DB$ are also (straight-lines which are) incommensurable in square, fulfilling the (other) aforementioned
(conditions) [Prop.~10.76]. And since
by whatever (area) the (sum of the squares) on $AD$ and $DB$
exceeds the (sum of the squares) on $AC$ and $CB$, 
twice the (rectangle contained) by $AD$ and $DB$ also exceeds
twice the (rectangle contained) by $AC$ and $CB$  by this (same area) [Prop.~2.7]. And the (sum of the) squares on
$AD$ and $DB$ exceeds the (sum of the) squares on $AC$ and $CB$
by a rational (area). For both are rational (areas). Thus, twice the
(rectangle contained) by $AD$ and $DB$ also exceeds
twice the (rectangle contained) by $AC$ and $CB$ by a rational (area).
The very thing is impossible. For both are medial (areas) [Prop.~10.26].

Thus,  only one straight-line, which is  incommensurable in square
with the whole, and (with the whole) makes the  squares on them
(added) together rational, and twice the (rectangle contained) by them medial, can be
attached to a minor (straight-line). (Which is) the very thing it was required to
show.}
\end{Parallel}
{\footnotesize\noindent$^\dag$ This proposition is equivalent to 
Prop.~10.45, with minus signs instead of plus signs.}

%%%%
%10.83
%%%%
\pdfbookmark[1]{Proposition 10.83}{pdf10.83}
\begin{Parallel}{}{}
\ParallelLText{
\begin{center}
{\large \ggn{83}.}
\end{center}\vspace*{-7pt}

{\greekfont Τῇ μετὰ ῥητοῦ μέσον τὸ ὅλον ποιούσῃ μία μόνον προσαρμόζει
εὐθεῖα δυνάμει ἀσύμμετρος οὖσα τῇ ὅλῃ, μετὰ δὲ τῆς
ὅλης ποιοῦσα τὸ μὲν συγκείμενον ἐκ τῶν ἀπ
α ὐτῶν τετραγώνων μέσον, τὸ δὲ δὶς ὑπ α ὐτῶν ῥητόν.}\\~\\

\epsfysize=0.25in
\centerline{\epsffile{Book10/fig079g.eps}}

{\greekfont Ἔστω ἡ μετὰ ῥητοῦ μέσον τὸ ὅλον ποιοῦσα ἡ ΑΒ,
καὶ τῇ ΑΒ προσαρμοζέτω ἡ ΒΓ; αἱ ἄρα ΑΓ, ΓΒ δυνάμει
εἰσὶν ἀσύμμετροι ποιοῦσαι τὰ προκείμενα; λέγω, ὅτι
τῇ ΑΒ ἑτέρα οὐ προσαρμόσει τὰ α ὐτὰ ποιοῦσα.}

{\greekfont Εἰ γὰρ δυνατόν, προσαρμοζέτω ἡ ΒΔ; καὶ αἱ ΑΔ, ΔΒ ἄρα
εὐθεῖαι δυνάμει εἰσὶν ἀσύμμετροι ποιοῦσαι τὰ
προκείμενα. ἐπεὶ οὐν, ᾧ ὑπερέθει τὰ ἀπὸ τῶν
ΑΔ, ΔΒ τῶν ἀπὸ τῶν ΑΓ, ΓΒ, τούτῳ ὑπερέθει καὶ τὸ
δὶς ὑπὸ τῶν ΑΔ, ΔΒ τοῦ δὶς ὑπὸ τῶν ΑΓ, ΓΒ
ἀκολούθως τοῖς πρὸ α ὐτοῦ, τὸ δὲ δὶς ὑπὸ τῶν ΑΔ, ΔΒ
τοῦ δὶς ὑπὸ τῶν ΑΓ, ΓΒ ὑπερέθει ῥητῷ; ῥητὰ γάρ
ἐστιν ἀμφότερα; καὶ τὰ ἀπὸ τῶν ΑΔ, ΔΒ ἄρα τῶν
ἀπὸ τῶν ΑΓ, ΓΒ ὑπερέθει ῥητῷ; ὅπερ ἐστὶν
ἀδύνατον; μέσα γάρ ἐστιν ἀμφότερα.}

 {\greekfont Οὐκ ἄρα τῇ ΑΒ
ἑτέρα προσαρμόσει εὐθεῖα δυνάμει ἀσύμμετρος οὖσα τῇ
ὅλῃ, μετὰ δὲ τῆς ὅλης ποιοῦσα τὰ προειρημένα;
μία ἄρα μόνον προσαρμόσει; ὅπερ ἔδει δεῖχαι.}}

\ParallelRText{
\begin{center}
{\large Proposition 83}
\end{center}

Only  one straight-line, which is incommensurable
in square with the whole, and (together) with the whole makes the sum of the squares on them
medial, and twice the (rectangle contained) by them rational, can be
attached to that (straight-line) which with a rational (area) makes a
medial whole.$^\dag$

\epsfysize=0.25in 
\centerline{\epsffile{Book10/fig079e.eps}}

Let $AB$ be a (straight-line) which with a rational (area) makes a medial
whole,  and let $BC$ be (so) attached to $AB$. Thus, $AC$ and $CB$
are (straight-lines which are) incommensurable in square, fulfilling the (other) proscribed (conditions) [Prop.~10.77]. I say that another (straight-line)
fulfilling the same (conditions) cannot be attached to $AB$.

For, if possible, let $BD$ be (so) attached (to $AB$). Thus, $AD$ and
$DB$ are also straight-lines (which are) incommensurable in square,
fulfilling the (other) prescribed  (conditions) [Prop.~10.77]. Therefore, analogously to the
(propositions) before this, since by whatever (area)
the (sum of the squares) on $AD$ and $DB$ exceeds the (sum of the
squares) on $AC$ and $CB$,  twice the (rectangle contained)
by $AD$ and $DB$ also exceeds twice the (rectangle contained) by
$AC$ and $CB$ by this (same area). And twice the (rectangle contained) by $AD$ and
$DB$ exceeds twice the (rectangle contained) by $AC$ and
$CB$ by a rational (area). For they are (both) rational (areas).
Thus, the (sum of the squares) on 
$AD$ and $DB$ also exceeds the (sum of the squares) on $AC$ and $CB$
by a rational (area). The very thing is impossible. For both are
medial (areas) [Prop.~10.26].

 Thus, another
straight-line cannot be attached to $AB$, which is incommensurable
in square with the whole, and fulfills the (other) aforementioned (conditions) with the
whole. Thus, only one (such straight-line) can be (so) attached. (Which is)
the very thing it was required to show.}
\end{Parallel}
{\footnotesize\noindent$^\dag$
This proposition is equivalent to Prop.~10.46, with minus signs instead of
plus signs.}

%%%%
%10.84
%%%%
\pdfbookmark[1]{Proposition 10.84}{pdf10.84}
\begin{Parallel}{}{}
\ParallelLText{
\begin{center}
{\large \ggn{84}.}
\end{center}\vspace*{-7pt}

{\greekfont Τῇ μετὰ μέσου μέσον τὸ ὅλον ποιούσῃ μία μόνη προσαρμόζει εὐθεῖα δυνάμει ἀσύμμετρος οὖσα τῇ ὅλῃ, μετὰ δὲ τῆς ὅλης ποιοῦσα τό τε συγκείμενον ἐκ τῶν ἀπ α ὐτῶν τετραγώνων
μέσον τό τε δὶς ὑπ α ὐτῶν μέσον καὶ ἔτι ἀσύμμετρον
τῷ συγκειμένῳ ἐκ τῶν ἀπ α ὐτῶν.}

{\greekfont Ἔστω ἡ μετὰ μέσου μέσον τὸ ὅλον ποιοῦσα ἡ ΑΒ,
προσαρμόζουσα δὲ α ὐτῇ ἡ ΒΓ; αἱ ἄρα ΑΓ, ΓΒ δυνάμει εἰσὶν
ἀσύμμετροι ποιοῦσαι τὰ προειρημένα. λέγω, ὅτι τῇ ΑΒ ἑτέρα
οὐ προσαρμόσει ποιοῦσα προειρημένα.}\\~\\~\\~\\

\epsfysize=1.6in
\centerline{\epsffile{Book10/fig081g.eps}}

{\greekfont Εἰ γὰρ δυνατόν, προσαρμοζέτω ἡ ΒΔ, ὥστε καὶ τὰς ΑΔ, ΔΒ δυνάμει
ἀσυμμέτρους εἶναι ποιούσας τά τε ἀπὸ τῶν ΑΔ, ΔΒ τετράγωνα
ἅμα μέσον καὶ τὸ δὶς ὑπὸ τῶν ΑΔ, ΔΒ μέσον καὶ ἔτι τὰ
ἀπὸ τῶν ΑΔ, ΔΒ ἀσύμμετρα τῷ δὶς ὑπὸ τῶν ΑΔ, ΔΒ; καὶ
ἐκκείσθω ῥητὴ ἡ ΕΖ, καὶ τοῖς μὲν ἀπὸ τῶν ΑΓ, ΓΒ ἴσον
παρὰ τὴν ΕΖ παραβεβλήσθω τὸ ΕΗ πλάτος ποιοῦν τὴν ΕΜ, τῷ
δὲ δὶς ὑπὸ τῶν ΑΓ, ΓΒ ἴσον παρὰ τὴν ΕΖ παραβεβλήσθω τὸ ΘΗ
πλάτος ποιοῦν τὴν ΘΜ; λοιπὸν
ἄρα τὸ ἀπὸ τῆς ΑΒ ἴσον ἐστὶ τῷ ΕΛ; ἡ ἄρα ΑΒ δύναται
τὸ ΕΛ. πάλιν τοῖς ἀπὸ τῶν ΑΔ, ΔΒ ἴσον
παρὰ τὴν ΕΖ παραβεβλήσθω τὸ ΕΙ πλάτος ποιοῦν τὴν ΕΝ.
 ἔστι
δὲ καὶ τὸ ἀπὸ τῆς ΑΒ ἴσον τῷ ΕΛ;
λοιπὸν ἄρα τὸ δὶς ὑπὸ τῶν ΑΔ, ΔΒ ἴσον [ἐστὶ]
τῷ ΘΙ. καὶ ἐπεὶ μέσον ἐστὶ τὸ συγκείμενον ἐκ τῶν
ἀπὸ τῶν ΑΓ, ΓΒ καί ἐστιν ἴσον τῷ ΕΗ, μέσον ἄρα ἐστὶ
καὶ τὸ ΕΗ. καὶ παρὰ ῥητὴν τὴν ΕΖ παράκειται πλάτος
ποιοῦν τὴν ΕΜ;
 ῥητὴ ἄρα ἐστὶν ἡ ΕΜ καὶ ἀσύμμετρος τῇ
ΕΖ μήκει. πάλιν, ἐπεὶ μέσον ἐστὶ τὸ δὶς ὑπὸ τῶν ΑΓ, ΓΒ καί
ἐστιν ἴσον τῷ ΘΗ, μέσον ἄρα καὶ τὸ ΘΗ. καὶ παρὰ ῥητὴν τὴν
ΕΖ παράκειται πλάτος ποιοῦν τὴν ΘΜ; 
ῥητὴ ἄρα ἐστὶν ἡ ΘΜ
καὶ ἀσύμμετρος τῇ ΕΖ μήκει. καὶ ἐπεὶ ἀσύμμετρά ἐστι
τὰ ἀπὸ τῶν ΑΓ, ΓΒ τῷ δὶς ὑπὸ τῶν ΑΓ, ΓΒ, ἀσύμμετρόν
ἐστι καὶ τὸ ΕΗ τῷ ΘΗ; ἀσύμμετρος ἄρα ἐστὶ καὶ ἡ ΕΜ τῇ
ΜΘ μήκει. καί εἰσιν ἀμφότεραι ῥηταί; αἱ ἄρα ΕΜ, ΜΘ ῥηταί
εἰσι δυνάμει μόνον σύμμετροι; ἀποτομὴ ἄρα ἐστὶν ἡ ΕΘ, προσαρμόζουσα δὲ α ὐτῇ ἡ ΘΜ. ὁμοίως δὴ δείχομεν, ὅτι ἡ ΕΘ πάλιν ἀποτομή
ἐστιν, προσαρμόζουσα δὲ α ὐτῇ ἡ ΘΝ. τῇ ἄρα ἀποτομ ῇ
ἄλλη καὶ ἄλλη προσαρμόζει ῥητὴ δυνάμει μόνον σύμμετρος
οὖσα τῇ ὅλῃ; ὅπερ ἐδείθθη ἀδύνατον. οὐκ ἄρα τῇ ΑΒ ἑτέρα
προσαρμόσει εὐθεῖα.}

{\greekfont Τῇ ἄρα ΑΒ μία μόνον προσαρμόζει εὐθεῖα δυνάμει
ἀσύμμετρος οὖσα τῇ ὅλῃ, μετὰ δὲ τῆς ὅλης ποιοῦσα
τά τε ἀπ α ὐτῶν τετράγωνα ἅμα μέσον καὶ τὸ δὶς ὑπ
α ὐτῶν μέσον καὶ ἔτι τὰ ἀπ α ὐτῶν τετράγωνα ἀσύμμετρα
τῷ δὶς ὑπ α ὐτῶν; ὅπερ ἔδει δεῖχαι.}}

\ParallelRText{
\begin{center}
{\large Proposition 84}
\end{center}

Only one straight-line, which is incommensurable
in square with the whole, and (together) with the whole makes the sum of the
squares on them medial, and twice the (rectangle contained) by
them medial, and, moreover, incommensurable with the sum of
the (squares) on them, can be attached to that (straight-line) which with
a medial (area) makes a medial whole.$^\dag$

Let $AB$ be a (straight-line) which with a medial (area) makes a medial
whole, $BC$ being (so) attached to it. Thus, $AC$ and $CB$ are incommensurable in square, fulfilling the (other) aforementioned (conditions)
[Prop.~10.78]. I say that a(nother) (straight-line)
fulfilling the aforementioned (conditions) cannot be attached to $AB$.

\epsfysize=1.6in 
\centerline{\epsffile{Book10/fig081e.eps}}

For, if possible, let $BD$ be (so) attached. Hence, $AD$ and
$DB$ are also (straight-lines which are) incommensurable in square, making the
squares on $AD$ and $DB$ (added) together medial, and twice the (rectangle contained) by $AD$ and $DB$ medial, and, moreover, the (sum of the squares)
on $AD$ and $DB$ incommensurable with twice the (rectangle contained)
by $AD$ and $DB$ [Prop.~10.78].
And let the rational (straight-line) $EF$ be laid down. And let $EG$, equal
to the (sum of the squares) on $AC$ and $CB$, be applied to
$EF$, producing $EM$ as breadth. And let $HG$, equal to twice the (rectangle contained) by $AC$ and $CB$, be applied to $EF$, producing $HM$ as breadth. Thus, the remaining (square) on  $AB$ is equal
to $EL$ [Prop.~2.7].  Thus, $AB$ is the square-root
of $EL$. Again, let $EI$, equal to the (sum of the squares) on $AD$ and $DB$, be applied to $EF$, producing $EN$ as breadth. And the (square) on $AB$
 is also equal to $EL$. Thus, the remaining twice the
(rectangle contained) by $AD$ and $DB$ [is] equal to $HI$ [Prop.~2.7].  And since the sum of the (squares) on $AC$ and $CB$ is medial, and is equal to $EG$, $EG$ is thus also medial.
And it is applied to the rational (straight-line) $EF$, producing $EM$ as breadth. $EM$ is thus rational, and incommensurable in length with $EF$
[Prop.~10.22]. Again, since twice the
(rectangle contained) by $AC$ and $CB$ is medial, and is equal to $HG$,
$HG$ is thus also medial. And it is applied to the rational (straight-line) $EF$,
producing $HM$ as breadth. $HM$ is thus rational, and incommensurable in
length with $EF$ [Prop.~10.22]. 
And since the (sum of the squares) on
 $AC$ and $CB$ is incommensurable
with twice the (rectangle contained) by $AC$ and $CB$, $EG$ is also incommensurable with $HG$. Thus, $EM$ is also incommensurable in length
with $MH $ [Props.~6.1, 10.11]. And they are both  rational (straight-lines). Thus, $EM$ and $MH $ are rational (straight-lines which are) commensurable in square
only. Thus, $EH$ is an apotome [Prop.~10.73], with $HM$ attached to it. So, similarly, we can show that $EH$ is again an apotome, with $HN$ attached to it. Thus, different rational (straight-lines),
which are commensurable in square only with the whole, are attached to
an apotome. The very thing was shown (to be) impossible [Prop.~10.79]. Thus, another straight-line cannot
be (so) attached to $AB$.

Thus, only one straight-line, which is incommensurable
in square with the whole, and (together) with the whole makes the 
squares on them (added) together medial, and twice the (rectangle contained) by
them medial, and, moreover, the (sum of the)
squares on them incommensurable with the (rectangle contained)
by them, can be attached to $AB$. (Which is) the very thing it
was required to show.}
\end{Parallel}
{\footnotesize\noindent$^\dag$ This proposition is equivalent to 
Prop.~10.47, with minus signs instead of
plus signs.}

%%%%%%%%%
% Definitions III
%%%%%%%%%
\pdfbookmark[1]{Definitions III}{def10.3}
\begin{Parallel}{}{}
\ParallelLText{
\begin{center}
\large{{\greekfont Ὅροι τρίτοι}.}
\end{center}\vspace*{-7pt}

\ggn{11}.~{\greekfont Ὑποκειμένης ῥητῆς καὶ ἀποτομῆς, ἐὰν μὲν ἡ ὅλη
τῆς προσαρμοζούσης μεῖζον δύνηται τῷ ἀπὸ συμμέτρου ἑαυτῇ
μήκει, καὶ ἡ ὅλη σύμμετρος ᾖ τῇ ἐκκειμένῃ ῥητῇ μήκει,
καλείσθω ἀποτομὴ πρώτη.}

\ggn{12}.~{\greekfont Ἐὰν δὲ ἡ προσαρμόζουσα σύμμετρος ᾖ τῇ ἐκκειμένῃ
ῥητῇ μήκει, καὶ ἡ ὅλη τῆς προσαρμοζούσης μεῖζον δύνηται
τῷ ἀπὸ συμμέτρου ἑαυτῇ, καλείσθω ἀποτομὴ δευτέρα.}

\ggn{13}.~{\greekfont Ἐὰν δὲ μ ηδετέρα σύμμετρος ᾖ τῇ ἐκκειμένῃ
ῥητῇ μήκει, ἡ δὲ ὅλη τῆς προσαρμοζούσης μεῖζον δύνηται
τῷ ἀπὸ συμμέτρου ἑαυτῇ, καλείσθω ἀποτομὴ τρίτη.}

\ggn{14}.~{\greekfont Πάλιν, ἐὰν ἡ ὅλη τῆς προσαρμοζούσης μεῖζον
δύνηται τῷ ἀπὸ ἀσυμμέτρου ἑαυτῇ [μήκει], ἐὰν
μὲν ἡ ὅλη σύμμετρος ᾖ τῇ ἐκκειμένῃ ῥητῇ
μήκει, καλείσθω ἀποτομὴ τετάρτη.}

\ggn{15}.~{\greekfont Ἐὰν δὲ ἡ προσαρμόζουσα, πέμπτη.}

\ggn{16}.~{\greekfont Ἐὰν δὲ μ ηδετέρα, ἕκτη.}}

\ParallelRText{
\begin{center}
{\large Definitions III}
\end{center}

11.~Given a rational (straight-line) and an apotome, if the square
on the whole is greater than the (square on a straight-line) attached (to the apotome) by the (square)
on (some straight-line) commensurable in length  with (the whole), and the
whole is commensurable in length with the 
(previously) laid down rational (straight-line), then let the (apotome) be
called a first apotome.

12.~And if the attached (straight-line)
is commensurable in length with the (previously) laid down rational
(straight-line), and the square on the whole is greater than (the
square on) the attached (straight-line) by the (square) on (some
straight-line) commensurable (in length) with (the whole), then let
the (apotome) be called a second apotome.

13.~And if neither of (the whole
or the attached straight-line) is commensurable
in length with the (previously) laid down rational (straight-line),
and the square on the whole is greater than (the
square on) the attached (straight-line) by the (square) on (some
straight-line) commensurable (in length) with (the whole), then let
the (apotome) be called a third apotome.

14.~Again, if the square
on the whole is greater than (the square on) the attached (straight-line) by the (square)
on (some straight-line) incommensurable [in length]  with (the whole), and the
whole is commensurable in length with the
(previously) laid down rational (straight-line), then let the (apotome) be
called a fourth apotome.

15.~And if the
attached (straight-line is commensurable),   a fifth (apotome).

16.~And if neither (the whole
nor the attached straight-line is commensurable), a sixth (apotome).}
\end{Parallel}

%%%%
%10.85
%%%%
\pdfbookmark[1]{Proposition 10.85}{pdf10.85}
\begin{Parallel}{}{}
\ParallelLText{
\begin{center}
{\large \ggn{85}.}
\end{center}\vspace*{-7pt}

{\greekfont Εὑρεῖν τὴν πρώτην ἀποτομήν.}

\epsfysize=0.65in
\centerline{\epsffile{Book10/fig085g.eps}}

{\greekfont Ἐκκείσθω ῥητὴ ἡ Α, καὶ τῇ Α μήκει σύμμετρος ἔστω ἡ ΒΗ;
ῥητὴ ἄρα ἐστὶ καὶ ἡ ΒΗ. καὶ ἐκκείσθωσαν δύο τετράγωνοι
ἀριθμοὶ οἱ ΔΕ, ΕΖ, ὧν ἡ ὑπεροθὴ ὁ ΖΔ μὴ ἔστω τετράγωνος;
οὐδ ἄρα ὁ ΕΔ πρὸς τὸν ΔΖ λόγον ἔχει, ὃν τετράγωνος
ἀριθμὸς πρὸς τετράγωνον ἀριθμόν. καὶ πεποιήσθω ὡς ὁ
ΕΔ πρὸς τὸν ΔΖ, οὕτως τὸ ἀπὸ τῆς ΒΗ τετράγωνον πρὸς τὸ ἀπὸ
τὴς ΗΓ τετράγωνον; σύμμετρον ἄρα ἐστὶ τὸ ἀπὸ τῆς ΒΗ τῷ
ἀπὸ τῆς ΗΓ. ῥητὸν δὲ τὸ ἀπὸ τῆς ΒΗ; ῥητὸν ἄρα καὶ τὸ
ἀπὸ τῆς ΗΓ; ῥητὴ ἄρα ἐστὶ καὶ ἡ ΗΓ. καὶ ἐπεὶ ὁ ΕΔ πρὸς 
τὸν ΔΖ λόγον οὐκ ἔχει, ὃν τετράγωνος ἀριθμὸς πρὸς τετράγωνον
ἀριθμόν, οὐδ ἄρα τὸ ἀπὸ τῆς ΒΗ πρὸς τὸ ἀπὸ τῆς ΗΓ
λόγον ἔχει, ὃν τετράγωνος ἀριθμὸς πρὸς τετράγωνον
ἀριθμόν; ἀσύμμετρος ἄρα ἐστὶν ἡ ΒΗ τῇ ΗΓ μήκει.
καί εἰσιν ἀμφότεραι ῥηταί; αἱ ΒΗ, ΗΓ ἄρα ῥηταί εἰσι
δυνάμει μόνον σύμμετροι; ἡ ἄρα ΒΓ ἀποτομή ἐστιν. λέγω
δή, ὅτι καὶ πρώτη.}

{\greekfont Ὧι γὰρ μεῖζόν ἐστι τὸ ἀπὸ τῆς ΒΗ τοῦ ἀπὸ τῆς ΗΓ, ἔστω
τὸ ἀπὸ τῆς Θ. καὶ ἐπεί ἐστιν ὡς ὁ ΕΔ πρὸς τὸν ΖΔ,
οὕτως τὸ ἀπὸ τῆς ΒΗ πρὸς τὸ ἀπὸ τῆς ΗΓ, καὶ ἀναστρέψαντι
ἄρα ἐστὶν ὡς ὁ ΔΕ πρὸς τὸν ΕΖ, οὕτως τὸ ἀπὸ τῆς ΗΒ
πρὸς τὸ ἀπὸ τῆς Θ. ὁ δὲ ΔΕ πρὸς τὸν ΕΖ λόγον ἔχει, ὃν
τετράγωνος ἀριθμὸς πρὸς τετράγωνον ἀριθμόν; ἑκάτερος
γὰρ τετράγωνός ἐστιν; καὶ τὸ ἀπὸ τῆς ΗΒ ἄρα πρὸς τὸ
ἀπὸ τῆς Θ λόγον ἔχει, ὃν τετράγωνος ἀριθμὸς πρὸς
τετράγωνον ἀριθμόν; σύμμετρος ἄρα ἐστὶν ἡ ΒΗ τῇ Θ μήκει.
καὶ δύναται ἡ ΒΗ τῆς ΗΓ μεῖζον τῷ ἀπὸ τῆς Θ; ἡ ΒΗ ἄρα
τῆς ΗΓ μεῖζον δύναται τῷ ἀπὸ συμμέτρου ἑαυτῇ μήκει.
καί ἐστιν ἡ ὅλη ἡ ΒΗ σύμμετρος τῇ ἐκκειμένῃ ῥητῇ μήκει
τῇ Α. ἡ ΒΓ ἄρα ἀποτομή ἐστι πρώτη.}

{\greekfont Εὕρηται ἄρα ἡ πρώτη ἀποτομὴ ἡ ΒΓ; ὅπερ ἔδει εὑρεῖν.}}

\ParallelRText{
\begin{center}
{\large Proposition 85}
\end{center}

To find a first apotome.

\epsfysize=0.65in
\centerline{\epsffile{Book10/fig085e.eps}}

Let the rational (straight-line) $A$ be laid down. And let $BG$
be commensurable in length with $A$. $BG$ is thus also a rational (straight-line).
And let two square numbers
$DE$ and $EF$ be laid down, and let their difference $FD$ be  not 
square [Prop.~10.28~lem.~I]. Thus,
$ED$ does not have to $DF$ the ratio which (some) square number
(has) to (some) square number. And let it be contrived that
as $ED$ (is) to $DF$, so the square on $BG$ (is) to the square on $GC$
[Prop.~10.6.~corr.]. Thus, the (square) on $BG$
is commensurable with the (square) on $GC$ [Prop.~10.6]. And the (square) on $BG$ (is)
rational. Thus, the (square) on $GC$ (is) also rational. Thus, $GC$ is also
rational. And since $ED$ does not have to $DF$ the ratio which
(some) square number (has) to (some) square number, 
the (square) on $BG$ thus does not have to the (square) on $GC$ the ratio
which (some) square number (has) to (some) square number either. Thus,
$BG$ is incommensurable in length with $GC$ [Prop.~10.9]. And they are both
rational (straight-lines). Thus, $BG$ and $GC$ are rational (straight-lines
which are) commensurable in square only. Thus, $BC$ is an apotome
[Prop.~10.73]. So, I say that (it is)
also a first (apotome).

Let the (square) on $H$ be that (area) by which the (square) on $BG$
is greater than the (square) on $GC$ [Prop.~10.13~lem.]. And since as $ED$ is to
$FD$, so the (square) on $BG$ (is) to the (square) on $GC$, thus, via
conversion, as $DE$ is to $EF$, so the (square) on $GB$ (is) to the
(square) on $H$ [Prop.~5.19~corr.]. 
And $DE$ has to $EF$ the ratio which (some) square-number
(has) to (some) square-number. For each is a square (number). 
Thus, the (square) on $GB$ also has to the (square) on $H$ the ratio which
(some) square number (has) to (some) square number. Thus, $BG$
is commensurable in length with $H$ [Prop.~10.9].
And the square on $BG$ is greater than (the square on) $GC$ by the
(square) on $H$. Thus, the square on $BG$ is greater than (the square on) $GC$ by the
(square) on (some straight-line) commensurable in length with ($BG$).
And the whole, $BG$, is commensurable in length with the (previously)
laid down rational (straight-line) $A$. Thus, $BC$ is a first
apotome [Def.~10.11].

Thus, the first apotome $BC$ has been found. (Which is) the
very thing it was required to find.}
\end{Parallel}
{\footnotesize\noindent$^\dag$ See footnote to Prop.~10.48.}

%%%%
%10.86
%%%%
\pdfbookmark[1]{Proposition 10.86}{pdf10.86}
\begin{Parallel}{}{}
\ParallelLText{
\begin{center}
{\large \ggn{86}.}
\end{center}

{\greekfont Εὑρεῖν τὴν δευτέραν ἀποτομήν.}

{\greekfont Ἐκκείσθω ῥητὴ ἡ Α καὶ τῇ Α σύμμετρος μήκει ἡ ΗΓ. ῥητὴ
ἄρα ἐστὶν ἡ ΗΓ. καὶ ἐκκείσθωσαν δύο τετράγωνοι ἀριθμοὶ
οἱ ΔΕ, ΕΖ, ὧν ἡ ὑπεροθὴ ὁ ΔΖ μὴ ἔστω τετράγωνος. καὶ
πεποιήσθω ὡς ὁ ΖΔ πρὸς τὸν ΔΕ, οὕτως τὸ ἀπὸ τῆς ΓΗ τετράγωνον
πρὸς τὸ ἀπὸ τῆς ΗΒ τετράγωνον. σύμμετρον ἄρα ἐστὶ τὸ ἀπὸ
τῆς ΓΗ τετράγωνον τῷ ἀπὸ τῆς ΗΒ τετραγώνῳ. ῥητὸν δὲ τὸ
ἀπὸ τῆς ΓΗ. ῥητὸν ἄρα [ἐστὶ] καὶ τὸ ἀπὸ τῆς ΗΒ;
ῥητὴ ἄρα ἐστὶν ἡ ΒΗ. καὶ ἐπεὶ τὸ ἀπὸ τῆς ΗΓ τετράγωνον
πρὸς τὸ ἀπὸ τῆς ΗΒ λόγον οὐκ ἔχει, ὃν τετράγωνος ἀριθμὸς πρὸς
τετράγωνον ἀριθμόν, ἀσύμμετρός ἐστιν ἡ ΓΗ τῇ ΗΒ
μήκει. καί εἰσιν ἀμφότεραι ῥηταί; αἱ ΓΗ, ΗΒ ἄρα  ρηταί εἰσι
δυνάμει μόνον σύμμετροι; ἡ ΒΓ ἄρα ἀποτομή ἐστιν.
λέγω δή, ὅτι καὶ δευτέρα.}\\~\\~\\~\\

\epsfysize=0.65in
\centerline{\epsffile{Book10/fig086g.eps}}

{\greekfont Ὧι γὰρ μεῖζόν ἐστι τὸ ἀπὸ τῆς ΒΗ τοῦ ἀπὸ τῆς
ΗΓ, ἔστω τὸ ἀπὸ τῆς Θ. ἐπεὶ οὖν ἐστιν ὡς τὸ ἀπὸ
τῆς ΒΗ πρὸς τὸ ἀπὸ τῆς ΗΓ, οὕτως ὁ ΕΔ ἀριθμὸς
πρὸς τὸν ΔΖ ἀριθμόν, ἀναστρέψαντι ἄρα ἐστὶν ὡς τὸ ἀπὸ
τῆς ΒΗ πρὸς τὸ ἀπὸ τῆς Θ, οὕτως ὁ ΔΕ πρὸς τὸν ΕΖ.  καί
ἐστιν ἑκάτερος τῶν ΔΕ, ΕΖ τετράγωνος; τὸ ἄρα ἀπὸ τῆς ΒΗ
πρὸς τὸ ἀπὸ τῆς Θ λόγον ἔχει, ὃν τετράγωνος ἀριθμὸς
πρὸς τετράγωνον ἀριθμόν; σύμμετρος ἄρα ἐστὶν ἡ ΒΗ τῇ Θ
μήκει. καὶ δύναται ἡ ΒΗ τῆς ΗΓ μεῖζον τῷ ἀπὸ τῆς Θ;
ἡ ΒΗ ἄρα τῆς ΗΓ μεῖζον δύναται τῷ ἀπὸ συμμέτρου
ἑαυτῇ μήκει. καί ἐστιν ἡ προσαρμόζουσα ἡ ΓΗ τῇ
ἐκκειμένῃ ῥητῇ σύμμετρος τῇ Α. ἡ ΒΓ ἄρα ἀποτομή
ἐστι δευτέτα.}

{\greekfont Εὕρηται ἄρα δευτέρα ἀποτομὴ ἡ ΒΓ;
ὅπερ ἔδει δεῖχαι.}}

\ParallelRText{
\begin{center}
{\large Proposition 86}
\end{center}\vspace*{-7pt}

To find a second apotome.

Let the rational (straight-line) $A$, and $GC$ (which is)
commensurable in length with $A$, be laid down. Thus, $GC$ is a rational (straight-line).
And let the two square numbers $DE$ and $EF$ be laid down, and let
their
difference $DF$ be not square [Prop.~10.28~lem.~I]. 
And let it be contrived that as $FD$ (is) to $DE$, so the square on $CG$ (is) to the square on $GB$ [Prop.~10.6~corr.].
Thus, the square on $CG$ is commensurable with the square on $GB$ [Prop.~10.6]. And the (square) on $CG$ (is) rational.
Thus, the (square) on $GB$ [is] also rational. Thus, $BG$ is a rational
(straight-line).  And since the square on $GC$ does not have to the
(square) on $GB$ the ratio which (some) square number (has) to
(some) square number, $CG$ is incommensurable in length with $GB$
[Prop.~10.9]. And they are both rational (straight-lines). Thus, $CG$ and $GB$ are rational (straight-lines which are)
commensurable in square only. Thus, $BC$ is an apotome [Prop.~10.73]. So, I say that it is also a second
(apotome).

\epsfysize=0.65in
\centerline{\epsffile{Book10/fig086e.eps}}

For let the (square) on $H$ be that (area) by which the (square) on $BG$
is greater than the (square) on $GC$ [Prop.~10.1 lem.]. Therefore, since
as the (square) on $BG$ is to the (square) on $GC$, so the number $ED$
(is) to the number $DF$, thus, also, via conversion, as the (square) on $BG$
is to the (square) on $H$, so $DE$ (is) to $EF$ [Prop.~5.19~corr.]. And $DE$ and $EF$ are each
square (numbers). Thus, the (square) on $BG$ has to the (square)
on $H$ the ratio which (some) square number (has) to (some) square number.
Thus, $BG$ is commensurable in length with $H$ [Prop.~10.9]. 
And the square on $BG$ is greater than (the square on) $GC$ by the (square)
on $H$.
Thus, the square on $BG$ is greater
than (the square on) $GC$ by the (square) on (some straight-line)
commensurable in length with ($BG$). And the attachment $CG$
is commensurable (in length) with the (prevously) laid down rational (straight-line) $A$.
Thus, $BC$ is a second apotome [Def.~10.12].$^\dag$

Thus, the second apotome $BC$ has been found. (Which is) the very thing
it was required to show.}
\end{Parallel}
{\footnotesize\noindent$^\dag$ See footnote to Prop.~10.49.}

%%%%
%10.87
%%%%
\pdfbookmark[1]{Proposition 10.87}{pdf10.87}
\begin{Parallel}{}{}
\ParallelLText{
\begin{center}
{\large \ggn{87}.}
\end{center}\vspace*{-7pt}

{\greekfont Εὑρεῖν τὴν τρίτην ἀποτομήν.}

\epsfysize=1.in
\centerline{\epsffile{Book10/fig087g.eps}}

{\greekfont Ἐκκείσθω ῥητὴ ἡ Α, καὶ ἐκκείσθωσαν τρεῖς ἀριθμοὶ
οἱ Ε, ΒΓ, ΓΔ λόγον μὴ ἔθοντες πρὸς ἀλλήλους, ὃν
τετράγωνος ἀριθμὸς πρὸς τετράγωνον ἀριθμόν, ὁ δὲ
ΓΒ πρὸς τὸν ΒΔ λόγον ἐθέτω, ὃν τετράγωνος ἀριθμὸς πρὸς
τετράγωνον ἀριθμόν, καὶ πεποιήσθω ὡς μὲν ὁ Ε πρὸς τὸν
ΒΓ, οὕτως τὸ ἀπὸ τῆς Α τετράγωνον πρὸς τὸ ἀπὸ τῆς
ΖΗ τετράγωνον, ὡς δὲ ὁ ΒΓ πρὸς τὸν ΓΔ, οὕτως τὸ ἀπὸ
τῆς ΖΗ τετράγωνον πρὸς τὸ ἀπὸ τὴς ΗΘ. ἐπεὶ οὖν ἐστιν
ὡς ὁ Ε πρὸς τὸν ΒΓ, οὕτως τὸ ἀπὸ τῆς Α τετράγωνον πρὸς
τὸ ἀπὸ τῆς ΖΗ τετράγωνον, σύμμετρον ἄρα ἐστὶ τὸ ἀπὸ
τῆς Α τετράγωνον τῷ ἀπὸ τῆς ΖΗ τετραγώνῳ. ῥητὸν δὲ
τὸ ἀπὸ τῆς Α τετράγωνον. ῥητὸν ἄρα καὶ τὸ ἀπὸ τῆς ΖΗ;
ῥητὴ ἄρα ἐστὶν ἡ ΖΗ. καὶ ἐπεὶ ὁ Ε πρὸς τὸν ΒΓ
λόγον οὐκ ἔχει, ὃν τετράγωνος ἀριθμὸς πρὸς τετράγωνον
ἀριθμόν, οὐδ ἄρα τὸ ἀπὸ τῆς Α
τετράγωνον πρὸς τὸ ἀπὸ τῆς ΖΗ [τετράγωνον] λόγον ἔχει,
ὅν τετράγωνος ἀριθμὸς πρὸς τετράγωνον ἀριθμόν; ἀσύμμετρος
ἄρα ἐστὶν ἡ Α τῇ ΖΗ μήκει. πάλιν, ἐπεί ἐστιν ὡς ὁ ΒΓ
πρὸς τὸν ΓΔ, οὕτως τὸ ἀπὸ τῆς ΖΗ τετράγωνον πρὸς τὸ ἀπὸ
τῆς ΗΘ, σύμμετρον ἄρα ἐστὶ τὸ ἀπὸ τῆς ΖΗ τῷ ἀπὸ τῆς ΗΘ.
ῥητὸν δὲ τὸ ἀπὸ τῆς ΖΗ; ῥητὸν ἄρα καὶ τὸ ἀπὸ τῆς ΗΘ; ῥητὴ
ἄρα ἐστὶν ἡ ΗΘ. καὶ ἐπεὶ ὁ ΒΓ πρὸς τὸν ΓΔ λόγον οὐκ
ἔχει, ὃν τετράγωνος ἀριθμὸς πρὸς τετράγωνον ἀριθμόν, οὐδ
ἄρα τὸ ἀπὸ τῆς ΖΗ πρὸς τὸ ἀπὸ τῆς ΗΘ λόγον ἔχει, ὃν
τετράγωνος ἀριθμὸς πρὸς τετράγωνον ἀριθμόν;
ἀσύμμετρος ἄρα ἐστὶν ἡ ΖΗ τῇ ΗΘ μήκει. καί εἰσιν ἀμφότεραι
ῥηταί; αἱ ΖΗ, ΗΘ ἄρα ῥηταί εἰσι δυνάμει μόνον σύμμετροι;
ἀποτομὴ ἄρα ἐστὶν ἡ ΖΘ. λέγω δή, ὅτι καὶ τρίτη.}

{\greekfont Ἐπεὶ γάρ ἐστιν ὡς μὲν ὁ Ε πρὸς τὸν ΒΓ, οὕτως τὸ ἀπὸ τῆς
Α τετράγωνον πρὸς τὸ ἀπὸ τῆς ΖΗ, ὡς δὲ ὁ ΒΓ πρὸς τὸν ΓΔ,
οὕτως τὸ ἀπὸ τῆς ΖΗ πρὸς  τὸ ἀπὸ τῆς ΘΗ, δι ἴσου ἄρα
ἐστὶν ὡς ὁ Ε πρὸς τὸν ΓΔ, οὕτως τὸ ἀπὸ τῆς Α πρὸς τὸ
ἀπὸ τῆς ΘΗ. ὁ δὲ Ε πρὸς τὸν ΓΔ λόγον οὐκ ἔχει, ὃν
τετράγωνος ἀριθμὸς πρὸς τετράγωνον ἀριθμόν; οὐδ ἄρα τὸ
ἀπὸ τῆς Α πρὸς τὸ ἀπὸ τῆς ΗΘ λόγον ἔχει, ὃν τετράγωνος
ἀριθμὸς
πρὸς τετράγωνον ἀριθμόν; ἀσύμμετρος ἄρα ἡ Α τῇ ΗΘ
μήκει. οὐδετέρα ἄρα τῶν ΖΗ, ΗΘ σύμμετρός ἐστι τῇ ἐκκειμένῃ
ῥητῇ τῇ Α μήκει. ᾧ οὖν μεῖζόν ἐστι τὸ ἀπὸ τῆς ΖΗ
τοῦ ἀπὸ τῆς ΗΘ, ἔστω τὸ ἀπὸ τῆς Κ. ἐπεὶ οὖν ἐστιν ὡς
ὁ ΒΓ πρὸς τὸν ΓΔ, οὕτως τὸ ἀπὸ τῆς ΖΗ πρὸς τὸ ἀπὸ τῆς
ΗΘ, ἀναστρέψαντι ἄρα ἐστὶν ὡς ὁ ΒΓ πρὸς τὸν ΒΔ, οὕτως
τὸ ἀπὸ τῆς ΖΗ τετράγωνον πρὸς τὸ ἀπὸ τῆς Κ. ὁ δὲ
ΒΓ πρὸς τὸν ΒΔ λόγον ἔχει, ὃν τετράγωνος ἀριθμὸς πρὸς
τετράγωνον ἀριθμόν. καὶ τὸ ἁπὸ τῆς ΖΗ ἄρα πρὸς τὸ ἀπὸ
τῆς Κ λόγον ἔχει, ὃν τετράγωνος ἀριθμὸς πρὸς τετράγωνον ἀριθμόν.
σύμμετρός ἄρα ἐστὶν ἡ ΖΗ τῇ Κ μήκει,
καὶ δύναται ἡ ΖΗ τῆς ΗΘ μεῖζον τῷ ἀπὸ συμμέτρου ἑαυτῇ.
καὶ οὐδετέρα τῶν ΖΗ, ΗΘ σύμμετρός ἐστι τῇ ἐκκειμένῃ ῥητῇ
τῇ Α μήκει; ἡ ΖΘ ἄρα ἀποτομή ἐστι τρίτη.}

{\greekfont Εὕρηται ἄρα ἡ τρίτη ἀποτομὴ ἡ ΖΘ; ὅπερ
ἔδει δεῖχαι.}}

\ParallelRText{
\begin{center}
{\large Proposition 87}
\end{center}

To find a third apotome.

\epsfysize=1.in
\centerline{\epsffile{Book10/fig087e.eps}}

Let the rational (straight-line) $A$ be laid down. And let the
three numbers, $E$, $BC$, and $CD$, not having to one another
the ratio which (some) square number (has) to (some) square number,
be laid down. And let $CB$ have to $BD$ the ratio which (some)
square number (has) to (some) square number. And let it be
contrived that as $E$ (is) to $BC$, so the square on $A$ (is) to the square
on $FG$, and as $BC$ (is) to $CD$, so the square on $FG$ (is) to the (square)
on $GH$ [Prop.~10.6~corr.]. Therefore,
since  as $E$ is to $BC$, so the square on $A$ (is) to the square on
$FG$, the square on $A$ is thus commensurable with the square on
$FG$ [Prop.~10.6]. And the square on $A$
(is) rational. Thus, the (square) on $FG$ (is) also rational. Thus,
$FG$ is a rational (straight-line). And since  $E$ does not have to $BC$
the ratio which (some) square number (has) to (some) square number,
the square on $A$ thus does not have to the [square] on $FG$ the ratio
which (some) square number (has) to (some) square number either. Thus,
$A$ is incommensurable in length with $FG$ [Prop.~10.9]. Again, since as $BC$ is to $CD$, so
the square on $FG$ is to the (square) on $GH$, the square on $FG$
is thus commensurable with the (square) on $GH$ [Prop.~10.6]. And the (square) on $FG$ (is)
rational. Thus, the (square) on $GH$ (is) also rational. Thus, $GH$ is
a rational (straight-line). And since $BC$ does not have to $CD$ the
ratio which (some) square number (has) to (some) square number, the
(square) on $FG$ thus does not have to the (square) on $GH$ the
ratio which (some) square number (has) to (some) square number either.
Thus, $FG$ is incommensurable in length with $GH$ [Prop.~10.9]. And both are rational (straight-lines).
$FG$ and $GH$ are thus rational (straight-lines which are) commensurable
in square only. Thus, $FH$ is an apotome [Prop.~10.73]. So, I say that (it is) also a third
(apotome).

For since as $E$ is to $BC$, so the square on $A$ (is) to the (square)
on $FG$, and as $BC$ (is) to $CD$, so the (square) on $FG$
(is) to the (square) on $HG$, thus, via equality, as $E$ is to
$CD$, so the (square) on $A$ (is) to the (square) on $HG$ [Prop.~5.22]. And $E$ does not have to $CD$
the ratio which (some) square number (has) to (some) square number. Thus,
the (square) on $A$ does not have to the (square) on $GH$ the ratio
which (some) square number (has) to (some) square number either. 
$A$ (is) thus incommensurable in length with $GH$ [Prop.~10.9]. Thus, neither of $FG$ and $GH$
is commensurable in length with the (previously) laid down rational
(straight-line) $A$. Therefore, let the (square) on $K$ be that (area) by which the (square) on
$FG$ is greater than the (square) on $GH$ [Prop.~10.13~lem.]. Therefore, since as $BC$ is to
$CD$, so the (square) on $FG$ (is) to the (square) on $GH$, thus,
via conversion, as $BC$ is to $BD$, so the square on $FG$ (is) to the square on
$K$ [Prop.~5.19~corr.]. And $BC$ has to $BD$ the ratio which
(some) square number (has) to (some) square number. Thus, the (square)
on $FG$ also has to the (square) on $K$ the ratio which (some)
square number (has) to (some) square number.
$FG$ is thus commensurable in length with $K$ [Prop.~10.9]. And the square on $FG$ is (thus) greater
than (the square on) $GH$ by the (square) on (some straight-line)
commensurable (in length) with ($FG$). And neither of $FG$ and
$GH$ is commensurable in length with the (previously) laid down
rational (straight-line) $A$. Thus, $FH$ is a third
apotome [Def.~10.13].

Thus, the third apotome $FH$ has been found. (Which is) very thing
it was required to show.}
\end{Parallel}
{\footnotesize\noindent$^\dag$ See footnote
to Prop.~10.50.}

%%%%
%10.88
%%%%
\pdfbookmark[1]{Proposition 10.88}{pdf10.88}
\begin{Parallel}{}{}
\ParallelLText{
\begin{center}
{\large\ggn{88}.}
\end{center}\vspace*{-7pt}

{\greekfont Εὑρεῖν τὴν τετάρτην ἀποτομήν.}

\epsfysize=0.65in
\centerline{\epsffile{Book10/fig088g.eps}}

{\greekfont Ἐκκείσθω ῥητὴ ἡ Α καὶ τῇ Α μήκει σύμμετρος ἡ ΒΗ;
ῥητὴ ἄρα ἐστὶ καὶ ἡ ΒΗ. καὶ ἐκκείσθωσαν δύο ἀριθμοὶ οἱ ΔΖ, ΖΕ,
ὥστε τὸν ΔΕ ὅλον πρὸς ἑκάτερον τῶν ΔΖ, ΕΖ λόγον μὴ
ἔχειν, ὃν τετράγωνος ἀριθμὸς πρὸς τετράγωνον ἀριθμόν. καὶ
πεποιήσθω ὡς ὁ ΔΕ πρὸς τὸν ΕΖ, οὕτως τὸ ἀπὸ τῆς ΒΗ
τετράγωνον πρὸς τὸ ἀπὸ τῆς ΗΓ; σύμμετρον ἄρα ἐστὶ τὸ
ἀπὸ τῆς ΒΗ τῷ ἀπὸ τῆς ΗΓ.
 ῥητὸν δὲ τὸ ἀπὸ τῆς ΒΗ; 
 ῥητὸν ἄρα καὶ τὸ ἀπὸ τῆς ΗΓ; 
 ῥητὴ ἄρα ἐστὶν ἡ ΗΓ.
καὶ ἐπεὶ ὁ ΔΕ πρὸς τὸν ΕΖ λόγον οὐκ ἔχει, ὃν τετράγωνος ἀριθμὸς πρὸς τετράγωνον ἀριθμόν, οὐδ ἄρα τὸ ἀπὸ τῆς
ΒΗ πρὸς τὸ ἀπὸ τῆς ΗΓ λόγον ἔχει, ὃν τετράγωνος ἀριθμὸς
πρὸς τετράγωνον ἀριθμόν;
ἀσύμμετρος ἄρα ἐστὶν ἡ ΒΗ τῇ ΗΓ μήκει. καί εἰσιν ἀμφότεραι
ῥηταί; αἱ ΒΗ, ΗΓ ἄρα ῥηταί εἰσι δυνάμει μόνον σύμμετροι;
ἀποτομὴ ἄρα ἐστὶν ἡ ΒΓ.  [λέγω δή, ὅτι καὶ τετάρτη.]}

{\greekfont Ὧι οὖν μεῖζόν ἐστι τὸ ἀπὸ τῆς ΒΗ τοῦ ἀπὸ τῆς
ΗΓ, ἔστω τὸ ἀπὸ τῆς Θ. ἐπεὶ οὖν ἐστιν ὡς ὁ ΔΕ πρὸς
τὸν ΕΖ, οὕτως τὸ ἀπὸ τῆς ΒΗ πρὸς τὸ ἀπὸ τῆς
ΗΓ, καὶ ἀναστρέψαντι ἄρα ἐστὶν ὡς ὁ ΕΔ πρὸς τὸν ΔΖ,
οὕτως τὸ ἀπὸ τῆς ΗΒ πρὸς τὸ  ἀπὸ τῆς Θ.
ὁ δὲ ΕΔ πρὸς τὸν ΔΖ λόγον οὐκ ἔχει, ὃν τετράγωνος
ἀριθμὸς πρὸς τετράγωνον ἀριθμόν; οὐδ ἄρα τὸ ἀπὸ τῆς ΗΒ
πρὸς τὸ ἀπὸ τῆς Θ λόγον ἔχει, ὃν τετράγωνος ἀριθμὸς
πρὸς τετράγωνον ἀριθμόν;
 ἀσύμμετρος ἄρα ἐστὶν
ἡ ΒΗ τῇ Θ μήκει. καὶ δύναται ἡ ΒΗ τῆς ΗΓ μεῖζον τῷ
ἀπὸ τῆς Θ; ἡ ἄρα ΒΗ τῆς ΗΓ μεῖζον δύναται
τῷ ἀπὸ ἀσυμμέτρου ἑαυτῇ.
καὶ ἐστιν ὅλη ἡ ΒΗ
 σύμμετρος τῇ ἐκκειμένῃ ῥητῇ μήκει τῇ Α.
ἡ ἄρα ΒΓ ἀποτομή ἐστι τετάρτη.}

{\greekfont Εὕρηται ἄρα ἡ τετάρτη ἀποτομή; ὅπερ ἔδει δεῖχαι.}}

\ParallelRText{
\begin{center}
{\large Proposition 88}
\end{center}

To find a fourth apotome.

\epsfysize=0.65in
\centerline{\epsffile{Book10/fig088e.eps}}

Let the rational (straight-line) $A$, and $BG$ (which is)
commensurable in length with $A$, be laid down.
Thus, $BG$ is also a rational (straight-line). And let the two numbers $DF$
and $FE$ be laid down such that the whole, $DE$, does not have to
each of $DF$ and $EF$ the ratio which (some) square number (has)
to (some) square number. And let it be contrived that as $DE$
(is) to $EF$, so the square on $BG$ (is) to the (square) on $GC$
[Prop.~10.6~corr.]. The (square) on $BG$ is
thus commensurable with the (square) on $GC$ [Prop.~10.6]. And the (square) on $BG$ (is)
rational. Thus, the (square) on $GC$ (is) also rational. Thus, $GC$
(is) a rational (straight-line). And since $DE$ does not have to $EF$
the ratio which (some) square number (has) to (some) square number,
the (square) on $BG$ thus does not have to the (square) on $GC$ the
ratio which (some) square number (has) to (some) square number either.
Thus, $BG$ is incommensurable in length with $GC$ [Prop.~10.9]. And they are both rational (straight-lines). Thus, $BG$ and $GC$ are rational (straight-lines which are)
commensurable in square only. Thus, $BC$ is an apotome [Prop.~10.73]. [So, I say that
(it is) also a fourth (apotome).]

Now, let the (square) on $H$ be that (area) by which the (square)
on $BG$ is greater than the (square) on $GC$ [Prop.~10.13~lem.]. Therefore, since as
$DE$ is to $EF$, so the (square) on $BG$ (is) to the (square)
on $GC$, thus, also, via conversion, as $ED$ is to $DF$, so the
(square) on $GB$ (is) to the (square) on $H$ [Prop.~5.19~corr.]. And $ED$ does not have to
$DF$ the ratio which (some) square number (has) to (some) square number.
Thus, the (square) on $GB$ does not have to the (square) on $H$
the ratio which (some) square number (has) to (some) square number either.
Thus, $BG$ is incommensurable in length with $H$ [Prop.~10.9]. And the square on $BG$ is greater
than (the square on) $GC$ by the (square) on $H$. Thus, the
square on $BG$ is greater than (the square) on $GC$ by the (square)
on (some straight-line) incommensurable (in length) with ($BG$). And the whole,
$BG$, is commensurable in length with the the (previously) laid
down rational (straight-line) $A$. Thus, $BC$ is a fourth apotome
[Def.~10.14].$^\dag$

Thus, a fourth apotome has been found. (Which is) the very thing it was required to show.}
\end{Parallel}
{\footnotesize\noindent$^\dag$ See footnote to Prop.~10.51.}

%%%%
%10.89
%%%%
\pdfbookmark[1]{Proposition 10.89}{pdf10.89}
\begin{Parallel}{}{}
\ParallelLText{
\begin{center}
{\large \ggn{89}.}
\end{center}\vspace*{-7pt}

{\greekfont Εὐρεῖν τὴν πέμπτην ἀποτομήν.}

\epsfysize=0.6in
\centerline{\epsffile{Book10/fig089g.eps}}

{\greekfont Ἐκκείσθω ῥητὴ ἡ Α, καὶ τῇ Α μήκει σύμμετρος ἔστω ἡ ΓΗ;
ῥητὴ ἄρα [ἐστὶν] ἡ ΓΗ. καὶ ἐκκείσθωσαν δύο ἀριθμοὶ
οἱ ΔΖ, ΖΕ, ὥστε τὸν ΔΕ πρὸς ἑκάτερον τῶν ΔΖ, ΖΕ λόγον
πάλιν μὴ ἔχειν, ὃν τετράγωνος ἀριθμὸς πρὸς τετράγωνον
ἀριθμόν; 
καὶ πεποιήσθω ὡς ὁ ΖΕ πρὸς τὸν ΕΔ, οὕτως τὸ ἀπὸ τῆς
ΓΗ πρὸς τὸ ἀπὸ τῆς ΗΒ. 
ῥητὸν ἄρα καὶ τὸ ἀπὸ τῆς ΗΒ; 
ῥητὴ ἄρα ἐστὶ καὶ ἡ ΒΗ. 
καὶ
ἐπεί ἐστιν ὡς ὁ ΔΕ πρὸς τὸν ΕΖ, οὕτως τὸ ἀπὸ τῆς ΒΗ
πρὸς τὸ ἀπὸ τῆς ΗΓ, ὁ δὲ ΔΕ πρὸς τὸν ΕΖ λόγον οὐκ
ἔχει, ὃν τετράγωνος ἀριθμὸς πρὸς τετράγωνον ἀριθμόν,
οὐδ ἄρα τὸ ἀπὸ τῆς ΒΗ πρὸς τὸ ἀπὸ τῆς ΗΓ λόγον
ἔχει, ὃν τετράγωνος ἀριθμὸς πρὸς τετράγωνον ἀριθμόν;
ἀσύμμετρος ἄρα ἐστὶν ἡ ΒΗ τῇ ΗΓ μήκει. καί εἰσιν
ἀμφότεραι ῥηταί; αἱ ΒΗ, ΗΓ ἄρα ῥηταί εἰσι δυνάμει
μόνον σύμμετροι; ἡ ΒΓ ἄρα ἀποτομή ἐστιν. λέγω δή,
ὅτι καὶ πέμπτη.}

{\greekfont Ὧι γὰρ μεῖζόν ἐστι τὸ ἀπὸ τῆς ΒΗ τοῦ ἀπὸ τῆς
ΗΓ, ἔστω τὸ ἀπὸ τῆς Θ. ἐπεὶ οὖν ἐστιν ὡς τὸ ἀπὸ
τὴς ΒΗ πρὸς τὸ ἀπὸ τῆς ΗΓ, οὕτως ὁ ΔΕ πρὸς τὸν ΕΖ,
ἀναστρέψαντι ἄρα ἐστὶν  ὡς ὁ ΕΔ πρὸς τὸν ΔΖ, οὕτως
τὸ ἀπὸ τῆς ΒΗ πρὸς τὸ ἀπὸ τῆς Θ, ὁ δὲ ΕΔ πρὸς τὸν ΔΖ
λόγον οὐκ ἔχει, ὃν τετράγωνος ἀριθμὸς πρὸς τετράγωνον
ἀριθμόν; οὐδ ἄρα τὸ ἀπὸ τῆς ΒΗ πρὸς τὸ ἀπὸ τῆς Θ
λόγον ἔχει, ὃν τετράγωνος ἀριθμὸς πρὸς τετράγωνον ἀριθμόν;
ἀσύμμετρος ἄρα ἐστὶν ἡ ΒΗ τῇ Θ μήκει. καὶ δύναται
ἡ ΒΗ τῆς ΗΓ μεῖζον τῷ ἀπὸ τῆς Θ; ἡ ΗΒ ἄρα τῆς ΗΓ
μεῖζον δύναται τῷ ἀπὸ ἀσυμμέτρου ἑαυτῇ μήκει. καί
ἐστιν ἡ προσαρμόζουσα ἡ ΓΗ σύμμετρος τῇ ἐκκειμένῃ ῥητῇ
τῇ Α μήκει; ἡ ἄρα ΒΓ ἀποτομή ἐστι πέμπτη.}

{\greekfont Εὕρηται ἄρα ἡ πέμπτη ἀποτομὴ ἡ ΒΓ;
ὅπερ ἔδει δεῖχαι.}}

\ParallelRText{
\begin{center}
{\large Proposition 89}
\end{center}

To find a fifth apotome.

\epsfysize=0.6in
\centerline{\epsffile{Book10/fig089e.eps}}

Let the rational (straight-line) $A$ be laid down, and let
$CG$ be commensurable in length with $A$. Thus, $CG$ [is] a rational
(straight-line). And let the two numbers $DF$ and $FE$ be laid
down such that $DE$ again does not have to each of $DF$ and $FE$
the ratio which (some) square number (has) to (some) square number. 
And let it be contrived that as $FE$ (is) to $ED$, so the
(square) on $CG$ (is) to the (square) on $GB$. Thus, the (square) on
$GB$ (is) also rational [Prop.~10.6]. 
Thus, $BG$ is also rational.
And since as $DE$ is to $EF$, so the (square) on $BG$ (is) to
the (square) on $GC$. And $DE$ does not have to $EF$ the
ratio which (some) square number (has) to (some) square number.
The (square) on $BG$ thus does not have to the (square) on $GC$
the ratio which (some) square number (has) to (some) square number either.
Thus, $BG$ is incommensurable in length with $GC$ [Prop.~10.9]. And they are both rational (straight-lines). $BG$ and $GC$ are thus rational (straight-lines which are)
commensurable in square only. Thus, $BC$ is an apotome [Prop.~10.73]. So, I say that (it is) also a fifth
(apotome).

For, let the (square) on $H$ be that (area) by which the (square) on $BG$
is greater than the (square) on $GC$ [Prop.~10.13~lem.]. Therefore, since
as the (square) on $BG$ (is) to the (square) on $GC$, so $DE$  (is) to
$EF$, thus, via conversion, as $ED$ is to $DF$, so the (square)
on $BG$ (is) to the (square) on $H$ [Prop.~5.19~corr.]. And $ED$ does not have to
$DF$ the ratio which (some) square number (has) to (some) square number.
Thus, the (square) on $BG$ does not have to the (square) on $H$
the ratio which (some) square number (has) to (some) square number
either. Thus, $BG$ is incommensurable in length with $H$ [Prop.~10.9]. And the square on $BG$ is greater
than (the square on) $GC$ by the (square) on $H$.
Thus, the square on $GB$ is greater
than (the square on) $GC$ by the (square) on (some straight-line)
incommensurable in length with ($GB$). And the attachment $CG$ is
commensurable in length with the (previously) laid down rational (straight-line) $A$. Thus, $BC$ is a fifth apotome [Def.~10.15].$^\dag$

Thus, the fifth apotome $BC$ has been found. (Which is) the very thing it
was required to show.}
\end{Parallel}
{\footnotesize\noindent$^\dag$ See footnote to Prop.~10.52.}

%%%%
%10.90
%%%%
\pdfbookmark[1]{Proposition 10.90}{pdf10.90}
\begin{Parallel}{}{}
\ParallelLText{
\begin{center}
{\large \ggn{90}.}
\end{center}

{\greekfont Εὑρεῖν τὴν ἕκτην ἀποτομήν.}

{\greekfont Ἐκκείσθω ῥητὴ ἡ Α καὶ τρεῖς ἀριθμοὶ οἱ Ε, ΒΓ, ΓΔ
λόγον μὴ ἔθοντες πρὸς ἀλλήλους, ὃν τετράγωνος ἀριθμὸς
πρὸς τετράγωνον ἀριθμόν; ἔτι δὲ καὶ ὁ ΓΒ πρὸς τὸν ΒΔ λόγον
μὴ ἐθετώ, ὃν τετράγωνος ἀριθμὸς πρὸς τετράγωνον ἀριθμόν;
καὶ πεποιήσθω ὡς μὲν ὁ Ε πρὸς τὸν ΒΓ, οὕτως τὸ ἀπὸ τῆς
Α πρὸς τὸ ἀπὸ τῆς ΖΗ, ὡς δὲ ὁ ΒΓ πρὸς τὸν ΓΔ, οὕτως τὸ ἀπὸ
τῆς ΖΗ πρὸς τὸ ἀπὸ τῆς ΗΘ.}\\

\epsfysize=1.1in
\centerline{\epsffile{Book10/fig090g.eps}}

{\greekfont Ἐπεὶ οὖν ἐστιν ὡς ὁ Ε πρὸς τὸν ΒΓ, οὕτως τὸ ἀπὸ τῆς
Α πρὸς τὸ ἀπὸ τῆς ΖΗ, σύμμετρον ἄρα τὸ ἀπὸ τῆς Α τῷ
ἀπὸ τῆς ΖΗ. ῥητὸν δὲ τὸ ἀπὸ τῆς Α; ῥητὸν ἄρα καὶ
τὸ ἀπὸ τῆς ΖΗ; ῥητὴ ἄρα ἐστὶ καὶ ἡ ΖΗ. καὶ ἐπεὶ ὁ Ε
πρὸς τὸν ΒΓ λόγον οὐκ ἔχει, ὃν τετράγωνος ἀριθμὸς
πρὸς τετράγωνον ἀριθμόν, οὐδ ἄρα τὸ ἀπὸ τῆς Α πρὸς τὸ
ἀπὸ τῆς ΖΗ λόγον ἔχει, ὃν τετράγωνος ἀριθμὸς πρὸς τετράγωνον
ἀριθμόν; ἀσύμμετρος ἄρα ἐστὶν ἡ Α τῆ ΖΗ μήκει. πάλιν,
ἐπεί ἐστιν ὡς ὁ ΒΓ πρὸς τὸν ΓΔ, οὕτως τὸ ἀπὸ τῆς ΖΗ
πρὸς τὸ ἀπὸ τῆς ΗΘ, σύμμετρον ἄρα τὸ ἀπὸ τῆς ΖΗ τῷ
ἀπὸ τῆς ΗΘ. ῥητὸν δὲ τὸ ἀπὸ τῆς ΖΗ; ῥητὸν ἄρα καὶ τὸ
ἀπὸ τῆς ΗΘ; ῥητὴ ἄρα καὶ ἡ ΗΘ. καὶ ἐπεὶ ὁ ΒΓ πρὸς
τὸν ΓΔ λόγον οὐκ ἔχει, ὃν τετράγωνος ἀριθμὸς πρὸς τετράγωνον
ἀριθμόν, οὐδ ἄρα τὸ ἀπὸ τῆς ΖΗ πρὸς τὸ ἀπὸ τῆς ΗΘ
λόγον ἔχει, ὃν τετράγωνος ἀριθμὸς πρὸς τετράγωνον ἀριθμόν;
ἀσύμμετρος ἄρα ἐστὶν ἡ ΖΗ τῇ ΗΘ μήκει. καί εἰσιν
ἀμφότεραι ῥηταί; αἱ ΖΗ, ΗΘ ἄρα ῥηταί εἰσι δυνάμει μόνον
σύμμετροι; ἡ ἄρα ΖΘ ἀποτομή ἐστιν. λέγω δή, ὅτι καὶ
ἕκτη.}

{\greekfont Ἐπεὶ γάρ ἐστιν ὡς μὲν ὁ Ε πρὸς τὸν ΒΓ, οὕτως τὸ ἀπὸ τῆς
Α πρὸς τὸ ἀπὸ τῆς ΖΗ, ὡς δὲ ὁ ΒΓ πρὸς τὸν ΓΔ, οὕτως
τὸ ἀπὸ τῆς ΖΗ πρὸς τὸ ἀπὸ τῆς ΗΘ, δι ἴσου ἄρα ἐστὶν
ὡς ὁ Ε πρὸς τὸν ΓΔ, οὕτως τὸ ἀπὸ τῆς Α πρὸς τὸ ἀπὸ τῆς
ΗΘ. ὁ δὲ Ε πρὸς τὸν ΓΔ λόγον οὐκ ἔχει, ὃν τετράγωνος
ἀριθμὸς πρὸς τετράγωνον ἀριθμόν; οὐδ ἄρα τὸ ἀπὸ τῆς
Α πρὸς τὸ ἀπὸ τῆς ΗΘ λόγον ἔχει, ὃν τετράγωνος ἀριθμὸς
πρὸς τετράγωνον ἀριθμόν; ἀσύμμετρος ἄρα ἐστὶν ἡ 
Α τῇ ΗΘ μήκει; οὐδετέρα ἄρα τῶν ΖΗ, ΗΘ σύμμετρός ἐστι
τῇ Α ῥητῇ μήκει. ᾧ οὖν μεῖζόν
ἐστι τὸ ἀπὸ τῆς ΖΗ τοῦ ἀπὸ τῆς ΗΘ, ἔστω τὸ ἀπὸ
τῆς Κ. ἐπεὶ οὖν ἐστιν ὡς ὁ ΒΓ πρὸς τὸν ΓΔ, οὕτως
τὸ ἀπὸ τῆς ΖΗ πρὸς τὸ ἀπὸ τῆς ΗΘ, ἀναστρέψαντι ἄρα
ἐστὶν
ὡς ὁ ΓΒ πρὸς τὸν ΒΔ, οὕτως τὸ ἀπὸ τῆς ΖΗ πρὸς τὸ
ἀπὸ τῆς Κ. ὁ δὲ ΓΒ πρὸς τὸν ΒΔ λόγον οὐκ ἔχει,
ὃν τετράγωνος ἀριθμὸς πρὸς τετράγωνον ἀριθμόν;
οὐδ ἄρα τὸ ἀπὸ τῆς ΖΗ πρὸς τὸ ἀπὸ τῆς Κ
λόγον ἔχει, ὃν τετράγωνος ἀριθμὸς πρὸς τετράγωνον
ἀριθμόν; ἀσύμμετρος ἄρα ἐστὶν ἡ ΖΗ τῇ Κ μήκει.
καὶ δύναται ἡ ΖΗ τῆς ΗΘ μεῖζον τῷ ἀπὸ τῆς Κ; ἡ
ΖΗ ἄρα τῆς ΗΘ μεῖζον δύναται τῷ ἀπὸ  ἀσυμμέτρου
ἑαυτῇ μήκει. καὶ οὐδετέρα τῶν ΖΗ, ΗΘ σύμμετρός ἐστι τῇ
ἐκκειμένῃ ῥητῇ μήκει τῇ Α. ἡ ἄρα ΖΘ ἀποτομή ἐστιν
ἕκτη.}

{\greekfont Εὕρηται ἄρα ἡ ἕκτη ἀποτομὴ ἡ ΖΘ; ὅπερ ἔδει
δεῖχαι.}}

\ParallelRText{
\begin{center}
{\large Proposition 90}
\end{center}

To find a sixth apotome.

Let the rational (straight-line) $A$, and the three numbers
$E$, $BC$, and $CD$, not having to one another the ratio which (some)
square  number (has) to (some) square number, be laid down. Furthermore, let $CB$
also not have to $BD$ the ratio which (some) square number (has) to
(some) square number. And let it be contrived that as $E$
(is) to $BC$, so the (square) on $A$ (is) to the (square) on $FG$,
and as $BC$ (is) to $CD$, so the (square) on $FG$ (is) to the (square)
on $GH$ [Prop.~10.6~corr.].

\epsfysize=1.1in
\centerline{\epsffile{Book10/fig090e.eps}}

Therefore, since as $E$ is to $BC$, so the (square) on $A$ (is)
to the (square) on $FG$, the (square) on $A$ (is) thus commensurable
with the (square) on $FG$ [Prop.~10.6]. 
And the (square) on $A$ (is) rational. Thus, the (square) on $FG$
(is) also rational. Thus, $FG$  is also a rational (straight-line). And
since $E$ does not have to $BC$ the ratio which (some) square
number (has) to (some) square number, the (square) on $A$
thus does not have to the (square) on $FG$ the ratio which (some)
square number (has) to (some) square number either. Thus,
$A$ is incommensurable in length with $FG$ [Prop.~10.9]. Again, since as $BC$ is to $CD$, so
the (square) on $FG$ (is) to the (square) on $GH$, the (square) on
$FG$ (is) thus commensurable with the (square) on $GH$ [Prop.~10.6]. And the (square) on $FG$ (is) rational.
Thus, the (square) on $GH$ (is) also rational. Thus, $GH$
(is) also rational. And since $BC$ does not have to $CD$ the ratio
which (some) square number (has) to (some) square number, the
(square) on $FG$ thus does not have to the (square) on $GH$
the ratio which (some) square (number) has to (some) square (number)
either. Thus, $FG$ is incommensurable in length with $GH$ [Prop.~10.9]. And both are rational (straight-lines).
Thus, $FG$ and $GH$ are rational (straight-lines which are) commensurable
in square only. Thus, $FH$ is an apotome [Prop.~10.73]. So, I say that (it is) also a sixth (apotome).

For since as $E$ is to $BC$, so the (square) on $A$ (is) to the (square)
on $FG$, and as $BC$ (is) to $CD$, so the (square) on $FG$ (is) to
the (square) on $GH$, thus, via equality, as $E$ is to $CD$, so the
(square) on $A$ (is) to the (square) on $GH$ [Prop.~5.22]. And $E$ does not have to $CD$
the ratio which (some) square number (has) to (some) square number.
Thus, the (square) on $A$ does not have to the (square) $GH$ the
ratio which (some) square number (has) to (some) square number either. $A$ is thus incommensurable in length with $GH$ [Prop.~10.9]. Thus, neither of $FG$ and $GH$
is commensurable in length with the rational (straight-line) $A$. Therefore, let the (square)
on $K$ be that (area) by which the (square) on $FG$ is greater than the (square) on $GH$ [Prop.~10.13~lem.].
 Therefore, since as $BC$ is to $CD$, so the (square) on $FG$ (is) to the
(square) on $GH$, thus, via conversion, as $CB$ is to $BD$, so the (square)
on $FG$ (is) to the (square) on $K$ [Prop.~5.19~corr.]. And $CB$ does not have
to $BD$ the ratio which (some) square number (has) to (some) square
number. Thus, the (square) on $FG$ does not have to the (square) on
$K$ the ratio which (some) square number (has) to (some) square number
either. $FG$ is thus incommensurable in length with $K$ [Prop.~10.9]. And the square on $FG$ is greater
than (the square on) $GH$ by the (square) on $K$.
Thus, the square on
$FG$ is greater than (the square on) $GH$ by the (square)
on (some straight-line) incommensurable in length with ($FG$).
And neither of $FG$ and $GH$ is commensurable in length
with the (previously) laid down rational (straight-line) $A$. Thus,
$FH$ is a sixth apotome [Def.~10.16].

Thus, the sixth apotome $FH$ has been found. (Which is) the
very thing it was required to show.}
\end{Parallel}
{\footnotesize\noindent$^\dag$ See footnote to Prop.~10.53.}

%%%%
%10.91
%%%%
\pdfbookmark[1]{Proposition 10.91}{pdf10.91}
\begin{Parallel}{}{}
\ParallelLText{
\begin{center}
{\large \ggn{91}.}
\end{center}\vspace*{-7pt}

{\greekfont Ἐὰν θωρίον περιέθηται ὑπὸ ῥητῆς καὶ ἀποτομῆς πρώτης,
ἡ τὸ θωρίον δυναμένη ἀπορομή ἐστιν.}

{\greekfont Περιεθέσθω γὰρ θωρίον τὸ ΑΒ ὑπὸ ῥητῆς τῆς ΑΓ καὶ ἀποτομῆς
πρώτης τῆς ΑΔ; λέγω, ὅτι ἡ τὸ ΑΒ θωρίον δυναμένη ἀποτομή
ἐστιν.}

{\greekfont Ἐπεὶ γὰρ ἀποτομή ἐστι πρώτη ἡ ΑΔ, ἔστω α ὐτῇ προσαρμόζουσα
ἡ ΔΗ; αἱ ΑΗ, ΗΔ ἄρα ῥηταί εἰσι δυνάμει μόνον σύμμετροι.
καὶ ὅλη ἡ ΑΗ σύμμετρός ἐστι τῇ ἐκκειμένῃ ῥητῇ τῇ ΑΓ,
καὶ ἡ ΑΗ τῆς ΗΔ μεῖζον δύναται τῷ ἀπὸ συμμέτρου ἑαυτῇ
μήκει; ἐὰν ἄρα τῷ τετάρτῳ μέρει τοῦ ἀπὸ τῆς ΔΗ ἴσον
παρὰ τὴν ΑΗ παραβληθῇ ἐλλεῖπον εἴδει τετραγώνῳ, εἰς σύμμετρα
α ὐτὴν διαιρεῖ. τετμήσθω ἡ ΔΗ δίθα κατὰ τὸ Ε, καὶ τῷ ἀπὸ
τῆς ΕΗ ἴσον παρὰ τὴν ΑΗ παραβεβλήσθω ἐλλεῖπον εἴδει
τετραγώνῳ, καὶ ἔστω τὸ ὑπὸ τῶν ΑΖ, ΖΗ; σύμμετρος ἄρα
ἐστὶν ἡ ΑΖ τῇ ΖΗ. καὶ διὰ τῶν Ε, Ζ, Η σημείων τῇ ΑΓ
παράλληλοι ἤθθωσαν αἱ ΕΘ, ΖΙ, ΗΚ.}\\~\\~\\~\\~\\

\epsfysize=1.3in
\centerline{\epsffile{Book10/fig091g.eps}}


{\greekfont Καὶ ἐπεὶ σύμμετρός ἐστιν ἡ ΑΖ τῇ ΖΗ μήκει, καὶ ἡ ΑΗ
ἄρα ἑκατέρᾳ τῶν ΑΖ, ΖΗ σύμμετρός ἐστι μήκει. ἀλλὰ
ἡ ΑΗ σύμμετρός ἐστι τῇ ΑΓ; καὶ ἑκατέρα
ἄρα τῶν ΑΖ, ΖΗ σύμμετρός ἐστι τῇ ΑΓ μήκει. καί ἐστι
ῥητὴ ἡ ΑΓ; ῥητὴ ἄρα καὶ ἑκατέρα τῶν ΑΖ, ΖΗ; ὥστε καὶ
ἑκάτερον τῶν ΑΙ, ΖΚ ῥητόν ἐστιν. καὶ ἐπεὶ σύμμετρός
ἐστιν ἡ ΔΕ τῇ ΕΗ μήκει, καὶ ἡ ΔΗ ἄρα ἑκατέρᾳ τῶν
ΔΕ, ΕΗ σύμμετρός ἐστι μήκει. ῥητὴ δὲ ἡ ΔΗ καὶ
ἀσύμμετρος τῇ ΑΓ μήκει; ῥητὴ ἄρα καὶ ἑκατέρα τῶν
ΔΕ, ΕΗ καὶ ἀσύμμετρος τῇ ΑΓ μήκει; ἑκάτερον ἄρα τῶν
ΔΘ, ΕΚ μέσον ἐστίν.}

{\greekfont Κείσθω δὴ τῷ μὲν ΑΙ ἴσον τετράγωνον τὸ ΛΜ,
τῷ δὲ ΖΚ ἴσον τετράγωνον ἀφῃρήσθω κοινὴν γωνίαν
ἔθον α ὐτῷ τὴν ὑπὸ ΛΟΜ τὸ ΝΧ; περὶ τὴν
α ὐτὴν ἄρα διάμετρόν ἐστι τὰ ΛΜ, ΝΧ τετράγωνα. ἔστω
α ὐτῶν διάμετρος ἡ ΟΡ, καὶ καταγεγράφθω τὸ σχῆμα.
ἐπεὶ οὖν ἴσον ἐστὶ τὸ ὑπὸ τῶν ΑΖ, ΖΗ περιεχόμενον ὀρθογώνιον τῷ ἀπὸ τῆς ΕΗ τετραγώνῳ, ἔστιν ἄρα ὡς
ἡ ΑΖ πρὸς τὴν ΕΗ, οὕτως ἡ ΕΗ πρὸς τὴν ΖΗ.
ἀλλ ὡς μὲν ἡ ΑΖ πρὸς τὴν ΕΗ, οὕτως τὸ ΑΙ πρὸς
τὸ ΕΚ, ὡς δὲ ἡ ΕΗ πρὸς τὴν ΖΗ, οὕτως ἐστὶ τὸ ΕΚ πρὸς
τὸ ΚΖ; τῶν ἄρα ΑΙ, ΚΖ μέσον ἀνάλογόν ἐστι τὸ ΕΚ. ἔστι
δὲ καὶ τῶν ΛΜ, ΝΧ μέσον ἀνάλογον τὸ ΜΝ,  ὡς
ἐν τοῖς ἔμπροσθεν ἐδείθθη, καί ἐστι τὸ [μὲν] ΑΙ τῷ ΛΜ
τετραγώνῳ ἴσον, τὸ δὲ ΚΖ τῷ ΝΧ; καὶ τὸ ΜΝ ἄρα τῷ ΕΚ ἴσον
 ἐστίν. ἀλλὰ τὸ μὲν ΕΚ τῷ ΔΘ ἐστιν ἴσον,
τὸ δὲ ΜΝ τῷ ΛΧ; τὸ ἄρα ΔΚ ἴσον ἐστὶ τῷ ΥΦΘ γνώμονι
καὶ τῷ ΝΧ. ἔστι δὲ καὶ τὸ ΑΚ ἴσον τοῖς ΛΜ, ΝΧ τετραγώνοις;
λοιπὸν ἄρα τὸ ΑΒ ἴσον ἐστὶ τῷ ΣΤ. τὸ δὲ ΣΤ τὸ ἀπὸ
τῆς ΛΝ ἐστι τετράγωνον; τὸ ἄρα ἀπὸ τῆς ΛΝ τετράγωνον ἴσον
ἐστὶ τῷ ΑΒ; ἡ ΛΝ ἄρα δύναται τὸ ΑΒ. λέγω δή, ὅτι
ἡ ΛΝ ἀποτομή ἐστιν.}

{\greekfont Ἐπεὶ γὰρ ῥητόν ἐστιν ἑκάτερον τῶν ΑΙ, ΖΚ, καί ἐστιν ἴσον
τοῖς ΛΜ, ΝΧ, καὶ ἑκάτερον ἄρα τῶν ΛΜ, ΝΧ ῥητόν ἐστιν,
τουτέστι τὸ ἀπὸ ἑκατέρας τῶν ΛΟ, ΟΝ; καὶ ἑκατέρα ἄρα τῶν
ΛΟ, ΟΝ ῥητή ἐστιν. πάλιν, ἐπεὶ μέσον ἐστὶ τὸ ΔΘ καί ἐστιν
ἴσον τῷ ΛΧ, μέσον ἄρα ἐστὶ καὶ τὸ ΛΧ. ἐπεὶ οὖν τὸ μὲν
ΛΧ μέσον ἐστίν, τὸ δὲ ΝΧ ῥητόν, ἀσύμμετρον ἄρα ἐστὶ
τὸ ΛΧ τῷ ΝΧ; ὡς δὲ τὸ ΛΧ πρὸς τὸ ΝΧ, οὕτως
ἐστὶν ἡ ΛΟ πρὸς τὴν ΟΝ; ἀσύμμετρος ἄρα ἐστὶν ἡ  ΛΟ
τῇ ΟΝ μήκει. καί εἰσιν ἀμφότεραι ῥηταί; αἱ ΛΟ, ΟΝ ἄρα
ῥηταί εἰσι δυνάμει μόνον σύμμετροι; ἀποτομὴ
ἄρα ἐστὶν ἡ ΛΝ. καὶ δύναται τὸ ΑΒ θωρίον;
ἡ ἄρα τὸ ΑΒ θωρίον δυναμένη ἀποτομή ἐστιν.}

{\greekfont Ἐὰν ἄρα θωρίον περιέθηται ὑπὸ ῥητῆς καὶ τὰ
ἑχῆς.}}

\ParallelRText{
\begin{center}
{\large Proposition 91}
\end{center}

If an area is contained by a rational (straight-line)
and a first apotome then the square-root of the area is an apotome.

For let the area $AB$ be contained by the rational (straight-line) $AC$ and the first apotome $AD$. I say that the square-root of area $AB$ is an apotome.

For since $AD$ is a first apotome, let $DG$ be its attachment. Thus,
$AG$ and $DG$ are rational (straight-lines which are) commensurable in
square only [Prop.~10.73]. And the whole, $AG$,
is commensurable (in length) with the (previously) laid down rational (straight-line)
$AC$, and the square on $AG$ is greater than (the square on) $GD$
by the (square) on  (some straight-line) commensurable in length with ($AG$) [Def.~10.11]. Thus, if (an area) equal to the
fourth part of the (square) on $DG$ is applied to $AG$, falling short
by a square figure, then it divides ($AG$) into (parts which are) commensurable
(in length) [Prop.~10.17]. Let $DG$
be cut in half at $E$. And let (an area) equal to the (square) on
$EG$ be applied to $AG$, falling short by a square figure. And let
it be the (rectangle contained) by $AF$ and $FG$. $AF$ is
 thus commensurable (in length) with $FG$. And let $EH$, $FI$, and $GK$
be drawn through points $E$, $F$, and $G$ (respectively),
parallel to $AC$.

\epsfysize=1.3in
\centerline{\epsffile{Book10/fig091e.eps}}

And since $AF$ is commensurable in length with $FG$, $AG$ is thus also
commensurable in length with each of $AF$ and $FG$ [Prop.~10.15]. But $AG$ is commensurable (in length)
with $AC$. Thus, each of $AF$ and $FG$ is also commensurable in length with $AC$ 
[Prop.~10.12]. 
And $AC$ is a rational (straight-line). Thus,  $AF$ and $FG$ (are) each
also rational (straight-lines). Hence, $AI$ and
$FK$ are also each rational (areas) [Prop.~10.19]. And since $DE$ is commensurable
in length with $EG$, $DG$ is thus also commensurable
in length with each of $DE$ and $EG$ [Prop.~10.15]. And $DG$ (is) rational, and
incommensurable in length with $AC$. $DE$ and $EG$ (are) thus each
rational, and incommensurable in length with $AC$ [Prop.~10.13]. Thus, $DH$ and $EK$ are
each medial (areas) [Prop.~10.21].

So let the square $LM$, equal to $AI$, be laid down. And let the square $NO$,
equal to $FK$, be
subtracted (from $LM$), having with it the common angle $LPM$. Thus,
the squares $LM$ and $NO$ are about the same diagonal [Prop.~6.26]. Let $PR$ be their (common) diagonal, and
let the (rest of the) figure be drawn. Therefore, since
the rectangle contained by $AF$ and $FG$ is equal to the square $EG$,
thus as $AF$ is to $EG$, so $EG$ (is) to $FG$ [Prop.~6.17]. But, as $AF$ (is) to $EG$, so $AI$ (is)
to $EK$, and as $EG$ (is) to $FG$, so $EK$ is to $KF$ [Prop.~6.1]. Thus, $EK$ is the mean proportional
to $AI$ and $KF$ [Prop.~5.11]. And $MN$ is also the mean proportional to
$LM$ and $NO$, as shown before [Prop.~10.53~lem.].  And $AI$ is equal to the
square $LM$, and $KF$ to $NO$. Thus, $MN$ is also equal to $EK$.
But, $EK$ is equal to $DH$, and $MN$ to $LO$ [Prop.~1.43]. Thus, $DK$ is equal to the
gnomon $UVW$ and  $NO$. And $AK$ is also equal to 
(the sum of) the squares $LM$ and $NO$.  Thus, the remainder $AB$ is equal
to $ST$. And $ST$ is the square on $LN$. Thus, the square on $LN$
is equal to $AB$. Thus, $LN$ is the square-root of $AB$. So, I say that
$LN$ is an apotome.

For since $AI$ and $FK$ are each rational (areas), and are equal to
$LM$ and $NO$ (respectively), thus  $LM$ and $NO$---that is to say, the
(squares) on each of $LP$ and $PN$ (respectively)---are  also
each rational (areas). Thus, $LP$ and $PN$ are also each rational (straight-lines). 
Again, since $DH$ is a medial (area), and is equal to $LO$, $LO$ is thus
also a medial (area). Therefore, since $LO$ is medial, and $NO$ rational, 
$LO$ is thus incommensurable with $NO$. And as $LO$ (is) to $NO$,
so $LP$ is to $PN$ [Prop.~6.1]. $LP$ is
thus incommensurable in length with $PN$ [Prop.~10.11]. And they are both rational
(straight-lines). Thus, $LP$ and $PN$ are rational (straight-lines which are)
commensurable in square only. Thus, $LN$ is an apotome [Prop.~10.73]. And it is the square-root of area $AB$. Thus, the square-root of area $AB$ is an apotome.

Thus, if an area is contained by a rational (straight-line), and
so on \ldots.}
\end{Parallel}

%%%%
%10.92
%%%%
\pdfbookmark[1]{Proposition 10.92}{pdf10.92}
\begin{Parallel}{}{}
\ParallelLText{
\begin{center}
{\large \ggn{92}.}
\end{center}\vspace*{-7pt}

{\greekfont Ἐὰν θωρίον περιέθηται ὑπὸ ῥητῆς καὶ ἀποτομῆς δευτέρας,
ἡ τὸ θωρίον δυναμένη μέσης ἀποτομή ἐστι πρώτη.}\\

\epsfysize=1.3in
\centerline{\epsffile{Book10/fig092g.eps}}

{\greekfont Θωρίον γὰρ τὸ ΑΒ περιεθέσθω ὑπὸ ῥητῆς τῆς ΑΓ καὶ ἀποτομῆς
δευτέρας τῆς ΑΔ; λέγω, ὅτι ἡ τὸ ΑΒ θωρίον δυναμένη μέσης
ἀποτομή ἐστι πρώτη.}

{\greekfont Ἔστω γὰρ τῇ ΑΔ προσαρμόζουσα ἡ ΔΗ; αἱ ἄρα ΑΗ, ΗΔ
ῥηταί εἰσι δυνάμει μόνον σύμμετροι, καὶ ἡ προσαρμόζουσα
ἡ ΔΗ σύμμετρός ἐστι τῇ ἐκκειμένῃ
ῥητῇ τῇ ΑΓ, ἡ δὲ ὅλη ἡ ΑΗ τῆς προσαρμοζούσης τῆς
ΗΔ μεῖζον δύναται τῷ ἀπὸ συμμέτρου ἑαυτῇ μήκει.
ἐπεὶ οὖν ἡ ΑΗ τῆς ΗΔ μεῖζον δύναται τῷ ἀπὸ
συμμέτρου ἑαυτῇ, ἐὰν ἄρα τῷ τετάρτῳ μέρει τοῦ ἀπὸ τῆς
ΗΔ ἴσον παρὰ τὴν ΑΗ παραβληθῇ ἐλλεῖπον εἴδει
τετραγώνῳ, εἰς σύμμετρα α ὐτὴν διαιρεῖ. τετμήσθω οὖν ἡ
ΔΗ δίθα κατὰ τὸ Ε; καὶ τῷ ἀπὸ τῆς ΕΗ ἴσον παρὰ τὴν ΑΗ
παραβεβλήσθω ἐλλεῖπον εἴδει τετραγώνῳ, καὶ ἔστω τὸ ὑπὸ
τῶν ΑΖ, ΖΗ; σύμμετρος ἄρα ἐστὶν ἡ ΑΖ τῇ ΖΗ μήκει. καὶ
ἡ ΑΗ ἄρα ἑκατέρᾳ τῶν ΑΖ, ΖΗ σύμμετρός ἐστι μήκει.
ῥητὴ δὲ ἡ ΑΗ καὶ ἀσύμμετρος τῇ ΑΓ μήκει; καὶ ἑκατέρα
ἄρα τῶν ΑΖ, ΖΗ ῥητή ἐστι καὶ ἀσύμμετρος τῇ ΑΓ μήκει;
ἑκάτερον ἄρα τῶν ΑΙ, ΖΚ μέσον ἐστίν. πάλιν, ἐπεὶ
σύμμετρός ἐστιν ἡ ΔΕ τῇ ΕΗ, καὶ ἡ ΔΗ ἄρα ἑκατέρᾳ τῶν
ΔΕ, ΕΗ σύμμετρός ἐστιν. ἀλλ ἡ ΔΗ σύμμετρός ἐστι τῇ
ΑΓ μήκει [ῥητὴ ἄρα καὶ ἑκατέρα τῶν ΔΕ, ΕΗ καὶ σύμμετρος
τῇ ΑΓ μήκει]. ἑκάτερον ἄρα τῶν ΔΘ, ΕΚ ῥητόν ἐστιν.}

{\greekfont Συνεστάτω οὖν τῷ μὲν ΑΙ ἴσον τετράγωνον τὸ ΛΜ, τῷ
δὲ ΖΚ ἴσον ἀφῃρήσθω τὸ ΝΧ περὶ τὴν α ὐτὴν γωνίαν ὂν
τῷ ΛΜ τὴν ὑπὸ τῶν ΛΟΜ; περὶ τὴν α ὐτὴν ἄρα ἐστὶ
διάμετρον τὰ ΛΜ, ΝΧ τετράγωνα. ἔστω α ὐτῶν διάμετρος
ἡ ΟΡ, καὶ καταγεγράφθω τὸ σχῆμα. ἐπεὶ οὖν τὰ ΑΙ, ΖΚ μέσα
ἐστὶ καί ἐστιν ἴσα τοῖς ἀπὸ τῶν ΛΟ, ΟΝ, καὶ τὰ ἀπὸ τῶν
ΛΟ, ΟΝ [ἄρα] μέσα ἐστίν; καὶ αἱ ΛΟ, ΟΝ ἄρα μέσαι εἰσὶ
δυνάμει μόνον σύμμετροι. καὶ ἐπεὶ τὸ ὑπὸ τῶν ΑΖ,
ΖΗ ἴσον ἐστὶ τῷ ἀπὸ τῆς ΕΗ, ἔστιν ἄρα ὡς ἡ ΑΖ πρὸς
τὴν ΕΗ, οὕτως ἡ ΕΗ πρὸς τὴν ΖΗ; ἀλλ ὡς μὲν ἡ ΑΖ
πρὸς τὴν ΕΗ, οὕτως τὸ ΑΙ πρὸς τὸ ΕΚ; ὡς δὲ ἡ ΕΗ πρὸς τὴν
ΖΗ, οὕτως [ἐστὶ] τὸ ΕΚ πρὸς τὸ ΖΚ; τῶν ἄρα ΑΙ, ΖΚ μέσον
ἀνάλογόν ἐστι τὸ ΕΚ. ἔστι δὲ καὶ τῶν ΛΜ, ΝΧ τετραγώνων
μέσον ἀνάλογον τὸ ΜΝ; καί ἐστιν ἴσον τὸ μὲν ΑΙ τῷ ΛΜ,
τὸ δὲ ΖΚ τῷ ΝΧ;  καὶ  τὸ ΜΝ ἄρα ἴσον ἐστὶ τῷ ΕΚ. ἀλλὰ
τῷ μὲν ΕΚ ἴσον [ἐστὶ] τὸ ΔΘ, τῷ δὲ ΜΝ ἴσον τὸ ΛΧ; ὅλον
ἄρα τὸ ΔΚ ἴσον ἐστὶ τῷ ΥΦΘ γνώμονι καὶ τῷ ΝΧ. ἐπεὶ
οὖν ὅλον τὸ ΑΚ ἴσον ἐστὶ τοῖς ΛΜ, ΝΧ, ὧν τὸ ΔΚ
ἴσον ἐστὶ τῷ ΥΦΘ γνώμονι καὶ τῷ ΝΧ, λοιπὸν ἄρα τὸ ΑΒ
ἴσον ἐστὶ τῷ ΤΣ. τὸ δὲ ΤΣ ἐστι τὸ ἀπὸ τῆς ΛΝ;
τὸ ἀπὸ τῆς ΛΝ ἄρα ἴσον ἐστὶ τῷ ΑΒ θωρίῳ; ἡ ΛΝ
ἄρα δύναται τὸ ΑΒ θωρίον. λέγω [δή], ὅτι ἡ ΛΝ μέσης
ἀποτομή ἐστι πρώτη.}

{\greekfont Ἐπεὶ γὰρ ῥητόν ἐστι τὸ ΕΚ καί ἐστιν ἴσον τῷ ΛΧ,
ῥητὸν ἄρα ἐστὶ τὸ ΛΧ, τουτέστι τὸ ὑπὸ τῶν
ΛΟ, ΟΝ. μέσον δὲ ἐδείθθη τὸ ΝΧ; ἀσύμμετρον ἄρα
ἐστὶ τὸ ΛΧ τῷ ΝΧ; ὡς δὲ τὸ ΛΧ πρὸς τὸ ΝΧ,
οὕτως ἐστὶν ἡ ΛΟ πρὸς ΟΝ; αἱ ΛΟ, ΟΝ ἄρα ἀσύμμετροί
εἰσι μήκει. αἱ ἄρα ΛΟ, ΟΝ μέσαι εἰσὶ δυνάμει μόνον σύμμετροι
ῥητὸν περιέθουσαι; ἡ ΛΝ ἄρα μέσης ἀποτομή ἐστι πρώτη;
καὶ δύναται τὸ ΑΒ θωρίον.}

{\greekfont Ἡ ἄρα τὸ ΑΒ θωρίον δυναμένη μέσης ἀποτομή ἐστι
πρώτη; ὅπερ ἔδει δεῖχαι.}}

\ParallelRText{
\begin{center}
{\large Proposition 92}
\end{center}

If an area is contained by a rational (straight-line)
and a second apotome then the square-root of the area is a first apotome of a medial (straight-line).

\epsfysize=1.3in
\centerline{\epsffile{Book10/fig092e.eps}}

For let the area $AB$ be contained by the rational (straight-line) $AC$ and the second apotome $AD$. I say that the square-root of area $AB$ is the first apotome of a medial (straight-line).

For let $DG$ be an attachment to $AD$. Thus, $AG$ and $GD$
are rational (straight-lines which are) commensurable in square only [Prop.~10.73],
and the attachment $DG$ is commensurable (in length) with the
(previously) laid down rational (straight-line) $AC$, and the square on
the whole, $AG$, is greater than (the square on) the attachment, $GD$,
by the (square) on (some straight-line) commensurable in length with ($AG$) 
[Def.~10.12]. Therefore, since the square on $AG$ is
greater than (the square on) $GD$ by the (square) on (some
straight-line) commensurable (in length) with ($AG$), thus if  (an area)
equal to the fourth part of the (square) on $GD$ is applied
to $AG$, falling short by a square figure, then it divides ($AG$) into (parts which
are) commensurable (in length) [Prop.~10.17]. 
Therefore, let $DG$ be cut in half at $E$. And let (an area) equal to
the (square) on $EG$ be applied to $AG$, falling short
by a square figure. And let it be the (rectangle contained) by $AF$ and $FG$.
Thus, $AF$ is commensurable in length with $FG$. $AG$ is thus
also commensurable in length with each of $AF$ and $FG$ [Prop.~10.15]. And $AG$ (is) a rational (straight-line), and incommensurable in length with $AC$. $AF$ and $FG$
are thus also each rational (straight-lines), and incommensurable in length with
$AC$ [Prop.~10.13].
Thus, $AI$
and $FK$ are each medial (areas) [Prop.~10.21]. 
Again, since $DE$ is commensurable (in length) with $EG$, thus
$DG$ is also commensurable (in length) with each of $DE$ and $EG$ [Prop.~10.15]. But, $DG$ is commensurable
in length with $AC$ [thus, $DE$ and $EG$ are also each rational, and
commensurable in length with $AC$]. Thus, $DH$ and $EK$ are each
rational (areas) [Prop.~10.19].

Therefore, let the square $LM$, equal to $AI$, be constructed.
And let $NO$, equal to $FK$, which is about the
same angle $LPM$ as $LM$, be subtracted (from $LM$).
Thus, the squares $LM$ and $NO$ are about the same diagonal
[Prop.~6.26]. Let $PR$ be their (common) diagonal, and
let the (rest of the) figure be drawn. Therefore, since $AI$ and $FK$
are medial (areas), and are equal to the (squares) on $LP$ and $PN$
(respectively), [thus] the (squares) on $LP$ and $PN$ are also medial.
Thus, $LP$ and $PN$ are  also medial (straight-lines which are)
commensurable in square only.$^\dag$ And since the (rectangle contained) by
$AF$ and $FG$ is equal to the (square) on $EG$, thus as $AF$ is to $EG$, so $EG$ (is) to $FG$ [Prop.~10.17]. But, as
$AF$ (is) to $EG$, so $AI$ (is) to $EK$. And as $EG$ (is) to $FG$, so
$EK$ [is] to $FK$ [Prop.~6.1]. Thus, $EK$ is the mean proportional to $AI$ and $FK$ [Prop.~5.11].
And $MN$ is also the mean proportional to the squares $LM$ and $NO$ [Prop.~10.53~lem.]. 
And $AI$ is equal to $LM$, and $FK$ to $NO$. 
Thus, $MN$ is also equal to
$EK$. But, $DH$ [is] equal to $EK$, and $LO$ equal to $MN$ [Prop.~1.43]. Thus, the whole  (of) $DK$ is equal to
the gnomon $UVW$ and $NO$. Therefore, since the whole (of) $AK$
is equal to $LM$ and $NO$, of which $DK$ is equal to the
gnomon $UVW$ and $NO$,  the remainder $AB$ is thus equal to $TS$.
And $TS$ is the (square) on $LN$. Thus, the (square) on $LN$
is equal to the area $AB$. $LN$ is thus the square-root of area $AB$.
[So], I say that $LN$ is the first apotome of a medial (straight-line).

For since $EK$ is a rational (area), and is equal to $LO$, $LO$---that is to
say, the (rectangle contained) by $LP$ and $PN$---is thus
a rational (area). And $NO$ was shown (to be) a medial (area). Thus,
$LO$ is incommensurable  with $NO$. And as $LO$ (is) to $NO$,
so $LP$ is to $PN$ [Prop.~6.1]. Thus,  $LP$ and
$PN$ are incommensurable in length [Prop.~10.11]. 
$LP$ and $PN$ are thus medial (straight-lines which are) commensurable
in square only, and which contain a rational (area). Thus, $LN$ is the
first apotome of a medial (straight-line) [Prop.~10.74]. And it is the square-root of area $AB$.

Thus, the square root of area $AB$ is the first apotome of a medial (straight-line). (Which is) the very thing it was required to show.}
\end{Parallel}
{\footnotesize\noindent$^\dag$ There is an error in the argument here. It should just say that $LP$ and $PN$ are commensurable in square,
rather than in square only, since $LP$ and $PN$ are only shown
to be incommensurable in length later on.}

%%%%
%10.93
%%%%
\pdfbookmark[1]{Proposition 10.93}{pdf10.93}
\begin{Parallel}{}{}
\ParallelLText{
\begin{center}
{\large \ggn{93}.}
\end{center}\vspace*{-7pt}

{\greekfont Ἐὰν θωρίον περιέθηται ὑπὸ ῥητῆς καὶ ἀποτομῆς τρίτης,
ἡ τὸ θωρίον δυναμένη μέσης ἀποτομή ἐστι δευτέρα.}

{\greekfont Θωρίον γὰρ τὸ ΑΒ περιεθέσθω ὑπὸ ῥητῆς τῆς ΑΓ καὶ ἀποτομῆς
τρίτης τῆς ΑΔ; λέγω, ὅτι ἡ τὸ ΑΒ θωρίον δυναμένη μέσης
ἀποτομή ἐστι δευτέρα.}\\

\epsfysize=1.3in
\centerline{\epsffile{Book10/fig093g.eps}}

{\greekfont Ἔστω γὰρ τῇ ΑΔ προσαρμόζουσα ἡ ΔΗ; αἱ ΑΗ, ΗΔ ἄρα ῥηταί εἰσι
δυνάμει μόνον σύμμετροι, καὶ οὐδετέρα τῶν ΑΗ, ΗΔ σύμμετρός
ἐστι μήκει τῇ ἐκκειμένῃ ῥητῇ τῇ ΑΓ, ἡ δὲ ὅλη ἡ ΑΗ
τὴς προσαρμοζούσης τῆς ΔΗ μεῖζον δύναται τῷ ἀπὸ
συμμέτρου ἑαυτῇ. ἐπεὶ οὖν ἡ ΑΗ τῆς ΗΔ μεῖζον
δύναται τῷ ἀπὸ συμμέτρου ἑαυτῇ, ἐὰν ἄρα τῷ τετάρτῳ
μέρει τοῦ ἀπὸ τῆς ΔΗ ἴσον παρὰ τὴν ΑΗ παραβληθῇ ἐλλεῖπον
εἴδει τετραγώνῳ, εἰς σύμμετρα α ὐτὴν διελεῖ. τετμήσθω
οὖν ἡ ΔΗ δίθα κατὰ τὸ Ε, καὶ τῷ ἀπὸ τῆς ΕΗ ἴσον
παρὰ τὴν ΑΗ παραβεβλήσθω ἐλλεῖπον εἴδει τετραγώνῳ, καὶ
ἔστω τὸ ὑπὸ τῶν ΑΖ, ΖΗ. καὶ ἤθθωσαν διὰ τῶν Ε, Ζ, Η σημείων
τῇ ΑΓ παράλληλοι αἱ ΕΘ, ΖΙ, ΗΚ; σύμμετροι ἄρα εἰσὶν αἱ ΑΖ, ΖΗ;
σύμμετρον ἄρα καὶ τὸ ΑΙ τῷ ΖΚ. καὶ ἐπεὶ αἱ ΑΖ, ΖΗ σύμμετροί
εἰσι μήκει, καὶ ἡ ΑΗ ἄρα ἑκατέρᾳ τῶν ΑΖ, ΖΗ σύμμετρός
ἐστι μήκει. ῥητὴ δὲ ἡ ΑΗ καὶ ἀσύμμετρος τῇ ΑΓ μήκει; ὥστε
καὶ αἱ ΑΖ, ΖΗ. ἑκάτερον ἄρα τῶν ΑΙ, ΖΚ μέσον ἐστίν. πάλιν, ἐπεὶ
σύμμετρός ἐστιν ἡ ΔΕ τῇ ΕΗ μήκει, καὶ ἡ ΔΗ ἄρα ἑκατέρᾳ
τῶν ΔΕ, ΕΗ σύμμετρός ἐστι μήκει. ῥητὴ δὲ ἡ ΗΔ καὶ
ἀσύμμετρος τῇ ΑΓ μήκει; ῥητὴ ἄρα καὶ ἑκατέρα τῶν ΔΕ,
ΕΗ καὶ ἀσύμμετρος τῇ ΑΓ μήκει; ἑκάτερον ἄρα τῶν ΔΘ, ΕΚ
μέσον ἐστίν. καὶ ἐπεὶ αἱ ΑΗ, ΗΔ δυνάμει μόνον σύμμετροί
εἰσιν, ἀσύμμετρος ἄρα ἐστὶ μήκει ἡ ΑΗ τῇ ΗΔ. ἀλλ
ἡ μὲν ΑΗ τῇ ΑΖ σύμμετρός ἐστι μήκει ἡ δὲ ΔΗ τῇ ΕΗ;
ἀσύμμετρος ἄρα ἐστὶν ἡ ΑΖ τῇ ΕΗ μήκει. ὡς δὲ ἡ ΑΖ
πρὸς τὴν ΕΗ, οὕτως ἐστὶ τὸ ΑΙ πρὸς τὸ ΕΚ;
ἀσύμμετρον ἄρα ἐστὶ τὸ ΑΙ τῷ ΕΚ.}

{\greekfont Συνεστάτω οὖν τῷ μὲν ΑΙ ἴσον τετράγωνον τὸ ΛΜ, τῷ δὲ
ΖΚ ἴσον ἀφῇρήσθω τὸ ΝΧ περὶ τὴν α ὐτὴν γωνίαν ὂν τῷ
ΛΜ; περὶ τὴν α ὐτὴν ἄρα διάμετρόν ἐστι τὰ ΛΜ, ΝΧ. ἔστω
α ὐτῶν διάμετρος ἡ ΟΡ, καὶ καταγεγράφθω τὸ σχῆμα. ἐπεὶ
οὖν τὸ ὑπὸ τῶν ΑΖ, ΖΗ ἴσον ἐστὶ τῷ ἀπὸ τῆς ΕΗ,
ἔστιν ἄρα ὡς ἡ ΑΖ πρὸς τὴν ΕΗ, οὕτως ἡ ΕΗ πρὸς τὴν
ΖΗ. ἀλλ ὡς μὲν ἡ ΑΖ πρὸς τὴν ΕΗ, οὕτως ἐστὶ τὸ ΑΙ
πρὸς τὸ ΕΚ; ὡς δὲ ἡ ΕΗ πρὸς τὴν ΖΗ, οὕτως ἐστὶ τὸ ΕΚ πρὸς τὸ
ΖΚ; καὶ ὡς ἄρα τὸ ΑΙ πρὸς τὸ ΕΚ, οὕτως τὸ ΕΚ πρὸς τὸ ΖΚ;
τῶν ἄρα ΑΙ, ΖΚ μέσον ἀνάλογόν ἐστι τὸ ΕΚ. ἔστι δὲ καὶ τῶν
ΛΜ, ΝΧ τετραγώνων μέσον ἀνάλογον τὸ ΜΝ; καί ἐστιν ἴσον
τὸ μὲν ΑΙ τῷ ΛΜ, τὸ δὲ ΖΚ τῷ ΝΧ; καὶ τὸ ΕΚ ἄρα  ἴσον
ἐστὶ τῷ ΜΝ. ἀλλὰ τὸ μὲν ΜΝ ἴσον ἐστὶ τῷ ΛΧ,
τὸ δὲ ΕΚ ἴσον [ἐστὶ] τῷ ΔΘ; καὶ ὅλον ἄρα τὸ
ΔΚ ἴσον ἐστὶ τῷ ΥΦΘ γνώμονι καὶ τῷ
ΝΧ. ἔστι δὲ καὶ τὸ ΑΚ ἴσον τοῖς ΛΜ, ΝΧ; λοιπὸν ἄρα τὸ
ΑΒ ἴσον ἐστὶ τῷ ΣΤ, τουτέστι τῷ ἀπὸ τῆς ΛΝ τετραγώνῳ;
ἡ ΛΝ ἄρα δύναται τὸ ΑΒ θωρίον. λέγω, ὅτι ἡ ΛΝ μέσης
ἀποτομή ἐστι δευτέρα.}

{\greekfont Ἐπεὶ γὰρ μέσα ἐδείθθη τὰ ΑΙ, ΖΚ καί ἐστιν ἴσα τοῖς ἀπὸ
τῶν ΛΟ, ΟΝ, μέσον ἄρα καὶ ἑκάτερον τῶν ἀπὸ τῶν ΛΟ, ΟΝ;
μέση ἄρα ἑκατέρα τῶν ΛΟ, ΟΝ. καὶ ἐπεὶ σύμμετρόν ἐστι
τὸ ΑΙ τῷ ΖΚ, σύμμετρον ἄρα καὶ τὸ ἀπὸ τῆς ΛΟ τῷ ἀπὸ
τῆς ΟΝ. πάλιν, ἐπεὶ ἀσύμμετρον ἐδείθθη τὸ ΑΙ τῷ ΕΚ,
ἀσύμμετρον ἄρα ἐστὶ καὶ τὸ ΛΜ τῷ ΜΝ, τουτέστι τὸ
ἀπὸ τῆς ΛΟ τῷ ὑπὸ τῶν ΛΟ, ΟΝ; ὥστε καὶ ἡ ΛΟ
ἀσύμμετρός ἐστι μήκει τῇ ΟΝ; αἱ ΛΟ, ΟΝ ἄρα μέσαι
εἰσὶ δυνάμει μόνον σύμμετροι. λέγω δή, ὅτι καὶ μέσον περιέθουσιν.}

{\greekfont Ἐπεὶ γὰρ μέσον ἐδείθθη τὸ ΕΚ καί ἐστιν ἴσον τῷ ὑπὸ τῶν
ΛΟ, ΟΝ, μέσον ἄρα ἐστὶ καὶ τὸ ὑπὸ τῶν ΛΟ, ΟΝ; ὥστε αἱ
ΛΟ, ΟΝ μέσαι εἰσὶ δυνάμει μόνον σύμμετροι μέσον περιέθουσαι.
ἡ ΛΝ ἄρα μέσης ἀποτομή ἐστι δευτέρα; καὶ δύναται τὸ ΑΒ
θωρίον.}

{\greekfont Ἡ ἄρα τὸ ΑΒ θωρίον δυναμένη μέσης ἀποτομή ἐστι δευτέρα;
ὅπερ ἔδει δεῖχαι.}}

\ParallelRText{
\begin{center}
{\large Proposition 93}
\end{center}

If an area is contained by a rational (straight-line)
and a third apotome then the square-root of the area is a second apotome
of a medial (straight-line).

For let the area $AB$ be contained by the rational (straight-line)
$AC$ and the third apotome $AD$. I say that the square-root of area $AB$
is the second apotome of a medial (straight-line).

\epsfysize=1.3in
\centerline{\epsffile{Book10/fig093e.eps}}

For let $DG$ be an attachment to $AD$. Thus, $AG$ and $GD$ are rational (straight-lines which are) commensurable in square only [Prop.~10.73], and neither of $AG$ and $GD$
is commensurable in length with the (previously) laid down rational (straight-line) $AC$, and the square on the whole, $AG$, is greater than
(the square on) the attachment, $DG$, by the (square) on (some straight-line)
commensurable (in length) with ($AG$) [Def.~10.13]. 
Therefore, since the square on $AG$ is greater than (the square on)
$GD$ by the (square) on (some straight-line) commensurable (in length)
with ($AG$), thus if (an area) equal to the fourth part of the square on $DG$
is applied to $AG$, falling short by a square figure, then it divides ($AG$)
into (parts which are) commensurable (in length) [Prop.~10.17]. Therefore, let $DG$ be
cut in half at $E$. And let (an area) equal to the (square) on $EG$
be applied to $AG$, falling short by a square figure.
And let it be the (rectangle contained) by $AF$ and $FG$. And let
$EH$, $FI$, and $GK$ be drawn through points $E$, $F$, and $G$
(respectively), parallel to $AC$. Thus, $AF$ and $FG$ are
commensurable (in length). $AI$ (is) thus also commensurable with
$FK$ [Props.~6.1, 10.11]. 
 And since $AF$ and $FG$ are commensurable
in length, $AG$ is thus also commensurable in length with each of
$AF$ and $FG$ [Prop.~10.15]. And
$AG$ (is) rational, and incommensurable in length with $AC$. 
Hence, $AF$ and $FG$ (are) also (rational, and incommensurable in length
with $AC$) [Prop.~10.13]. 
Thus, $AI$ and $FK$ are each medial (areas) [Prop.~10.21]. Again, since $DE$ is commensurable
in length with $EG$, $DG$ is also commensurable
in length with each of $DE$ and $EG$ [Prop.~10.15]. And $GD$ (is) rational, and incommensurable in length with $AC$. Thus, $DE$ and $EG$ (are)
each also rational, and incommensurable in length with $AC$ [Prop.~10.13]. $DH$ and $EK$ are thus
each medial (areas) [Prop.~10.21]. And since 
$AG$ and $GD$ are commensurable in square only, $AG$ is thus
incommensurable in length with $GD$. But, $AG$ is commensurable in
length with $AF$, and $DG$ with $EG$. Thus, $AF$ is incommensurable
in length with $EG$ [Prop.~10.13]. 
And as $AF$ (is) to $EG$, so $AI$ is to $EK$ [Prop.~6.1]. Thus, $AI$ is incommensurable with $EK$ [Prop.~10.11].

Therefore, let the square $LM$, equal to $AI$, be constructed. And
let $NO$, equal to $FK$, which is about the same angle as $LM$, 
be subtracted (from $LM$). Thus, $LM$ and $NO$ are about the same
diagonal [Prop.~6.26]. Let $PR$ be their (common) diagonal,
and let the (rest of the) figure be drawn. Therefore, since
the (rectangle contained) by $AF$ and $FG$ is equal to the (square) on
$EG$, thus as $AF$ is to $EG$, so $EG$ (is) to $FG$
[Prop.~6.17]. But, as $AF$ (is) to $EG$, so
$AI$ is to $EK$ [Prop.~6.1]. And as $EG$ (is)
to $FG$, so $EK$ is to $FK$ [Prop.~6.1]. And thus
as $AI$ (is) to $EK$, so $EK$ (is) to $FK$ [Prop.~5.11]. Thus, $EK$
is the mean proportional to $AI$ and $FK$. And $MN$ is also the
mean proportional to the squares $LM$ and $NO$ [Prop.~10.53~lem.]. And $AI$
is equal to $LM$, and $FK$ to $NO$. Thus, $EK$ is also equal to $MN$.
But, $MN$ is equal to $LO$, and $EK$ [is] equal to $DH$ [Prop.~1.43]. And thus the whole
of $DK$ is equal to the gnomon $UVW$ and $NO$. And $AK$ (is) also
equal to $LM$ and $NO$. Thus, the remainder $AB$ is equal to $ST$---that 
is to say, to the square on $LN$. Thus, $LN$ is the square-root of area $AB$.
I say that $LN$ is the second apotome of a medial (straight-line).

For since $AI$ and $FK$ were shown (to be)
medial (areas), and are equal to the (squares) on $LP$ and $PN$ (respectively), the (squares) on each of $LP$ and $PN$ (are) thus also
medial. Thus, $LP$ and $PN$ (are) each medial (straight-lines).
And since $AI$ is commensurable with $FK$ [Props.~6.1, 10.11], the (square) on $LP$
(is) thus also commensurable with the (square) on $PN$.  Again, since
$AI$ was shown (to be) incommensurable with $EK$, $LM$ is thus
also incommensurable with $MN$---that is to say, the (square) on 
$LP$ with the (rectangle contained) by $LP$ and $PN$.
Hence, $LP$ is also incommensurable in length with $PN$ [Props.~6.1, 10.11]. Thus, $LP$ and $PN$
are medial (straight-lines which are) commensurable in square only.
So, I say that they also contain a medial (area).

For since $EK$ was shown (to be) a medial (area), and is equal to the
(rectangle contained) by $LP$ and $PN$, the (rectangle contained) by
$LP$ and $PN$ is thus also medial. Hence, $LP$ and $PN$
are medial (straight-lines which are) commensurable in square only, and
which contain a medial (area). Thus, $LN$ is the
second apotome of a medial (straight-line) [Prop.~10.75]. And it is the square-root of area
$AB$.

Thus, the square-root of area $AB$ is the second apotome of a medial (straight-line). (Which is) the very thing it was required to show.}
\end{Parallel}

%%%%
%10.94
%%%%
\pdfbookmark[1]{Proposition 10.94}{pdf10.94}
\begin{Parallel}{}{}
\ParallelLText{
\begin{center}
{\large \ggn{94}.}
\end{center}\vspace*{-7pt}

{\greekfont Ἐὰν θωρίον περιέθηται ὑπὸ ῥητῆς καὶ ἀποτομῆς τετάρτης,
ἡ τὸ θωρίον δυναμένη ἐλάσσων ἐστίν.}

{\greekfont Θωρίον γὰρ τὸ ΑΒ περιεθέσθω ὑπὸ ῥητῆς τῆς ΑΓ καὶ ἀποτομῆς 
τετάρτης τῆς ΑΔ; λέγω, ὅτι ἡ τὸ ΑΒ θωρίον δυναμένη ἐλάσσων
ἐστίν.}

{\greekfont Ἔστω γὰρ τῇ ΑΔ προσαρμόζουσα ἡ ΔΗ; αἱ ἄρα ΑΗ, ΗΔ ῥηταί
εἰσι δυνάμει μόνον σύμμετροι, καὶ ἡ ΑΗ σύμμετρός ἐστι
τῇ ἐκκειμένῃ ῥητῇ τῇ ΑΓ μήκει, ἡ δὲ ὅλη ἡ ΑΗ τῆς
προσαρμοζούσης τῆς ΔΗ μεῖζον δύναται τῷ ἀπὸ ἀσυμμέτρου
ἑαυτῇ μήκει. ἐπεὶ οὖν ἡ ΑΗ τῆς ΗΔ μεῖζον δύναται
τῷ ἀπὸ ἀσυμμέτρου ἑαυτῇ μήκει, ἐὰν ἄρα τῷ τετάρτῳ
μέρει τοῦ ἀπὸ τῆς ΔΗ ἴσον παρὰ τὴν ΑΗ παραβληθῇ
ἐλλεῖπον εἴδει τετραγώνῳ, εἰς ἀσύμμετρα α ὐτὴν διελεῖ.
τετμήσθω οὖν ἡ ΔΗ δίθα κατὰ τὸ Ε, καὶ τῷ ἀπὸ τῆς
ΕΗ ἴσον παρὰ τὴν ΑΗ παραβεβλήσθω ἐλλεῖπον εἴδει
τετραγώνῳ, καὶ ἔστω τὸ ὑπὸ τῶν ΑΖ, ΖΗ; ἀσύμμετρος
ἄρα ἐστὶ μήκει ἡ ΑΖ τῇ ΖΗ. ἤθθωσαν οὖν διὰ τῶν Ε, Ζ, Η
παράλληλοι ταῖς ΑΓ, ΒΔ αἱ ΕΘ, ΖΙ, ΗΚ. ἐπεὶ οὖν ῥητή ἐστιν
ἡ ΑΗ καὶ σύμμετρος τῇ ΑΓ μήκει, ῥητὸν ἄρα ἐστὶν ὅλον
τὸ ΑΚ.
 πάλιν, ἐπεὶ ἀσύμμετρός ἐστιν ἡ ΔΗ τῇ ΑΓ μήκει, καί
 εἰσιν ἀμφότεραι ῥηταί, μέσον
 ἄρα ἐστὶ τὸ ΔΚ. πάλιν, ἐπεὶ ἀσύμμετρός ἐστὶν ἡ
 ΑΖ τῇ ΖΗ μήκει,
ἀσύμμετρον ἄρα καὶ τὸ ΑΙ τῷ ΖΚ.}\\~\\~\\~\\~\\~\\~\\~\\~\\

\epsfysize=1.3in
\centerline{\epsffile{Book10/fig093g.eps}}

{\greekfont Συνεστάτω οὖν τῷ μὲν
ΑΙ ἴσον τετράγωνον τὸ ΛΜ, τῷ δὲ ΖΚ ἴσον ἀφῃρήσθω
περὶ τὴν α ὐτὴν γωνίαν τὴν ὑπὸ τῶν ΛΟΜ τὸ ΝΧ. περὶ
τὴν α ὐτὴν ἄρα διάμετόν ἐστι τὰ ΛΜ, ΝΧ τετράγωνα. ἔστω
α ὐτῶν διάμετος ἡ ΟΡ, καὶ καταγεγράφθω τὸ σχῆμα.
ἐπεὶ οὖν τὸ ὑπὸ τῶν ΑΖ, ΖΗ ἴσον ἐστὶ τῷ ἀπὸ τῆς
ΕΗ, ἀνάλογον ἄρα ἐστὶν ὡς ἡ ΑΖ πρὸς τὴν ΕΗ, οὕτως
ἡ ΕΗ πρὸς τὴν ΖΗ. ἀλλ ὡς μὲν ἡ ΑΖ πρὸς τὴν ΕΗ, οὕτως
ἐστὶ τὸ ΑΙ πρὸς τὸ ΕΚ, ὡς δὲ ἡ ΕΗ πρὸς τὴν ΖΗ, οὕτως
ἐστὶ τὸ ΕΚ πρὸς τὸ ΖΚ; τῶν ἄρα ΑΙ, ΖΚ μέσον ἀνάλογόν
ἐστι τὸ ΕΚ. ἔστι δὲ καὶ τῶν ΛΜ, ΝΧ τετραγώνων μέσον ἀνάλογον
τὸ ΜΝ, καί ἐστιν ἴσον τὸ μὲν ΑΙ τῷ ΛΜ, τὸ δὲ ΖΚ τῷ ΝΧ;
καὶ τὸ ΕΚ ἄρα ἴσον ἐστὶ τῷ ΜΝ. ἀλλὰ τῷ μὲν ΕΚ ἴσον
ἐστὶ τὸ ΔΘ, τῷ δὲ ΜΝ ἴσον ἐστὶ  τὸ ΛΧ; ὅλον ἄρα τὸ ΔΚ ἴσον ἐστὶ
τῷ ΥΦΘ γνώμονι καὶ τῷ ΝΧ. ἐπεὶ οὐν ὅλον τὸ ΑΚ ἴσον ἐστὶ
τοῖς ΛΜ, ΝΧ τετραγώνοις, ὧν τὸ ΔΚ ἴσον ἐστὶ  τῷ ΥΦΘ γνώμονι καὶ
τῷ ΝΧ τετραγώνῳ, λοιπὸν ἄρα
τὸ ΑΒ ἴσον ἐστὶ τῷ ΣΤ, τουτέστι τῷ ἀπὸ τῆς ΛΝ τετραγώνῳ;
ἡ ΛΝ ἄρα δύναται τὸ ΑΒ θωρίον. λέγω, ὅτι ἡ ΛΝ ἄλογός
ἐστιν ἡ καλουμένη ἐλάσσων.}

{\greekfont Ἐπεὶ γὰρ ῥητόν ἐστι τὸ ΑΚ καί ἐστιν ἴσον τοῖς ἀπὸ τῶν
ΛΟ,  ΟΝ τετράγωνοις, τὸ ἄρα συγκείμενον ἐκ τῶν ἀπὸ
τῶν ΛΟ, ΟΝ ῥητόν ἐστιν. πάλιν, ἐπεὶ τὸ ΔΚ μέσον ἐστίν,
καί ἐστιν ἴσον τὸ ΔΚ τῷ δὶς ὑπὸ τῶν ΛΟ, ΟΝ, τὸ
ἄρα δὶς ὑπὸ τῶν ΛΟ, ΟΝ μέσον ἐστίν. καὶ ἐπεὶ ἀσύμμετρον
ἐδείθθη τὸ ΑΙ τῷ ΖΚ, ἀσύμμετρον ἄρα καὶ τὸ ἀπὸ
τῆς ΛΟ τετράγωνον τῷ ἀπὸ τῆς ΟΝ τετραγώνῳ. αἱ ΛΟ, ΟΝ
ἄρα δυνάμει εἰσὶν ἀσύμμετροι ποιοῦσαι τὸ μὲν
συγκείμενον ἐκ τῶν ἀπ α ὐτῶν τετραγώνων ῥητόν, τὸ δὲ δὶς
ὑπ α ὐτῶν μέσον. ἡ ΛΝ ἄρα ἄλογός ἐστιν ἡ καλουμένη
ἐλάσσων; καὶ δύναται τὸ ΑΒ θωρίον.}

{\greekfont Ἡ ἄρα τὸ ΑΒ θωρίον δυναμένη ἐλάσσων ἐστίν; ὅπερ
ἔδει δεῖχαι.}}

\ParallelRText{
\begin{center}
{\large Proposition 94}
\end{center}

If an area is contained by a rational (straight-line)
and a fourth apotome then the square-root of the area is a minor (straight-line).

For let the area $AB$ be contained by the rational (straight-line)
$AC$ and the fourth apotome $AD$. I say that the square-root of area $AB$
is a minor (straight-line).

For let $DG$ be an attachment to $AD$. Thus, $AG$ and $DG$ are rational (straight-lines which are) commensurable in square only [Prop.~10.73], and $AG$ is commensurable
in length with the (previously) laid down rational (straight-line) $AC$, and
the square on the whole, $AG$, is greater than (the square on) the attachment,
$DG$, by the square on (some straight-line) incommensurable
in length with ($AG$) [Def.~10.14]. Therefore,
since the square on $AG$ is greater than (the square on)  $GD$ by the
(square) on (some straight-line) incommensurable in length with ($AG$),
thus if (some area), equal to the fourth part of the (square) on $DG$,
is applied to $AG$, falling short by a square figure, then it divides
($AG$) into (parts which are) incommensurable (in length) 
[Prop.~10.18]. Therefore, let $DG$ be
cut in half at $E$, and let (some area), equal to the (square) on
$EG$, be applied to $AG$, falling short by a square figure, and
let it be the (rectangle contained) by $AF$ and $FG$. Thus, $AF$
is incommensurable in length with $FG$. Therefore, let $EH$, $FI$, and $GK$
be drawn through $E$, $F$, and $G$ (respectively), parallel
to $AC$ and $BD$. Therefore, since $AG$ is rational, and
commensurable in length with $AC$, the whole (area) $AK$
is thus rational [Prop.~10.19]. Again, since $DG$ is incommensurable in length with
$AC$, and both are rational (straight-lines), $DK$ is thus
a medial (area) [Prop.~10.21]. Again, since
$AF$ is incommensurable in length with $FG$, $AI$ (is) thus also
incommensurable with $FK$ [Props.~6.1, 10.11].

\epsfysize=1.3in
\centerline{\epsffile{Book10/fig093e.eps}}

Therefore, let the square $LM$, equal  to $AI$, be constructed. And let $NO$, equal to $FK$, (and) about the same angle, $LPM$,  be subtracted (from $LM$). Thus, the squares $LM$ and $NO$ are about the
same diagonal [Prop.~6.26]. Let $PR$ be their (common) diagonal, and let the (rest of the) figure be drawn. Therefore, since
the (rectangle contained) by $AF$ and $FG$ is equal to the (square) on $EG$, thus, proportionally, as $AF$ is to $EG$, so $EG$ (is) to $FG$
[Prop.~6.17]. 
But, as $AF$ (is) to $EG$, so
$AI$ is to $EK$, and as $EG$ (is) to $FG$, so $EK$ is to $FK$ [Prop.~6.1]. Thus, $EK$ is the mean proportional to
$AI$ and $FK$ [Prop.~5.11]. And $MN$ is also the mean proportional to the squares $LM$
and $NO$ [Prop.~10.13~lem.], and $AI$ is equal to $LM$, and $FK$ to $NO$. $EK$ is thus also equal to $MN$. But, $DH$ is equal to $EK$, and $LO$ is equal to $MN$ [Prop.~1.43]. 
Thus, the whole of $DK$ is equal to the gnomon $UVW$ and $NO$.
Therefore, since the whole of $AK$ is equal to the (sum of the) squares $LM$ and
$NO$, of which $DK$ is equal to the gnomon $UVW$ and the square
$NO$, the remainder $AB$ is thus equal to $ST$---that is to say, to the
square on $LN$. Thus, $LN$ is the square-root of area $AB$. I say
that $LN$ is the irrational (straight-line which is) called minor.

For since $AK$ is rational, and is equal to the (sum of the) squares $LP$ and $PN$,
the sum of the (squares) on $LP$ and $PN$ is thus rational. Again, since
$DK$ is medial, and $DK$ is equal to twice the (rectangle contained)
by $LP$ and $PN$, thus twice the (rectangle contained) by $LP$
and $PN$ is medial. And since $AI$ was shown (to be) incommensurable
with $FK$, the square on $LP$ (is) thus also incommensurable with
the square on $PN$. Thus, $LP$ and $PN$ are (straight-lines which are) incommensurable in square,
making the sum of the squares on them rational, and twice the (rectangle
contained) by them medial. $LN$ is thus the irrational (straight-line)
called minor [Prop.~10.76]. And it is
the square-root of area $AB$.

Thus, the square-root of area $AB$ is a minor (straight-line). (Which is)
the very thing it was required to show.}
\end{Parallel}

%%%%
%10.95
%%%%
\pdfbookmark[1]{Proposition 10.95}{pdf10.95}
\begin{Parallel}{}{}
\ParallelLText{
\begin{center}
{\large \ggn{95}.}
\end{center}\vspace*{-7pt}

{\greekfont Ἐὰν θωρίον περιέθηται ὑπὸ ῥητῆς καὶ ἀποτομῆς πέμπτης,
ἡ τὸ θωρίον δυναμένη [ἡ] μετὰ ῥητοῦ μέσον τὸ ὅλον
ποιοῦσά ἐστιν.}

\epsfysize=1.3in
\centerline{\epsffile{Book10/fig093g.eps}}

{\greekfont Θωρίον γὰρ τὸ ΑΒ περιεθέσθω ὑπὸ ῥητῆς τῆς ΑΓ καὶ ἀποτομῆς
πέμπτης τῆς ΑΔ; λέγω, ὅτι ἡ τὸ ΑΒ θωρίον δυναμένη [ἡ]
μετὰ ῥητοῦ μέσον τὸ ὅλον ποιοῦσά ἐστιν.}

{\greekfont Ἔστω γὰρ τῇ ΑΔ προσαρμόζουσα ἡ ΔΗ; αἱ ἄρα ΑΗ, ΗΔ ῥηταί
εἰσι δυνάμει μόνον σύμμετροι, καὶ ἡ προσαρμόζουσα ἡ ΗΔ
σύμμετρός ἐστι μήκει τῇ ἐκκειμένῃ ῥητῇ τῇ ΑΓ,
ἡ δὲ ὅλη ἡ ΑΗ τῆς προσαρμοζούσης τῆς ΔΗ μεῖζον δύναται
τῷ ἀπὸ ἀσυμμέτρου ἑαυτῇ. ἐὰν ἄρα τῷ τετάρτῳ μέρει
τοῦ ἀπὸ τῆς ΔΗ ἴσον παρὰ τὴν ΑΗ παραβληθῇ ἐλλεῖπον
εἴδει τετραγώνῳ, εἰς ἀσύμμετρα α ὐτὴν διελεῖ. τετμήσθω
οὖν ἡ ΔΗ δίθα κατὰ τὸ Ε σημεῖον, καὶ τῷ ἀπὸ τῆς
ΕΗ ἴσον παρὰ τὴν ΑΗ παραβεβλήσθω ἐλλεῖπον εἴδει
τετραγώνῳ καὶ ἔστω τὸ ὑπὸ τῶν ΑΖ, ΖΗ; ἀσύμμετρος
ἄρα ἐστὶν ἡ ΑΖ τῇ ΖΗ μήκει. καὶ ἐπεὶ ἀσύμμετρός
 ἐστὶν ἡ ΑΗ τῇ ΓΑ μήκει, καί εἰσιν ἀμφότεραι
ῥηταί, μέσον ἄρα ἐστὶ τὸ ΑΚ. πάλιν, ἐπεὶ ῥητή
ἐστιν ἡ ΔΗ καὶ σύμμετρος τῇ ΑΓ μήκει, ῥητόν ἐστι τὸ ΔΚ.}

{\greekfont Συνεστάτω οὖν τῷ μὲν ΑΙ ἴσον τετράγωνον τὸ ΛΜ,
τῷ δὲ ΖΚ ἴσον τετράγωνον ἀφῃρήσθω τὸ ΝΧ περὶ
τὴν α ὐτὴν γωνίαν τὴν ὑπὸ ΛΟΜ; περὶ τὴν α ὐτὴν ἄρα
διάμετρόν ἐστι τὰ ΛΜ, ΝΧ τετράγωνα. ἔστω α ὐτῶν διάμετρος
ἡ ΟΡ, καὶ καταγεγράφθω τὸ σχῆμα. ὁμοίως
δὴ δείχομεν, ὅτι ἡ ΛΝ δύναται τὸ ΑΒ θωρίον. λέγω, ὅτι
ἡ ΛΝ ἡ μετὰ ῥητοῦ μέσον τὸ ὅλον ποιοῦσά ἐστιν.}

{\greekfont Ἐπεὶ γὰρ μέσον ἐδείθθη τὸ ΑΚ καί ἐστιν ἴσον τοῖς ἀπὸ
τῶν ΛΟ, ΟΝ, τὸ ἄρα συγκείμενον ἐκ τῶν ἀπὸ
τῶν ΛΟ, ΟΝ μέσον ἐστίν. πάλιν, ἐπεὶ ῥητόν ἐστι τὸ ΔΚ καί
ἐστιν ἴσον τῷ δὶς ὑπὸ τῶν ΛΟ, ΟΝ, καὶ α ὐτὸ ῥητόν
ἐστιν. καὶ ἐπεὶ ἀσύμμετρόν ἐστι τὸ ΑΙ τῷ ΖΚ, ἀσύμμετρον
ἄρα ἐστὶ καὶ τὸ ἀπὸ τῆς ΛΟ τῷ ἀπὸ τῆς ΟΝ; αἱ
ΛΟ, ΟΝ ἄρα δυνάμει εἰσὶν ἀσύμμετροι ποιοῦσαι τὸ  μὲν συγκείμενον ἐκ τῶν ἀπ α ὐτῶν τετραγώνων μέσον, τὸ δὲ
δὶς ὑπ α ὐτῶν ῥητόν. ἡ λοιπὴ ἄρα ἡ ΛΝ ἄλογός ἐστιν ἡ
καλουμένη μετὰ ῥητοῦ μέσον τὸ ὅλον ποιοῦσα; καὶ δύναται
τὸ ΑΒ θωρίον.}

{\greekfont Ἡ τὸ ΑΒ  ἄρα θωρίον δυναμένη μετὰ ῥητοῦ μέσον τὸ ὅλον
ποιοῦσά ἐστιν; ὅπερ ἔδει δεῖχαι.}}

\ParallelRText{
\begin{center}
{\large Proposition 95}
\end{center}

If an area is contained by a rational (straight-line)
and a fifth apotome then the square-root of the area is that (straight-line) which with a rational (area) makes a medial whole.

\epsfysize=1.3in
\centerline{\epsffile{Book10/fig093e.eps}}

For let the area $AB$ be contained by the rational (straight-line)
$AC$ and the fifth apotome $AD$. I say that the square-root of area $AB$
is  that (straight-line) which with a rational (area) makes a medial whole.

For let $DG$ be an attachment to $AD$. Thus, $AG$ and $DG$ are rational (straight-lines which are) commensurable in square only [Prop.~10.73], and the attachment $GD$ is
commensurable in length the the (previously) laid down rational (straight-line) $AC$, and the square on the whole, $AG$, is greater than (the square on)
the attachment, $DG$, by the (square) on (some straight-line)
incommensurable (in length) with ($AG$) [Def.~10.15]. 
Thus, if (some area), equal to the fourth part of the (square) on $DG$,
is applied to $AG$, falling short by a square figure, then it divides ($AG$)
into (parts which are) incommensurable (in length) [Prop.~10.18]. Therefore, let $DG$
be divided in half at point $E$, and let (some area), equal to the
(square) on $EG$, be applied to $AG$, falling short by a
square figure, and let it be the (rectangle contained) by $AF$ and $FG$.
Thus, $AF$ is incommensurable in length with $FG$. And since
$AG$ is incommensurable in length with $CA$, and both are
rational (straight-lines), $AK$ is thus a medial (area) [Prop.~10.21]. Again, since $DG$ is rational,
and commensurable in length with $AC$, $DK$ is a rational (area)
[Prop.~10.19].

Therefore, let the square $LM$, equal to $AI$, be constructed. 
And let the square $NO$, equal to $FK$, (and) about the same angle, $LPM$,
be subtracted (from $NO$). Thus, the squares $LM$ and
$NO$ are about the same diagonal [Prop.~6.26]. 
Let $PR$ be their (common) diagonal, and let (the rest of) the figure be drawn.
So, similarly (to the previous propositions), we can show that $LN$ is the square-root of area $AB$.
I say that $LN$ is that (straight-line) which with a rational (area) makes a medial whole.

For since $AK$ was shown (to be) a medial (area), and is equal to (the sum of) the squares on $LP$ and $PN$, the sum of the (squares) on $LP$
and $PN$ is thus medial. Again, since $DK$ is rational, and
is equal to twice the (rectangle contained) by $LP$ and $PN$,  (the latter) is also
rational. And since $AI$ is incommensurable with $FK$, the (square) on
$LP$ is thus also incommensurable with the (square) on $PN$. Thus,
$LP$ and $PN$ are (straight-lines which are) incommensurable in square, making the
sum of the squares on them medial, and twice the (rectangle contained) by
them rational. Thus, the remainder $LN$ is the irrational (straight-line)
called that which with a rational (area) makes a medial whole [Prop.~10.77].
And it is the square-root of area $AB$.

Thus, the square-root of area $AB$ is that (straight-line) which with a rational (area) makes a medial whole. (Which is) the very thing it was required to show.}
\end{Parallel}

%%%%
%10.96
%%%%
\pdfbookmark[1]{Proposition 10.96}{pdf10.96}
\begin{Parallel}{}{}
\ParallelLText{
\begin{center}
{\large \ggn{96}.}
\end{center}\vspace*{-7pt}

{\greekfont Ἐὰν θωρίον περιέθηται ὑπὸ ῥητῆς καὶ ἀποτομῆς ἕκτης, ἡ τὸ
θωρίον δυναμένη μετὰ μέσου μέσον τὸ ὅλον ποιοῦσά ἐστιν.}\\

\epsfysize=1.3in
\centerline{\epsffile{Book10/fig093g.eps}}

{\greekfont Θωρίον γὰρ τὸ ΑΒ περιεθέσθω ὑπὸ ῥητῆς τῆς ΑΓ καὶ ἀποτομῆς 
ἕκτης τῆς ΑΔ; λέγω, ὅτι ἡ τὸ ΑΒ θωρίον δυναμένη [ἡ]
μετὰ μέσου μέσον τὸ ὅλον ποιοῦσά ἐστιν.}

{\greekfont Ἔστω γὰρ  τῇ ΑΔ προσαρμόζουσα ἡ ΔΗ; αἱ ἄρα ΑΗ, ΗΔ ῥηταί
εἰσι δυνάμει μόνον σύμμετροι, καὶ οὐδετέρα α ὐτῶν σύμμετρός
ἐστι τῇ ἐκκειμένῃ ῥητῇ τῇ ΑΓ μήκει, ἡ δὲ ὅλη ἡ ΑΗ
τῆς προσαρμοζούσης τῆς ΔΗ μεῖζον δύναται τῷ ἀπὸ ἀσυμμέτρου
ἑαυτῇ μήκει. ἐπεὶ οὖν ἡ ΑΗ τῆς ΗΔ μεῖζον δύναται τῷ ἀπὸ
ἁσυμμέτρου ἐαυτῇ μήκει, ἑὰν ἄρα τῷ τετάρτῳ μέρει
τοῦ ἀπὸ τῆς ΔΗ ἴσον παρὰ τὴν ΑΗ παραβληθῇ ἐλλεῖπον
εἴδει τετραγώνῳ, εἰς ἀσύμμετρα α ὐτὴν διελεῖ. τετμήσθω
οὖν ἡ ΔΗ δίθα κατὰ τὸ Ε [σημεῖον], καὶ τῷ ἀπὸ τῆς
ΕΗ ἴσον παρὰ τὴν ΑΗ παραβεβλήσθω ἐλλεῖπον εἴδει τετραγώνῳ,
καὶ ἔστω τὸ ὑπὸ τῶν ΑΖ, ΖΗ; ἀσύμμετρος ἄρα ἐστὶν
ἡ ΑΖ τῇ ΖΗ μήκει. ὡς δὲ ἡ ΑΖ πρὸς τὴν ΖΗ, οὕτως ἐστὶ τὸ ΑΙ πρὸς τὸ ΖΚ; ἀσύμμετρον
ἄρα ἐστὶ τὸ ΑΙ τῷ ΖΚ. καὶ ἐπεὶ αἱ ΑΗ, ΑΓ ῥηταί
εἰσι δυνάμει μόνον σύμμετροι, μέσον ἐστὶ τὸ ΑΚ. πάλιν, ἐπεὶ αἱ ΑΓ, ΔΗ ῥηταί εἰσι καὶ ἀσύμμετροι μήκει, μέσον ἐστὶ καὶ τὸ
ΔΚ. ἐπεὶ οὖν αἱ ΑΗ, ΗΔ δυνάμει μόνον σύμμετροί εἰσιν,
ἀσύμμετρος ἄρα ἐστὶν ἡ ΑΗ τῇ ΗΔ μήκει. ὡς δὲ ἡ ΑΗ πρὸς τὴν ΗΔ, οὕτως ἐστὶ τὸ ΑΚ πρὸς τὸ ΚΔ; ἀσύμμετρον ἄρα ἐστὶ
τὸ ΑΚ τῷ ΚΔ.}

{\greekfont Συνεστάτω οὖν τῷ μὲν ΑΙ ἴσον τετράγωνον τὸ ΛΜ, τῷ δὲ
ΖΚ ἴσον ἀφῃρήσθω περὶ τὴν α ὐτὴν γωνίαν τὸ ΝΧ; περὶ τὴν
α ὐτὴν ἄρα διάμετρόν ἐστι τὰ ΛΜ, ΝΧ τετράγωνα. ἔστω
α ὐτῶν διάμετρος ἡ ΟΡ, καὶ καταγεγράφθω τὸ σχῆμα.
ὁμοίως δὴ τοῖς ἐπάνω δείχομεν, ὅτι ἡ ΛΝ δύναται τὸ ΑΒ
θωρίον. λέγω, ὅτι ἡ ΛΝ [ἡ] μετὰ μέσου μέσον τὸ ὅλον ποιοῦσά
ἐστιν.}

{\greekfont Ἐπεὶ γὰρ μέσον ἐδείθθη τὸ ΑΚ καί ἐστιν ἴσον τοῖς ἀπὸ τῶν
ΛΟ, ΟΝ, τὸ ἄρα συγκείμενον ἐκ τῶν ἀπὸ τῶν ΛΟ, ΟΝ
μέσον ἐστίν. πάλιν, ἐπεὶ μέσον ἐδείθθη τὸ ΔΚ καί ἐστιν  ἴσον
τῷ δὶς ὑπὸ  τῶν ΛΟ, ΟΝ, καὶ τὸ δὶς ὑπὸ τῶν ΛΟ, ΟΝ
μέσον ἐστίν. καὶ ἐπεὶ ἀσύμμετρον ἐδείθθη τὸ ΑΚ τῷ ΔΚ, ἀσύμμετρα [ἄρα] ἐστὶ καὶ τὰ ἀπὸ τῶν ΛΟ, ΟΝ τετράγωνα
τῷ δὶς ὑπὸ τῶν ΛΟ, ΟΝ. καὶ ἐπεὶ ἀσύμμετρόν ἐστι τὸ ΑΙ τῷ
ΖΚ, ἀσύμμετρον ἄρα καὶ τὸ ἀπὸ τῆς ΛΟ τῷ
ἀπὸ τῆς ΟΝ; αἱ ΛΟ, ΟΝ ἄρα δυνάμει εἰσὶν
ἀσύμμετροι ποιοῦσαι τό τε συγκείμενον ἐκ τῶν ἀπ
α ὐτῶν τετραγώνων μέσον καὶ τὸ δὶς ὑπ α ὐτῶν
μέσον ἔτι τε τὰ ἀπ α ὐτῶν τετράγωνα ἀσύμμετρα τῷ δὶς ὑπ
α ὐτῶν. ἡ ἄρα ΛΝ ἄλογός ἐστιν ἡ καλουμέμ η μετὰ μέσου
μέσον τὸ ὅλον ποιοῦσα; καὶ δύναται τὸ ΑΒ θωρίον.}

{\greekfont Ἡ ἄρα τὸ θωρίον δυναμένη μετὰ μέσου μέσον τὸ ὅλον
ποιοῦσά ἐστιν; ὅπερ ἔδει δεῖχαι.}}

\ParallelRText{
\begin{center}
{\large Proposition 96}
\end{center}

If an area is contained by a rational (straight-line)
and a sixth apotome then the square-root of the area is that (straight-line) which with a medial (area) makes a medial whole.

\epsfysize=1.3in
\centerline{\epsffile{Book10/fig093e.eps}}

For let the area $AB$ be contained by the rational (straight-line)
$AC$ and the sixth apotome $AD$. I say that the square-root of area $AB$
is  that (straight-line) which with a medial (area) makes a medial whole.

For let $DG$ be an attachment to $AD$. Thus, $AG$ and $GD$ are rational (straight-lines which are) commensurable in square only [Prop.~10.73], and neither of them is
commensurable in length with the (previously) laid down rational (straight-line) $AC$, and the square on the whole, $AG$, is greater than (the square on)
the attachment, $DG$, by the (square) on (some straight-line)
incommensurable in length with ($AG$) [Def.~10.16]. Therefore, since the square on $AG$
is greater than (the square on) $GD$ by the (square) on (some straight-line)
incommensurable in length with ($AG$), thus if (some area), equal to
the fourth part of square on $DG$, is applied to $AG$, falling short
by a square figure, then it divides ($AG$) into (parts which are) incommensurable
(in length) [Prop.~10.18]. Therefore, let $DG$ be cut in  half at [point] $E$. And let (some area), equal to the
(square) on $EG$, be applied to $AG$, falling short by a square figure.
And let it be the (rectangle contained) by $AF$ and $FG$. $AF$ is
thus incommensurable in length with $FG$. And as $AF$ (is) to $FG$,
so $AI$ is to $FK$ [Prop.~6.1].  Thus, $AI$
is incommensurable with $FK$ [Prop.~10.11]. 
And since $AG$ and $AC$ are rational (straight-lines which are)
commensurable in square only, $AK$ is a medial (area) [Prop.~10.21]. Again, since $AC$ and $DG$
are rational (straight-lines which are) incommensurable in length, $DK$
is also a medial (area) [Prop.~10.21]. Therefore,
since $AG$ and $GD$ are commensurable in square only,  $AG$
is thus incommensurable in length with $GD$. And as $AG$ (is) to $GD$,
so $AK$ is to $KD$ [Prop.~6.1]. Thus, $AK$
is incommensurable with $KD$ [Prop.~10.11].

Therefore, let the square $LM$, equal to $AI$, be constructed.
And let $NO$, equal to $FK$, (and) about the same angle, be subtracted (from $LM$). Thus, the squares $LM$ and $NO$ are about the
same diagonal [Prop.~6.26]. 
Let $PR$ be their (common) diagonal, and let (the rest of) the figure be drawn.
So, similarly to the above, we can show that $LN$ is the square-root of
area $AB$. I say that $LN$ is  that (straight-line) which with a medial (area) makes a medial whole.

For since $AK$ was shown (to be) a medial (area), and is equal to the
(sum of the) squares on $LP$ and $PN$, the sum of the (squares) on
$LP$ and $PN$ is medial. Again, since $DK$ was shown (to
be) a medial (area), and is equal to twice the (rectangle contained)
by $LP$ and $PN$, twice the (rectangle contained) by $LP$ and
$PN$ is also medial. And since $AK$ was shown (to be) incommensurable
with $DK$, [thus] the (sum of the) squares on $LP$ and $PN$ is also
incommensurable with twice the (rectangle contained) by $LP$ and $PN$.
And since $AI$ is incommensurable with $FK$, the (square) on
$LP$ (is) thus also incommensurable with the (square) on $PN$. Thus, $LP$ and $PN$
are (straight-lines which are) incommensurable in square, making the sum of the squares on them
medial, and twice the (rectangle contained) by  medial, and, furthermore,
the (sum of the) squares on them incommensurable with twice the
(rectangle contained) by them. Thus, $LN$ is the irrational (straight-line) called  that  which with a medial (area) makes a medial whole [Prop.~10.78]. And it is the square-root of area $AB$.

Thus, the square-root of area ($AB$) is that (straight-line) which with a medial (area) makes a medial whole. (Which is) the very thing it was required to show.\\}
\end{Parallel}

%%%%
%10.97
%%%%
\pdfbookmark[1]{Proposition 10.97}{pdf10.97}
\begin{Parallel}{}{}
\ParallelLText{
\begin{center}
{\large \ggn{97}.}
\end{center}\vspace*{-7pt}

{\greekfont Τὸ ἀπὸ ἀποτομῆς παρὰ ῥητὴν παραβαλλόμενον πλάτος ποιεῖ
ἀποτομὴν πρώτην.}

{\greekfont Ἔστω ἀποτομὴ ὴ ΑΒ, ῥητὴ δὲ ἡ ΓΔ, καὶ τῷ ἀπὸ τῆς ΑΒ ἴσον 
παρὰ τὴν ΓΔ παραβεβλήσθω τὸ ΓΕ πλάτος ποιοῦν τὴν ΓΖ;
λέγω, ὅτι ἡ ΓΖ ἀποτομή ἐστι πρώτη.}

\epsfysize=1.6in
\centerline{\epsffile{Book10/fig097g.eps}}

{\greekfont Ἔστω γὰρ τῇ ΑΒ προσαρμόζουσα ἡ ΒΗ; αἱ ἄρα ΑΗ, ΗΒ ῥηταί
εἰσι δυνάμει μόνον σύμμετροι. καὶ τῷ μὲν ἀπὸ τῆς ΑΗ ἴσον
παρὰ τὴν ΓΔ παραβεβλήσθω τὸ ΓΘ, τῷ δὲ ἀπὸ τῆς ΒΗ τὸ ΚΛ. ὅλον
ἄρα τὸ ΓΛ ἴσον ἐστὶ τοῖς ἀπὸ τῶν ΑΗ, ΗΒ; ὧν τὸ ΓΕ ἴσον
ἐστὶ τῷ ἀπὸ τῆς ΑΒ; λοιπὸν ἄρα τὸ ΖΛ ἴσον ἐστὶ τῷ δὶς
ὑπὸ τῶν ΑΗ, ΗΒ. τετμήσθω ἡ ΖΜ δίθα κατὰ τὸ Ν σημεῖον,
καὶ ἤθθω διὰ τοῦ Ν τῇ ΓΔ παράλληλος ἡ ΝΧ; ἑκάτερον ἄρα
τῶν ΖΧ, ΛΝ ἴσον ἐστὶ τῷ ὑπὸ τῶν ΑΗ, ΗΒ. καὶ ἐπεὶ τὰ  ἀπὸ
τῶν ΑΗ, ΗΒ ῥητά ἐστιν, καί ἐστι τοῖς ἀπὸ τῶν ΑΗ, ΗΒ ἴσον τὸ
ΔΜ, ῥητὸν ἄρα ἐστὶ τὸ ΔΜ. καὶ παρὰ ῥητὴν τὴν ΓΔ παραβέβληται
πλάτος ποιοῦν τὴν ΓΜ; ῥητὴ ἄρα ἐστὶν ἡ ΓΜ καὶ σύμμετρος
τῇ ΓΔ μήκει. πάλιν, ἐπεὶ μέσον ἐστὶ τὸ δὶς ὑπὸ τῶν ΑΗ, ΗΒ,
καὶ τῷ δὶς ὑπὸ τῶν ΑΗ, ΗΒ ἴσον τὸ ΖΛ, μέσον ἄρα τὸ
ΖΛ. καὶ παρὰ ῥητὴν τὴν ΓΔ παράκειται πλάτος ποιοῦν τὴν ΖΜ;
ῥητὴ ἄρα ἐστὶν ἡ ΖΜ καὶ ἀσύμμετρος τῇ ΓΔ μήκει.
καὶ ἐπεὶ τὰ μὲν ἀπὸ τῶν ΑΗ, ΗΒ ῥητά ἐστιν, τὸ δὲ δὶς ὑπὸ
τῶν ΑΗ, ΗΒ μέσον, ἀσύμμετρα ἄρα ἐστὶ τὰ ἀπὸ τῶν ΑΗ, ΗΒ
τῷ  δὶς ὑπὸ τῶν ΑΗ, ΗΒ. καὶ τοῖς μὲν ἀπὸ τῶν ΑΗ, ΗΒ
ἴσον ἐστὶ τὸ ΓΛ, τῷ δὲ δὶς ὑπὸ τῶν ΑΗ, ΗΒ τὸ ΖΛ; ἀσύμμετρον
ἄρα ἐστὶ τὸ ΔΜ τῷ ΖΛ. ὡς δὲ τὸ ΔΜ πρὸς τὸ ΖΛ,
οὕτως ἐστὶν ἡ ΓΜ πρὸς τὴν ΖΜ. ἀσύμμετρος ἄρα ἐστὶν
ἡ ΓΜ
 τῇ ΖΜ μήκει. καί εἰσιν ἀμφότεραι
ῥηταί; αἱ ἄρα ΓΜ, ΜΖ ῥηταί εἰσι δυνάμει μόνον σύμμετροι;
ἡ ΓΖ ἄρα ἀποτομή ἐστιν. λέγω δή, ὅτι καὶ πρώτη.}

{\greekfont Ἐπεὶ γὰρ τῶν ἀπὸ τῶν ΑΗ, ΗΒ μέσον ἀνάλογόν ἐστι τὸ
ὑπὸ τῶν ΑΗ, ΗΒ, καί ἐστι τῷ μὲν ἀπὸ τῆς ΑΗ ἴσον
τὸ ΓΘ, τῷ δὲ ἀπὸ τῆς ΒΗ ἴσον τὸ ΚΛ, τῷ δὲ ὐπὸ
τῶν ΑΗ, ΗΒ τὸ ΝΛ, καὶ τῶν ΓΘ, ΚΛ ἄρα μέσον ἀνάλογόν
ἐστι τὸ ΝΛ; ἔστιν ἄρα ὡς τὸ ΓΘ πρὸς τὸ ΝΛ, οὕτως
τὸ ΝΛ πρὸς τὸ ΚΛ. ἀλλ ὡς μὲν τὸ ΓΘ πρὸς τὸ ΝΛ, οὕτως
ἐστὶν ἡ ΓΚ πρὸς τὴν ΝΜ; ὡς δὲ τὸ ΝΛ πρὸς τὸ ΚΛ,
οὕτως ἐστὶν  ἡ ΝΜ πρὸς τὴν ΚΜ; τὸ ἄρα ὑπὸ τῶν ΓΚ, ΚΜ
ἴσον ἐστὶ τῷ ἀπὸ τῆς ΝΜ, τουτέστι τῷ τετάρτῳ μέρει
τοῦ ἀπὸ τῆς ΖΜ. καὶ επεὶ σύμμετρόν ἐστι τὸ ἀπὸ τῆς
ΑΗ τῷ ἀπὸ τῆς ΗΒ, σύμμετρόν [ἐστι] καὶ τὸ ΓΘ τῷ ΚΛ.
ὡς δὲ τὸ ΓΘ πρὸς τὸ ΚΛ, οὕτως ἡ ΓΚ
πρὸς τὴν ΚΜ; σύμμετρος ἄρα ἐστὶν ἡ ΓΚ τῇ ΚΜ.
ἐπεὶ οὖν δύο εὐθεῖαι ἄνισοί εἰσιν αἱ ΓΜ, ΜΖ, καὶ
τῷ τετάρτῳ μέρει τοῦ ἀπὸ τῆς ΖΜ ἴσον παρὰ τὴν ΓΜ
παραβέβληται ἐλλεῖπον εἴδει τετραγώνῳ τὸ ὑπὸ τῶν
ΓΚ, ΚΜ, καί ἐστι σύμμετρος ἡ ΓΚ τῇ ΚΜ, ἡ ἄρα ΓΜ τῆς
ΜΖ μεῖζον δύναται τῷ ἀπὸ συμμέτρου ἑαυτῇ μήκει.
καί ἐστιν ἡ ΓΜ σύμμετρος τῇ ἐκκειμένῃ ῥητῇ τῇ ΓΔ
μήκει; ἡ ἄρα ΓΖ ἀποτομή ἐστι πρώτη.}

{\greekfont Τὸ ἄρα ἀπὸ ἀποτομῆς παρὰ ῥητὴν παραβαλλόμενον
πλάτος ποιεῖ ἀποτομὴν πρώτην; ὅπερ ἔδει δεῖχαι.}}

\ParallelRText{
\begin{center}
{\large Proposition 97}
\end{center}

The (square) on an apotome, applied
to a rational (straight-line), produces  a first apotome as breadth.

Let $AB$ be an apotome, and $CD$ a rational (straight-line). And
let $CE$, equal to the (square) on $AB$, be applied to
$CD$, producing $CF$ as breadth. I say that $CF$ is a first apotome.

\epsfysize=1.6in
\centerline{\epsffile{Book10/fig097e.eps}}

For let $BG$ be an attachment to $AB$. Thus, $AG$ and $GB$
are rational (straight-lines which are) commensurable in square only
[Prop.~10.73]. And let $CH$, equal to the
(square) on $AG$, and $KL$, (equal) to the (square) on $BG$, be applied to $CD$. Thus, the whole of $CL$ is equal to the (sum of the
squares) on $AG$ and $GB$, of which $CE$ is equal to the (square) on
$AB$. The remainder $FL$ is thus equal to twice the (rectangle contained)
by $AG$ and $GB$ [Prop.~2.7]. Let $FM$
be cut in half at point $N$. And let $NO$ be drawn
through $N$, parallel to $CD$. Thus, $FO$ and $LN$ are each equal to
the (rectangle contained) by $AG$ and $GB$. And since the (sum of the
squares) on $AG$ and $GB$ is rational, and $DM$ is equal to the (sum of the squares) on
$AG$ and $GB$, $DM$ is thus rational. And it has
been applied to the rational (straight-line) $CD$, producing $CM$ as breadth. Thus, $CM$
is rational, and commensurable in length with $CD$ [Prop.~10.20]. Again, since twice the
(rectangle contained) by $AG$ and $GB$ is medial, and $FL$
(is) equal to twice the (rectangle contained) by $AG$ and $GB$, $FL$
(is) thus a medial (area). And it is applied to the rational (straight-line)
$CD$, producing $FM$ as breadth. $FM$ is thus rational, and
incommensurable in  length with $CD$ [Prop.~10.22]. And since the (sum of the
squares) on $AG$ and $GB$ is rational, and twice the (rectangle contained)
by $AG$ and $GB$ medial, the (sum of the squares) on
$AG$ and $GB$ is thus incommensurable with twice the (rectangle
contained) by $AG$ and $GB$. And $CL$ is equal to the (sum of the
squares) on $AG$ and $GB$, and $FL$ to twice the (rectangle contained)
by $AG$ and $GB$. $DM$ is thus incommensurable with $FL$.
And as $DM$ (is) to $FL$, so $CM$ is to $FM$ [Prop.~6.1]. $CM$ is thus incommensurable in length
with $FM$ [Prop.~10.11]. And both are rational 
(straight-lines). Thus, $CM$ and $MF$ are rational (straight-lines
which are) commensurable in square only. $CF$ is thus an apotome
[Prop.~10.73]. So, I say that
(it is) also a first (apotome).

For since the (rectangle contained) by $AG$ and $GB$ is the
mean proportional to the (squares) on $AG$ and $GB$ [Prop.~10.21 lem.], and $CH$ is equal
to the (square) on $AG$, and $KL$ equal to the (square) on 
$BG$, and $NL$ to the (rectangle contained) by $AG$ and $GB$,
$NL$ is thus also the mean proportional to $CH$ and $KL$. Thus,
as $CH$ is to $NL$, so $NL$ (is) to $KL$. But, as $CH$ (is) to
$NL$, so $CK$ is to $NM$, and as $NL$ (is) to $KL$, so
$NM$ is to $KM$ [Prop.~6.1]. 
Thus, the (rectangle contained) by $CK$ and $KM$
is equal to the (square) on $NM$---that is to say, to the fourth
part of the (square) on $FM$ [Prop.~6.17]. 
And since the (square) on $AG$ is commensurable with the
(square) on $GB$, $CH$ [is] also commensurable with
$KL$. And as $CH$ (is) to $KL$, so $CK$ (is) to $KM$ [Prop.~6.1]. $CK$ is thus
commensurable (in length) with $KM$ [Prop.~10.11].  Therefore, since $CM$ and
$MF$ are two unequal straight-lines, and the (rectangle
contained) by $CK$ and $KM$, equal to the fourth part of the
(square) on $FM$, has been applied to $CM$, falling short by a square figure, and $CK$ is commensurable
(in length) with $KM$,  the square on $CM$ is thus greater than
(the square on) $MF$ by the (square) on (some straight-line)
commensurable in length with ($CM$) [Prop.~10.17]. 
And $CM$ is commensurable in length with the (previously) laid down rational (straight-line) $CD$. Thus, $CF$ is a first apotome [Def.~10.15].

Thus, the (square) on an apotome, applied to a rational (straight-line),
produces  a first apotome as breadth. (Which is) the very thing it was required to show.}
\end{Parallel}

%%%%
%10.98
%%%%
\pdfbookmark[1]{Proposition 10.98}{pdf10.98}
\begin{Parallel}{}{}
\ParallelLText{
\begin{center}
{\large \ggn{98}.}
\end{center}\vspace*{-7pt}

{\greekfont Τὸ ἀπὸ μέσης ἀποτομῆς πρώτης παρὰ ῥητὴν παραβαλλόμενον
πλάτος ποιεῖ ἀποτομὴν δευτέραν.}

{\greekfont Ἔστω μέσης ἀποτομὴ πρώτη ἡ ΑΒ, ῥητὴ δὲ ἡ ΓΔ,
καὶ τῷ ἀπὸ τῆς ΑΒ ἴσον παρὰ τὴν ΓΔ παραβεβλήσθω
τὸ ΓΕ πλάτος ποιοῦν τὴν ΓΖ; λέγω, ὅτι ἡ ΓΖ ἀποτομή
ἐστι δευτέρα.}

{\greekfont Ἔστω γὰρ τῇ ΑΒ προσαρμόζουσα ἡ ΒΗ; αἱ ἄρα ΑΗ, ΗΒ μέσαι
εἰσὶ δυνάμει μόνον σύμμετροι ῥητὸν περιέθουσαι. καὶ τῷ
μὲν ἀπὸ τῆς ΑΗ ἴσον παρὰ τὴν ΓΔ παραβεβλήσθω τὸ ΓΘ
πλάτος ποιοῦν τὴν ΓΚ, τῷ δὲ ἀπὸ τῆς ΗΒ ἴσον τὸ ΚΛ
πλάτος ποιοῦν τὴν ΚΜ; ὅλον ἄρα τὸ ΓΛ ἴσον ἐστὶ τοῖς
ἀπὸ τῶν ΑΗ, ΗΒ; μέσον ἄρα καὶ τὸ ΓΛ. καὶ παρὰ ῥητὴν
τὴν ΓΔ παράκειται πλάτος ποιοῦν τὴν ΓΜ; ῥητὴ ἄρα ἐστὶν
ἡ ΓΜ καὶ ἀσύμμετρος τῇ ΓΔ μήκει. καὶ ἐπεὶ τὸ ΓΛ
ἴσον ἐστὶ τοῖς ἀπὸ τῶν ΑΗ, ΗΒ, ὧν τὸ ἀπὸ τῆς
ΑΒ ἴσον ἐστὶ τῷ ΓΕ, λοιπὸν ἄρα τὸ δὶς ὑπὸ τῶν
ΑΗ, ΗΒ ἴσον ἐστὶ τῷ ΖΛ. ῥητὸν δέ [ἐστι] τὸ δὶς ὑπὸ
τῶν ΑΗ, ΗΒ; ῥητὸν ἄρα τὸ ΖΛ. καὶ παρὰ ῥητὴν τὴν ΖΕ
παράκειται πλάτος ποιοῦν τὴν ΖΜ; ῥητὴ ἄρα ἐστὶ καὶ ἡ ΖΜ
καὶ σύμμετρος τῇ ΓΔ μήκει. ἐπεὶ οὖν τὰ μὲν ἀπὸ
τῶν ΑΗ, ΗΒ, τουτέστι τὸ ΓΛ, μέσον ἐστίν, τὸ δὲ δὶς ὑπὸ
τῶν ΑΗ, ΗΒ, τουτέστι τὸ ΖΛ, ῥητόν ἀσύμμετρον ἄρα ἐστὶ
τὸ ΓΛ τῷ ΖΛ. ὡς δὲ τὸ ΓΛ πρὸς τὸ ΖΛ, οὕτως
ἐστὶν ἡ ΓΜ πρὸς τὴν ΖΜ; ἀσύμμετρος ἄρα ἡ ΓΜ τῇ ΖΜ
μήκει. καί εἰσιν ἀμφότεραι ῥηταί; αἱ ἄρα ΓΜ, ΜΖ ῥηταί
εἰσι δυνάμει μόνον σύμμετροι; ἡ ΓΖ ἄρα ἀποτομή ἐστιν.
λέγω δή, ὅτι καὶ δευτέρα.}\\~\\~\\~\\~\\~\\~\\~\\~\\

\epsfysize=1.6in
\centerline{\epsffile{Book10/fig097g.eps}}

{\greekfont Τετμήσθω γὰρ ἡ ΖΜ δίθα κατὰ τὸ Ν, καὶ ἤθθω διὰ τοῦ
Ν τῇ ΓΔ παράλληλος ἡ ΝΧ; ἑκάτερον ἄρα τῶν
ΖΧ, ΝΛ  ἴσον ἐστὶ τῷ ὑπὸ τῶν ΑΗ, ΗΒ. καὶ ἐπεὶ
τῶν ἀπὸ τῶν ΑΗ, ΗΒ τετραγώνων μέσον ἀνάλογόν ἐστι τὸ
ὑπὸ τῶν ΑΗ, ΗΒ, καί ἐστιν ἴσον τὸ μὲν ἀπὸ τῆς ΑΗ τῷ
ΓΘ, τὸ δὲ ὑπὸ τῶν ΑΗ, ΗΒ τῷ ΝΛ, τὸ δὲ ἀπὸ τῆς ΒΗ τῷ
ΚΛ, καὶ τῶν ΓΘ, ΚΛ ἄρα μέσον ἀνάλογόν ἐστι τὸ ΝΛ; ἔστιν
ἄρα ὡς τὸ ΓΘ πρὸς τὸ ΝΛ, οὕτως τὸ ΝΛ πρὸς τὸ ΚΛ.
ἀλλ ὡς μὲν τὸ ΓΘ πρὸς τὸ ΝΛ, οὕτως ἐστὶν ἡ ΓΚ πρὸς τὴν
ΝΜ, ὡς δὲ τὸ ΝΛ πρὸς τὸ ΚΛ, οὕτως
ἐστὶν ἡ ΝΜ πρὸς τὴν ΜΚ; ὡς ἄρα ἡ ΓΚ πρὸς τὴν ΝΜ, οὕτως
ἐστὶν ἡ ΝΜ πρὸς τὴν ΚΜ; τὸ ἄρα ὑπὸ τῶν ΓΚ, ΚΜ
ἴσον ἐστὶ τῷ ἀπὸ τῆς ΝΜ, τουτέστι τῷ τετάρτῳ μέρει τοῦ
ἀπὸ τῆς ΖΜ [καὶ ἐπεὶ σύμμετρόν ἐστι τὸ ἀπὸ τῆς ΑΗ
τῷ ἀπὸ τῆς ΒΗ, σύμμετρόν ἐστι καὶ τὸ ΓΘ τῷ ΚΛ, τουτέστιν
ἡ ΓΚ τῇ ΚΜ]. ἐπεὶ οὖν δύο εὐθεῖαι ἄνισοί εἰσιν αἱ ΓΜ, ΜΖ,
καὶ τῷ τετάτρῳ μέρει τοῦ ἀπὸ τῆς ΜΖ ἴσον παρὰ τὴν μείζονα
τὴν ΓΜ παραβέβληται ἐλλεῖπον εἴδει τετραγώνῳ τὸ ὑπὸ τῶν
ΓΚ, ΚΜ καὶ εἰς σύμμετρα α ὐτὴν διαιρεῖ, ἡ ἄρα ΓΜ τῆς
ΜΖ μεῖζον δύναται τῷ ἀπὸ συμμέτρου ἑαυτῇ μήκει. καί
ἐστιν ἡ προσαρμόζουσα ἡ ΖΜ σύμμετρος μήκει τῇ ἐκκειμένῃ
ῥητῇ τῇ ΓΔ; ἡ ἄρα ΓΖ ἀποτομή ἐστι δευτέρα.}

{\greekfont Τὸ ἄρα ἀπὸ μέσης ἀποτομῆς πρώτης παρὰ ῥητὴν παραβαλλόμενον πλάτος ποιεῖ ἀποτομὴν δευτέραν; ὅπερ ἔδει δεῖχαι.}}

\ParallelRText{
\begin{center}
{\large Proposition 98}
\end{center}

The (square) on a first apotome of
a medial (straight-line), applied to a rational (straight-line), produces
 a second apotome as breadth.
  
Let $AB$ be a first apotome of a medial (straight-line), and
$CD$ a rational (straight-line). And let $CE$, equal to the (square) on
$AB$, be applied to $CD$, producing
$CF$ as breadth. I say that $CF$ is a second apotome.

For let $BG$ be an attachment to $AB$. Thus,
$AG$ and $GB$ are medial (straight-lines which are)
commensurable in square only, containing a rational (area)
[Prop.~10.74]. And let $CH$, equal to
the (square) on $AG$, be applied to $CD$, producing $CK$
as breadth, and $KL$, equal to the (square) on $GB$, producing $KM$ as
breadth. Thus, the whole of $CL$ is equal to the (sum of the squares)
on $AG$ and $GB$. Thus, $CL$ (is) also a medial (area) [Props.~10.15, 10.23~corr.]. And it is
applied to the rational (straight-line) $CD$, producing $CM$ as breadth.
$CM$ is thus rational, and incommensurable in length with
$CD$ [Prop.~10.22]. And since $CL$ is equal to the
(sum of the squares) on $AG$ and $GB$, of which the (square)
on $AB$ is equal to $CE$, 
the remainder, twice  the (rectangle contained) by $AG$ and $GB$, is thus equal to $FL$ [Prop.~2.7]. And twice the (rectangle contained)
by $AG$ and $GB$ [is] rational. Thus, $FL$ (is) rational. And it
is applied to the rational (straight-line) $FE$, producing $FM$ as breadth.
$FM$ is thus also rational, and commensurable in length with $CD$
[Prop.~10.20]. Therefore, since the (sum of the squares) on $AG$ and $GB$---that is to say, $CL$---is medial, 
and twice the (rectangle contained) by $AG$ and $GB$---that is to say,
$FL$---(is) rational, $CL$ is thus incommensurable with $FL$. 
And as $CL$ (is) to $FL$, so $CM$ is to $FM$ [Prop.~6.1]. Thus, $CM$ (is) incommensurable
in length with $FM$ [Prop.~10.11]. And they are both rational (straight-lines). Thus, $CM$ and $MF$ are rational (straight-lines which are) commensurable in square only. $CF$ is thus an
apotome [Prop.~10.73]. So, I say that (it is)
also a second (apotome).

\epsfysize=1.6in
\centerline{\epsffile{Book10/fig097e.eps}}

For let $FM$ be cut in half at $N$. And let $NO$ be drawn
through (point) $N$, parallel to $CD$. Thus, $FO$ and
$NL$ are each equal to the (rectangle contained) by $AG$ and $GB$.
And since
the (rectangle contained) by $AG$ and $GB$ is the mean
proportional to the squares on $AG$ and $GB$ [Prop.~10.21~lem.], and the (square) on $AG$
is equal to $CH$, and the (rectangle contained) by $AG$ and
$GB$ to $NL$, and the (square) on $BG$ to $KL$,  $NL$ is thus also the
mean proportional to $CH$ and $KL$. Thus, as $CH$ is to $NL$, so
$NL$ (is) to $KL$ [Prop.~5.11]. But, as $CH$ (is) to $NL$, so $CK$ is to
$NM$, and as $NL$ (is) to $KL$, so $NM$ is to $MK$
[Prop.~6.1]. Thus, as $CK$ (is) to $NM$, so $NM$
is to $KM$ [Prop.~5.11]. The (rectangle contained) by  $CK$ and $KM$ is thus
equal to the (square) on $NM$ [Prop.~6.17]---that is to say, to the fourth part of the
(square) on $FM$ [and since the (square) on $AG$ is commensurable
with the (square) on $BG$, $CH$ is also commensurable with $KL$---that is to say, $CK$ with $KM$].  Therefore, since $CM$ and $MF$ are two
unequal straight-lines, and the (rectangle contained) by $CK$ and
$KM$, equal to the fourth part
of the (square) on $MF$,  has been applied to the greater $CM$,
falling short by a square figure, and divides it into commensurable
(parts), the square on $CM$ is thus greater than (the square on)
$MF$ by the (square) on (some straight-line) commensurable
in length with ($CM$) [Prop.~10.17]. The
attachment $FM$ is also commensurable in length with the (previously)
laid down rational (straight-line) $CD$. $CF$ is thus a second apotome [Def.~10.16].

Thus, the (square) on a first apotome of
a medial (straight-line), applied to a rational (straight-line), produces
 a second apotome as breadth. (Which is) the very thing it was required to
show.}
\end{Parallel}

%%%%
%10.99
%%%%
\pdfbookmark[1]{Proposition 10.99}{pdf10.99}
\begin{Parallel}{}{}
\ParallelLText{
\begin{center}
{\large \ggn{99}.}
\end{center}\vspace*{-7pt}

{\greekfont Τὸ ἀπὸ μέσης ἀποτομῆς δευτέρας παρὰ ῥητὴν παραβαλλόμενον
πλάτος ποιεῖ ἀποτομὴν τρίτην.}

{\greekfont Ἔστω μέσης ἀποτομὴ δευτέρα ἡ ΑΒ, ῥητὴ δὲ ἡ ΓΔ, καὶ τῷ
ἀπὸ τῆς ΑΒ ἴσον παρὰ τὴν ΓΔ παραβεβλήσθω τὸ ΓΕ πλάτος
ποιοῦν τὴν ΓΖ; λέγω, ὅτι ἡ ΓΖ ἀποτομή ἐστι τρίτη.}\\~\\

\epsfysize=1.6in
\centerline{\epsffile{Book10/fig097g.eps}}

{\greekfont Ἔστω γὰρ τῇ ΑΒ προσαρμόζουσα ἡ ΒΗ; αἱ ἄρα ΑΗ, ΗΒ μέσαι
εἰσὶ δυνάμει μόνον σύμμετροι μέσον περιέθουσαι. καὶ τῷ μὲν
ἀπὸ τῆς ΑΗ ἴσον παρὰ τὴν ΓΔ παραβεβλήσθω τὸ ΓΘ πλάτος ποιοῦν
τὴν ΓΚ, τῷ δὲ ἀπὸ τῆς ΒΗ ἴσον παρὰ τὴν ΚΘ παραβεβλήσθω
τὸ ΚΛ πλάτος ποιοῦν τὴν ΚΜ; ὅλον ἄρα τὸ ΓΛ ἴσον ἐστὶ
τοῖς ἀπὸ τῶν ΑΗ, ΗΒ [καί ἐστι μέσα τὰ ἀπὸ τῶν ΑΗ, ΗΒ];
μέσον ἄρα καὶ τὸ ΓΛ. καὶ παρὰ ῥητὴν τὴν ΓΔ παραβέβληται
πλάτος ποιοῦν τὴν ΓΜ; ῥητὴ ἄρα ἐστὶν ἡ ΓΜ καὶ ἀσύμμετρος
τῇ ΓΔ μήκει. καὶ ἐπεὶ ὅλον τὸ ΓΛ ἴσον ἐστὶ τοῖς ἀπὸ
τῶν ΑΗ, ΗΒ, ὧν τὸ ΓΕ ἴσον ἐστὶ τῷ ἀπὸ τῆς ΑΒ,
λοιπὸν ἄρα τὸ ΛΖ ἴσον ἐστὶ τῷ δὶς ὑπὸ τῶν ΑΗ, ΗΒ. τετμήσθω
οὖν ἡ ΖΜ δίθα κατὰ τὸ Ν σημεῖον, καὶ τῇ ΓΔ παράλληλος
ἤθθω ἡ ΝΧ; ἑκάτερον ἄρα τῶν ΖΧ, ΝΛ ἴσον ἐστὶ τῷ
ὑπὸ τῶν ΑΗ, ΗΒ. μέσον δὲ τὸ ὑπὸ τῶν ΑΗ, ΗΒ; μέσον
ἄρα ἐστὶ καὶ τὸ ΖΛ. καὶ παρὰ ῥητὴν τὴν ΕΖ παράκειται πλάτος
ποιοῦν τὴν ΖΜ; ῥητὴ ἄρα καὶ ἡ ΖΜ καὶ ἀσύμμετρος τῇ ΓΔ
μήκει. καὶ ἐπεὶ αἱ ΑΗ, ΗΒ δυνάμει μόνον εἰσὶ σύμμετροι,
ἀσύμμετρος ἄρα [ἐστὶ] μήκει ἡ ΑΗ τῇ ΗΒ; ἀσύμμετρον
ἄρα ἐστὶ καὶ τὸ ἀπὸ τῆς ΑΗ τῷ ὑπὸ τῶν ΑΗ, ΗΒ.
ἀλλὰ τῷ μὲν ἀπὸ τῆς ΑΗ σύμμετρά ἐστι τὰ ἀπὸ τῶν
ΑΗ, ΗΒ, τῷ δὲ ὑπὸ τῶν ΑΗ, ΗΒ τὸ δὶς ὑπὸ τῶν
ΑΗ, ΗΒ; ἀσύμμετρα ἄρα ἐστὶ τὰ ἀπὸ τῶν
ΑΗ, ΗΒ τῷ δὶς ὑπὸ τῶν ΑΗ, ΗΒ. ἀλλὰ τοῖς μὲν ἀπὸ
τῶν ΑΗ, ΗΒ ἴσον ἐστὶ τὸ ΓΛ, τῷ δὲ δὶς ὑπὸ τῶν ΑΗ, ΗΒ
ἴσον ἐστὶ τὸ ΖΛ; ἀσύμμετρον ἄρα ἐστὶ τὸ ΓΛ τῷ ΖΛ.
ὡς δὲ τὸ ΓΛ πρὸς τὸ ΖΛ, οὕτως ἐστὶν ἡ ΓΜ πρὸς τὴν ΖΜ;
ἀσύμμετρος ἄρα ἐστὶν ἡ ΓΜ τῇ ΖΜ μήκει. καί εἰσιν
ἀμφότεραι ῥηταί; αἱ ἄρα ΓΜ, ΜΖ ῥηταί εἰσι δυνάμει μόνον
σύμμετροι; ἀποτομὴ ἄρα ἐστὶν ἡ ΓΖ. λέγω δή, ὅτι καὶ τρίτη.}

{\greekfont Ἐπεὶ γὰρ σύμμετρόν ἐστι τὸ ἀπὸ τῆς ΑΗ τῷ ἀπὸ τῆς ΗΒ, σύμμετρον ἄρα καὶ τὸ ΓΘ τῷ ΚΛ; ὥστε καὶ ἡ ΓΚ τῇ ΚΜ. καὶ
ἐπεὶ τῶν ἀπὸ τῶν ΑΗ, ΗΒ μέσον ἀνάλογόν ἐστι τὸ ὑπὸ
τῶν ΑΗ, ΗΒ, καί ἐστι τῷ μὲν ἀπὸ τῆς ΑΗ ἴσον τὸ ΓΘ, τῷ
δὲ ἀπὸ τῆς ΗΒ ἴσον τὸ ΚΛ,  τῷ δὲ ὑπὸ τῶν ΑΗ, ΗΒ ἴσον
τὸ ΝΛ, καὶ τῶν ΓΘ, ΚΛ ἄρα μέσον ἀνάλογόν ἐστι τὸ ΝΛ; ἔστιν
ἄρα ὡς τὸ ΓΘ πρὸς τὸ ΝΛ, οὕτως τὸ ΝΛ πρὸς τὸ ΚΛ. ἀλλ
ὡς μὲν τὸ ΓΘ πρὸς τὸ ΝΛ, οὕτως ἐστὶν ἡ  ΓΚ πρὸς τὴν ΝΜ,
ὡς δὲ τὸ ΝΛ πρὸς τὸ ΚΛ, οὕτως ἐστὶν ἡ ΝΜ πρὸς τὴν ΚΜ; ὡς
ἄρα ἡ ΓΚ πρὸς τὴν ΜΝ, οὕτως ἐστὶν ἡ ΜΝ πρὸς τὴν
ΚΜ; τὸ ἄρα ὑπὸ τῶν ΓΚ, ΚΜ ἴσον ἐστὶ τῷ [ἀπὸ
τῆς ΜΝ, τουτέστι τῷ] τετάρτῳ μέρει τοῦ ἀπὸ τῆς ΖΜ. ἐπεὶ
οὖν δύο εὐθεῖαι ἄνισοί εἰσιν αἱ ΓΜ, ΜΖ, καὶ τῷ τετάρτῳ
μέρει τοῦ ἀπὸ τῆς ΖΜ ἴσον παρὰ τὴν ΓΜ παραβέβληται ἐλλεῖπον
εἴδει τετραγώνῳ καὶ εἰς σύμμετρα α ὐτὴν διαιρεῖ, ἡ ΓΜ ἄρα
τῆς ΜΖ μεῖζον δύναται τῷ ἀπὸ συμμέτρου ἑαυτῇ. καὶ
οὐδετέρα τῶν ΓΜ, ΜΖ σύμμετρός ἐστι μήκει τῇ ἐκκειμένῃ
ῥητῇ τῇ ΓΔ; ἡ ἄρα ΓΖ ἀποτομή ἐστι τρίτη.}

{\greekfont Τὸ ἄρα ἀπὸ μέσης ἀποτομῆς δευτέρας παρὰ ῥητὴν παραβαλλόμενον πλάτος ποιεῖ ἀποτομὴν τρίτην; ὅπερ ἔδει δεῖχαι.}}

\ParallelRText{
\begin{center}
{\large Proposition 99}
\end{center}

The (square) on a second apotome of a medial
(straight-line), applied to a rational (straight-line), produces a
third apotome  as breadth.

Let $AB$ be the second apotome of a medial (straight-line), and
$CD$ a rational (straight-line). And let $CE$, equal to the (square) on  $AB$,
be applied to $CD$, producing $CF$ as breadth. I say that $CF$
is a third apotome.

\epsfysize=1.6in
\centerline{\epsffile{Book10/fig097e.eps}}

For let $BG$ be an attachment to $AB$. Thus, $AG$ and $GB$ are medial
(straight-lines which are) commensurable in square only, containing
a medial (area) [Prop.~10.75]. And let $CH$,
equal to the (square) on $AG$, be applied to $CD$, producing
$CK$ as breadth. And let $KL$, equal to the (square) on $BG$, 
be applied to $KH$, producing $KM$ as breadth. Thus, the whole
of $CL$ is equal to the (sum of the squares) on $AG$ and $GB$ [and
the (sum of the squares) on $AG$ and $GB$ is medial]. $CL$ (is) thus
also medial [Props.~10.15, 10.23~corr.]. And it has been applied to the rational (straight-line)
$CD$, producing $CM$ as breadth. Thus, $CM$ is rational, and
incommensurable in length with $CD$ [Prop.~10.22]. And since the whole of $CL$ is equal
to the (sum of the squares) on $AG$ and $GB$, of which $CE$ is
equal to the (square) on $AB$, the remainder $LF$ is thus equal to
twice the (rectangle contained) by $AG$ and $GB$ [Prop.~2.7]. Therefore, let $FM$ be
cut in half at point $N$. And let $NO$ be drawn parallel to $CD$.
Thus, $FO$ and $NL$ are each equal to the (rectangle contained) by
$AG$ and $GB$. And the (rectangle contained) by $AG$ and $GB$ (is)
medial. Thus, $FL$ is also medial. And it is applied to the
rational (straight-line) $EF$, producing $FM$ as breadth.
$FM$ is thus rational, and incommensurable in length with $CD$ [Prop.~10.22]. And since $AG$ and $GB$
are commensurable in square only, $AG$ [is] thus incommensurable
in length with $GB$. Thus, the (square) on $AG$ is also incommensurable
with the (rectangle contained) by $AG$ and $GB$
[Props.~6.1, 10.11]. But,  the (sum of the squares) on $AG$
and $GB$ is commensurable with the (square) on $AG$, and twice the
(rectangle contained) by $AG$ and $GB$ with the (rectangle contained)
by $AG$ and $GB$. The (sum of the squares) on $AG$ and $GB$
is thus incommensurable with twice the (rectangle contained)
by $AG$ and $GB$ [Prop.~10.13]. But,
$CL$ is equal to the (sum of the squares) on $AG$ and $GB$, 
and $FL$ is equal to twice the (rectangle contained) by $AG$ and $GB$. Thus,
$CL$ is incommensurable with $FL$. And as $CL$ (is) to $FL$, so
$CM$ is to $FM$ [Prop.~6.1].  $CM$ is thus
incommensurable in length with $FM$ [Prop.~10.11]. And they are both rational (straight-lines). Thus, $CM$ and $MF$ are rational (straight-lines which are)
commensurable in square only. $CF$ is thus an apotome [Prop.~10.73]. So, I say that (it is) also a
third (apotome).

For since the (square) on $AG$ is commensurable with the (square) on 
$GB$, $CH$ (is) thus also commensurable with $KL$. Hence,
$CK$ (is) also  (commensurable in length) with $KM$ [Props.~6.1, 10.11]. 
And since the (rectangle contained) by $AG$ and $GB$ is the mean
proportional to the (squares) on $AG$ and $GB$ [Prop.~10.21~lem.], and $CH$
is equal to the (square) on $AG$, and $KL$ equal to the (square) on  $GB$,
and $NL$ equal to the (rectangle contained) by $AG$ and $GB$, $NL$ is
thus also the mean proportional to $CH$ and $KL$. Thus, as $CH$ is to
$NL$, so $NL$ (is) to $KL$. But, as $CH$ (is) to $NL$, so $CK$
is to $NM$, and as $NL$ (is) to $KL$, so $NM$ (is) to $KM$
[Prop.~6.1]. Thus, as $CK$ (is) to $MN$, so $MN$
is to $KM$ [Prop.~5.11]. Thus, the (rectangle contained) by $CK$ and $KM$
is equal to the [(square) on $MN$---that is to say, to the] fourth part of the
(square) on $FM$ [Prop.~6.17]. Therefore,
since $CM$ and $MF$ are two unequal straight-lines, and
(some area), equal to the fourth part of the (square) on  $FM$, has been
applied to $CM$, falling short by a square figure, and divides it into
commensurable (parts),  the square on $CM$ is thus greater than (the square on) $MF$ by the (square) on (some straight-line) commensurable
(in length) with ($CM$) [Prop.~10.17]. 
And neither of $CM$ and $MF$ is commensurable in length with the
(previously) laid down rational (straight-line) $CD$. $CF$ is thus a
third apotome [Def.~10.13].

Thus, the (square) on a second apotome of a medial
(straight-line), applied to a rational (straight-line), produces  a
third apotome as breadth. (Which is) the very thing it was required to show.}
\end{Parallel}

%%%%%
%10.100
%%%%%
\pdfbookmark[1]{Proposition 10.100}{pdf10.100}
\begin{Parallel}{}{}
\ParallelLText{
\begin{center}
{\large \ggn{100}.}
\end{center}\vspace*{-7pt}

{\greekfont Τὸ ἀπὸ ἐλάσσονος παρὰ ῥητὴν παραβαλλόμενον πλάτος ποιεῖ
ἀποτομὴν τετάρτην.}

\epsfysize=1.6in
\centerline{\epsffile{Book10/fig097g.eps}}

{\greekfont Ἔστω ἐλάσσων ἡ ΑΒ, ῥητὴ δὲ ἡ ΓΔ, καὶ τῷ ἀπὸ
τῆς ΑΒ ἴσον παρὰ ῥητὴν τὴν ΓΔ παραβεβλήσθω τὸ ΓΕ
πλάτος ποιοῦν τὴν ΓΖ; λέγω, ὅτι ἡ ΓΖ ἀποτομή ἐστι
τετάρτη.}

{\greekfont Ἔστω γὰρ τῇ ΑΒ προσαρμόζουσα ἡ ΒΗ; αἱ ἄρα ΑΗ, ΗΒ δυνάμει
εἰσὶν ἀσύμμετροι ποιοῦσαι τὸ μὲν συγκείμενον ἐκ τῶν ἀπὸ
τῶν ΑΗ, ΗΒ τετραγώνων ῥητόν, τὸ δὲ δὶς ὑπὸ τῶν ΑΗ, ΗΒ
μέσον. καὶ τῷ μὲν ἀπὸ τῆς ΑΗ ἴσον παρὰ τὴν ΓΔ παραβεβλήσθω
τὸ ΓΘ πλάτος ποιοῦν τὴν ΓΚ, τῷ δὲ ἀπὸ τῆς ΒΗ
ἴσον τὸ ΚΛ πλάτος ποιοῦν τὴν ΚΜ; ὅλον ἄρα τὸ ΓΛ ἴσον
ἐστὶ τοῖς ἀπὸ τῶν ΑΗ, ΗΒ. καί ἐστι τὸ συγκείμενον ἐκ τῶν
ἀπὸ τῶν ΑΗ, ΗΒ ῥητόν; ῥητὸν ἄρα ἐστὶ καὶ τὸ ΓΛ.
καὶ παρὰ ῥητὴν τὴν ΓΔ παράκειται πλάτος ποιοῦν τὴν ΓΜ; ῥητὴ ἄρα
καὶ ἡ ΓΜ καὶ σύμμετρος τῇ ΓΔ μήκει. καὶ ἐπεὶ ὅλον
τὸ ΓΛ ἴσον ἐστὶ τοῖς ἀπὸ τῶν ΑΗ, ΗΒ, ὧν τὸ ΓΕ ἴσον
ἐστὶ τῷ ἀπὸ τῆς ΑΒ, λοιπὸν ἄρα τὸ ΖΛ ἴσον ἐστὶ τῷ δὶς
ὑπὸ τῶν ΑΗ, ΗΒ. τετμήσθω οὖν ἡ ΖΜ δίθα κατὰ τὸ Ν σημεῖον,
καὶ ἤθθω δὶα τοῦ Ν ὁποτέρᾳ τῶν ΓΔ, ΜΛ παράλληλος ἡ ΝΧ;
ἑκάτερον ἄρα τῶν ΖΧ, ΝΛ ἴσον ἐστὶ τῷ ὑπὸ τῶν ΑΗ, ΗΒ. καὶ
ἐπεὶ τὸ δὶς ὑπὸ τῶν ΑΗ, ΗΒ μέσον ἐστὶ καί ἐστιν ἴσον τῷ
ΖΛ, καὶ τὸ ΖΛ ἄρα μέσον ἐστίν. καὶ παρὰ ῥητὴν τὴν ΖΕ παράκειται
πλάτος ποιοῦν τὴν ΖΜ; ῥητὴ ἄρα ἐστὶν ἡ ΖΜ καὶ ἀσύμμετρος
τῇ ΓΔ μήκει. καὶ ἐπεὶ τὸ μὲν συγκείμενον ἐκ τῶν ἀπὸ
τῶν ΑΗ, ΗΒ ῥητόν ἐστιν, τὸ δὲ δὶς ὑπὸ τῶν ΑΗ, ΗΒ μέσον,
ἀσύμμετρα [ἄρα] ἐστὶ τὰ ἀπὸ τῶν ΑΗ, ΗΒ 
 τῷ δὶς ὑπὸ τῶν ΑΗ, ΗΒ.
ἴσον δέ [ἐστι] τὸ ΓΛ τοῖς ἀπὸ τῶν ΑΗ, ΗΒ, τῷ δὲ δὶς ὑπὸ
τῶν ΑΗ, ΗΒ ἴσον τὸ ΖΛ; ἀσύμμετρον ἄρα [ἐστὶ] τὸ ΓΛ τῷ ΖΛ.
ὡς δὲ τὸ ΓΛ πρὸς τὸ ΖΛ, οὕτως ἐστὶν
  ἡ ΓΜ πρὸς τὴν ΜΖ; ἀσύμμετρος ἄρα
  ἐστὶν ἡ ΓΜ τῇ ΜΖ μήκει.
καί εἰσιν ἀμφότεραι ῥηταί; αἱ ἄρα ΓΜ, ΜΖ ῥηταί
εἰσι δυνάμει μόνον σύμμετροι; ἀποτομὴ ἄρα ἐστὶν ἡ ΓΖ. λέγω
[δή], ὅτι καὶ τετάρτη.}

{\greekfont Ἐπεὶ γὰρ αἱ ΑΗ, ΗΒ δυνάμει εἰσὶν ἀσύμμετροι, ἀσύμμετ-ρον
ἄρα καὶ τὸ ἀπὸ τῆς ΑΗ τῷ ἀπὸ τῆς ΗΒ.  καί ἐστι τῷ μὲν
ἀπὸ τῆς ΑΗ ἴσον τὸ ΓΘ, τῷ δὲ ἀπὸ τῆς ΗΒ ἴσον
τὸ ΚΛ; ἀσύμμετρον ἄρα ἐστὶ τὸ ΓΘ τῷ ΚΛ. ὡς δὲ τὸ ΓΘ πρὸς
τὸ ΚΛ, οὕτως ἐστὶν ἡ ΓΚ πρὸς τὴν ΚΜ; ἀσύμμετρος ἄρα ἐστὶν
ἡ ΓΚ τῇ ΚΜ μήκει. καὶ ἐπεὶ τῶν ἀπὸ τῶν ΑΗ, ΗΒ μέσον ἀνάλογόν ἐστι τὸ ὑπὸ τῶν ΑΗ, ΗΒ, καί ἐστιν ἴσον τὸ
μὲν ἀπὸ τῆς ΑΗ τῷ ΓΘ, τὸ δὲ ἀπὸ τῆς ΗΒ τῷ ΚΛ, τὸ δὲ ὑπὸ
τῶν ΑΗ, ΗΒ τῷ ΝΛ, τῶν ἄρα ΓΘ, ΚΛ μέσον ἀνάλογόν ἐστι
τὸ ΝΛ; ἔστιν ἄρα ὡς τὸ ΓΘ πρὸς τὸ ΝΛ, οὕτως τὸ ΝΛ πρὸς
τὸ ΚΛ. ἀλλ ὡς μὲν τὸ ΓΘ πρὸς τὸ ΝΛ, οὕτως ἐστίν ἡ ΓΚ
πρὸς τὴν ΝΜ, ὡς δὲ τὸ ΝΛ πρὸς τὸ ΚΛ, οὕτως ἐστὶν ἡ ΝΜ
πρὸς τὴν ΚΜ; ὡς ἄρα ἡ ΓΚ πρὸς τὴν ΜΝ, οὕτως ἐστὶν
ἡ ΜΝ πρὸς τὴν ΚΜ; τὸ ἄρα ὑπὸ τῶν ΓΚ, ΚΜ ἴσον ἐστὶ
τῷ ἀπὸ τῆς ΜΝ, τουτέστι τῷ τετάρτῳ μέρει τοῦ ἀπὸ τῆς
ΖΜ. ἐπεὶ οὖν δύο εὐθεῖαι ἄνισοί εἰσιν αἱ ΓΜ, ΜΖ, καὶ
τῷ τετράρτῳ μέρει τοῦ ἀπὸ τῆς ΜΖ ἴσον παρὰ τὴν ΓΜ παραβέβληται
ἐλλεῖπον εἴδει τετραγώνῳ τὸ ὑπὸ τῶν ΓΚ, ΚΜ καὶ
εἰς ἀσύμμετρα α ὐτὴν διαιρεῖ, ἡ ἄρα ΓΜ τῆς ΜΖ μεῖζον
δύναται τῷ ἀπὸ ἀσυμμέτρου ἑαυτῇ. καί ἐστιν ὅλη
ἡ ΓΜ σύμμετρος μήκει τῇ ἐκκειμένῃ ῥητῇ τῇ ΓΔ;
ἡ ἄρα ΓΖ ἀποτομή ἐστι τετάρτη.}

{\greekfont Τὸ ἄρα ἀπὸ ἐλάσσονος καὶ τὰ ἑχῆς.}}

\ParallelRText{
\begin{center}
{\large Proposition 100}
\end{center}

The (square) on a minor
(straight-line), applied to a rational (straight-line), produces  a
fourth apotome as breadth.

\epsfysize=1.6in
\centerline{\epsffile{Book10/fig097e.eps}}

Let $AB$ be a minor (straight-line), and $CD$ a rational (straight-line).
And let $CE$, equal to the (square) on $AB$, be applied to the rational (straight-line) $CD$, producing $CF$ as breadth. I say that
$CF$ is a fourth apotome.

For let $BG$ be an attachment to $AB$. Thus, $AG$ and $GB$ are incommensurable in square, making the sum of the squares on $AG$ and $GB$ rational, and twice the (rectangle contained) by $AG$ and $GB$ medial [Prop.~10.76]. And let $CH$,
equal to the (square) on $AG$, be applied to $CD$, producing
$CK$ as breadth, and $KL$, equal to the (square) on $BG$, producing
$KM$ as breadth. Thus, the whole of $CL$ is equal to the (sum of the
squares) on $AG$ and $GB$. And the sum of the (squares) on $AG$ and
$GB$ is rational. $CL$ is thus also rational. And it is applied to the
rational (straight-line) $CD$, producing $CM$ as breadth. Thus, $CM$
(is) also rational, and commensurable in length with $CD$ [Prop.~10.20]. And since the whole of $CL$
is equal to the (sum of the squares) on $AG$ and $GB$, of which
$CE$ is equal to the (square) on $AB$, the remainder $FL$ is thus
equal to twice the (rectangle contained) by $AG$ and $GB$ [Prop.~2.7]. Therefore, let $FM$ be cut in
half at point $N$. And let $NO$ be drawn through $N$, parallel
to either of $CD$ or $ML$. Thus, $FO$ and $NL$ are each
equal to the (rectangle contained) by $AG$ and $GB$. And since
twice the (rectangle contained) by $AG$ and $GB$ is medial, and is equal
to $FL$, $FL$ is thus also medial. And it is applied to the rational (straight-line) $FE$, producing $FM$ as breadth. Thus, $FM$ is  rational, and incommensurable in length with $CD$ [Prop.~10.22]. And since the sum of the (squares)
on $AG$ and $GB$ is rational, and twice the (rectangle contained) by $AG$
and $GB$ medial, the (sum of the squares) on $AG$ and $GB$
is [thus] incommensurable with twice the (rectangle contained)
by $AG$ and $GB$. And $CL$ (is) equal to the (sum of the squares)
on $AG$ and $GB$, and $FL$ equal to twice the (rectangle contained) by $AG$ and $GB$. $CL$ [is] thus incommensurable with $FL$. And as
$CL$ (is) to $FL$, so $CM$ is to $MF$ [Prop.~6.1]. 
$CM$ is thus incommensurable in length with $MF$ [Prop.~10.11]. And both are rational (straight-lines). 
Thus, $CM$ and $MF$ are rational (straight-lines which are) commensurable
in square only. $CF$ is thus an apotome [Prop.~10.73]. [So], I say that (it is) also
a fourth (apotome).

For since $AG$ and $GB$ are incommensurable in square, the (square)
on $AG$ (is) thus also incommensurable with the (square) on $GB$.
And $CH$ is equal to the (square) on $AG$, and $KL$ equal to the (square) on $GB$. Thus, $CH$ is incommensurable with $KL$. And as $CH$ (is) to
$KL$, so $CK$ is to $KM$ [Prop.~6.1]. 
$CK$ is thus incommensurable in length with $KM$ [Prop.~10.11]. And since the (rectangle contained)
by $AG$ and $GB$ is the mean proportional to the (squares) on $AG$
and $GB$ [Prop.~10.21~lem.], and the
(square) on $AG$ is equal to $CH$, and the (square) on $GB$ to $KL$,
and the (rectangle contained) by $AG$ and $GB$ to $NL$, $NL$
is thus the mean proportional to $CH$ and $KL$. Thus, as $CH$
is to $NL$, so $NL$ (is) to $KL$. But, as $CH$ (is) to $NL$, so
$CK$ is to $NM$, and as $NL$ (is) to $KL$, so $NM$ is to $KM$ [Prop.~6.1]. Thus, as $CK$ (is) to $MN$, so
$MN$ is to $KM$ [Prop.~5.11]. The (rectangle
contained) by $CK$ and $KM$ is thus equal to the (square) on $MN$---that is to say, to the fourth part of the (square) on $FM$ [Prop.~6.17]. Therefore, since $CM$ and $MF$ are two unequal straight-lines, and the (rectangle contained) by $CK$ and
$KM$, equal to the fourth part of the (square) on $MF$, has been applied
to $CM$, falling short by a square figure, and divides it into incommensurable (parts), the square on $CM$ is thus greater than (the
square on) $MF$ by the (square) on (some straight-line) incommensurable
(in length) with ($CM$) [Prop.~10.18]. 
And the whole of $CM$ is commensurable in length with the (previously)
laid down rational (straight-line) $CD$. Thus, $CF$ is a fourth apotome
[Def.~10.14].

Thus, the (square) on a minor, and so on \ldots}
\end{Parallel}

%%%%%
%10.101
%%%%%
\pdfbookmark[1]{Proposition 10.101}{pdf10.101}
\begin{Parallel}{}{}
\ParallelLText{
\begin{center}
{\large \ggn{101}.}
\end{center}\vspace*{-7pt}

{\greekfont Τὸ ἀπὸ τῆς μετὰ ῥητοῦ μέσον τὸ ὅλον ποιούσης
παρὰ ῥητὴν παραβαλλόμενον πλάτος ποιεῖ ἀποτομὴν πέμπτην.}\\

\epsfysize=1.6in
\centerline{\epsffile{Book10/fig097g.eps}}

{\greekfont Ἔστω ἡ μετὰ ῥητοῦ μέσον τὸ ὅλον ποιοῦσα ἡ ΑΒ,
ῥητὴ δὲ ἡ ΓΔ, καὶ τῷ ἀπὸ τῆς ΑΒ ἴσον παρὰ τὴν ΓΔ
παραβεβλήσθω τὸ ΓΕ πλάτος ποιοῦν τὴν ΓΖ; λέγω, ὅτι ἡ ΓΖ
ἀποτομή ἐστι πέμπτη.}

{\greekfont Ἔστω γὰρ τῇ ΑΒ προσαρμόζουσα ἡ ΒΗ; αἱ ἄρα ΑΗ, ΗΒ εὐθεῖαι
δυνάμει εἰσὶν ἀσύμμετροι ποιοῦσαι τὸ μὲν συγκείμενον
ἐκ τῶν ἀπ α ὐτῶν τετραγώνων μέσον, τὸ δὲ δὶς ὑπ
α ὐτῶν ῥητόν, καὶ τῷ μὲν ἀπὸ τῆς ΑΗ ἴσον παρὰ τὴν ΓΔ
παραβεβλήσθω τὸ ΓΘ, τῷ δὲ ἀπὸ τῆς ΗΒ ἵσον τὸ ΚΛ;
ὅλον ἄρα τὸ ΓΛ ἴσον ἐστὶ τοῖς ἀπὸ τῶν ΑΗ, ΗΒ.
τὸ δὲ συγκείμενον ἐκ
τῶν ἀπὸ τῶν ΑΗ, ΗΒ ἅμα μέσον ἐστίν; μέσον ἄρα ἐστὶ
τὸ ΓΛ. καὶ παρὰ ῥητὴν τὴν ΓΔ παράκειται πλάτος ποιοῦν
τὴν ΓΜ; ῥητὴ ἄρα ἐστὶν ἡ ΓΜ καὶ ἀσύμμετρος τῇ ΓΔ. καὶ
ἐπεὶ ὅλον τὸ ΓΛ ἴσον ἐστὶ τοῖς ἀπὸ τῶν ΑΗ, ΗΒ, ὧν
τὸ ΓΕ ἴσον ἐστὶ τῷ ἀπὸ τῆς ΑΒ, λοιπὸν ἄρα τὸ ΖΛ ἴσον
ἐστὶ τῷ δὶς ὑπὸ τῶν ΑΗ, ΗΒ. τετμήσθω οὖν ἡ ΖΜ δίθα
κατὰ τὸ Ν, καὶ ἤθθω διὰ τοῦ Ν ὁποτέρᾳ τῶν ΓΔ, ΜΛ παράλληλος
ἡ ΝΧ; ἑκάτερον ἄρα τῶν ΖΧ, ΝΛ ἴσον ἐστὶ τῷ ὑπὸ τῶν
ΑΗ, ΗΒ, καὶ ἐπεὶ τὸ δὶς ὑπὸ τῶν ΑΗ, ΗΒ ῥητόν ἐστι καί
[ἐστιν] ἴσον τῷ ΖΛ, ῥητὸν ἄρα ἐστὶ τὸ ΖΛ. καὶ παρὰ ῥητὴν
τὴν ΕΖ παράκειται πλάτος ποιοῦν τὴν ΖΜ; ῥητὴ ἄρα ἐστὶν ἡ ΖΜ
καὶ σύμμετρος τῇ ΓΔ μήκει. καὶ ἐπεὶ τὸ μὲν ΓΛ μέσον
ἐστίν, τὸ δὲ ΖΛ ῥητόν, ἀσύμμετρον ἄρα ἐστὶ τὸ ΓΛ τῷ ΖΛ. ὡς δὲ τὸ ΓΛ πρὸς τὸ ΖΛ, οὕτως ἡ ΓΜ πρὸς τὴν ΜΖ; ἀσύμμετρος
ἄρα ἐστὶν ἡ ΓΜ τῇ ΜΖ μήκει. καί εἰσιν ἀμφότεραι ῥηταί;
αἱ ἄρα ΓΜ, ΜΖ ῥηταί εἰσι δυνάμει μόνον σύμμετροι; ἀποτομὴ
ἄρα ἐστὶν ἡ ΓΖ. λὲγω δή, ὅτι καὶ πέμπτη.}

{\greekfont Ὁμοίως γὰρ δείχομεν, ὅτι τὸ ὑπὸ τῶν ΓΚΜ ἴσον
ἐστὶ τῷ ἀπὸ τῆς ΝΜ, τουτέστι τῷ τετάρτῳ μέρει τοῦ
ἀπὸ τῆς ΖΜ. καὶ ἐπεὶ ἀσύμμετρόν ἐστι τὸ ἀπὸ τῆς
ΑΗ τῷ ἀπὸ τῆς ΗΒ, ἴσον δὲ τὸ μὲν ἀπὸ τῆς ΑΗ τῷ
ΓΘ, τὸ δὲ ἀπὸ τῆς ΗΒ τῷ ΚΛ, ἀσύμμετρον ἄρα
τὸ ΓΘ τῷ ΚΛ. ὡς δὲ τὸ ΓΘ πρὸς τὸ ΚΛ, οὕτως
ἡ ΓΚ πρὸς τὴν ΚΜ; ἀσύμμετρος ἄρα ἡ ΓΚ τῇ ΚΜ μήκει.
ἐπεὶ οὖν δύο εὐθεῖαι ἄνισοί εἰσιν αἱ ΓΜ, ΜΖ, καὶ τῷ
τετάρτῳ μέρει τοῦ ἀπὸ τῆς ΖΜ ἴσον παρὰ τὴν ΓΜ παραβέβληται
ἐλλεῖπον εἴδει τετραγώνῳ καὶ εἰς ἀσύμμετρα α ὐτὴν διαιρεῖ,
ἡ ἄρα ΓΜ τῆς ΜΖ μεῖζον δύναται τῷ ἀπὸ ἀσύμμέτρου
ἑαυτῇ. καί ἐστιν ἡ προσαρμόζουσα ἡ ΖΜ σύμμετρος τῇ
ἐκκειμένῃ ῥητῇ τῇ ΓΔ; ἡ ἄρα ΓΖ ἀποτομή ἐστι
πέμπτη; ὅπερ ἔδει δεῖχαι.}}

\ParallelRText{
\begin{center}
{\large Proposition 101}
\end{center}\vspace*{-7pt}

The (square) on that (straight-line) which
with a rational (area) makes a medial whole, applied to a rational (straight-line), produces a fifth apotome as breadth.

\epsfysize=1.6in
\centerline{\epsffile{Book10/fig097e.eps}}

Let $AB$ be that (straight-line) which with a rational (area) makes a
medial whole, and $CD$ a rational (straight-line). And let $CE$, equal
to the (square) on $AB$, be applied to $CD$, producing $CF$
as breadth. I say that $CF$ is a fifth apotome.

Let $BG$ be an attachment to $AB$. Thus, the straight-lines $AG$
and $GB$ are incommensurable in square, making the sum of the
squares on them medial, and twice the (rectangle contained) by them
rational [Prop.~10.77]. And let $CH$, equal 
to the (square) on $AG$, be applied to $CD$, and $KL$, equal
to the (square) on $GB$. The whole of $CL$
is thus equal to the (sum of the squares) on $AG$ and $GB$. 
And the sum of the (squares) on $AG$ and $GB$ together is medial. Thus, $CL$
is medial. And it has been applied to the rational (straight-line) $CD$,
producing $CM$ as breadth. $CM$ is thus rational, and incommensurable
(in length) with $CD$ [Prop.~10.22]. And since the whole
of $CL$ is equal to the (sum of the squares) on $AG$ and $GB$, of
which $CE$ is equal to the (square) on $AB$, the remainder $FL$ is
thus equal to twice the (rectangle contained) by $AG$ and $GB$ [Prop.~2.7]. Therefore, let $FM$ be
cut in half at $N$. And let $NO$ be drawn through $N$, parallel
to either of $CD$ or $ML$. Thus, $FO$ and $NL$ are each equal
to the (rectangle contained) by $AG$ and $GB$. And since twice the
(rectangle contained) by $AG$ and $GB$ is rational, and [is]
equal to $FL$, $FL$ is thus rational. And it is applied to the rational
(straight-line) $EF$, producing $FM$ as breadth. Thus, $FM$ is rational,
and commensurable in length with $CD$ [Prop.~10.20]. And since $CL$ is medial, and
$FL$ rational, $CL$ is thus incommensurable with $FL$. 
And as $CL$ (is) to $FL$, so $CM$ (is) to $MF$ [Prop.~6.1]. $CM$ is thus incommensurable in length
with $MF$ [Prop.~10.11]. And both are
rational. Thus, $CM$ and $MF$ are rational (straight-lines which are)
commensurable in square only. $CF$ is thus an apotome [Prop.~10.73].  So, I say that (it is) also
a fifth (apotome).

For, similarly (to the previous propositions), we can show that the
(rectangle contained) by $CKM$ is equal to the (square) on $NM$---that is
to say, to the fourth part of the (square) on $FM$.  And since the (square)
on $AG$ is incommensurable with the (square) on $GB$, and  the
(square) on $AG$ (is)
equal to $CH$, and the (square) on $GB$ to $KL$, $CH$ (is) thus
incommensurable with $KL$. And as $CH$ (is) to $KL$, so
$CK$ (is) to $KM$ [Prop.~6.1]. 
Thus, $CK$ (is) incommensurable in length with $KM$ [Prop.~10.11]. Therefore, since $CM$ and $MF$
are two unequal straight-lines, and (some area), equal to the
fourth part of the (square) on $FM$, has been applied to $CM$, falling
short by a square figure, and divides it into incommensurable
(parts), the square on $CM$ is thus greater than (the square on) $MF$
by the (square) on (some straight-line) incommensurable (in length)
with ($CM$) [Prop.~10.18]. And the attachment
$FM$ is commensurable with the (previously) laid down
rational (straight-line) $CD$. Thus, $CF$ is a fifth apotome [Def.~10.15]. (Which is) the very thing it was required to show.}
\end{Parallel}

%%%%%
%10.102
%%%%%
\pdfbookmark[1]{Proposition 10.102}{pdf10.102}
\begin{Parallel}{}{}
\ParallelLText{
\begin{center}
{\large \ggn{102}.}
\end{center}\vspace*{-7pt}

{\greekfont Τὸ ἀπὸ τῆς μετὰ μέσου μέσον τὸ ὅλον ποιούσης παρὰ ῥητὴν
παραβαλλόμενον πλάτος ποιεῖ ἀποτομὴν ἕκτην.}\\

\epsfysize=1.6in
\centerline{\epsffile{Book10/fig097g.eps}}

{\greekfont Ἔστω ἡ μετὰ μέσου μέσον τὸ ὅλον ποιοῦσα ἡ ΑΒ, ῥητὴ δὲ ἡ
ΓΔ, καὶ τῷ ἀπὸ τῆς ΑΒ ἴσον παρὰ τὴν ΓΔ παραβεβλήσθω
τὸ ΓΕ πλάτος ποιοῦν τὴν ΓΖ; λέγω, ὅτι ἡ ΓΖ ἀποτομή
ἐστιν ἕκτη.}

{\greekfont Ἔστω γὰρ τῇ ΑΒ προσαρμόζουσα ἡ ΒΗ; αἱ ἄρα ΑΗ, ΗΒ
δυνάμει εἰσὶν ἀσύμμετροι ποιοῦσαι τό τε συγκείμενον
ἐκ τῶν ἀπ α ὐτῶν τετραγώνων μέσον καὶ τὸ δὶς ὑπὸ
τῶν ΑΗ, ΗΒ μέσον καὶ ἀσύμμετρον τὰ ἀπὸ τῶν ΑΗ, ΗΒ
τῷ δὶς ὑπὸ τῶν ΑΗ, ΗΒ. παραβεβλήσθω οὖν παρὰ τὴν ΓΔ
τῷ μὲν ἀπὸ τῆς ΑΗ ἴσον τὸ ΓΘ πλάτος ποιοῦν τὴν ΓΚ,
τῷ δὲ ἀπὸ τῆς ΒΗ τὸ ΚΛ; ὅλον ἄρα τὸ ΓΛ ἴσον
ἐστὶ τοῖς ἀπὸ τῶν ΑΗ, ΗΒ; μέσον ἄρα [ἐστὶ] καὶ τὸ ΓΛ.
καὶ παρὰ ῥητὴν τὴν ΓΔ παράκειται πλάτος ποιοῦν τὴν ΓΜ;
ῥητὴ ἄρα ἐστὶν ἡ ΓΜ καὶ ἀσύμμετρος τῇ ΓΔ μήκει.
ἐπεὶ οὖν τὸ ΓΛ ἴσον ἐστὶ τοῖς ἀπὸ τῶν
ΑΗ, ΗΒ, ὧν τὸ ΓΕ ἴσον τῷ ἀπὸ τῆς ΑΒ, λοιπὸν
ἄρα τὸ ΖΛ ἴσον ἐστὶ τῷ δὶς ὑπὸ τῶν ΑΗ, ΗΒ. καί
ἐστι τὸ δὶς ὑπὸ τῶν ΑΗ, ΗΒ μέσον; καὶ τὸ ΖΛ ἄρα
μέσον ἐστίν. καὶ παρὰ ῥητὴν τὴν ΖΕ παράκειται πλάτος
ποιοῦν τὴν ΖΜ; ῥητὴ ἄρα ἐστὶν ἡ ΖΜ καὶ ἀσύμμετρος
τῇ ΓΔ μήκει. καὶ ἐπεὶ τὰ ἀπὸ τῶν ΑΗ, ΗΒ ἀσύμμετρά
ἐστι τῷ δὶς ὑπὸ τῶν ΑΗ, ΗΒ, καί ἐστι τοῖς μὲν ἀπὸ
τῶν ΑΗ, ΗΒ ἴσον τὸ ΓΛ, τῷ δὲ δὶς ὑπὸ τῶν ΑΗ,
ΗΒ ἴσον τὸ ΖΛ, ἀσύμμετρος ἄρα [ἐστὶ] τὸ ΓΛ τῷ ΖΛ.
ὡς δὲ τὸ ΓΛ πρὸς τὸ 
ΖΛ, οὕτως ἐστὶν ἡ ΓΜ πρὸς τὴν
ΜΖ; ἀσύμμετρος ἄρα
ἐστὶν ἡ ΓΜ τῇ ΜΖ
μήκει. καί εἰσιν ἀμφότεραι ῥηταί. αἱ ΓΜ, ΜΖ ἄρα ῥηταί
εἰσι δυνάμει μόνον σύμμετροι; ἀποτομὴ ἄρα ἐστὶν ἡ ΓΖ.
λέγω δή, ὅτι καὶ ἕκτη.}

{\greekfont Ἐπεὶ γὰρ τὸ ΖΛ ἴσον ἐστὶ τῷ δὶς ὑπὸ τῶν ΑΗ, ΗΒ, τετμήσθω
δίθα ἡ ΖΜ κατὰ τὸ Ν, καὶ ἤθθω διὰ τοῦ Ν τῇ ΓΔ παράλληλος
ἡ ΝΧ; ἑκάτερον ἄρα τῶν ΖΧ, ΝΛ ἴσον ἐστὶ τῷ ὑπὸ
τῶν ΑΗ, ΗΒ. καὶ ἐπεὶ αἱ ΑΗ, ΗΒ δυνάμει εἰσὶν ἀσύμμετροι,
ἀσύμμετρον ἄρα ἐστὶ τὸ ἀπὸ τῆς ΑΗ τῷ ἀπὸ τῆς
ΗΒ. ἀλλὰ τῷ μὲν ἀπὸ τῆς ΑΗ ἴσον ἐστὶ τὸ ΓΘ, τῷ δὲ
ἀπὸ τῆς ΗΒ ἴσον ἐστὶ τὸ ΚΛ; ἀσύμμετρον ἄρα ἐστὶ
τὸ ΓΘ τῷ ΚΛ. ὡς δὲ τὸ ΓΘ πρὸς τὸ ΚΛ, οὕτως ἐστὶν
ἡ ΓΚ πρὸς τὴν ΚΜ; ἀσύμμετρος ἄρα ἐστὶν ἡ ΓΚ τῇ ΚΜ.
καὶ ἐπεὶ τῶν ἀπὸ τῶν ΑΗ, ΗΒ μέσον ἀνάλογόν ἐστι
τὸ ὑπὸ τῶν ΑΗ, ΗΒ, καὶ ἐστι τῷ μὲν ἀπὸ
τῆς ΑΗ ἴσον τὸ ΓΘ, τῷ δὲ ἀπὸ τῆς ΗΒ ἴσον
τὸ ΚΛ, τῷ δὲ ὑπὸ τῶν ΑΗ, ΗΒ ἴσον τὸ ΝΛ, καὶ τῶν ἄρα
ΓΘ, ΚΛ μέσον ἀνάλογόν ἐστι τὸ ΝΛ; ἔστιν ἄρα ὡς τὸ ΓΘ
πρὸς τὸ ΝΛ, οὕτως τὸ ΝΛ πρὸς τὸ ΚΛ. καὶ διὰ τὰ α ὐτὰ ἡ ΓΜ
τῆς ΜΖ μεῖζον δύναται τῷ ἀπὸ ἀσυμμέτρου ἑαυτῇ. καὶ
οὐδετέρα α ὐτῶν σύμμετρός
ἐστι τῇ ἐκκειμένῃ ῥητῇ τῇ ΓΔ; ἡ ΓΖ ἄρα ἀποτομή
ἐστιν ἕκτη; ὅπερ ἔδει δεῖχαι.}}

\ParallelRText{
\begin{center}
{\large Proposition 102}
\end{center}

The (square) on that (straight-line)
which with a medial (area) makes a medial whole, applied to
a rational (straight-line), produces a sixth apotome as breadth.

\epsfysize=1.6in
\centerline{\epsffile{Book10/fig097e.eps}}

Let $AB$ be that (straight-line) which with a medial (area) makes a
medial whole, and $CD$ a rational (straight-line). And let $CE$,
equal to the (square) on $AB$, be applied to $CD$, producing
$CF$ as breadth. I say that $CF$ is a sixth apotome.

For let $BG$ be an attachment to $AB$. Thus, $AG$ and $GB$
are incommensurable in square, making the sum of the squares on them
medial, and twice the (rectangle contained) by $AG$ and $GB$ medial,
and the (sum of the squares) on $AG$ and $GB$ incommensurable
with twice the (rectangle contained) by $AG$ and $GB$ [Prop.~10.78]. Therefore, let $CH$, equal to
the (square) on $AG$, be applied to $CD$, producing $CK$
as breadth, and $KL$, equal to the (square) on $BG$. Thus, the whole of
$CL$ is equal to the (sum of the squares) on $AG$ and $GB$. $CL$ [is]
thus also medial. And it is applied to the rational (straight-line) $CD$,
producing $CM$ as breadth. Thus,  $CM$ is rational, and incommensurable
in length with $CD$ [Prop.~10.22]. Therefore,
since $CL$ is equal to the (sum of the squares) on $AG$ and $GB$,
of which $CE$ (is) equal to the (square) on $AB$, the remainder $FL$
is thus equal to twice the (rectangle contained) by $AG$ and $GB$ [Prop.~2.7]. And twice the (rectangle contained)
by  $AG$ and $GB$ (is) medial. Thus, $FL$ is also medial.
And it is applied to the rational (straight-line) $FE$, producing $FM$
as breadth. $FM$ is thus rational, and incommensurable
in length with $CD$ [Prop.~10.22]. 
And since the (sum of the squares) on $AG$ and $GB$ is incommensurable
with twice the (rectangle contained) by $AG$ and $GB$, and $CL$ equal to the
(sum of the squares) on $AG$ and $GB$, and $FL$ equal to twice the
(rectangle contained) by $AG$ and $GB$, $CL$ [is] thus
incommensurable with $FL$. And as $CL$ (is) to $FL$, so $CM$ is
to $MF$ [Prop.~6.1]. Thus, $CM$ is incommensurable
in length with $MF$ [Prop.~10.11]. And they
are both rational. Thus, $CM$ and $MF$ are rational (straight-lines which are) commensurable in square only. $CF$ is thus an apotome [Prop.~10.73]. So, I say that (it is) also a
sixth (apotome).

For since $FL$ is equal to twice the (rectangle contained) by $AG$
and $GB$, let $FM$ be cut in half at $N$, and let $NO$
be drawn through $N$, parallel to $CD$. Thus, $FO$ and $NL$
are each equal to the (rectangle contained) by $AG$ and $GB$. And since
$AG$ and $GB$ are incommensurable in square, the (square) on $AG$
is thus incommensurable with the (square) on $GB$. But, $CH$ is equal to
the (square) on $AG$, and $KL$ is equal to the (square) on $GB$. 
Thus, $CH$ is incommensurable with $KL$. And as $CH$ (is) to
$KL$, so $CK$ is to $KM$ [Prop.~6.1].  Thus,
$CK$ is incommensurable (in length) with $KM$ [Prop.~10.11]. And since the (rectangle contained)
by $AG$ and $GB$ is the mean proportional to the (squares) on 
$AG$ and $GB$ [Prop.~10.21~lem.], and $CH$ is equal to the (square) on $AG$, and $KL$ 
equal to the (square) on $GB$, and $NL$ equal to the (rectangle contained)
by $AG$ and $GB$, $NL$ is thus also the mean proportional to
$CH$ and $KL$. Thus, as $CH$ is to $NL$, so $NL$ (is) to $KL$.
And for the same (reasons as the preceding propositions), the square on 
$CM$ is greater than (the square on) $MF$ by the (square) on (some
straight-line) incommensurable (in length) with ($CM$) [Prop.~10.18]. And neither
of them is commensurable with the (previously) laid down rational 
(straight-line) $CD$. Thus, $CF$ is a sixth apotome 
[Def.~10.16]. (Which is) the very thing it was required to show.\\}
\end{Parallel}

%%%%%
%10.103
%%%%%
\pdfbookmark[1]{Proposition 10.103}{pdf10.103}
\begin{Parallel}{}{}
\ParallelLText{
\begin{center}
{\large\ggn{103}.}
\end{center}\vspace*{-7pt}

{\greekfont Ἡ τῇ ἀποτομ ῇ μήκει σύμμετρος ἀποτομή ἐστι καὶ τῇ τάχει ἡ
α ὐτή.}

{\greekfont Ἔστω ἀποτομὴ ἡ ΑΒ, καὶ τῇ ΑΒ μήκει σύμμετρος ἔστω ἡ ΓΔ;
λέγω, ὅτι καὶ ἡ ΓΔ ἀποτομή ἐστι καὶ τῇ τάχει ἡ α ὐτὴ τῇ
ΑΒ.}

\epsfysize=0.7in
\centerline{\epsffile{Book10/fig103g.eps}}

{\greekfont Ἐπεὶ γὰρ ἀποτομή ἐστιν ἡ ΑΒ, ἔστω α ὐτῇ προσαρμόζουσα
ἡ ΒΕ; αἱ ΑΕ, ΕΒ ἄρα ῥηταί εἰσι δυνάμει μόνον σύμμετροι.
καὶ τῷ τῆς ΑΒ πρὸς τὴν ΓΔ λόγῳ ὁ α ὐτὸς γεγονέτω ὁ τῆς
ΒΕ πρὸς τὴν
 ΔΖ; καὶ ὡς ἓν ἄρα πρὸς ἕν, πάντα [ἐστὶ] πρὸς πάντα; ἔστιν ἄρα καὶ ὡς ὅλη ἡ ΑΕ πρὸς ὅλην τὴν
ΓΖ, οὕτως ἡ ΑΒ πρὸς τὴν ΓΔ. σύμμετρος δὲ ἡ ΑΒ τῇ
ΓΔ μήκει; σύμμετρος ἄρα καὶ ἡ ΑΕ μὲν τῇ ΓΖ, ἡ δὲ ΒΕ τῇ
ΔΖ. καὶ αἱ ΑΕ, ΕΒ ῥηταί εἰσι δυνάμει μόνον σύμμετροι; καὶ
αἱ
ΓΖ, ΖΔ ἄρα ῥηταί εἰσι δυνάμει μόνον σύμμετροι
[ἀποτομὴ ἄρα ἐστὶν ἡ ΓΔ. λέγω δή, ὅτι καὶ τῇ τάχει
ἡ α ὐτὴ τῇ ΑΒ].}

{\greekfont Ἐπεὶ οὖν ἐστιν ὡς ἡ ΑΕ πρὸς τὴν ΓΖ, οὕτως ἡ ΒΕ πρὸς τὴν
ΔΖ, ἐναλλὰχ ἄρα ἐστὶν ὡς ἡ ΑΕ πρὸς τὴν ΕΒ, οὕτως ἡ ΓΖ
πρὸς τὴν ΖΔ. ἤτοι δὴ ἡ ΑΕ τῆς ΕΒ μεῖζον δύναται τῷ ἀπὸ
συμμέτρου ἑαυτῇ ἢ τῷ ἀπὸ ἀσυμμέτρου. εἰ μὲν οὖν
ἡ ΑΕ τῆς ΕΒ μεῖζον δύναται τῷ ἀπὸ συμμέτρου ἑαυτῇ,
καὶ ἡ ΓΖ τῆς ΖΔ μεῖζον δυνήσεται τῷ ἀπὸ συμμέτρου
ἑαυτῇ. καὶ εἰ μὲν σύμμετρός ἐστιν ἡ ΑΕ τῇ ἐκκειμένῃ ῥητῇ
μήκει, καὶ ἡ ΓΖ, εἰ δὲ ἡ ΒΕ, καὶ ἡ ΔΖ, εἰ δὲ οὐδετέρα
τῶν ΑΕ, ΕΒ, καὶ οὐδετέρα τῶν ΓΖ, ΖΔ. εἰ δὲ ἡ ΑΕ [τῆς ΕΒ]
μεῖζον δύναται τῷ ἀπὸ ἀσυμμέτρου ἑαυτῇ, καὶ ἡ ΓΖ τῆς
ΖΔ μεῖζον δυνήσεται τῷ ἀπὸ ἀσυμμέτρου ἑαυτῇ. καὶ εἰ μὲν
σύμμετρός ἐστιν ἡ ΑΕ τῇ ἐκκειμένῃ ῥητῇ μήκει, καὶ ἡ ΓΖ,
εἰ δὲ ἡ ΒΕ, καὶ ἡ ΔΖ, εἰ δὲ οὐδετέρα τῶν ΑΕ, ΕΒ, οὐδετέρα
τῶν ΓΖ, ΖΔ.}

{\greekfont Ἀποτομὴ ἄρα ἐστὶν ἡ ΓΔ καὶ τῇ τάχει ἡ α ὐτὴ
τῇ ΑΒ; ὅπερ ἔδει δεῖχαι.}}

\ParallelRText{
\begin{center}
{\large Proposition 103}
\end{center}

A (straight-line) commensurable in length
with an apotome is an apotome, and  (is) the same in order.

Let $AB$ be an apotome, and let $CD$ be commensurable in length with $AB$. I say that $CD$ is also an apotome, and (is) the same in order as $AB$.

\epsfysize=0.7in
\centerline{\epsffile{Book10/fig103e.eps}}

For since $AB$ is an apotome, let $BE$ be an attachment to it. Thus,
$AE$ and $EB$ are rational (straight-lines which are) commensurable in
square only [Prop.~10.73]. And let it 
be contrived that the (ratio)
of $BE$ to $DF$ is the same as the ratio of $AB$ to $CD$
[Prop.~6.12]. Thus, also, as one is to one, (so)
all [are] to all [Prop.~5.12]. And thus as the whole $AE$ is to the whole $CF$, so $AB$
(is) to $CD$. And $AB$ (is) commensurable in length with $CD$.
$AE$ (is) thus also commensurable (in length) with $CF$, and $BE$
with $DF$ [Prop.~10.11]. And $AE$ and $BE$
are rational (straight-lines which are) commensurable in square only.
Thus, $CF$ and $FD$ are also rational (straight-lines which are) commensurable in square only [Prop.~10.13].
[$CD$ is thus an apotome. So, I say that (it is) also the same in order as $AB$.]

Therefore, since as $AE$ is to $CF$, so $BE$ (is) to $DF$, thus, alternately,
as $AE$ is to $EB$, so $CF$ (is) to $FD$ [Prop.~5.16]. So,  the square on $AE$
is greater than (the square on) $EB$ either by the (square) on (some straight-line)
commensurable,  or by the (square) on (some straight-line)
incommensurable, (in length) with ($AE$). Therefore, if the
(square) on $AE$ is greater than (the square on) $EB$ by the (square)
on (some straight-line) commensurable (in length) with ($AE$) then the
square on $CF$ will also be greater than (the square on) $FD$ by the (square)
on (some straight-line) commensurable (in length) with ($CF$)
[Prop.~10.14]. And if $AE$ is commensurable
in length with a (previously) laid down rational (straight-line) then so (is)
$CF$ [Prop.~10.12], and if $BE$ (is commensurable), so (is) $DF$, and if neither of
$AE$ or $EB$ (are commensurable), neither  (are) either of $CF$ or $FD$
[Prop.~10.13]. And if the (square) on
$AE$ is greater [than (the square on) $EB$] by the (square) on
(some straight-line) incommensurable (in length) with ($AE$) then the
(square) on $CF$ will also be greater than (the square on) $FD$ by the (square)
on (some straight-line) incommensurable (in length) with ($CF$)
[Prop.~10.14].  And if $AE$ is commensurable
in length with a (previously) laid down rational (straight-line), so (is)
$CF$ [Prop.~10.12], and if $BE$ (is commensurable), so (is) $DF$, and if neither of
$AE$ or $EB$ (are commensurable),  neither (are)  either of $CF$ or $FD$
[Prop.~10.13].

Thus, $CD$ is an apotome, and (is) the same in order as $AB$ [Defs.~10.11---10.16].
(Which is) the very thing it was required to show.}
\end{Parallel}

%%%%%
%10.104
%%%%%
\pdfbookmark[1]{Proposition 10.104}{pdf10.104}
\begin{Parallel}{}{}
\ParallelLText{
\begin{center}
{\large \ggn{104}.}
\end{center}\vspace*{-7pt}

{\greekfont Ἡ τῇ μέσης ἀποτομ ῇ σύμμετρος μέσης ἀποτομή ἐστι
καὶ τῇ τάχει ἡ α ὐτή.}\\

\epsfysize=0.7in
\centerline{\epsffile{Book10/fig103g.eps}}

{\greekfont Ἔστω μέσης ἀποτομὴ ἡ ΑΒ, καὶ τῇ ΑΒ μήκει σύμμετρος 
ἔστω ἡ ΓΔ; λέγω, ὅτι καὶ ἡ ΓΔ μέσης ἀποτομή ἐστι
καὶ τῇ τάχει ἡ α ὐτὴ τῇ ΑΒ.}

{\greekfont Ἐπεὶ γὰρ μέσης ἀποτομή ἐστιν ἡ ΑΒ, ἔστω α ὐτῇ προσαρμόζουσα
ἡ ΕΒ. αἱ ΑΕ, ΕΒ ἄρα μέσαι εἰσὶ δυνάμει μόνον σύμμετροι.
καὶ γεγονέτω ὡς ἡ ΑΒ πρὸς τὴν ΓΔ, οὕτως ἡ ΒΕ πρὸς
τὴν ΔΖ; σύμμετρος ἄρα [ἐστὶ] καὶ ἡ ΑΕ τῇ ΓΖ, ἡ δὲ ΒΕ τῇ ΔΖ.
αἱ δὲ ΑΕ, ΕΒ μέσαι εἰσὶ δυνάμει μόνον σύμμετροι; καὶ αἱ
ΓΖ, ΖΔ ἄρα μέσαι εἰσὶ δυνάμει μόνον σύμμετροι;
μέσης ἄρα ἀποτομή ἐστιν ἡ ΓΔ. λέγω δή, ὅτι καὶ
τῇ τάχει ἐστὶν ἡ α ὐτὴ τῇ ΑΒ.}

{\greekfont Ἐπεὶ [γάρ] ἐστιν ὡς ἡ ΑΕ πρὸς τὴν ΕΒ, οὕτως
ἡ ΓΖ πρὸς τὴν ΖΔ [ἀλλ ὡς μὲν ἡ ΑΕ πρὸς τὴν ΕΒ, οὕτως
τὸ ἀπὸ τῆς ΑΕ πρὸς τὸ ὑπὸ τῶν ΑΕ, ΕΒ, ὡς δὲ ἡ ΓΖ πρὸς
τὴν ΖΔ, οὕτως τὸ ἀπὸ τῆς ΓΖ πρὸς τὸ ὑπὸ τῶν ΓΖ, ΖΔ], ἔστιν
ἄρα καὶ ὡς τὸ ἀπὸ τῆς ΑΕ πρὸς τὸ ὑπὸ τῶν ΑΕ, ΕΒ, οὕτως
τὸ ἀπὸ τῆς ΓΖ πρὸς τὸ ὑπὸ τῶν ΓΖ, ΖΔ [καὶ ἐναλλὰχ
ὡς τὸ ἀπὸ τῆς ΑΕ πρὸς τὸ ἀπὸ τῆς ΓΖ, οὕτως τὸ ὑπὸ
τῶν ΑΕ, ΕΒ πρὸς τὸ ὑπὸ τῶν ΓΖ, ΖΔ]. σύμμετρον δὲ τὸ
ἀπὸ τῆς ΑΕ τῷ ἀπὸ τῆς ΓΖ; σύμμετρον ἄρα ἐστὶ καὶ τὸ
ὑπὸ τῶν ΑΕ, ΕΒ τῷ ὑπὸ τῶν ΓΖ, ΖΔ.
εἴτε οὖν ῥητόν ἐστι τὸ ὑπὸ τῶν ΑΕ, ΕΒ, ῥητὸν ἔσται
καὶ τὸ ὑπὸ τῶν ΓΖ, ΖΔ, εἴτε μέσον [ἐστὶ] τὸ ὑπὸ
τῶν ΑΕ, ΕΒ, μέσον [ἐστὶ] καὶ τὸ ὑπὸ τῶν ΓΖ, ΖΔ.}

{\greekfont Μέσης ἄρα ἀποτομή ἐστιν ἡ ΓΔ καὶ τῇ τάχει ἡ
α ὐτὴ τῇ ΑΒ; ὅπερ ἔδει δεῖχαι.}}

\ParallelRText{
\begin{center}
{\large Proposition 104}
\end{center}

A (straight-line) commensurable (in length)
with an apotome of a medial (straight-line) is an apotome of a
medial (straight-line), and (is) the same in order.

\epsfysize=0.7in
\centerline{\epsffile{Book10/fig103e.eps}}

Let $AB$ be an apotome of a medial (straight-line), and let $CD$
be commensurable in length with $AB$. I say that $CD$ is also an
apotome of a medial (straight-line), and (is) the same in order as $AB$.

For since $AB$ is an apotome of a medial (straight-line), let $EB$ be
an attachment to it. Thus, $AE$ and $EB$ are medial (straight-lines which are) commensurable in square only [Props.~10.74, 10.75]. And let it be contrived that as
$AB$ is to $CD$, so $BE$ (is) to $DF$ [Prop.~6.12]. 
Thus, $AE$ [is] also commensurable (in length) with $CF$, and $BE$
with $DF$ [Props.~5.12, 10.11]. And $AE$ and $EB$ are medial (straight-lines which are) commensurable in square only. $CF$ and $FD$
are thus also medial (straight-lines which are) commensurable in square only
[Props.~10.23, 10.13]. 
Thus, $CD$ is an apotome of a medial (straight-line) [Props.~10.74, 10.75].
So, I say that it is also the same in order as $AB$.

\mbox{[}For] since as $AE$ is to $EB$, so $CF$ (is) to $FD$ [Props.~5.12, 5.16] [but
as $AE$ (is) to $EB$, so the (square) on $AE$ (is) to the (rectangle
contained) by $AE$ and $EB$, and as $CF$ (is) to $FD$, so the
(square) on $CF$ (is) to the (rectangle contained) by $CF$ and $FD$],
thus as  the (square) on $AE$ is to the (rectangle contained) by $AE$
and $EB$, so the (square) on $CF$ also (is) to the (rectangle contained)
by $CF$ and $FD$ [Prop.~10.21~lem.] [and,
alternately, as the (square) on $AE$ (is) to the (square) on $CF$, so
the (rectangle contained) by $AE$ and $EB$ (is) to the (rectangle contained)
by $CF$ and $FD$]. And the (square) on $AE$ (is) commensurable
with the (square) on $CF$. Thus, the (rectangle contained) by $AE$ and
$EB$ is also commensurable with the (rectangle contained) by $CF$ and
$FD$ [Props.~5.16, 10.11].
Therefore, either the (rectangle contained) by $AE$ and $EB$ is rational, and
the (rectangle contained) by $CF$ and $FD$ will also be rational [Def.~10.4], or
the (rectangle contained) by $AE$ and $EB$ [is] medial, and the
(rectangle contained) by $CF$ and $FD$ [is] also medial [Prop.~10.23~corr.].

Therefore, $CD$ is the apotome of a medial (straight-line), and is
the same in order as $AB$ [Props.~10.74, 10.75]. (Which is) the very thing it was required to show.}
\end{Parallel}

%%%%%
%10.105
%%%%%
\pdfbookmark[1]{Proposition 10.105}{pdf10.105}
\begin{Parallel}{}{}
\ParallelLText{
\begin{center}
{\large \ggn{105}.}
\end{center}\vspace*{-7pt}

{\greekfont Ἡ τῇ ἐλάσσονι σύμμετρος ἐλάσσων ἐστίν.}

{\greekfont Ἔστω γὰρ ἐλάσσων ἡ ΑΒ καὶ τῇ ΑΒ σύμμετρος ἡ ΓΔ; λέγω,
ὅτι καὶ ἡ ΓΔ ἐλάσσων ἐστίν.}\\~\\

\epsfysize=0.7in
\centerline{\epsffile{Book10/fig103g.eps}}

{\greekfont Γεγονέτω γὰρ τὰ α ὐτά; καὶ ἐπεὶ αἱ ΑΕ, ΕΒ δυνάμει εἰσὶν
ἀσύμμετροι, καὶ αἱ ΓΖ, ΖΔ ἄρα δυνάμει εἰσὶν ἀσύμμετροι.
ἐπεὶ οὖν ἐστιν ὡς ἡ ΑΕ πρὸς τὴν
ΕΒ, οὕτως ἡ ΓΖ πρὸς τὴν ΖΔ, ἔστιν ἄρα καὶ ὡς τὸ ἀπὸ
τῆς ΑΕ πρὸς τὸ ἀπὸ τῆς ΕΒ, οὕτως τὸ ἀπὸ τῆς ΓΖ πρὸς τὸ
ἀπὸ τῆς ΖΔ. συνθέντι ἄρα ἐστὶν ὡς τὰ ἀπὸ τῶν ΑΕ, ΕΒ πρὸς
τὸ ἀπὸ τῆς ΕΒ, οὕτως τὰ ἀπὸ τῶν ΓΖ, ΖΔ πρὸς τὸ ἀπὸ τῆς
ΖΔ [καὶ ἐναλλάχ];  σύμμετρον δέ ἐστι τὸ ἀπὸ τῆς ΒΕ τῷ ἀπὸ
τῆς ΔΖ; σύμμετρον ἄρα καὶ τὸ συγκείμενον ἐκ τῶν ἀπὸ
τῶν ΑΕ, ΕΒ τετραγώνων τῷ συγκειμένῳ ἐκ τῶν ἀπὸ τῶν
ΓΖ, ΖΔ τετραγώνων. ῥητὸν δέ ἐστι τὸ συγκείμενον ἐκ τῶν
ἀπὸ τῶν ΑΕ, ΕΒ τετραγώνων; ῥητὸν ἄρα ἐστὶ καὶ τὸ συγκείμενον
ἐκ τῶν ἀπὸ τῶν ΓΖ, ΖΔ τετραγώνων. πάλιν, ἐπεί ἐστιν ὡς τὸ
ἀπὸ τῆς ΑΕ πρὸς τὸ ὑπὸ τῶν ΑΕ, ΕΒ, οὕτως τὸ ἀπὸ τῆς
ΓΖ πρὸς τὸ ὑπὸ τῶν ΓΖ, ΖΔ, σύμμετρον δὲ τὸ ἀπὸ τῆς ΑΕ
τετράγωνον τῷ ἀπὸ τῆς ΓΖ τετραγώνῳ, σύμμετρον ἄρα ἐστὶ
καὶ τὸ ὑπὸ τῶν ΑΕ, ΕΒ τῷ ὑπὸ τῶν ΓΖ, ΖΔ. μέσον δὲ τὸ
ὑπὸ τῶν ΑΕ, ΕΒ; μέσον ἄρα καὶ τὸ ὑπὸ τῶν ΓΖ, ΖΔ; αἱ
ΓΖ, ΖΔ ἄρα δυνάμει εἰσὶν ἀσύμμετροι ποιοῦσαι τὸ μὲν συγκείμενον ἐκ τῶν ἀπ α ὐτῶν τετραγώνων ῥητόν, τὸ δ
ὑπ α ὐτῶν μέσον.}

{\greekfont Ἐλάσσων ἄρα ἐστὶν ἡ ΓΔ; ὅπερ ἔδει δεῖχαι.}}

\ParallelRText{
\begin{center}
{\large Proposition 105}
\end{center}

A (straight-line) commensurable (in length) with a minor (straight-line) is a minor (straight-line).

For let $AB$ be a minor (straight-line), and (let) $CD$ (be) commensurable
(in length) with $AB$. I say that $CD$ is also a minor (straight-line).

\epsfysize=0.7in
\centerline{\epsffile{Book10/fig103e.eps}}

For let the same things be contrived (as in the former proposition).
And since $AE$ and $EB$ are (straight-lines which are) incommensurable in square [Prop.~10.76], $CF$ and $FD$ are thus also
(straight-lines which are)  incommensurable in square [Prop.~10.13].
Therefore, since as $AE$ is to $EB$, so $CF$ (is) to $FD$ [Props.~5.12, 5.16], thus also
as the (square) on $AE$ is to the (square) on $EB$, so the (square)
on $CF$ (is) to the (square) on $FD$ [Prop.~6.22]. Thus, via composition,
as the (sum of the squares) on $AE$ and $EB$ is to the (square) on $EB$,
so the (sum of the squares) on $CF$ and $FD$  (is) to the (square) on $FD$  [Prop.~5.18], [also alternately]. And the
(square) on $BE$ is commensurable with the (square) on $DF$ [Prop.~10.104]. The sum of the squares on $AE$ and $EB$ (is) thus also commensurable with the sum of the squares
on $CF$ and $FD$ [Prop.~5.16, 10.11].  And the sum of the (squares) on
$AE$ and $EB$ is rational [Prop.~10.76].
Thus, the sum of the (squares) on $CF$ and $FD$ is also rational
[Def.~10.4]. Again, since as the (square) on $AE$
is to the (rectangle contained) by $AE$ and $EB$, so
the (square) on $CF$ (is) to the (rectangle contained) by
$CF$ and $FD$ [Prop.~10.21~lem.],
and the square on $AE$ (is) commensurable with the square on
$CF$, the (rectangle contained) by $AE$ and $EB$ is thus also commensurable
with the (rectangle contained) by $CF$ and $FD$. And the
(rectangle contained) by $AE$ and $EB$ (is) medial [Prop.~10.76].  Thus, the (rectangle contained) by
$CF$ and $FD$ (is) also medial [Prop.~10.23~corr.]. $CF$ and $FD$
are thus (straight-lines which are) incommensurable in square, making the sum of the
squares on them rational, and the (rectangle contained) by them
medial.

Thus, $CD$ is a minor (straight-line) [Prop.~10.76]. (Which is) the very thing it
was required to show.}
\end{Parallel}

%%%%%
%10.106
%%%%%
\pdfbookmark[1]{Proposition 10.106}{pdf10.106}
\begin{Parallel}{}{}
\ParallelLText{
\begin{center}
{\large \ggn{106}.}
\end{center}\vspace*{-7pt}

{\greekfont Ἡ τῇ μετὰ ῥητοῦ μέσον τὸ ὅλον ποιούσῃ σύμμετρος μετὰ
ῥητοῦ μέσον τὸ ὅλον ποιοῦσά ἐστιν.}

{\greekfont Ἔστω μετὰ ῥητοῦ μέσον τὸ ὅλον ποιοῦσα ἡ ΑΒ καὶ τῇ
ΑΒ σύμμετρος ἡ ΓΔ; λέγω, ὅτι καὶ ἡ ΓΔ μετὰ ῥητοῦ
μέσον τὸ ὅλον ποιοῦσά ἐστιν.}


{\greekfont Ἔστω γὰρ τῇ ΑΒ προσαρμόζουσα ἡ ΒΕ; αἱ ΑΕ, ΕΒ ἄρα δυνάμει
εἰσὶν ἀσύμμετροι ποιοῦσαι τὸ μὲν συγκείμενον ἐκ τῶν ἀπὸ
τῶν ΑΕ, ΕΒ τετραγώνων μέσον, τὸ δ ὑπ α ὐτῶν ῥητόν. καὶ
τὰ α ὐτὰ κατεσκευάσθω. ὁμοίως δὴ δείχομεν τοῖς πρότερον,
ὅτι αἱ ΓΖ, ΖΔ ἐν τῷ α ὐτῷ λόγῳ εἰσὶ ταῖς ΑΕ, ΕΒ,
καὶ σύμμετρόν ἐστι τὸ συγκείμενον ἐκ τῶν ἀπὸ τῶν ΑΕ,
ΕΒ τετραγώνων τῷ συγκειμένῳ ἐκ τῶν ἀπὸ τῶν
ΓΖ, ΖΔ τετραγώνων, τὸ δὲ ὑπὸ τῶν ΑΕ, ΕΒ τῷ ὑπὸ
τῶν ΓΖ, ΖΔ; ὥστε καὶ αἱ ΓΖ, ΖΔ δυνάμει εἰσὶν
ἀσύμμετροι ποιοῦσαι τὸ μὲν συγκείμενον ἐκ τῶν
ἀπὸ τῶν ΓΖ, ΖΔ τετραγώνων μέσον, τὸ δ ὑπ α ὐτῶν
ῥητόν.}\\~\\~\\~\\~\\~\\

\epsfysize=0.7in
\centerline{\epsffile{Book10/fig103g.eps}}

{\greekfont Ἡ ΓΔ ἄρα μετὰ ῥητοῦ μέσον τὸ  ὅλον ποιοῦσά
ἐστιν; ὅπερ ἔδει δεῖχαι.}}

\ParallelRText{
\begin{center}
{\large Proposition 106}
\end{center}

A (straight-line) commensurable (in length)
with a (straight-line) which with a rational (area) makes a medial
whole is a (straight-line) which with a rational (area) makes a medial
whole.

Let $AB$ be a (straight-line) which with a rational (area) makes
a medial whole, and (let) $CD$ (be) commensurable
(in length) with $AB$. I say that $CD$ is also a (straight-line)
which with a rational (area) makes a medial (whole).

For let $BE$ be an attachment to $AB$. Thus, $AE$ and
$EB$ are (straight-lines which are) incommensurable in square, making the sum of the
squares on $AE$ and $EB$ medial,  and the (rectangle contained)
by them rational [Prop.~10.77]. 
And let the same construction be made (as in the previous propositions).
So, similarly
to the previous  (propositions), we can  show that $CF$ and $FD$
are in the same ratio as $AE$ and $EB$, and the sum of the squares
on $AE$ and $EB$ is commensurable with the sum of the squares on
$CF$ and $FD$, and the (rectangle contained) by $AE$ and $EB$
with the (rectangle contained) by $CF$ and $FD$. Hence, $CF$ and
$FD$ are also (straight-lines which are) incommensurable in square, making the sum of the
squares on $CF$ and $FD$ medial, and the (rectangle contained) by them
rational.

\epsfysize=0.7in
\centerline{\epsffile{Book10/fig103e.eps}}

$CD$ is thus a (straight-line) which with a rational (area) makes a
medial whole [Prop.~10.77]. (Which is)
the very thing it was required to show.}
\end{Parallel}

%%%%%
%10.107
%%%%%
\pdfbookmark[1]{Proposition 10.107}{pdf10.107}
\begin{Parallel}{}{}
\ParallelLText{
\begin{center}
{\large \ggn{107}.}
\end{center}\vspace*{-7pt}

{\greekfont Ἡ τῇ μετὰ μέσου μέσον τὸ ὅλον ποιούσῃ σύμμετρος καὶ α ὐτὴ
μετὰ μέσου μέσον τὸ ὅλον ποιοῦσά ἐστιν.}\\~\\

\epsfysize=0.7in
\centerline{\epsffile{Book10/fig103g.eps}}

{\greekfont Ἔστω μετὰ μέσου μέσον τὸ ὅλον ποιοῦσα ἡ ΑΒ, καὶ τῇ
ΑΒ ἔστω σύμμετρος ἡ ΓΔ; λέγω, ὅτι καὶ ἡ ΓΔ μετὰ μέσου
μέσον τὸ ὅλον ποιοῦσά ἐστιν.}

{\greekfont Ἔστω γὰρ τῇ ΑΒ προσαρμόζουσα ἡ ΒΕ, καὶ τὰ α ὐτὰ
κατεσκευάσθω; αἱ ΑΕ, ΕΒ ἄρα δυνάμει εἱσὶν ἀσύμμετροι
ποιοῦσαι τό τε συγκείμενον ἐκ τῶν ἀπ α ὐτῶν τετραγώνων
μέσον καὶ τὸ ὑπ α ὐτῶν μέσον καὶ ἔτι ἀσύμμετρον τὸ
συγκέιμενον ἐκ τῶν ἀπ α ὐτῶν τετραγώνων τῷ ὑπ
α ὐτῶν. καί εἰσιν, ὡς ἐδείθθη, αἱ ΑΕ, ΕΒ σύμμετροι ταῖς
ΓΖ, ΖΔ, καὶ τὸ συγκείμενον ἐκ τῶν ἀπὸ τῶν ΑΕ, ΕΒ
τετραγώνων τῷ συγκειμένῳ ἐκ τῶν ἀπὸ τῶν ΓΖ, ΖΔ, τὸ δὲ
ὑπὸ τῶν ΑΕ, ΕΒ τῷ ὑπὸ τῶν ΓΖ, ΖΔ; καὶ αἱ ΓΖ, ΖΔ ἄρα
δυνάμει εἰσὶν ἀσύμμετροι ποιοῦσαι τό τε συγκείμενον
ἐκ τῶν ἀπ α ὐτῶν τετραγώνων μέσον καὶ τὸ ὑπ
ἀ ὐτῶν μέσον καὶ ἔτι ἀσύμμετρον τὸ συγκείμενον ἐκ τῶν
ἀπ α ὐτῶν [τετραγώνων] τῷ ὑπ α ὐτῶν.}

{\greekfont Ἡ ΓΔ ἄρα μετὰ μέσου μέσον τὸ ὅλον ποιοῦσά
ἐστιν; ὅπερ ἔδει δεῖχαι.}}

\ParallelRText{
\begin{center}
{\large Proposition 107}
\end{center}

A (straight-line) commensurable
(in length) with a (straight-line) which with a medial (area) makes a
medial whole is  itself also a (straight-line) which with a medial (area) makes
a medial whole.

\epsfysize=0.7in
\centerline{\epsffile{Book10/fig103e.eps}}

Let $AB$ be a (straight-line) which with a medial (area) makes a medial
whole, and let $CD$ be commensurable (in length) with $AB$.
I say that $CD$ is also a (straight-line) which with a medial (area) makes
a medial whole.

For let $BE$ be an attachment to $AB$. And let the same
construction be made (as in the previous propositions). Thus,
$AE$ and $EB$ are (straight-lines which are) incommensurable in square, making
the sum of the squares on them medial, and the (rectangle contained)
by them medial, and, further, the sum of the squares on them incommensurable with the (rectangle contained) by them [Prop.~10.78].  And, as was shown (previously),
$AE$ and $EB$ are commensurable (in length) with $CF$ and $FD$ (respectively),
and the sum of the squares on $AE$ and $EB$ with the sum of the squares on
$CF$ and $FD$, and the (rectangle contained) by $AE$ and $EB$ with
the (rectangle contained) by $CF$ and $FD$. Thus, $CF$ and
$FD$ are also (straight-lines which are) incommensurable in square, making the sum of the squares on
them medial, and the (rectangle contained) by them medial, and, further,
the sum of the [squares] on them incommensurable with the (rectangle
contained) by them.

Thus, $CD$ is a (straight-line) which with a medial (area) makes a medial
whole [Prop.~10.78]. (Which is) the very thing
it was required to show.}
\end{Parallel}

%%%%%
%10.108
%%%%%
\pdfbookmark[1]{Proposition 10.108}{pdf10.108}
\begin{Parallel}{}{}
\ParallelLText{
\begin{center}
{\large \ggn{108}.}
\end{center}\vspace*{-7pt}

{\greekfont Ἀπὸ ῥητοῦ μέσου ἀφαιρουμένου ἡ τὸ λοιπὸν θωρίον δυναμένη 
μία δύο ἀλόγων γίνεται ἤτοι ἀποτομὴ ἢ ἐλάσσων.}\\~\\

\epsfysize=1.5in
\centerline{\epsffile{Book10/fig108g.eps}}

{\greekfont Ἀπὸ γὰρ ῥητοῦ τοῦ ΒΓ μέσον ἀφῃρήσθω τὸ ΒΔ; λέγω,
ὅτι ἡ τὸ λοιπὸν δυναμένη τὸ ΕΓ μία δύο ἀλόγων γίνεται
ἤτοι ἀποτομὴ ἢ ἐλάσσων.}

{\greekfont Ἐκκείσθω γὰρ ῥητὴ ἡ ΖΗ, καὶ τῷ μὲν ΒΓ ἴσον
παρὰ τὴν ΖΗ παραβεβλήσθω ὀρθογώνιον παραλληλόγραμμον
τὸ ΗΘ, τῷ δὲ ΔΒ ἴσον ἀφῃρήσθω τὸ ΗΚ;
λοιπὸν ἄρα τὸ ΕΓ ἴσον ἐστὶ τῷ ΛΘ. ἐπεὶ οὖν
ῥητὸν μέν ἐστι τὸ ΒΓ, μέσον δὲ τὸ ΒΔ, ἴσον δὲ τὸ μὲν
ΒΓ τῷ ΗΘ, τὸ δὲ ΒΔ τῷ ΗΚ, ῥητὸν μὲν ἄρα ἐστὶ τὸ ΗΘ,
μέσον δὲ τὸ ΗΚ. καὶ παρὰ ῥητὴν τὴν ΖΗ παράκειται;  ῥητὴ
μὲν ἄρα ἡ ΖΘ καὶ σύμμετρος τῇ ΖΗ μήκει, ῥητὴ δὲ ἡ
ΖΚ καὶ ἀσύμμετρος τῇ ΖΗ μήκει; ἀσύμμετρος ἄρα ἐστὶν
ἡ ΖΘ τῇ ΖΚ μήκει. αἱ ΖΘ, ΖΚ ἄρα ῥηταί εἰσι δυνάμει
μόνον σύμμετροι; ἀποτομὴ ἄρα ἐστὶν ἡ ΚΘ, προσαρμόζουσα
δὲ α ὐτῇ ἡ ΚΖ. ἤτοι δὴ ἡ ΘΖ τῆς ΖΚ μεῖζον δύναται
τῷ ἀπὸ συμμέτρου ἢ οὔ.}

{\greekfont Δυνάσθω πρότερον τῷ ἀπὸ συμμέτρου. καί ἐστιν ὅλη ἡ ΘΖ
σύμμετρος τῇ ἐκκειμένῃ ῥητῇ μήκει τῇ ΖΗ;
ἀποτομὴ ἄρα πρώτη ἐστὶν ἡ ΚΘ. τὸ δ ὑπὸ ῥητῆς
καὶ ἀποτομῆς πρώτης περιεχόμενον ἡ δυναμένη
ἀποτομή ἐστιν. ἡ ἄρα τὸ ΛΘ, τουτέστι τὸ ΕΓ, δυναμένη
ἀποτομή ἐστιν.}

{\greekfont Εἰ δὲ ἡ ΘΖ τῆς ΖΚ μεῖζον δύναται τῷ ἀπὸ ἀσυμμέτρου
ἑαυτῇ, καί ἐστιν ὅλη ἡ ΖΘ σύμμετρος τῇ ἐκκειμένῃ ῥητῇ
μήκει τῇ ΖΗ, ἀποτομὴ τετάρτη ἐστὶν ἡ ΚΘ. τὸ δ ὑπὸ
ῥητῆς καὶ
 ἀποτομῆς τετάρτης περιεχόμενον ἡ δυναμένη
ἐλάσσων ἐστίν; ὅπερ ἔδει δεῖχαι.}}

\ParallelRText{
\begin{center}
{\large Proposition 108}
\end{center}

A
medial (area) being subtracted from a rational (area), one of two irrational (straight-lines)
arise (as) the square-root of the remaining area---either an apotome, or a minor (straight-line).

\epsfysize=1.5in
\centerline{\epsffile{Book10/fig108e.eps}}

For let the medial (area) $BD$ be subtracted from the
rational (area) $BC$. I say that  one of two irrational (straight-lines) arise
(as) the square-root of the
remaining (area), $EC$---either
an apotome, or a minor (straight-line).

For let the rational (straight-line) $FG$ be laid out,
and let the right-angled parallelogram $GH$, equal to $BC$, 
be applied to $FG$, and let $GK$, equal to $DB$, 
be subtracted (from $GH$).  Thus, the remainder $EC$
is equal to $LH$. Therefore, since $BC$
is a rational (area), and $BD$ a medial (area), and $BC$
(is) equal to $GH$, and $BD$ to $GK$, $GH$ is thus a rational (area),
and $GK$ a medial (area). And they are applied to the rational (straight-line)
$FG$. Thus, $FH$ (is) rational, and commensurable in length with $FG$
[Prop.~10.20], and $FK$ (is) also rational, and
incommensurable in length with $FG$ [Prop.~10.22]. Thus, $FH$ is incommensurable
in length with $FK$ [Prop.~10.13]. 
$FH$ and $FK$ are thus rational (straight-lines which are) commensurable
in square only. Thus, $KH$ is an apotome [Prop.~10.73], and $KF$ an attachment to it.
So, the square on $HF$ is  greater than (the square on) $FK$ by
the (square) on (some straight-line which is) either commensurable,
or not (commensurable), (in length with $HF$).

First, let the square (on it) be (greater) by the (square) on  (some straight-line
which is) commensurable
(in length with $HF$). And the whole of $HF$ is commensurable
in length with the (previously) laid down rational (straight-line) $FG$. 
Thus, $KH$ is a first apotome [Def.~10.1]. And
the square-root of an (area) contained by a rational (straight-line)
and a first apotome is an apotome [Prop.~10.91]. 
Thus, the square-root of $LH$---that is to say, (of) $EC$---is an apotome.

And if the square on  $HF$ is greater than (the square on) $FK$ by the
(square) on (some straight-line which is) incommensurable (in length) with ($HF$),
and (since) the whole of $FH$ is commensurable in length with
the (previously) laid down rational (straight-line) $FG$, $KH$ is a fourth
apotome [Prop.~10.14]. And the square-root
of an (area) contained by a rational (straight-line) and a fourth apotome
is a minor (straight-line)
[Prop.~10.94]. (Which is) the very thing it was required to show.}
\end{Parallel}

%%%%%
%10.109
%%%%%
\pdfbookmark[1]{Proposition 10.109}{pdf10.109}
\begin{Parallel}{}{}
\ParallelLText{
\begin{center}
{\large \ggn{109}.}
\end{center}\vspace*{-7pt}

{\greekfont Ἀπὸ μέσου ῥητοῦ ἀφαιρουμένου ἄλλαι δύο ἄλογοι γίνονται
ἤτοι μέσης ἀποτομὴ πρώτη ἢ μετὰ ῥητοῦ μέσον τὸ ὅλον ποιοῦσα.}

{\greekfont Ἀπὸ γὰρ μέσου τοῦ ΒΓ ῥητὸν ἀφῃρήσθω τὸ ΒΔ. λέγω, ὅτι
ἡ τὸ λοιπὸν τὸ ΕΓ δυναμένη μία δύο ἀλόγων γίνεται ἤτοι
μέσης ἀποτομὴ πρώτη ἢ μετὰ ῥητοῦ μέσον τὸ ὅλον
ποιοῦσα.}

{\greekfont Ἐκκείσθω γὰρ ῥητὴ ἡ ΖΗ, καὶ παραβεβλήσθω ὁμοίως
τὰ θωρία. ἔστι δὴ ἀκολούθως ῥητὴ μὲν ἡ ΖΘ καὶ ἀσύμμετρος
τῇ ΖΗ μήκει, ῥητὴ δὲ ἡ ΚΖ καὶ σύμμετρος τῇ ΖΗ μήκει; αἱ
ΖΘ, ΖΚ ἄρα ῥηταί εἰσι δυνάμει μόνον σύμμετροι; ἀποτομὴ
ἄρα ἐστὶν ἡ ΚΘ, προσαρμόζουσα δὲ ταύτῃ ἡ ΖΚ. ἤτοι δὴ ἡ ΘΖ
τῆς ΖΚ μεῖζον δύναται τῷ ἀπὸ συμμέτρου ἑαυτῇ ἢ
τῷ ἀπὸ ἀσυμμέτρου.}\\~\\~\\~\\~\\~\\~\\~\\

\epsfysize=2.in
\centerline{\epsffile{Book10/fig109g.eps}}

{\greekfont Εἰ μὲν οὖν ἡ ΘΖ τῆς ΖΚ μεῖζον δύναται τῷ ἀπὸ συμμέτρου
ἑαυτῇ, καί ἐστιν ἡ προσαρμόζουσα ἡ ΖΚ σύμμετρος τῇ ἐκκειμένῃ
ῥητῇ μήκει τῇ ΖΗ, ἀποτομὴ δευτέρα ἐστὶν ἡ ΚΘ. ῥητὴ δὲ
ἡ ΖΗ; ὥστε ἡ τὸ ΛΘ, τουτέστι τὸ ΕΓ, δυναμένη μέσης ἀποτομὴ
πρώτη ἐστίν.}

{\greekfont Εἰ δὲ ἡ ΘΖ τῆς ΖΚ μεῖζον δύναται τῷ ἀπὸ ἀσυμμέτρου, καί
ἐστιν ἡ προσαρμόζουσα ἡ ΖΚ σύμμετ-ρος τῇ ἐκκειμένῃ ῥητῇ
μήκει τῇ ΖΗ, ἀποτομὴ πέμπτη ἐστὶν ἡ ΚΘ; ὥστε ἡ τὸ
ΕΓ δυναμένη μετὰ ῥητοῦ μέσον τὸ ὅλον ποιοῦσά
ἐστιν; ὅπερ ἔδει δεῖχαι.}}

\ParallelRText{
\begin{center}
{\large Proposition 109}
\end{center}

A rational (area) being subtracted from
a medial (area), two other irrational (straight-lines) arise (as the square-root
of the remaining area)---either a first apotome of a medial (straight-line),
or that (straight-line) which with a rational (area) makes a medial whole.

For let the rational (area) $BD$ be subtracted from the medial
(area) $BC$. I say that one of two irrational (straight-lines) arise (as) the square-root of the remaining (area), $EC$---either a first apotome of a medial
(straight-line), or that (straight-line) which with a rational (area) makes
a medial whole.

For let the rational (straight-line) $FG$ be laid down, and let
similar areas (to the preceding proposition) have  been applied (to it). So,
accordingly, $FH$ is rational, and incommensurable in length with
$FG$, and $KF$ (is) also rational, and commensurable
in length with $FG$. Thus, $FH$ and $FK$ are rational (straight-lines which
are) commensurable in square only [Prop.~10.13].
$KH$ is thus an apotome [Prop.~10.73], and $FK$ an attachment to it. So, the square on $HF$ is greater than (the square on)
$FK$ either by the (square) on (some straight-line) commensurable (in length)
with ($HF$), or by the (square) on (some straight-line) incommensurable
(in length with $HF$).

\epsfysize=2.in
\centerline{\epsffile{Book10/fig109e.eps}}

Therefore, if the square on $HF$ is greater than (the square on) $FK$
by the (square) on (some straight-line) commensurable (in length) with ($HF$),
and (since) the attachment $FK$ is commensurable
in length with the (previously) laid down rational (straight-line) $FG$,
$KH$ is a second apotome [Def.~10.12]. 
And $FG$ (is) rational. Hence, the square-root of $LH$---that is to say, (of)
$EC$---is a first apotome of a medial (straight-line) [Prop.~10.92].

And if the square on $HF$ is greater than (the square on) $FK$
by the (square) on (some straight-line) incommensurable (in length with $HF$),
and (since) the attachment $FK$ is commensurable
in length with the (previously) laid down rational (straight-line) $FG$,
$KH$ is a fifth apotome [Def.~10.15]. Hence,
the square-root of $EC$ is that (straight-line) which with a rational (area)
makes a medial whole [Prop.~10.95]. 
(Which is) the very thing it was required to show.}
\end{Parallel}

%%%%%
%10.110
%%%%%
\pdfbookmark[1]{Proposition 10.110}{pdf10.110}
\begin{Parallel}{}{}
\ParallelLText{
\begin{center}
{\large \ggn{110}.}
\end{center}\vspace*{-7pt}

{\greekfont Ἀπὸ μέσου μέσου ἀφαιρουμένου ἀσυμμέτρου τῷ ὅλῳ αἱ λοιπαὶ
δύο ἄλογοι γίνονται ἤτοι μέσης ἀποτομὴ δευτέρα ἢ μετὰ
μέσου μέσον τὸ ὅλον ποιοῦσα.}

{\greekfont Ἀφῃρήσθω γὰρ ὡς ἐπὶ τῶν προκειμένων καταγραφῶν ἀπὸ μέσου
τοῦ ΒΓ μέσον τὸ ΒΔ ἀσύμμετρον τῷ ὅλῳ; λέγω, ὅτι
ἡ τὸ ΕΓ δυναμένη μία ἐστὶ δύο ἀλόγων ἤτοι μέσης
ἀποτομὴ δευτέρα ἢ μετὰ μέσου μέσον τὸ ὅλον
ποιοῦσα.}\\~\\~\\~\\

\epsfysize=2.2in
\centerline{\epsffile{Book10/fig109g.eps}}

{\greekfont Ἐπεὶ γὰρ μέσον ἐστὶν ἑκάτερον τῶν ΒΓ, ΒΔ, καὶ ἀσύμμετρον
τὸ ΒΓ τῷ ΒΔ, ἔσται ἀκολούθως ῥητὴ ἑκατέρα τῶν ΖΘ, ΖΚ
καὶ ἀσύμμετρος τῇ ΖΗ μήκει. καὶ ἐπεὶ ἀσύμμετρόν ἐστι
τὸ ΒΓ τῷ ΒΔ, τουτέστι τὸ ΗΘ τῷ ΗΚ, ἀσύμμετρος καὶ ἡ ΘΖ
τῇ ΖΚ; αἱ ΖΘ, ΖΚ ἄρα ῥηταί εἰσι δυνάμει μόνον σύμμετροι;
ἀποτομὴ ἄρα ἐστὶν ἡ ΚΘ [προσαρμόζουσα δὲ ἡ ΖΚ. ἤτοι
δὴ ἡ ΖΘ τῆς ΖΚ μεῖζον δύναται τῷ ἀπὸ συμμέτρου
ἢ τῷ ἀπὸ ἀσυμμέτρου ἑαυτῇ].}

{\greekfont Εἰ μὲν δὴ ἡ ΖΘ τῆς ΖΚ μεῖζον δύναται τῷ ἀπὸ συμμέτρου
ἑαυτῇ, καὶ οὐθετέρα τῶν ΖΘ, ΖΚ σύμμετρός ἐστι τῇ
ἐκκειμέμνῃ ῥητῇ μήκει τῇ ΖΗ, ἀποτομὴ τρίτη ἐστὶν
ἡ ΚΘ. ῥητὴ δὲ ἡ ΚΛ, τὸ δ ὑπὸ ῥητῆς
καὶ ἀποτομῆς τρίτης περιεχόμενον ὀρθογώνιον ἄλογόν
ἐστιν, καὶ ἡ δυναμένη α ὐτὸ ἄλογός ἐστιν, καλεῖται δὲ μέσης
ἀποτομὴ δευτέρα; ὥστε ἡ τὸ ΛΘ, τουτέστι τὸ ΕΓ, δυναμένη
μέσης ἀποτομή ἐστι δευτερά.}

{\greekfont Εἰ δὲ ἡ ΖΘ τῆς ΖΚ μεῖζον δύναται τῷ ἀπὸ
ἀσυμμέτρου ἑαυτῇ [μήκει], καὶ οὐθετέρα τῶν ΘΖ, ΖΚ
σύμμετρός ἐστι τῇ ΖΗ μήκει, ἀποτομὴ ἕκτη ἐστὶν ἡ ΚΘ.
τὸ δ ὑπὸ ῥητῆς καὶ ἀποτομῆς ἕκτης
ἡ δυναμένη ἐστὶ μετὰ μέσου μέσον τὸ ὅλον ποιοῦσα.
ἡ τὸ ΛΘ ἄρα, τουτέστι τὸ  ΕΓ, δυναμένη μετὰ μέσου μέσον
τὸ ὅλον ποιοῦσά ἐστιν; ὅπερ ἔδει δεῖχαι.}}

\ParallelRText{
\begin{center}
{\large Proposition 110}
\end{center}

A medial (area), incommensurable with the whole, being subtracted
from a medial (area),  the two remaining
irrational (straight-lines) arise (as) the (square-root of the area)---either a second apotome of a medial (straight-line), 
or that (straight-line) which with a medial (area) makes a medial whole.

For, as in the previous figures, let the medial (area) $BD$, incommensurable with the whole, be subtracted from the medial (area) $BC$. I say that
the square-root of $EC$ is one of two irrational (straight-lines)---either
a second apotome of a medial (straight-line), or that (straight-line) which
with a medial (area) makes a medial whole.

\epsfysize=2.2in
\centerline{\epsffile{Book10/fig109e.eps}}

For since $BC$ and $BD$ are each medial (areas), and $BC$ (is) incommensurable with $BD$, accordingly,  $FH$ and $FK$ will each be rational (straight-lines), and incommensurable
in length with $FG$ [Prop.~10.22]. And
since $BC$ is incommensurable with $BD$---that is to say, $GH$ with
$GK$---$HF$ (is) also incommensurable (in length) with $FK$ [Props.~6.1, 10.11]. Thus,
$FH$ and $FK$ are rational (straight-lines which are) commensurable
in square only. $KH$ is thus as apotome [Prop.~10.73], [and $FK$ an attachment (to it).
So, the square on $FH$ is  greater than (the square on)
$FK$ either by the (square) on (some straight-line) commensurable, or by the (square) on (some straight-line) incommensurable, (in length)
with ($FH$).]

So, if the square on $FH$ is greater than (the square on) $FK$ by the
(square) on (some straight-line) commensurable (in length) with ($FH$),
and (since) neither of $FH$ and $FK$ is commensurable in length with
the (previously) laid down rational (straight-line) $FG$, $KH$ is a
third apotome [Def.~10.3]. And $KL$ (is) rational.
And the rectangle contained by a rational (straight-line) and a third
apotome is irrational, and the square-root of it is that irrational (straight-line)
called a second apotome of a medial (straight-line) [Prop.~10.93]. Hence, the square-root of $LH$---that is to say, (of) $EC$---is a second apotome of a medial (straight-line).

And if the square on $FH$ is greater than (the square on) $FK$ by the
(square) on (some straight-line) incommensurable [in length] with ($FH$), 
and (since) neither of $HF$ and $FK$ is commensurable in length
with $FG$, $KH$ is a sixth apotome [Def.~10.16].
And the square-root of the (rectangle contained) by a rational (straight-line) and a sixth apotome is that (straight-line) which with a medial (area)
makes a medial whole [Prop.~10.96]. Thus,
the square-root of $LH$---that is to say, (of) $EC$---is 
that (straight-line) which with a medial (area) makes a medial whole.
(Which is) the very thing it was required to show.}
\end{Parallel}

%%%%%
%10.111
%%%%%
\pdfbookmark[1]{Proposition 10.111}{pdf10.111}
\begin{Parallel}{}{}
\ParallelLText{
\begin{center}
{\large \ggn{111}.}
\end{center}\vspace*{-7pt}

{\greekfont Ἡ ἀποτομὴ οὐκ ἔστιν ἡ α ὐτὴ τῇ ἐκ δύο ὀνομάτων.}

\epsfysize=2.4in
\centerline{\epsffile{Book10/fig111g.eps}}

{\greekfont Ἔστω ἀποτομὴ ἡ ΑΒ; λέγω, ὅτι ἡ ΑΒ οὐκ ἔστιν ἡ α ὐτὴ τῇ
ἐκ δύο ὀνομάτων.}

{\greekfont Εἰ γὰρ δυνατόν, ἔστω; καὶ ἐκκείσθω ῥητὴ ἡ ΔΓ, καὶ τῷ
ἀπὸ τῆς ΑΒ ἴσον παρὰ τὴν ΓΔ παραβεβλήσθω ὀρθογώνιον
τὸ ΓΕ πλάτος ποιοῦν τὴν ΔΕ. ἐπεὶ οὖν ἀποτομή ἐστιν ἡ ΑΒ, ἀποτομὴ
πρώτη ἐστὶν ἡ ΔΕ. ἔστω α ὐτῇ προσαρμόζουσα ἡ ΕΖ; αἱ
ΔΖ, ΖΕ ἄρα ῥηταί εἰσι δυνάμει μόνον σύμμετροι, καὶ
ἡ ΔΖ τῆς ΖΕ μεῖζον δύναται τῷ ἀπὸ συμμέτρου ἑαυτῇ,
καὶ ἡ ΔΖ σύμμετρός ἐστι τῇ ἐκκειμένῃ ῥητῇ μήκει
τῇ ΔΓ. πάλιν, ἐπεὶ ἐκ δύο ὀνομάτων ἐστὶν ἡ ΑΒ, ἐκ
δύο ἄρα  ὀνομάτων πρώτη ἐστὶν ἡ ΔΕ. διῃρήσθω εἰς
τὰ ὀνόματα κατὰ τὸ Η, καὶ ἔστω μεῖζον ὄνομα τὸ ΔΗ;
αἱ ΔΗ, ΗΕ ἄρα ῥηταί εἰσι δυνάμει μόνον σύμμετροι,
καὶ ἡ ΔΗ τῆς ΗΕ μεῖζον δύναται τῷ ἀπὸ
συμμέτρου ἑαυτῇ, καὶ τὸ μεῖζον ἡ ΔΗ σύμμετρός ἐστι
τῇ ἐκκειμένῃ ῥητῇ μήκει τῇ ΔΓ. καὶ ἡ ΔΖ ἄρα τῇ
ΔΗ σύμμετρός ἐστι μήκει; καὶ λοιπὴ ἄρα ἡ ΗΖ σύμμετρός
ἐστι τῇ ΔΖ μήκει. [ἐπεὶ οῦν σύμμετρός ἐστιν ἡ ΔΖ
τῇ ΗΖ, ῥητὴ δέ ἐστιν ἡ ΔΖ, ῥητὴ ἄρα ἐστὶ καὶ ἡ
ΗΖ. ἐπεὶ οὖν σύμμετρός ἐστιν 
 ἡ ΔΖ
τῇ ΗΖ μήκει] ἀσύμμετρος δὲ ἡ ΔΖ τῇ ΕΖ μήκει. 
ἀσύμμετρος ἄρα ἐστὶ καὶ ἡ ΖΗ τῇ ΕΖ μήκει.
αἱ
ΗΖ, ΖΕ ἄρα  ῥηταί [εἰσι] δυνάμει μόνον σύμμετροι; ἀποτομὴ
ἄρα ἐστὶν ἡ ΕΗ. ἀλλὰ καὶ ῥητή; ὅπερ ἐστὶν ἀδύνατον.}

{\greekfont Ἡ ἄρα ἀποτομὴ οὐκ ἔστιν ἡ α ὐτὴ τῇ ἐκ δύο
ὀνομάτων; ὅπερ ἔδει δεῖχαι.}}

\ParallelRText{
\begin{center}
{\large Proposition 111}
\end{center}

An apotome is not the same as a binomial.

\epsfysize=2.4in
\centerline{\epsffile{Book10/fig111e.eps}}

Let $AB$ be an apotome. I say that $AB$ is not the same as a binomial.

For, if possible, let it be (the same). And let a rational (straight-line)
$DC$ be laid down. And let the rectangle $CE$, equal to the (square) on $AB$, be applied to $CD$, producing $DE$ as breadth. Therefore,
since $AB$ is an apotome, $DE$ is a first apotome [Prop.~10.97]. Let $EF$ be an attachment to it.
Thus, $DF$ and $FE$ are rational (straight-lines which are) commensurable
in square only, and the square on $DF$ is greater than (the square on)
$FE$ by the (square) on (some straight-line) commensurable (in length)
with ($DF$), and $DF$ is commensurable in length with the (previously)
laid down rational (straight-line) $DC$ [Def.~10.10]. 
Again, since $AB$ is a binomial, $DE$ is thus a first binomial [Prop.~10.60]. Let ($DE$) be divided
into its (component) terms at $G$, and let $DG$ be the greater term.
Thus, $DG$ and $GE$ are rational (straight-lines which are)
commensurable in square only, and the square on $DG$ is greater
than (the square on) $GE$ by the (square) on (some straight-line)
commensurable (in length) with ($DG$), and the greater (term)
$DG$ is commensurable in length with the (previously)
laid down rational (straight-line) $DC$ [Def.~10.5].
Thus, $DF$ is also commensurable in length with $DG$ [Prop.~10.12].
The remainder $GF$ is thus commensurable in length with $DF$
[Prop.~10.15]. [Therefore, since
$DF$ is commensurable with $GF$, and $DF$ is rational, $GF$
is thus also rational. Therefore, since $DF$ is commensurable
in length with $GF$,] $DF$ (is) incommensurable in length with $EF$.
Thus, $FG$ is also incommensurable in length with $EF$ [Prop.~10.13]. $GF$ and $FE$ [are] thus
rational (straight-lines which are) commensurable in square only.
Thus, $EG$ is an apotome [Prop.~10.73].
But, (it is) also rational. The very thing is impossible.

Thus, an apotome is not the same as a binomial. (Which is) the very thing
it was required to show.}
\end{Parallel}

\begin{Parallel}{}{}
\ParallelLText{
\begin{center}
{\large {\greekfont [Πόρισμα.]}}
\end{center}\vspace*{-7pt}

{\greekfont Ἡ ἀποτομὴ καὶ αἱ μετ α ὐτὴν ἄλογοι οὔτε
τῇ μέσῃ οὔτε ἀλλήλαις εἰσὶν αἱ α ὐταί.}

{\greekfont Τὸ μὲν γὰρ ἀπὸ μέσης παρὰ ῥητὴν παραβαλλόμενον
πλάτος ποιεῖ ῥητὴν καὶ ἀσύμμετρον τῇ, παρ
ἣν παράκειται, μήκει, τὸ δὲ ἀπὸ ἀποτομῆς παρὰ ῥητὴν
παραβαλλόμενον πλάτος ποιεῖ ἀποτομὴν πρώτην, τὸ
δὲ ἀπὸ μέσης ἀποτομῆς πρώτης παρὰ ῥητὴν παραβαλλόμενον
πλάτος ποιεῖ ἀποτομὴν δευτέραν, τὸ δὲ ἀπὸ
μέσης ἀποτομῆς δευτέρας παρὰ ῥητὴν παραβαλλόμενον
πλάτος ποιεῖ ἀποτομὴν τρίτην, τὸ δὲ ἀπὸ ἐλάσσονος
παρὰ ῥητὴν παραβαλλόμενον πλάτος ποιεῖ
ἀποτομὴν τετάρτην, τὸ δὲ ἀπὸ τῆς μετὰ ῥητοῦ μέσον
τὸ ὅλον ποιούσης παρὰ ῥητὴν παραβαλλόμενον πλάτος
ποιεῖ ἀποτομὴν πέμπτην, τὸ δὲ ἀπὸ τῆς μετὰ μέσου
μέσον τὸ ὅλον ποιούσης παρὰ ῥητὴν παραβαλλόμενον
πλάτος ποιεῖ ἀποτομὴν ἕκτην. ἐπεὶ οὖν τὰ εἰρημένα
πλάτη διαφέρει τοῦ τε πρώτου καὶ ἀλλήλων, τοῦ μὲν πρώτου,
ὅτι ῥητή ἐστιν, ἀλλήλων δὲ, ἐπεὶ τῇ τάχει
οὐκ εἰσὶν αἱ α ὐταί, δῆλον, ὡς καὶ α ὐταὶ αἱ ἄλογοι
διαφέρουσιν ἀλλήλων. καὶ ἐπεὶ δέδεικται ἡ ἀποτομὴ
οὐκ οὖσα ἡ α ὐτὴ τῇ ἐκ δύο ὀνομάτων, ποιοῦσι
δὲ πλάτη παρὰ ῥητὴν παραβαλλόμεναι αἱ μετὰ τὴν ἀποτομὴν
ἀποτομὰς
ἀκολούθως ἑκάστη τῇ τάχει τῇ καθ α ὐτήν, αἱ δὲ μετὰ τὴν
ἐκ δύο ὀνομάτων τὰς ἐκ δύο ὀνομάτων καὶ α ὐταὶ τῇ
τάχει
 ἀκολούθως, ἕτεραι ἄρα εἰσὶν αἱ μετὰ
τὴν ἀποτομὴν καὶ ἕτεραι αἱ μετὰ τὴν ἐκ δύο
ὀνομάτων, ὡς εἶναι τῇ τάχει πάσας ἀλόγουςὸv{ιγ},}}

\ParallelRText{
\begin{center}
{\large [Corollary]}
\end{center}\vspace*{-7pt}

The apotome and the  irrational (straight-lines) after it are neither the
same as a medial (straight-line) nor (the same) as one another.

For the (square) on a medial (straight-line), applied to a rational (straight-line), produces as breadth a rational (straight-line which is) incommensurable in length with the
(straight-line) to which  (the area) is applied [Prop.~10.22]. And the (square) on an apotome, applied to a rational (straight-line), produces as breadth a first apotome [Prop.~10.97]. 
And the (square) on a first
apotome of a medial (straight-line), applied to a rational (straight-line),
produces as breadth a second apotome [Prop.~10.98].
And the (square) on a second
apotome of a medial (straight-line), applied to a rational (straight-line),
produces as breadth a third apotome [Prop.~10.99].
And (square) on a minor (straight-line), applied to a rational (straight-line),
produces as breadth a fourth apotome [Prop.~10.100]. And (square) on that (straight-line) which with a rational (area) produces a medial whole, applied to a rational (straight-line),
produces as breadth a fifth apotome [Prop.~10.101]. And (square) on that (straight-line) which with a medial (area) produces a medial whole, applied to a rational
 (straight-line), produces as breadth a sixth apotome [Prop.~10.102]. Therefore, since the aforementioned breadths differ from the first (breadth), and from one another---from the
first, because it is rational, and from one another since they are not the
same in order---clearly,  the  irrational (straight-lines)
themselves also differ from one another. And since it has been shown that an apotome
is not the same as a binomial [Prop.~10.111],
and (that) the (irrational straight-lines) after the apotome, being applied to
a rational (straight-line), produce as breadth, each according to its own (order), apotomes, and (that) the (irrational straight-lines) after
the binomial  themselves also (produce as breadth), according (to their) order, binomials, 
the (irrational straight-lines) after the apotome are thus different, and the
(irrational straight-lines) after the binomial (are also) different, so that
there are, in order, 13 irrational (straight-lines) in all:}
\end{Parallel}

\begin{Parallel}{}{}
\ParallelLText{

{\greekfont Μέσην,}\vspace*{-7pt}

{\greekfont Ἐκ δύο ὀνομάτων,}\vspace*{-7pt}

{\greekfont Ἐκ δύο μέσων πρώτην,}\vspace*{-7pt}

{\greekfont Ἐκ δύο μέσων δευτέραν,}\vspace*{-7pt}

{\greekfont Μείζονα,}\vspace*{-7pt}

{\greekfont μὴῬητὸν καὶ μέσον δυναμένην,}\vspace*{-7pt}

{\greekfont Δύο μέσα δυναμένην,}\vspace*{-7pt}

{\greekfont Ἀποτομήν,}\vspace*{-7pt}

{\greekfont Μἑσης ἀποτομὴν πρώτην,}\vspace*{-7pt}

{\greekfont Μἑσης ἀποτομὴν δευτέραν,}\vspace*{-7pt}

{\greekfont Ἐλάσσονα},\vspace*{-7pt}

{\greekfont Μετὰ ῥητοῦ μέσον τὸ ὅλον ποιοῦσαν,}\vspace*{-7pt}

{\greekfont Μετὰ μέσου μέσον τὸ ὅλον ποιοῦσαν.}}

\ParallelRText{

Medial,\vspace*{-7pt}

Binomial,\vspace*{-7pt}

First bimedial,\vspace*{-7pt}

Second bimedial,\vspace*{-7pt}

Major,\vspace*{-7pt}

Square-root of a rational plus a medial (area),\vspace*{-7pt}

Square-root of (the sum of) two medial (areas),\vspace*{-7pt}

Apotome,\vspace*{-7pt}

First apotome of a medial,\vspace*{-7pt}

Second apotome of a medial,\vspace*{-7pt}

Minor,\vspace*{-7pt}

That which with a rational (area) produces a medial whole,\vspace*{-7pt}

That which with a medial (area) produces a medial whole.}
\end{Parallel}

%%%%%
%10.112
%%%%%
\pdfbookmark[1]{Proposition 10.112}{pdf10.112}
\begin{Parallel}{}{}
\ParallelLText{
\begin{center}
{\large \ggn{112}.}
\end{center}\vspace*{-7pt}

{\greekfont Τὸ ἀπὸ ῥητῆς παρὰ τὴν ἐκ δύο ὀνομάτων παραβαλλόμενον
πλάτος ποιεῖ ἀποτομήν, ἧς τὰ ὀνόματα σύμμετρά ἐστι τοῖς
τῆς ἐκ δύο ὀνομάτων ὀνόμασι καὶ ἔτι ἐν τῷ α ὐτῷ λόγῳ,
καὶ ἔτι ἡ γινομένη ἀποτομὴ τὴν α ὐτὴν ἕχει τάχιν τῇ ἐκ
δύο ὀνομάτων.}\\

\epsfysize=0.9in
\centerline{\epsffile{Book10/fig112g.eps}}

{\greekfont Ἔστω ῥητὴ μὲν ἡ Α, ἐκ δύο ὀνομάτων δὲ ἡ ΒΓ, ἧς μεῖζον
ὄνομα ἔστω ἡ ΔΓ, καὶ τῷ ἀπὸ τῆς Α ἴσον ἔστω τὸ
ὑπὸ τῶν ΒΓ, ΕΖ; λέγω, ὅτι ἡ ΕΖ ἀποτομή ἐστιν, ἧς τὰ ὀνόματα
σύμμετρά ἐστι τοῖς ΓΔ, ΔΒ, καὶ ἐν τῷ α ὐτῷ λόγῳ,
καὶ ἔτι ἡ ΕΖ τὴν α ὐτὴν ἕχει τάχιν τῇ ΒΓ.}

{\greekfont Ἔστω γὰρ πάλιν τῷ ἀπὸ τῆς Α ἴσον τὸ ὑπὸ τῶν ΒΔ, Η.
ἐπεὶ οὖν τὸ ὑπὸ τῶν ΒΓ, ΕΖ ἴσον ἐστὶ τῷ ὑπὸ
τῶν ΒΔ, Η, ἔστιν ἄρα ὡς ἡ ΓΒ πρὸς τὴν ΒΔ,  οὕτως ἡ Η
πρὸς τὴν ΕΖ. μείζων δὲ ἡ ΓΒ τῆς ΒΔ; μείζων ἄρα ἐστὶ καὶ
ἡ Η τῆς ΕΖ. ἔστω τῇ Η ἴση ἡ ΕΘ; ἔστιν ἄρα ὡς ἡ ΓΒ πρὸς
τὴν ΒΔ, οὕτως ἡ ΘΕ πρὸς τὴν ΕΖ; διελόντι ἄρα ἐστὶν ὡς ἡ ΓΔ
πρὸς τὴν ΒΔ, οὕτως ἡ ΘΖ πρὸς τὴν ΖΕ. γεγονέτω ὡς
ἡ ΘΖ πρὸς τὴν ΖΕ, οὕτως ἡ ΖΕ πρὸς τὴν ΚΕ;
καὶ ὅλη ἄρα ἡ ΘΚ πρὸς ὅλην τὴν ΚΖ ἐστιν, ὡς ἡ ΖΚ
πρὸς ΚΕ; ὡς γὰρ ἓν τῶν ἡγουμένων πρὸς ἓν τῶν
ἑπομένων, οὕτως ἅπαντα τὰ ἡγούμενα πρὸς ἅπαντα
τὰ ἑπόμενα. ὡς δὲ ἡ ΖΚ πρὸς ΚΕ, οὕτως ἐστὶν ἡ ΓΔ
πρὸς τὴν ΔΒ; καὶ ὡς ἄρα ἡ ΘΚ πρὸς ΚΖ, οὕτως
ἡ ΓΔ πρὸς τὴν ΔΒ. σύμμετρον δὲ τὸ ἀπὸ τῆς ΓΔ τῷ
ἀπὸ τῆς ΔΒ; σύμμετρον ἄρα ἐστὶ καὶ τὸ ἀπὸ τῆς ΘΚ
τῷ ἀπὸ τῆς ΚΖ. καί ἐστιν ὡς τὸ ἀπὸ τῆς ΘΚ
πρὸς τὸ ἀπὸ τῆς ΚΖ,
οὕτως ἡ ΘΚ πρὸς τὴν ΚΕ, ἐπεὶ αἱ
τρεῖς
αἱ ΘΚ, ΚΖ, ΚΕ ἀνάλογόν εἰσιν. σύμμετρος
ἄρα ἡ ΘΚ τῇ ΚΕ μήκει. ὥστε καὶ ἡ ΘΕ τῇ ΕΚ
σύμμετρός ἐστι μήκει.
καὶ ἐπεὶ τὸ ἀπὸ τῆς Α ἴσον
ἐστὶ τῷ ὑπὸ τῶν ΕΘ, ΒΔ, ῥητὸν δέ ἐστι τὸ ἀπὸ
τῆς Α, ῥητὸν ἄρα ἐστὶ καὶ τὸ ὑπὸ τῶν ΕΘ, ΒΔ. καὶ παρὰ
ῥητὴν τὴν ΒΔ παράκειται; ῥητὴ ἄρα ἐστὶν ἡ ΕΘ καὶ σύμμετρος
τῇ ΒΔ μήκει; ὥστε καὶ ἡ σύμμετρος α ὐτῇ ἡ ΕΚ ῥητή ἐστι
καὶ σύμμετρος τῇ ΒΔ μήκει. ἐπεὶ οὖν ἐστιν ὡς ἡ ΓΔ
πρὸς ΔΒ, οὕτως ἡ ΖΚ πρὸς ΚΕ, αἱ δὲ ΓΔ, ΔΒ δυνάμει
μόνον εἰσὶ σύμμετροι, καὶ αἱ ΖΚ, ΚΕ δυνάμει
μόνον εἰσὶ σύμμετροι. ῥητὴ δέ ἐστιν ἡ ΚΕ; ῥητὴ ἄρα
ἐστὶ καὶ ἡ ΖΚ. αἱ ΖΚ, ΚΕ ἄρα ῥηταὶ δυνάμει μόνον εἰσὶ
σύμμετροι; ἀποτομὴ ἄρα ἐστὶν ἡ ΕΖ.}

{\greekfont Ἤτοι δὲ ἡ ΓΔ τῆς ΔΒ μεῖζον δύναται τῷ ἀπὸ συμμέτρου
ἑαυτῇ ἢ τῷ ἀπὸ ἀσυμμέτρου.}

{\greekfont Εἰ μὲν οὖν ἡ ΓΔ τῆς ΔΒ μεῖζον δύναται τῷ ἀπὸ
συμμέτρου [ἑαυτῇ], καὶ ἡ ΖΚ τῆς ΚΕ μεῖζον
δυνήσεται τῷ ἀπὸ συμμέτρου ἑαυτῇ. καὶ εἰ μὲν
σύμμετρός ἐστιν ἡ ΓΔ τῇ ἐκκειμένῃ ῥητῇ μήκει,
καὶ ἡ ΖΚ; εἰ δὲ ἡ ΒΔ, καὶ ἡ ΚΕ; εἰ δὲ οὐδετέρα τῶν ΓΔ,
ΔΒ, καὶ οὐδετέρα τῶν ΖΚ, ΚΕ.}

{\greekfont Εἰ δὲ ἡ ΓΔ τῆς ΔΒ μεῖζον δύναται τῷ ἀπὸ ἀσυμμέτρου
ἑαυτῇ, καὶ ἡ ΖΚ τῆς  ΚΕ μεῖζον δυνήσεται τῷ ἀπὸ
ἀσυμμέτρου ἑαυτῇ. καὶ εἰ μὲν ἡ ΓΔ σύμμετρός ἐστι
τῇ ἐκκειμένῃ ῥητῇ μήκει, καὶ ἡ ΖΚ; εἰ δὲ ἡ ΒΔ, καὶ ἡ ΚΕ;
εἰ δὲ οὐδετέρα τῶν ΓΔ, ΔΒ, καὶ οὐδετέρα τῶν ΖΚ, ΚΕ;
ὥστε ἀποτομή ἐστιν ἡ ΖΕ, ἧς τὰ ὀνόματα τὰ ΖΚ, ΚΕ σύμμετρά
ἐστι τοῖς τῆς ἐκ δύο ὀνομάτων ὀνόμασι τοῖς ΓΔ, ΔΒ
καὶ ἐν τῷ α ὐτῷ λόγῳ, καὶ τὴν α ὐτῆν τάχιν ἔχει τῇ
ΒΓ; ὅπερ ἔδει δεῖχαι.}}

\ParallelRText{
\begin{center}
{\large Proposition 112}$^\dag$
\end{center}

The (square) on a rational (straight-line), applied
to a binomial (straight-line), produces  as breadth an apotome whose terms are commensurable (in length) with the terms of the
binomial, and, furthermore, in the same ratio. Moreover,
the created apotome will have the same order as the binomial.

\epsfysize=0.9in
\centerline{\epsffile{Book10/fig112e.eps}}

Let $A$ be a rational (straight-line), and $BC$ a binomial
(straight-line), of which let $DC$ be the greater term. And let the
(rectangle contained) by $BC$ and $EF$ be equal to the (square) on $A$.
I say that $EF$ is an apotome whose terms are commensurable (in length)
with $CD$ and $DB$, and in the same ratio, and, moreover, that $EF$
will have the same order as $BC$.

For, again, let the (rectangle contained) by $BD$ and $G$  be equal
to the (square) on $A$. Therefore, since the (rectangle contained) by $BC$
and $EF$ is equal to the (rectangle contained) by $BD$ and $G$, thus as
$CB$ is to $BD$, so $G$ (is) to $EF$ [Prop.~6.16].
And $CB$ (is) greater than $BD$. Thus, $G$ is also greater than $EF$
[Props.~5.16, 5.14].
Let $EH$ be equal to $G$. Thus, as $CB$ is to $BD$, so $HE$
(is) to $EF$. Thus, via separation, as $CD$ is to $BD$, so $HF$ (is) to $FE$
[Prop.~5.17]. Let it be contrived that as
$HF$ (is) to $FE$, so $FK$ (is) to $KE$. And, thus, the
whole $HK$ is to the whole $KF$, as $FK$ (is) to $KE$.
For as one of the leading (proportional magnitudes is) to one of
the following, so all of the leading (magnitudes) are to all of the
following [Prop.~5.12]. And as $FK$
(is) to $KE$, so $CD$ is to $DB$ [Prop.~5.11].  And, thus, as $HK$ (is) to $KF$, so
$CD$ is to $DB$ [Prop.~5.11].
And the (square) on  $CD$ (is) 
 commensurable with the (square) on $DB$ [Prop.~10.36]. The (square) on $HK$ is
 thus also commensurable with the (square) on $KF$ [Props.~6.22, 10.11].
  And as the (square) on $HK$ is to the (square) on $KF$, so $HK$ (is)
  to $KE$, since the three (straight-lines) $HK$, $KF$, and $KE$ are
  proportional [Def.~5.9]. $HK$ is thus commensurable
  in length with $KE$ [Prop.~10.11]. 
  Hence, $HE$ is also commensurable in length with $EK$
  [Prop.~10.15]. And since the (square) on $A$
  is equal to the (rectangle contained) by $EH$ and $BD$, and the (square)
  on $A$ is rational, the (rectangle contained) by $EH$ and $BD$ is thus
  also rational. And it is applied to the rational (straight-line) $BD$.
  Thus, $EH$ is  rational, and commensurable in length with $BD$ [Prop.~10.20]. And, hence, the (straight-line)
  commensurable (in length) with it, $EK$, is also rational [Def.~10.3], and commensurable
  in length with $BD$ [Prop.~10.12].  Therefore,
  since as $CD$  is to $DB$, so $FK$ (is) to $KE$, and $CD$ and $DB$
  are (straight-lines which are) commensurable in square only, $FK$ and $KE$ are also
  commensurable in square only [Prop.~10.11]. 
  And $KE$ is rational. Thus, $FK$ is also rational. $FK$ and
  $KE$ are thus rational (straight-lines which are) commensurable in square
  only. Thus, $EF$ is an apotome [Prop.~10.73].
  
  And the square on $CD$ is greater than (the square on) $DB$ either
  by the (square) on (some straight-line) commensurable, or by the
  (square) on (some straight-line) incommensurable, (in length) with ($CD$).
  
  Therefore, if the square on $CD$ is greater than (the square on) $DB$
  by the (square) on (some straight-line) commensurable  (in length) with [$CD$]  then the square on $FK$ will also be greater than (the square on) $KE$
  by the (square) on (some straight-line) commensurable  (in length) with ($FK$) [Prop.~10.14]. And if $CD$ is commensurable in length with a (previously) laid down rational (straight-line), (so) also (is) $FK$ [Props.~10.11, 10.12]. And if $BD$ (is commensurable), (so) also (is) $KE$ [Prop.~10.12]. And if neither of
  $CD$ or $DB$ (is commensurable), neither also (are) either of $FK$ or $KE$.
  
And if the square on $CD$ is greater than (the square on) $DB$
  by the (square) on (some straight-line) incommensurable  (in length) with ($CD$)  then the square on $FK$ will also be greater than (the square on) $KE$
  by the (square) on (some straight-line) incommensurable  (in length) with ($FK$) [Prop.~10.14]. And if $CD$ is commensurable in length with a (previously) laid down rational (straight-line), (so) also  (is) $FK$ [Props.~10.11, 10.12]. And if $BD$ (is commensurable), (so) also (is) $KE$ [Prop.~10.12]. And if neither of
  $CD$ or $DB$ (is commensurable), neither also (are) either of $FK$ or $KE$.
  Hence, $FE$ is an apotome whose terms, $FK$ and $KE$, are
  commensurable (in length) with the terms, $CD$ and $DB$, of the binomial,
  and in the same ratio. And ($FE$) has the same order as $BC$ [Defs.~10.5---10.10]. 
  (Which is) the very thing it was required to show.}
  \end{Parallel}
{\footnotesize\noindent$^\dag$ Heiberg considers this proposition, and
the succeeding ones, to be  relatively early interpolations into the original
text.}

%%%%%
%10.113
%%%%%
\pdfbookmark[1]{Proposition 10.113}{pdf10.113}
\begin{Parallel}{}{}
\ParallelLText{
\begin{center}
{\large \ggn{113}.}
\end{center}\vspace*{-7pt}

{\greekfont Τὸ ἀπὸ ῥητῆς παρὰ ἀποτομὴν παραβαλλόμενον πλάτος ποιεῖ
τὴν ἐκ δύο ὀνομάτων, ἧς τὰ ὀνόματα σύμμετρά ἐστι
τοῖς τῆς ἀποτομῆς ὀνόμασι καὶ ἐν τῷ α ὐτῷ λόγῳ,
ἔτι δὲ ἡ γινομένη ἐκ δύο ὀνομάτων τὴν α ὐτὴν τάχιν
ἔχει τῇ ἀποτομ ῇ.}\\

\epsfysize=0.9in
\centerline{\epsffile{Book10/fig113g.eps}}

{\greekfont Ἔστω ῥητὴ μὲν ἡ Α, ἀποτομὴ δὲ ἡ ΒΔ, καὶ τῷ ἀπὸ
τῆς Α ἴσον ἔστω τὸ ὑπὸ τῶν ΒΔ, ΚΘ, ὥστε τὸ ἀπὸ
τῆς Α ῥητῆς παρὰ τὴν ΒΔ ἀποτομὴν παραβαλλόμενον
πλάτος ποιεῖ τὴν ΚΘ; λέγω, ὅτι ἐκ δύο ὀνομάτων ἐστὶν
ἡ ΚΘ, ἧς τὰ ὀνόματα σύμμετρά ἐστι τοῖς τῆς ΒΔ
ὀνόμασι καὶ ἐν τῷ α ὐτῷ λόγῳ, καὶ ἔτι ἡ ΚΘ τὴν
α ὐτὴν ἔχει τάχιν τῇ ΒΔ.}

{\greekfont Ἔστω γὰρ τῇ ΒΔ προσαρμόζουσα ἡ ΔΓ; αἱ ΒΓ, ΓΔ ἄρα
ῥηταί εἰσι δυνάμει μόνον σύμμετροι. καὶ τῷ ἀπὸ τῆς
Α ἴσον ἔστω καὶ τὸ ὑπὸ τῶν ΒΓ, Η. 
ῥητὸν δὲ τὸ ἀπὸ τῆς Α; ῥητὸν ἄρα καὶ τὸ ὑπὸ
τῶν ΒΓ, Η.
καὶ παρὰ ῥητὴν
τὴν ΒΓ παραβέβληται; ῥητὴ ἄρα ἐστὶν ἡ Η καὶ σύμμετρος
τῇ ΒΓ μήκει. ἐπεὶ οὖν τὸ ὑπὸ τῶν ΒΓ, Η ἴσον
ἐστὶ τῷ ὑπὸ τῶν ΒΔ, ΚΘ, ἀνάλογον ἄρα ἐστὶν ὡς ἡ
ΓΒ πρὸς ΒΔ, οὕτως ἡ ΚΘ πρὸς Η. μείζων δὲ ἡ ΒΓ
τῆς ΒΔ; μείζων ἄρα καὶ ἡ ΚΘ τῆς Η. κείσθω τῇ Η ἴση
ἡ ΚΕ; σύμμετρος ἄρα ἐστὶν ἡ ΚΕ τῇ ΒΓ μήκει. καὶ ἐπεί
ἐστιν ὡς ἡ ΓΒ πρὸς ΒΔ, οὕτως ἡ ΘΚ πρὸς ΚΕ, ἀναστρέψαντι
ἄρα ἐστὶν ὡς ἡ ΒΓ πρὸς τὴν ΓΔ, οὕτως ἡ ΚΘ πρὸς ΘΕ.
γεγονέτω ὡς ἡ ΚΘ πρὸς ΘΕ, οὕτως ἡ ΘΖ πρὸς ΖΕ; καὶ λοιπὴ ἄρα
ἡ ΚΖ πρὸς ΖΘ ἐστιν, ὡς ἡ ΚΘ πρὸς ΘΕ, τουτέστιν [ὡς] ἡ ΒΓ
πρὸς ΓΔ. αἱ δὲ ΒΓ, ΓΔ δυνάμει μόνον [εἰσὶ] σύμμετροι;
καὶ αἱ ΚΖ, ΖΘ ἄρα δυνάμει μόνον εἰσὶ σύμμετροι; καὶ
ἐπεί ἐστιν ὡς ἡ ΚΘ πρὸς ΘΕ, ἡ ΚΖ πρὸς ΖΘ, ἀλλ ὡς ἡ
ΚΘ πρὸς ΘΕ, ἡ ΘΖ πρὸς ΖΕ, καὶ ὡς ἄρα ἡ ΚΖ πρὸς ΖΘ, ἡ ΘΖ πρὸς
ΖΕ; ὥστε καὶ ὡς ἡ πρώτη πρὸς τὴν τρίτην, τὸ ἀπὸ τῆς
πρώτης πρὸς τὸ ἀπὸ τῆς δευτέρας; καὶ ὡς ἄρα ἡ ΚΖ πρὸς
ΖΕ, οὕτως τὸ ἀπὸ τῆς ΚΖ πρὸς τὸ ἀπὸ τῆς ΖΘ. σύμμετρον
δέ ἐστι τὸ ἀπὸ τῆς ΚΖ τῷ ἀπὸ τῆς ΖΘ; αἱ γὰρ ΚΖ, ΖΘ
δυνάμει εἰσὶ σύμμετροι; σύμμετρος ἄρα ἐστὶ καὶ ἡ ΚΖ
τῇ ΖΕ μήκει; ὥστε ἡ ΚΖ καὶ τῇ ΚΕ σύμμετρός [ἐστι]
μήκει. ῥητὴ δέ ἐστιν ἡ ΚΕ καὶ σύμμετρος τῇ ΒΓ μήκει.
ῥητὴ ἄρα καὶ ἡ ΚΖ καὶ σύμμετρος τῇ ΒΓ μήκει. καὶ ἐπεί
ἐστιν ὡς ἡ ΒΓ πρὸς ΓΔ, οὕτως ἡ ΚΖ πρὸς ΖΘ, ἐναλλὰχ
ὡς ἡ ΒΓ πρὸς ΚΖ, οὕτως ἡ ΔΓ πρὸς ΖΘ. σύμμετρος
δὲ ἡ ΒΓ τῇ ΚΖ; σύμμετρος ἄρα καὶ ἡ ΖΘ τῇ ΓΔ
μήκει. αἱ ΒΓ, ΓΔ δὲ ῥηταί εἰσι δυνάμει μόνον σύμμετροι;
καὶ αἱ ΚΖ, ΖΘ ἄρα ῥηταί εἰσι δυνάμει μόνον σύμμετροι;
ἐκ δύο ὀνομάτων ἐστὶν ἄρα ἡ ΚΘ.}

{\greekfont Εἰ μὲν οὖν ἡ ΒΓ τῆς ΓΔ μεῖζον δύναται τῷ ἀπὸ
συμμέτρου ἑαυτῇ, καὶ ἡ ΚΖ τῆς ΖΘ μεῖζον δυνήσεται
τῷ ἀπὸ συμμέτρ-ου ἑαυτῇ. καὶ εἰ μὲν σύμμετρός
ἐστιν ἡ ΒΓ τῇ ἐκκειμένῃ ῥητῇ μήκει, καὶ ἡ ΚΖ, εἰ δὲ
ἡ ΓΔ σύμμετρός ἐστι τῇ ἐκκειμένῃ ῥητῇ μήκει,
καὶ ἡ ΖΘ, εἰ δὲ οὐδετέρα τῶν ΒΓ, ΓΔ, οὐδετέρα
τῶν ΚΖ, ΖΘ.}

{\greekfont Εἰ δὲ ἡ ΒΓ τῆς ΓΔ μεῖζον δύναται τῷ ἀπὸ ἀσυμμέτρου
ἑαυτῇ, καὶ ἡ ΚΖ τῆς ΖΘ μεῖζον δυνήσεται τῷ ἀπὸ
ἀσυμμέτρ-ου ἑαυτῇ. καὶ εἰ μὲν σύμμετρός ἐστιν
ἡ ΒΓ τῇ ἐκκειμένῃ ῥητῇ μήκει, καὶ ἡ ΚΖ,
εἰ δὲ ἡ ΓΔ, καὶ ἡ ΖΘ, εἰ δὲ οὐδετέρα τῶν ΒΓ, ΓΔ, οὐδετέρα
τῶν ΚΖ, ΖΘ.}

{\greekfont Ἐκ δύο ἄρα ὀνομάτων ἐστὶν ἡ ΚΘ, ἧς τὰ ὀνόματα
τὰ ΚΖ, ΖΘ σύμμετρά [ἐστι] τοῖς τῆς ἀποτομῆς ὀνόμασι
τοῖς ΒΓ, ΓΔ καὶ ἐν τῷ α ὐτῷ λόγῳ, καὶ ἔτι ἡ ΚΘ
τῇ ΒΓ τὴν α ὐτὴν ἕχει τάχιν; ὅπερ ἔδει δεῖχαι.}}

\ParallelRText{
\begin{center}
{\large Proposition 113}
\end{center}

The (square) on a rational (straight-line),
applied to an apotome, produces as breadth a binomial whose terms
are commensurable with the terms of the apotome, and in  the same ratio.
Moreover, the created binomial has the same order as the apotome.

\epsfysize=0.9in
\centerline{\epsffile{Book10/fig113e.eps}}

Let $A$ be a rational (straight-line), and $BD$ an apotome. And let the
(rectangle contained) by $BD$ and $KH$ be equal to the (square) on $A$,
such that the square on the rational (straight-line) $A$, applied to the apotome
$BD$, produces $KH$ as breadth. I say that $KH$ is a binomial whose
terms are commensurable with the terms of $BD$, and in the same ratio, and,
moreover, that $KH$ has the same order as $BD$.

For let $DC$ be an attachment to $BD$. Thus, $BC$ and $CD$
are rational (straight-lines which are) commensurable in square only [Prop.~10.73]. And let the (rectangle contained)
by $BC$ and $G$ also be equal to the (square) on $A$. And the (square)
on $A$ (is) rational. The (rectangle contained) by $BC$ and $G$ (is) thus also rational. And it has been applied to the rational (straight-line) $BC$.
Thus, $G$ is rational, and commensurable in length with $BC$ [Prop.~10.20]. Therefore, since the (rectangle contained) by $BC$ and $G$ is equal to the (rectangle contained) by
$BD$ and $KH$, thus, proportionally, as $CB$ is to $BD$, so $KH$ (is) to
$G$ [Prop.~6.16]. And $BC$ (is) greater than $BD$.
Thus, $KH$ (is) also greater than $G$ [Prop.~5.16,
5.14]. Let $KE$ be made equal to $G$.
$KE$ is thus commensurable in length with $BC$. And since as
$CB$ is to $BD$, so $HK$ (is) to $KE$, thus, via conversion, as
$BC$ (is) to $CD$, so $KH$ (is) to $HE$ [Prop.~5.19~corr.]. 
Let it be contrived that as
$KH$ (is) to $HE$, so $HF$ (is) to $FE$. And thus the remainder $KF$ is to
$FH$, as $KH$ (is) to $HE$---that is to say, [as] $BC$ (is) to $CD$
[Prop.~5.19]. And $BC$ and $CD$
[are] commensurable in square only. $KF$ and $FH$ are thus also commensurable in square only [Prop.~10.11].
And since as $KH$ is to $HE$, (so) $KF$ (is) to $FH$, but as $KH$
(is) to $HE$, (so) $HF$ (is) to $FE$, thus, also as $KF$ (is) to $FH$, (so)
$HF$ (is) to $FE$ [Prop.~5.11].
And hence as the first (is) to the third,
so the (square) on the first (is) to the (square) on the second [Def.~5.9].  And thus as $KF$ (is) to $FE$, so the (square) on $KF$ (is) to the (square) on $FH$. And the (square) on $KF$
is commensurable with the (square) on $FH$. For $KF$ and $FH$
are commensurable in square. Thus, $KF$ is also commensurable
in length with  $FE$ [Prop.~10.11]. 
 Hence, $KF$ [is] also commensurable in length with $KE$
  [Prop.~10.15]. And $KE$ is rational,
and commensurable in length with $BC$.  Thus, $KF$ (is) also rational,
and commensurable in length with $BC$ [Prop.~10.12].  And since as $BC$ is to
$CD$, (so)
$KF$ (is) to $FH$, alternately, as $BC$ (is) to $KF$, so $DC$
(is) to $FH$ [Prop.~5.16]. And $BC$ (is) commensurable (in length) with $KF$. Thus, $FH$ (is) also commensurable in length with $CD$ [Prop.~10.11]. And
$BC$ and $CD$ are rational (straight-lines which are) commensurable
in square only. $KF$ and $FH$ are thus also rational (straight-lines which
are) commensurable in square only [Def.~10.3, Prop.~10.13]. Thus, $KH$ is a binomial [Prop.~10.36].

Therefore,  if the square on $BC$ is greater than (the square on) $CD$
by the (square) on (some straight-line) commensurable (in length) with ($BC$), 
then the square on $KF$ will also be greater than (the square on) $FH$
by the (square) on (some straight-line) commensurable (in length) with ($KF$) [Prop.~10.14]. And if $BC$ is commensurable
in length with a (previously) laid down rational (straight-line), (so) also (is) $KF$
[Prop.~10.12].
And if $CD$ is commensurable in length with a (previously) laid
down rational (straight-line), (so) also  (is) $FH$ [Prop.~10.12]. And if neither of $BC$ or $CD$
(are commensurable), neither also (are) either of $KF$ or $FH$ [Prop.~10.13].

And if the square on $BC$ is greater than (the square on) $CD$
by the (square) on (some straight-line) incommensurable (in length) with ($BC$) then
the square on $KF$ will also be greater than (the square on) $FH$
by the (square) on (some straight-line) incommensurable (in length) with ($KF$) [Prop.~10.14]. And if $BC$ is commensurable
in length with a (previously) laid down rational (straight-line), (so) also (is) $KF$
[Prop.~10.12].
And if $CD$ is commensurable, (so) also (is) $FH$ [Prop.~10.12]. And if neither of $BC$ or $CD$
(are commensurable), neither also (are) either of $KF$ or $FH$ [Prop.~10.13].

$KH$ is thus a binomial whose terms, $KF$ and $FH$, [are] commensurable
(in length) with the terms, $BC$ and $CD$,  of the apotome, and in the same ratio. 
Moreover, $KH$ will have the same order as $BC$ [Defs.~10.5---10.10]. 
  (Which is) the very thing it was required to show.}
\end{Parallel}

%%%%%
%10.114
%%%%%
\pdfbookmark[1]{Proposition 10.114}{pdf10.114}
\begin{Parallel}{}{}
\ParallelLText{
\begin{center}
{\large \ggn{114}.}
\end{center}\vspace*{-7pt}

{\greekfont Ἐὰν θωρίον περιέθηται ὑπὸ ἀποτομῆς καὶ τῆς ἐκ δύο
ὀνομάτων, ἧς τὰ ὀνόματα σύμμετρά τέ ἐστι τοῖς τῆς
ἀποτομῆς ὀνόμασι καὶ ἐν τῷ α ὐτῷ λόγῳ, ἡ τὸ θωρίον
δυναμένη ῥητή ἐστιν.}

{\greekfont Περιεθέσθω γὰρ θωρίον τὸ ὑπὸ τῶν ΑΒ, ΓΔ ὑπὸ ἀποτομῆς
τῆς ΑΒ καὶ τῆς ἐκ δύο ὀνομάτων τῆς ΓΔ, ἧς μεῖζον ὄνομα
ἔστω τὸ ΓΕ, καὶ ἔστω τὰ ὀνόματα τῆς ἐκ δύο ὀνομάτων
τὰ ΓΕ, ΕΔ σύμμετρά τε τοῖς τῆς ἀποτομῆς ὀνόμασι τοῖς
ΑΖ, ΖΒ καὶ ἐν τῷ α ὐτῷ λόγῳ, καὶ ἔστω ἡ τὸ ὑπὸ τῶν
ΑΒ, ΓΔ δυναμένη ἡ Η; λέγω, ὅτι ῥητή ἐστιν ἡ Η.}

{\greekfont Ἐκκείσθω γὰρ ῥητὴ ἡ Θ, καὶ τῷ ἀπὸ τῆς Θ ἴσον παρὰ
τὴν ΓΔ παραβεβλήσθω πλάτος ποιοῦν τὴν ΚΛ; ἀποτομὴ ἄρα
ἐστὶν ἡ ΚΛ, ἧς τὰ ὀνόματα ἔστω τὰ ΚΜ, ΜΛ
σύμμετρα τοῖς τῆς ἐκ δύο ὀνομάτων ὀνόμασι τοῖς
ΓΕ, ΕΔ καὶ ἐν τῷ α ὐτῷ λόγῳ. ἀλλὰ καὶ αἱ ΓΕ, ΕΔ
σύμμετροί τέ εἰσι ταῖς ΑΖ, ΖΒ καὶ ἐν τῷ α ὐτῷ λόγῳ;
ἔστιν ἄρα ὡς ἡ ΑΖ πρὸς τὴν ΖΒ, οὕτως ἡ ΚΜ πρὸς ΜΛ.
ἐναλλὰχ ἄρα ἐστὶν ὡς ἡ ΑΖ πρὸς τὴν ΚΜ, οὕτως
ἡ ΒΖ πρὸς τὴν ΛΜ; καὶ λοιπὴ ἄρα ἡ ΑΒ πρὸς
λοιπὴν τὴν ΚΛ ἐστιν ὡς ἡ ΑΖ πρὸς ΚΜ. σύμμετρος
δὲ ἡ ΑΖ τῇ ΚΜ; σύμμετρος ἄρα ἐστὶ καὶ ἡ ΑΒ τῇ ΚΛ.
καί ἐστιν ὡς ἡ ΑΒ πρὸς ΚΛ, οὕτως τὸ ὑπὸ τῶν ΓΔ, ΑΒ
πρὸς τὸ ὑπὸ τῶν ΓΔ, ΚΛ; σύμμετρον ἄρα ἐστὶ καὶ τὸ ὑπὸ
τῶν ΓΔ, ΑΒ τῷ ὑπὸ τῶν ΓΔ, ΚΛ. ἴσον δὲ τὸ ὑπὸ
τῶν ΓΔ, ΚΛ
τῷ ἀπὸ τῆς Θ; σύμμετρον ἄρα
ἐστὶ τὸ ὑπὸ τῶν ΓΔ, ΑΒ τῷ ἀπὸ τῆς Θ. τῷ δὲ ὑπὸ τῶν
ΓΔ, ΑΒ ἴσον ἐστὶ τὸ ἀπὸ τῆς Η; σύμμετρον ἄρα
ἐστὶ τὸ ἀπὸ τῆς Η τῷ ἀπὸ τῆς Θ. ῥητὸν δὲ τὸ ἀπὸ
τῆς Θ; ῥητὸν ἄρα ἐστὶ καὶ τὸ ἀπὸ τῆς Η; ῥητὴ
ἄρα ἐστὶν ἡ Η. καὶ δύναται τὸ ὑπὸ τῶν ΓΔ, ΑΒ.}\\~\\~\\~\\~\\~\\~\\

\epsfysize=1.8in
\centerline{\epsffile{Book10/fig114g.eps}}

{\greekfont Ἐὰν ἄρα θωρίον περιέθηται ὑπὸ ἀποτομῆς καὶ τῆς
ἐκ δύο ὀνομάτων, ἧς τὰ ὀνόματα σύμμετρά ἐστι τοῖς
τῆς ἀποτομῆς ὀνόμασι καὶ ἐν τῷ α ὐτῷ λόγῳ,
ἡ τὸ θωρίον δυναμένη ῥητή ἐστιν.}\\

\begin{center}
{\large {\greekfont Πόρισμα}.}
\end{center}\vspace*{-7pt}

{\greekfont Καὶ γέγονεν ἡμῖν καὶ διὰ τούτου φανερόν, ὅτι δυνατόν ἐστι
ῥητὸν θωρίον ὑπὸ ἀλόγων εὐθειῶν περιέθε-σθαι.
ὅπερ ἔδει δεῖχαι.}}

\ParallelRText{
\begin{center}
{\large Proposition 114}
\end{center}

If an area is contained by an apotome, and
a binomial whose terms are commensurable with, and in the same ratio as, the 
terms of the apotome then the square-root of the area is a rational (straight-line).

For let an area, the  (rectangle contained) by $AB$ and $CD$, 
be contained by the
apotome $AB$, and the binomial $CD$, of which let the greater term be $CE$. And let the terms of the binomial, $CE$ and $ED$, be commensurable with the terms of the apotome, $AF$ and $FB$ (respectively), and in the same ratio. And
let the square-root of the (rectangle contained) by $AB$ and $CD$ be $G$.
I say that $G$ is a rational (straight-line).

For let the rational (straight-line) $H$ be laid down. And let (some rectangle),
equal to the (square) on $H$, be applied to $CD$, producing
$KL$ as breadth. Thus, $KL$ is an apotome, of which let the
terms, $KM$ and $ML$, be  commensurable with the terms of the
binomial, $CE$ and $ED$ (respectively), and in the same ratio [Prop.~10.112]. But, $CE$ and $ED$ are also
commensurable with $AF$ and $FB$ (respectively), and in the same ratio. Thus,
as $AF$ is to $FB$, so $KM$ (is) to $ML$. Thus, alternately, as 
$AF$ is to $KM$, so $BF$ (is) to $LM$ [Prop.~5.16]. Thus, the remainder $AB$ is
also to the remainder $KL$ as $AF$ (is) to $KM$ [Prop.~5.19]. And $AF$ (is) commensurable with
$KM$ [Prop.~10.12]. $AB$ is thus also commensurable with $KL$ [Prop.~10.11]. 
And as $AB$ is to $KL$, so the (rectangle contained) by $CD$ and
$AB$ (is) to the (rectangle contained) by $CD$ and $KL$ [Prop.~6.1].  Thus, the (rectangle contained)
by $CD$ and $AB$ is also commensurable with the
(rectangle contained) by $CD$ and $KL$ [Prop.~10.11]. And the (rectangle contained) by
$CD$ and $KL$ (is) equal to the (square) on $H$. Thus, the
(rectangle contained) by $CD$ and $AB$ is  commensurable with the
(square) on $H$. And the (square) on $G$ is equal to the
(rectangle contained) by $CD$ and $AB$. The (square) on $G$ is thus commensurable with the (square) on $H$. And the (square) on $H$ (is)
rational. Thus, the (square) on $G$ is also rational. $G$ is thus rational.
And it is the square-root of the (rectangle contained) by $CD$ and $AB$.

\epsfysize=1.8in
\centerline{\epsffile{Book10/fig114e.eps}}

Thus, if an area is contained by an apotome, and
a binomial whose terms are commensurable with, and in the same ratio as, the 
terms of the apotome, then the square-root of the area is a rational (straight-line).\\

\begin{center}
{\large Corollary}
\end{center}\vspace*{-7pt}

And it has also been made clear to us,  through this,  that it is
possible for  a rational
area to be contained by irrational straight-lines. (Which is) the very thing it was required to show.}
\end{Parallel}

%%%%%
%10.115
%%%%%
\pdfbookmark[1]{Proposition 10.115}{pdf10.115}
\begin{Parallel}{}{}
\ParallelLText{
\begin{center}
{\large \ggn{115}.}
\end{center}\vspace*{-7pt}

{\greekfont Ἀπὸ μέσης ἄπειροι ἄλογοι γίνονται, καὶ οὐδεμία οὐδεμιᾷ
τῶν πρότερον ἡ α ὐτή.}

{\greekfont Ἔστω μέση ἡ Α; λέγω, ὅτι ἀπὸ τῆς Α ἄπειροι
ἄλογοι γίνονται, καὶ οὐδεμία οὐδεμιᾷ τῶν πρότερον
ἡ α ὐτή.}\\~\\~\\

\epsfysize=1.in
\centerline{\epsffile{Book10/fig115g.eps}}

{\greekfont Ἐκκείσθω ῥητὴ ἡ Β, καὶ τῷ ὑπὸ τῶν Β, Α ἴσον
ἔστω τὸ ἀπὸ τῆς Γ; ἄλογος ἄρα ἐστὶν ἡ Γ;
τὸ γὰρ ὑπὸ ἀλόγου καὶ ῥητῆς ἄλογόν ἐστιν. καὶ
οὐδεμιᾷ τῶν πρότερον ἡ α ὐτή; τὸ γὰρ ἀπ οὐδεμιᾶς
τῶν πρότερον παρὰ ῥητὴν παραβαλλόμενον πλάτος ποιεῖ
μέσην. πάλιν δὴ τῷ ὑπὸ τῶν Β, Γ ἴσον ἔστω τὸ ἀπὸ
τῆς Δ; ἄλογον ἄρα ἐστὶ τὸ ἀπὸ τῆς Δ. ἄλογος ἄρα
ἐστὶν ἡ Δ; καὶ οὐδεμιᾷ τῶν πρότερον ἡ α ὐτή; τὸ γὰρ
ἀπ οὐδεμιᾶς τῶν πρότερον παρὰ ῥητὴν παραβαλλόμενον
πλάτος ποιεῖ τὴν Γ. ὁμοίως δὴ τῆς τοιαύτης τάχεως ἐπ
ἄπειρον προβαινούσης φανερόν, ὅτι ἀπὸ τῆς μέσης 
ἄπειροι ἄλογοι γίνονται, καὶ οὐδεμία οὐδεμιᾷ τῶν πρότερον
ἡ α ὐτή; ὅπερ ἔδει δεῖχαι.}}

\ParallelRText{
\begin{center}
{\large Proposition 115}
\end{center}

An infinite (series) of irrational (straight-lines) can be created from a medial
(straight-line), and none of them is the same as any of the preceding
(straight-lines).

Let $A$ be a medial (straight-line). I say that an infinite (series) of irrational (straight-lines) can be created from $A$, and that none of them is the same as any of the preceding
(straight-lines).

\epsfysize=1.in
\centerline{\epsffile{Book10/fig115e.eps}}

Let the rational (straight-line) $B$ be laid down. And let the (square) on $C$
be equal to the (rectangle contained) by $B$ and $A$. Thus, $C$ is
irrational [Def.~10.4]. For an (area contained) by
an irrational and a rational (straight-line) is irrational [Prop.~10.20]. And ($C$ is) not the same as
any of the preceding (straight-lines). For the (square) on none of
the preceding (straight-lines), applied to a rational (straight-line),
produces a medial (straight-line) as breadth. So, again, let the (square)
on $D$ be equal to the (rectangle contained) by $B$ and $C$. Thus,
the (square) on $D$ is irrational [Prop.~10.20]. 
$D$ is thus irrational [Def.~10.4]. And ($D$ is)
not the same as any of the preceding (straight-lines). For the (square) on
none of the preceding (straight-lines), applied to a rational (straight-line),
produces $C$ as breadth. So, similarly, this arrangement being
advanced to infinity, it is clear that an infinite (series) of irrational (straight-lines) can be created from a medial
(straight-line), and that none of them is the same as any of the preceding
(straight-lines). (Which is) the very thing it was required to show.}
\end{Parallel}

\newpage
\thispagestyle{empty}
~\\